# Bullinger's Apocalypse Sermons

*Bullinger, Heinrich, 1504-1575.*

---

## SERMONS OF HENRI BULLINGER, upon the Apocalypse revealed to the holy Apostle and Evangelist S. John.

## Sermon 1: ☞ Of the author of the book of Apocalypse, of the argument and parts thereof: Finally of the sundry use, and most profitable commodity of the same.

### Paragraph 1

### Paragraph 2

The prophets of God of the old Testament were God’s messengers to the people, expositors of Moses or God’s law, and even ecclesiastical preachers, who applied the doctrine taken from the law of God to the people and times in which they lived, for the edification of congregations. And they all with one accord have chiefly handled two things in their sermons. For first, they have reproved the corrupt manners of all states in their time, alleging the rule and prescription of God’s law; and exhorting all men to repentance unto God most acceptable. And to many that were uncurable they threatened all kinds of plagues, which they setting forth with all beauty of speech, showed them plainly to be seen with the eye, if they apply them so might be made afraid, and healed. Secondly, they promised and set forth by the mouth of God, the Lord Christ, the true Messiah: Whom also they described lively, and with all his holy church, teaching the faith in Christ, and showing what good things are prepared for the faithful in Christ; and also what be the true duties of piety and godliness. Neither have they concealed such things as were needful to be known of Antichrist; admonishing us most diligently that we should beware of that wolf, or rather the deepest pit of all abominations, and that we should stand fast in the sincere faith of Christ. &c.

### Paragraph 3

I have undertaken, through God’s help and your prayers, to explain the Apocalypse,[note 1.2] which is a doctrine concerning the matters of Christ’s church revealed from heaven in Christ’s glory, a summary of all godly religion, an exposition and brief declaration of the Prophets, and a prophecy of the New Testament and the history of the church.

### Paragraph 4

However, since this book is despised by both good and well-learned men, and because most people are firmly convinced that it is an unprofitable study, I will speak somewhat about this matter here.

### Paragraph 5

First, many abhor this book for this cause only that it is full of visions, types, and figures, for they suppose this does not become the Evangelical and Apostolic doctrine of the New Testament. But by the same reasoning, a portion of Daniel should be cast away. Nevertheless, this is commended to us by Christ himself in Matthew. The greatest matter of all, namely the calling of the Gentiles, was shown to Saint Peter by a vision, as appears in the Acts.

### Paragraph 6

And the prophet Joel also said how the people of the New Testament should see visions. And so does the blessed Apostle Jerome, expound the same place in the Apocalypse of the Apostles, speaking of the people of the New Testament. And our Savior Christ in the Gospel, proposed and declared to the people many mysteries by parables, and in a manner by feigned fables, as they call them. And how much do you think these visions, types, and figures of John differ from the same? This kind of speaking does not darken matters, but makes them plain; and it contributes much to the efficacy and clarity, and to the firming of the memory. For by this means, matters are not only declared with words and heard with the ears, but they are set forth to be seen by the eyes, and after a sort are fixed in the memory.

Many for this cause attribute much to painting. But I suppose that I may much more rightly attribute very much to this manner of speaking and teaching, where the matter is uttered, not by a colored, dull, and dead painting, but as it were with a speaking and lively manner set forth to be seen by the eyes. Which is therefore proposed, that men might rightly and exactly understand the same. Although therefore this whole book in a manner consists of visions and figures, yet shall we, through the inspiration of God’s grace, show in our exposition that it pertains to the perspicuity and plainness, and not to the obscuring or darkening of most high and godly matters. I will bring my exposition out of the very scriptures, by comparing and joining thereunto the rule of faith and charity. I will search out the circumstances, the things that follow and go before; I will bring similitudes and dissimilitudes; I will add also thereunto the experience of things, and the faith of histories. Which manner of expounding the scriptures all interpreters have always granted to be sound and true.

If better things shall be revealed to others, I will gladly, after the precept of the apostle, give place unto my betters. For I offer these my doings to be weighed by the godly, upon this condition, that they should try all things, and that which they shall find to be good, to hold fast.

### Paragraph 7

Secondly, they object that both new and old men of no small authority have doubted this book and its author, and have also scorned it as full of fables and unworthy to be reckoned canonical. Let those who think so give me the same liberty I desire of them, which they themselves usurp and deem lawful. For if the book of the Apocalypse should therefore seem worthy of scorn because some notable men, both old and new, have doubted of its authority, why may it not regain its authority if I show that the best doctors of the church, both old and new, have held a right good opinion of this book? And here, to the intent that I would dissemble nothing at all, I am not ignorant that Doctor Martin Luther, a man notably learned, in his first edition of the New Testament in Dutch, with a sharp preface set before, has rejected this book as it were with a dagger. However, it offended good and well-learned men that he held this judgment, which they found lacked both wit and modesty. The same man, therefore, weighing all things more uprightly and diligently, when he corrected his Dutch Bible… My worshipful master also seems not to have set very much by his book, and to have ascribed it not to John the Apostle, but to John who they called a divine. But herein there is no doubt but that he followed plainly Erasmus of Rotterdam, who in his annotations upon the New Testament, in all the Greek copies, which I have seen, states that the title was not of John the Apostle, but of John the divine. Erasmus adds that amongst the Greeks and certain old writers, men doubted of this author, which thing he declares by the testimonies of Eusebius and Jerome, whose opinion shall be spoken of straightways. But the exemplar or Spanish copy which is set forth after the faith of the most ancient and approved Greeks, exhibits to us such a title of this book: That is, the Apocalypse of the holy Apostle and Evangelist St. John Divine. For the ancient writers say how St. John the Apostle and Evangelist, for his excellent writing of the Son of God, was commonly called Divine. Whereof it follows that this title does attribute to, and not take from, St. John this book. Certainly, Arethas was also a Greek author, a bishop of Caesarea: Of the ancients, he says, certain have plucked this Apocalypse from the tongue of that well-beloved John, ascribing it to another; but it is not so. For that same Gregory, who as well as he is called a Divine, accounts this amongst those scriptures which utterly lack a suspicion of counterfeiting, saying, as the Apocalypse of St. John teaches me. And the same man a little after: But if this book was written by the mouth of the Holy Ghost, St. Basil, Cyril, Papias and Hippolytus, fathers of the church, are men to be credited. Thus saith he. What shall we say that Erasmus confesses that the consent of the world, and the authority of the church, are of such force with him that he dare not refuse this book?

### Paragraph 8

Let us hear now the judgment of that most excellent man, Dr. John Oecolampadius, the most faithful pillar of the church of Basel, and excellently learned in prophetical matters and in all the canonical scriptures, concerning this book, which he left us written in the twelfth chapter of the second book of his commentaries upon Daniel. But St. Jerome, the paraphraser or expounder of the Prophets, says he—[?how much he attributes to this our author]—whom I marvel why certain with so rash a judgment do reject as a dreamer, and frantic, and an unprofitable writer to the church. Nevertheless, he propounds and sets forth very many of the most secret and hidden things of the old testament and of the Prophets. But those great men do betray what a sense of their own importance they have. Whose judgments I would verily rather scorn as profane than I would cast away such a treasure. I could here bring forth goodly testimonies of other new writers, but that I make haste to the judgments of the ancient fathers.

### Paragraph 9

The earliest of all writers after the Apostles, Justin and Irenaeus, the noble Martyrs of Christ, attribute this book to John the Apostle. For Eusebius, in chapter eighteen of the fourth book of the Ecclesiastical History, states that Justin mentions the Apocalypse of John, plainly saying that it is the Apostle’s. Jerome also writes in the life of blessed Justin that Justin explained the Apocalypse of Saint John, but that explanation does not survive to any great extent as far as I know. The same author writes that Irenaeus presented the Apocalypse of Saint John with a commentary, which is also unavailable. He himself, said to have lived around the year of our Lord 160, testifies plainly in the fifth book against the Valentinians that this revelation was shown to John the Apostle a little before his time. We cite certain words of his in chapter thirteen of this book. Tertullian, who lived around the year of our Lord two hundred and twenty, in the fourth book against Marcion, although he says that Marcion rejects the Apocalypse of John, nevertheless states that the order of bishops reckoned up to the very beginning affirms Saint John to be its author. In serious matters and when reasoning against heretics, he readily uses the testimonies of this book.

### Paragraph 10

The same things are also recorded concerning the blessed Martyr Saint Cyprian, under the name of John the Apostle, in his Epistles, treatises, and Sermons. Eusebius also, in chapter 18 of the fifth book of the Ecclesiastical History, shows that Apollonius, a very ancient writer, uses the testimonies of the Apocalypse of Saint John. And likewise Theophilus, Bishop of Antioch, who affirms this in chapter 23 of the fourth book of the Ecclesiastical History. Also Origen, a great man in the church of God, in chapter 25 of the sixth book of the same Eusebius. He wrote, says he, the Apocalypse, which rested upon the Lord’s breast, and so forth.

### Paragraph 11

I have thus far recounted the opinions of the most ancient Martyrs and Doctors of the Christian church, concerning the Apocalypse, meaning Justin, Irenaeus, Tertullian, Cyprian, Apollonius, Theophilus, and Origen. I will shortly bring yet more judgments, both of Greek and Latin writers, of most authority in the church, agreeing with the minds of those that we have alleged hitherto. However, I will first briefly touch upon such things as Dionysius of Alexandria left written about the same book in the twenty-fifth chapter of the seventh book of Eusebius, whom I suppose they have followed, as many after him have spoken against this book. He says that diverse of his predecessors did utterly reject and reject this book. Neither does he hide the cause why they so did, for that the kingdom of Christ is affirmed therein to be earthly. Whereunto doubtless they referred that precious city, and the rest which under terrestrial kinds, figured spiritual things. Which when we, in the treating thereof, have dissolved, declaring that it does not edify the earthly kingdom of Christ, but a spiritual and celestial one, no man, I think, will reject a good and godly book, because certain abusing the testimonies thereof, give unto it a wrong sense.

### Paragraph 12

Have heretics twisted many passages of scripture to defend their errors, should the authority of scripture itself then be called into question? I favor neither the Chiliasts nor the Millenaries in regard to this book; I give them no support.

### Paragraph 13

Eusebius speaks very well at the end of the third book, when speaking of Papias, the first author of the Millenaries. He thought, says he, that after the resurrection Christ should reign here corporally with his thousand years on earth. Which I suppose he thought because he did not understand well the Apostles' words, nor did he consider well those things that were spoken of him under figures, for that he was endowed with a small judgment [? understanding].

### Paragraph 14

But in the meantime, Dionysius himself, I say, would not reject this book. He adds immediately that he thinks it not yet to be the book of John the Apostle, but of someone else, but who that someone should be, he did not know. He gathers also by certain conjectures, by the phrase of speech, and handling of the book, and by the unlikeness of wit, that this book should be another’s than he who wrote the Gospel and Epistle. But seeing that the arguments of the story and Epistle are so diverse, that neither are like, and the argument of the book of Revelation most diverse of all: why should it seem marvelous, though it agrees not with them in all things?

### Paragraph 15

No one can deny that there is great agreement in doctrine. The Epistle to the Hebrews appeared to many to favor the Novatians or Catharites in chapter six and ten. The difference in style was noted to differ from the rest of Saint Paul’s Epistles. But if we were to judge the holy scriptures in this way, I do not know what would be firm and sure enough. Leaving this disputation in suspense, I will now proceed to bring forth the judgments of other old writers concerning this book.

### Paragraph 16

Eusebius, nicknamed Pamphilus, Bishop of Caesarea, living during the reign of Emperor Constantine, and a devoted student of ancient writers, whom many suspect sought to undermine the authority of this book to their own advantage, eloquently affirms in chapter eighty of the third book of his history that John was exiled to Patmos and wrote his Apocalypse there, while denouncing the tyranny of Domitian.

### Paragraph 17

And other historians do likewise. He again, in the twenty-third chapter of the third book, concerning the Apocalypse, says that the opinions of people are diverse, some approving and others reproving it. Again, when he should present his opinion concerning the canon of the New Testament in the twenty-fifth chapter, he joins the Apocalypse with the books of unquestioned authority, although he does not dissemble that he will show elsewhere what other people think of it. While he performs this, he incorporates many more and better, who judged the Apocalypse to be of Saint John the Apostle, and embraced it as a most godly book, than those who denied or reproved it.

### Paragraph 18

Epiphanius, Bishop of Salamis in Cyprus, a Greek author, also clearly ascribes this book to John the Apostle. Read what he has written against heretics in the fifteenth heresy. And Jerome attributes much to this Epiphanius.

Jerome himself ascribes this book to John the Apostle, to Paulinus: “The Apocalypse of John,” he says, “has so many sacraments as it has words.” Moreover, Philastrius, Bishop of Griria, whom Augustine says he saw with Ambrose at Milan, considers them heretics who reject the Apocalypse of John, and say that it is not of John the Apostle, but of Cerinthus, an heretic. Ambrose himself alleges in his books testimonies from the Apocalypse, under the name of John the Apostle.

### Paragraph 19

Augustine embraced this book as Apostolic and referred the same to his church, leaving certain treatises upon the same. Primasius, also Bishop of Utica in Africa, explained it as Apostolic. Of Bed and the rest of that sort, I speak nothing, since his opinion is known to all men. Andreas, also Bishop of Caesarea, wrote upon this book, as Arethas reports in his commentaries, an opinion I declared before.

### Paragraph 20

I believe I have sufficiently confirmed the authority of this book, against those who diminish it. But the strongest evidence seems to be the book itself and the way it is handled, proving that it proceeded from the Apostle. This we shall demonstrate in the treatise itself. But if those blessed fathers in their time did interpret the Apocalypse for their churches, why should it not be lawful for us also in our time to interpret it for our people, who live at the end of the world, where all things are now more fully realized than they were in their time? Indeed, these things serve most chiefly for us and for our time, especially as we travel and are exercised under Antichrist.

### Paragraph 21

Therefore, it is in vain for many to claim that this book is obscure and cannot be understood, and for that very reason should not be read in the church without any profit or benefit.

### Paragraph 22

For to speak nothing of this, nothing is found in holy scripture which does not have an excellent fruit. Neither must we immediately despair of true understanding, even though at first sight the holy scripture seems obscure, which is opened by God himself, and is obtained through prayer and godly exercises. Certainly, we are not ignorant that many would prefer that nothing be spoken of Antichrist, so that he might reign more carelessly, and they themselves be less subject to peril. But Christ commands us to trouble him. Let us therefore proceed in the work of the Lord.

### Paragraph 23

And where it offends them, that John makes little mention or none of Christ, whereas the manner of the Apostles is always to intimate Christ and the grace of redemption: we suppose this same book will be looked upon more thoroughly to prove the contrary. Whose argument now I will recite.

### Paragraph 24

The prophet Zechariah in chapter three [note 1.32] presents the entire mystery of Christ to all people’s eyes in a most evident figure. He sees Jesus, the high priest, clothed in shameful garments, and like a coal taken from the fire, suffering much opposition from the Devil. Immediately afterward, he sees the same person cast off the shameful clothing and put on white garments, to be glorified, and proclaimed king, priest, and savior of all.

### Paragraph 25

This is what the Apostle and Evangelist John explains. First, after presenting the Gospel, he describes Christ in humble clothing, detailing the great opposition he endured from the wicked until he was nailed to the cross. He also touches upon his glory there. However, the Apocalypse appended to that Gospel expands upon this, showing him in a white garment and in glory, truly demonstrating how, after this humbling, he was exalted and obtained a name above all names. And now, being in glory, he nevertheless works in the church, the savior of all the faithful within the church. In his Epistle, he praises this entire mystery of piety and urges it upon all people.

### Paragraph 26

The entire book is divided into six sections.

### Paragraph 27

For first is set the title with the beginning and some of the work, and with a brief narration: and all this in the first part of the first chapter.

### Paragraph 28

Secondly, from the midst of the first chapter to the fourth chapter, Christ’s reigning in glory is described, at the right hand of the Father, and it is declared how he is active in the Church through his Spirit and the ministry of his word. What things he teaches from heaven, and what the sincere doctrine of the church is, are revealed. Also, the restoration of fallen churches and their preservation are shown.

### Paragraph 29

3 From chapter four through twelve, Christ continues to earnestly warn his Church through seven seals and seven trumpets, describing the events that will befall the Church, all of which are most justly ordained by God himself through the Lamb, Christ.

### Paragraph 30

4 Moreover, from chapter twelve to fifteen, the conflict of the Church with the old Serpent and with the old and new beast is more fully described. Here, too, that monstrous Tyranny, both old and new, and very Antichrist himself is clearly portrayed in his colors; nevertheless, these things are afterward also again more plainly declared.

### Paragraph 31

5 And from the fifteenth Chapter to the twenty-second Chapter, are recited the pains and torments of Antichrist and Antichristians, and the destruction of the same and the condemnation of all the wicked. Also the Judge, Christ, is set forth, and the process of an external judgment is figured. There is also rehearsed the triumph, joy, and reward of Saints. Where also heaven itself is opened to be seen by our eyes, that now we may by faith look into the same. The depth of Hell is opened, that we may look into it also: and may take good heed that we be not thrown thither headlong.

### Paragraph 32

Finally, at the end of chapter twenty-two, follows the conclusion and commendation of the work, with its sealing.

### Paragraph 33

And here I will not conceal another division of this work, not to be disregarded [note 1.34]. I observe that commentators have followed this structure. First, they present the title and beginning. Then, they append the entire work itself, divided into seven visions. Indeed, the number seven is most frequent and, as it were, distinctive to this book. Finally, they add a conclusion that summarizes the work, contained within the last chapter. And these visions are contained within their boundaries.

### Paragraph 34

In the first three chapters, the initial vision is explained, exhibiting Christ to us reigning in glory, governing, ordering, correcting, and preserving his church.

### Paragraph 35

The second vision begins in the fourth chapter and extends to the eighth chapter. This section portrays God himself and his Christ, presenting them for observation, and it commends their most just governance of all things in the world, opening seven seals.

### Paragraph 36

The third vision features seven angels sounding with seven trumpets. This section extends to chapter twelve.

### Paragraph 37

The fourth vision depicts the struggle of the woman with the serpent, and presents to us the old seven-headed beast and the new two-horned beast, which are descriptions of Antichrist, as found in chapters twelve, thirteen, and fourteen.

### Paragraph 38

In the fifth vision, seven angels pour out seven vials of God’s wrath, continuing to the seventeenth chapter.

### Paragraph 39

From there begins the sixth vision, and it extends to chapter twenty-one, discussing God’s just judgment against Babylon, the great prostitute, and against those who oppose Christ, ultimately against all unrepentant and wicked people.

### Paragraph 40

The seventh and final vision presents to the eyes of all the faithful the glory and eternal blessedness of the saints. Indeed, this section of the work possesses great beauty and affinity with the rest of the things which in this book are all, in a manner, treated by the seventh number. Let the reader follow whichever he prefers.

### Paragraph 41

Now concerning these things, every person may perceive that this book is entirely apostolic and exceedingly profitable to us all, especially those whom the end of the world has overtaken. And this book will be easier to understand, for now all things are in a manner accomplished.

### Paragraph 42

Daniel was thought to have told of strange visions when before the monarchies, he prophesied the monarchies. But after those things were accomplished which he prophesied, he seemed unto many to have compiled a fictitious story. The self same, I am sure, you will judge also: this same book of Saint John. A few prophecies only we shall recite.

### Paragraph 43

First we have in this book a most full description of Christ reigning in glory, our king, I say, and Bishop; and how he governs the Church, and is the Savior of all the faithful. We have also a most gallant description of Christ’s Church, and how the same may be built, repaired, and maintained. Then we have a perfect description of Antichrist, of his members, and synagogue of his counsels, crafty devises, kingdom, cruelty, and destructions of the same; from which it bids us beware. Moreover, we have an abridgment of history from Christ’s time unto the world’s end.

### Paragraph 44

Finally, an absolute and certain prophecy of things to come, so that we need not the prophecies of Methodius, Cyril, Merlin, Bridget, Nolhard, and certain triflers.

### Paragraph 45

Furthermore, we have great consolation and comfort for the church in adversity, while we see the Lamb open the Seals, and that all things are accomplished by God’s providence, and that there is an end to evil. And the church shall endure forever, in spite of all the devils in hell. Lastly, we have a most plentiful and sure doctrine of the Judge and final judgment, of punishments and of rewards.

### Paragraph 46

All these things I say, the treatise itself will plainly demonstrate, for our edification through Jesus Christ our Lord.

## Sermon 2: OF THE TITLE OF THE whole work, and exposition thereof.

### Paragraph 47

. ☞ I said the whole book was contained in six parts.Three members of the first part. Now must we look on the first part: Which hath chiefly three members: The title, beginning, and brief narration. For this present we will only speak of the Title, which is thus.

### Paragraph 48

The revelation of Jesus Christ, [note 2.1] which God gave to him, to show his servants the things that must shortly come to pass. He sent and showed by his angel to his servant John, who bore witness to the word of God, and to the testimony of Jesus Christ, and to all things that he saw. Happy is he who reads, and those who hear the words of the prophecy, and keep those things written therein. For the time is at hand.

### Paragraph 49

This title is abundant and proclaims all the beneficial matters that ought to be stated at the start of books.

### Paragraph 50

First, the title or inscription of the entire work is set: the Apocalypse, or revelation of Jesus Christ, which truly was opened or revealed by Jesus Christ himself.

### Paragraph 51

This title immediately demonstrates that this work is not a human invention, but a divine doctrine: as that which was revealed by our Lord, king, and priest Jesus Christ, out of heaven, from the right hand of the Father, exercising there the office of the high priest, and still teaching us profitable things. And although it is also called the Revelation of John, it is ascribed to him for no other reason than that, as a scribe, he wrote and presented it.

### Paragraph 52

Again, it is declared even more plainly from whence this Revelation comes: even from God himself. For he says, “which God,” namely the Father, gave unto him, that is, to Christ. For in the holy and blessed Trinity, there is a distinction of persons. And although all things of the Father are also the Son’s, and all things of the Son are likewise the Father’s, yet Scripture mentions the Father giving to the Son, and the Son receiving from the Father.

### Paragraph 53

All the ancient writers have consistently explained this as occurring through the mystery of dispensation. For the Son received something from the Father, just as the Son of God himself says, "Glorify me with the glory that I had with you before the world existed."

### Paragraph 54

Moreover, the Son is the wisdom, word, and mouth of the Father, by whom God in times past and now spoke and speaks to the Fathers, Prophets, Apostles, and to the universal church. The Father, by dispensation, gave to his Son this office, that he should be Bishop. For no man has seen God at any time; the only-begotten Son, who is in the bosom of the Father, has revealed him to us. Let us know, therefore, that this same is a Divine Revelation, which God the Father, loving mankind, has revealed by the only Bishop Christ unto his Church. And so it joins together the Father and the Son, that nevertheless the holy distinction of persons remains safe.

### Paragraph 55

Now also is added, for what purpose God the Father has revealed, or granted the gift of revealing, namely, the office of priesthood to his Son, our Lord Jesus Christ: so that, being revealed, he might demonstrate it, and as it were set it before the eyes of his servants, that is, his worshippers and Christians, who are called the servants of God for their willing obedience. And just as a servant owes to his lord all that he has or is worth, so we owe to God ourselves wholly, and all that belongs to us, or else we are free, and not bound.

### Paragraph 56

It is also declared to whom this revelation is opened [note 2.7]. To all the servants of God. If, therefore, you desire to be called a servant of God, hear this book and remember it. And know that this book is prepared for you by God.

### Paragraph 57

After he summarizes in a few words what Christ revealed to John, the events that must shortly come to pass. Therefore, the destinies of the Church are recounted, what good and evil things will happen to the faithful, and likewise what punishments must be inflicted upon the wicked.

### Paragraph 58

And let no one gather from this teaching a sense of obligation, necessity, as though God did not work freely; or that the wicked do evil, not through their own fault, but by God’s compulsion. Good things must be done, because God, willingly binding himself to us by his promise, cannot fail to do what he does and promises; nevertheless, he works freely.

### Paragraph 59

Good things must be done in the godly community, for the nature of grace and faith is such, just as the characteristic of ungodliness is to despise and transgress. Therefore, they must also be punished. And because the world is as it is, there must needs be countless heresies and calamities. He says these things must shortly be accomplished, for certain things began in the very time of Saint John. And although many things are found to be done a thousand years after, Saint Peter, the apostle, says, “A thousand years before the Lord are as it were yesterday.” Therefore, this Apocalypse pertains to the times of the primitive and last church, and declares whatever things shall happen to it until the last judgment. Indeed, how it shall reign forever.

### Paragraph 60

Moreover, the manner of revealing is also addressed. For Christ revealed those things, sending by his angel, or his angel sent forth, to whom he gave commandment what he should say and do. Whereupon this angel is afterward also called Christ, because he represented the person of Christ. Therefore, we must not consider the angel in this book, but always Christ, the true author of all these things. Indeed, the divinity of Christ is commended to us when we hear that Christ is the Lord of angels. Paul reasoned more at length about this to the Hebrews. Moses in the twelfth chapter of Numbers sets forth chiefly three manners of prophecy or revelation. First, by vision, of which sort many are ascribed to Daniel, one notable to Peter in the tenth chapter of Acts, and likewise to Paul. And into this form falls also the Apocalypse. Secondly, by dream: of which sort were those of Pharaoh and Nebuchadnezzar, kings whom Joseph and Daniel interpreted. The Prophet Joel in the second chapter mentions visions and dreams. For in the New Testament also there are very many holy and prophetical dreams. Lastly, Moses sets forth a skillful exposition, as many were made to Moses and to the Apostles. Into whose fellowship the Apocalypse comes after a sort also, where visions are openly explained. Here appears the unspeakable goodness of God, which in many ways procures and works our salvation, and pleasantly prepared offers it unto us to enjoy. Unhappy is he who does not know these things.

### Paragraph 61

Besides this, much mention is made of him to whom Christ has opened this divine and most excellent revelation, given to John. He commends himself (for it was expedient for confuting his adversaries, seeing that Paul also many times maintained his authority against the false Apostles) by four epithets. First, he calls himself the servant of Christ. This is the oldest and noblest title, which the fathers, prophets, and apostles have used, for they dedicate and consecrate themselves to God. Secondly, John testified to the word of God among the apostles, most expressly declaring the divinity of Christ, especially where he testified and said, “In the beginning was the word.” Moreover, he was a witness of the witness of Jesus Christ: under which name the Lord himself in the Gospel, and John in the twelfth chapter of his Gospel, comprised the whole evangelical doctrine. And he was a seeing witness of all these things. For in the first chapter we have seen his glory, he says, and in the nineteenth chapter he saw water and blood gush out of the Lord’s side. In his epistle, he says, “We have seen and have heard.”

### Paragraph 62

Arethas notes that in some Greek manuscripts is added what is also found in the Greek manuscript of Spain. And whatever he has heard, and whatever may be, and whatever must be done afterward.

### Paragraph 63

That same John is the author of this book,[note 2.18] who, just as he saw the Lord in the flesh upon earth, so he saw him in spirit, revealing these things in heaven, and who presents to the church sights that are most certain and sure. This John was that beloved disciple of the Lord, who at the last supper rested upon his breast, to whom in his last will he bequeathed his mother on the cross, one virgin to another.

### Paragraph 64

He alone stood by at the altar of the cross when Christ died, a witness to the true death and to our purification. He lived until the time of Emperor Trajan, as Eusebius recounts in his chronicles, citing Irenaeus, noting the year as one hundred and three from the birth of Christ. Dorotheus, an ancient writer, affirms that John lived sixty years.

### Paragraph 65

It is also important to consider the benefit of this godly work or revelation, so that readers and hearers might be spurred to diligence. For this book is also called a prophecy. This is because it foretells future events, making it the prophecy of the New Testament. Moreover, a prophecy is an exposition that clarifies and explains the Old Prophets. It promises blessedness to those who read, hear, and keep the things written in this book. Blessedness encompasses the benefits of the present life, to the extent that the Lord allows them to the faithful, but chiefly pertains to the life to come. If the benefit of this book has already been discussed in the first sermon. Note that it is not enough simply to read or hear this book; it must be put into practice and kept diligently.

### Paragraph 66

For the Lord also said in the Gospel: Blessed are those who hear the word of God and keep it. Therefore, those who shape their lives according to this book are blessed. For they avoid the deception of Antichrist, remain in the faith of Christ, and live forevermore. &c.

### Paragraph 67

And he concludes the title with an exclamation, by which he stirs up the listeners greatly: “The time is at hand!” As though he should say: Let no one think here that strange things, which concern him not, are told here, which shall come to pass at length after many worlds; they belong to every one of us. For they are written of matters that chiefly concern us, and even of our own affairs. Thus he shows that this book is profitable for all men, ages, and worlds. God the Father, by his Son, teaches profitable things, and admonishing in due season, be praised world without end. Amen.

## Sermon 3: OF THE BEGINNING OF THIS book, and the Apostles salutation: wherein are declared the mysteries chiefly of Christ, ſecondly of our whole faith and redemption.

### Paragraph 68

.

### Paragraph 69

John to the seven congregations in Asia. Grace be with you and peace from him who is, and who was, and who is to come, and from the seven spirits which are sent forth before his throne. And from Jesus Christ, the faithful witness, the firstborn of the dead, and Lord over the kings of the earth. To him who loved us and washed us from our sins in his own blood, and made us kings and priests to God, to him be glory and dominion forevermore. Amen. Behold, he comes with clouds, and every eye shall see him, and they also who pierced him. And all kindreds of the earth shall wail over him. Even so. Amen. I am Alpha and Omega, the beginning and the ending, says the Lord Almighty, who is, and who was, and who is to come.

### Paragraph 70

Another section of the first part of this book contains a beginning or preface, wherein is the Apostle’s salutation, whereby he first describes the whole mystery of Christ, and secondly of our faith and redemption. For so the Apostles were accustomed in the beginning of their writings to include a brief summary of salutation. This practice is evident everywhere in Paul’s Epistles. By this description, he gains the goodwill and attention of all people.

### Paragraph 71

The opening greeting of the Apostles is nothing else but a blessing. Blessing is an old, customary practice, by which the patriarchs wished for their children all manner of good things, both of body and soul. This is described at length in Genesis. Also, the high priest had authority given to bless the people, as we read in the sixth of Numbers, especially commanding that God’s name be placed upon the people. Therefore, it is a superstition to say that God, from whom every good gift descends from above, blesses, that is, gives good things, but that ministers or men merely wish. And the Lord in the law promises that he will grant those things to the people which the high priests should wish for them. Therefore, neither words nor shaven crowns, but the truth and power of God give the gifts. We ought not therefore to doubt that God will grant to us also the apostolic blessing, so that, being reconciled and accepted by God, we might have peace. And first, Saint John repeats his name, lest we should doubt of the author, whom we see Christ to have used as scribe and interpreter to all congregations. But he does not repeat that he is that servant of God, or witness, or Apostle of Jesus Christ. That was sufficient to have heard at the first beginning. Therefore, he teaches modesty and humility, which have received great gifts. Afterward, he indicates to whom he writes, and to whom this book pertains, to the seven churches of Asia, the names of which he will mention shortly after. Arethas, bishop of Caesarea, by the seven churches, and by the seven numbers, says he signified the multitude of churches that are in all places. So also Primasius, bishop of Utica in Africa, explains the seven numbers. Therefore, this greeting, this book, and the whole doctrine of Jesus Christ, written by Saint John, pertains to the whole universal church of Christ throughout all the world, and in all times and ages. Whereupon it belongs to all of us as many as are of us in the church of Christ. For although the letters are addressed to the Romans and Galatians, it does not follow that they are not ours. And he writes expressly to the churches of Asia, not to the churches of Jerusalem or Judaea, that he might show that the kingdom of Christ has already come to the Gentiles. And as God from the beginning chose Israel, in whom he might set forth a perfect example of the church and commonwealth, so from the beginning of the New Testament, he chose those seven churches of Asia, which he might set forth to the whole Christian world. But if Rome had been placed in the first place among the churches, as Ephesus is, good God, how much would the Roman faction make of it, for the establishing of their supremacy?

### Paragraph 72

And the manner of the Apostles’ salutation wishes grace. Grace is the favor of God, and the reconciliation meant, whereby God the Father for Christ’s sake is made one with us, our sins pardoned, and we adopted for his children. Thereof arises the peace and tranquility of mind and the desire of concord with all men.

### Paragraph 73

And here he shows abundantly who grants the church his blessing, that is to say, grace, reconciliation, and peace—God, and God three in persons, the Father, the Son, and the Holy Ghost, one God in essence. But here he discerns the persons very well. From him that is, to wit, the Father; and from the seven spirits, that is, from the Holy Ghost; and from Jesus Christ—this is the diversity of persons. And the signification of the unity is, when after the properties of persons are declared, he adds, “I am Alpha and Omega.” &c. And that the Holy Ghost is set here in the midst, it disorders not the mystery of the Trinity, but appears to be an argument that he is the spirit as well of the Father as of the Son, and that he proceeds from both. As it is also proved by the words of our Lord, the 14th, 15th, and 16th of John. Here is also described the whole wholesome mystery first of Christ, then of the catholic faith, and of our redemption, so that herein you may find the chiefest articles of the Apostles' Creed, and have here a most goodly description of Christ our Lord. Hereof all men shall judge how truly some say that this book, contrary to the common opinion of the Apostles, makes little mention of Christ and of faith.

### Paragraph 74

The father, as source and origin, from whom the Son is generated, is first described: for it is he who is, who was, and who is to come. John took those words from Moses in the third and thirty-fourth chapters of Exodus, and from many testimonies of Isaiah. And he says nothing but that God the Father is an eternal essence, which exists by and of itself, and is and preserves life in all. And that this essence is such that it has been always without beginning. For this is what he joins to being, or existence, “was.” He adds, and he who shall come, which shall be, and shall remain even to the end, and to everlastingness without end. The Greeks derive [illegible] from running, for that concerning and running, he meddles with all matters; is everywhere present, bringing help to the godly, or restraining and punishing the wicked.

### Paragraph 75

And the Holy Spirit, where he is but one,[note 3.8] for the sevenfold—that is, all manner of grace and manifold gifts—is here called, as I may say, Septenary or of the seventh number. And from the seven spirits, John says that it is from that Spirit, which is endowed with the sevenfold grace. Those diverse gifts are, in a way, explained in the eleventh chapter [?of Revelation] and elsewhere in the scriptures. He is said to be in the sight of the throne, that is, before the throne of God, joined truly in governance with the Father and the Son. For the throne is often used to represent the kingdom. Therefore, the Holy Spirit is of the same glory, power, and majesty as God.

### Paragraph 76

Now he has come to Christ, whom by his attributes he describes most fully. You know that Jesus is the proper name of Christ, which Matthew explains as “Saviour.” Christ is the surname of his office and dignity, as one might say “anointed,” that is, bishop and king.

### Paragraph 77

First he calls Christ our Lord a faithful witness, as in chapters 49 and 50 of Isaiah. For he was sent by the Father into the world from heaven, an Apostle, to testify God’s will, what he would have done with humankind—namely, that he would save the world by his Son and through faith in him, who is obedient to God’s law. For he must do the will of his Father. This Christ is a faithful witness, meaning sure, constant, and true; of whose doctrine no man ought to doubt. No one has seen God at any time; the only-begotten Son, who is in the Father’s embrace, has revealed him. This one, therefore, is the overseer and universal guide of the church. Whoever dissents from him are to be shunned. “Hear him,” says the Father.

### Paragraph 78

2 He is the firstborn of the dead. For he died indeed for our sins; he rose again from the dead, and was made the firstborn of the dead, Lord and conqueror of death. In him we believe that we also shall rise again, and in what manner. As Paul writes in 1 Corinthians 15. And just as in the first part he shadows the humanity of Christ, wherein he taught also his deity, in that he was the faithful, true, and catholic bishop, and is yet at this day; so in the second, the articles of our belief concerning the death of Christ and his resurrection are confirmed. To these may also be added the article of the resurrection of the dead.

### Paragraph 79

[note 3.14]3 Christ is prince over the kings of the earth, a monarch truly, and Lord of all rulers: He has taken a name above all names, the Lord of angels and of all creatures, to whom all things are subject, as the Apostle expounds in Colossians and Philippians 2. And he does not abolish laws and magistrates which will be king of kings and Lord of lords. For if there were no kings, how could Christ be king of kings? The most sacred Emperors, Constantine, Theodosius, and Justinian, knew themselves to be clients of Christ: their kingdom was Christ's, and they were to be subjects. Christ acknowledges these as his, by whom he governs those he has redeemed with his blood. Those who proudly rule over the people, boast themselves to be lords of all things, and do not acknowledge Christ to be monarch over all, are stark mad. And here are comprehended such things as we confess in the articles of our faith, that Christ ascended into heaven and sits on the right hand of the Father: that is, that he has received the highest power over all things in heaven and earth, as in Ephesians 1 and Acts 2.

### Paragraph 80

4 Christ has loved us with incomparable love. For he himself says: “Greater love has no one than that someone should lay down his life for his friends.” The Apostle amplifies this in the fifth chapter to the Romans. And it was exceeding great love that moved Christ to come down from heaven and be incarnate and to redeem us by his death. With a free love he loved us, provoked by no merit of ours. For as John in his canonical Epistle speaks the same of the Father, “In this is love, not that we have loved God, but that he has loved us and sent his son as an atoning sacrifice for our sins,” so it is to be understood of the Son, that he has and does bear us great good will, not moved thereto through our love, with which we have embraced him. And of the free love to mankind, he gave himself unto death and washed us from our sins. For it is straightway added, by his blood.

Where three things seem to be observed by us. First, that Christ washes, purges, purifies, or cleanses the faithful: and that most fully, not partly. He alluded to the washings of the law, which he also expounded. For David says: “Purge me with hyssop, and I shall be made clean; wash me, and I shall be whiter than snow.” The same phrase of speech repeats Augustine in the first chapter. Micah also says: “The Lord will return and will have mercy on us; he will trample underfoot our iniquities. And you shall throw into the depths of the sea all their sins.” And the Lord says: “I will pour clean waters upon you, and you shall be cleansed from all your uncleanness.” The Lord Christ accomplishing these things, washes us, purges, and cleanses thoroughly, as well from the fault as the pain. He cleanses us from our sins, not from one, but from all. This which thing is proved both by former testimonies, and again in the first and second Epistle of John. Lastly, the manner of purifying is set forth, by blood. For without the shedding of blood no remission was made. Therefore through the mediation of death and bloodshedding, full remission of all sins was obtained for the faithful. They who bring forth any other manner of forgiveness of sins, are injurious to the death and blood of the Son of God. And here we may plainly see set forth an article of the Apostolic creed: I believe the forgiveness of sins.

### Paragraph 81

In the fifth place, the effect of our redemption and purification is shown. For Christ has accomplished that as many of us as believe in the Father through the Son of God should be kings and priests to God and to his Father. Arethas and the copy of the complete text read not “kings” but “kingdom,” which is not read amiss. For we are the kingdom of God, because God by his Spirit, not the flesh nor the world, ought to reign in us. And when we submit to the Spirit’s governance, we are the kingdom of God. Which thing Paul discusses at length in chapter six to the Romans. Moreover, we are made kings—that is, free—by Christ, so that we should not serve the devil, the flesh, and the world, according to that saying of Zechariah, being delivered from the hands of our enemies, that we might serve him without fear in holiness and righteousness before him all the days of our life. And Christ has consecrated priests with his Spirit and blood, that we should offer up to God spiritual sacrifices, ourselves pure, prayers and praises, and almsdeeds. For that these are spiritual offerings, Peter and Paul testify. And John took these things from Exodus: For we of the Gentiles who have believed have succeeded in the place of the people of Israel, rejected Christ through unbelief. And these things give light to that article of the Creed, “I believe in the holy catholic church, the communion of saints.” For we, as many of us as believe, are the fellowship of God’s people, sanctified through Christ to the service of God. Of whom are these things said?

### Paragraph 82

In the sixth place, in the description of Christ, he shows that glory and rule are due to God alone through Christ, for the church forevermore. We give glory to God when we ascribe our salvation and all goodness to his goodness, not to our own strength and merits. We give him rule when we acknowledge him to be Lord and head in the church, working by himself, not by the saints in heaven, to whom he has granted power. Not by the Pope, whom he has not constituted Vicar on earth. The whole glory and rule is Christ’s.

### Paragraph 83

Seventhly, in the description follows the coming of Christ for judgment,[note 3.26] and the manner of his coming. Just as a cloud took him from the sight of the Apostles, even so shall he come in clouds to judge the living and the dead. The scripture testifies. And he adds that the eyes of all men shall see the judge,[note 3.27] even of those who have displeased him. From this we gather two things: first, that the judgment shall be universal, wherein men rising shall see Christ with their own eyes. Another thing is that Christ shall come to judge in the same flesh,[note 3.28] in which he was wounded and pierced, hung upon the cross, was buried, and rose again. This passage is taken from Zechariah, and is also cited in Saint John’s Gospel.[note 3.29] And it behooves that his body be shown to the whole world full of prints and marks, so that from this may be judged the godly and also the ungodly: those who then believed in such a redeemer, these who then rejected and scorned such a one. Of these we understand that is added: And they shall wail, for indeed they have neglected their own salvation, which the wise man discourses at large.[note 3.30] Moreover, lest any should doubt about those things that are spoken of the judgment, and of the lamentation of the wicked (as Saint Peter said,[note 3.31] the contemners and mockers of the judgment would be), he adds a kind of confirmation, even so. Amen.

### Paragraph 84

In them also is explained the article of the creed of Christ who will judge the living and the dead. He concludes this passage with these words: “I am Alpha and Omega,” with the phrase “that which follows (the beginning and end)” omitted in some copies. It seems that interpretation of “I am Alpha and Omega” crept in from the margin. It is a saying of Saint John the Apostle, “I am Alpha and Omega.” Heretics, such as Basilides and Valentine, were greatly delighted in letters. But against those lettered heretics, John speaks plainly by the mouth of Christ: “I am Alpha and Omega.” If anything ought to be ascribed to letters, I am all of this—the whole everlasting virtue, essence, and eternity. For the meaning is that God is the beginning and end, that is, eternal, unspeakable, best, and greatest. These things are repeated: “He that is, which was…” and so on, which were explained before. There is added, “almighty.” For hereby is declared the unity and majesty of God, of whom the Trinity was also explained before. Hereby also the authority of this book is confirmed, the author of which is shown to be that eternal and almighty God. To whom be glory.

## Sermon 4: ¶ Of the Narration of this book, where also is discoursed of the place and time, and of the author of this Revelation.

### Paragraph 85

.

### Paragraph 86

I, John, your brother and companion in tribulation, and in the kingdom and patience which are in Jesus Christ, was on the island of Patmos for the word of God and for the testimony of Jesus Christ. I was in the spirit on the Lord’s Day and heard behind me a great voice, as if of a trumpet, saying, “I am Alpha and Omega, the first and the last.” And I saw another vision, and it was commanded that I write in a book and send it to the congregations which are in Asia, to Ephesus, and to Smyrna, and to Pergamum, and to Thyatira, and to Sardis, and to Philadelphia, and to Laodicea.

### Paragraph 87

Note 4.1: The final section of the first part shows us a brief account wherein the Apostle John declares the time and place of this Revelation, and by whose authority he sent it to the seven churches in Asia.

### Paragraph 88

And again now for the third time the name of John is repeated. He saw undoubtedly that there would be some who, to take away the use and benefit of this book, would doubt the author. Against whom he repeats and reiterates his name so often, lest we should doubt and lose the great value of so worthy a book.

### Paragraph 89

He adds certain things to his name, which instruct concerning the state of the Apostle, and certain profitable matters. First, he calls himself a brother, namely of those churches, and of all ours. As where I have admonished you that in the seventh number are comprised all churches of all times throughout the whole world. We are all, so many believe, the children of one heavenly father. And therefore all spiritual brethren in Christ, coheirs with Christ, and heirs of God. Which thing Paul taught after Christ.

And seeing our dignity is so great, let us be ashamed of our misdeeds, lest our memory be put out of this most noble and celestial family. It is a shame that a brother of Christ, of John, and all the Apostles should degenerate. And so forth.

But why have they not so instantly urged this brotherhood, as the Monks have promoted their forged fraternities, the Rosaries of the Virgin Mary and of Saints? Because that was free, and cost nothing. But the Monks sell theirs dearly. They are therefore deceivers and seducers.

### Paragraph 90

After he calls himself a partaker in affliction, or oppression and persecution, as he who was even now banished by the Emperor Domitian and lived in exile. And he joins together and does not separate himself in the evil common to all the faithful brethren. And truly it is one and the same persecution that vexed the Apostles and torments us at this day. Let us therefore rejoice that we have the Apostles and all the Martyrs of Christ as fellows in our trouble and affliction, that we are broken and bruised with the heavy burden of evils. Let us therefore be patient and long-suffering. For it is not enough to be afflicted and vexed with all kinds of evils (for many without any fruit or praise at all endure most grievous pains). But it becomes us also to be patient in adversity. Therefore John at this present joins with all in patience. For the Lord said in the Gospel, “In your patience you shall possess your souls.”

### Paragraph 91

After he adds unto tribulation and patience a kingdom, and that an heavenly not a terrestrial kingdom.[note 4.5] And he brings in the kingdom for the comfort of the patient people. For also the Apostle Paul said, a certain and sure saying. For if we die with Christ, we shall live also with him: If we suffer, we shall reign with him. &c. Let us always here with comfort our selves in adversity. For we are thrust down, that we might once be exalted again, 2 Corinthians 4. And all these things are concluded in Christ Jesus, by whom we are both the children and brethren of God, and suffer many things patiently, and are made partakers of his kingdom. For even for these things must we thank him and his mercies, and not our own desert.

### Paragraph 92

Let us note here also what and how great the humility of the greatest and worthy Apostle of God was, considering his state: It was not pleasant, but difficult, yet in patience joyful through Christ. But where are those now who glory in the name of Apostles? Who, meanwhile, swelling with pride, are addicted to filthy pleasures? I warn you, that we flee from them as from apostates.

### Paragraph 93

And now he shows the place where this divine revelation was made to him, where also he was commanded by God to write the same. The place was the Isle of Patmos. It is accounted among the islands called Sporades by Pliny in the fourth book and twelfth chapter. It lay opposite Asia and the city of Ephesus, and was in the sight both of Europe and Africa, so that it seemed to be as it were a middle seat, or holy chair, out of which Christ preached through John from heaven to the whole world. And indeed the councils of God are wonderful, and his goodness is unspeakable, which reveals so great mysteries, as it were in the Roman prison or Babylonian captivity, to his faithful.

### Paragraph 94

Neither does he hide the cause of his coming into the island, saying, “I was there for the word of God and the testimony of Jesus Christ.” The word of God is the very essence of God, called by John with a unique expression the word or sermon of God, as appears in the first chapter of John. The testimony of Jesus Christ is the Gospel itself, which Jesus testified, and which his disciples have testified concerning Jesus. Therefore, for the confession and preaching of Jesus Christ and his wholesome Gospel (for he also explained how he became a partaker of affliction), John was apprehended in Asia and led by soldiers to Rome, that he might plead his cause before Emperor Domitian, who, by his cruel nature, condemned the innocent. And he was put into a cauldron of hot boiling oil. When he escaped unharmed from this, he was carried to Patmos. He answered no other matter before the emperor than Paul did twenty-seven years before, during Nero’s reign. This occurred in the thirteenth or fourteenth year of Domitian’s reign and twenty-four years after the destruction of the city of Jerusalem and sixty-six years after the birth of our Lord. Domitian, who sought to be called a god, was slain by his own men after many murders and cruel acts, and died himself a shameful death in the fifteenth year of his reign. The authorities for this are Suetonius in the life of Domitian, Tertullian in *Prescription Against Heretics*, Eusebius in his chronicles and in the third book of the ecclesiastical history, chapters seventeen and eighteen. And to this is added the common consent of all writers.

### Paragraph 95

Moreover, he notes the time also,[note 4.9] in which these mysteries began to be revealed to him, [illegible], in that solemn day of the Lord, namely Sunday. For so the ancient fathers have called one of the Sabbaths, that is to say, the first day of the week, wherein Christ rose again from the dead, as recorded in Matthew 28 and Mark 16. And this day the churches have chosen for themselves in place of the Sabbath day, as holy in remembrance of the Lord’s resurrection, wherein they might keep their sacred and solemn assemblies. For that this day was observed and consecrated for assemblies in the congregation of Corinth appears manifestly in the 16th chapter of the first Epistle to the Corinthians, where the Apostle commands them to lay apart their collections on one of the Sabbaths. The same day also the faithful did celebrate their service with Paul, as recorded in Acts 20. Where Sozomenus reports in the 8th chapter of the first book of the *ecclesiastical history*, that great Constantine made certain holy days, and even the Lord’s day, which the Heathens call Sunday; it is to be understood that he renewed rather the custom of the Apostles and the catholic church, than that he newly instituted it. And freely the churches have received that day; for we read not that it was anywhere commanded. And the congregations saw how it was altogether necessary that there should be a certain time in which the saints should meet and come together. They chose therefore the day of the resurrection, neither did they maliciously contend among themselves for these things, as the histories testify was done in the church afterward. And at this day, verily, the superstitious holy days being abrogated, it is better to observe certain and moderate days, and to keep peace and quietness in the church.

### Paragraph 96

But where this Apostle knew that the faithful on that day served God in all assemblies, where he could not be present in body, in spirit and contemplation he was with them. And as he was thus in the spirit and contemplation of divine matters, and in holy prayers, he heard a voice, whereof I will speak hereafter. But here we are presently taught what is the religion of the Sunday, and how it is fitting to serve it. Finally, worldly men are reproved who pollute and break it with profane works and affairs. David, in his time, when he suffered persecution from Saul, lamented chiefly that he might not come to the Lord’s tabernacle. Our people, however, count it a great felicity never to enter into the fellowship of the Saints. And to abuse the Sunday, in gaming, drinking, dancing, and worldly business.

### Paragraph 97

These things declared in this manner, he comes now to the Revelation, setting forth before the clear commandment of God, by which he was commanded both to write the things revealed and also to send them to the seven churches of Asia. A chief characteristic of the Revelation’s manner and majesty is that he heard a voice, and one notable as the sound of a trumpet. For so we read it was done in the giving of the law at Mount Sinai. Now it is declared who that voice was, and who the author of the Revelation. Truly, it was the eternal God, who calls himself Alpha and Omega, that is, the beginning and the end. Or as it is said in Essay, first and last, as we shall see elsewhere.

### Paragraph 98

Now follows the commandment which has several parts. For first, the Lord commands John to see and to write such things as he saw, that is to wit, the Apocalypse. And he should write neither on a tablet nor on the wall, but in a book: Verily for the edifying and profit of the church present, and of all posterity. After this, he also commanded to send those writings to seven congregations, and verily to all the churches of the whole world for all times and ages. Therefore, all these things belong to the provision of congregations, and that of all that be, have been, or shall be.

### Paragraph 99

Here we learn how great is the authority of the scriptures. It was not written nor compiled in books, but by God’s commandment. There are notable testimonies of the book of Moses, in the thirty-fourth chapter of Exodus and the thirty-first chapter of Deuteronomy. And to say nothing of the rest of the prophets, was not Jeremiah commanded to write his sermons again, which King Jehoiakim had cut in pieces and burnt?

Doubts Saint Peter bears manifest witness that the prophets received the mysteries of God to no other end than that they should reveal them to us; which indeed might only be done by the scriptures. Now John is most clearly commanded to write. What will we say that he is also commanded to send his writings to the congregations? Whereof again we gather that God wills right well to the congregations, and even to every one of us. Let us beware and take heed that we do not put from us unworthily so great benefits of God, to whom be praise and glory.

## Sermon 5: ¶ THE BEGINNING OF THE work is made, and a most goodly description to us exhibited of Christ king and bishop in glory, and nevertheless woorkyng in the Church.

### Paragraph 100

.

### Paragraph 101

And I turned myself to look at the voice that spoke with me. And when I turned, I saw seven golden lampstands. And in the midst of the seven lampstands, I saw one like the Son of Man, clothed in a linen garment that reached to the ground, and girded about the chest with a golden belt. His head and his hair were white, like white wool and snow. And his eyes were like a flame of fire, and his feet were like burnished bronze, as if they had been heated in a furnace. And his voice was like the sound of many waters. And he had in his right hand seven stars, and out of his mouth went out a sharp two-edged sword, and his face shone as the sun in its strength.

### Paragraph 102

Such things as have been treated of hitherto in this book serve in place of a prologue or preface, as they term it. Now at last the matter itself will be presented to us. [Note 5.1] Therefore, what follows is the second part of this book, which extends to the fourth chapter. In this, Christ is described to us with his universal church. For first, we are shown the most sacred image of Christ our Lord, teaching what he is on the right hand of his Father in glory, and how he, sitting on the right hand of his Father, continually works in his church, never absent, present always. Of what sort the church is here on earth, is figured in those seven congregations. Here, therefore, are shown the excellent gifts of churches, and again, shameful errors: How the Lord Christ confirms those who are sliding and ready to fall, establishes those that stand, strengthens the weakhearted, restrains the foolish and bold, and preserves things that are corrupt. Finally, how faithful ministers of the church must work and travel with the people committed to their care. For here is exceedingly taught what is the repairing and preservation of churches. Where also a brief summary of the whole ecclesiastical doctrine, brought in to an abridgement, shall be set before us. For here is repeated from heaven concerning Christ in glory, the doctrine of true religion, which he had set forth more plentifully when he was yet here on earth; and here it is most aptly applied to churches, after consideration of the same.

### Paragraph 103

And in most goodly order the words are knit together (as likewise the whole book is written with plain words and hanging right well together; they are deceived who think it to be loose, unbound writings). John heard a voice behind him crying: whereupon he turned to see the voice speaking, that is to wit, him who spake. For Arethas also admonishes that there is a trope. For no man sees, but hears the voice. After turning to see, he saw a figure of Christ our Savior. Therefore, when the Lord speaks, let us turn also with all our heart,[note 5.2] that we may likewise deserve to see the mysteries of the kingdom of God, for he gladly reveals himself to those who turn and desire heavenly things. And from those who neglect the mysteries of the kingdom of God, all things of salvation are hidden.

### Paragraph 104

Furthermore, John depicts to us the image of Christ, our universal king and chief bishop, seated in glory:[note 5.3] and within this description are contained the principal matters concerning Christ. For such a glimpse of Christ is here given us as our frail flesh may perceive in this world. But we shall see him ultimately in the world to come as he is, in the fullness of his majesty, wherein shall be joy and everlasting life; but this in this corrupt world is yet granted to no man.

### Paragraph 105

So much, therefore, is permitted to us who live yet in this world and are still visible, as is profitable, and as our weakness may comprehend. But this is not small or insignificant; it is great, abundant, and full of spiritual joy, if we behold these mysteries of God with a faithful eye and a mind eager for godly matters. And there is no doubt that these are certain and true things that are revealed to us here. For they are revealed by the very Son of God. Let us not wish to see more, or desire greater things than these are, but take pleasure in those which Christ has granted us. And let us know for certain that a wonderful benefit of God is given to us in this vision. For who would not desire to see Christ in glory, sitting on the right hand of the Father? Who does not desire to know what our Savior does in heaven? And how, being in heaven, he is nevertheless present with his church on earth? But this sacred and holy image instructs in all these points most fully the faithful of Christ. However, this image of Christ is not to be depicted with colors, since colors cannot attain to its majesty, but with the ecclesiastical doctrine, which has the promise of the Spirit of Christ, and is therefore more evident and fitting for its true expression. Let us also print this image, not upon any dead surface with colors that will perish and fade, but in our hearts through the living Spirit of God, which may also keep it in our minds, never to be erased. And the things spoken in the second and third chapters of this book are derived from this description of Christ, that the majesty of the thing might invite us to singular diligence. The matter is very plain.

### Paragraph 106

First we are taught who it is, whose image is to be exhibited to us: not the Son of Man himself in his own substance, but like the Son of Man. The Son of Man, according to the Gospel, is called Christ himself, very God. Here he did not show himself to be seen as John in his own substance, but in the form of an angel, that represents Christ. This thing is found more than once in the book. We shall therefore refer all these things unto Christ, not to the angel, which is the minister of Christ in the mystery. And we shall see Christ in his own substance, what time our base body shall depart from hence, and being raised from the dead shall be glorified. In the meantime, the soul from the death of the body till it rises again, shall clearly have the fruition of the sight of Christ: wherein, as I said before, shall be the chief joy and felicity. We shall now therefore see Christ as it were in a glass, and so much as shall suffice us. The Lord open to us the eyes of our mind.

### Paragraph 107

He also tells where he saw Christ, in the midst of seven candlesticks. Soon we will see that the candlesticks must be understood as representing the churches. Christ is then in the midst of the church. He sits truly on the right hand of the Father, and according to the nature of his divine body, he is only in one place, and no more. As Saint Augustine declares abundantly in the fifty-seventh Epistle to Donatus. Yet, because he is also very God, he is likewise in the midst of the church, as he promised in the Gospel: “Wherever two or three are gathered in my name, I am in the midst of them.” And again: “Behold, I am with you until the end of the world.” Therefore, by his power, Christ remains and works in the present church, and is not absent. (Leave therefore to inquire, what Christ does on the right hand of his Father, whether he sits continually?) And he is truly in the midst of the churches, fixed to a place, but shows himself equally to all, as helpful. For he accepts no persons, nor sleeps. He is not painted, he is not idle, paying no regard to matters of the church, but is chiefly and only attentive to the salvation of the same. Such a one he promised himself to be in chapters 14, 15, and 16 of John. And saying that Christ is in the midst of the church, what Vicar shall he have? Shall he have that enemy which is directly against him? For a Vicar is in the place of one who is absent: but Christ is in the middle of the church, present, not absent.

### Paragraph 108

Christ is described most plentifully here, and many things are ascribed to him. It is declared in what manner Christ is in the midst of the church. First, it is shown what garment he has on: namely, both priestly and princely. By this, it is figured what manner of one Christ is in heaven and on earth: namely, bishop and king, intercessor, mediator, and sacrifice, a most perfect sanctification and justification, a redeemer and deliverer of the faithful to his father, evermore working the salvation of his faithful. As Paul teaches, Romans 8, and Hebrews 7. A *poderes* [?] is found among the apparel of Aaron, and it is a priestly garment. Jerome writes to Fabiola concerning the priestly garment. The second vesture of linen is a coat down to the foot, of double linen. Josephus calls it *bessina*. It is called in Hebrew *ketheneth*, in Greek [illegible]. This clings closely to the body, and is so narrow and straight-cut that there is no wrinkle at all in the garment, and came down to the legs. This was truly white and clean. For the Lord Christ is an undefiled priest, as in Hebrews 7. Neither does he wear again a foul vesture, as he did in Zechariah 3, nor a purple, as in John 19. But a bright one, as he who has obtained a name above all names. But his girdle or belt is worn by soldiers and triumphant persons, and it signifies in Christ royal dignity. For Christ is king, deliverer, and redeemer of the faithful. His victory is ours. He has overcome Satan, Hell, sin, and death.

### Paragraph 109

But the belt or girdle of Christ is not set in the customary place, namely, about the loins. For as Arethas has also admonished, there are no desires to be restrained in Christ. Therefore, he is not girded after the manner of sinners, but about the breasts: so that we should understand by the girding that he is king of kings, devoid of all affections; most righteous and holy in judgments and government. But yet, in the meantime, he is furnished for the defense of his church, as we have read it written in the 9th Psalm. The Lord has put on strength and girded himself, and so forth. Christ might seem to have girded himself not after the manner that priests or kings use, for that he has obtained a more excellent priesthood and kingdom, enduring forever. To accomplish these things, it behooved him to use a temple and palace, not transitory, but heaven itself. Hebrews 6 and 9. Yet, in the meantime, the effect penetrates into the church itself, so that he may be present in the church also.

### Paragraph 110

But the head of Christ appears white, and his hair white, like the finest wool and the whitest snow. Such a head is also ascribed to the Father of our Lord Jesus Christ, in the seventh chapter of Daniel. For they are of the same essence. And hereby is signified wisdom and age, and also the divinity and deity of Christ. And because Christ is God, therefore he is the head of the church, ministering to the body spiritual wisdom, and all celestial gifts. Ephesians 5. The Pope of Rome, that most wicked man of sin, what a head is he then? Without life, without brains, utterly foolish. As he is described in the eleventh chapter of Zechariah. It is a shame that we will not see these things, being ever blind. Christ is everlasting, omnipotent, and knows all things. And he may be the health and head of the body. “In the beginning was the Word, and the Word was with God,” says he. Christ himself: “Before Abraham was, I am,” says he. Therefore, the heretics lie, denying Christ to be very God, of the same substance with the Father. He is the dominion of God, all things are subject to him. Ephesians 1. And he himself fulfills all things, just as he is present with his church.

### Paragraph 111

Now his eyes are not darkened nor blind, but keen and bright. For Christ knows all things. Christ’s eyes are watchful; nothing is hidden from him, he sees all things that are done, both good and evil. And he sees in order that he may judge and require. He is light in darkness, and the sight of Christ is to good men joyful in perils. Finally, the judgments of Christ are righteous. The prophet David says, “The eyes of the Lord are upon the just, and his ears are to their prayer.” Again, “The face of the Lord is against them that do evil.” And just as the head is not plucked from the body, so Christ cannot be absent from his church. And considering that his eyes are quicksighted, and that the Lord foresees all our things, and has the charge over us, how is he absent from his church? What need is there of any vicar?

### Paragraph 112

And the festivals of the Lord are like copper, or like brass and frankincense burning in a furnace. For Chalcolibaum is a word compounded of brass and frankincense. This Erasmus notes, and that Swidas shows also the same, that there is a kind of copper more precious than gold: which he says is made of saltpeter and of a stone. Pliny, in the thirty-fourth and second chapter, calls it a kind of brass, which was dug out of the veins of the earth, and in times past had great price. It seems to me to be the same which in the first and tenth of Ezekiel is called Hasmal, a present remedy against poisons. For if wine intoxicated is put into a cup thereof, it will hiss. And so the death and poison are detected.

### Paragraph 113

The most blatant, fiery spectacle signifies the fellowship and the ways of the Lord, blameless, his judgments right and just. And that he so walks in the church and governs all things, that in the meantime all uncleanness be detected and consumed, but he himself remains always most holy and pure. For fire purges. God is a consuming fire.

### Paragraph 114

But the voice of Christ [note 5.15] is as it were the noise of many waters, not so much because all nations and people commend and praise him, but because the Gospel and word of God have come into the whole world. This voice also, mighty kings could less assuage and appease than they could stop the gushing of waters or contain the winds in sacks [note 5.16]. Therefore, by the power of preaching, the Lord is always present in his church.

### Paragraph 115

The hand is an instrument of all instruments, especially the right hand. [note 5.17] In this, Christ holds seven stars, namely, seven prelates or pastors of churches in Asia. And even all the bishops throughout the whole world, Christ by his power gives to us as pastors, and instructs, comforts, confirms, and defends them, so that they should preach his word. By this, he joins himself to the church; Christ works through them in the church and preserves them.

### Paragraph 116

The same is more vividly expressed in the words that follow. For a sharp, two-edged sword comes out of the Lord’s mouth. This is the word of God, as Jerome declared very well in the sixth chapter to the Ephesians and fourth chapter to the Hebrews. And this word or sword does not hang upon a wall, nor sticks fast in the sheath, nor hangs by the fringe, but comes out of the mouth. He does not say it came forth or it shall come forth, but it comes forth, as the thing that is in continual operation, or perpetual preaching throughout the world. And it is two-edged, sharp and piercing, as in the heart of the godly unto salvation, so in the heart of the wicked to pain and condemnation. And yet at each day comes forth that sword from the mouth of Christ by the mouths of ministers. The word of Christ is indeed condemned by the world, and is called by many a fable, but it is a sword, and that a sword out of Christ’s mouth. All the unfaithful do find and shall find this, however they resist. With this sword Christ kills the wicked. The effect of this sword is greater than was the sword of Alexander, Pompey, Julius Caesar, or Marius, Attila, or [illegible]. Neither does it matter, though the world now acknowledge it not. It shall accomplish their greatest evil in time to come. Doubtless with this spirit of his, the Lord always continues to comfort and govern the church, so that he is never absent from the same.

### Paragraph 117

Finally, the countenance of Christ shines as the sun does in its greatest strength around noon, when it is bright, clear, and pleasant. By the countenance we recognize men chiefly. Therefore, by the countenance we know Christ. The countenance of Christ is light. Christ, therefore, is light. And that truly a divine and eternal light, enlightening all that they may also be made the children of light, and then the faces of the saints may shine in that day as brightly as the sun and as the face of Christ shone. Matthew 13 and 17. And furthermore, he thus communicates this light unto us (1 John and 1 John). How is it to be thought that Christ should be absent from his church? Thou see how he is present.

### Paragraph 118

And so our Lord Christ has revealed himself to us, making himself visible for salvation, and has opened himself wholly, so that we may understand what he does for us and how he is in his church. In these things are all the mysteries of the Gospel comprehended. For what can you say of Christ that you have not here encompassed? Let us therefore remember them and write them in our minds, that we may embrace Christ, king and bishop, and that we never let him depart from our arms. To him be glory.

## Sermon 6: ¶ How John was affected towards the vision to him exhibited, the comfort of John, and the exposition of the vision, applied unto consolation.

### Paragraph 119

.

### Paragraph 120

And when I saw him, I fell at his feet as if dead. And he laid his right hand upon me, saying to me, “Fear not. I am the first and the last, and I am alive, and was dead. And behold, I am alive forevermore, and have the keys of hell and of death. Write therefore the things which you have seen, and the things which are, and the things which shall be fulfilled hereafter. And the mystery of the seven stars which you saw in my right hand, and the seven golden candlesticks. The seven stars are the messengers of the seven congregations. And the seven candlesticks which you saw are the seven congregations.

### Paragraph 121

It follows how blessed John was moved by that heavenly and wondrous vision, and how he received comfort, along with an explanation of the vision applied to his comfort, with a command to write down all these things diligently.

### Paragraph 122

What time he had fully seen this divine and heavenly vision of Christ our Lord, sitting on the right hand of God in glory, his strength failing him, he falls down on the earth: and lying at the feet of the Lord, is like a dead body. We read that the same happened to Daniel in the tenth chapter. And other men of God also have been frightened by the visions of Angels. The women also in the New Testament trembled at the sepulchre, seeing Angels. Peter was amazed at the greatness of the miracle. Luke 5. And falling at the knees of the Lord, cries out, “Go from me, Lord, for I am a sinful man.” For godly visions betray our infirmity. Neither are we prepared or sufficiently purged to behold those celestial things with eyes and minds sick and not yet well purified. Therefore, the elect must be glorified in another life, that they may be made partakers of celestial glory. In the mean season here, all the godly are humbled and abased by high visions and revelations. For they do not advance themselves proudly through the glory of revelation, but perceiving their natural corruption, they crave pardon and the augmentation of the celestial grace and light. For unless we be illumined with the spirit of God, we shall lie like dead people, however we receive with our corporal senses the visions of heaven.

### Paragraph 123

But those who humble themselves before the Lord find a most present consolation at the Lord’s hand. Therefore, there came to John immediately both in word and in a full consolation. For the angel representing the person of Christ lays his right hand upon John, which is a token of friendship, protection, and of present help. For when pressing this manner of speaking in Dutch, we say that laying on the hand signifies that Christ is good to John, ready to help him. Which he immediately makes plain by the addition of words, saying, “Fear not.” This saying is common everywhere in the story of the Gospel, and therefore is most gospel-like, that is to say, most fortunate. For God commands the humbled to be of good hope and to live assured under the protection of the highest. Which truly we understand to be spoken not to John alone but to all of us also, that we in like manner, albeit that we feel the infirmity of our flesh, should yet hope well of the goodness and mercy of God. Here follows that cause more fully declared, why John should not be afraid. For the vision showed was not exhibited for the terror of him, but that John might perceive how great and mighty he is who is prepared for the defense of him and all the faithful. As though he should say, “Where you see how great he is, who has taken upon himself to defend you, who finally protects and governs the whole Church, there is no cause why you should be afraid. But rather execute boldly what he commands you. Write what he commands to be written. Be not afraid of men, fear God rather. For if good men are so sore afraid at the sight of him, where shall the enemies and contemners of God appear?”

### Paragraph 124

Therefore, consequently, he explains the vision, teaching who he is, who was seen like the Son of Man walking among the golden candlesticks. And he applies this explanation to comfort, so that both John and every faithful person may perceive how mighty Christ is, and what the faithful have obtained through him. For the angel shapes his speech so that we may seem to hear all things spoken not by the mouth of the angel, but by Christ himself. [note 6.6] And this explanation has its parts. First, he declares (as I said just now) whose image it was that was shown. Then, a commandment to write this book is added. After that, the mystery of the stars is opened. Finally, the secret of the candlesticks is revealed: and all these things, right plainly and briefly.

### Paragraph 125

First you have seen, says the Lord, a vision whereat you were dismayed: but fear not. For you have not seen any evil or fearful spirit, foretelling any misfortune: but my form, which am your redeemer and Lord. I am first and last. And this manner of speaking he took out of the Prophecies of Isaiah, as it is to be seen in the 41st, 43rd, 45th, and 48th chapters. And he signifies himself to be coequal and of the same substance with the Father in all things, very God, eternal, and incomprehensible. For just as the Father attributes things to himself, the Son also claims them. But there is no order or certain time to be understood in “first and last,” but plainly everlastingness. Therefore Christ here signifies that he is very God, equal and of the same essence with the Father from all eternity, as is also much confirmed in John 1:5, 10, 14, and 17. This contends against the heretics, who at that time also, as today the Servetarians, deny the eternal deity of Christ the Lord. And thus, when the true God is acknowledged and believed by us, he may be for our salvation. If Christ is not very God, he is not our salvation. For I am God, says the truth: And besides me there is no God, no salvation.

### Paragraph 126

Secondly, he says, “I am alive, and was dead,” by which he signifies that he took on the true human nature. Which many also at the same time denied, just as there are some at this day who plainly diminish the humanity of Christ. Against all such heresies, the Lord himself confesses that he was dead. Whereby it is now manifest that he is very man, as he is also very God, of the same essence with his Father in deity, as he is also of the same substance with us in humanity, like unto us in all things, sin excepted. For he took not the nature of angels, but the seed of Abraham. And it behooved indeed that the Son of Man should be incarnate, that both he might suffer and shed blood. For the testament in the dead is finally ratified; neither is there any remission made without shedding of blood. The Lord therefore dies and sheds blood, to the intent he might give full remission of sins and confirm the new testament. Yet even he that was thought to be dead, now lives, and is that same living one, who having vanquished death, on the third day rose again from the dead, and restores life to all believers, and inspires into them his own very life.

### Paragraph 127

And therefore adds immediately: behold I am living world without end. For now Christ dies no more, death shall not rule over him. But rather he is the life of all his faithful, who in rising again, brought again life: and that life everlasting, enduring, I say, world without end. As he himself declares more at large in John, chapter 5, verses 6 and 10. And the Apostle to the Romans, 4:1; Corinthians, 15; and 2 Timothy, 1.

### Paragraph 128

Furthermore, where many were inclined to doubt this life obtained and restored by Christ, the Lord himself confirms what he said through another, and says, “Amen.” As though he should say, “This is altogether true and certain that I say.”

### Paragraph 129

Finally, he adds,[note 6.11] “And I have the keys of Hell and of death.” By these words, he greatly comforts and expresses his power, declaring how great he is and what we have from him. Here, we must speak by way of analogy. The common explanation says well: he who has the keys to any house lets in whom he will and keeps back whom he will from entering. Therefore, Christ possesses the keys of death and Hell, for he delivers whom he will from perpetual condemnation of death, and whom he will,[note 6.12] he allows to remain justly in the same danger of damnation. Indeed, Isaiah in chapter 22, speaking of Eliakim, whom he says should be made steward in the court of Hezekiah, says, “I will lay the key of the house of David upon his shoulder, which shall open, and no man shall shut; shall shut, and no man shall open.” Therefore, keys in Scripture represent the charge and governance of a house. Eliakim would govern all things in the court of Hezekiah uprightly. Whatever he determined, no one would infringe; what he abrogated, no one would restore. Christ, of whom Eliakim was a figure, himself also has the chief governance in the house or kingdom of God, so that he may quicken whom he wills and pluck back from Hell and damnation. Again, he may condemn whom he wills through his just judgment. For he has full power over death and Hell.[note 6.13] For he has overcome and weakened both. These things greatly comfort the faithful and retain them in all godly duties. It is chiefly to be observed that he does not say he *had* the keys or *will have* them, but “I have,” he says, “I have.” He did not give his power to the Bishop of Rome, but retains it himself and will keep it forever. Nor did he give to the Apostles full power of life and death, of salvation and damnation,[note 6.14], thus divesting himself; but he gave the keys of opening and shutting heaven, as it were, to his ministers and servants, through the preaching of the Gospel. By this, he promised life to all who believed, and Christ himself should give that life for the truth of the promise. To whomever they should threaten damnation, Christ himself should condemn for the truth of his word. Therefore, we see that the Lord still keeps and exercises the power, and his ministers the ministry (by preaching, not by absolute power). Therefore, the Pope is Antichrist,[note 6.15] who usurps and takes upon himself this full power and authority in Heaven and on Earth, and in the midst of the earth, or beyond all the earth, namely those unfortunate islands. I mean purgatory. By this crafty device, he has subtly emptied the purses, coffers, garners, and wine cellars,[note 6.16] of foolish people who stray from the articles of their belief, namely, that I believe in the forgiveness of sins, the resurrection of the flesh, and everlasting life. The beast dares usurp the two horns of the Lamb,[note 6.17], namely, the authority of King and Bishop, and therefore hangs two keys[illegible] under his triple crown, so that even by these emblems all the world may perceive that this is he who, having subdued three kings or horns, has arisen and claims all power in heaven and on earth, signified by the two keys. And surely the blindness of our time is wonderful and to be lamented, that having eyes, it sees nothing. Let those who are wise remember that Christ still has the keys of death[illegible] and Hell, his ministers the denouncing of life and death.

### Paragraph 130

And now when he had declared these great and wholesome matters, and had comforted the mind of John, he adds the commandment: write the vision exhibited, finally write those things also which must be done shortly after this. He places in the midst those that be—that is, which are in deed, and true, and be not false. And these things are to give authority to this book, finally to the whole scripture, which is revealed with like truth of the same Author. And as John is commanded to write without fear, so are we commanded to preach and publish the same boldly, though the world be never so mad thereat.

### Paragraph 131

He adds further, the exposition that remains, and says: “The mystery of the seven stars, &c.” The reason seems almost unpersisted [?]. Therefore, we must understand this to be the mystery or sacrament of the seven stars and candlesticks, so that it may be as it were a proposition, and the exposition should follow immediately. The seven stars are seven messengers, &c. And by sacrament understand a secret mystery, and the very exposition of the mystery. As if to note here the goodness of Almighty God, which declares to us even the hardest places of Scripture. Where, then, are those who accuse the Scripture of obscurity, and contend that it cannot be understood? Let us also mark here the common manner of speaking of the whole scripture: seven stars are seven messengers. The seven lights are seven churches. For signs receive the names of the things they are, although they remain in their own substance, and are not changed into another. This is also granted by those contentious persons who, in the words of the supper, “This is my body,” will acknowledge no figurative speech at all.

### Paragraph 132

Stars are called angels. Angels are God’s messengers, pastors of churches, so called in the second and third chapters of Malachi. For God sends preachers as ambassadors to the people, and wills them to be heard in the same way as himself. Luke 10; John 13. Let no man therefore wait until the Lord himself comes down from heaven again and preach to us. Even now he preaches to us by his messengers, who preach his word, that is to say, the word of Christ. If you despise them, you despise Christ. Preachers are called stars because of their bright and heavenly doctrine, and for their purity of life. [note 6.22] Beware therefore you preachers, that you be not wandering plants, lest you have no light at all, neither in doctrine nor in conversation of life. For then you shall be likened to stars that fall out of heaven, as will happen hereafter in this book to false teachers.

### Paragraph 133

But those stars are not in the head, or in the feast, or on the back or sides: but in the right hand of Christ. This thing has indeed a great consolation for the pastors, being in God’s right hand, in God’s protection, so that no man shall take them out of his hand. God himself also upholds pastors, and furnishes them with necessary goods of the Church. Therefore is the whole government and glory his. Wherefore the Apostle says also: he that waters and plants is nothing, but God that gives increase.

### Paragraph 134

Now as concerning the candlesticks, there was certainly in the Tabernacle of Moses with seven sockets, set in seven candles. In Solomon’s temple there were ten candlesticks. The one represented a figure of Christ: and even upon it, and the ten, signified the universality of churches, which are lit by the only light Christ, and have of this one, whatever light they have. And the candlesticks are of gold. The mystery whereof Arethas expounds: They are all gold, says he, for the purity and preciousness of faith lying hid in them. Indeed, the candlesticks of themselves give no light, but are receptacles of light. So from us arises no light, but darkness. But in that everlasting light, set a light in the candlestick, the light shines: if Christ illumine the church with faith and understanding, then faith shows itself in open confession and the purity of life in conversation. And this the Lord requires of his church in the fifth chapter of Matthew. So let your light shine, and so forth. And the apostle in the second letter to the Philippians: Shine like lights in the world amidst a crooked and perverse generation.

### Paragraph 135

And hitherto we have handled the consolation of Christ and the exposition of that great and celestial vision, whereby we have learned the mysteries of the faith of Christ and of his Church: to the end we should know that Christ is the Lord reigning in his Church, and applying all things to the salvation of his faithful. That he sends preachers, teaches by them, and keeps and defends them. To him be glory, and so forth.

## Sermon 7: ¶ Of the Epistles revealed out of the throne of God from Christ by an Angel, and received and sent of John. Where also a part of the Epistle to the Epheſians is expounded.

### Paragraph 136

.

### Paragraph 137

To the messenger of the congregation of Ephesus write: These things says he who holds the seven stars in his right hand and walks among the seven golden lampstands. I know your works, your labor, and your patience, and how you cannot bear those who are evil. You examine those who claim to be apostles but are not, and you have found them liars. You have endured, showing patience. For my name’s sake you have labored and have not fainted. Nevertheless, I have something against you, because you have left your first love.

### Paragraph 138

Your kindness has observed a certain image of the Lord Christ, [note 7.1] sitting on the right hand of the Father in glory; yet so, that in no wise does he forsake or neglect his Church. Now follows a more full and plain explanation of how our Savior Christ in Heaven executes the office of the high Bishop, and teaches the whole church by his ministers, rebukes, comforts, and retains it in its duty; finally, turns away all things hurtful, and advances it to greater things. For here follow seven Epistles, to the seven Congregations: that is to say, [note 7.2] unto all the churches in the [illegible] whole world. For this most ample and wholesome doctrine [illegible] must not be restrained to a few, since Christ is Bishop universal.[note 7.3]

### Paragraph 139

But great is the authority of these epistles. For they are revealed from the throne of God, by the Son of God speaking by an angel, which prescribes what is to be written in those epistles. John receives and writes the same, through Christ’s commandment, and sends them to the seven congregations. And verily they pertain no less unto us than if now the bearer entering into the church should deliver these letters unto us.

### Paragraph 140

Moreover, in these seven Churches is figured unto us the nature, manners, vices, medicines, rebukes, praises of all Churches in all times, and whatsoever is wont to happen about them. Then by examples of most excellent means mixed, of hypocritical also and wicked. And these our Lord doth evidently instruct, reprove, rebuke, and blame, praise, correct, move, exhort, comfort, the same he threatens, and promises them also joyful things, &c. This is no light or common example, but of the Son of God, the high and most blessed Bishop: teaching us how we should deal with congregations after the capacity and disposition of every one.

### Paragraph 141

And not without cause he chooses seven of the most noble cities of Asia. It is certain that Asia was among the first regions inhabited, and from there people were dispersed into other parts of the world. It is also certain that the Devil set up his throne in Asia, and there reigned in men through idolatry, murder, ambition, avarice, uncleanness, and base pleasures. For the proverb says, “The laughter of Jove.” It is known what the apostle wrote to the Ephesians in the fourth and fifth chapters. Therefore, our Savior Christ would overthrow that throne of the Devil and set up his throne of righteousness and holiness. Therefore, he goes first and chiefly to those of Asia, that by their example the whole world might be corrected and amended.

### Paragraph 142

And amongst other cities of Asia and Ionia, Ephesus was most famous, called in the old time the light of Asia. And amongst the twelve cities of Ionia, it was accounted the principal port. Neither was there any other Church richer or more beautiful anywhere in Asia than the temple of Diana at Ephesus. It stood in the midst of the city, a great wonder of Greek magnificence, as writes Pliny. This temple is said to have been built in two hundred and twenty years by a king of all Asia, and set in a marshy ground so that it would not feel earthquakes or the opening of the earth. The length of the temple was 425 feet, the breadth 220 feet. It had 107 pillars, dedicated by so many kings. Seek the rest out of the summary by the famous Dr. Joachim Vadian. The apostle Paul first illuminated this city with the name of the Gospel. His epistle to the same city remains, and a plentiful story in the Acts of the Apostles. After Paul was executed, John went to Ephesus, and from thence preached to all Asia; from thence also he was brought to Rome to the emperor Domitian. To Ephesus he returned after his exile, and there at length, as ecclesiastical histories testify, he slept in the Lord.

### Paragraph 143

And before every epistle, especially that of the Ephesians, is placed a command, write. This command gives authority to the writing, so that we may not ask, whether this writing ought to be credited? and why it should be believed? For here is the express commandment of God, and the divine authority, concerning which to inquire curiously is thought not without cause unlawful. Moses wrote by the commandment of God. And by the same commandment of God wrote also the prophets and Apostles. Which writings are not believed to be authentic? Certainly John said truly and wittily: he who knows God hears us; he who is not of God hears us not. 1 John 5. Curious questions cease, where the mind of the godly, or of any poor sheep, knows the voice of his Lord and Shepherd.

### Paragraph 144

And let no man think that this letter, being written to one angel—that is, a bishop or pastor—appertains nothing to the Church. For to the end of the letter is added an acclamation: “He that hath an ear let him hear what the Spirit says to the congregations.” Therefore, the pastor is named, but the sheep are not excluded. All ranks and states in the church know what is said to them. Ignorance says: “That which is written to the Romans concerns me nothing.” Yet nevertheless it is addressed to the angel [note 7.9], to the intent that pastors may be admonished concerning the state of the Church.

### Paragraph 145

The argument of the first epistle is this. Christ declares that he rules over his church, that he takes charge and government of the same. Some things he praises, and some he blames. In the meantime, he exhorts to repentance, threatening grievous things and promising most joyful ones. And also, he applies this epistle to all churches and communicates it to all congregations in the whole world. But the epistle is exhortative, for it instructs the churches, exhorts, and directs.

### Paragraph 146

And first, he shows who he is, from whom the epistle proceeds, that he may give authority thereunto; and may declare also that he is the head of his church, the Bishop, Duke, and governor. That part is taken from the image set forth in the first chapter. And follows the prophetic manner of speaking: This says he who holds the seven stars in his right hand. For the Prophets say likewise: “Thus says the Lord God of Israel,” “Thus says the Lord of hosts.” “Thus says the Lord,” who brought you out of Egypt, &c. And two special things he repeats from the former description, whereby he will be known and do us understand how he being Lord and Bishop rules and maintains in his church. First, he affirms that he holds in his hand the seven stars. The hand is a token of working, protection, or deliverance. The stars we have heard are the ministers, and the ministry of the word, or the church. Therefore Christ holds the ministry in the church, and the ministers work the salvation of the faithful. After, he affirms that he walks, not sleeping or doing nothing, in the midst of seven golden candlesticks. In the midst said, to the end we should understand, that he gives himself indifferently to all men, and rules over all with like power and government. Fulwel wrote hereof. What, says he, is to walk or to be in the midst of congregations, but to assist them, keep, instruct, help them, and by all means to watch over them? For which cause he says also in the last of Saint Matthew: “Behold, I am with you always, even unto the end of the world.” Hereof you have a most excellent figure in the law: wherein amongst other things which pertained to the ministry of the high priest, he had charge of oil and of seven candles; for those must he snuff and pour in oil, when it wanted. So Christ, the high and true Bishop, has the charge of the seven candles, that is to say, of all congregations; and is careful that they want not that oil, which is mentioned in the 44th Psalm. He watches that they want not the fire and light of the truth. Finally, he snuffs and purges by faith, whatever thing so ever has need to be purged in them. Thus far he. Which things when they hear, which make the Bishop of Rome head of the church, it is marvelous if by and by they understand not their folly and madness. Here the Lord adds also, that he knows the works, to wit, all both good and evil, as well of the Bishop as of his church. For the Lord knows all things, and is head Bishop of the Catholic or universal Church, which also remembers the thoughts of all men in the world at one instant; who sees what is done, and what is not done, and what things are needful, nothing escapes him. And such in deed ought he to be, that is head universal of his church. And this sentence is repeated, “I know thy works,” in the beginning of every epistle. And verily it is full of comfort, when we hear that Christ knows all our doings. For we believe also that he has a faithful care of all our matters.

### Paragraph 147

Now this great Bishop commends certain things within the congregation of Ephesus. For good works in deed are approved by Christ, and he praises them, so that he might encourage those who run in his way. First, he acknowledges the labor and patience of both the Bishop and the Church. Labor encompasses thought and care in the way of God, mortifying the flesh, the study of good works, but chiefly the cross and persecution, which history testifies to have been extreme and cruel in the time of Domitian. And unless the persecuted have patience, they cannot endure the labor. Holy patience keeps us in work and holy labor.

### Paragraph 148

But lest patience should be stretched to those things wherein impatience is considered praiseworthy,[note 7.13] he adds the second point that he praises in them, that you cannot bear evil men. And by these evil men he means not weaklings, or those who err without malice: but the prophet David also says, Psalm 119, “I have hated the wicked; your law I have loved.” What we should do with the weak in the faith, or with those who err through ignorance rather than obstinate stubbornness, the Apostle has taught us in Romans 14. The example of our Savior has also taught, bringing again that strayed sheep upon his shoulder into the sheepfold. Therefore, the Lord speaks here of the obstinate, of the deceivers who delight to err themselves and to draw others with them into errors; no Christian patience bids us to bear with such men.

### Paragraph 149

And in the words following he declares of what sort those evil men were. And you have examined those who say they are Apostles, but are not, and have found them liars. Observe, he speaks of the false apostles, of whom in John’s time there was exceedingly great abundance. For they were Zarbians, mixing law with grace, and attributing justification to the law and to our own righteousness. Which the holy and great council at Jerusalem condemned, as it appears in the fifteenth chapter of the Acts of the Apostles. Such a false apostle was Hebion. Eusebius mentions in the book of Ecclesiastical History the twenty-seventh chapter. Here was added Cerinthus, that heretic, not apostle. There were more also, of whom some denied the humanity of Christ, some his deity. Against whom John wrote in his Gospel and in his Epistle, and Irenaeus in the first book against heretics. These the Lord denies to be apostles, or apostolic men, which the apostles have also denied, Acts 15. And therefore the apostle Jerome in his canonical epistle says, “Who is a liar, but he who denies Jesus to be Christ?” But if such trouble were in the churches while the apostles were yet living, and there were then so many deceivers, what marvel is it, though in the dregs of the world, that is, in this our last time, there are not a few sounding where they are now that twisted dissensions and troubles in the defense of their error? The Gospelers themselves say they are at variance. God is God of concord; how then should I believe that God is among those who cause strife? So might the Sophists also have reasoned in the apostles’ time.

### Paragraph 150

And here we have a fitting way, in what manner the churches should proceed, while troublesome persons, like false Apostles, attempt to divide the Church further. [note 7.16] For such ringleaders must be tried and examined: and they must be tried according to the Christian belief and doctrine of the Apostles, and inquiry must be made whether they are Apostles and true men, or false Apostles and liars. When we have found them to be false Apostles and liars, and they persist stubbornly in their wickedness, they are not to be tolerated; as indeed the Ephesians did not deem safe to bear with such deceivers. And we must know that pastors ought to act one way, the Christian magistrate another way, and the people a third way, not to tolerate open heretics. For the pastor not only does not tolerate them, by being watchful and wary of those wolves, but assails them with wholesome doctrine, and rejects them from the flocks of Christ. But the magistrate, because he is a Christian magistrate, and by his duty also, not only as a private person but also as a magistrate, ought to serve Christ; he ought also with the sword of justice to drive away poison from the church, and to punish manifest blasphemies. [note 7.17] And the people are commanded neither to hear them, nor receive them, nor to have anything to do with heretics, and so not to tolerate them. They may therefore be ashamed of their faultiness and the persistence of their perverse patience, which think it no shame to maintain heretics, and to flatter the manifest enemies of Christ and the Church. Psalm 15. He is praised who does not value the wicked; that is to say, in whose sight the wicked man is vile. Therefore, he is rightly blamed, whoever flatters the ungodly. And the hatred in deed is rather against wickedness than against the person of the wicked, which itself is commanded to be loved. The Devil at this day raises up the old heresies of Hiero[?]m, Ireney[?], and others in Ter[illegible]paniarde, and in the Anabaptists, Libertines, and other monsters, so that the thing itself and the danger there[illegible] commands us to watch and to drive away the most ravenous wolves from the holy Church of Christ, which nevertheless set forth nothing more than patience and charity, for this intent verily that they might be spared, and might unpunished teach what they list against Christ, and work against his church, yea tear it in pieces with their wicked teeth.

### Paragraph 151

But when these evil men are not tolerated, but oppose those who deceive and are deceived, [note 7.18] a great conflict arises, involving labor, thought, anxiety, watchfulness, and injuries to be suffered for the name of Christ and the defense of the truth. For unless we are diligent and patient here, the deceivers will prevail. In this, the church of the Ephesians behaved admirably, so that the Lord now commends exceedingly the courage, patience, and constancy of the pastor and his church. For these things should not be understood to refer to that patience whereby evil men are tolerated and permitted to proceed in their malice and deceitfulness. Otherwise, this passage would contradict what preceded it. This is what the learned interpreter seems to have intended to avoid. For he reads, “You have patience, and have suffered,” where the Greek says, “You have suffered and have patience.” He altered the order, and would not place “you have suffered” before “you have patience,” lest anyone misunderstand that they had tolerated the false apostles. Instead, he placed “patience” first and “suffered” afterward, so that we might understand that they did not tolerate evil men, but the evil wrought by evil men. Thus, they labored with patience for Christ’s name, namely, to maintain it against harmful heresies. And he adds, “You have not fainted, though weary and broken with labors.” For we are taught to overcome through patient constancy, which is rightly called the fulfillment of every good work.

### Paragraph 152

All and every of these things we must apply to ourselves and understand how we may also, this day, please Christ our Redeemer, King and Bishop. Truly, we walk in the same steps wherein we see the congregation of the Ephesians to have walked.

### Paragraph 153

What thing did he condemn in that church? That they had abandoned their first love. When they first received the Gospel by Paul, and afterward by John and other godly men, there was seen a great fervor in the words and deeds of the faithful; which thing may be gathered both from the acts of the Apostles and also from Paul’s epistle to the Ephesians. They loved God and their neighbors with a most fervent zeal. They burned in reforming of manners. But in process of time, this heat was cooled, and they grew colder in true godliness. This great mischief he rebukes in them, and then desires to have it redressed. And here let us note that not only revolting from idolatry and other great crimes are imputed to the church, but also if we slack in any thing in holy zeal. So that hereof we may learn how holy and blameless we ought to be before God. Doubtless we cannot here excuse ourselves before the divine majesty, who has been for thirty years past more fervent in this congregation than we are today. &c. Our Lord God enlighten our minds, that we may please him. To whom be glory.

## Sermon 8: ¶ The second part of the Epistle to the Epheſians where is spoken of Penance and of the Nicolaites.

### Paragraph 154

.

### Paragraph 155

Remember therefore from whence you have fallen, and repent, and do your first works. Or else I will come to you shortly, and will remove your lampstand out of its place, except you repent. But you have this, because you hate the deeds of the Nicolaitanes, which deeds I hate also. Let him who has ears, hear what the Spirit says to the congregations. To him that overcomes will I give to eat of the tree of life, which is in the midst of the Paradise of God.

### Paragraph 156

The accusations our Savior Christ uses against his servants who sin do not, I am certain, aim to cause those overwhelmed with reproaches to be ashamed, despair, and perish; but rather that they should be restored and live. Therefore, the Lord Jesus immediately adds an exhortation to repentance, that they may be saved, and also describes the true and lawful penalty.

### Paragraph 157

For we heard what he rebuked in the congregations of the Ephesians; let us hear now what the Lord requires of them, and how he seeks to have the error corrected truly by repentance, to which he exhorts. For he has said that the Lord strikes and heals; chiefly in this case. Which doctrine surely is proper and perpetual to the church of Christ.

### Paragraph 158

He speaks primarily of three things in this matter, in his counsel or exhortation to amendment. [note 8.2] First, he reminds them of from what they have fallen: that is to say, with how great love they had burned previously, and how they could now be deceived. He shows how fortunate and blessed a state they had maintained until now, and how unfortunate and shameful their current state is. For acknowledging the transgression is the beginning of repentance, if, illuminated by faith, we consider what benefits we have lost and in what misery we now find ourselves. He who believes he has lost nothing will not be moved to make any search or inquiry; he who thinks he has fallen from no happiness will not consider how he may be restored. Therefore, in the matter of a godly life, acknowledgment and confession of sins must come first, by which we may lament before God our poverty and misery. Indeed, they do not fall from happiness who were never in any happiness. [note 8.4] Therefore, we say that holy men may fall, and also be restored by repentance. Then, after the acknowledgment of our failure, follows repentance: that is, the turning again of our mind, so that we do not proceed forever like madmen and fools, walking in the way of vanity and unrighteousness. [note 8.5] For repentance is a conversion and change, whereby we are turned away from evil and converted to good, returning to our senses and working righteousness.

### Paragraph 159

Finally, what follows explains this repentance. Repeat your former deeds: be fervent again in love, perform the works of faith, which are fruits worthy of repentance. See, there is no need for new laws [note 8.6] or lengthy debates concerning amendment. Briefly, repeat your former deeds, not of the flesh, but those which you began when you first received the Gospel and were born again in Christ. Ultimately, this is the true reformation: to do the former works of God, not the latter, which the error of the world has devised. This truly is, and will be, the true repentance: namely, the acknowledgment of sin, a turning to God and to good, and a turning away from the Devil and from evil, and the working of the first righteousness, the righteousness of faith in Christ. There are many and varied discussions about repentance and its parts, concerning the contrition of the heart, the confession of the mouth, and the satisfaction of that work. But just as none is briefer than that of Jesus Christ, so truly is there none better or more certain.

### Paragraph 160

To these his exhortations and godly cautions, he adds most grievous threatenings, if a great danger could not be avoided. And he speaks few words indeed, but understands a great evil that cannot be spoken or declared. “Unless you repent, I will remove your lampstand from its place.” The lampstand, as the Lord himself has explained, is the church. It stands in its place while it leans toward Christ, and is preserved by Christ so long as the preaching of the truth is maintained within it. It is indeed the church of Christ. She is removed from her place when she is without the preaching of the truth, and no longer leans upon Christ, nor is defended by him, but is forsaken, and is no longer indeed the church of Christ. This is done by Christ himself, through his just judgment, when our unthankfulness and a life that cannot repent drives God to depart from us, to abandon us to our error and darkness, and to leave us to deceivable men, and so on. This understanding Arethas acknowledges, who says that to remove the church is when it is left bare and destitute of God’s grace, because of which nakedness it wavers in doubtful perplexity and in storms cast upon it by wicked men. And truly we see how at this day the church of the Ephesians is removed from its place, and no longer enjoys the wholesome doctrine of Christ, nor stands upon the wholesome rock Jesus Christ, but is oppressed by the pestiferous doctrine, or rather madness, of Muhammad, and lies in sorrow under the feet of the Turks. We see at this day in Germany (and it is the greater pity) many lampstands removed from their place, not without the great triumph of Satan and the loss of souls. Moreover, it is also to be observed in this threatening that, without terror, he says, “I will come to you shortly.” For it is a phrase of speech. For we also say, “I will come to you by and by,” that is, I will come to avenge and punish, and perhaps sooner than you look for. Most certainly, whenever I happen to come, I will take punishment of you. Let no man therefore think to escape unpunished in a life that cannot repent.

### Paragraph 161

Again, where the Lord repeats, “except you repent,” he plainly testifies that the bosom of God’s mercy and clemency is readily open, if we do penance, however we have offended him before. In the meantime, we learn here openly and most certainly that we can by no counsels or consultations, by no armies nor policies, prevail against our perils unless we repent. Therefore, unless we will have our churches subverted and given over to be seduced and destroyed by the devil and his seducers, let us repent in time and receive again the first love.

### Paragraph 162

Again he commends the singular virtue in this congregation, especially for that they have hated the doings of the Nicolaitanes, which God himself also hates. Let us mark every word here. He does not say you did slyly or avoid and abstain, but that you have hated. The force of hatred is great, moving one to persecute, that you hate.

### Paragraph 163

Moreover, he does not say, “You have hated the Nicolaitans,” but “the works of the Nicolaitans.” For we ought to hate no man’s person in itself, but the vice in the man, so that when we forsake it, we should love the man with all our heart. It must needs be a great evil, which God himself confesses that he hates. Here all congregations shall understand that they ought also by all means to hate the heresy and abomination of the Nicolaitans. Albeit that at this day the name be extinguished, yet the heresy and abomination of the Nicolaitans remains.

### Paragraph 164

This Nicolas was from Antioch, one of the seven deacons, mentioned in chapter 6 of the Acts. He is said to have revolted from the purity of faith, as Judas did. And where he was previously a Gentile (for it is said that he was a proselyte), he returned in some respects to gentility, like a dog to its vomit. The Nicolaitans are also Gnostics and associates of Carpocrates, filthy and most wicked people. Clement somewhat excuses Nicolas in Eusebius in book 3, chapter 29 of the Ecclesiastical History. But that excuse seems insufficient or just, since all the ancients accuse him with one voice, and namely, the judgment of God in this present and in the following epistles. Irenaeus condemns him from this same place, in book 1 against the Valentinians, chapter 27 and following. Tertullian, at the end of Prescribed Heretics, touches cleverly on the facts of the Nicolaitans and detests them. Nevertheless, he does not explain them, but passes them over. And I do not know how cleverly Epiphanius has uttered and declared the wicked and abominable acts, neither to be thought nor told, and most beastly filthiness, such as has not been heard of in the heresy, 25, 26, 27, and 31 and following. Philastrius and Augustine also have touched on the Nicolaitans, either of them in their lists of heresies. Shameful details will not allow me to recite. It is enough if we know what the Lord himself explained in the epistle to Pergamum, calling the doctrine of the Nicolaitans the doctrine of Balaam the false prophet. But who does not know what counsel he gave to King Balaac of Moab and Midian, and how he prostituted fair women to the young men of Israel, by whose acquaintance they both defiled themselves with fornication and ate also of meats offered up to idols, becoming partakers of Baal-peor. Let him who will read Josephus in book 4 of Antiquities, chapter 6. And without doubt, the sacrifices of the Nicolaitans seem to differ nothing from the secrets of Priapus, or Gerecinthia, or the mother of God, and the nightly service of Bacchus. Irenaeus openly signifies that the Carpocratites, who are also called Gnostics, did not abhor images, but painted and fashioned into images Jesus and Paul, along with images of certain philosophers. And that the image of Jesus, as they say, was made expressly by Pilate, who commanded that the face of Jesus be painted vividly. But however that was, it is certain that the actions of the Nicolaitans were in ill repute because of their fornications and adulteries. And that the Nicolaitans abstained not from images, nor from meats offered to idols. Against this error, Paul also wrote many things.

### Paragraph 165

Hereof let us learn to abhor and fly fornication, and never to think of restoring brothels or other places of whoredom. Fie for shame. Let us learn hereby to keep virginity, single life, and lawful marriages, and avoid those called Nicolaitans. Let us learn hereby to keep ourselves well from idols, idolatry, and from all strange kinds of worshipings. All those God hates.

### Paragraph 166

And with an acclamation he pierced the ears of all, moving all to attentiveness and holy obedience. And he applied this doctrine to all times and to all congregations throughout the world. He used his accustomed speech, repeated so often in the Gospel: “He that hath ears to hear, let him hear.” It is not in our strength to hear and obey God. For God prepares our ears, and with his grace frames and draws our hearts. And let those to whom the grace of God is given, beware lest through their negligence, vanity, and lightness they lose it not. Let them show such diligence as God in his word requires and prescribes. They that do have ears to hear. He says therefore, take heed to what God now speaks, and whose hearts he now stirs and moves, that you lose not this grace through your negligence. Be diligent, attentive, and circumspect, stirring up within yourselves the gift of God.

### Paragraph 167

Now also he urges diligence by authority. The Spirit of God speaks and reveals these things, not the spirit of men or of error, for God speaks by his scripture, which is recognized as the spirit both of the Father and of the Son. Moreover, he applies all things and every thing to all congregations, where he says, “What the Spirit says to the congregations,” not to the congregation. It is now manifest, and beyond all controversy, that those seven churches represent a figure of all churches throughout the whole world, and that all are instructed in those seven.

### Paragraph 168

Moreover, to ensure that nothing is lacking in the just exhortation to repentance, faith, and diligence, he adds a most ample promise and employs allegorical language, so that it might have the greater grace. [Note 8.16] To those who overcome, he promises to give the fruit of the tree of life, planted in the paradise of God. He alludes to the second chapter of Genesis and transfers the meaning from earthly things to heavenly ones. The paradise [Note 8.17] of God (which some understand to be the church) is that everlasting blessedness and felicity of which the Lord spoke to the thief, saying, “This day you will be with me in paradise.” Herein is the tree of life, Christ, communicating his eternal life to us, which we enjoy and experience while being conveyed into heaven by him and living with him. Finally, this is that ambrosia or godly drink which the heavenly Father gives us to drink. But this great and wonderful good does not happen to everyone, but only to him who overcomes. For Adam did not overcome; he had died in defeat. Therefore, if we overcome the flesh, the Devil, and the world, and that through Christ, we shall also live in the world to come with Christ.

### Paragraph 169

The Complutensian Bible has it, which is in the midst of the Paradise of my God. And Arethas explains it [note 8.18], and says: Let no man be offended by this. All humble things agree to the dispensation of the incarnation, which was made for our cause, since he himself in the Gospel says: I ascend unto my Father, and your Father, to my God, and to your God. &c.

### Paragraph 170

And thus far, concerning the Epistle of Jesus Christ by John to the Ephesians, and what benefit our churches, and each of us, may receive from it. May the Lord enlighten the eyes of our minds.

## Sermon 9: ¶ The second Epistle of Jesus Christ by John to them of Smyrna is expounded. And is an exhortati [illegible][illegible] to patience, and consolation in afflictions.

### Paragraph 171

.

### Paragraph 172

And unto the angel of the congregation of Smyrna write. These things saith he that is first and last, which was dead and is alive. I know thy works, and thy labors and poverty, but thou art rich. And I know the blasphemy of them which call themselves Jews, and are not: but are the congregation of Satan. Fear none of those things which thou shalt suffer. Behold, the Devil shall cast some of you into prison to tempt you, and you shall have tribulation for ten days. Be faithful unto death, and I will give thee a crown of life. Let him that hath ears, hear what the Spirit saith to the congregations; he that overcometh shall not be hurt of the second death.

### Paragraph 173

Jesus Christ, from the right hand of the Father, through the ministry of an angel by the Apostle and Evangelist John, exhorts the congregations of Smyrna, who are afflicted with every kind of evil for the word of God, to endurance and comfort while they now suffer under the cross, promising great things to those who overcome. Truly, there can be no exhortation or consolation of this manner or in this matter that is briefer. For it is wisely couched in the eternal wisdom of the Father, and is applicable to all times, and it agrees very well with all who mourn under the cross. For just as Christ at the right hand of the Father is the catholic or universal Bishop, so too is his doctrine general, which he himself also applies to all congregations in the end of this Epistle and in others. And he declares that he loves his church, and is present in it by his power and aid.

### Paragraph 174

And truly it is to be marveled at that nothing is condemned in this church, since some fault is found in conduct with all others. Therefore, the church of Smyrna was exceedingly excellent, though not without imperfections. For the Lord, in his goodness, does not hold small faults against us—of which the Prophet speaks, who will say, “My heart is clean? And from my hidden sins cleanse me”—so long as there is a fervent desire and zeal for godliness within us, and we are void of great evils.

### Paragraph 175

First, it is shown to whom this heavenly letter is sent,[note 9.4] to the pastor of the church of Smyrna, and to the whole flock. For the captain is said to have sought peace, or fled, or taken peace, when the whole army together with him has done this. And the histories bear witness that Polycarp was that same messenger or pastor of the church of Smyrna, ordained by the Apostles themselves,[note 9.5], namely by John the Apostle, Bishop there, and that he lived in the mystery of this congregation for eighty-six years. For so many years he declared himself before the Lieutenant Herod, when he was brought to execution. For in the fourth persecution of the church, Aurelius Antoninus and Aurelius Commodus, being Emperors, he was taken and brought before the governor. And at length, for the open and sincere confessing of Christ, he was burned. He had this very much in his mouth: that nothing ought to be received as true unless it were known to be set forth by the Apostles. Irenaeus affirms that when he was a child he saw this old father, a man of great years and reverence, in the third book and third chapter against heresies, where he tells many things of him besides. As also does Eusebius in the fourth book of the church history, in the fourteenth and fifteenth chapters. And Jerome in the register of the famous writers of the church. Eusebius in his Chronicles notes that he suffered martyrdom in the year of our Lord 190. Whereby it appears that he was made Bishop of Smyrna in the year of our Lord 84, or thereabout. For we said even now that he had been in that ministry for eighty-six years. And therefore, had he been Bishop of Smyrna many years before the setting forth of the Apocalypse, which was written in the ninety-seventh year, would God all pastors would set before their eyes this good Polycarp to be followed, of whom there remains a notable Epistle to the Philippians.

### Paragraph 176

After again is the author of the Epistle [note 9.6] declared, which is set forth with two titles, taken out of the former visions of John and description of Christ. Thus saith the first and the last. &c. Whereby is signified the eternal divinity of Christ, which wants beginning and ending. And of himself is everlasting. There is added that he was dead, and lives again, that is to wit, has risen from the dead. And this beginning accords right well to the matter. For they perceive that whosoever are afflicted for Christ and his Gospel, have a Lord and patron more mighty and more faithful, which in no wise can be overcome. Who can also in death keep his, like as he raised Christ from the dead, to the intent we might have an opportunity, that we shall live with Christ, even in death it seems.

### Paragraph 177

And now he comes to the matter itself, and that which he repeats in all his Epistles, he says here also: “I know your works, both good and evil. Do not think that I neither know nor care about your affairs. You are laid out in my hands; I know, see, and care for you and all yours.” These things greatly encourage us when we know that we have God watching over us, who has a care for us. And they also comfort us greatly, those who suffer, knowing that he who loves us and in no matter neglects us, has us always as it were before his eyes.

### Paragraph 178

And here particularly he declares what he knew: First in deed the afflictions, which verily they suffered in the present persecution of the Emperor Domitian. And affliction is as it were a general word, encompassing four kinds of suffering. For he recounts, concerning their substance the spoiling of their goods and their poverty; in their name, insults, reproaches, or blasphemies; by imprisonment and bonds, yea and death also. For through these afflictions godly men are tested for the sake of the wicked. And in these may be included all other kinds of tribulation. The which the Epistle of Jesus Christ recites in a godly order. There is nothing therefore of these matters which the Lord Christ does not know.

### Paragraph 179

Poverty has the first place. We should not take this to mean here a spiritual poverty, signifying humility of mind—though it is certain that the church of Smyrna possessed that virtue—but rather a poverty and lack of all things due to the confiscation of their goods. For in times of persecution, by virtue of royal proclamations, the goods of faithful Christians are confiscated for the king’s use, or permitted to soldiers, nobles, or others to take at their pleasure. The faithful, thrust out of their homes, are either driven into exile or reduced to begging. Would God we lacked examples of this today. Let us learn from this to bear and suffer such circumstances patiently, being persuaded that God knows our necessity. And because it is difficult for an honest man to hunger and want with his family, for comfort and consolation he adds, “But you are rich.”

### Paragraph 180

This appears to the world as a paradox, or something unbelievable. What will they say, is he rich who has nothing and is reduced to the state of beggars? There are undoubtedly goods and riches of the mind much better than material possessions.

### Paragraph 181

For this may be had, without the true felicity, of rich men of this world, who live a most miserable life. Again you shall see a poor man, concerning worldly goods, but furnished with the richness of the mind, for this cause only to be happy and most blessed. He covets nothing, he is content with his vocation: Neither would he change his state with the most wealthy and rich kings. Contrariwise you shall see rich men but of an evil conscience, and therefore thoughtful and burdened with cares, and never merry. You shall see poor men, but with merry hearts to lead a joyful life. Why then should it seem marvelous, if he that is stripped of his worldly goods for Christ, and enriched with the gifts of the mind, is glad and rejoices in God, and takes a good part in all chances, and for the same cause is judged to be verily rich? Doubtless the wise men of this world know also, that the only wise man is truly rich. Which is eloquently discoursed of Cicero. Arethas saith, in spiritual matters having a treasure hid in the field of thy heart, which is Christ, by reason of whom thou art rich also: Since thou hast him thy protector, who also when he was rich, for us became poor. &c.

### Paragraph 182

Secondly, blasphemy is mentioned, by which we understand all manner of slanderous remarks and insults, whereby the reputation and standing of the faithful are harmed. To this sort belong those who call the adherents of this faith heretics and schismatics; they deem them wicked people, despisers of God and his saints, and enemies of God’s service, and therefore plagues of the common wealth. If such people are tolerated, the common wealth must necessarily be destroyed. And these things often trouble good men more grievously than the loss of their possessions. For people value a good reputation as much as great wealth. Therefore, the Lord in the Gospel of Saint Matthew, chapter 10, heals this disease with many words and exhorts people not to do anything unworthy of the name of Christians to avoid such infamy.

### Paragraph 183

In the meantime, he declares also what motivated the authors of this mischief, whom he also blames severely, so that the faithful might understand how greatly God disapproves of the enemies of all godliness, and so that they might also care less for their hatred and persecution. They claim to be Jews, when they are nothing of the sort. Thus also Paul dealt with the Jews in his second letter to the Corinthians. The Jews are called confessors, honorers, and faithful servants of God. But these blaspheme God’s name, they attack the true faith, and oppress those who profess and worship God. Therefore, they are not Jews. What then? The synagogue, congregation, or assembly of Satan. Thus, the very Son of God plucks off the vermin from these young men, to the comfort of all those who suffer persecution. It should not grieve those who present themselves with proud titles that they are condemned by such harlots, the children of the Devil. Christ gives them a true title and calls them not the holy and catholic Church of God, but the conspiracy and school of Satan, as in whom, not the spirit of God, but of Satan, inspires lies, jugglings, deceits, blasphemies, fires, and deaths. Therefore, let it not grieve you at this day, even if it is your fortune to be condemned for the Gospel, by those who call themselves most holy, most shining, most reverent, and most irreproachable prelates and patrons of the old church, religion, and catholic faith, which have on their side, councils, fathers, so many successions of bishops, the prescription of so long a time, and the consent of so many realms. They desire nothing less than to be called such, but rather are the champions of Antichrist, and the professed enemies and tramplers underfoot of all Christian piety, for whom everlasting destruction is prepared.

### Paragraph 184

After this, he offers a most evident exhortation and consolation, placing the Son [note 9.13] before them, and saying, “Fear nothing of all that you shall suffer.” The Son of God himself feared the cross and death, and it is a natural thing to fear evils and death. Therefore, we are not commanded to be inhuman, nor to say, like stones, that the same things grieve us not, even though they torment us exceedingly. Rather, the faithful are encouraged to stand strong in the faith, and to do nothing unworthy of it for fear of punishment. We are therefore commanded boldly and freely to despise or suppress fear, and to seek strength by the Spirit of God, and to exercise it in temptations.

### Paragraph 185

There follow reasons whereby he may obtain that which he hath persuaded, may confirm, comfort and exhort them with patience and constancy. He prophesies therefore to the godly, what things they shall suffer. And touches also the third kind of affliction, imprisonment and bonds, which we understand to mean all punishments whereby our bodies are tormented. But to be warned before of the evil is a great benefit. We are more easily overcome by unprovoked perils. And therefore the Lord in the Gospel after Augustine, the tenth chapter, and after John in the fifteenth and sixteenth chapters, tells his disciples of many evils that should come unto them, and adds thereto: These things have I spoken to you, that when the time shall come, you might remember that I have told you before. So now also faithfully warns the faithful in this Epistle.

### Paragraph 186

And he touches the author of these evils, saying: The Devil will cast some of you into prison. Therefore we perceive that these evils arise from the common enemy of mankind and of the salvation of the faithful. Whereof we may conjecture that he seeks to intercept our salvation, and that we ought therefore to stand more earnestly against him. The soldiers when they hear that their old enemy is at hand, do not become sluggish, but cheerful. But the Devil inspires evil men, corrupts princes and magistrates, which attempt persecution against the church. We read that Satan afflicted Job, that is to have provoked the Chaldeans and Sabeans to kill his servants and drive away his cattle. Here therefore they may see with what authority they are encouraged, who at this day persecute the church of Christ, for the profession of the truth. The godly have that which may comfort them: For they hear that the same filthy beast is set against them, which so often has been vanquished by Christ the Prince of the faithful, and of the faithful through Christ's aid, may without any difficulty be overcome. And verily the Lord permits to the Devil and devilish men power over his servants. If you may well ask why, hear: That you may be tempted. God does not permit his people to Satan that they should perish, but that they should be tempted and tried. Therefore to a good end are delivered to the fire, that we might be purged from our sins, that the virtue of our faith might shine, and God might be glorified, and we made the purer. Who therefore will hereafter be impatient, when we hear that we for a great good are put to evil? We read in the third book of Wisdom: As gold is tried in the fire, so are the faithful proved. This parable has Saint Peter expounded at large in the fourth chapter of the first Epistle. Where he that will may have it more abundantly.

### Paragraph 187

Moreover, the time of tribulation is also appointed, as it were, for ten days. The number ten signifies a multitude. For Jacob says to his father-in-law, “Ten times you have changed my wages.” Genesis 31. And Numbers 14 says he was tempted ten times, that is, often and many times. Job also affirms himself in chapter 29 to have been wronged ten times. Therefore, the Lord says at this present, “You shall be diversely and greatly troubled with evils.” Nevertheless, because he does not specify months, years, or ages, but days, he prophesies that the evils will not be continuous, but that there will always be spaces between to breathe in. Indeed, the shortness of persecution is first shown in Essay 26. Secondly, Saint Peter in 1 Peter 1 comforts the faithful. It is the part of the faithful not to prescribe to God, but whether we are put to pain for a long time or a short time, to take it patiently. Let us think rather that in the long continuance of evils, there is also some end foreseen by the Lord, and that in the same time of breathing, we must repair the evils and return to battle.

### Paragraph 188

Finally, the godly are encouraged by a most ample and large promise, in which is also included the fourth and most grievous kind of affliction, and even bitter death itself, through fire, hanging, sword, water, etc. But if you are not afraid of death, but overcome it, and offer yourself to God, then I will give you, says the Lord, a crown of life. Hereunto is annexed the state of the Epistle, and some of all. Therefore be faithful, cheerful, constant, even to death. For the Lord also says in the Gospel: Whoever perseveres to the end shall be saved. And we read that the Apostle has said, if we die with Christ, we shall live with him. And truly, the crown of life is nothing other than eternal life, that everlasting, celestial, and unspeakable joy. And the Lord alludes to conflicts, after which, when fortunately finished, the victors are crowned. Blessed is the man, says the Apostle James, who suffers temptation, because when he is tried, he shall receive a crown of life, which the Lord has promised to those whom he loves. Similar things the Apostle Paul has also written in the first to the Corinthians, chapter nine, and in the second to Timothy, chapter four. Therefore let it be hard hereafter for no man to lose this temporal life, since, as it is lost for Christ, we shall receive eternal life; otherwise, we will, or will not, we must die. Let us therefore be content rather to die blessedly than to live miserably, so that we may please God.

### Paragraph 189

Finally, just as in the end of the first Epistle, he communicated and applied the same wholly to all times and churches, lest any should suppose that these things concern them nothing. So in the end of this Epistle also, he both declares the Spirit to be the author of all these things, and exhorts all men to hear and obey diligently, and affirms it to be written unto all congregations in the world for edifying. Moreover, the promise of life he communicates, saying: “He who overcomes shall not be hurt by the second death.” This is spoken to all men and women, if you overcome. Therefore, we must overcome the world, the Devil, the flesh, and all temptation. And we must overcome by him, through faith, by his Spirit dwelling in us, and that we should walk in the way wherein he has commanded us to walk. If you overcome, you shall not be hurt in the second death. Thomas of Aquin says that the first death is of sin, the second of pain. We understand plainly by the first death the natural separation of the soul from the body, which also comes to us because of sin, as appears in the third chapter of Genesis. This same [illegible] has come to good and evil. For we are all earth, and to earth we shall return. And immediately follows the second death and the second life: Those who believe in Christ overcome, feel nothing of the second death, but live, as the Lord himself assures us in the third and fifth chapters of John. He shall not come into judgment, but has passed from death to life. But the wicked or unbelievers are conveyed straight ways from corporal death to everlasting death. Not that their souls can die, that is, cease to be, or that their bodies will not rise again, but that being deprived of that celestial and divine life of Christ, they feel everlasting torment, which state is rightly called death. These things are unknown to worldly men, who know no other life but this temporal. But God’s truth teaches us that there is another life and death after this, namely, the celestial and death infernal, or full of perpetual sorrows. That same is full of consolation, that we hear how the faithful, after paying the debt of this temporal life once, shall no more feel any torments. What then do the monks and friars prate of purgatory, babes? Let us praise our Savior Christ, who has delivered us from death and given us the hope of life everlasting, to whom be glory, praise, and so forth.

## Sermon 10: ¶ The first part of the third Epistle of the cu ſtancie and cofeſſion of Christ in the time of perſecutio.

### Paragraph 190

.

### Paragraph 191

And to the messenger of the congregation in Pergamum write: This says he who has the sharp sword with two edges. I know your works and where you dwell, even where Satan’s throne is, and you hold fast to my name and have not denied my faith. And that in the days in which Antipas was my faithful witness, who was slain among you where Satan dwells.

### Paragraph 192

The third epistle among those seven heavenly ones proceeding from the right hand of God,[note 10.1] is written to the pastor and congregation of Pergamum. The argument is as follows: First, he commends the constancy of their faith in cruel persecutions. Then, he rebukes those who clung to the sect of the Nicolaitans. Afterward, he exhorts them to repentance. And this doctrine he applies afterward to all congregations throughout the world. Finally, he promises most ample rewards to the faithful.[note 10.2] We understand from this that the congregation of Pergamum is presented as a type or a mirror to all churches, showing how they ought to walk before the Lord: first, whenever persecution arises; secondly, when heresies break out. For by its example, he teaches all to suffer adversity patiently and openly to profess the true faith, and also by Scripture to reprove heresies, and in fleeing from them to despise them.

### Paragraph 193

However, all the Epistles have certain things in common, and especially three. It plainly expresses to whom the Epistle is sent, as in this present case, to the messenger of the congregation of Pergamum, namely to the pastor, whoever he was (perhaps Antipas), and to the whole congregation, as was said before. It is shown moreover who speaks here, or who is the author of this Epistle: even the Lord himself. This lends authority to the writing. For it is not to be thought that the word of God is not as it is spoken because it is written by a man, dictated by a man, or written with ink, either on paper or parchment. For these things do not make the word of God cease to be the word of God, any more than water ceases to be water if it runs out of a conduit of wood, lead, brass, or stone. For water always remains water. The diversion of the conduit pipes does not change it from being water, as it is in substance. So says Paul, that he may be bound, but the word of God is not bound. A man may be stoned, hanged, or burned, being a preacher of God’s word; the word of God that was put in the mouth of the preacher is not burned. The Lord puts it in the mouth of another, that the truth should not be extinguished, but continually might sound in the church. Finally, without cause, at the beginning of every Epistle, Christ intimates that he knows all things of the church. I said before that this is as it were the foundation of the fear of God, and of his true service. For imagine a man who thinks that God neither sees what men do, nor knows what they think in their hearts. Shall not that man fall into all ungodliness? He will cry, let us do what we list, since God knows not what we do. Who will not cast off the hope of reward, and the practice of good works, after he is once persuaded that God knows not our works? But if he knew them not, how would he judge the world?

### Paragraph 194

Nevertheless, in every epistle, be certain of special and peculiar things. Of this sort, in the epistle of Pergamum [note 10.6], is that from the first vision and description of Christ, at the beginning of the epistle, he takes up the sworn, sharp, and two-edged sword, which we heard come out of the mouth of Christ. By this is signified the judicial power of equity and justice, and also the deliverance of the good and the punishment of the evil, for the sword is given to the magistrate as an authority to punish the evil and defend the good. Christ himself defends his own, and he hews his adversaries in pieces. The sword is the very word of God, most sharp, two-edged, and piercing the hearts, for it animates the godly and discourages the wicked. Christ therefore governs his Church as a Judge and defender, most rightful and just, which has his sword not in his hands, but in his mouth, and with his spirit and word comforts and preserves the faithful, but fears and brings death to the unbelievers. Full rightly therefore is this beginning applied to the cause that follows, concerning the cross of the faithful and expelling the Nicolaitans. For it is Christ, by whose virtue these things are successfully brought to pass.

### Paragraph 195

Furthermore, the specific actions of this congregation are recorded.[note 10.7] He commends the church for its remarkable steadfastness in faith and its public declaration of that faith, even amidst perilous trials, temptations, and persecutions. It appears to be a straightforward account showing that the Lord is aware of their suffering and how severely they are afflicted, yet praise is interwoven throughout. This commendation serves as an encouragement, urging them to continue in the actions they have already undertaken.

### Paragraph 196

He says he is not ignorant of where the church of Pergamum dwells: even there, truly, where Satan has fixed his seat or throne. That is to say, I know in what case you are, in what dangers, and with whom you are contending. He says not, I know that you sit in the seat of Satan, but, I know that you dwell there where Satan has his seat. Christ therefore is not ignorant of the labors, sorrows, and temptations of the faithful. And the knowledge of Christ has a certain peculiar quality. For Christ so knows the matters of the faithful that he is both touched by them and has also a consideration or respect for his servants. And we see how Christ also places his throne there where the Devil has his seat, just beside it. At length he thrusts him out of his seat.

### Paragraph 197

And for two causes Pergamos seems to be called the seat, throne, and kingdom of the devil. For first as Arethas has admonished, in superstition and worshipping of idols it excelled all Asia, which nevertheless was most corrupt. Pergamos was the most ancient and famous city of Asia or of Mysia and Phrygia, renowned by king Attalus and Eumenes. For the same was the princely palace of king Attalus, which came into the hands of the Romans by the legacy of kings, who were most addicted to idolatry. Strabo speaks much hereof in the 13th book. Moreover this place was also, as Pliny seems to signify in the 5th book, the 30th Chapter, most noble and frequented, by reason the lieutenant or governor there inhabited, who at the commandment of the emperor Domitian, persecuted the true faith of Christ, imprisoning, scourging and afflicting all that professed Christ. By good reason therefore is Pergamos called the seat or throne of the devil. For he is a liar, and the father of lying, and a murderer from the beginning: which the Lord also testifies in the 8th of John. For because therefore at Pergamos reigned heathens, lying, idolatry, superstition, the oppression and murder of good men, it is rightly called the seat or throne of the devil. This appears to be a slander not to be dissembled or suffered. For Rome seemed to herself established forever, and which the gods favored, who had sent them victory over most great nations, and given the empire of the whole world: in the which city justice and religion might seem to be observed. And therefore that this seat of justice and religion should be called the seat of Satan, might be thought both blasphemy and treason. But this doth the only begotten son of God from the right hand of his father pronounce against Rome, against Pergamos, and against all the supporters of Rome. Who shall accuse him of temerity, of rashness, or of bitter speaking? Light persons are doubtless angry, and very sharp wits will be offended, in case they be called by their own names, and be called as they are in deed. For such is the glory of virtue, that all men covet the same even the open enemies of virtue, so that no man will seem to be void of virtue: and such is the corruption and darkness of man’s mind, that he would be that he is not, and would not be that he is. Thereof comes all this impatience in the whole world: when a mattock is called a mattock, and a fig a fig, as the proverb is. Is an harlot therefore no harlot, because she will not be called an harlot? Yes verily is she an harlot, and a shameful harlot, and though she deny never so oft that she is a whore, yet is she an whore nevertheless, and remains a whore. So the seat or throne of Satan is at this day Rome itself, which will seem to be the seat of Christ and the seat Apostolical. For the work and instruction of the devil therein abounds. Finally all cities, towns, and places, wherein truth, godliness, religion and virtue are exiled, wherein the preaching of God’s truth and correction of most corrupt manners have no place, wherein filthiness and uncleanness, bawdy songs and not spiritual Psalms, wherein craft and deceit, surfeiting, murder, adultery, oppression of good people and of godly religion triumphs, be the seats of Satan, however they be called the most Christian and Catholic cities, and worshippers of the right and Christian faith. This thing Jesus Christ the very son of God says, cries, affirms, repeats, and even with a majesty pronounces. For by and by after the murder of Antipas, he adds: where Satan dwells. And these things are doubtless true, which Christ says and pronounces in the Church: and most false be the things which this most sinful world here alleges against the words of Christ.

### Paragraph 198

But the Lord greatly praises this very thing, that in so precarious and unfortunate a place they have stood steadfastly until now, and could not be subdued in the very stronghold of Satan. [note 10.12] Here we learn that it is lawful, as occasion serves, to dwell among a perverse nation; yet so that we are not made conformed to them in any way, either in manners or superstition. And since it is dangerous to dwell among the ungodly, as it were to touch pitch with our hands, you shall commit nothing offensive against the Lord if you go to a safer place, where there is less danger and more opportunity for all godliness. Indeed, when you may conveniently pass to such places, you risk perilously upon rocky ground, where you may chance at last to suffer shipwreck.

### Paragraph 199

And he allows chiefly two things in this church, first that they hold the name of Christ. For the Greek word is not to be touched lightly, but to hold fast, so that it cannot with force be plucked away that you hold. And so they held Christ most deeply fixed in their minds. The name of Christ is the wholesome working of our redemption and sanctification, besides which there is no other name, as Saint Peter says, whereby we may be saved. They cleaved therefore unto Christ, as we read of the apostles in John 6. And it is necessary that every one of us hold fast the mystery of salvation rooted in our hearts.

Secondly, it is not enough to retain the mystery of salvation in our heart, unless we profess it also with full and open mouth. Whereupon he adds straightway, “I have not denied my faith.” Behold how he calls it faith now, which of late he called the name of Christ. And he calls it properly his faith, that is, not devised or invented by men, but set forth by Christ himself by the word of his truth. This true, right, and catholic faith must we confess and not deny, and profess it expressly as well in words as in works.

### Paragraph 200

Christ and his Gospel are denied in many ways. They are denied by silence, when we remain silent when we ought to speak primarily for the glory of God. Christ is denied through dissimulation, as when Peter says, “I do not know what you are saying.” For he knew very well what the servant girl said; but fear caused him to dissemble. He is denied when Christ and his truth are plainly and with explicit words denied. He is denied with a figurative confession, when in deed we confess something, but so darkly and so vaguely that it is unknown what we profess. He is denied when we pretend in our heart that we keep the true doctrine, and deny it in our works, by bowing before idols, attending profane churches, participating in the ceremonies of Antichrist, kneeling on the ground and worshipping that which our conscience gives us, and the faith taught by the apostles instructs us is not God. And truly, all this denial arises from fear and our corrupt affections. If there were as assuredly a reward proposed by men for confessing him, as you are greatly afraid to suffer pain if you do confess, there would seem no difficulty at all in sincerely professing Christ. Wherefore, when you deny or dissemble, you do so from fear. But such timorous and fearful deniers the Lord shuts out of his kingdom. Therefore, the world being despised, the name of the Lord must be confessed boldly and without fear, according to the doctrine of Christ, Matthew 10; Mark 8.

### Paragraph 201

This declaration of faith from the congregation of Pergamum is amplified and highly commended because of the time.[note 10.16] For it is a great thing to profess Christ not in peace, but in very troubled times. It is clear that the church of Pergamum confessed Christ in the midst of persecution, during which the holy martyr Antipas was executed. Therefore, the profession was noble. It is commonly said, but these men saw Antipas slain, and yet could not be deterred from the true faith. These things are indeed expressed briefly, but carry a profound meaning for all churches to follow. Some others read it differently in my days. But the Complutensian copy is better, which reads, “in the days wherein Antipas, &c.” As though he should say, “And you have confessed my name in those days in which Antipas was my faithful witness, who for the same reason was also slain.”

### Paragraph 202

Antipas is commended,[note 10.17] and as it were canonized by the very Son of God. He is praised as a witness, that is, a martyr: and indeed a faithful witness, by testifying, teaching, confessing, and keeping his faith to the Lord, even to the end. Acts 13. Perhaps he was pastor of this Church, or some other man of singular constancy among the faithful. Certainly faith, not torment, makes martyrs. And because this martyr is praised by Christ, we understand that the agonies and conflicts of martyrs should be preached in the church of Christ, and many be excited and exhorted to follow their steps. Therefore we affirm that the holy martyrs of God are honored, but not to be worshipped or called upon. We condemn all those that speak against holy martyrs, and associate them with those that slew them. But concerning the worshipping of Saints, I have spoken elsewhere more at length; we learn hereof also, that they do not die forever, those who die in this world for the name of Christ: neither that the martyrs be polluted with worldly reproach, considering how they are commended by the mouth of God. To Christ therefore, king of martyrs, be honor, praise, and glory world without end. Amen.

## Sermon 11: ¶ The latter part of the third Epistle is expounded, wherein is spoken of the Nicolaitans, which are damned. And exhortation is made to repentance.

### Paragraph 203

.

### Paragraph 204

But I have a few things against you: [note 11.1] that you hold there, that you maintain the doctrine of Balaam, which he taught in Balak, to put occasion of sin before the children of Israel, that they should eat of meat dedicated to idols, and commit fornication. Even so, I hate those who maintain the doctrine of the Nicolaitans. Repent, or else I will come to you shortly, and will fight against them with the sword of my mouth. Let him that has ears, hear what the Spirit says to congregations. To him that overcomes, I will give to eat manna that is hidden, and give him a white stone, and in the stone a new name write, which no man knows saving he that receives it.

### Paragraph 205

In the first part of this letter, the Lord praises many things within the church of Pergamum,[note 11.2] in the second part he will rebuke a few. And he says a few things, not that the error of the Nicolaitans is a minor offense, but that the sin is in others rather than in the true church itself: namely, in those who, though not truly members of the church’s body, outwardly joined with the church and wished to be considered members of it. After this, he speaks gently, lest by excessively exasperating the sin and error among the faithful, he should trouble their minds and discourage them utterly. There is a measure in all things, as the common saying is. And if in a church so commendable, something is found worthy of rebuke, what shall we say of those that are less commendable?[note 11.3] Indeed, why should we not always see in all churches something to be blamed: not so much because the saints are always troubled with the weakness of the flesh, as because hypocrites and corrupt persons always join themselves to the church of God, such as were here the Nicolaitans, and as Judas the thief and traitor was among the apostles. In Christ, the church is without any spot or wrinkle, as the Lord says in John 13. And in the coming age it shall be most fully made perfect, which Saint Augustine also affirms.

### Paragraph 206

And the Lord Jesus rebukes in the church of Pergamum, not that they uphold the Nicolaitan or Galaamitical doctrine, but that they harbor individuals who do. They were therefore culpable because they did not detest the Nicolaitans as much as the Ephesians did:[note 11.4] of whom we heard in the first epistle, that they could not tolerate the wicked. Wherefore, lest the leaven of error should spread further through the entire mass of dough, the old leaven must be purged. It must be tested whether you favor or adhere to heresies. Furthermore, the Lord requires that we should not nurture them, but that we should pursue them with a holy hatred. Of this there is speaking in the first epistle.

### Paragraph 207

Moreover, he describes the heresy of the Nicolaitans to the intent we may see, wherefore he blames it, wherefore he condemns it, and wherefore it ought to be hated. And he describes it carefully by the example of Scripture, that modest concerns or sensitivities might not be hurt or offended. I told you before how the ancient writers report the Nicolaitans to be most filthy. But all things are most aptly and carefully declared of Christ. They are taken from the 22nd, 23rd, 24th, and 25th chapters of the fourth book of Moses, called Numbers. He calls the Nicolaitan doctrine the doctrine of Balaam, and that by a comparison. In Balaam the soothsayer, these wicked acts are manifest; from which it may easily appear what sort of doctrine he had. First, he took the reward or price of iniquity, as St. Peter terms it, and would curse those whom God has blessed, doing clean contrary to his own mind. Secondly, he gives the king most pestilent counsel, which Scripture therefore calls a slander or offense. For he taught the king a way whereby he might entice the people of God into certain destruction, into the most unclean feeding of meats offered to idols, and into most filthy whoredom. All this shall be counted the doctrine of Balaam, which, in hope of filthy lucre and against God’s word and his own conscience, teaches idolatry, unclean eating, and fornication; or reproves not, but counsels rather, when he knows the thing to be filthy. Even so did the Nicolaitans, in speaking evil of the truth and of Christian purity, gave naughty counsel to many, that they should be partakers of meats offered up to idols, and couple with harlots, as in the first epistle I declared more at large.

### Paragraph 208

Here we observe, by the example of our Savior Christ, how heresies should be refuted, not with angry or abusive words, but rather by passages and examples from holy Scripture. As is fitting here, the heresy of the Nicolaitans is condemned. And once condemned by the Lord, it remains condemned forever; we need no new deliberations to condemn impurity. Again, even if all the councils in the world decree the contrary, this remains true and certain, which the Lord Christ pronounces here: accursed be he who decides otherwise.

### Paragraph 209

And here it seems good now to consider whether the Balaamitic and Nicolaitane doctrine in the church is completely extinguished. The names of Balaam and Nicolaitans we certainly abhor, but the thing itself, as well in the affairs of spiritual as temporal rulers, is openly present. For there are men in high authority, in various kinds of learning rightly excellent, most expert in the laws both of God and men, who nevertheless, blinded by the reward of iniquity, curse both the persons and things which they know that God has blessed. Of these Saint Peter also made mention in the second chapter of the second Epistle. The same do suggest evil counsels to Kings and Princes, tending to the destruction both of the preaching of the Gospel and the safeguard of the Church. The same being given to idolatry, and drowned in fleshly pleasures, eat of the sacrifices of the dead, and even feed of idol offerings, and in fornications run at riot. Consider, I pray you, what is the life and sustenance of most popish priests, what opinion they have of holy matrimony, and how much they abhor adultery and whoredom. They dare be bold to condemn matrimony and to judge whoredom better, so that they may enjoy the sacrifices of the dead and in many ways take their pleasure. If any, for the avoiding of whoredom, be joined in lawful matrimony, he is considered unworthy to sacrifice or to come at the altar; but whoremongers are admitted thick and threefold. And all they for the most part are the most beastly bondslaves of the belly, of whom you may believe that the holy Apostle of Christ, Saint Paul, has spoken: whose God is the belly, and glory in reproach of them that seek earthly things. And who will not acknowledge these and affirm them to be very Nicolaitans, maintaining the doctrine of Balaam the enchanter? Among the secular [? populace] you shall find men of all sorts who value the doctrine of Balaam and the wantonness of Zuryk more than they do modesty, gravity, and Christian sincerity. They love the liberty and wantonness of the flesh. They will not have youth and free people to be restrained by virtuous laws. They will even at this day banquet and mask with the maidens of Madian, and follow their fleshly lust. For they maintain surfeiting, drunkenness, and whoredom. And these are also very Nicolaitans, and have neither few nor obscure adherents to favor their sect. They want not their worldly reasons, many and great, to maintain the same.

### Paragraph 210

[note 11.9]But let us hear what Christ himself, sitting at the right hand of his Father, judges of them. Those or that same which these men think, teach, and do, I hate, says the Lord. What thing can be spoken more grievously than that God hates the doctrine of the Nicolaitans? For the whole scripture of both Testaments condemns this Nicolaitism.

### Paragraph 211

[note 11.10] After this description and reprehension of the Nicolaites, he proceeds, as in the former epistles, to exhort them to amendment or repentance. For when he says, “repent,” he understands or comprehends all penance or repentance. We said that to be a conversion unto God, whereby we amend evil things for good, relinquishing that which is evil and in its stead placing that which is good: and that through faith in sincere love and fear of God. You shall amend, therefore, if you abstain from meats offered up to idols and from fornication, and receive the true religion of Christ instituted, and possess your body in honor, not in the lust of cupidity, as Paul says in 1 Thessalonians 4. The church of Pergamum repeated, that if they did not disable or wink at the filthiness of the Nicolaites, but stoutly withstood it. The Nicolaites repented, if laying aside their filthiness, they received again the purity of faith and life. And to all and singular is said, “repent.”

### Paragraph 212

The Lord now also drives them to repentance with severe threats: “Unless you amend,” he says, “I will come to you shortly,” a manner of speaking that has been discussed before. He adds, “And I will fight with you with the sword of my mouth.” With whom? With the impenitent, and especially with the Nicolaitans. He did not threaten utter destruction or devastation to the church, where there was great hope that the old leaven would be purged, but he threatens the impenitent people. Just as a judge, magistrate, or soldier uses the sword, so does Christ his word. And the word, indeed, wounds or slays no one, but it reveals God’s word, so does the execution of God’s power. Therefore, Christ, even as he reveals with his word, reveals that he will judge idolaters, those who worship false gods, swine, dogs, and adulterers, and not only judge them, but punish them. And as he threatens, he does.

Thus he fights with the sword of his mouth. We have an example in the Israelites, of whom 25,000 men were destroyed for following the doctrine of Balaam. Afterward, the Moabites and Midianites were also destroyed, nor were the corrupt women spared. Moses discusses this at length in Numbers 31. We also see at this day the sword of God going through the world, overthrowing these now, and then, for no other causes than those for which the Lord slew and destroyed Balaam and his adherents. Therefore, let us fear the Lord and walk in his commandments, for he will strike far more grievously with his sword when he pronounces judgment, “Go, you cursed, into everlasting fire.” Matthew 25. He does not expressly say, “I will cut you with the sword of my mouth.” For we are often severed and cut by the word of God to our great profit, discipline, and amendment. At present, he says he will fight: behold, he will fight, namely, against his enemies. Therefore, he threatens destruction. And we doubt nothing but that the impenitent of those times and of all times shall be destroyed, for (as I said even now) we lack no examples at this day.

### Paragraph 213

Again, so that this notable and wholesome doctrine might not seem to belong only to a few men of Pergamum,[note 11.13], he applies it to all churches throughout the world. We have discussed this application on one or two occasions in the previous letters.

### Paragraph 214

Finally, in his manner, intending to move us all more strongly to repentance and obedience, he proposes a most ample promise: and that to those who strive and overcome the flesh, the world, and the devil, not to idlers, nor to those who wallow in wickedness. We are therefore encouraged by that promise, which is of three kinds. First, he promises to those who fight bravely and overcome, and do their duty, manna, and that secret or hidden [illegible]. That external manna, known to all men, is not the true manna. For the ungrateful Israelites despised it, considering it a very light food, and preferring the meat pots of Egypt, full of meat, onions, leeks, and garlic, that they might eat their fill. They do not see the celestial manna figured by this outward manna, yielding all sweetness and spiritual pleasure. The faithful see that this hidden manna is Christ, as he himself explains in John 6. Christ, therefore, gives himself to those who overcome, gives himself in food, which truly nourishes. He who has once, with true faith, tasted Christ will desire no other food to be given to him. For in Christ he has all things, in Christ he is complete, and with all good things fully satisfied. Oh, that our subtle disputers understood these things, they would reason about nothing at all concerning the merits and intercession of saints and such other things, about which, while they reason after their customary manner, they declare that they have not yet tasted how good and sweet the Lord is.

### Paragraph 215

After he promises to give to the victors a white stone, that is, absolution and remission of all sins, and that without any doubt. For Christ truly absolves us from our sins and from the punishment due for them, and from condemnation. And he referred to the custom of ancient judges, who condemned with black stones and acquitted with white. For these verses are well known in Ovid’s *Metamorphoses*, book 15. The practice long ago was to condemn with black stones and acquit with a white stone, and here we give warning that the remission of sins is not granted to people living for their works or merit, but that faith is the victory that overcomes the world. Which John himself testifies. And that faith indeed fights bravely in our hearts, but in the meantime acknowledges in all things the grace of God, neither does it nullify the merit of Christ. For as it is not sluggish, so is it also fearful.

### Paragraph 216

[note 11.16]Finally, he promises that he will write a new name on the stone, and that it is known only to him who possesses it. Christ will not only grant us remission of our sins, but also the glory and unspeakable joy of his heavenly communion. Isaiah and other prophets have made mention of this new name. Conquerors had famous names. If we overcome, we shall enjoy celestial glory. That glory is so immeasurable that it can only be perceived by feeling, not by speaking. For whatever you may say, however great, famous, or excellent, something greater will be given to the overcomers. For the Apostle Paul cites from Isaiah: What no eye has seen, nor ear heard, God has prepared for those who love him. And even in this present world, we are given a quiet conscience and unspeakable joy, which those who experience it truly feel. Those who have not tasted it cannot believe how great it truly is. Therefore, Paul says, “And the peace of God which surpasses all understanding,” etc. May our Savior Christ grant us such minds. Amen.

## Sermon 12: ¶ The Epistle of Thyatirena is expounded, wherein are sundry virtues commended, and the vice of Ieſabell reprehended.

### Paragraph 217

.

### Paragraph 218

And unto the messenger of the congregation of Thyatira write. This says the Son of God, who has eyes like a flame of fire, and whose feet are like brass. I know your works and your love, service, and faith, and your patience, and your deeds, which are more at the last than at the first. Nevertheless, I have a few things against you, that you allow Jezebel, who calls herself a prophetess, to teach and deceive my servants, to lead them into fornication, and to eat foods offered up to idols.

### Paragraph 219

The fourth epistle written to the Thyatireans,[note 12.1] is more plentiful than the rest, and with manifold fruits replenished. For it commends and praises in that church excellent virtues, and singular gifts not a few. Straightforwardly he reproves them for suffering too gently the Jezebelism, which he describes and shows how filthy it is. He threatens them severely, unless with perfect repentance, they amend their sins and wickedness. Furthermore, he warns them that they should expect no new revelations, but that they persevere and abide in those which they had learned hitherto, and in which they now are. Hither also with most large promises he allures them, and finally communicates and commends this doctrine to all churches. And there is a wonderful likeness and correspondence in all epistles, as may be seen also in all the books of the prophets, in the story of the evangelists,[note 12.2] and in Paul’s epistles. Whereof it may easily be gathered that the doctrine of the vexation is most absolute, perfect, and plain, and agreeable to itself in all things. In so much that if all the writings of all other Apostles and Prophets did remain, we should have had no more in those many and most plentiful books than we now have in the holy Bible. God provided well for us and for our infirmity by this brief way. Here are seven Epistles set in the 2nd Chapter: but it is marvelous to see how like they be all, teaching in a manner all one thing.

### Paragraph 220

This fourth [section] is chiefly profitable for those congregations which are sound in the purity of doctrine and are pure moreover in holiness of life, but do not with a fervent zeal enough persecute open heresies. There are other fruits and benefits, which we shall speak of in order. But just as in all the other epistles that go before, first is set forth to whom the epistle is sent and from whom it comes: so also in this epistle to Thyatira, both the superscription, as they term it, and the subscriptio are expressly set. It is sent to the messenger of the church of Thyatira, and so to the whole church, as I have told you before often. And Thyatira is a noble and famous city of Lydia, in Asia, on the river Hermus: where we read that the woman was born, who sold purple, which was converted to Christ by Saint Paul in the 16th chapter of the Acts. It was a populous city and much frequented, so that it is no marvel though many diverse, unclean, curious, and heretics did associate and join themselves to the church of God. Geographers write many things of that famous city of Asia.

### Paragraph 221

And the author of the epistle is the Lord Christ himself, the high king and Bishop, who uses the Apostle’s pen, or blessed John for his scribe or secretary, by whom he will have those things published throughout the whole world. And he establishes the epistle’s authority, while repeating certain members of the former image and description, he shows himself in such a way to be seen of the church, to be viewed in faith, that they help the matter wonderfully. He sees here heresies and the secrets of hearts, and treads under his most pure and clean feet whatever advances itself against God’s glory and truth.

### Paragraph 222

He calls himself therefore the Son of God,[note 12.7] whom before we heard to be the Son of Man. He is therefore and remains both, even in glory, as well the Son of God as man. In the divine nature, of the same substance with the Father, in the human nature, communicating with us in all things except sin, the other nature is not swallowed up in glory, but two distinct and several natures without any mixture abide in one person undivided: which indeed is one Christ, very God and very man, to be worshipped world without end. Hereof we have testimonies in Luke 1, 1 John 1, and Romans 1. And which of the heretics or persecutors will make war with the living Son of God?

### Paragraph 223

After he attributes to himself eyes, casting out fire and flame. For nothing escapes the knowledge and judgment of Christ our Judge; he beholds the reins and hearts. Moreover, he lights some, and some he commits to everlasting fire, therein to burn forever. Now then, if any do imagine with themselves that they can hide heresies and malice in their hearts, they are deceived. For in the eyes of Christ, darkness itself is light also. The same Lord has also fetched most purged and clean ones; he treads down all ungodliness. And wherever he walks with his shining feet of brass, he consumes immediately all heresies and corrupt life. Therefore, this prelate, most pure, and most fit and apt to purge, finally best furnished to bolt out the secrets of hearts, shows to the congregations these things that follow: he himself walks and is conversant in the midst of the church, both King and Priest.

### Paragraph 224

And just as he has testified in all his letters, that he knew the works of that church, so he repeats it here as well, so that we should never fall into wicked complacency, as though the almighty and all-knowing God did not know us and all that belongs to us—a matter about which I have spoken sufficiently before.

### Paragraph 225

Now he clearly sets forth every work of this congregation and commends five most notable gifts or brightest virtues. First, Charity, which encompasses the love of God and our neighbor: whereby it is brought to pass that we prefer nothing in the world before God, neither harm our neighbor, but rather heap upon him all duties and benefits. This we owe to God and all our brethren in the congregation. Of Charity is spoken elsewhere most abundantly, as in the gospel and epistle of Saint John. Secondly, he praises Diaconian, that is, the Ministry. The which may be explained two ways. For either he understands, as Arethas supposes, ministries toward the poor and needy, that is to wit, duties and pains taken about the poor, by leading, relieving, succoring, speaking faithfully in their cause, in giving them meat, drink, clothing, and visiting them. For so this word Diaconia is used in the second epistle to the Corinthians, and so forth. Or else he means the ministry of the word, by which in teaching, exhorting, comforting, and rebuking, we advance greatly God’s glory and the health of souls. The Thyatirenians were doubtless diligent in either of both. And accuse us grievously, which addict ourselves to our own affairs, do neglect our poor brethren: who finally make the ministry of God’s word odious, by our railing and slandering, especially with them that be ignorant as yet, and have heard nothing of God’s word.

### Paragraph 226

He also commends faith among the Thyatireans. Thomas Aquinas in his commentary on this book admonishes that faith does not come from charity, because it is presented here first; but that charity and good works spring from faith. John has recited charity before faith, because faith derives its value from charity and works. Nevertheless, it seems here that “faith” is not so much to be taken as trust in God, but rather as fidelity, truth, and promises kept. For faithfulness beautifies all other gifts. Suppose you have servants, male and female, who are fortunate in doing their tasks, but are at the same time untrustworthy, fickle, and deceitful: what benefit would there be in their being equipped with various gifts? Imagine again that a preacher or senator is not so furnished with wisdom and experience, but is nevertheless faithful, and with all his heart strives to do all things uprightly and to favor the just cause: will not fidelity here supply his lack? Great therefore is faith, that is, fidelity and truth. Not without cause the Apostle required this of ministers in the 14th chapter of the first Epistle to the Corinthians, saying: That which is chiefly required of stewards is that a man be found trustworthy. This faith is also required of us at this day. This faith, good brethren, is rare, and therefore evils have overflowed everywhere. Let us heartily pray to the Lord that he will grant it to us, and that we may expel from our hearts unfaithfulness and deceitfulness.

### Paragraph 227

Hereunto is added patience, which was also praised in the earlier churches. This is a necessary virtue. For impatience causes us to murmur and complain against God, and prevents us from standing firm in the confession of faith, while we refuse to suffer patiently the things that the enemies of faith threaten to inflict upon us. But why do you defile yourself with theft? Why do you rush into the wars of a foreign prince? Why do you practice usury and lewdness? Because you lack patience in your poverty, which you seek to relieve with wicked deeds.

### Paragraph 228

To be brief, the Lord now recounts all kinds of good works [note 12.13]. In this, he chiefly praises that they often surpassed themselves, in doing more and greater things. And this is worthy praise. For the husbandman, that is to say, the faithful father, prunes and cuts the vines, so that they may bear more plentiful fruit. It does not become the godly to remain stagnant, but to proceed in godliness.

### Paragraph 229

It is most shameful to become worse the longer one endures. As a finger, the longer it is, the less it is. This is what is pointed out to children in schools who learn nothing. Let us be ashamed of our sluggishness. Let me say, weigh these things carefully in our minds, and consider often that God allows them, requires them, and that they are the true seals of the faithful walking in truth, and of those who boast of faith only as a vain name without the substance. If you feel yourself not to be utterly void of these gifts, praise God, and know that none of all these things is of yourself, but of grace. And pray for the increase of these gifts. If you are destitute of these virtues, mourn and lament before the Lord, humbly ask him for forgiveness, and require the abundance of God’s gifts.

### Paragraph 230

Secondly, he rebukes certain things within that same congregation, namely that they allowed Jezebel to teach, and so forth. He calls this a minor matter, not that Jezebel’s doctrine is inherently insignificant, but because, though it is found among others, the church tolerated it gently—that is, did not punish it with greater severity. Regarding this phrasing, I have spoken of it previously. We do not condone Jezebel’s shameful acts, nor do we approve of them. But when we might have imposed a more severe punishment, we permitted them to abound and increase. Although, therefore, we possess many good gifts, the Lord is displeased with us for allowing ungodliness to reign.

### Paragraph 231

But if the Lord should condemn that very allowance, how much more blameworthy should we consider the wickedness itself to be, namely, the Jezebelic practice? And how base and filthy it is, I will briefly declare.

### Paragraph 232

[note 12.16]Just as he previously refuted the Nicolaitans using the example of Balaam from Scripture, now he brings forth the example of Jezebel to refute the Cataphrygians, or Montanists. Arethas understands the entire passage concerning the Nicolaitans, which I do not dare to agree with, given the whole composition of the epistle. I grant that the Montanists were partakers in filthiness with the Nicolaitans. But Jezebel possesses a peculiar characteristic of her own.

### Paragraph 233

Iezabel, as sacred history testifies in the third and fourth books of Kings, chapters 16 and 17 and following, was the daughter of Hethbahal, king of Sidon, who, when married to Ahab, introduced the worship of Baal into the kingdom of Israel, building a magnificent temple in Samaria and establishing a large college of Baal’s priests. For Elijah is recorded to have slain 450 Baalites, even of the kings’ chaplains, as it were canons or prebendaries, and 400 ministers or country chaplains who served in hills, woods, and groves. This woman, therefore, founded this religion and sought to govern the prophesying at her pleasure. For pursuing Elijah intensely, she slew many of the prophets, truly because they would not teach according to the woman’s desire. Moreover, through Baal’s religion, whoredom and all uncleanness increased. King Jehu objects to King Joram, his son, the whoredoms of his mother. So Iezabel also augmented the eating of meats offered unto idols and all idolatry throughout the whole kingdom. Even at that time, when the Lord in a solemn sacrifice by miracle on Mount Carmel through the ministry of Elijah, had declared to that whole realm that the religion of Baal was most vain and false, and that the religion of the only God of Israel was most sincere and true; for Iezabel nevertheless persecuted the truth and established falsehood. Yea, moreover, she took upon herself governance in civil matters.

### Paragraph 234

She seized the king’s seal and forged letters, sending them in the king’s name to have Naboth put to death, a truly good and innocent man. Such was Jezebel’s wickedness.

### Paragraph 235

Now following the example of that defiled woman,[note 12.18] there were women in the church of Thyatira who claimed authority in religion and teaching within the congregation, presuming to have the spirit of prophecy. They taught in fact, but corrupt doctrine, seducing those whom God by his doctrine had prepared to be his servants. But these false prophetesses corrupted their minds, and brought forth a new doctrine and prophecy, and many things not found in the scriptures, but drawn from their own devilish dreams and deceitfulness. And among other things, they communed with the Nicolaitans, in word, and participated in meals offered to idols, as has been spoken of before. And the Lord seems plainly to speak of the Cataphrygians or Montanists, whose foundation was laid in the time of Saint John, but after, in the course of time, and especially in the empire of Antoninus, sixty years after the Apocalypse was published, broke out more strongly and abundantly. They say how Montanus had prophetesses Priscilla and Maximilla, who had visions and brought in wonderful revelations into the church. Eusebius treats of them at length in the fifth book of ecclesiastical history, chapter 16. And Epiphanius in the 48th heresy, in Panarion. Certainly, John, or Christ himself by John, seeking at the very beginning to uproot and destroy the roots of this heresy, by the example of that wicked woman Jezebel, has condemned that same heresy. Scripture also elsewhere prohibits a woman from ruling, teaching, or ministering in the congregation. Immediately, the Lord himself will refute the new prophecies, when he admonishes us that he will reveal no other new kind of doctrine besides that which he has committed or delivered to his church. Now also fornication, and the eating of meats offered to idols, are condemned elsewhere in scripture most severely, as before is said.

### Paragraph 236

But since those things so greatly afflicted and troubled the church of God in the time of the Apostles, it is not difficult to see how rash they are who, at this day (as I showed you before), for the sake of hating the true religion restored, accuse it of sects, which abound so plentifully, as though that corruption proved that the gospel we preach were not the gospel. For the gospel that was preached by John and the rest of the Apostles was the truest and purest gospel, even though false gospel proponents like the Nicolaitans, Cataphrygians, and other sects innumerable crept in. Nevertheless, the gospel refutes and condemns all such sects, and maintains the Christian truth and unity of the catholic church. Praise be to our God, the Lord. Amen.

## Sermon 13: ¶ The Lord threatens sore the impenitent, as he that renders to every man after his works.

### Paragraph 237

.

### Paragraph 238

And I gave her opportunity to repent of her fornication, but she did not repent. Behold, I will cast her into a bed, and those who commit fornication with her, into great adversity, unless they turn from their deeds. And I will put her children to death. And all congregations shall know that I am he who searches the minds and hearts, and I will give to each of you according to your works.

### Paragraph 239

To the former errors and sins of Jezebel he adds another sin, nothing slight, namely, the abuse and even the contempt of God’s long-suffering. God does not immediately destroy those who are in error and sin, even the most grievous. But sinners are wont for the most part to abuse God’s long patience into an opportunity and pretext for sinning more openly, saying: “If God did so much abhor these offenses, he would have destroyed us long ago. But now he nourishes us benignly, therefore he does not greatly dislike it.” But this is an abuse of God’s long-suffering. For the Lord says at this present, “I have given Jezebel a time to repent, to leave her fornication, and to turn to the Lord.” However, she has not converted. Which thing the Lord takes in most evil part, that his grace should be truly despised and set at naught. Therefore, Paul to the Romans: “Do you despise the riches of God, his goodness, long-suffering, and lenity, not knowing that the goodness of God provokes you to repentance?” And so forth. If then the Lord has not suddenly oppressed us in our sins, let us not take that as license to sin, but let us rather amend. Peter says, “The Lord is patient toward us, while he will destroy none, but receive all to repentance.” 2 Peter 3. Certainly, Jezebel herself, when after the death of her husband Ahab, and the mortal fall of her son Ocozias, she did not amend, nor within the twelve years of her son Joram, wherein he is said to have reigned, did repeat her [illegible]—felt the wrath of God so much more grievous, for that it was long before it came.

### Paragraph 240

And in the text following, the Lord Jesus in truth threatens the followers of Jezebel most severely—that is to say, the Cataphrigians or Montanists—unless they repent in time. For he opens again the gates of his grace to the penitent, recounting how he will punish the impenitent. By this, he truly seeks to drive them to repentance through threats. For in recounting the kinds or degrees of punishments, he also shows diverse kinds of those who are in error, and declares to each his judgment, which they may avoid through repentance. And it is thought that he rehearsed those kinds for this consideration, lest any man should perhaps think himself blameless and free, even if he is but a small partaker with Jezebel.

### Paragraph 241

First, the Lord threatens Jezebel herself, that he will cast her into a bed. He speaks of the first authors of the evil and of the doctrine, upon whom he menaces to send a sickness. For the bed in many times in the scripture is taken for the very diseases with which those lying in bed are afflicted. And we Germans say that he is taken with a most grievous and deadly disease. And the Lord afflicts the archheretics with sickness of body and soul. In the meantime also he weakens the force of their terror, to the intent it might by little and little vanish away.

### Paragraph 242

Secondly, he threatens great affliction to those who have dealings with Jezebel: that is to say, those who cleave to false doctrine, receive errors, delight in heresies, and seek to promote them. To these, he threatens most grievous afflictions, namely of body and soul, both in this present life and in the life to come. He seems to have said somewhat more than if he had merely recited certain kinds of punishment.

### Paragraph 243

Finally, he threatens death to the children born of this union and fornication, namely, illegitimate sons and bastards. These are chiefly the children of heretics, who stir up new heresies and revive those already condemned, weakened, and fading away. The Lord destroyed these with temporal and eternal death. And ecclesiastical history testifies that the Lord has indeed severely punished not only the heresy of the Cataphrigias but all heresies in general. Certain things concerning the Cataphrigias or Montanists are mentioned by Eusebius, book 1 of the ecclesiastical history, chapter 16.

### Paragraph 244

The Lord seems to me here to have alluded to the old story of Jezebel and Ahab for them, as it were cast in a bed, from day to day, ever since they began to worship Baal; he vexed them with sickness and brought them low. The people who received the religion of Baal were put to much sorrow, evils, and afflictions. Finally, their children he brought to a shameful death. Their partakers also were slain, those who would have had Baal’s religion safe and sound, and even restored again. For after the death of Ahab his father, not many days after, Jehoaz the son of Ahab and Jezebel, bruised with an unhappy fall and cast in bed, died (2 Kings 1:4). And Jehoram, another son of Ahab and Jezebel, was slain, struck through with an arrow by Jehu (2 Kings 9). Athalia, the daughter of Ahab and Jezebel, the wife of Jehoram king of Judah, the son of Jehoshaphat, being ensnared by the oath of Jehoiada, fell down before the gates of the temple (2 Kings 11). And Jehoaz king of Judah, the son of Athalia and Jehoram, was slain also by the power of Jehu. And after were put to death by the same Jehu, seventy sons of Ahab. And all the priests of Baal were slain together in the temple, and before the altar of Baal, and not one of so great a number escaped. Yea, the temple, the idol, and the service of Baal were quite and clean overthrown. This old, marvelous, and wonderful history the Lord calls to memory, signifying that he lives yet a revenger and a punisher, who will neither overpass the just limit nor touch it out of time. For he adds, “And all congregations shall know,” etc.

### Paragraph 245

It is important and worthy of remembrance, and full of comfort, that in this recounting of punishments, he inserts a mention of repentance as though to say, let no one think that he must be destroyed and perish through a certain fatal necessity. For if anyone will repent, the gates of God’s grace are opened, his sins will be forgiven, and he will be received into favor and delivered from all those evils. And in this manner the Prophets have also taught, Jeremiah in chapter eighteen, and Ezekiel in chapter eighteen.

### Paragraph 246

But since the punishment is not immediately executed upon the unrepentant, some will exclaim that God is asleep, that he sees or hears nothing. Therefore, the Lord himself answers them, saying: “And all congregations shall know…” [note 13.8] For I will surely, at the last, execute my vengeance in due season. For then all men will learn that I neither sleep nor neglect my servants at any time, nor will I allow those who deserve evil of me and of my Church to escape unpunished. Furthermore, Christ testifies that he searches the depths and hearts of all men.

### Paragraph 247

And he means that he knows all thoughts and devices of the heart, finally the appetite itself and all the desires of man, so that he can judge them truly, for nothing, however secret, is hidden from Christ. Therefore, he is very God. For it is the property of God, and belongs to him alone, to know the hearts of the children of men, as Solomon testifies in the third book of Kings. Christ therefore sees the secret and filthy works both of the Nicolaitans and all other beastly men, which Paul says are unworthy to come to light or to be expressly declared to men. Ephesians 5.

### Paragraph 248

[note 13.10] Christ does not only know all the thoughts of men, whatever they may be, but also gives to each man according to his own works. And so the Apostle Paul teaches, saying: “The just judgment of God will be revealed, which will reward each man according to his deeds—praise, honor, and immortality to those who continue in good works and seek eternal life, but indignation and wrath, tribulation, and anguish to those who are disobedient and disobey the truth, and follow iniquity.” 2 Romans. For works are the tests of faith and infidelity. And works, whether good or evil, are judged by God and by godly men, according as they proceed from faith or from infidelity. Therefore, whatever any of us sows, that he will also reap. For God is the most just rewarder of good and the revenger of evil. This sentence, as it is most true, is the foundation of the true and godly religion. Glory be to God.

## Sermon 14: ¶ That the doctrine of pity is so fully ſ [illegible][illegible] forth to the church, that there needs no new Revelations. And of the most large promeſſes of Christ made unto the church.

### Paragraph 249

.

### Paragraph 250

And to you I say, and to others who are of Thyatira: Whoever does not have this doctrine, and has not known the devices of Satan, as they say, I will place no other burden upon you than what they already have. Hold fast until I come. And whoever overcomes and keeps my works until the end, to him I will give power over nations, and he shall rule them with a rod of iron. And as a potter breaks the vessels, so will he break them. Even as I received from my father, so I will give it to him—the morning star. Let him who has ears hear what the Spirit says to the congregations.

### Paragraph 251

He now speaks to the Cataphrygians[note 14.1] and also to the faithful of the church of Thyatira who rightly believed in Christ, and heals their diseases. This demonstrates the unspeakable mercy of God, which does not cease to speak to those still entangled with heresy and to heal their pestilential diseases. He admonishes all men that they should not look for new revelations, but rather know that God through Christ and his Apostles has set forth a most perfect doctrine, to which he will add nothing. Therefore, they should keep fast in memory the things they had already learned and in which they were now being exercised[note 14.2]. For the Cataphrygians, also called Montanists, boasted of a new comforter and a new revelation, as though all things had not been fully set forth by the Apostles, but that many things were still to be revealed. As also the maintainers of the Popish church stubbornly affirm today. And just as the Cataphrygians covered their trivialities under the pretense of the Holy Spirit, so too do the Papists cloak the vain constitutions of men and set them forth under a false color of the Holy Spirit, as though the Lord spoke of their decrees when he said, “I have yet many things to say unto you, which now you cannot bear.” Nevertheless, the faithful people of Thyatira, who did not have the doctrine of Jezebel but rather detested it as doubtful, said that the Devil was a certain dependence [?] and had a thousand crafts, which could also transform him into an angel of light. And that they were but simple men, who, being ignorant of his wonderful crafts and subtleties, did not know what they should chiefly follow, while false prophets also make their boast of the Holy Spirit, shine in miracles, and avouch their doctrine to be true. You will find this even today, when people will say, “I am a plain, simple man, and do not know to which side I should cleave, since the doctors of both sides affirm with great cost that they have the truth on their side, and therefore some will say, they shall agree better before I will believe any of them.” &c.

### Paragraph 252

The Lord, therefore, answering both parties, shows what they should do. To you, he says, who follow the doctrine of Jezebel, I say also to the rest of the Thyatirans, who do not follow Jezebel’s doctrine, yet nevertheless complain in such disagreements and wondrous schemes of the devil, that they do not see what is best: To you all I say, if you are simple in deed, as you pretend, if you will with all your heart embrace the truth, give yourselves to the simplicity of the apostolic faith, cleaving fast to such things as you have once learned from the Apostles, neither seeking nor receiving any new religions, or additions, constitutions, or anything else moreover, than that which you have learned from the Apostles. For these things which you have received are sufficient to obtain salvation.

### Paragraph 253

And these words of the Lord must be weighed more diligently, so that we may perceive the great benefit in them: that is, I will lay upon you no other weight or burden besides what you already have. The Lord affirms that he will add nothing more to the evangelical doctrine set forth by the Apostles, as to that which is most perfect. Certainly, if the doctrine of Moses were so perfect that the Lord himself prohibited any man from adding or taking away anything from it, but only doing what was commanded, as we read in the fourth and twelfth chapters of Deuteronomy. Who would doubt that the doctrine of Christ, the Son of God, would be lacking in anything? He therefore now affirms that he will lay nothing upon them more than he had already laid, and which they were bearing at that time.

### Paragraph 254

A burden in the sermons of the Prophets is taken as doctrine of grave and weighty matters. The Apostles also call the law a yoke and burden. Wherefore, when the Lord says that he will not lay any other burden upon the church, he means that he will not reveal any other doctrine, nor further charge them with other rites or ceremonies than those he had already ordained and imposed. These words of Christ accord very well with what is read in the apostles’ epistle, the Synodical Act of xv. For by the common consent of the congregation, and after the mind of the Holy Spirit, they say they will impose nothing moreover upon the church than such things as they had already received of Paul, and a few things that they added for a declaration of the same. Whereupon Paul said to the Galatians: If an angel from heaven preach unto you another gospel besides that which is preached, let him be accursed.

### Paragraph 255

What then? Hold fast, namely that which you have received, suffering it not to be plucked out of your hands: Hold fast I say, with tooth and nail, until I come; that is to say, unto the last judgment. Therefore he testifies expressly that this doctrine shall be perpetual and unchangeable, and therefore to be kept most steadfastly by all of us, and not to be shaken from, though all the world cry out and persuade the contrary. Arethas, Bishop of Caesarea, required of them nothing else, he says, but that they would keep safely the godly pledge of faith until his coming. This if we shall do, we may easily avoid the crafts of the devil and deceivable clouds. For whatsoever they shall bring forth, whatsoever they shall forge and refine, or disguise with the counterfeited color of the Holy Spirit, we shall always have recourse to the simple doctrine of Christ set forth by the apostles, wherein alone we shall rest, rejecting all things that shall not accord with the same.

### Paragraph 256

And this wholesome doctrine of Christ confounds all traditions, and subverts all constitutions made since the time of the Apostles. The godly may always object this saying of Christ to the traditioners. I will lay no other burden upon you, besides that you have: That same hold fast until the last judgment. They shall allege that same also, that the Apostles deny that they will add anything more. Acts 15. Christ spake this in the time of John, in the year of our Lord 87. Therefore whatever laws, traditions, decrees have been made since that time, we know they were not imposed by Christ, which saith so expressly that he will lay no other burden on the faithful. Where then become the decrees and constitutions of worshipping images in the church, for the consecration and celebrating of masses? What shall we say to the decretals of the Bishop of Rome? They are all overthrown and stricken down as it were with a thunderbolt, by this only sentence of Christ. I will impose no other burden than that you have, keep that until the judgment. Behold, he saith, unto the judgment, lest any should imagine in the meantime that another thing had pleased the Holy Spirit. Let us therefore persevere in the same.

### Paragraph 257

Hereunto he adds, as was his custom, most ample promises, that through hope of so great rewards he might draw them from errors and join them to the true religion. And just as in the former epistles he has said, “He who overcomes,” so here he repeats the same, admonishing us not to sleep but to watch and fight valiantly. He overcomes who keeps the works of Christ until the end. The works of Christ, by a subtle opposition, are set against the inventions and works of men.

The works of Christ signify both doctrine and faith, and whatever good works flow from them, the service of worshiping God, and the observance of God’s word. For in the twenty-eighth chapter of Saint Matthew, the Lord says to his disciples, “Teach them to keep those things which I have commanded you.” He speaks with emphasis, “those things which I have commanded you,” not such as you shall have invented of your own brain. For the Lord quotes from the prophet in the fifteenth chapter of the same Saint Matthew, saying, “In vain do they worship me, teaching the doctrines of men.” Therefore these works have no promise; but the works of Christ, which he himself has ordained, and which are done by his Spirit and by true faith, while we forsake our errors and cleave to the truth, they have a most ample promise.

### Paragraph 258

And promises two notable things. The first: just as my father has promised me victory, and fulfilled it, that I overcome all my enemies, and triumphed over them, shattering them like clay vessels without any difficulty, so will I give unto you also power and victory against all the ungodly. And that same promise shall at the last be fully accomplished in the final judgment, in which all the enemies of godliness shall be cast under the feet of Christ, as it is declared in the Psalms, especially in the second and one hundredth Psalm. And in this world also Christ affirms that his servants shall spiritually rule over his enemies, just as Christ, although he was tormented and died, nevertheless overcame his enemies. The holy and ecclesiastical histories bear witness of these things sufficiently.

### Paragraph 259

The latter: I will give him the morning star. And he understood the knowledge of Christ increasing daily more and more, and so even Christ himself, [note 14.11] in like case as the day in the rising of the morning star waxes brighter and brighter. In this sense the Apostle Saint Peter is said to have used this allegory in the first chapter of his Second Epistle, or at the least he promised a clarity most bright. For Daniel says how the faithful in the resurrection shall shine like the firmament. The which thing also the Lord Christ mentions in Matthew 13. And the Apostle alluding hereunto said that one star was brighter than another: so likewise in the resurrection one shall be made brighter than another. These promises are most great, neither can I think that any greater can be given us. God grant us grace that we may be made partakers of so great things.

### Paragraph 260

Finally, he applies this epistle to all the churches and all ages of the world. Concerning which, since we have spoken of it more than once, there is no need for me to be tedious to anyone by repeating it again. To the Lord our God be praise and glory.

## Sermon 15: ¶ He blames certain things in the congregation of Sardis: notwithstanding he shows straight ways a remedy, whereby they may be healed, and be safe.

### Paragraph 261

.

### Paragraph 262

And write to the messenger of the congregation of Sardis, this says he who has the spirits of God and the seven stars. I know your works: you have a reputation that you live, but you are dead. Be watchful and strengthen what remains, which is ready to die. For I have not found your works perfect before God. Remember therefore how you have received and heard, and hold fast and repent. If you do not watch, I will come on you as a thief, and you will not know the hour when I will come upon you.

### Paragraph 263

In one congregation of Sardis there were two sorts of people, professing on either side the name of Christ. But some in deed answered but little to the holy profession, living more licentiously than became them. And the others in holiness of life set forth the doctrine of our Savior that they professed. The Lord Jesus accuses the first sort in this Epistle by John, and shows also a medicine for the disease. And he exhorts the latter to perseverance, commending their integrity. Therefore this Epistle is divided into two parts, very fit and profitable for our time.

### Paragraph 264

The first part of the Epistle contains those things which we have now recited. He does not proceed in this matter in any other order than we have seen him proceed hitherto. For first, he shows to whom it is dedicated and sent: namely, to the pastor of the congregation of Sardis, and therefore also to the whole church. Sardis is said to have been the head city of Lydia or of Maonia, the metropolitan city of Croesus, the most rich king of Lydia, whom Herodotus writes that King Cyrus overcame, a town most famous, and steeped and stained with pride, that it was a den of vice and addicted to voluptuousness. For Strabo in the thirteenth book of Geography testifies that all the maidens thereof were harlots, who mentions more of the same city. Certainly, it seems to have kept its old ways even at such a time as it had received the name of the Lord, and therefore to have been more given to fornication and all manner of filthy lust. The Lord seems to have blamed this in them, as Saint Paul likewise persecuted the same vice in the Corinthians. The world can hardly believe that simple fornication is a sin, whereupon in that great council of the Apostles, both they and the elders and the whole assembly with one mind decreed that the Gentiles should abstain from fornication. The devil at this day goes about many times to defile the church again with fornication, to set up brothels, and that by authority and openly whoredom might be practiced. For so being cast out, he takes seven worse spirits, enterprising to possess that place again out of which he was exiled by the preaching of the Gospel. We must therefore resist him, lest the Lord Jesus himself do accuse us, as he does here accuse them of Sardis most grievously. Then the Lord Jesus is declared to be author of the Epistle, not without praise. For he is said to have the seven spirits of God, that is to have the seven formed spirit, whom he also pours out upon the faithful, or else he is one only spirit, and not seven: but seven, that is to wit, his graces are many and diverse, as I declared in the first chapter. For he also has in his right hand seven stars, to wit, the whole multitude of all preachers and ministers, keeping and instructing them. And this beginning agrees not amiss with this argument, which he treats in this Epistle. For of the spirit of Christ is life; of the want of the spirit, death. Christ preserves the ministers, however angry some are in the church, for accusing their wickedness. Privily therefore he warns them to crave the spirit, to nourish the spiritual life, and to trust in Christ, who will defend the ministers and advance them.

### Paragraph 265

As he testifies in all his other letters, he repeats it here also: “I know your works.” Of this I have spoken before. The Lord is ignorant of nothing that is done in the church, which is also the searcher of hearts. And especially he condemns this in this church,[note 15.5] that she thought herself alive where she was dead. He speaks not of physical, but of spiritual life and death. For Christ lives by his spirit in his saints and faithful, and shows lively works through them, as the Lord teaches in John 6 and many other places in the Gospel of Saint John. The Apostle said also that he did not now live,[note 15.6] but that Christ lived in him. The same Apostle said that widows living in luxury, though alive, were dead. Therefore, those are dead who do not have Christ living in them by faith and spirit, who do not have the virtue of Christ working in them—that is, who do not produce lively works.[note 15.7] For the Lord is reported to have said also in the Gospel: “Let the dead bury their dead.” The Sardinians therefore had the name of living men, that is to say, they were called Christians, spiritual, regenerated, and holy worshippers of God; but they were dead, namely, hypocrites, in whom no spirit nor Christian life appeared. The flesh, the world, and corruption still lived in them. But such churches displease Christ. There are many such today. But how does Christ reject them? Truly, he condemns such, but not to destroy them—for so the world condemns—but that they should repent. For he does not will the death of a sinner, but rather that he should turn and live. And therefore, consequently, he prepares a medicine for the disease.

### Paragraph 266

First, he prescribes to the stars, or bishops, what they should do in this situation. Then, he tells the entire congregation their duty. From this, we learn how similar ailments in churches are to be healed. This belongs to the pastors, whom he commanded to watch diligently over the flock and to strengthen those who remained, not yet truly lost, but close to destruction, unless they are helped in time with sound and wholesome doctrine. He no doubt alluded to that pastoral care and responsibility which the Lord describes in chapter thirty-four of Ezekiel. The flock is strengthened by the word of God; by it, they are turned away from death and preserved in life, and so forth.

### Paragraph 267

Now he also adds the reason why he commands that the flock be confirmed, lest they slip into death. For I have not found your works full or perfect before God. The Greek copies of Complutensian and Arethas read, “my God.” By works he understands all things that are done, words, deeds, and the entire conduct of people. The works, even those of the saints, are always imperfect if we consider human weakness. For as long as we live here, the flesh fights against the spirit; so that Job said he feared all his works and therefore fled to the clemency of the judge. Nevertheless, they are perfect and full in respect of Christ. For he is our fullness, and in him we are complete, as in John 1, Ephesians 1, and Colossians 1:12. And he makes us partakers of his fullness by faith. Those of Sardis were destitute of true faith, therefore every work of theirs must needs be imperfect before God, who allows nothing but that which comes from the Son and is most pure. Therefore the Lord commands that faith be taught diligently and instilled, that they may be made perfect in Christ. This is the best medicine for the deadly disease of Christ’s church.

### Paragraph 268

Here follows the duty of the people,[note 15.10] how they may be healed by apostolic repentance. The chief point is to remember the Lord’s words, in what we have heard and received them. We are not commanded to devise new forms of religion and repentance, but we are sent to the old tradition, not of men, but which we have in the Scriptures of the Evangelists and Apostles. These, I say, we ought to remember. For through custom of sinning, we forget God’s word. And truly, the beginning of Peter’s repentance was to have remembered the words of the Lord. Therefore, those who will not be reproved and instructed by God’s word shall never come to, or attain, true repentance.

Furthermore, it is necessary that we keep and retain the words of God, that is, the true doctrine of Christ, lest we forget it straightaway, or that we set it in vain contemplation and not in effectual work. The doctrine of Christ must be kept and performed in work. For in the last place it follows: and repent. True repentance consists in work: that in mind and body we should turn away from evil and turn unto God, and do good, being sorry for our wicked deeds past; this is the true apostolic repentance.

### Paragraph 269

To which repentance, now, after the divine, prophetic, and apostolic manner, he draws people by means of threats. These threats are to be understood as applying both to ministers and to the people in the congregation. Again, the Lord uses parables, as we read in Matthew 24. With these, he exhorts to watchfulness and sobriety. Since that passage is explained at length, I need not use many words about it here. To the Lord be praise and thanks giving forevermore.

## Sermon 16: ¶ He allows and commends those that covet to live godly in the Church of Sardis, exhorting them that they would so hold on and proceed.

### Paragraph 270

.

### Paragraph 271

But you have a few names in Sardis who have not defiled their garments; and they shall walk with me in white, for they are worthy. He who overcomes shall be clothed in white array, and I will not blot out his name from the book of life; and I will confess his name before my Father and before the angels. He who has ears let him hear what the Spirit says to the congregations.

### Paragraph 272

The second part of this heavenly letter is contained in these points, wherein the innocence, holiness, and integrity of the faithful in the congregation of Sardis, in true religion, are praised and commended. He exhorts them with a very generous promise to persevere. Finally, he again proposes to them abundant rewards: even to the corrupt, if they repent; and to the faithful, if they continue as they are.

### Paragraph 273

The Complutensian book reads thus: But you have a few names in Sardis. This is as much as to say that not all are corrupt and dead among you, although indeed those are very few. And so Arethas reads it in Greek, and the common translation in Latin, as found in other copies, which Erasmus follows: you have a few names also at Sardis; that is, even in Sardis you have names [note 16.2], but few. And he uses “names” to signify notable men. This manner of speaking is also in our language. For we say, “there is no man of name,” meaning no excellent or noble personage. He signifies therefore that there are in the same church noble personages, and that noble in soundness of faith and holiness of life, but very few [note 16.3], if they should be referred or compared to the number of hypocrites or the dead, which in truth are a great deal more. Neither ought we here to marvel. For the Lord says also in the Gospel that many are called, few are chosen; and that the greater part of this world walks in that broad and wide way of perdition, Matthew 20:7. Which also Saint Peter repeated in the second chapter of his second epistle. Those who seek to defend their error by a multitude are rather to be hissed at than confuted. You shall hear very often at this day: “We are but a few in number, we are innumerable, and therefore our matter is the better.”

### Paragraph 274

But that same excellent thing is chiefly to be observed, that although they were but few good, yet nevertheless the Lord commends and extols those few, doubtless for the example and imitation of all other churches. The words in deed are short, but the praise most ample and large. That they had not defiled their garments: which is as much as if he had said, you have not polluted your souls with strange opinions or spots of heresy. For you have remained sincere in the true faith: your bodies also, and the whole conversation of your life, you have not defiled with filthy lusts, with fleshly pleasures and voluptuousness. Doubtless this is the greatest praise and most certain sign of perfect godliness: wherewith I would wish that more of us were marked. But the manner of speech here requires also an exposition. The allegory of garments is often and much used in holy scripture. The use of apparel invented of God himself, and showed to our forefathers, has this chief property, to hide the privy parts of our body, to beautify and set forth the body, and keep of heat and cold. And therefore Christ himself is called the garment of Christians, and in the gospel in deed the wedding garment. Whereupon the Apostle advises us to put on the new man, which is made after God even Christ himself. Romans 13, Ephesians 4, Colossians 3. For Christ covers not only our privy parts, but all the filthiness also of the soul, he adorns and beautifies us, and drives from us all injury, and all evil. And we defile this garment, when neither in faith nor in holiness of life we do answer to our profession. For Christ is our garment, and Christianity, sincere faith, and holiness of life are our apparel: and even faith and our conversation is our garment. Forasmuch therefore as the Sardinians were of a sincere faith, and incorrupt manners, they are said to have kept their garments clean and undefiled. The Lord also gives now a reward unto virtue. And they shall walk with me, says he, in white array. These excellent things verily doth he rehearse to retain the Sardinians in their duty, to nourish them to greater things and to move other also to sincerity and integrity. Saints walk with Christ in white array, that is to say, have fruition of the same glory, wherein we believe Christ to shine. For he desires his father, that he will grant to the faithful, that wherever he is, they may be with him, and see his glory, etc., in John 17. And with Saint Matthew, in the transformation or clarifying, the face of Christ appeared bright like the sun, his apparel and rest of his body as light. So appeared Christ unto John in the first chapter of this book, clothed in white array. Now therefore says he, the godly that have not defiled their garment, shall accompany me, having put on light also. He adds another thing, for they be worthy. This is the greatest praise, when the Captain says, that the soldier is worthy of honor and glory. The greatest shame or ignominy is, when it is said with us, thou art unworthy. The first kind of speech shows him to be most excellent in all kind of virtue, which is said to be worthy of eternal light; by the latter is signified, that he which is accounted unworthy of a good and excellent thing, is marvelous negligent and ungracious. But here we need not to reason of the merit and desert of worthiness. God pronounces his to be worthy of glory; the godly refer all the goodness that is in them unto grace, and still complain of their unworthiness. Not to reprove God of lying, but to praise and commend the excellent goodness that is in him, acknowledging in deed that he rewards good works, and dignifies the worthiness of saints: but they are nothing proud hereof, but acknowledge all this to come of grace. This appears in the doctrine of the Gospel, Luke 17, Matthew 25, where saints commended of God, for the works of mercy, seem to acknowledge nothing thereof.

### Paragraph 275

However, he declares at length the most ample promises of God,[note 16.8] whereby he may not only retain in their duty the saints and undefiled Sardinians, but might also reduce all others who go astray at all times into the way of repentance, integrity, and holiness. And three things he promises: first, indeed, white apparel, that is to say, glorifying and everlasting light, and the glorious company of Christ, of which I have spoken already. Secondly, and I will not put out his name out of the book of life. For just as cities have books wherein the names of their citizens are written,[note 16.9] so too is God in the scriptures said to have, after the manner of men, a book of life, or of his elect. What that book is, and whose name is written therein, none of us can tell, since none has looked therein. We must learn from the scriptures who are the citizens of the kingdom of God. For that their names are written in the book of life, no one need doubt. And Jerome says: as many as have believed, he has given them power that they may be made the children of God. Paul says: He who does not have the spirit of Christ is none of his. And the spirit cries in the minds of the godly, “Abba, Father.” The same apostle says: God has predestined us that he might adopt us as his children through Jesus Christ. Moreover: he has chosen us in Christ before the foundations of the world were laid. Therefore, all believers are written in that celestial number. Whoever therefore does not believe, or does not persevere in the faith, either they are not written in the book of life, or else they are put out again from the book of life. Finally, the Son acknowledges the believers and those who persevere in the true faith before his heavenly Father and his angels. And here he repeats the evangelical doctrine from the 10th chapter of Matthew and the 8th of Mark. And it is undoubtedly a great matter in that universal judgment to be known by the Son of God, by the high judge, to be greeted and friendly spoken to by him, and that to our great praise. If any prince would in a great assembly of people know you, yes, embrace and commend you, how happy and fortunate would you think yourself? But then the very Son of God, king of kings, and lord of lords, will embrace you. Let us think of these things in time and amend our manners.

### Paragraph 276

For all these things pertain to us, the final and resounding declaration of Saint John confirms: “Let him who has ears, hear,” and so on, as we have discussed elsewhere. To the Lord be praise and glory.

## Sermon 17: ¶ The Lord comendeth the virtues, namely the constancy of the congregation of Philadelphia. &c.

### Paragraph 277

.

### Paragraph 278

And write unto the angel of the congregation of Philadelphia: this says he who is holy and true, who has the key of David, who opens and no one can shut, and who shuts and no one can open. I know your works. Behold, I have set before you an open door, and no one can shut it; for you have little strength. And you have kept my word, and you have not denied my name. Behold, I shall give some of the congregation of Satan, who call themselves Jews, and are not, but do lie. Behold, I will make them, that they shall come and worship before your feet.

### Paragraph 279

In all other congregations, the Lord at the least found some fault; in the only church of Philadelphia, he finds nothing to blame. Not that any person in this life is so perfect that he does not need God’s grace. For David cries out, “Enter not, Lord, into judgment with your servant, for no living person shall be justified in your sight.” But Saint John and Saint Paul also acknowledge all people to be subject to sin, a matter which Saint Augustine also discusses learnedly against the Pelagians. Therefore, when he finds nothing to blame in this congregation, it is not to be understood as though it were not tainted with daily faults, but rather that he imputes nothing because the sincerity and integrity of faith covers and hides whatever vice there may be. For there is no condemnation to those who are grafted in Christ Jesus. And although other churches also possess the true faith, this one excels especially, &c. It might be referred chiefly to the Bishop of the same church.

### Paragraph 280

In this sixth epistle, he commends the sincere faith and steadfastness of faith, and admonishes to persevere, proposing abundant rewards. And it contains much learning and diverse material, which will become apparent in the treatise itself.

### Paragraph 281

And the Lord herein follows the same pattern, which we see he has followed in others. For it is one and the same doctrine with all churches and in all times. First, therefore, is shown to whom the epistle is written or dedicated: to the pastor and whole congregation of Philadelphia. [Note 17.2] Philadelphia was a city of Lydia, neither very famous nor yet obscure. We read how it has been often shaken by earthquakes and repaired again. Strabo mentions it in the 12th book of Geography, and so have other authors also. Yet it made itself famous by its virtues. After this, the Lord Christ is signified to be the Author of this epistle, who at other times also has told St. John what he should write.

### Paragraph 282

And to Christ are attributed three things, or rather Christ attributes three things to himself: that he is holy, true, and has the key of David. He has borrowed these from the image of the first chapter. Christ is holy because he is pure and clean from all defilement and all unrighteousness, very God, a consuming fire, doing no wrong to anyone, having nothing at all that may be blamed. For the Seraphim rightly say, “Holy, holy, holy, Lord God of hosts,” Isaiah 6:3. Christ is also the holiness of the saints, a sanctification that sanctifies all who are sanctified. He loves holiness in saints. Therefore, Christ is most truly called holy. Antichrist, the Pope, has taken upon himself this title, and sits on this beast so filthily, as if you were to call a privy or a jakes a rose garden. Spit upon that vile and filthy beast, which allows himself to be called the most holy father, and worship Christ, the holy one of all holy, unless you understand by that holiness not every holiness, but papal holiness—that is to say, stinking and swimming full of all abominations. Christ is likewise called true because he is eternal and faithful, evermore constant and incorruptible. He can neither deceive nor be deceived. He steadfastly keeps his promises. All his words are without doubt and true. Although flesh, which can endure no delay, often begins to doubt, yet not one point or jot of them falls away. The truth of the Lord endures forever. You stand upon a most sure foundation if you lean unto Christ, who in John 14 also calls himself the truth.

### Paragraph 283

Finally, he adds, “He has the key of David.” I spoke of the key in the first chapter. He alludes to chapter 22 of Isaiah, signifying the divine and almighty power of Christ, by which he brings us purified into the kingdom of heaven. This work, indeed, neither devils nor any power can hinder. It casts down the unholy into hell, and there is no one who can prevent it. Therefore, he says aptly and expressly, “He has,” not “He had” or “He will have,” but “He has now.” For he alone has this power, which he communicates to no other. The Pope of Rome lies, claiming to have this power. The only Son of God surpasses all others in this prerogative.

The apostles, as ministers and preachers, have received the keys of knowledge and of utterance, of learning, instruction, and introduction, by which they also, in threatening, exclude unbelievers from the kingdom of God and bind them in their sins. Almighty God, who has the ultimate authority, ratifies the judgment of the minister, which he pronounces not of himself, but according to Christ’s words. But these things agree well with what follows about the opened door, which no one can shut, and with the entire matter.

### Paragraph 284

Now the Lord proceeds to tell what he wishes: and as he has said in all his letters, so he repeats in this also, that he knows all things concerning this and every other congregation.

### Paragraph 285

And he praises the perseverance in faith within this congregation,[note 17.8] indicating that it proceeds from the grace of Christ. “You have little power, and seemingly no strength, which the world considers power—riches, worldly wisdom, fortunate success, many friends, and things of that kind. Therefore, you can attribute nothing to yourself, nothing to your own strength, not even this: that you are a church, and that the truth of the gospel is freely preached among you.” For I have opened this door. And by my strength I keep it open, so that no one can shut this door—that is, to prohibit, hinder, or take away, by any means, the preaching and grace once granted. To open the door is a common phrase used by the Apostle in 1 Corinthians 16 and 2 Corinthians 2. He opens the door, which provides an opportunity and prepares the way to enter. Therefore, the word opened the door of life, so that the faithful might enter; the unbelievers could not stop this way. For the hand of Christ held the door open. And these things indeed declare,[note 17.10] which explains why, in cities, towns, and villages not greatly furnished with power or strength, the course of the gospel proceeds with such fortunate success. And where many attempt to shut the door through schemes, trickery, threats, and persecutions, they cannot. These things are not accomplished through our skill and wisdom, but by the grace of God.

### Paragraph 286

Nevertheless, if anyone desires to understand those things and what follows particularly concerning the pastor or bishop of the church, I will not object. For where he was humble and instructed with no worldly wisdom, yet furnished with God’s grace, he opened the way of salvation, which those who seek to abolish the preaching of the gospel now cannot shut. The power of Christ preserved him.

### Paragraph 287

And now he more clearly urges or commends the faithful constancy in faith of the pastor and congregation: [note 17.11] “You have kept my word, and have not denied my name.” When the Lord opened the door, lit the lamp, and gave spiritual gifts, the pastor and the congregation received them and kept them, and so did not deny them, nor trample them underfoot. This is excellent praise. Would God there were many such churches found today. Here also, O you church of Christ, you may learn all things—and each individual thing—about the duty of pastors, of the church, and of all godly men and women. You had no merit at all; God shone upon you through his grace. You had no worthiness, no deserving, no power nor authority: Christ has revealed himself to you through his mercy. Embrace him therefore who offers himself to you, hold fast, and never at any time let him go. &c.

### Paragraph 288

And note that the Lord says, my word, not every man’s word, but mine. What the word of Christ is, it is known to all men. For that which is written in the Gospel, and first indeed by the Prophets, and after by the Apostles was set forth in holy writ, is the word of Christ. It is not Christ's word that contends with this, although it be set forth by councils and holy fathers. Christ does not acknowledge that word; he acknowledges his own. And this must be observed and kept. The word of Christ is observed when it is not corrupted with additions, detractions, and wrappings, but is kept sincere in its natural sense. It is not kept when it is corrupted or depraved with men's inventions and perverse interpretations. The word of Christ is kept when it is commended not with the mouth alone, but is also expressed with godly works in the whole life, and beautified with holiness. It is not observed when, without repentance, men live most filthily. Finally, the word of Christ is observed and kept when it is not with any loathing of ours or impatience cast away, denied and forsaken. And therefore he affirms immediately, “You have not denied my name.” I have spoken elsewhere largely of confessing and denying of Christ's name. These things truly did the Philadelphians, and with these virtues through faith pleased the Lord. By these also may we commend ourselves to our Savior.

### Paragraph 289

Furthermore, the Lord shows with how great a reward he would honor the constancy of the faithful in faith. “You now have many enemies,” he says, “because of your pure religion; but if you continue to hold fast, I will cause those same enemies to become your friends and ultimately companions in your faith. Those who have hitherto condemned you as evildoers and heretics will come to you with great humility, asking forgiveness, ready to receive your religion and to worship him whom they have blasphemed. And they shall come in the most humble manner, with the greatest humility possible. For so Isaiah foretold that it would come to pass in chapter 49, to which the Lord alludes at this present.

### Paragraph 290

In the meantime, he addresses the Jews, singular enemies of the faith [note 17.15], whom he calls the Synagogue of Satan. For their teacher was none other than the Devil, as indeed they still have no better today. He calls them false Jews and liars. For neither did they confess the Lord, nor glorify God, nor believe in Christ their Messiah. But those who are Jews in deed are not such, as the Apostle Saint Paul said in the second chapter to the Romans. The power of God compelled many of them, forsaking their Jewishness, to embrace the Christian religion.

### Paragraph 291

Therefore, if we desire or seek to retain also in our churches the pure word of God, and to receive our enemies humbly, we shall not achieve these things through wars or injustices, through railing and reproachful words, but by steadfast faith. But if either we do not profess our faith purely, or do not adorn it with virtues, what is it to be wondered at, though enemies remain enemies still, and continue to hate us every day more intensely than before, and at length oppress us, and extinguish the light of God’s word with many? Let us learn, dear brethren, by godliness, constancy, and holiness to win our brethren. The Lord Jesus grant us his grace to accomplish the same.

## Sermon 18: ¶ He exhorts them to persevere in the true faith propoundyng most ample rewards.

### Paragraph 292

.

### Paragraph 293

And they will know that I have loved you, because you have kept the words of my patience. Therefore, I will keep you from the hour of temptation that is coming upon the whole world to test those who dwell on the earth. Behold, I am coming shortly. Hold fast to what you have, so that no one takes away your crown.

### Paragraph 294

A commendable virtue is praised in the congregation of Philadelphia: they have kept the word of Christ, not every word, but the word of Christ, and have not denied it. And he has begun to recount great rewards, which both he has given to this church, and is also ready to give to any other church similarly devoted to godly religion. For we are drawn by rewards.

### Paragraph 295

First, he says, I will convert your enemies, so that they may become your friends and brethren, coming into the congregation to worship Christ, whom they have blasphemed until now, and submitting themselves humbly and lowly. As we read of Saint Paul, who in the 15th chapter of the first epistle to the Corinthians says that he is unworthy to be called an apostle, and so forth. And this is a wonderful benefit. For God is glorified by those who are converted, and the truth is proclaimed; lies and superstition are confounded. The saints cannot but be exceedingly glad. The faithful also are delivered out of the Devil’s clutches and are saved.

### Paragraph 296

Then follows another benefit of God. God’s enemies will discover and find that the church, and every member of it, are God’s well-beloved children. God’s enemies suppose the faithful to be wicked, enemies of God, heretics, church robbers, hated by God, and unworthy to live. But they will understand that nothing is dearer to God than the church, for whom he gave his son, whom he also chose as his spouse, and has made a partaker of his kingdom.

### Paragraph 297

[note 18.4]But of this love of God, by which he, provoked by no merit of ours, but by his only grace and innate goodness, has joined himself to the church, all virtues proceed. That chiefly which immediately follows: that the church has kept the word of patience. John in his canonical epistle says, “not that we have loved God, but that he has loved us, and so forth.” Therefore, where the observance of the word of patience is joined as the cause of love, it must be religiously explained, that the favor of God and all our gifts are truly of grace, but yet that he, through that same grace, seems to requite and reward us for our efforts. The saints are not proud because of this, but humbly acknowledge and preach grace everywhere and in all things.

### Paragraph 298

Again he commends the perseverance of the faithful in the true religion. You have kept, he says, [note 18.5] the word of my patience. The word of patience is the Gospel of eternal salvation, which Paul also calls the word of the cross, and that for two reasons. First, because he describes the cross and patience of Christ by which we are saved. And again, he urges us also to bear the cross and patiently suffer with Christ, Matthew 16:24; 2 Timothy 2:3. Neither must anyone expect perseverance from one who is impatient. The Lord says in Luke 12:34, "In your patience you will possess your souls." Therefore, either the pastor or the church of Philadelphia has kept the word of patience, namely, by retaining the patience of Christ in their hearts through faith, and by showing patience in words and sayings, and enduring much hardship in body. This indeed is the best way to keep churches and every faithful person safe and sound. Let them keep the word of Christ’s patience, and commit the rest to the Lord.

### Paragraph 299

For it follows: And I will keep you again from the hour of temptation, &c.[note 18.6] The hour of temptation is explained in two ways. Either he speaks of heresies and heretics, by whose talk, and crafty juggling, lewdness, and deceptiveness tempts the faith, simplicity, and integrity of the faithful. Of this, the Lord treats much in the 13th Chapter of Deuteronomy. Or else he speaks truly of the persecutions which the emperors of Rome have inflicted, among whom Trajan, a most mighty prince, issued severe proclamations against the Christians. Of this, Pliny also made mention in the 10th book of Epistles, the holder and one. But Christ preserved the church of Philadelphia, and keeps also at this day the faithful by his word and power in the perils of heretics and heresies, and finally of persecutions also: so that the faithful may stand sure in all controversies and receive nothing of heretics that is strange from God’s word, and also give no place in persecutions. Christ causes many times that the burden of persecution presses not so heavily. Therefore let us always be steadfast in God’s word, and permit the defense to our Lord God. He will not neglect us, &c.

### Paragraph 300

But because in temptations and afflictions the Lord often appears to our flesh to delay, and to seem to neglect us, we say that the Lord anticipates and adds, “Behold, I come shortly.” Shortly, I say, that is to say, in time—neither late nor hastily. We say this neither too soon nor too late, but in due time and season. If the Lord therefore seems to be slow, do not despair, for he will come at the opportune moment when he sees fit. Do not prescribe to him the manner and means of deliverance, but wait for the Lord’s appointed time. Read what wholesome things Paul has written concerning this matter at the end of the tenth chapter to the Hebrews, where a passage also from the second chapter of Habakkuk is cited.

### Paragraph 301

And now he exhorts in few but most evident words to perseverance in piety, wherein they had excelled hitherto. And he says two things are to be held fast: that which thou hast. They had the gospel of Christ, and the word of eternal life, the true faith and godly religion. These things he commands to hold fast, and to persist in the religion one has received. And whilst he commanded them to keep that which they had, he signifies by the way that no new or other doctrine is to be looked for, but that this one received doth suffice. Let us not therefore think in the government of the church upon other laws, upon other traditions, than of the Gospel of Jesus Christ. This is sufficient for the church. After reasoning as it were of the loss, he says: Therefore must thou watch diligently and hold strongly the gospel, for if this is taken away thou art spoiled of thy crown. The crown is a token of virtue and victory.

### Paragraph 302

Rulers and those worthy of empire are crowned. The virgin loses her crown, that is, is defiled. Therefore, heretics, false prophets, and seducers take away the crown when they seduce and corrupt. Thus says the Lord: “You have gained honor and glory; see that no one takes it from you.” So we read that Paul spoke in 2 Colossians. Let no one take the victory from you. In Ezekiel 18, the Lord testifies that he will not impute righteousness to the righteous if he forsakes and abandons his righteousness. Let us therefore pray that we may always persevere in the word of the Lord.

## Sermon 19: ¶ He proceeds in recityng most great rewards.

### Paragraph 303

.

### Paragraph 304

He who overcomes, I will make a pillar in the temple of my God, and he shall go no more out. And I will write upon him the name of my God, and the name of the city of my God, New Jerusalem, which comes down out of Heaven from my God, and I will write upon him my new name. Let him who has ears hear what the Spirit says to the congregations.

### Paragraph 305

Our Lord proceeds in recounting much more ample rewards, which he would give to those who overcome: and so he tempers his words, that we may easily perceive that this promise does not only pertain to the congregation of Philadelphia, but to all the churches in the whole world, yes, and to every of the faithful. And as we have oftentimes repeated already (for I am not ashamed to repeat, saying that the Lord himself so greatly urges the victory), again we say, that those things are promised not to such as fight lightly or negligently (for diverse fight, and by and by run away) but to those that overcome and persevere to the end. For our life is a warfare upon earth, which Job also has confessed. The soldier has a sure purpose to overcome his enemies. Our enemies are the devil, the world, and the flesh. Against these we must earnestly fight, of no other intent but that we may overcome. The Apostle in the 6th to the Ephesians described the armor of the faithful. In victory the saints consider sincerity and integrity, that we lose nothing of the known truth; but let us retain the pure word of God and sincere faith, and let us keep our bodies and our souls clean from all pollution, and that to our lives’ end. He propounds most ample rewards by promise. Whereby he alludes to the manner of Greeks and Romans, who decreed images to those who deserved well of the common wealth, in which also they wrote their virtues, for whose cause they were set up either in the court or marketplace or else where. They seemed by this means to deliver to their posterity as it were by hand the glory of their elders, which they made also as it were everlasting. Otherwise the use of this vocable Column, or pillar is diverse. Jeremiah is called by God a pillar for his constancy. The Apostles are called by Paul in the 2nd to the Galatians chief pillars, for their excellence, and that the churches leaned upon them, for the preaching of the truth. The church itself also is called the pillar and base of truth, for as much as it is grounded upon the sure rock Christ. In the Temple of Solomon stood two columns or pillars, figures of the everlasting kingdom of Christ, and of the holy church. In this place a pillar is taken for a man, in glory and beauty excellent. For he says not that he will erect a pillar for a godly man: but I will, says he, make him a pillar, that is to say, I will beautify him with honors and glory everlasting.

### Paragraph 306

But where shall this pillar be established? where shall the glory of the saints be renowned? [note 19.5] Not in a royal court or marketplace, but in the temple of my God. And the temple of God is heaven itself, and in this world, the holy Church. Therefore, he will be glorious in the church of the saints, both those engaged in warfare [militaunt] and those victorious [triumphaunt]. Although the faithful may hear evil in this world, this world will perish, Christ will reign forever, and the saints will reign with him. Their glory will remain forever and ever. And where he says, “of my God,” Arethas explains and says: This saying “of my God” does not diminish the divine nature seen in Christ, but establishes, as I may say, the consubstantiality. For it declares the union of two natures that are in the person of our Lord Jesus, namely his deity and humanity, even after understanding, yet not confusedly, to be inseparable. For they answer one another mutually, because of the assumption of the human nature, the divine to the human, and likewise the human to the divine properties. &c.

### Paragraph 307

Moreover, the lasting and steadfast nature of the glory of the saints and faithful is signified, where it is added [note 19.6], “and he shall go no more out.” For often pillars are broken and cast down, and the renown of great men perishes and fades away. But Christ promises to those who overcome that they shall never be cast out of the fellowship of saints, neither that the glory of the faithful should be obscured at any time. And thus far concerning the pillar itself.

### Paragraph 308

Here follows a description of the inscription on the pillar, noting what sort it will be. Three things chiefly are written on the Saints: the name of God, the name also of the city of God, and the new name of God or of Christ. Which we shall discuss in order. First, the name of God is ascribed to the godly, that is, they themselves are called by the name of God, and are the children and heirs of God. This is discussed at length in 1 John, chapters 1-8, and in Romans 8. And what can you devise to be more honorable than to be, and be called the son, not of a king nor emperor, but of the living God? But this same noble grace the Lord grants to those who overcome. This is found in 1 John 3 and 5. Secondly, to the overcomers is inscribed the name of the city of God: that is to say, the godly man is written in the number of the citizens of the city of God, and is truly a citizen of the city of God, I say even of the city of God. It was a great matter in times past to be a citizen of Rome, but it is far greater to be a citizen of the city of God. The citizens enjoy all privileges and commodities, finally the glory of the city. But this is greater and more than that it can be declared in few words.

### Paragraph 309

But the church is the city of God: and the city of God is the church. This is presented here with three descriptions or titles, from which it is easy to judge what the church is, or what we should think of it. The church is the city of God. For just as a city is the fellowship of citizens, so the church is the communion of saints. The Prince of it is Christ, the head of the church. The entire world itself was a representation of this church, and the very setting up of the tents, in the midst of which was seen the Tabernacle, was a sign of the deity present, as if a cohabiter. &c. For the Lord is in the midst of the church, as we have read in Leviticus 26 and 2 Corinthians 6. Secondly, the church is called New Jerusalem. For the old [?] was a representation of the new. This earthly church is New Jerusalem, that is to say, spiritual. Paul also affirms this in Galatians 4. For in the third place, this newness is explained. It is not built by humans, but comes down from heaven above. For unless we are born from above by spirit and by immortal seed, namely the word of God, we cannot be members of the church. And we are born by a spiritual regeneration, the children of Christ and of the church, as the Lord himself discusses at length in John 3.

### Paragraph 310

1 Peter 3:1. And Paul, in his first letter to the Corinthians, chapters 3 and 4, will speak further of the New Jerusalem toward the end of this book. But understand this: what is the church of Christ, the fellowship of the faithful, regenerated by the word of God, and so forth.

### Paragraph 311

Finally, in those who overcome, a new name is written, and indeed the new name of Christ. Not only that they should be called Christians of Christ, but because the name is a brief description of everything and nature, and a new name is promised, it follows that we should understand that people will be renewed, chiefly by glorifying. He promises therefore a glorifying to the godly. This is spoken elsewhere in Matthew 17:1, 1 Corinthians 15, Philippians 3, and 1 John 3. These most ample rewards the saints may truly expect if they fight to overcome.

### Paragraph 312

Hereunto is annexed the customary affirmation, by which both this doctrine is applied and communicated to all churches throughout the world. And it is declared that it came not of men, as a vain thing, but of the very spirit of God, most true. This spirit the Lord grant us.

## Sermon 20: ¶ The Lord blames sore the church of Laodicea.

### Paragraph 313

.

### Paragraph 314

And unto the angel of the congregation which is in Laodicea, write: This saith Amen, the faithful and true witness, the beginning of the creation of God. I know your works, that you are neither cold nor hot; I would that you were either cold or hot. So then, because you are lukewarm, and neither cold nor hot, I will spit you out of my mouth. Because you say, ‘I am rich and increased with goods, and have need of nothing,’ and do not know that you are wretched, miserable, poor, blind, and naked.

### Paragraph 315

The seventh and final Epistle from our Savior Christ [note 20.1] is written by John to the Bishop of Laodicea. It is a great rebuke of that people, finding nothing worthy of praise in them. Nevertheless, it is a faithful admonition or exhortation to repentance. And following his customary practice, he indicates to whom he is writing and from whom the Epistle originates. The Epistle was commissioned by Christ to the Bishop of Laodicea and to the entire congregation. Therefore, something must be said about the Laodiceans, so that the rest may better understand and consider them.

### Paragraph 316

Laodicea, the chief city of Caria, as Strabo and Pliny state, stands by the river Lycus. Antiochus Theos built the city and named it after his wife. It was the wealthiest city in Asia, as Vadiane also notes in his Epitome. It had a very profitable trade in woolen cloth. It appears that Paul also preached the gospel there. For he mentions Laodicea, and some believe he wrote his first letter to Timothy from there. Certainly, it appears that the people of Laodicea had received the gospel, but in a corrupted form. For they sought to blend the world and the church together, and to join Christ and Mammon, as is said today.

### Paragraph 317

Therefore they did not abandon their greed, and their immoderate commerce (for to conduct trade moderately, without deception, no religion forbids), nor their excessive riot and pride. They did not seem to lack anything, but to possess and appear to possess all things, because they were wealthy. Against these, the Lord contends severely, declaring them to be very wretched, and more destitute than simple beggars. For just as in the church of Philadelphia he found no fault, so in this one he commends nothing at all.

### Paragraph 318

You will find many today who speak like this, and it is ever on their lips: “I have learned to be a follower of the Gospel, and to be a soldier, to drink, to play the whoremonger, and to live at pleasure.” You will find churches like these, serving both Christ and Mammon, or commerce, Bacchus, Venus, and the god of battle. Both they and all such people here are refuted and called to repentance. This demonstrates that God’s mercy is greatest, not abandoning or rejecting churches so corrupt, and people so full of such great wickedness.

### Paragraph 319

Woe to those who despise this boundless mercy and goodness of God and his long-suffering, and persist in their wrongdoing.

### Paragraph 320

Christ is described here again most abundantly, as he is in the previous titles. Indeed, it can be gathered from all that has been said that this is the best and most complete description of Christ, so that no further explanation from human sources is needed. He reveals himself with a new name, calling him “Ho Amen,” that is, Amen. This is a Hebrew word, frequently used in the Gospels, especially in John, and by Paul in 2 Corinthians 1:20. Paul says that Christ, the Son of God, whom he and Sylvanus and Timothy preach among you, is not fickle, but in him is truth. For all the promises of God are in him, and in him they are confirmed to the praise of God through us. And the Lord explains why he calls himself that Amen: “For I am the witness,” that is, the trustworthy, faithful, constant, and true one. For Christ has been given to us by the Father to testify about the will of the Father. And his testimony, as he himself repeats often in the Gospel of John, is firm, constant, sure, certain, and true, without falsehood, doubt, or inconsistency. These things agree well with the argument in which he rebukes the Laodiceans for their sin and urges them to repentance. It is a difficult matter for the flesh to hear such doctrine; but when the certainty, assuredness, or truthfulness of the teacher is perceived, it will commonly move people’s minds, if they are not entirely hopeless and despairing.

### Paragraph 321

He adds moreover another thing,[note 20.7] which declares his dignity. For he calls himself the beginning of the creatures of God. Neither ought the Arians to seek here any defense for themselves. For it is not fitting by any one place, much less by a little word, to subvert the whole scripture, and to strive with the articles of the creed, the lively tradition of the Apostles. Our Savior Christ is considered after his deity and after his humanity. After his deity, he has no beginning, but is rather the beginning (actively as it is commonly said, not passively of all things and creatures. Neither is he a creature: for all things are made by him. Which thing both the Evangelical and Apostolic scriptures prove, John 1 and Colossians 1, and Hebrews 1, where you have passages expounding this same one. After his humanity he is called the beginning of the creature of God (namely man, which is called a creature by reason of his excellency, and for that he is the lord of creatures, for whom all things were made) as he is called the firstborn of the dead. For in Christ humankind is repaired, that it has not perished: God looked upon the countenance of his Christ when he first made man. For Christ is the beginning, that is to say, the preserver of the human nature: as it has elsewhere been told you at large. Hitherto we have had the description of Christ, which is called Amen, and the beginning of the creature of God, by whom verily all things are made, which is very and true God, witness of the divine will of God. &c.

### Paragraph 322

Now he tells the church what opinion he has of her, and what she is, that is to say, blames her. And as he has impressed upon all the former that he knows all their works, so he does to this also. First, he shows that he knows this of the church of Laodicea, and especially of its bishop, that he is neither cold nor hot. He adds, I would it were better if you were altogether cold, or thoroughly hot. But now you are lukewarm, or blood warm. An allegory taken of men’s meat, or of cold, hot, or warm water, and it is applied in a manner proverbially. He is cold, who openly follows the world, being wrapped in heathenish errors and sins of this world, and boasts nothing, or will seem to have nothing to do with the true religion. He is hot, whose breast inflames with the Holy Spirit, contemns the world, loves the true religion exceedingly, and lives a holy life. He is warm, or between both, which has neither forsaken the world, his errors, and sins, nor has fully received Christ, his truth and righteousness, but serves partly the world, partly Christ: In outward things he shows himself to be a Christian, in resorting to holy assemblies and receiving the sacraments, but inwardly he is so besieged by the world that he lives a worldly life rather than a Christian. Such a mixture the Lord does not allow, which elsewhere forbids plowing with an ox and an ass, and making a garment of linen and wool; to pour new wine into old bottles, and to patch an old garment with new cloth.

### Paragraph 323

[note 20.12]In religious practices, a blending of different traditions is often less acceptable to God. For there are those who attempt to combine various religions, compiling many into one. Muhammad constructed his religion from Jewish and Christian teachings. Many today create a mishmash of Roman Catholicism and the Gospel, or bake a “chuchurnullis,” as the Germans call it, a cake made of diverse ingredients. If a Roman Catholic observes such a service, he does not recognize it as his own; and if an adherent of the Gospel sees it, he recognizes it as belonging to no one. For it is a mixture of the pure and the corrupt, where the pure element has no greater strength, and the corrupt element generally prevails. The masses used by many today are of this sort, neither entirely Roman Catholic nor wholly Evangelical. The Lord’s Supper does not appear in them; the Roman Catholic mass is also cut off and altered within them. If we believe that Christ established the finest rule for religion and living, why do we not follow only that Master? But we value the favor of men more, which we are unwilling to lose. We do not value the favor of Christ enough to recall the Apostle’s saying: “If I pleased men, I should not be the servant of Christ.”

### Paragraph 324

But hear what the Lord says to these hypocrites. [note 20.13] It would be better, he says, if you were either cold or hot. It would be better if you were a sinner or a heathen than a hypocrite and a mongrel. For then you might be more easily helped, according to that saying of the Lord, “If you were blind, you would have no sin.” [note 20.14] Now, when you seem to yourselves to be just and sufficiently taught and furnished with godly doctrines and practices that please God, you leave no place for further instruction, but, despising the word of God and Christ’s institution, you prefer your mixtures before all the justifications of God. The Lord also in the Gospel says to the Pharisees: “Truly I say to you, that tax collectors and prostitutes go before you into the kingdom of God.” [note 20.15]

### Paragraph 325

The other part is clear enough; it would be better if they were fervent, namely with the Spirit of God, which thing the Apostle requests in the twelfth chapter to the Romans.

### Paragraph 326

Furthermore, he threatens to plague them if they continue, as they have begun to be neutral; [illegible]. “I will spew thee out of my mouth.” By this manner of speaking, two things are signified: both the loathsomeness which God conceives of this neutrality or lukewarmness, and the vomiting out, which punishes the same. Warm water provokes a vomit. To this, he appears to have alluded, as likewise to that old phrase of speaking: “The land has vomited the Canaanites, and the same shall vomit you up also.” Therefore, these composers or mongers, with their temperature and mixture, so displease God that they engender in him a loathsomeness, are an abomination to him, so that finally he shakes them off—the same we understand of those who join together Christ and Mammon. And the phrase of speech is to be noted, “now therefore,” or “so forasmuch as,” or “now seeing it is so,” etc. Moreover, the longsuffering of God is here noted, which does not plague immediately unless there appears nowhere any hope of amendment.

### Paragraph 327

He explains more fully the sin of the Laodiceans, and what causes their lukewarmness: because they love riches, in which they trust, believing they lack nothing. They consider themselves wise, and that they see all things, and are sufficiently supplied with spiritual and temporal possessions. It is less, they say, “we are rich.” More follows, “I am increased with goods”: that is to say, “I have acquired so many riches that I want nothing.”

### Paragraph 328

That same he now refutes, and shows that they are utterly deceived, and are to be miserable people. For he rebukes them severely, and says, “You do not know what you are.” That ignorance is a great evil, and the beginning of desperate blindness, when a man thinks he has something that he does not. For such persevere in their error, and admit no counselor. Therefore the Lord says, “You do not know that you are poor, wretched, weary, and worn with evils.” For they are toiled with many labors that serve this world. They are miserable. You do not see your own misery; others who see it are deeply sorry. You do not see in what condition you are. This kind of speech signifies a man very wretched and desperate, whose misery others see, but he himself sees nothing, like a poor or a beggar. You think yourself very rich, but you are a stark beggar. Covetous rich men are poor; they are poor also in virtues. The people of Laodicea were blind, as the Pharisees were called blind in John 9. Well-sighted in worldly matters, yet heavily blind as betels. Naked, or destitute of good works. Void of your wedding garment. They nevertheless were richly arrayed with garments of the finest wool. But before God they appeared naked. Let the gallants of this world, or proud peacocks rather, who are so well-sighted and gorgeously appareled, mark these things well. The Lord give them understanding.

## Sermon 21: ¶ The Lord gives wholesome counsel to the Laodiceans, admonishing them to repent.

### Paragraph 329

.

### Paragraph 330

I advise you to buy from me gold refined by fire, so that you may be wealthy; and white clothing, so that you may be clothed, and the shame of your nakedness may not be revealed; and anoint your eyes with ointment, so that you may see. As many as I love, I rebuke and chasten. Be zealous, therefore, and repent.

### Paragraph 331

Since God does not desire the death of a sinner, but rather that he should repent and live: Therefore, after severely rebuking the church of Laodicea, he gives her wholesome counsel, admonishing and exhorting her to repentance, and indicates what true repentance is.

### Paragraph 332

The Lord uses the word of counseling,[note 21.2] not of commanding, to attempt to confound the madness of those who, unless they are forcefully drawn, do not think themselves admonished, allured, or called by the Lord. And while they look for such a drawing, they neglect all of God’s counsel and fall from true salvation. God counsels his elect with things that are wholesome; the chosen obey good counsel. God touches their hearts inwardly, and outwardly by preaching of the word, and by sundry admonitions he pulls and drives man from evil to good. This counsel of God is not to be despised, and another violent vocation to be imagined. God’s word must be heard. “Today,” says the Prophet, “if you hear his voice, do not harden your hearts.” When the Lord counsels with his word, and the hearers harden their minds, they do so through their own fault and are made authors of their own destruction. But those who receive God’s counsel receive it not by the force of free will, but by the grace of God, which works in us to will and to perform.

### Paragraph 333

Therefore, when the Lord advises beneficial things, the elect pray that they may receive them; and they receive them through grace, obeying the counsels of God.

### Paragraph 334

And part of the wholesome counsel is this: “I know your works, says the Lord, that you have tried and have found Me faithful, as gold tried in the fire; that you may be rich, may be clothed, and may receive eye salve to anoint your eyes.” He sets these things as a medicine against the diseases which He revealed before, calling the church of the Laodiceans poor, naked, and blind. Now therefore He teaches them how they may be rich, may be clothed, and may receive their sight again, if they truly obtain gold tried in the fire, or are purified.

### Paragraph 335

[note 21.4]And gold tried in the fire is gold most purified and clean, having no base metals or impurities, but pure and clean gold. This shadows the word of God, of which the Prophet sang: “The word of the Lord is a pure word, silver tried in the fire, seven times purged in a vessel of earth.” Certainly, the word of God is light, coming from the eternal and most pure light, having no part of human defilement or affections, savoring of no errors, teaching nothing that is corrupt. However, of itself it will profit a man nothing unless it is received with a true and sincere faith. Therefore, I do not separate faith from the word, and say that the pure and sincere faith is signified by gold. Whereof Saint Peter said that the faith of our hearts should be purified. For although there be in us spots and infirmities, yet is faith, by reason of the subject upon which it rests, most pure.

### Paragraph 336

The word of promise, and even Christ himself, is the object of faith, which is purity itself. Therefore, the Lord advises the congregation of Laodicea to be tested by fire, and advises them to hear God’s word and believe it in deed. For the Lord uses the word of being for receiving, hearing, and obeying.

### Paragraph 337

For no man should suppose that there is bargaining before God, as there is with men: as though the spiritual gifts of God might be bought for money. This is contrary to the entire Scripture, and especially against the judgment of Saint Peter pronounced against Simon Magus. But this our exposition the Prophet Isaiah approves in the fifty-fifth chapter, where among other things it says, “Come without money, and without price, or exchange.” And immediately following: “Listen carefully, and give your attention.” Therefore, the Roman Antichrist has no claim to this, I mean the Pope, that great merchant, who sells all things in the church, even those things which he does not possess, the greatest deceiver in the world. Moreover, just as it is plainly expressed in Isaiah concerning from whom such graces or gifts are to be bought, so also Christ here says expressly, “I counsel you to buy of me.” Behold, he says of me, not of the Pope, of monks, friars, or priests. For Christ alone has the things which we may require. He alone satisfies, he alone grants those gifts. And therefore he says in the Gospel of Saint John: “Let him who is hungry or thirsty come to me.” To me, I say, let him come. John 4:6 and 7. And Saint Peter says, “Lord, to whom shall we go? You have the words of eternal life.” As though he should say: If we will live, we can go to no other, but unto you. You are the life and fountain of all goodness.

### Paragraph 338

Moreover, the use and benefit of this pure gold—that is, the word of God’s truth and pure faith, tried and most purified—is threefold. First, that you may be rich; secondly, that you may purchase clothing; thirdly, that you may buy the eye salve, to heal the blindness of your eyes. For the word of God and faith in him are the foundation of true piety. Without the word and faith, nothing is sound.

### Paragraph 339

The first fruit is wealth or riches, namely spiritual. For the word and faith are not a false imagination and a vain dream of excellent things. For he who believes the word feels joy in his heart and enjoys spiritual gifts; and possessing Christ through faith possesses all goodness. Whereupon also the Apostle in the first chapter of the first epistle to the Corinthians said: “I give thanks to my God always for you, for the grace of God that is given to you in Christ Jesus, because you are in all things enriched by him, in every word, and in all knowledge, as the testimony of Christ is confirmed in you. In so much that you are not destitute in any gift.” Consider these things well, those who think worldly goods to be true riches. These people shall be judged by the wisdom of God, as it is manifest in the twelfth chapter of Saint Luke. And besides this, those who are destitute of the light of God’s word and lack faith cannot rightly use these earthly riches. Therefore, heavenly riches are the true riches.

### Paragraph 340

The second fruit is the clothing and beautiful apparel with which we are covered, so that our shameful nakedness should not be revealed. Before their fall, our parents were naked, but without any shame or disgrace. After the fall, they were ashamed. Because sin brings shame, and a lack of all good works; and an ungodly lifestyle is a most shameful nakedness.

### Paragraph 341

The Laodiceans were afflicted with this. But Christ, who is learned through the word of truth and perceived by true faith, is the white apparel of the faithful, their righteousness and innocence. He covers all our blemishes, he abolishes our shameful nakedness, decks us with all kinds of virtues, so that we may appear honorable and comely before God in holy conduct. For Christ is the wedding garment. The Apostle counsels us to put on Christ, and that we be clothed with righteousness, temperance, and all goodness. The passages are in Romans 13, Ephesians 4, and Colossians 3. Away with the cowl of our Lady, under which gather for the most part wicked and impenitent persons. The most pure virgin does not cover such [?], she loves righteousness and repentance.

### Paragraph 342

Finally, with this gold is bought an eye salve, which is a medicine for the eyes, [note 21.10] as physicians are accustomed to apply to sore and blurred eyes, to prevent blindness. The commandment of the Lord, says David, is bright, giving light to the eyes. Faith also rightly informs human judgment, so that we may rightly discern virtues and vices. The lack of God’s word and true faith brings about blindness.

### Paragraph 343

For all these things the Lord exhorts the people of Laodicea to seek God’s word and believe it faithfully. Thus it will come to pass that, being enriched with all spiritual gifts, they may live a blameless life within the church, may possess Christ, and rightly judge all matters of salvation. And in these things also consists true repentance: forgiveness of sins and amendment of life. &c.

### Paragraph 344

But lest they should say, “We hear these things in vain,” as those who have heard before that we shall be spewed out of the Lord’s mouth; yea, and are so sharply rebuked with bitter words and sentences that we are compelled to despair—he anticipates that objection and says, “Whomsoever I love, I rebuke and chasten.” The first word signifies to accuse and reprove openly, which is done with sharper words; the latter is referred to discipline, whereby children are kept in awe with the rod, lest they forget themselves through wantonness. The Lord therefore, alluding to the words of Solomon in chapter 3, signifies that a sharp rebuke or severe chastening is not always a sign that God is angry, but often a token that he is pleased and loves us. Therefore he says, “First, I rebuked you sharply out of love, and so sought your salvation.” Therefore, it is now also a wholesome sign if the preachers rebuke the church with sharp words; and again, it is an unlucky sign if a fox’s tail is stroked over intolerable faults. It is a token of love also, if a man suffers sundry mishaps. Which thing the Apostle discourses at large in chapter 12 to the Hebrews.

### Paragraph 345

Upon these things he infers the substance,[note 21.12] and saith: Where you see God so earnestly seeking your salvation, I pray you do not always remain in a murmuring, neither hot nor cold. Be zealous, take unto yourself a fervent zeal to follow and grasp your salvation. For now he sets the fervor of faith, conceived by the word and spirit of God, against this neutrality or lukewarmness. After he adds, and repent, in forsaking your evil conversation, and being of Christ tried gold: that is, purified and purged, whereby you may be rich, be arrayed in white, and may have a medicine wherewith to anoint your eyes, that you may see. To God be glory.

## Sermon 22: ¶ He draws them also hereby unto repentance.

### Paragraph 346

.

### Paragraph 347

Behold, I stand at the door and knock. If anyone hears my voice and opens the door, I will come in to him, and will share a meal with him, and he with me. To the one who overcomes, I will grant the privilege of sitting with me on my throne, just as I overcame and have sat with my Father on his throne. Let the one who has ears hear what the Spirit says to the congregations.

### Paragraph 348

Hereby also the Lord allures the Laodiceans to repentance, showing that every time is suitable for conversion, and that God is ever ready to receive sinners, and constantly urges them to repent and live. And this matter he explains in an allegorical and beautiful speech, taken from the fifth chapter of the book of Song of Songs. For he imagines the Lord standing at the door and knocking, indeed promising to those who open it the greatest intimacy and unspeakable joys.

### Paragraph 349

First, therefore, is declared the benevolence of God toward sinners, and his most ready will always to receive them, yes, and his infinite study to move men to repentance, that they might live. For the Lord stands at the door, and knocks. The word “stands” signifies that God is always prepared, always watches over our salvation. For he sits not idly, nor lies on one side like a sluggard: He stands busily to his work. And “I stand,” he says, not “I stood” or “shall stand,” but “I stand” evermore ready, evermore loving and gentle. What does he do? He knocks, and that indeed at the door, desiring to be let in. For just as he who knocks at a door seriously covets to be let in, so God desires earnestly to be received by us. And God uses sundry kinds of knocking. For he warns and excites with his word by the Prophets, again by signs and wonders, and also by sundry chances and movements [?]. These things may be seen in the city of Jerusalem. He sends to them his Prophets and Apostles. He shows divers wonders. He brings on the sorrowful chances that they might admonish them—such as are reported in Luke 13, of the Galileans and of those whom the tower of Siloam had overwhelmed. We may see the like at this day, how the Lord knocks. Therefore he said truly, “Jerusalem, Jerusalem,” and so forth, Matthew 23. These are doubtless the parts and doings of God who will not that a sinner should die, but rather convert and live.

### Paragraph 350

Then we must see what is required of us, namely that we should hear the knocking and noise of the knocker, and also open and receive those who desire to come in. Here are refuted those who speak of humankind as though it were a block, and an image of some unknown design, saying: It is neither in the runner nor in the willer, and so on. Some entirely abstain from doing good, saying, “If I am chosen, that is enough.” But Scripture requires everywhere hearing and obedience. We know that people are only saved, and that in Christ: in Christ to be those who believe; that faith is of hearing, hearing by the word of God. Therefore, the prophet says, “This day if you hear his voice.” This is recited by the Apostle in Hebrews 4. The Apostle also, in 2 Timothy 2, says, “In a great house there are not only vessels of gold, but of earth also.” If anyone purges himself [?], and so on. And therefore the Lord says, “I knock.” It shall be your part, not to despise him who knocks, but to open to him. And he mentions indeed two things: to hear, which is required both in John 8 and John 9 of the children of God and the true sheep; and to open, that is to receive the Lord, or believe, to obey, and to shape oneself after the will of God, and to do penance. Nevertheless, we must be on our guard here, so that we do not think that humankind has power of itself to receive the Lord. The Lord illuminates his elect, and by him we can do all things; without him we can do nothing. Other passages must be consulted with this, such as John 15:2, 1 Corinthians 3, and Philippians 2. Therefore, those who open, open by the grace of God. Those who do not open are wrapped in their sins, not through their own fault, but through their own fault they do not open, and not through any fault of God.

### Paragraph 351

Let us hear moreover what the Lord promises to those who open [their hearts], that is to say, to such as receive Christ with true faith.[note 22.4] The Lord promises them two things chiefly. First, “I will go into him,” he says. Scripture signifies that Christ dwells everywhere through faith in the hearts of the faithful, and with a most tight bond to be joined unto them. “He who eats my flesh and drinks my blood abides in me, and I in him.” These things are spoken of the Lord in the sixth chapter of John. And in the fourteenth chapter, he says, “He who loves me will keep my word, and my Father and I will come unto him and will make abode with him.” Paul says that he no longer lives, but that Christ lives in him. He affirms that Christ through faith dwells in the hearts of the faithful. And so the Lord enters the hearts of those who let him in. Not the least part of felicity consists in this conjunction. For to be united with God is blessedness, which begins here and is made perfect in another life. And therefore, in the second place, the Lord says, “And I will sup with him, and he with me.” Whereby he notes not only again a most dear friendship and familiarity (for the table is consecrated to amity) but rather the fruition of eternal glory. For by the supper are signified the celestial joys, greatest and unspeakable, which the godly receive immediately after death, but more fully in the end of times, when the bodies shall arise again. Therefore it is not applied to a dinner, but to a supper, as it is also in the fourteenth chapter of Luke. If we receive Christ, we shall have him dwelling with us continually while we live in this world. And in the world to come we shall have the full fruition of all the celestial joys. These things are certain and true. For otherwise in the life to come there shall be no riotous banquets, such as the Turks do imagine.

### Paragraph 352

He adds another general promise, by which he exhorts and encourages the study of godly religion and repentance. For to him who overcomes is promised the kingdom of heaven. And he says to him who overcomes—of which I have spoken in other epistles—not to him who flees or is a coward, etc. He presents the example of conquering Christ. For we must overcome, as he has overcome. He indeed overcame most perfectly; we, with our limited strength, fight and overcome. And truly, the true victory in us is the lively virtue of Christ: that is to say, by him they overcome, whoever overcomes. And just as he, having overcome death and vanquished the world and the devil, ascended into heaven and sat on the right hand of the Father, so he promises us, upon overcoming, that he will give us the seat of his Father: not that we, sitting on the right hand of God, should judge over all flesh, being made Christ’s, but that being made partakers of everlasting glory and delivered from all judgment, we may appear in glory when he shall come to judge the living and the dead. We read of a similar promise made to the disciples in Matthew 19 and Luke 22. And this glory will assuredly come to us, just as Christ himself did truly ascend into heaven and sit in celestial glory.

### Paragraph 353

And here we must note a special thing, that Christ gives here that which in the twentieth chapter of Matthew he denies he can give to James and John, namely, to sit in celestial glory. Therefore this passage explains that. For Christ, after declaring his deity, gives what he later, in his humanity, denies he can give. This passage then proves that Christ is very God, giver of eternal life, and so forth.

### Paragraph 354

He then, in his characteristic way, adds an affirmation, by which he applies this letter to all congregations, and asserts that it is inspired by the Spirit of Christ. Concerning this, we have spoken previously.

### Paragraph 355

And we have treated thus far of the second part of this work, wherein are declared the most excellent points of our religion, namely, who and of what sort is Christ, sitting in the glory of the Father, how he is present in his church, and governs it as king and priest, by his Spirit, by his word, and by the Sacraments. Also, what and of what sort is the church of Christ? What is the true and right doctrine of the church? What opinions are wicked? What is to be done with erroneous doctrines and seducers? How may the church, having fallen and been afflicted, be restored? What is true repentance, and what are the duties of the godly, and many other things of like sort. To God the Father be praise, thanksgiving, and glory, through Jesus Christ our Lord.

## Sermon 23: ¶ The second vision is showed to S. John, wherein he sees God in his Throne with Elders, whom he describes gallantly.

### Paragraph 356

.

### Paragraph 357

After this, I looked, and behold, a door was open in heaven, and the first voice which I heard was like a trumpet speaking with me, which said, “Come up here, and I will show you things which must be fulfilled hereafter.” And immediately I was in the spirit, and behold, a throne was set in heaven, and one sat on the throne. And he who sat appeared like a jasper stone and a sardine stone; and there was a rainbow around the throne, in appearance like an emerald. And around the throne were twenty-four seats. And upon the seats were twenty-four elders sitting, clothed in white raiment, and had on their heads crowns of gold.

### Paragraph 358

The third part of this work reaches from the beginning of chapter 4 to the beginning of chapter 12. It contains a notable vision, most wholesome, and of much fruit. The first vision, which we heard explained in chapter 3, exhibits a figure of Christ and of his church, and how the Lord reigns in them, and how the church behaves or ought to conduct herself. In the second vision, John declares how, by a most just and most holy government, God governs all things through Christ, and what happens to the church in the world and of the world.[note 23.1] In these are rehearsed the most sorrowful destinies of the church, calamities, plagues, and destructions, famines, persecutions, revolts, heresies, conflicts, and other evils of the same sort. &c. Who also, and what, and how just God is in all his judgments, is described here: That he is the Author of all; that God through the most astute and excellent government of Christ rules all things; that the holy angels also and all creatures acknowledge him, and give glory to God. For so it teaches us also in all our doings,[note 23.2] and even in the very grievous calamities and persecutions, whereof it shall prophesy moreover, to acknowledge the providence and good will of God towards us, and his most just government. If we do this with quiet minds, we shall bear even the most heavy burdens patiently; we shall cease with curious questions to inquire why God permits Antichrist to arise, to increase, and to reign, to oppress the religion and saints of God? Then shall cease also the blasphemous mutinying of those who are not afraid to say, “God is indeed the Lord, he is almighty, he does what he will, and as he will; we are bond servants, and rather worse than bondmen. We are forced to bear whatever he will lay upon us,” &c. As though God were unjust, and after a tyrannical fear terrible, and ruled after a carnal lust. It is most shameful to think thus, much more to speak it. This vision shall declare that God by his providence governs all things, and that his ways are just in all his ways, and his works are holy.

### Paragraph 359

And first, John is prepared to receive this vision, yea, and we also are prepared in him. For when he had seen the door in heaven to be wide open, he heard withal, “Come up hither,” and so forth. It is surely a benefit beyond expression that the Lord opens heaven for us miserable, mortal men, and allows us to see what is done therein, or what he himself does there, and what his works or judgments are toward men. Let no man hereafter say that God does in heaven whatever he pleases, not passing judgment upon us who creep upon earth, and who also must suffer what we would not. For now he makes as it were an account of his works and, being assured, admits you as a witness to the matter.

### Paragraph 360

And here he declares with a godly voice, what John should do, and how he should behave himself. Christ bids John ascend into heavenly places, not in body, but in mind. Therefore, our mind must be lifted up into the contemplation of heavenly things, and be purged as much as possible from earthly affections, that we may behold heavenly things with a heavenly contemplation. What shall we say, but that the example of John follows immediately? And immediately I was in the spirit: that is, in a spiritual contemplation, or ravished by the spirit into the faithful consideration of those things which were shown to me.

### Paragraph 361

Now also is assembled an explanation of things that ought to be said: I will show you what things must be done hereafter. For after the scepter of God, rightly ordering or governing all things through Christ, immediately are declared the destinies of the church by seven seals and seven trumpets, within which are everywhere interwoven most comforting consolations and full of efficacy.

### Paragraph 362

And first of all, before the Seals and symbolic representations are presented, there is a depiction or glimpse of God, and his most righteous judgment and governance in all things; and this is shown throughout chapters 4 and 5, so that it might prepare us for the things that will follow in chapters 6, 7, and 8. And it may seem to others, and to human judgment, to be grievous, harsh, and unjust. And the depiction or vision was as follows.

In heaven itself appeared a seat or throne of Majesty. He who sits therein holds in his right hand a book, closed with Seals. Beside him who sat was a Lamb, who takes the book and opens its Seals. And from this throne also proceeds a sevenfold Spirit, wonderfully expressing his virtues. Before the throne appears a sea of glass, bright and smooth like crystal. The throne itself rests like a wagon upon four beasts full of eyes and wings; beneath it appear all around, encircling the throne. A rainbow like an emerald goes around it. About the throne, by a circle, appear twenty-four seats, and so many elders sitting in them, crowned and in white array. This is the order of this second vision. Their roles will be explained: what the Lamb, what the beasts, what the Elders, and the other parts did. It suffices now to have touched the chiefest points of the vision, and to have given a shadowing of some of the same.

### Paragraph 363

Secondly, we must see what each thing signifies. For a great part of the entire mystery depends on this. Regarding the manner of vision, John introduces nothing new concerning the revelation of Christ. For we read that such visions were exhibited to the prophets in a similar manner, as to Isaiah in chapter 6, to Ezekiel in chapters 1 and 11, and to Daniel in chapter 7, and so on. [note 23.9] And a throne signifies imperial majesty and judicial administration. And because the throne is not on earth, but is seen in heaven, we shall think that God’s providence and administration of judgment are celestial, sound, most holy, and entirely free from all corruption. And upon this same throne is one sitting, sitting I say, [note 23.10] not lying or standing. For God, the judge of all, is of a quiet mind, and is not moved by affections like men. There is no affection, injury, or unrighteousness to be thought upon in the universal government of all things. Elihu, in chapter 34 of Job, says: Far be wickedness from God, and iniquity from the Almighty. For he will render to a person according to his work, and according to the ways of each, he will reward them. For truly God will not condemn in vain, nor will the Almighty subvert judgment, and so on. And Arethas, Bishop of Caesarea, an old commentator, admonishes [note 23.11] that the shape of a man was not attributed to him who sits in the seat, on purpose. For although afterward mention is made of a right hand holding the book, here no shape of a man is exhibited. But he says simply one sitting, he gives him no name. The reason is ready: for God, by his nature, cannot be defined, as he who is invisible and unmeasurable. According to the manner of men, human members are attributed to him, but to be understood by a trope. Moreover, when the same God appeared to the people of Israel at Sinai, they heard a voice only, but the Israelites saw no shape, as Moses witnesses in chapter 4 of Deuteronomy. Doubtless, that they should not express the incomprehensible with an image, and should commit idolatry, the great sin and wickedness. Paul in Acts 17 denies that the deity is like the forging of men, to the Romans. He ascribes to the greatest folly idols made after the shape of men, which should represent God. Of this we have spoken elsewhere. In the meantime, two precious stones are mentioned, which by their colors do, in a way, shadow the nature of our God, and admonish the godly of greater and more excellent things. A jasper is a green stone like an emerald. Green signifies the perpetuity of God, and that he quiets and keeps all things in life. But the sardine looks with a fine color like a bright red. For God dwells in inaccessible light; he is a consuming fire, and also charity itself. For the nature of stones, read Pliny, and so on.

### Paragraph 364

But a rainbow environs the Throne round about; a rainbow is generally of diverse colors, but here it is of one color, and that of an emerald, namely green. The rainbow is a token of a perpetual grace and covenant made after the flood, as is declared in the 9th chapter of Genesis. And truly, the Throne of the high judge might put us wretched men in fear. Therefore, the rainbow puts us in remembrance of God’s grace, and that God who by his providence governs all things has bound himself in league to mankind, to whom, indeed, he wishes well. That league is still green, and always of force. The goodness of God towards men is perpetual. For though heaven should fall, and out of this Throne proceed most grievous thunderbolts, and calamities should fall upon us like a storm, yet is God in league with us, and loves us dearly.

### Paragraph 365

Concerning the throne, twenty-four seats are seen, and twenty-four elders sit in them, as senators of the most mighty kingdom of God, and fathers of the heavenly hierarchy. This number is composed of twelve and twelve. But the twelve patriarchs signify the entire people of Israel, and the old church before Christ. The Christian church was planted and sprang up from the twelve Apostles, after the incarnation of Christ, so that this number twelve encompasses the whole church of the new people. Therefore, the entire universality of saints is assembled in heaven, triumphing with Christ their king. And therefore they are clothed in white raiment, namely, purged by Christ, and pure and clean from all corruption. They are also crowned, because they have overcome and now reign in eternal glory, truly kings and priests through Christ. The description also of their behavior admonishes that in them there is nothing lacking, but that they are truly blessed; and therefore they are shown sitting, not that they are judges or judge for Christ, but because they rest from their labors and are of most quiet and pure affections, sitting with the high judge. But what do these do? They give God no counsel as to what he should do, or by what means or way he may do this or that, but they approve his judgments. For they know all his works to be just and holy, which will immediately follow. What then shall we do? Shall it be fitting for us to inquire about the judgments of God, or prescribe what he should do or not do? I think not. You do not have in this universality of saints, all patriarchs, all judges and kings, all princes, and the whole people of God; you have among these, King Solomon himself, and the most excellent and wisest princes of the world; you have the Apostles, and men apostolic, martyrs, and the wise men of the whole universal world. Will you condemn their judgments? Following therefore their example, do not busy yourself with curious questions; praise the just judgments of God, and know that the Lord is just in all his ways, and holy in all his works. To whom be glory.

## Sermon 24: ¶ Here is described the procedynge of the holy spirit, and operation, the almighty knoweledge of God, and howe the Throne of God is borne up or sustained of the four beasts, and what the beasts doe.

### Paragraph 366

.

### Paragraph 367

And out of the Seat proceeded lightning and thunderings and voices, and there were seven lamps of fire, burning before the Seat, which are the seven spirits of God. And before the Seat there was a Sea of glass like unto Crystal. And in the midst of the Seat, and round about the Seat, were four Beasts full of eyes before and behind. And the first beast was like a Lion, the second beast like a calf, and the third beast had a face like a man, and the fourth beast was like a flying Eagle. And the four beasts had each one of them six wings, and round about without and within, they were full of eyes. And they had no rest day neither night, saying, “Holy, holy, holy is the Lord God Almighty, which was, and is, and is to come.”

### Paragraph 368

Our Lord Jesus Christ, as the faithful pastor of his church, will reveal the destinies and wonderful calamities that will come upon the church. Therefore, to stop the mouths of those who question and are inquisitive about the judgments of God, and to persuade all to have patience in these storms of evils, he sets forth a treatise beforehand, wherein he shows that all things are done or permitted to be done by God through his most just providence, and are governed or ordered by the Lamb, with a judgment most righteous and holy. For whoever believes and remembers this, whatever circumstances he encounters, he submits himself humbly, lowly, and obediently to his God, and always cries out, “The Lord is righteous in all his ways and holy in all his works.” And this is the true state of the first part of this vision, which is presented in chapters 4 and 5. It is moreover most elegant, most pleasant, and most full of consolation. All things are more vividly set forth and perceived in such fitting and heavenly representations than they can be understood in bare words.

### Paragraph 369

First, a throne is described, and indeed a heavenly throne, so that in the works, providence, and judgments of God, we should not imagine anything carnal or corrupt. Secondly, he who sits on the throne is represented to us by two colors, green and red. For God is an eternal essence belonging to all their griefs or being. The same burns in love toward humankind and wills well unto man; but to the disobedient and rebels, he is a consuming fire. And the throne is surrounded by a rainbow greener than grass, comforting us so that we should not be dismayed at the sight of that godly throne, but should always remember that he who sits in the throne, judge and governor of all, is most true and keeps his promises. And to be that same league friend. Around the throne sit twenty-four elders, who have already been signified what they are, as it were shadowed, immediately in the end of the fourth chapter, and in the fifth shall be declared what they do, or what they say. Doubtless all the saints in heaven are onlookers of the judgments and works of God. For the judgments of God are not such that they should flee the light and knowledge of saints.

### Paragraph 370

Now from the throne came flashes of lightning, and thunder, and other things. In the throne are he who sits and the Lamb, that is, the Father and the Son, and from both proceeds the Holy Spirit. For by interpretation it follows immediately, which are the seven spirits of God. For the lightning, the thunderings, and other things mentioned signify, or are tokens of the Holy Spirit, which elsewhere is also said to be represented by fire, and water, and wind, and by fiery tongues. But no one will think that the Holy Spirit, which is one in substance and of the simple divine nature, should be divided into seven parts. For I told you in the first chapter how the seven spirits of God are put for the sevenfold, most full, and most perfect spirit of God.

### Paragraph 371

We have in the beginning of this vision the whole mystery of the blessed Trinity, so much as is needful for us to know, [note 24.4] and believe, and profess. There is one Seat, in that one seat are contained the sitter, the Lamb, and the Spirit: therefore there is one divine essence and nature, and thereof is one power and majesty, one rule, because there is one throne; briefly, there is one God, true and eternal, forevermore blessed, as Moses also in the sixth chapter of Deuteronomy, and all the prophets and Apostles have everywhere taught. However, in this only and undivided substance is seen a most plain distinction of persons. For there is he that sits in the throne, and the Lamb, and from both proceeds the Holy Spirit. This mystery of the Trinity we profess in the Creed. This appears openly in the incarnation of our Lord, while the angel says to the virgin, “The Holy Spirit shall come upon thee, and the power of the Most High shall overshadow thee; and that which shall be born of thee shall be called the Son of God.” Likewise, in the baptism of Christ, a voice is heard from heaven upon the Lord: “This is my well-beloved Son.” The Holy Spirit also appears in the likeness of a dove. Whereupon the Lord commanded us also to be baptized in the name of the Father, and of the Son, and of the Holy Spirit. This profession is certain and true, and so set forth by the most manifest scriptures and lively preaching of the Apostles, like as Tertullian declares against the heretic Praxeas. We ought rather believe and cleave unto these things than to the monstrous and blasphemous Spanish sophistry of Servetus, a man most corrupt.

### Paragraph 372

But especially here the entire mystery of the Holy Spirit is declared to us, and that in few words, which in the Gospel of John is expressed more fully. First, his procession is noted, which in times past was rashly affirmed to be set forth in no part of Scripture. John here says, “out of the throne proceeded lightnings, &c.” And immediately after: “which are the seven spirits of God.” This word in Greek signifies a proceeding or going out, but John says that it proceeded or went forth. Therefore, the ancient council of Constantinople rightly decreed: that is, concerning the Holy Spirit, the Lord, that quickener, proceeds from the Father, &c. But because the Lord himself in the Gospel, speaking of the Holy Spirit, says, “he shall glorify me: for he shall take of mine, and shall show it unto you. All things whatsoever the Father has, are mine. Therefore I said, that he shall take of mine, and shall show it unto you,” no one will understand the Spirit to proceed from the Father only, and not also from the Son, concerning which there was also a lengthy contention between the Greeks and Latins. For if he proceeds from the Father, he proceeds from the Son also. For even for the same cause at present he is shown to proceed out of the Throne. But in the Throne is not only he that sits, but the Lamb also, of whom in chapter 5 we shall add, that the Lamb has seven eyes, which are the seven spirits of God, sent into the whole world. Albeit therefore that in John 15, the Holy Spirit is said to proceed from the Father, yet it is stated beforehand: “who I (says the Son) will send unto you from my father.” To be brief, if there is one substance and nature of the Father and of the Son, I do not see how the Holy Spirit should proceed from the Father without also proceeding from the Son. Let us rather leave these scrupulous disputations to idle wits; let us believe that the Spirit proceeds from both.

### Paragraph 373

Moreover, the power or effect and operation of the Holy Spirit is here also set forth and declared clearly. [note 24.7] For first, he illuminates, when he shines upon the obedient, and terrifies the rebellious with severe threats. Secondly, he governs, when he contends against this ungodly world and rebukes it for its sins, thundering out God’s terrible judgments. Two Apostles in Mark are called the sons of thunder, or thunderers. He utters moreover wholesome voices of doctrine, exhortation, and consolation, by men, for the favor of men. Finally, where the operation of the Holy Spirit cannot sufficiently well be expressed, yet by the seventh number he encompasses and accomplishes his fullness, and says that seven fiery lamps are burning before the Seat, burning, I say, not quenched, or smoking. For the grace of the Holy Spirit is bright and full of efficacy, as was spoken of before: and where these things are found in the Throne, how should any man think that the judgments proceeding from thence should be in any part corrupt, defiled, or to be blamed? By the Holy Spirit all things are preserved, and by his providence all things are wrought.

### Paragraph 374

Hereunto is added another thing, a glassy sea before the seat, in clearness and brightness representing crystal. This signifies this frail world, which is subject to God, and as it were in his sight. Also, in other places of holy scripture, because of its instability, tossing, and turmoil, it bears the figure of this variable and most unconstant world. Certainly, the state of this world is more fickle than glass. Something of this shall follow in the 15th Chapter. But whatever things are done in the world through a marvelous variety, all the same shine as in a glass before the Throne, so that God sees them all as it were in a crystal: whose eyes or knowledge the least things that be can not escape. For we shall not think that such things as are done in the world are done rashly, and by a certain fortune to hap or chance besides the knowledge of God, or to be of God unknown.

### Paragraph 375

After this, he returns again to the throne, intending to complete what he had begun to describe, and to show that all the works of God, done by his creatures, are most holy. Royal seats, chairs, or thrones of kings are typically borne up and adorned with beasts, as Solomon’s seat was with lions, as can be seen in the third book of Kings, chapter 10. In other places, the most excellent beasts draw the triumphant chariots of princes. In the same manner, by a figure of speech, beasts are assigned to the throne of God. For God is carried upon Cherubim in his prophets. And Ezekiel, in chapter 10, names Cherubim, beasts, and the entire text demonstrates that the passage must be understood as God’s chariot, drawn by beasts, in which he himself was carried out of the city of Jerusalem. There is much mention in poets of the chariot of the Gods, taken happily by the first writers from the holy scriptures. For Satan, the ape of God, always seeks to defame the word of truth. But we, omitting the triflings of poets, will consider the sober description of this carriage of God, or rather of God’s throne. Almighty God sits in this seat. Sitting in the Scriptures signifies government. Here it is signified that God sits in all his creatures; that is to say, governs his creatures, and by his most wise providence works all in all, using every creature according to his good and just pleasure, according to the nature of each. We shall say then that by those beasts are understood all the creatures of God, dispersed throughout the four quarters of the world, which are comprehended in the whole world.

### Paragraph 376

And first, it is shown in what place of the throne the beasts were [note 24.12]: namely, in the midst of the throne, and in the circuit of the same. You will ask, if they are in the throne, how should they be about the throne? If they are about the throne, how are they in the midst of the throne? The thing must be conceived as I admonished before, that we understand that under the throne, the midst of the beasts does with their hind parts reach to the midst of the throne inwardly, and so as it were to have borne up the throne. And with their fore parts, I mean, with their breasts, and heads and wings, to have stood forth, and so to have compassed the throne, and as it were environed it round about. For so they might seem to be in the midst of the same throne, and round about the same.

### Paragraph 377

After,[note 24.13] the manner of beasts is described diligently, and they were four in number. For in times past also the number was expressed by Ezekiel. And the parts of the world are suitably signified by the fourth number, comprehending the universality of things. Some here have suggested the four monarchies of the world, &c. And every beast had his face and his body, six wings, and was full of eyes within, as also their bodies were full of eyes. The first was like a lion in shape and fashion, the second like a calf, the third like a man, and the fourth like a flying eagle. By these appear to be signified all creatures, visible and invisible, reasonable and unreasonable, and the most excellent. For afterward in chapter five we shall hear that all creatures jointly worship the Lamb and him that sits on the Throne. And truly, God uses them all—the sun, the moon, the stars, the air, the fire—and briefly, all living things. And such creatures as he has chosen, in order to work something by them, he makes them efficacious, instructing each according to their state and condition, that they should lack no wisdom, reason, strength, power, patience, labor, quickness, nor swiftness. The face of man signifies wit and wisdom, as also the eyes signify foresight, watchfulness, subtleties, and keenness in doing things. The lion’s face betokens force and strength, and boldness or magnanimity. As the sight of an ox or a calf betokens enduring of labor. The eagle and the six wings signify swiftness. For example, God chose the Assyrians or Babylonians, who should destroy Nineveh. These therefore, as it is in Nahum, the Lord prepared and furnished, that they were swifter than eagles, and the rest as you may read in the first and second chapters of Nahum. And so all creatures are ministers of God’s judgments, coming out of his judicial throne.

### Paragraph 378

Then it is also addressed what those beasts do. They go around the throne, always waiting for God’s command, that they may apply it cheerfully, speedily, and stoutly. Neither do they have any rest—mark how he says “have,” not “shall have” or “have had,” but “have”—any rest; that is to say, they are in continual service of God. But let us not understand from this that they are burdened with any painfulness. And also they honor God with continual praise. Arethas: it signifies, he says, no laborious thing. And they have no rest, but a continual occupation, about the singing of godly praises, and so forth.

### Paragraph 379

Finally, here is also presented the hymn and praise of all creatures. In old times David sang, “Praise the sun and moon, &c.” The same hymn is found in Isaiah 6. And what do all creatures commend in God, whose service God uses, and whose power and operation they feel? Primarily, holiness. These things chiefly concern the substance of the matter. For they teach God to be holy, unspotted, just, good, omnipotent, doing all things, eternal, the beginning of things, and preserver. For they say, “Holy Lord God Almighty, which was, &c.” Which words we indeed explained in the first chapter. Who would not gather from this that the works and judgments of him are most holy and just? Who therefore shall hereafter reprove the judgments and works of the Lord? Just is the Lord in all his ways, and holy in all his works. This testimony of all creatures makes us willing, ready, cheerful, and untroubled, that we should willingly quiet ourselves in the judgments of God, and murmur at him in nothing, questioning why he should do this or that? But wholly submit ourselves unto God, believing all his works to be good, and to be done for the profit of the godly, and for the most just punishment of the wicked. Holy is God the Father, holy is God the Son, and holy is God the Holy Ghost, holy is one God in Trinity, blessed forevermore. Holy are all his works, and his ways undefiled. And we read more correctly three times holy than nine times, following the example of the Complutensian book. For the former reading, the prophet Isaiah approves. To God Almighty be praise and glory.

## Sermon 25: ¶ Here is declared what the Elders did about the Throne, and how they sang unto God a song of praise.

### Paragraph 380

.

### Paragraph 381

And when these beasts gave glory and honor, and thanks to the one who sits on the throne, who lives forever and ever, the twenty-four elders fell down before him who sat on the throne, and worshipped him who lives forever, and cast their crowns before the throne, saying, “You are worthy, Lord, to receive glory and honor and power, for you have created all things, and for your will’s sake they exist and were created.”

### Paragraph 382

This most godly vision, well and rightly understood, and held in faithful memory, instructs us rightly in judging rightly the works of God, that we should fear God, be patient, and submit ourselves wholly to God, and give all glory unto him. For this is the very fruit that comes unto us, and the end of all things that here are spoken.

### Paragraph 383

And by repeating what the beasts did, he indicates what the twenty-four elders did [note 25.2]. By this, we are clearly taught what we owe to God, and how we should judge his works, and how we ought to conduct ourselves in this matter.

### Paragraph 384

Those beasts, that is to say, the whole number of creatures, whose ministry God uses in the government of things, ascribe three things unto God sitting, that is to say, ruling and governing all things, to God I say living forever, that is to say, eternal, living, and giving or inspiring life into all things. First in deed glory, [illegible], which is a majesty, or great estimation, a reputation, worship, or good opinion: when we think well of God, protesting that there is nothing better than he, greater, more worthy, more just, more holy and more excellent. This glory are we always commanded to give him, and to esteem nothing in this world dearer and more precious than God. Secondly, they give to him honor [illegible]; in Greek, this signifies honor and price, and the due and bounden duty that we owe to any. We owe unto God reverence and submission, as to the supreme good, and the only and true lord of all. Paul in Romans 13, speaking of obedience due to the magistrate, says: “To whom you owe fear, give fear. And to whom you owe honor, give honor.” In the third place follows benediction, which he called [illegible], that is, thanksgiving, and praise. For we are commanded to praise all the works of the Lord, and to give thanks for the same. Job is said to have blessed or thanked God for the most grievous affliction that he sent him. For he said: “Like as it pleased the Lord, so hath it been done: the name of the Lord be blessed.” Whilst the beasts do attribute all these things to him that sits on the throne, by their example they teach us what we should do, namely, to give all these and singular things unto God. Which if we do, all murmuring shall cease, and disputations commenced of searching and examining the works of God through curiosity. With the laud and praise of the beasts is joined the hymn or song of the twenty-fourth [note 25.5]. Elders. This is the church triumphant, the company of all Saints, Patriarchs, Prophets, Apostles, Martyrs, &c., as I declared to you before. Mortal men have not here an example of some one saint, or wise man: but of all holy, godly, wise and worthy men. They have put off their flesh, and want affections and errors: they are therefore of uncorrupted judgements, so that there can be no more clear or pure examples ministered to us. Three or four things are taught us concerning these Elders, which they did or performed, not to every body, but to him that sits on the throne, and lives forever and ever. For so are the titles of God repeated, whereof is spoken before. We told you also that the seats of the Elders were set round about the Throne, in which they sat clothed with white raiment, crowned with crowns of gold, living with him that lives forever.

### Paragraph 385

They first rise from their seats or chairs,[note 25.6] and fall upon their knees or face before God. And in falling or kneeling, they show a submission and humility of mind, that we might learn with great humility and reverence to submit our souls and bodies to our God, submitting, I say, ourselves and all our things to his good will and pleasure. But if the blessed souls, now purified and already enjoying the sight of God, fall down before the Lord, what should a wretched, miserable, mortal, and sinful man do? Let him be ashamed of rebellion and slothfulness, which sees such great submission in the most noble and godly souls of heavenly dwellers.

### Paragraph 386

Then the saints worship, and worship in deed none other, but him that sits on the throne, and lives forever, the Father, the Son, and the Holy Ghost, God three and one, everlasting and almighty. Therefore let us also worship this God, following the example of all saints. We worship God with external adoration, if we uncover our heads, kneel, and bow before him. In spirit and truth and with inward worship, if we depend wholly on him, consecrate ourselves wholly unto him, and wholly look upon him as one, the only, soul, incomprehensible, most wise, best, and greatest, most righteous and most merciful. And they that thus fall down before the throne of God, and so worship him, do not contend with God about his works; they do not expostulate with God impatiently, why he does this, and permits that?

### Paragraph 387

To all these things is added that they pluck the crowns from their heads and cast them away before the throne, in honor of him who sits there. This is not only a remarkable modesty but also an humble humility without parallel. Primasius, an interpreter of the Apocalypse, affirms that whatever virtue and dignity they possess is truly assigned to God. For to him is rightly attributed whatever is won or obtained; from him, the one who overcomes is aided. Thus he says. They also testify and signify that they would not assume any godly power, that they would not reign, that they would not, as counselors of God, give counsel to God or prescribe even the smallest thing to him, but to submit to God all power, all rule, and the entire government, themselves and all others to be governed. For they have experience and see no one in the world, neither in heaven nor on earth, who is wiser, mightier, greater, or who governs all things more faithfully, more diligently, more safely, and better. Let us therefore, O brothers, submit to the judgment of the saints, and let us consent with them in all things.

### Paragraph 388

Yes, and with clear words they testify why they cast off their crowns. Not that they are ungrateful to God, and do not value his gifts highly; but because they openly acknowledge that all glory is due to him alone. Therefore, they are in complete agreement with the beasts and all of God’s creatures, and, singing a hymn to the highest ruler, they confess that he is worthy to receive glory. And he receives it, not as if he lacked it before, but so that it may seem unworthy if either they or any other creature were to claim as their own those things that belong to God alone. These things belong to no creature. And they praise God highly, whom they call their Lord and God. Some copies add, “who is holy.” For they agree in all things with the beasts, which also cried, “Holy, holy, holy, Lord God Almighty.” To him they also gave glory and honor, as mentioned before. So also the elders ascribe the same things to him now. Especially, they attribute power to God and take it from themselves. Why, then, do the Papists attribute power and activity to the saints in heaven? Which, nevertheless, here they plainly attribute to God alone. John and Peter, while they lived, did not take kindly to the people seeming to attribute some godly power to them. For when they had healed a man who had been lame before the temple, and the people were amazed, they said, “Men of Israel, why do you marvel at this? Or why do you look at us as if through our own power or holiness we had brought about this healing? The God of our fathers has done this.” But how much less should we now think that, being delivered from all corruption, they would desire any godly power to be given or divine honor attributed.

### Paragraph 389

They also add or give a reason why they submit themselves, and all that they possess, to God, and attribute to him glory, honor, and power. For you have created all things, and by your will they are, and were created. This glory of God is wonderful and immeasurable. How great you are, and that all power and glory are due to you, appears from the making and creation of the universal world. No one was with you at the creation, no one gave you counsel about what or how you should do, no one helped you in the slightest. Who then should approach you to share in power? Who should glory before you, God and maker of all things? You alone made all things, alone preserve all, and alone govern all. You will, and they were made; you said, and they were created. It was enough to have said, it was enough to have willed. Indeed, all things today have their being through your will, without any effort or toil on your part. You govern all things in the best and most excellent order. This testifies to the wonderful course of the stars, the pleasant change of things, the most sweet and plentiful fruits. Who then would not gladly submit himself and all that he has to you and to your government, who would not commit all his possessions to you? Who would not acknowledge the power and glory to be yours? Let us mark these things with attentive minds, that we may also appear before God as we see the saints in heaven appear. God grant us this.

## Sermon 26: ¶ Of him that sits in throne, and holds the book in his right hand sealed with vii. seals: What that sealed book is.

### Paragraph 390

.

### Paragraph 391

And I saw in the right hand of him who sat on the throne a book written inside and on the back, sealed with seven seals. And I saw a mighty angel proclaiming with a loud voice, “Who is worthy to open the book and to loose its seals?” And no one in heaven, nor on earth, nor under the earth, was able to open the book and to look upon it. And I wept greatly because no one was found worthy to open and read the book, nor even to look upon it.

### Paragraph 392

He now proceeds to describe more fully him that sits on the Throne: of whom he had touched certain things before. This section contains considerable weight for our matter. For now he will show that which is principal in this treatise: that all things which are done in the world through God’s providence are most justly and holily governed by Christ. All the saints of God, and all creatures, acknowledging this, for an example to us, do praise and celebrate him that lives forever.

### Paragraph 393

And it shall behoove us to weigh every word, since in each one are great mysteries, and nothing is spoken in vain. And truly, God Almighty sits in a throne. And by sitting is signified not only the power of judging, ruling, and governing, but also a quiet mind (not troubled with any evil affections, after the manner of judges of this world) and great equity in all things. Secondly, a book is seen in the right hand of him that sits, of which book we must speak more at large.

### Paragraph 394

Here appears an allusion, as there is in many other places of the Scripture, to the princes of this world, who have books of laws, of privileges, of institutes—regarding what has been done and what is to be done—finally, of secrets, of acts, of the condemned, and of citizens, of life and of death. For so is both the book and books assigned to God. Moses says in Exodus 32:33, “Put me out of the book of life,” and so forth. In the Psalms there is much mention of these books of God: in Psalms 56, 69, and 139. In the seventh of Daniel, books are opened, of which mention is also made in the twentieth of the Apocalypse. We read in Malachi 3 of a book of remembrance before God. Therefore, this book of God contains all the counsels of God, all his works and judgments. For we shall hear by and by [note 26.2] that all things that are done in the world come out of this book, as it were out of a fountain or wellspring.

### Paragraph 395

And three things are chiefly spoken of this book. First, that it lies not in the Throne, or in the bosom of him that sits, or under the Throne, or that it hangs before or behind the Throne: but it is in the right hand of God. Hereby is signified the operation or power of God, and the same most just and most mighty. For the book is not seen in the left hand. God therefore works, and contains or ministers all his works and judgments most holily. Secondly, that book is written within and without, or on the backside. For in God’s providence and judgments, all things are contained, both good and evil, fortunate and unfortunate, sharp and gentle, sweet and sour, visible and invisible, secret and apart, and all things in general.

### Paragraph 396

Finally, the book is sealed with seven seals. For it is most strongly closed and fastened. [note 26.4] For the judgments and works of God are firm, true, just, and such as cannot be withstood. The use of seals among men is diverse, yet it may be considered in two points. [note 26.5] First, seals are set because of fidelity, truth, and righteousness. And great deliberation is had in setting the seals. For they are not put to unjust, vain, or false matters. Therefore, seals are tokens of certainty and testimonies of a right. It seems an unworthy thing to speak against sealed writings. By the seals therefore that are set to the book of God is signified that the judgments and works of God are most firm, true, and just, whatever is done by his providence, and are ordained by Christ. It shall therefore be a shame to find fault with the judgments of God, or to speak evil of his works. Again, by seals, secrets are kept, that they be not seen by every man, but by them only to whom they are appointed. The judgments therefore and works of God are for the most part hid, and not open to all men, saving to such as the Lord hath appointed, namely to the faithful and obedient. [note 26.6] But there are only seven seals, for in them the fullness of times, and of things to be done in these times throughout the world and church, and of the judgments and mysteries of God are comprehended.

### Paragraph 397

Now therefore, the opening of the book [note 26.7] and its unsealing, are nothing else but the revealing of God’s judgments and the declaring or uttering of his most secret counsels. Finally, the most holy and just operation, dispensation, and execution of his will. Nothing in that opening is done against the truth, faith, love, and justice of God.

### Paragraph 398

And with many words, and also most diligently and well is treated here of the opening of the seals, who truly might be thought worthy to open to the church the secret judgments of God, and to execute and minister his holy works: that is to say, to whom the kingdom is given and government of the divine providence. For an Angel, and that not of the common sort, but a strong and worthy one, with a loud voice cries, to make us all attentive, and that we should note diligently, who he is that should both open the book and unloose or undo the seals. And he holds the hearer, beholder, or reader in suspense before he will show him, to the intent truly to commend him to us exceedingly. No man, says he, in the whole universal world, neither among the angels and saints in heaven, nor among earthly men and under the earth, was found who could either open or unseal the book.

### Paragraph 399

Let us observe that no one can open the book and its seals except Christ alone. Why then is the administration of things attributed or communicated to saints? No one can reveal to us the counsels and judgments of God, nor can any man govern those judgments and works of God that he works in the world, save only Christ the Lord. Why then are such great benefits sought from saints, and credited to them, if either the sick are restored to health, or if a mortal man receives any other gift or benefit? Many will say, “I received this as a gift of God, but through the meditation, power, and merit of this or that saint, to whom God granted authority to rule over such a disease and might heal those who call upon the name of the saint, or the name of God by the saint.” These things are now refuted by the words of the Lord and Saint John, who say that no one in heaven or on earth is found who could open the book. Yet nevertheless, about the throne sat the twenty-four elders, representing the type of all saints in glory, and not one of them was found who could open the book. Therefore, they are greatly mistaken who attribute the governance of things in the church to the Pope, a most corrupt and filthy man. Only Christ received all power in heaven and on earth, as we shall presently understand more fully.

### Paragraph 400

John wept, for he understood that a weighty matter consisted in the opening of this godly book; and yet he saw no one at all who could either open or unseal it. Neither did he as yet fully understand the matter. And he bore the figure of those who understand not the judgments of God, nor know that all things are through God’s providence holily governed by Christ. For in them nothing else remains but mourning and heaviness. Certainly, without Christ and his opening, whereby he reveals to us the divine mysteries and judgments, no man can rightly judge of the same. For unless we understand the seals to be opened by Christ, and that all things are done by his order, which loved us and gave himself for us, what thing shall be left in us but sighing?

### Paragraph 401

He presented three things, to open, read, and look upon. No living person can open [it], for no person is fit for so great a responsibility, except the Son of God. No person reads or fully understands God’s judgments, but the Son, and to whomever he has revealed as much as any person has. No person can look upon [it], that is to say, can behold the works and judgments of God, but will be offended, unless he is endowed with the Spirit and purified by it. Therefore, we must ask grace from him, that we may understand so much of God’s judgments as will suffice, and may judge them well.

### Paragraph 402

Arethas, Bishop of Caesarea, an expositor of this book, says that neither those who lack a body, nor those who are embodied, nor those who have departed, leaving their bodies behind, have received a perfect understanding of divine matters. And immediately afterward: neither was there anyone who could open it, nor even one who could attentively consider it—that is to say, could not carefully observe the judgments of God. And the context of the entire passage sufficiently proves that Saint John speaks here of the judgments truly, but especially of the governance of things. The Lord Jesus be glorified forever. Amen.

## Sermon 27: ¶ Here is lively described the Lamb in the throne of God, receiving the book of the hand of him that sits and opening it.

### Paragraph 403

.

### Paragraph 404

And one of the Elders said to me, “Do not weep.” Behold, the Lion from the tribe of Judah, the root of David, has prevailed to open the book and to break the seven seals thereof. And I looked, and behold, in the midst of the throne, and of the four living creatures, and in the midst of the Elders, stood a lamb as though he had been slain, which had seven horns and seven eyes, which are the seven spirits of God sent into all the world, and he came and took the book out of the right hand of him who sat on the throne.

### Paragraph 405

Since John had wept, that no one was worthy even to look upon the book belonging to the one who sat on the throne, one of the twenty-four elders comforted him. His name is not expressed, and it seems vain and curious to seek it. Nevertheless, some commentators suppose him to be the patriarch Jacob, since shortly after his oracle or prophecy is recited. And so the author proceeds in a most orderly fashion to the description also of the Son of God, by whom the heavenly Father, as all Scripture affirms everywhere, governs all things. Hitherto he has described the one who sits on the throne, and before that, the Holy Spirit. Therefore, these are wholesome and most profitable doctrines for the church, by which the true faith is confirmed.

### Paragraph 406

The comfort of this elder, and truly the heavenly and most godly doctrine aims at this: that we should understand that all the complaints, weeping, grudging, and the various distresses of our mind cannot be quenched, appeased, and quieted unless we see and believe that to Christ (as is most plainly and manifestly set forth here) has been given by the Father all power in heaven and on earth. Therefore, to be constituted both the only Redeemer and also the head, Prince, and governor of all, which under the seal of faith and truth should govern all things that are by God’s providence ordained, and even now despises them, and reveals to us so much of God’s judgments as suffice us. If we believe this with a faithful and sincere mind, we shall have quiet consciences in all the works of God, even those that are hard to endure and seem to some men most unreasonable. For we know that he by whom all things are governed is of our nature and kind, yes even our own brother. And he who truly favors us with all his heart has suffered death for us, and loves nothing better in all the world than humankind. Moreover, he who has overcome death, the Devil, and Hell, and has overcome them for us. Who will now suspect his government, permission, or operation? You have a brother in the Prince’s Court, whom you are assured will favor you from the bottom of his heart. You hear how he has given to him of the Prince the government and judgment of the whole country; would you hesitate or be reluctant to submit yourself to him? No, rather you trust and hope to obtain anything from your brother.

### Paragraph 407

Therefore let us remember how the Scripture does not teach this only, but everywhere, that Jesus Christ, the Son of God, and in truth of the same substance with us after his humanity, in dying for us, deserved to have a name given him, which is above all names, and that all things should be subject to his rule, whatever may be in the visible or invisible world. For so John testifies in the first chapter. And Paul also to the Philippians, Colossians 1, and to the Hebrews in the first chapter. He is said at this present to have overcome and obtained to open the book and loose the seals thereof. Therefore, by knowledge of him and through faith in him, we obtain that with a joyful mind we may look upon the book, the judgments, and all the works of God, and quietly and patiently bear the opening thereof and the rule of all together.

### Paragraph 408

But so that we may judge more rightly of Christ as ruler of all—although he has already depicted him vividly—now he proceeds to portray him in his most godly and excellent qualities, so that we should not be afraid of his rule, nor unwilling to submit ourselves wholly and quietly to his governance.

### Paragraph 409

First, it is said that a Lion of the tribe of Judah has overcome,[note 27.4] namely, that same Christ of old: to have overcome the Devil, sin, death, the world, hell, and all power of the adversary. And he overcame in dying, and so achieved the highest dignity, and was made Lord of all. The Devil is also called a lion by Saint Peter. Solomon and the Prophets call tyrants lions.[note 27.5] Our author therefore calls Christ a lion, not of the common sort, but of the tribe of Judah. For he alludes to the prophecy of the patriarch Jacob, which is in Genesis 49. He prophesies there that Shiloh shall come, with plenty and good fortune, which like a lion that has taken its prey, neither is there any man that can drive him from it, can defend those that are his, whom he has caught out of the dragon’s claws, so that no hostile power dares once hiss against him. Christ therefore is declared a victor or conqueror greatest,[note 27.6] most mighty, and most invincible. Which belongs to him alone. Yet you shall find kings who are every hour overcome by wicked lusts, which will suffer themselves to be called invincible. Briefly, this first note in the description of Christ shows that Jesus Christ, governor of all, is that very same whom the patriarchs and prophets have prophesied to come into the world, a prince most invincible.

### Paragraph 410

Secondly, Christ is called the root of David, wherein he appears to have alluded to that saying of Isaiah in the 11th chapter. Then shall a bud come forth of the stock of Jesse, and a flower shall ascend out of the roots thereof. Namely, Mary, the daughter of David, from whom that most sacred flower, Christ, sprang and came, was the stock of Jesse. And of the very roots of David, or of the virgin—I mean of the most true human nature—Jesus Christ was born very man into the world. For he took not the angelic nature, but the seed of Abraham. He is therefore our brother, of the same substance with us, after his humanity. These things do comfort us exceedingly and confute heretics strongly, who deny that Christ has a very human body. We have more hereof in Matthew 1 and Luke 1:2-3. After it is expressly spoken of the same our Lord, that he is in the midst of the throne, in the midst of the four beasts, and in the midst of the twenty-four Elders: and is therefore excepted out of the number of creatures, out of the number of angels, and out of the number of saints. For he is greater than these, to wit, of the same substance with the Father, in glory and power equal. For the Father is in the midst of the throne, from thence proceeds the Holy Ghost: even there is found also now the Lamb, Christ, not only very man, but also very God. And is a distinct person. For the blessed Trinity knows no confusion. The Father is God, the Son is God, the Holy Ghost is God: yet are all three but one God, the Father in his subsistence, the Son in his, and the Holy Ghost in his, not making three Gods, but three properties and persons in one indivisible and eternal essence. And whereas Christ is mentioned to be in the midst of the beasts and in the midst of the Elders, he is doubtless signified after the divine nature to be everywhere, to be the life and preservation of all creatures, also in the midst of his chosen and of his Church. Therefore, like as we believe Jesus Christ to be very man, so let us also believe him to be very God, of the same substance with God the Father. Therefore let Servetus perish with Arius and Mahomet, and as many as deny Christ to be the Son of God, coequal with the Father in all things. Furthermore, he is now also called a Lamb, not that he is a sheep by nature, but for that by a lamb is prefigured the innocent redeemer of the world and the only wholesome sacrifice of all the faithful. A lamb is a token of innocence, and from the beginning appointed for sacrifices. Abel offered up a lamb; after the law was offered a daily sacrifice, in the morning a lamb, and at evening a lamb. For Christ is the expiation of those who were in the beginning of the world and who in the end shall be. The Paschal lamb in Exodus 12 whose blood prohibited the destroying angel from the houses and tents, represented the figure of Christ, by whose precious blood we are reconciled to God. This exposition of the Paschal lamb, Jerome, Peter in 1 Peter 1, and Paul in 1 Corinthians 5, have brought. Isaiah accords with them in the 53rd chapter. And so expounded by the Apostle Philip in Acts 8. Finally, John the Baptist, which with the finger stretched out and pointing to Christ, exclaimed: “Behold the Lamb of God, which takes away the sins of the world.” Let us therefore believe that the same Jesus Christ, unto whom all power is given of the Father, is our deliverer, our expiation, reconciliation, innocence, sanctification, justification and everlasting salvation: as he whom we shall hear in the 13th chapter to have been slain from the beginning of the world, for so much as his only death and one oblation made from the beginning of the world, and continually to the world’s end, doth sanctify all those that are sanctified, which the Apostle also affirms in Hebrews 10.

### Paragraph 411

Nevertheless, this Lamb or Savior of the world is said to stand in the midst of the throne, truly because now he executes the office of a universal king, priest, and governor, always ready and prepared to save. Thus, Saint Stephen also in Acts 7 sees him standing. Or else, in other places, we read that Christ sits on the right hand of the Father. But this passage does not contradict that, considering that to sit signifies both rest and reign.

### Paragraph 412

Moreover, this our Lamb appears on the throne of divine majesty [note 27.12] as if he were slain, not because he was not truly slain and dead—for that is described in detail shortly thereafter—but because he did not remain in death, but rose again on the third day, so that he might declare himself to be the life and resurrection of the faithful. Or, indeed, because after his humanity he is seen to be slain, but after his deity, he is immortal and subject to no reproach. Therefore, in the old law, one goat in Leviticus 16 is slain, but the other is not killed, but is led forth into the desert by the work of a man appointed for this purpose. Nevertheless, some interpreters explain it in this way: he is said to be as if he were slain, because, following Saint Chrysostom and Saint Augustine, he has reserved the scars of his death’s wounds as a token of his victory, &c.

### Paragraph 413

Furthermore, this Lamb, Christ Jesus our Lord, has seven horns, not that in deed he carries so many horns like a goat of Judea. An horn, as appears by Daniel and by the song of Zacharias in the first chapter of Luke, signifies power and kingdom. The number seven is the number of fullness. It is therefore signified that Christ is endowed with every kind of power, divine, human, imperial, pontifical, royal, briefly, most absolute. In the thirteenth chapter, we shall hear that the beast has taken to him two horns, as it were of the lamb, whereof I shall speak in its place. Daniel in the seventh chapter says, “And rule was given him, and honor and kingdom, that all nations and tongues might worship him, whose rule is an everlasting rule, which shall not perish nor decay at any time.” Now he has seven eyes also. These he expounds, and says [note 27.14], which are the seven spirits of God, sent into the whole world. I showed you before that the seven spirits are called a sevenfold spirit. Here therefore is signified the fullness of the Spirit, which the Lord pours out upon all flesh. Here is signified the universal knowledge of the Son, in whose sight are present whatever things are done in heaven and on earth, openly and privately. For the Spirit of Christ, that unmeasurable force, incomprehensible and most divine, searches and pierces all things; nothing is hidden from his eyes, which view the whole world.

### Paragraph 414

And such is Christ, as we have heard described hitherto,[note 27.15] whom the patriarchs have before said should come, a victor and triumphant conqueror alone, truly invincible, very man of our own substance, and also our very brother, yet very God nevertheless, of the same substance with the Father and the Holy Ghost, the reconciler, redeemer, and the only salvation of the world: he has suffered for us, and the same has risen again from the dead, and ascended into heaven, having all power in heaven and on earth, who sees all things, communicates his spirit unto men, and is the most faithful keeper and defender of all mankind. This Christ Jesus our Lord came and received; he did not convey or steal it away, but took that book of divine providence, of the judgments of God, of the universal government of all things, that he might open and loose the Seals thereof: that is to say, that he might reveal to us who are redeemed with his blood the judgments of God, and might dispose and order all things in heaven and on earth. Therefore, if we know that the governor of all things is given to us as redeemer, King, Bishop, and our only salvation, who will from henceforth willingly submit himself to his government? And saying we now understand certainly how that under the seal of faith and truth all things are done by Christ, who dares hereafter more curiously inquire of his works and judgments, unto whose credit and government we should now commit all things,[note 27.16], in case they were in our power? Nevertheless, we shall observe that the Son does not so receive these things of the Father, that the Father is deprived thereof. For in the fifth chapter of John’s Gospel, the Lord says: “My Father works until now, and I work,” and so forth. Certainly, the Son is called the Word, mouth, and arm of the Father, and so forth, lest after his humanity the Son should seem less than the Father. For very godly Arethas says, where the Lamb received the book from the right hand of him that sits on the Throne, it must be understood on behalf of his humanity, as also that he was slain. For concerning his deity, none of all those things that may worthily be spoken or thought of God is severally assigned to three persons, saving the manner of bringing forth, of him that begets and of him that is begotten, and of him that proceeds, and so forth.

### Paragraph 415

This description of Christ is singular, most excellent, very evangelical, and full of consolation; therefore, it should chiefly be stored in the depths of our hearts. We also see that those who were not afraid to say that, besides the apostolic manner, few things were taught in this book concerning Christ and our redemption, were deceived in their judgment. Let us pray to the Lord that he would deign to illuminate our minds. Amen.

## Sermon 28: ¶ Here is described adoration and praise geuyng or an Himne song unto Christ of the beasts and Elders.

### Paragraph 416

.

### Paragraph 417

And when he had taken the book, the four living creatures and the twenty-four elders fell down before the Lamb, having harps and golden vials full of fragrant incense (which are the prayers of the saints), and they sang a new song, saying: “You are worthy to take the book, and to open the seals thereof: for you were slain, and you have redeemed us by your blood, out of every tribe and tongue, and people and nation, and you have made us to our God, kings and priests, and we shall reign on the earth.”

### Paragraph 418

We have heard that the Lamb has received the book from the hand of him who sits on the throne, so that he might open it and break the seals of it; that is, we understand that Christ is the only and eternal Savior and Lord, to whom all power is given in heaven and earth; that he alone and forever saves, that he reveals to us the mysteries and judgments of God, and that he finally governs and disposes all things in the world. It follows moreover how all the creatures of God behaved themselves toward this Son of God, the monarch and governor of all things. This is set forth with a marvelous figurative and plentiful speech in the vision of the four beasts, twenty-four elders, and so on. Certainly, that we might from their gestures, words, and works understand what is fitting for us to do in the judgments of God.

### Paragraph 419

For this example is truly manifold, and even consists of six parts, such as you will scarcely find presented in any other matter. [note 28.2] And this matter is of very great force. First, we have heard in the fourth chapter that the four beasts cried out before the throne of him who sat: “Holy, holy, holy, Lord God Almighty.” Second, we understand that the twenty-four elders fell down, worshipped, and cast away their crowns, and sang a hymn. Now follows the third degree of this example. For as first the beasts and elders did these things separately, so now jointly with one accord, the beasts and elders fall down together before the Lamb. Let us therefore fall down also in all the judgments and works of God, before the Lamb, governor of all, and let us worship. For although it is not here added, “and they worshipped,” they are to be understood to have fallen down for the purpose of worship, for to fall down is to worship. This is also evident from what follows. For they offer prayers to the Lamb, that is to say, sing a hymn, which is a part of godly worship. Moreover, it follows immediately that every creature sang a hymn to him who sits on the throne and to the Lamb, and so forth. And truly, two things especially and diligently Jerome treats in this example. First, he depicts the behavior of the beasts and elders. After that, he adds the hymn, praise-giving, or song. And so much as pertains to their behavior: before all things, they fall down before the Lamb, as I said even now.

### Paragraph 420

And this passage is sufficiently powerful to demonstrate the deity of our Savior Christ. These points should be compared with what is written on the same topic in chapter four. The twenty-four elders fell down before him who sits on the throne, and worship him who lives forever and ever; and now it is said that the very same elders have fallen down before the Lamb. Consequently, he who lives forever and the Lamb are worshipped with equal glory, cult, and honor, and that the Son is coequal with the Father, to be worshipped eternally. Thus, the abominable and detestable error of Arius and Servetus is now openly perceived, refuted not only by beasts but also by the entire congregation of saints in heaven. Idle men reason subtly, perverting and twisting God’s word according to their habitually bold inclinations, at their pleasure. We will rather follow the examples of all saints and creatures in the world, and will worship the Lamb with him who sits on the throne, blessed forevermore.

### Paragraph 421

Again, there are objections to us: the Elders lying prostrate on the pavement, holding in their hands harps and vials. [Note 28.6] A harp in the Psalms and holy history is an instrument of music, consecrated to divine praises. Concerning the vial, writers of vessels discuss much about its shape and fashion; I understand it to be simply a cup or a bowl, such as we read there were many in the tabernacle and temple, appointed both for drink offerings and also for sweet odors and incense. Nevertheless, these things in the holy heavenly dwellers are not to be taken corporally, but spiritually, after a figuration. For what the Spirit of God understood, the revealer of secrets, Saint John himself adds, which are the prayers of Saints. Therefore, it is signified that Saints offer prayers to God, which are much more acceptable to him [Note 28.7] than the sweet melody of musical instruments is to man, or the pleasant savor of sweet gums or of incense. Arethas, the expositor, says that because they have harps, it shows a concord and agreement in giving God thanks. And from this we learn again what we should do in the contemplation and understanding of the judgments and works of God. The Lord is to be praised and blessed because he is good, and his mercy endures forever. But if thanks must be given to God, if his works and judgments are to be praised, why do certain men expostulate with God, blaming or bringing into suspicion his judgments? Let us learn moreover that organs and those corporal incenses do no longer become the church of God.

### Paragraph 422

Saint Irenaeus, in book 4 against heresies, chapters 33 and 34, demonstrates that the prayers and thanksgiving of the saints are the same offering that Malachi prophesied would be offered throughout the world. Shortly afterward, Tertullian followed the same interpretation against the Jews and in book 4 against Marcion; other doctors of the Church have followed them. But those pleasant, sophistical triflers, I mean the papist divines, as if triumph in these things, yet they lead a shadowing and a most vain triumph. For they apply these things to their sacrifice, wherein they pretend under the guise of bread and wine to offer up to God the Father the body and blood of Christ, a propitiatory sacrifice for the living and the dead. But Irenaeus and Tertullian speak not of such a sacrifice, but of the oblation of prayers, which the mass-priest offers not up alone, but the whole congregation of Christ sanctified in his blood, giving thanks in the Lord’s Supper to God the Father for their free redemption. These holy fathers never knew the sale Masses of these Canaanites.

### Paragraph 423

[note 28.9] Regarding this same passage from the writings of John, the Roman Catholics attempt to use it to prove and establish the practice of praying to saints in heaven. “Behold,” they say, “the saints are said to pray openly in heaven.” But they argue that they do not pray for themselves, and therefore, as intercessors and patrons, they pray for their clients and worshippers on earth. I reply that the saints indeed pray in heaven, but you, by adding your own explanations of the kind and manner of that prayer, distort it, fabricate, and maliciously and falsely misrepresent it. John explains himself here, and does not require your additions. For he adds, “and they sing a new song.” Indeed, he recites the entire form of this song, lest anyone corrupt what he had said about prayers. And that form contains praise and blessing or thanksgiving, and not intercession or invocation. It is certain, even according to the Apostle’s doctrine, in 1 Timothy 2 and Philippians 4, that there are two chief kinds of prayer: invocation and praise or thanksgiving. But the passage itself plainly shows that John speaks here of the latter, not the former. [note 28.10] And just as this passage explains certain types, shadows, or mysteries of God’s law, we can more effectively refute the intercession of saints in heaven for their worshippers. For in the law, only one golden altar of incense was permitted. And that altar represented a figure of Christ. For one Christ is the mediator and intercessor between God and man. It was not lawful for the people of God to burn incense, but only upon this altar. It was not lawful for any man to prepare or make for himself an odor of those kinds of gums from which the divine incense consisted, and to offer it to God, as appears in Exodus 30. Why, then, do these people not understand that prayers belong to God alone, and that the saints in heaven would not offer such incense? David, in Psalm 141, says, “Let my prayer be directed as incense in your sight, the lifting up of my hands an evening sacrifice.” The Devil desires to have such incense made to him, as appears in Matthew 4 and in Augustine’s *City of God*. But our heavenly saints are not devils. Why do they not understand that this altar of incense now stands in heaven on the right hand of the Father, and there makes intercession for us? And that for his sake the Father is reconciled to us, and we are accepted of God, and that through him alone we must offer up our prayers to God, which are otherwise abominable? Why do they not see the heavenly saints at this present attributing all things to the only Lamb alone, and claiming nothing for themselves? Finally, that they make no mention of their worshippers, but plainly testify that the only Lamb was and is worthy, who should take the book, etc.?

### Paragraph 424

And the praise or thanksgiving of the heavenly saints he has called a new song,[note 28.11] which in the Scriptures is no new thing. For the saints say that they will sing on Earth to God a new song, Psalm 33, 96, 98, 149, Isaiah 42. And new songs are called these new ballads or verses in meter, which are made of some new benefit or noble act done. And because the mind of man is greatly delighted with new benefits, they sing a new song, which with a joyful mind praises God and gives him thanks with their inward affections. Finally, they sing a new song, which with purified minds and renewed by the Spirit does laud God; which thing was chiefly given to those heavenly saints. From this we learn again how our minds ought to be affected and furnished in the prayers and praises of God. This same, says Arethas, I call a new song, by whose benefit we, who being enlightened in all parts of the earth, departing from the antiquity of the law written, and walking in the newness of life, are taught by the Holy Spirit to sing a giving of thanks.

### Paragraph 425

To these things now is added the Hymn of Saints, so that we might also have a form whereby to praise God. And in the Hymn they sing that all things are subject to Christ, and all things ordered by his rule, that he humbled himself to death, and was therefore exalted above all things. Now also the virtues or effects and wonderful benefits of his death are commended to us, that, esteeming the governor of the benefits done to us, we may believe also that his rule will be wholesome for us, and therefore may submit ourselves to him willingly in faith and patience. Which indeed is the chiefest end of those things which here are treated with such diligence.

### Paragraph 426

First, they commend the majesty and dignity of Christ, stating that he alone is found throughout the entire world, possessing dominion over all, the sole savior of the world, the revealer of divine mysteries, and the governor of all. For this is to take and to open the book, which we have often repeated. Secondly, they add the reason why this glory should belong only to the Lamb or Son of God: because, they say, you were slain. And they understand by the lesser the greater, namely, his entire incarnation and the entire mystery of our redemption, death, resurrection, and ascension into heaven, and the remainder. Therefore, he is the true and only mediator of God and men; he is the only savior, as he alone was incarnate and crucified for us; he is the only governor, who, through his humility, deserved to be exalted, Philippians 2. And he is a most suitable governor of all things, as from whom all men may, as their most faithful savior and even their brother, hope well, whatever events befall them through his governance, etc.

### Paragraph 427

Meanwhile, they commend highly the power or effect of Christ’s death. For if this is rightly understood, we are more ready to submit ourselves to that governor whom we know to be our savior, who loves us dearly and desires all to be saved. And the chiefest effect of Christ’s death is redemption; redemption includes deliverance from captivity. We were prisoners and servants of sin, of death, and utter slaves of the devil and hell. And the Son of God came and took flesh, and shed his blood (for so also is the manner of redeeming us expressed by the elders), and he has washed us from our sins, and being purified, he has ransomed us from the power of death, hell, sin, and Satan, so that now we belong to God. Therefore, they say expressly, “You have redeemed us to God.” We therefore belong to God; the devil has no more power over us; we are the freemen of Christ, delivered through his blood. 1 Peter, 1; Hebrews, 9. And since we now belong to God, namely, justified freely by his grace, through the blood of Christ, as the apostle also says in the 3rd chapter to the Romans, we ought to serve God truly in the newness of spirit, not the flesh and the devil, in the oldness of the letter and of our flesh. Which the same apostle discusses more at length in the 6th chapter to the Romans.

### Paragraph 428

Those whom he has redeemed declare it openly, men truly of every tribe, and so forth. In this recounting, he imitates Daniel in chapter 7, signifying universality. For the Lord died for all; but that all are not made partakers of this redemption is through their own fault. For the Lord excludes no one, but only him who, through his own unbelief, excludes himself.

### Paragraph 429

Following redemption is another effect of Christ's death,[note 28.17] for it makes men justified to God as kings and priests. For those who are justified, work righteousness. I have explained this passage concerning the priesthood and kingdom of Christians in the first chapter, where you may find it.

### Paragraph 430

The saints also affirm that they will reign on earth, namely through the power of Christ; not physically, as the Millenarians imagine, and the Turks following the same, imagining physical things in this world and joys in a terrestrial paradise. For the entire Scripture promises better things. Nor must the godly be so devoted to physical things that they should hope for nothing beyond physical matters. The saints speak here of the final judgment, in which it will be shown to the entire world and to all who dwell on earth that the saints, who at one time seemed to the world to have been wicked, ungodly, peacebreakers, heretics, and parricides, and for that reason were slain, are righteous, holy kings and priests of God. Thus I say they will reign on earth. This matter is explained more fully in chapters 3 and 5 of the Book of Wisdom.

### Paragraph 431

Let the faithful consider these things, I say, when they are oppressed by the wicked for the truth and righteousness, through the permission of Christ, the governor of all, in this world; nevertheless, let them glorify the Lord God and praise him without ceasing. To him be glory forever.

## Sermon 29: ¶ Here is described the commendation and Hymn said unto Christ of the Angels and all creatures. &c.

### Paragraph 432

.

### Paragraph 433

And I looked, and I heard the voice of countless angels around the throne, around the beast and the elders, and I heard thousands of thousands saying with a mighty voice: “Worthy is the Lamb who was slain to receive power, and riches, and wisdom, and strength, and honor, and glory, and blessing.” And every creature that is in heaven, and on the earth, and under the earth, and in the sea, and all that are in them, heard. Saying: “Blessing, honor, glory, and power be to him who sits on the throne, and to the Lamb forever and ever.” And the four living creatures said, “Amen.” And the twenty-four elders fell on their faces and worshiped him who lives forever and ever.

### Paragraph 434

Fourthly, the angels of God now come to the elders and the beasts [note 29.1], that is, to God’s most excellent creatures, and together with them offer praise to God and the Lamb. This is surely an example for us, as I often say and repeat, so that we may understand what is fitting for us as well.

### Paragraph 435

Of Angels David in Psalm 4, speaking among other things: which says he makes his Angels spirits, and his ministers a flame of fire. He testifies therefore that the Angels were made or created by God. By their substance he calls them spirits, and by a parable a flame of fire, which is pure, bright, most swift, piercing, and burning. Therefore, in their sort and manner, the Angelical spirits are altogether such: whom by their office he calls ministers, to wit, of God and man. For Paul also to the Hebrews, bringing this same place of David, asks, “Are they not all ministering spirits, sent forth to serve those who will inherit salvation?” truly understanding men. These things teach us to judge rightly of Angels, and that no man should worship ministers, or any, however excellent creatures they may be, for their godly gifts. Neither, in deed, can the Angels or Saints abide being worshipped. Here, without doubt, they attribute all glory to God and to the Lamb, to God three and one, that all we should do the like. Here is also declared the place wherein the Angels were, about the throne, about the beasts, and about the Elders. Therefore, they guarded all these places round about as it were a guard. Daniel in times past saw things not much unlike these in chapter 7. Certainly, they stand like ministers and servants, ready to do service. Angels are said also to be about the godly upon earth, and to attend upon the salvation and ministry of men. In Psalm 34, David sings, that being afflicted he called upon the Lord, and the Lord heard him, and from all his troubles he delivered him. The Angel of the Lord pitches his tents about those who fear him, and he delivers them. And not much unlike things you may read in Psalm 91. And you shall note here, that those who are afflicted do call upon the Lord, and not the Angels: and that the Lord does hear, and deliver, and for the working thereof uses the ministry of Angels, as his ministers. And just as no man of sound mind reveres, calls upon, and worships the son, because God by him gives great benefits to men, so no man honors, calls upon, and worships Angels, because God uses their ministry in delivering men.

### Paragraph 436

Now he also states the number of angels, but a certain number for an uncertain one, thousands of thousands for innumerable. He alludes in the meantime to that saying of Daniel in chapter 7: “Ten thousands ministered to him, and ten thousand times ten assisted him.” We are accustomed to consider the power of kings as a great and innumerable army. What, then, shall we think of the power of our God, who is the God of hosts, and whom not only innumerable legions of angels, but all creatures serve? And what an excellent praise is sung together by so many blessed spirits?

### Paragraph 437

After this, the proper duty of the angels is addressed: they sing praises to God and commend God’s work, and that with a loud voice. It matters little whether you sing praises to the Lord with a low or a high voice; but because those who cry with a loud voice are often deeply moved, whether overwhelmed with great sorrow or rejoicing with great gladness, therefore we shall praise God with a loud voice, if with a steadfast spirit and with the inward affection of the heart we praise God.

### Paragraph 438

[note 29.5]The evangelical hymn is now included, which agrees in all things with the Hymns of the Beasts and Elders. For they celebrate the Lamb, that is to say the Son, which as he is the savior alone, so hath he deserved to receive all power and glory, and to govern all things: as is said before.

### Paragraph 439

And seven things the Angles attribute to the label, that is to Jesus Christ our Lord, on the right hand of the Father. First, power, namely godly, almighty, vivificative, and conservative. Of this I spoke also before. Secondly, riches. For he is rich, as the Apostle says, for all who call upon him. Romans 10. And Primasius says: Christ himself is the treasure of all good things. For *schaddais* is the sufficiency of all goods of the mind and body. And if it be lawful to attribute a profane word unto God, he is very Saturn, fulfilling all creatures. And the Angles commend Christ so, who would think that men should strive so to acquire things for themselves, as though they themselves could fill their own desires? Then they attribute to Christ wisdom, namely godly and great. For the Son is the wisdom of the Father, a topic much treated by Solomon. By this wisdom he can rule all things by the most agreeable and best government. Who shall say? Thus it should have been done. The wisdom of God has most beautifully and well made all things from the beginning, so that our reason can justly blame nothing: what thing shall we blame now in the universal government of Christ? They ascribe unto Christ also strength to execute truly such things as he has most wisely ordained, finally strength to defend his, and to subdue the adversaries. For he is almighty. Such things as follow, honor, glory, and thanksgiving, are declared before, what they are, and of what force; saving that the thing he called first [illegible] he calls now [illegible], blessing, praise, and giving of thanks.

### Paragraph 440

This hymn, sung in praise of Christ, [note 29.6] teaches that Christ is truly God, of the same substance and coequal with the Father, greater than angels, indeed the Lord of angels, whom the angels themselves also worship, as Paul declared in the first letter to the Hebrews. Therefore, those who place angels above Christ are refuted. The heretics called Angelicals—that is, worshippers of angels—are refuted. The angels themselves here rebuke their error: rightly are they considered heretics by Augustine. If riches, glory, and honor are due to Christ alone, and he surpasses them, why are these communicated to creatures? Otherwise, we admonish all faithful people to esteem angels highly, to acknowledge and marvel at the benefits of God working through them, and to love them as brethren, companions, and coheirs of the same salvation; much less to despise or blame them. Of this I will speak more at another time.

### Paragraph 441

Hitherto he has recited the excellent praises, the celebratory verses or hymns, of God’s excellent creatures, especially of the Elders and generally of all, [note 29.7], moreover of Angels also, spoken to Christ our redeemer and prince. And yet not content with these, he adds moreover in the fifth place, the agreement, and praise, and submission of all the creatures in the world, so that if we are not moved by the excellent example of the excellent creatures, the Elders and Angels, at the last we might be ashamed, seeing all creatures of their own accord do their duty. For man sins; he is Lord of all, and all things were created for him. How, I pray you, shall he offend more heinously against God, who has made him Lord over all, than if by his hardness, ingratitude, and malice he not only does not do his own duty, but is rather inferior to all creatures; as he who alone contends with God, and does not attribute to him due praise? Therefore this example greatly excites man that he should submit himself unto God, and give God the whole glory; and in no wise strive with God, nor complain of anything. But mark, I pray you, with how diligent a division of things he encompasses all creatures, excluding none, the Devil only excepted, when he recounts the creatures that are in Heaven, that are in Earth, that are under the Earth, and in the Sea; finally, he adds, and all that are therein. Therefore, if all created things do celebrate and worship him that sits in the Throne and the Lamb, and submit themselves unto him, is it not a shame, yea and a soul shame, that man alone, Lord of all, should revolt to the sworn enemy of God, the Devil, and with him to expostulate with God, to teach and blame and find fault with his judgments and governments, to complain of his works and will?

### Paragraph 442

You marvel, I well know, how all creatures, even those without reason and feeling, can praise God. However, this figure of personification—the attributing of human qualities to non-human things—is very common among the prophets, and especially with David, who says, “Praise him, sun and moon; praise him, bright stars.” “Praise the Lord from the earth, sea monsters and all deep places; fire and hail, snow and ice.” And by such ways of speaking, the prophets would encourage and stir up people to praise God, saying that creatures which have no life, do after their manner praise God; see that you after your manner do praise God in hymns and spiritual psalms. Indeed, David shows a plain reason why he commands bodies that have no life to praise God: “Let them praise the name of the Lord,” he says, “because he commanded, and they were created.” As though he should say: they are his creatures, and in that they remain, they have it of him; therefore let them make the name of God glorious, as of their maker and preserver. And he signified also the manner of praising, where he adds, “He hath ordained them, that they should endure forever; he hath appointed a decree, and it has not passed away.” As if he should have said: Where they neglect no part of those things for which they were made, but are ready in their place, order, and time, and do their duty exceedingly well, do they not preach unto men the wonderful wisdom and power of God? For in another psalm also, David says, “The heavens show forth the glory of God, and the firmament declares the works of his hands.” Thus I say, the creatures without life do praise and commend the name of God to men, when they are moved, work wonderfully, and obediently do the thing to which they are appointed.

### Paragraph 443

The Hymn of all creatures, similar to that which was of the beasts, Elders, and Angels, is here also suitably described, though briefly. But where it contains nothing that has not already been declared, I will not weary the gentle listeners by repeatedly speaking the same things. However, one thing seems chiefly to be observed: that they join him who sits on the throne and the Lamb together, thus acknowledging the Son to be coequal with the Father, and that both are to be worshipped with like honor, and with like praises celebrated and commended. They attribute uniquely to the Lamb dominion or kingdom, for he received the book of the Father, as has been declared before: namely, all power and authority to govern all things.

### Paragraph 444

The four beasts sing to it Amen,[note 29.10], either confirming the hymn of the creatures, or declaring their consent with them. This is so that we may, with one mind, pray together and praise God, blessed forevermore. With these are moreover refuted the dispensations of men. The Lord allows the concord and agreement of men, and requires it utterly, especially in prayers and godly praises. For he commands in the Gospel to lay down your offering, which you would offer, if you remember any discord between you and your brother, to go to him, and to renew friendship, and then to return to your offering: which in the prophets is called an abomination, if it be offered by minds possessed with rancor and malice. &c.

### Paragraph 445

Finally, the elders fall down again and worship him who lives forever, no doubt so that all of us on earth might be moved to obedience. For if these things are done in heaven by the blessed spirits, what, I pray you, is fitting for us to do here on earth? And note that they are said to worship him who lives forever, who nevertheless also fell down before the Lamb and before the Throne, from which the Spirit proceeded, and where he who sits is seated; from which we gather that the Father, the Son, and the Holy Spirit are indeed distinct in persons, yet these three are not three Gods, but one God living forever.

And truly, this notable vision and treatise may serve as an effective remedy against various poisons of heresies, especially those of the Arians and Servetans, or rather, [?perdetans], moreover against diverse and curious disputations and temptations concerning the works, judgments, and providence of God. If we are wise, we will obediently submit ourselves to the living God with all the creatures and saints of God, worshipping him, and with the prophet crying: “You are a just Lord in all your ways, and holy in all your works. You have created us; all things are yours. You govern all things in the best order. You love humankind. You have given us your Son. You, by your Son our redeemer, govern all things uprightly. We worship you, the Father, the Son, and the Holy Spirit, one very God. To you is due the kingdom, honor, and glory forever and ever. Amen.”

## Sermon 30: ¶ Two seals are opened, and the direct course of Gods word is, and a cruel course of wars against the disobedient.

### Paragraph 446

.

### Paragraph 447

And I saw when the Lamb opened one of the seals, and I heard one of the four beasts say, as it were the voice of thunder: Come and see. And I saw, and behold a white horse: And he that sat on him had a bow, and a crown was given unto him: And he went forth conquering, and to overcome. And when he opened the second seal, I heard the second beast say: Come and see. And there went out another horse that was red, and power was given to him that sat there on, to take peace from the earth, and that they should kill one another. And there was given unto him a great sword.

### Paragraph 448

Up to this point, the Apostle has prepared the audience to hear with a calm mind about God’s judgments and the church’s destined fate, and to patiently endure all adversity, and to worship him in all things, and to give glory to his name. Consequently, he explains in a very orderly fashion God’s judgments and the church’s destiny, showing how the Son of God governs God’s ordinances and his eternal providence. And this is, as it were, a prophecy for all times and ages until the end of the world. [note 30.2] For we should not think that only the events of one age or two are recounted here, but of all. First, all things are generally described by parts; afterward, particularly, when we come to the opening of the seventh seal. The sum is, the Lord sends forth the preaching of the truth into the world, which when people refuse and despise, they are destroyed with wars and other innumerable calamities.

### Paragraph 449

[note 30.3] But above all things, Saint John is stirred and eager (and in him, all of us) to be attentive. And one, that is to say, the first of the beasts, excites him. One of the Sabbath is set for the first day of the week, which is truly Sunday. The voice of the beast is like unto thunder, signifying that great and weighty matters are being treated. For most great and terrible things follow, which shake the whole world. Therefore let us not play the sleepy sluggards, let us not be blind and deaf. Doubtless the slothfulness of our time is such that we little consider the works of God and what is done in our time. The storks, swallows, and doves, and the rest of living things pass us by, which carefully observe their time. Therefore, let us be stirred up here, that we should not be slothful, but should mark what things are declared and shown to us by the Lord.

### Paragraph 450

And when John had carefully recorded what was done, he sees the Lamb, Christ, our Redeemer, open one seal, that is to say, the first. Immediately a white horse came forth, on whom he who sat had a bow bent, and an arrow in it. To him was given a crown, and he went forth conquering, that he might overcome. This is the vision: the explanation of which is easy. For the Lord says that he will declare the destinies of the church.

### Paragraph 451

Horses of sundry colors are also brought forth by Zechariah in the first chapter. They signify the variable course and state of the people of Israel. The white color is consecrated to innocence, purity, victory, and felicity. Therefore, the white horse signifies the fortunate proclamation of God’s word, or the prosperous preaching of the Gospel. For upon the horse sits a horseman, who guides the horse and has a bow. Certainly, Christ promotes the course of preaching the Gospel. And Psalm 45 attributes to the same shafts or arrows, for he strikes his enemies far off and brings them into his subjection. Briefly, with the word of his mouth, he subdued to himself people and nations. Isaiah, in chapter 49, bringing in Christ speaking, says: “And he put my mouth as a sharp sword, the shadow of his hand covered me, and he put me as a sharpened arrow, he hid me in his quiver.” Through Christ, therefore, proceeds the preaching of the word; he gives strength to the preaching; he draws back his bent bow. Whatever force the word has, that is wholly due to the Horseman.

### Paragraph 452

A crown is given to him [note 30.8], namely, a kingdom and all power of ruling. For David prophesying beforehand said, “The Lord will send forth the rod of his power out of Zion, to rule among your enemies.” Moreover, a crown is given to him, that he may crown those who serve him faithfully. And it is a figure of speech, and he went forth conquering, that he might overcome: for he who went forth is a conqueror, and for this purpose went forth, that he might overcome. For it signifies that Christ will advance the preaching of his word throughout the world, no one being able to resist, and even in defiance of the gates of Hell. For the word of the Lord endures forever.

### Paragraph 453

And this place teaches that the Church shall always be in the world, and likewise always the truth preached, even though the enemies’ hearts burst. But if we read over the history of the Church, we shall better understand all things, and shall perceive that this prophecy has always been most certain. Christ was once through the ministry of the word showed to the world by the Apostles, and the matter proceeded most favorably, however much the most powerful people of this world resisted it. The thing is wonderful, if one considers those five hundred years which immediately after the incarnation of our Lord are accounted. In them went forth the conqueror that he might overcome: and overcame in deed, the whole world receiving Christ and worshipping him. Since those years, as before also, certain seeds of errors began to be sown abroad. The bishops began to contend for supremacy, and who should be the universal head of the Church on earth: they began to reason about the use of images in the Church, and brought them into churches in deed, as also they called the bishop of Rome the supreme and general head of the Church on earth. And mighty princes, and in a manner the whole state of learned men conspired in these opinions: but he has vanquished, which went forth that he might vanquish. He had in his church innumerable, who bowed not their knees before this Baal. A thousand years after the incarnation of Christ, the bishops began to profanely pollute the Lord’s Supper and other undefiled doctrines of faith: but what, I pray you, did they prevail by so many councils, determinations, and earnest endeavors? He that went forth to overcome, has overcome. That white Horse has stoutly advanced to the salvation of many. For how great battles the godly and learned have sustained against the Popes and bishops in these last five hundred years, history bears witness. At this day also appears throughout the whole world, how favorably yet that white horse goes forward, which has pierced even until our time. The Gospel is believed, and that faith cannot be extinguished with any waters or fires.

### Paragraph 454

You make exception that they were heretics who resisted the bishop and See of Rome in these last 500 years, such as Bertram, John Scot, nicknamed Duns, Berengarius, Arnoldus Brixianus, Waldo, Wycliffe, and Hus, Luther, and Zwinglius, and others of the same sort; moreover, some of these were also overcome and put to death by the Pope. I answer that, as men, they might err in many things; but in those things wherein with the Scripture they consent against the See of Rome, I affirm that they erred not, but spoke the truth. Whereupon it is certain that Christ overcame by them. When Micah, Elijah, Zechariah, Amos, Jeremiah, and others preached by the word of God against idols and worshippers of idols, they also were condemned as seditious and heretics; indeed, some of them were removed from the way. But was the truth vanquished? Antichrist is said to have good fortune, and to punish and afflict the strong and the people of God; but men being ministers may be oppressed, the ministry never decays. Saint Paul says that he is bound for the Gospel’s sake, but the Gospel is not to be bound. Therefore, he has overcome hitherto, and shall overcome still, which went forth that he might conquer. They stumble upon this conqueror as at the stone of offense, whoever and whatever they may be, who seek to interrupt the plain course of the Gospel.

### Paragraph 455

Moreover, concerning the time when the second seal should be opened, the second beast [note 30.10] exhorts John again to attentiveness, and that we should consider what is presented to us. And now comes forth the red horse, whose color is somewhat like fire; there sits also on him a rider, to whom power is given to disturb peace on earth, and that men should kill one another. For there is given him a great sword. The red horse signifies the state of wars, full of fire and blood. He who sits on this horse is Mars, or rather the father of lies, that is, the Devil, who was a murderer from the beginning. He gathers to him the dregs of men to make civil commotions, for the wars, destruction, burning, slaughter, and desolation. You see from whence the breaking of peace is, which God hates. And we hear how it is given him: Mark given, by the just judgment of God to be permitted, that troubling all peace, he should take it away, and set men together by the ears, that one may wound and kill another. For so we read in the first book of Job, how Satan had power given him by God against Job. To bloody soldiers is given a great sword, great power to hurt, a wonderful force of fighting, as also Nahum explains it. Neither is it a rare thing in the scriptures for monarchs, tyrants, and mighty men of war to be called a sword. For so Ezekiel called Nebuchadnezzar; and Isaiah called Sennacherib, king of Assyria, a whetstone.

### Paragraph 456

And the chiefest righteousness is to give every man his own. Therefore, this passage justly ascribes that which is good to God and that which is evil to the Devil. But you say, if God permits, the same that he does not prohibit he does. He does not prohibit war, for because justice will not suffer him so to do; but he commands him by war to punish the wicked and to try the good. But in permitting wars, God offends nothing, saying that for most just causes he permits the same. For they would not embrace peace offered them by the preachers of the Gospel, therefore were they worthy to be entangled with wars. The Jews did not know the day of Christ’s visitation, therefore they were justly visited by the Romans and destroyed. And this thing is in the world perpetually, that they who will not obey the Gospel must obey the Captain of the wars; they who will not hear Christ must hear Antichrist. You may not contend with God, which he does this and permits that. Worship God rather, as you have been taught in the fourth and fifth chapters.

### Paragraph 457

Let us examine stories and see if such wars are found, wherein men have slain themselves with mutual wounds, and have killed one another like beasts. If you will read Herodotus, Orosius, and other good historians, you may find that the Roman Emperors were troubled with most grievous wars, for no other cause than that they refused peace offered to them by the gospel. For no other cause was Rome itself at the last taken by the West Goths, and burnt and destroyed by the East Goths. The Lord had given them Christian princes, but they loved idols more. For Simmachus, governor of the city, was so bold as to demand the restoration of idolatry. I speak nothing now of Attila, nothing of the Persian and African wars. And when there was a great strife among the bishops about supremacy, the Saracens arose and became mighty. After the thousand years began the holy war, which was as bloody as it was long-lasting. Never was such a war made in all the world. Boniface VIII instituted first the year of Jubilee, a most wicked man, who also exhibited himself to be seen by the people, both as Pope and Emperor. But in the same year of one thousand three hundred, Ottoman arose in Asia, the scourge of God, the founder of the emperors of the Turks who reign at this day. For so when Solomon built places of idolatry, his enemies arose, who wonderfully vexed and afflicted the kingdom of Solomon. What wars are made nowadays, and what are the causes of wars, all wise men see. We will not receive the peaceable gospel; it is reasonable therefore that Turkish armies should invade us, that we may both feel Antichrist to be a stout warrior, and may all abhor and detest him.

### Paragraph 458

But what else remains here, than that, being converted to God through Christ, we may serve the Lord in sincere faith and holy purity? For unless we repent, the groundwork is laid at the tree’s root, &c.

## Sermon 31: ¶ Here is opened the iii. and iiii. Seal, and is declared what the world shall suffer of honger and Pestilence.

### Paragraph 459

.

### Paragraph 460

And when he opened the third seal, I heard the third beast say, “Come and see.” And I looked, and behold, a black horse; and he who sat on him had a pair of balances in his hand. And I heard a voice in the midst of the four beasts saying, “A measure of wheat for a penny, and three measures of barley for a penny; and injure not the oil and wine.”

### Paragraph 461

And when he opened the fourth Seal, I heard the voice of the fourth living creature say, “Come and see.” And I looked, and behold a pale horse, and the name of the one who sat on him was Death, and Hell followed after him, and power was given to them over the fourth part of the earth, to kill with sword, and with famine, and with death of the beasts of the earth.

### Paragraph 462

Christ, exalted above all things, and Lord of all in heaven and on earth, opens the seals of the divine book, that is to say, disposes and governs with great righteousness the ordinances and judgments of God. First, indeed, he gives a prosperous course to the preaching of the gospel, sending always faithful ministers preaching the Gospel of the kingdom of God, peace, and concord. But because evil men despise the Evangelical peace, they are certainly worthy to be afflicted with cruel wars. Therefore, the Lamb opens the second seal, and there rush out cruel wars, slaughters, seditions, and robberies.

### Paragraph 463

But before the third Seal is opened, the third beast, resembling the face of a man, exhorts us to take most diligent heed: that when we see these things come to pass which are spoken of here, we should consider from whence they come, and for what causes they are sent, and that they may be turned away by true repentance. Some refer these things absolutely to chance and fortune, others to natural causes, with no respect had at all for God and the divine operation: whereas we know that God uses natural causes according to his good will and pleasure. Let us watch therefore, look and consider, and know that the righteous God works all things for the salvation of the elect, and the overthrow of his enemies. That black horse with its rider, showing a balance in his hand, signifies the unfortunate or sorrowful time of scarcity, famine, and penury of all things.

For it is a worthy and fitting punishment that those who do not esteem the bread of life, nor have any consideration of the food of souls, but both reject it themselves, and by their tyrannical proclamations bring it to pass that it is not received by others, and finally who, for the bread of life, despoil the godly of their goods, and most wickedly waste them in all kinds of riot, should be driven to necessities at excessive prices: yea, and cannot find necessities, but should pine for hunger. We know that the color black is used in mourning and grief; and that when flesh and blood are consumed for want of food, the skin grows black and ill-favored: and therefore this horse is black.

### Paragraph 464

The rider of this horse holds in his hand a balance, [?a token] with two scales hanging at either end of the beam, which we call a pair of weights. Arethas says that a balance is a token of right and equity. For you have sat, says David, upon your throne which judges righteousness; therefore a balance is the judgment of the just judgment of God. Arethas has not alleged these things amiss, but we ought rather to prefer the exposition of Saint John himself. For a voice is heard from the midst of the beasts, which expounds to us the balance. For it sounds, a measure of wheat for a penny, and three measures of barley for a penny. And this measure, called Choinix, signifies a diet or daily meat, as Erasmus has in his proverb, “Sit not upon your measure.” The same in his annotations upon this place: Choinix, says he, is a measure of wheat, or other breadcorn, which is sufficient for one day’s meat. Budaeus thinks that it weighs four pounds, Pollux three. The word therefore signifies that a very little meat shall cost a great price, and yet not be gotten for money, which happens in the time of famine. What the Roman penny is worth Budaeus shows; we understand by it plainly a great price. Therefore two things are signified: scarcity or dearth of corn, and famine. Scarcity raises the price beyond reason. Famine has nothing to buy, though he has never so much money lying by him, but hungers, wants, pines, and at the last miserably consumes to naught, wherein verily dearth and famine do differ. The Germans discern them by several words, calling dearth scarcity, and famine hunger. Yet they are for the most part indivisible.

### Paragraph 465

And we read in the old story of the Bible,[note 31.7] that for the contempt of the preaching of God’s law, and the bringing in of a strange kind of worshiping God, the Israelites in the times of Elijah and Elisha were most grievously punished with hunger and poverty. These things are plentifully declared in the third book of Kings, chapters 17 and 18. Also in the fourth of Kings, chapters 6 and 7. Moreover, in the time of Emperor Claudius, while the Apostles preached the Gospel faithfully, and the Jews and Gentiles stoutly repulsed it, famine most grievously afflicted the Roman Empire: which thing Luke rehearses in the Acts of the Apostles, chapter 11. These things were indeed done before this revelation was exhibited to John. Since that time, historians recite sundry and innumerable famines, dearths, and poverties in diverse countries, sent of God for contempt of the truth. Nauclerus mentions a famine in the year of our Lord 1339, wherein mothers also devoured their own children. What has chanced in our memory in those wars of Milan and elsewhere, it is no need to rehearse. They are yet fresh in memory, and written in the histories of Galeazzo Capella. We felt some part hereof also in the year of our Lord 1529 and the years following. The just Lord punishes, and more will punish the great ingratitude and contempt of his godly word, as he did in the destruction of Jerusalem. Would God it would please the world, so blind, through repentance to convert unto God when he punishes, and with free and willing minds embrace the word of truth: for so should there be more felicity and less misery.

### Paragraph 466

However, for comfort at the end of this seal, it is added, “See that you hurt not the oil and the wine.” He names the kinds most necessary for human use, and means that God mercifully reserves some things that are chiefly necessary for humankind, especially for the sake of the elect, so that all should not perish and suffer generally. By this, we understand that the Lord does not forget his mercy, even in the midst of affliction and plagues that he sends. Thus, in times past, intending to punish Egypt and other nations with famine, he sent Joseph beforehand, by whom he might preserve the house of Jacob and other people innumerable. You see herein most clearly that it is from God that sometimes the grain is blasted, and the vines and olives perish; and that it is from him that the grain increases, and wine also. So he has also previously protested in the law, Leviticus 26, and Deuteronomy 28.

### Paragraph 467

We have now come to the fourth Seal, at the opening of which, and to observe its effect, we are stirred up by the Eagle, the fourth beast. Of whom we have spoken before on one or two occasions. And the pale Horse comes forth, in Greek *chloros*, a color resembling withered grass and herbs. Solomon in the twelfth chapter of Ecclesiastes calls the color appearing in dead bodies and their faces a golden liquid. All poets call death pale. And the rider is indeed expressly called Death. We understand the course of plague and all diseases, and even death itself; and Hell follows, that is to say, a pit or a grave. For *Sheol* in Hebrew signifies a pit or a grave. But if you will needs understand it of the place of the damned, doubtless they are carried headlong into Hell, so many as here, consumed by sickness, die without faith and repentance. Therefore Hell follows death rightly. But if you would rather understand Hell as a grave, it signifies that it will be full of corpses and sepulchers.

### Paragraph 468

And indeed plagues and pestilences most mortal have severely afflicted the Roman Empire [note 31.11]. Orosius testifies in his seventh book in the Acts of L. Aurelius Verus and Decius, emperors, the most cruel persecutors of our faith. Evagrius in chapter 29 of book 4 of Ecclesiastical History tells of a marvelous plague that lasted about 50 years. And all men know with what a pestilence and sudden death Italy was wasted in the time of Emperor Maurice and Bishop Gregory of Rome. Time would fail me if I were to recount from histories all the plagues and calamities of all times. What is done at this day, and has been done in our memory, you yourselves know best. New diseases have arisen, whose names were never known to our elders. With these evils and calamities God wastes the world, and has always done so, intending that by plagues he might call us again to repentance. Thus we must always judge of calamities. If anyone judges otherwise, they are not reformed, therefore they are punished here and after this will burn in perpetual torments.

### Paragraph 469

To these moreover is added another thing also, and power was given them, and so forth. For when people will not reform with simple calamities, the evils or plagues of God are doubled. They are listed in the same order and number as the prophets, Jeremiah in chapter 15 and Ezekiel in chapter 14. For they are these: sword, famine, death or pestilence, and wild beasts; so they are recited in the Law also. With these, as if sent from the four parts of the world, God, most righteous, executes his judgments.

### Paragraph 470

Let us observe this chiefly: that power is given by God to kill, and that over a fourth part of the earth. For we learn that God alone quickens and slays, and that he works this justly by his instruments, finally that all his things are numbered and done in order. Whereupon he pours out his fury upon a third part of the world. For he knows whom he should punish and whom he should nurture tenderly.

### Paragraph 471

Certainly, accounts attest how in dire circumstances, when all things are brought to an extremity of misfortune, God has unleashed sword, pestilence, famine, and wild beasts to afflict people. And Arethas appropriately cites the words of his predecessor, Saint Andrew, Bishop of Caesarea, from the ecclesiastical history of Eusebius, book 9, chapter 8. Indeed, within the last fifty years, historians have reported many such occurrences, and we have witnessed some ourselves.

### Paragraph 472

Therefore, if we desire to be free from such great evils, let us serve God in truth and value his word, which he has sent to heal us. It is reasonable that those who reject sound doctrine should be afflicted with various diseases of soul and body, and so forth.

### Paragraph 473

You will say, but these evils invade even the best of things. And indeed they do. Why God permits this, Augustine explains at length in the first book of *The City of God*. Certainly, for the godly, all things work for the best. The thieves suffered the same death on the cross that Christ did, and he as they: but the consideration of them is very different. The apostles and countless martyrs die by the sword, just as soldiers do in war, but with unlike circumstances. The godly become partakers of the suffering of the Son of God. The ungodly are punished for their wickedness, and their suffering is without glory; rather, this is the beginning, unless they acknowledge him who strikes them, of everlasting torments. The Lord preserve us from evil.

## Sermon 32: ¶ The fifth Seal is opened, and the persecution of the faithful set before our eyes, and also the state of Martyrs in an other world.

### Paragraph 474

.

### Paragraph 475

And when he had opened the fifth seal, I saw under the altar the souls of those who had been killed for the word of God, and for the testimony they had borne. They cried out with a loud voice, saying, “How long, O Lord, holy and true, will you delay to judge and to avenge our blood on those who dwell on the earth?” And white robes were given to each of them, and it was said to them that they should rest for a little season, until the number of their fellow servants and brethren, who were to be killed as they had been, should be fulfilled.

### Paragraph 476

The fifth seal being opened by the Lamb, he exhibits to our eyes, or rather compels us to see, the continual persecutions of the church; and he shows us diligently what is the state of those who die in persecutions. Truly, the Lord Christ sends forth ministers and preachers for the salvation of men. And they, unthankful, overwhelm the faithful messengers of God with all kinds of injuries, and at length most cruelly slay them. Concerning this matter, the talk of men among themselves is diverse; but the very Son of God at this present boldly instructs his church, declaring what the godly shall suffer.

### Paragraph 477

And first in expounding the same, we shall speak generally of the persecutions, by which both the ministers and the entire faithful church are diversely afflicted. The Lord Christ has shown us before in the Gospel many things concerning the persecutions to come, truly that he might prepare the minds of all the faithful for battle and patience. The passages are in Matthew 10:1-24, Luke 12:1-21, and John 14:15-16. And also the Acts of the Apostles tell of many things which the godly suffered in that most holy primitive Church. Would anyone have been considered sane if, at that time, they had said: It appears that the apostolic church is not the church, for that it is subject to all the mockeries, injuries, and slaughters of all men? Why, then, do we not acknowledge at this day that those are greatly deceived who measure the church by the outward peace and tranquility of things? Paulus Orosius in the seventh book of histories recounts ten grievous persecutions raised against the church from the time of the Apostles until the emperor Constantine, a time that did not fully accomplish the space of three hundred years. The first was stirred up by Nero, a monstrous man, of whom also Tacitus mentions in his Chronicles. This man put out of the way Peter and Paul, the most holy Apostles of Christ. The second destruction of the church came with Domitian, who in that persecution most grievously afflicted both our John and the whole church also, and when he was brought to Rome, banished him to the Isle of Patmos. The third was raised by Trajan, of whom Pliny, governor of Asia, makes mention in the tenth book of Epistles. In this persecution Ignatius, a holy bishop, was cast to the beasts and devoured. And Marcus Aurelius molested the church with the fourth persecution, and consumed with fire Polycarp, a bishop most worthy. Septimus Severus moved the fifth persecution, which Eusebius pursues in the sixth book of the Ecclesiastical History. Julius Maximinus killed Pamphilus, a martyr, and Sextus raged cruelly against the church. And Decius Trajan began the seventh persecution, and executed very many who professed Christ. And Licinius Valerian, emperor, beheaded Cyprian, the good bishop of Carthage, and was the eighth persecutor of the church. Aurelian began the ninth persecution, which he but little advanced, for God, most just, took him away immediately. But Diocletian and Maximian shed more Christian blood than any other of the Roman emperors. Read, I pray you, the beginning of the eighth book of the Ecclesiastical History of Eusebius. Compare those things with our time, and judge and conjecture what will shortly come to pass, and what our state will be. Persecutions are again renewed after Constantine, under Constantius and Julian. But the most terrible and grievous of all have boiled up under Antichrist, and have endured now for the space of five hundred years and more. What is done at this day, all the world sees. The ground is wet with the blood of martyrs, which things John foresaw.

### Paragraph 478

And the causes of persecution arise partly from the rule of Christ, which opens here the fifteenth Seal, and partly from men. The Lord sends to his people the cross and suffering to quicken those who are slow, to cleanse those covered with rust, and to refine corrupt gold. For so the Scripture declares in the eleventh chapter of Daniel and the Apostle Peter in the fourth chapter. Christ therefore does not destroy, but tests, and permits many things to tyrants against the Church. Godly men also bring upon themselves the Lord’s heavy hand, while indeed they believe rightly in the Son of God and depend only on him, but nevertheless are entangled with various and evil affections and commit acts that do not befit them. This may be seen declared at length in the beginning of the eighth book of the ecclesiastical history of Eusebius, which I recently cited. And the tyrants who persecuted had a different motivation: as Sennacherib and Antiochus did, than our bishops and princes do today. For these now are moved by hatred of religion and are spurred forward by Satan. They will by all means have their idolatrous religion maintained and the religion of the Gospel utterly destroyed. They cannot tolerate having their idols or other sins reproved; for this cause they are enraged at the faithful and those who frankly speak against and blame their idols and wickedness. And thus does the persecution arise, boil up, and proceed.

### Paragraph 479

When the faithful see such increase [?] and feel deeply oppressed, they wonder how long the Lord will overlook this. Many cry out that the Lord is neglecting his affairs. The Lord seems to deal unjustly with his servants; he appears utterly to forget them. There is no doubt that many, by murmuring, offend the Lord grievously. Now therefore, we are taught to have hope and patience.

### Paragraph 480

And at this present, heaven is opened to us, and showed us what is to be seen,[note 32.4] namely, the souls of those who have been slain in persecutions, and their state is declared. Moreover, that God does not forget to avenge; why also he delays it, and for how long. These things are spoken for the consolation of all the faithful who are now afflicted by persecution. Far different things are exhibited here than those painters, instructed or rather corrupted by monks and friars, have shown us: namely, a great company of monks and nuns covered in heaven with our Lady’s cloak, as though the greatest part of them were to be saved. John shows us never a friar, but rather many martyrs, whom the friars at this day make before other men. Hereof, therefore, as of the doctrine of truth, we shall learn what state or degree is most plentiful in heaven, not that we should think no man but only martyrs are to be saved (for as many as truly believe in Christ, and crucify their flesh with its desires, shall be associated with the holy martyrs, and rejoice with Christ forever), but that chiefly the holy martyrs are saved, whom the mad world supposes to be lost.

### Paragraph 481

But all things here must be examined by us with the greatest diligence. For this passage is both very clear and full of most wholesome doctrines. First, John sees and shows us, as it were pointing with his finger, the souls, and those of those who were slain—that is, the spiritual and immortal substances which remain alive when the body is lost and consumed. The body may be killed, but the soul cannot be killed. Our Savior has vividly expressed this in Matthew 10. In Luke 12, he says, “Do not fear those who kill the body, and after that can do nothing more,” and so on. Therefore, tyrants could rightly kill the bodies of martyrs; they had no power over their souls. This passage clearly demonstrates that the souls of men are not only immortal but also living or watchful, not sleeping, and remain and live in heaven. For some believe that the souls, when they depart from the body, sleep—which is a most vain notion.

### Paragraph 482

Now the reason for the martyrs’ deaths is also shown: they were put to death for the word of God and for the testimony they had. They were not put to death for their wickedness or evil deeds, but for the true religion, by which they confessed and preached that word of God which was in the beginning and was made flesh, and the Gospel which they had committed to them, the testimony of God and eternal life, which also they ministered and preached. Of the word of God and the testimony of Jesus Christ I have spoken in the first chapter. For no other cause are countless bishops, kings, and princes slain today. If they were adulterers, usurers, blasphemers, and wicked doers, they would be regarded with some measure of tolerance; but now, because they profess the only Son of God and preach the Gospel, they are murdered without mercy. Here we have certainly defined who are true martyrs in deed, not those who suffer torments, but those who are tormented for God’s word. For the cause makes the martyr.

### Paragraph 483

But where are the souls of those who are slain for the word of God? They are shown to us under the Altar [note 32.7]; the Altar, as in the eighth chapter, is set in heaven, before the throne of God. Therefore, the souls of all the saints are in heaven, before the throne of God, which was also signified before in the vision of the twenty-four Elders. The Lord has said also, “Where I am, there shall also my servant be.” But the Lord is in heaven; therefore, the souls of the faithful, whose bodies have been slain [note 32.8], or buried without slaughter, are nowhere else but in heaven. Nevertheless, it does not lack a singular mystery that they are laid under the Altar, as under a shadow, through whose benefit the souls may be well at ease. I told you before, and here again repeat, that the altar signifies Christ. For he is also the golden altar, intercessor, and propitiation for our sins. Through the propitiation and mediation of Christ, we are received into celestial joys. And Christ is our life and salvation. Under him we lie hid, as under a cover or a shadow. Thomas Aquinas, expounding this passage from John, says that by the altar is signified Christ, in whom and by whom we should offer to the Father whatever good we do; and through him, whatever is pleasing to God is made acceptable. Under this Altar, namely under Christ, are the souls, not only in the state of life (that is, while we live here on earth) but also in the state of our homeland (that is, in heaven), as under him of whom they are covered, as under a shadow against all evil. Thus says Thomas. But I suppose that there is another thing also signified, that martyrs are made conformable to the altar, that is, to the passion of Christ, and therefore to rest now under the Altar, Christ. For those who are partakers with him in passion, do communicate also with him in glory [note 32.9]. For just as the bosom of Abraham is called a receptacle, and that port and haven of salvation, into which the souls of those who had the faith of Abraham are received, so we understand the altar to be a place of blessedness in heaven, wherein they rest who with true faith have acknowledged Christ as the altar, propitiation, sanctification, and satisfaction; and have moreover, in suffering, offered themselves to God in Christ, through patience, an acceptable sacrifice to God. Under this Altar was gathered the first martyr, Abel; and after him, as many as have died for the faith, and shall be gathered, whoever in bearing the cross, through tribulation, enter with Christ into glory.

### Paragraph 484

Now it is also declared what they do under the altar. I say, the very martyrs cry out, not the beasts, as they have done until now, and they cry out with a loud voice. Let no one imagine that the blessed souls in heaven are complaining, are sorrowful, accuse, or are troubled. These things are feigned for another purpose, so that we may understand that God does not forget his people, that he does not withhold vengeance, that he sees, feels, and regards the injuries and deaths of his servants. For when vengeance does not follow immediately, many think that God is asleep and has no regard for them. Therefore, we hear that the holy martyrs cry out, and with a loud voice. He appears to have alluded to that same in Genesis 4: “The voice of your brother’s blood cries to me, for vengeance.”

For the theologians call certain sins “crying,” as those which are red in the Scriptures cry out to God, as is presently the shedding of blood; the sin of Sodom in Genesis 9; the oppression of widows and orphans in Exodus 22; wages for work detained, Deuteronomy 24; and James 5. How long so ever therefore God delays vengeance, even if it is for many years, the blood of the just is not forgotten before God. Paul in Hebrews 12 cries out and says that the blood of Abel speaks. In Luke 18, the Lord says that the afflicted cry out both day and night for deliverance. Would God they would weigh these things, whose feet are swift to shed blood. God did not show mercy to his people in times past because much innocent blood was shed among them through the means of Manasseh, their king, as appears in the 4th book of Kings. Therefore, dear brethren, let us consider well at this day what we do, and let us not shed innocent blood rashly.

### Paragraph 485

Certainly the words are expressed of Saint John,[note 32.12] which the martyrs cried to the Lord: how long, say they, Lord, which art holy and true, &c. They put God in remembrance, not as ignorant or inconstant, but as knowing and steadfastly mindful of holiness and truth. For in as much as the Lord is holy, he hates all profane and unclean persons, and spares them not. For as much as he is true, he maintains and defends his chosen, and punishes and oppresses his enemies as he has promised by his word. Sins therefore you are such, say they, O God, why doest thou not judge and avenge our blood, of them which in earth, as in their kingdom, exercise tyranny and oppress every good man? All this signifies no other thing than that God for his own sake, which is holy and true, will never forget the injuries of his servants. Therefore we understand these things to be spoken by a figure called Prosopopeia: that is, the feigning of a person; not that the saints in Heaven do expostulate with God, but that we by such a figure might understand that God has care of martyrs, because he is holy and true. Saint Augustine in the 68th question upon the New Testament: Saying the Lord, says he, has taught us to pray for our enemies, what is the cause that the souls of those that are slain cry out as doth the blood of Abel, and require that they may be avenged? And he makes answer: Saints be not impatient, that they should urge that thing to be done now, which they know shall come to pass in the time prefixed, which neither can be prevented, nor yet delayed; but by this saying he would show, how God will avenge the blood of his servants, least because he seems now so patient, that wicked war should be thought unpunished, which is made against the saints: that both he might drive fear into them that persecute the servants of God, and might also exhort the sufferers unto patience. Thus says he. And this in deed seems the plainest sense of all others, especially if we consider the things that follow in the Lord’s answer, and it was said unto them that they should rest, &c.

### Paragraph 486

Primasius, Bishop of Utica, explaining this passage in the writings of Saint John,[note 32.13] states that it should not be understood as the saints seeking vengeance with a carnal understanding and stubbornness, for we know that through the abundance of charity even enemies are beloved in this case. Rather, it is evident that they prayed against the kingdom of sin, and earnestly desired the things of that kingdom, as we say, “thy kingdom come.” It is not lawful to think that they would desire anything against the pleasure of God, since their desires depend upon his will, and so forth. And Saint Gregory asks, what is it that the souls request when they ask for vengeance, but that they desire the last day of judgment and the resurrection of slain bodies? Arethas notes here also from the commentaries of Saint Andrew, Bishop of Caesarea: moreover, the saints appear through this to wish for the end of the world. Therefore, they are commanded to patiently abide until the fulfillment of their brethren,[note 32.14] lest they should be fulfilled without them, as the holy Apostle says.

### Paragraph 487

However, Thomas Aquinas in his explanation of the Apocalypse shows that vengeance is required of God in two ways. First, with an evil and malicious affection, which Scripture utterly condemns. Secondly, by a zeal for righteousness, and according to God’s will, judgment is required against those who are incurable. After this, he adds: Therefore, the blessed souls demand vengeance of their enemies, although they do not intend it chiefly, because of a zeal for righteousness and an affection of godly love, they resent it, as does God himself at the wickedness of the persecutors, who attack God himself, and seek to hinder his religion, and torment those who worship him; therefore, they would have their malice and power ended. Thus far he. But since Scripture consistently witnesses that the saints in heaven are free from griefs and affections, and live a new life far removed from all pain and disturbance, and have submitted their wills to the will of God, whom they may follow in all things, approving all his judgments, sayings, and doings, and reverencing them: I suppose we need not reason more subtly about this at present, but simply understand that by this figurative speech (as crying is also elsewhere attributed to the blood of martyrs shed) is signified that the blood of the oppressed shall never be forgotten by God, and that before him a just judgment and vengeance is prepared, to be executed in his time against the enemies and contemners of God: but chiefly against the persecutors of the word and the murderers of saints. This is further explained by what follows.

### Paragraph 488

For by the same reasoning that was given in response to the complaints of martyrs [note 32.16], so that we may understand the state and glory of the saints in heaven, who have offered their bodies for God’s testament: and that God has not forgotten the blood spilled, but that he will at length requite those bloodshedders when he sees fit. But where he has reserved this time to himself, when he will reward those who drink the blood, it is not our part to inquire curiously thereof, but rather to be in readiness, that if he will that we also should suffer for the testimony of Jesus Christ, we should run speedily and cheerfully through afflictions unto glory, doubting nothing that we shall be joined to the blessed martyrs in heaven, and that the just judge in that day will render to all the enemies of God, the Church, and God’s word, according to their deeds. And although the time of persecution seems long to the flesh, yet it is here and elsewhere in the scriptures called short [note 32.17]. But these things must be seen and considered in parts.

### Paragraph 489

Undoubtedly, the condition of souls in heaven is in every respect most fortunate. This is represented by the white garments. The glory of the blessed, who are now in light and feel nothing of darkness, is signified; of this garment I have spoken before. And it is expressly stated that white garments are given to each of them. For every soul receives its reward; and the body also, at the end of the world, shall receive its own garment. Saint Gregory says that saints as yet have only the blessedness of souls; but at the end of the world, they shall receive two robes or garments, for they shall be clothed with both the perfect joy of souls and the incorruption of bodies. This will be examined more carefully toward the end of Chapter 7, where this passage will be explained more fully. After it was said to the blessed souls that they should rest, therefore they are altogether at peace and feel no disturbance, which will be declared more plainly in Chapter 12. Nevertheless, this may be referred to delay and breathing, as though it were said: It was signified to the souls that they should yet wait and abide. For it follows: yet a little while. Therefore God signifies that after a little time he will deliver his own and punish the adversaries.

### Paragraph 490

And the indication of the time appears to be drawn from the second chapter of Habakkuk, a passage also cited by the Apostle in Hebrews 11. For yet a little while, because he who is to come will come, and will not delay. And the righteous shall live by faith, and so forth. In Isaiah 26, we read these words (after he has shown the resurrection and the final end of the world): “Go, therefore, my people, and enter into your chambers, and shut your doors after you; hide yourselves a little while, even for a moment, until the indignation passes.” And likewise Saint Peter called this entire time of affliction, which is a judgment, a short time, so that we might find comfort in it. 1 Peter 1 and 2 Peter 3.

### Paragraph 491

To these matters is joined another, which completes the time until the fulfillment of three companions and brethren who would also be slain for the word of God. Therefore, let us no longer inquire after the end of persecution, or why the Lord delays vengeance, or how long [?]. For we hear that the number of the elect must be accomplished. But since that time is known to God alone, let us not be curious, but let us think of those things that concern our duty, so that if the situation requires it, we may also die bravely for the testament of our Lord God, that we may be associated with our brethren and companions, and have the blessed sight of our redeemer. The number shall be accomplished at the end of the world, at the last judgment. Until then, persecution shall continue; but then assuredly the Lord will requite it, as the prophet Malachi has witnessed in the third and fourth chapters.

### Paragraph 492

Hereof we learn also what we should judge of the holy Martyrs in Heaven, and blessed souls: the same that we learn here of God’s word. Brethren and friends, fellow servants (as also angels, in the 19th and 22nd verses of the Apocalypse) they are expressly called, not Lords and founders. For although the words must be understood of us yet living, there is a relation. For if we be their brethren and fellow servants, they are verily our brethren also, and God’s servants with us, and even fellow servants. Now though we grant that they pray in Heaven, what, I pray you, do they pray for here, but that God would avenge and punish? And what do they obtain? As we read that Christ said to his mother, requiring wine at the marriage, “Woman, what have I to do with thee? My hour is not yet come,” so likewise are the Martyrs here commanded to tarry and abide the time of God appointed. Which we believe that the Saints do. What can we say then of their intercession and praying for sinners unto God? There is one only mediator given, even the Lord Christ; let us go unto him in all our necessities, he alone shall suffice all, and in all.

### Paragraph 493

These things have been spoken up to this point concerning the persecutions of all times, so that in the meantime they have provided the most comforting consolations to all who suffer persecution until the end of the world. They have also silenced curious questions and left us safe and whole in the will of God, in which, resting alone, we may know that it is best for us.

### Paragraph 494

Therefore, it behooves us to gather certain sure sayings, with which to comfort ourselves as with the most certain sentences of God pronounced. First, that God is true and just: and therefore not to neglect his, but to cherish with fatherly care. And if he casts us into any danger or difficulty, that verily shall turn to great profit for the godly. If he shall take us away by torments, that he delivers us from evils, from miseries, and corruption of this world, and renders for the same everlastingness. Secondly, it is certain that God is just and true, that he will requite the wicked after their deserts. Again, if he makes men fortunate in this world, that in deed appertains to their destruction. Where he is slow to punish, that is done through God’s long suffering: but that God recompenses this slowness with the weightiness of the punishment, in case they be incurable. Whereas these things undoubtedly are most certain, what remains but that we should commit ourselves and all ours to the Lord our God? He knows the time and means whereby to avenge his, and to plague his enemies. To him be glory forevermore. Amen.

## Sermon 33: ¶ The ſixte seal is opened, and the corrupting of the sincere doctrine is exhibited.

### Paragraph 495

.

### Paragraph 496

And I saw, when he had opened the sixth seal, and behold there was a great earthquake. And the sun was as black as a sackcloth made of hair: and the moon became as blood; and the stars of Heaven fell to the earth, even as a fig tree casts its figs when it is shaken by a mighty wind: and Heaven vanished away as a scroll when it is rolled together.

### Paragraph 497

The opening of the sixth seal of the lamb [note 33.1] generally reveals to us a scene visible to all men, depicting the corruption of doctrine within the church, with a mournful and terrible effect. Nor is there anything said here in the sixth seal, as also in the five preceding seals, other than what our Lord Jesus Christ foretold previously in Matthew 24: that the Gospel should be preached throughout the world, and that wars, famines, pestilences, and grievous persecutions would arise, along with false prophets who would deceive people and drown them in great sorrows.

### Paragraph 498

Nevertheless, these things must be explained with religious care. For it is not to be thought, simply because the Lamb opens the sixth seal and the Son immediately appears dark, that Christ is the author of corrupt and evil doctrine. For Christ is the one who sows good seed in the field; the enemy sows darnel, as the Lord himself explains in Matthew 13. Christ teaches sound doctrine through the Apostles and sincere preachers; and when it seems vile to the world and cannot please, by his just judgment he leaves the contemners to their affections, as the Apostle Saint Paul says, he sends upon them the power of delusion, that they may believe lies, and so be judged—all those who would rather believe a lie than the truth. And the deception through corrupt doctrine is a more harmful evil than the bloody persecutions. Indeed, the seducers and false prophets have done more harm to the church than cruel tyrants. Finally, people are more grievously punished when they are abandoned to be seduced by deceivers than when they are subjected to being torn into pieces by their murderers. Therefore, it is a most grievous plague of God, and utterly to be abhorred, when the simple truth is despised and delivered to lying deceivers, who, it may be spoken with reverence, seem to all to be [illegible] and be peace. For where the gospel is purely preached to many, these men say, “I do not understand what these men teach us from the Gospel; but I can see that the old [?have] all to be raided, and these new both to be peace and be[illegible].” Therefore, you shall have teachers who will perform to you in deed the same that you talk. Would God we lacked examples, and did not see certain nations that formerly had the free and pure preaching of the gospel, and now stripped of all truth, sigh under the pleasure and boldness of most wicked deceivers, who trample God’s word under foot, condemn it as heresy, and stop the mouths of the wretched people full of human dung. This is the punishment of the truth despised.

### Paragraph 499

And this passage may not be explained as referring to a single age, for the sins mentioned are things rehearsed in general; but rather to the entire time stretching from the age of the Apostles to the last judgment. It contains therefore the corrupt doctrine of Valentine, Marcion, Manichaeus, Arius, Macedonius, Nestorius, Eutyches, Donatus, Pelagius, the followers of Priscilla, and finally of all heretics, and the confused mixture of Mahomet composed of the same, and chiefly the sophistry and most corrupt doctrine of Antichrist and of his ministers.

### Paragraph 500

But when the Lamb opened the sixth seal, there was not heard now as before the voice of the Beasts, Elders, or Martyrs, but a terrible earthquake. An earthquake in the Scriptures signifies a wonderful commotion of all things, troubles, tumults, and great alterations. And truly, greater darkness arises from nothing else than by altering godly religion and receiving of wicked doctrine. For so arise sects, seditions, and wars. You may see many examples of this in the history of the ancient people, who are recorded to have been grievously shaken so often as they have changed their religion and kind of doctrine. By this earthquake, therefore, is signified that exceedingly great trouble shall arise from this, for that a new and a strange kind of doctrine should be brought into the world, amiss inciting men.

### Paragraph 501

Here is what you may say to those who blame the Gospel and its preachers for all the current troubles, seditions, and upheavals in the world. Elijah once made an answer for us that will suffice for all times; the passage is in the third book of Kings, chapter 18. “I have not troubled Israel, but you and your fathers’ house, who have forsaken God, &c.” The story of Jeremiah in chapter 44 is also relevant, where all the evils then afflicting the wicked were wrongly attributed to sincere doctrine and to the prophet Jeremiah. Learn also what to say to those who claim that God has allowed his Church to lie and rot in errors for so many years [?].

### Paragraph 502

And the corrupt doctrine is described by parts, even from the top to the toe,[note 33.7] and the effect also of the corrupt doctrine is annexed. And first of all the sun, a planet most bright, not only waxes dark, but black also. And immediately is added an Image or a parable, like a hairy sack which is woven or made of hairs or of bristles. The sun lights and gives life to the world. And through Christ, which is the life of the world, we are illumined and quickened. He casts abroad from him the bright beams of the Evangelical truth. And just as Christ is not darkened in himself, so neither is the truth of the Gospel, which by nature is without pollution. By reason of the black clouds that override it, the light of the Sun waxes black and is hindered; and from the traditions of men, and the corrupting of the scripture arises darkness and blackness in matters of religion. The Gospel of itself is bright and wholesome. Christ is light, full redemption, health, and life most perfect. But when men had rather seek of others doctrine, life, and salvation, than of Christ and his wholesome Gospel, most thick and gross darkness arises in the minds of those men. For there is established another doctrine, righteousness, intercession, redemption, salvation, and life than that of Christ. They that receive that doctrine,[note 33.8] seem to have put on them a shirt of hair which pricks, burns, and vexes continually. For there is no rest, quietness, security, or spiritual pleasure and refreshment from corrupt doctrine, but only tediousness. Christ pure, and sincerely received, is to man a joy unspeakable, and a most bright and joyful light.

### Paragraph 503

After this, it is added that the whole moon, not a part only, is become bloody. For an image is again annexed, as blood. The moon receives light from the sun and is subject to courses, or changes, while one while it increases, and another while it decreases, and signifies the church. The church, established upon the rock, is not unstable; but because of its variable fortunes, is subject to many diverse chances. For now the church triumphs, and immediately being oppressed she mourns; now she increases in number, and immediately she is diminished. And the church is illuminated by Christ. But while the Sun itself is darkened, the moon cannot but be most obscure. Blood in the scriptures betokens great wickedness, chiefly idolatry and false worship of God. The Lord in Leviticus 17 says that he will account strange worship as blood. Therefore, when faith and knowledge are darkened in Christ’s church, it cannot be chosen but that blood shall arise in the universal church; that is to wit, the corrupt worship of God, which the Lord esteems as murder; there must needs innumerable sins and wickedness spring thereof. For the lively doctrine of Christ being corrupted, all things must of necessity be most corrupt, and swarm full of superstitions and iniquities.

### Paragraph 504

To these things is added another thing, which supports what has been said: stars fall from Heaven[note 33.11] to the earth. Daniel called stars preachers in chapter 12. As also Saint Peter, in 2 Peter 2. Therefore, the preachers of churches depart from the heavenly doctrine of Christ, which was brought and revealed from heaven, and leads people to heaven, and keeps them in heavenly conversation; and they embrace earthly things, that is, the doctrine of men. By this, it comes to pass that both the sun is obscured and the moon is made bloody. Stars shine; preachers should proclaim Christ, the true light, to the whole world, but they have neglected this, being devoted to their own traditions.

To these is also added an image,[note 33.12] which falls as the fig tree casts off its figs when shaken by a strong wind. Here is signified the corruption of preachers, and a great number of them. For the fig tree was made to produce sweet fruit; so also the ministry of the word was ordained for the salvation of men. However, the figs were unripe, therefore they remain green or untimely fruit. This signifies that the preachers were not ripe in true knowledge of Christ, and therefore are to be shaken down by every wind of doctrine, both that which they have admitted and that which they have proclaimed as earthly. The abundance of false teachers is signified by the fact that the unripe figs fall down in great abundance.

### Paragraph 505

From these things now follows another, and Heaven departed, seeming to flee from human sight and vanish. Again, an image or likeness, like a scroll folded or rolled up, is added. Heaven in the Gospel signifies many times the kingdom of God. Therefore, the kingdom winds itself up on earth, and the church seems to hide herself, not that there will be no church at the last (for the church will always endure until the end of the world), but because, at the end of the world, the church will lie hidden, and it will not be considered the true church, which is indeed the true church. The letters and words are not wiped out of the book, but are not seen, rather they are hidden when it is rolled up. It is manifest today what John meant by this parable. For all, in a manner, judge the newly arisen Roman church to be the true church, which in truth is not the church of Christ; and the church which is the spouse of Christ is judged heretical. Therefore, the church is wrapped up and as rolled together. The Lord unfold and preserve it. Amen.

## Sermon 34: ¶ The effect of corrupt doctrine is expounded, and that the Angels let that the wind blow not.

### Paragraph 506

.

### Paragraph 507

And all mountains and islands were moved out of their places. And the kings of the earth, and the great men, and the rich men, and the chief captains, and the mighty men, and every slave and every free person hid themselves in caves and in the rocks of the hills, and said to the hills and rocks, “Fall upon us and hide us from the presence of him who sits on the throne, and from the wrath of the Lamb,” for the great day of his wrath has come. And who can endure it?

### Paragraph 508

And after this, I saw four angels standing on the four corners of the earth, holding back the four winds of the earth, so that the winds would not blow on the earth, or on the sea, or on any tree.

### Paragraph 509

Now follows the effect of the corrupt doctrine in men. And hills and islands are moved out of their place; wherein is also a respect had to the earthquake, as though by the earthquake they were removed from their place. And mountains and islands do betoken realms, nations, and people, so steadfast in faith, that as mountains and islands are immovable, and are not shaken by the storms of the sea, so these might seem to be immutable. Nevertheless, at the alteration and corrupting of doctrine, they are now also removed out of their place, and quite overthrown. And such as read histories shall find everywhere that such have been deceived by the crafts of heretics, by the power of Muhammad, and by the hypocrisy of the pope, whom you would not have thought should have been abused, in so much that whole cities and realms have clean revolted. For seducing is of efficacy namely in such as now begin to slip and slide from the rock of the church.

### Paragraph 510

And those who, being unsettled, are removed from the firm foundation, retreat into caves and rocks of hills. For it is impossible for one who does not hold Christ with steadfast faith to find rest. For like a raging sea, he is tossed to and fro. For since he does not have a sure and certain way of life, nor commits himself to be ruled only by the Scriptures, so that he might hold the certainty, he is content to be led by anyone he encounters. Therefore, we see those to whom Christ alone is not all seeking salvation in pilgrimages, in hermitages, in monasticism, in stricter discipline, in satisfactions, and I know not what other follies, or rather blasphemies. And these are truly said to hide themselves in dens and caves of stone, and think they may lie hidden safely within them, make satisfaction for their sins, and please the Lord.

### Paragraph 511

But in describing many kinds of men, he encompasses all states in the world. For of all sorts of men, there have been found not a few, nor of low status, who have taken upon themselves the hermitic and monastic life, and have bound themselves to a strict kind of living. Here, therefore, are reckoned up kings, judges, great men or princes, rich men, merchants, captains over thousands or chieftains, valiant men in this world, bond servants, and free men, whom we call at this day gentlemen. But how many kings and princes and noble gentlemen are depicted in the churches of abbeys, painted in pictures and hung on trees, who have lived some time a monastic life?

### Paragraph 512

But entering monasteries, woods, and wilderness, and adopting a stricter life, along with various satisfactions, pilgrimages, and similar disciplines, have not yet led to a quietness of mind; indeed, they are now more afraid than they were before, and have fallen into utter despair. For in these things wherein they sought quietness, they have found none; no, besides Christ, there is found no quiet nor rest. Those who live in such strictness under the unhappy papistry understand this very well. And the words which John the Apostle recites here are those of people in the greatest distress, even to desperation, where they cry out, “Let the hills fall upon us,” and so forth. For so this word is also used by Hosea in chapter 10 and by Luke in chapter 23. This signifies a conscience most afflicted and most entangled, finding nowhere any comfort or consolation, but desiring nothing other than present destruction, so as to be delivered from the present evil and intolerable grief of mind. Such a thing is that of Turnus with Virgil in the tenth book of the Aeneid.

### Paragraph 513

Moreover, the causes of this fear, despair, and hiding are the appearance of the one who sits on the throne, the wrath of the Lamb, and because they perceive that they cannot stand before God in the day of wrath and God’s vengeance. Therefore, they flee from the face of God, they flee from the Lamb, that they might avoid the vengeance, if they could escape it. The fear of God is commended to us in the scriptures, and those who do not fear God are condemned; but the scripture speaks of a fear joined with true faith and love. For John says, “Love casts out fear.” Even so, the same scripture preaches to us God as just, and shows him to be angry with sin; nevertheless, it declares him to be benign and merciful to those who acknowledge their sins and ask for forgiveness, that his only begotten Son is given by God to humankind, by whose mediation we may come to the throne of God, which otherwise no man may attain to. It preaches Christ, the Son of God, to be the Lamb, that is a propitiation for the sins of the whole world; and that the same calls all to him, excludes no one, but promises and offers to all, all things of life and salvation. But whereas corrupt preachers, friars, and popish priests have forsaken this simple and most pure doctrine, wholesome and full of consolation, and show that God is like Rhadamanthus, a judge inexorable, and set forth Christ rather as one angry than favorable, they do alienate, without doubt, the minds of men from God; that now they may say expressly, “Who is worthy to come into the sight of God?” No man shall be saved before this God most severe, and his Son a judge most righteous. They turn them therefore to sundry means of salvation, they choose mediators and intercessors by whose mediation and merits they may redeem to themselves the favor of the angry deity. But those who rely on anything other than God’s only mediation and intercession of the Son are disappointed of their purpose, and at the length fall into that same desperation. When they perceive that the monastic life and their merits cannot stand before God, they flee from the face of God; and tormented with the pricks of their conscience, they know not what they may do, whither they may turn, where is the true salvation. Therefore, we judge those who through Christ acknowledge the Father as a Father, and through Christ have access to the Father, as favoring them and loving them, to be rightly most blessed; acknowledging truly in the fear of God their sins, but yet with a true faith hoping for remission of sins, knowing that they are through Christ reconciled to God the Father. The monastic, heretical, expiatory, and Pharisaical faction fully acknowledges this doctrine, therefore they are tormented with sorrows that cannot be uttered. If I speak not here of the monasteries or monks of this our time, in whom we see almost no conscience at all, nor other intent than to be addicted to idleness, voluptuousness, and to bear rule. In times past were found men full of conscience, entering into cells and woods, for no other cause than that they might so be saved. Of such spoke the Lord in the gospel: “When they shall say,” he says, “Christ is in the wilderness, go not forth,” and so forth. And I doubt not, but that some simple also at this day for this intent take upon them the monastic life; but they shall find also, as John said here they should prove and try by experience.

### Paragraph 514

Furthermore, this passage might seem to call for an explanation of the signs preceding the final judgment, and of the terror of the wicked, as the Lord preached in a similar manner in Luke 21. But the final judgment will be discussed more fully in chapters 11 and 19 of this book, and elsewhere. And while I do not condemn such an interpretation, it seems to me that the general destinies of the church are here presented together, in which, where corrupt doctrine occupies not the last place, nothing should be spoken of it in general; many things will be addressed in particular in chapter 8 and those that follow, unless this present passage is explained in the same way as it is. Moreover, those things that follow will be better connected, which will not be part of the final judgment, as the matter itself will demonstrate.

### Paragraph 515

And the things that follow in the 7th chapter pertain to the explanation of the sixth seal, or to the discussion of it. And it chiefly recounts three things: how the angels restrain the winds so they do not blow; an innumerable company to be sealed in the midst of corrupt doctrine, which should not perish. And what the state is of those who have departed out of this world either by martyrdom, or else being either undefiled by the corruption so full of enormity, or delivered and purged from the same; which are added for consolation. For this book of Revelation is wonderfully evangelical, most full not only of prophecies, but also of admonitions, exhortations, and most comforting consolations.

### Paragraph 516

First, it is to be explained what is spoken of the restraint of the winds by the angels, so that they should not blow. Wind, as also appears in the scriptures, is used in both a good and an evil sense. For wind is called vain and false doctrine, and a hope conceived of erroneous doctrine, as in Hosea 12, 5, and 22 of Jeremiah. So leaven is called the Pharisaic doctrine, and the hypocrisy springing from it. Paul in 4 Ephesians forbids that we be carried about with every wind of doctrine. And the Holy Spirit is symbolized by wind in chapter 3 of John. And in chapter 2 of Acts, wind is subtle, it pierces, is felt, and is not seen; great is its force, it cools, it dries, gathers clouds, which bring rain and make the earth fertile. Therefore, rightly, by wind is signified the Spirit of God, and the sound doctrine which is of the Spirit of God. Therefore, it is one wind, the Spirit of God which inspires; and there are four winds, that is to say, many from the corners of heaven and parts of the earth, that is to say, preachers dispersed throughout the whole world. Therefore, the doctrine of the Gospel, inspired from all parts of the world, blows or is preached, so that there are many winds, yet all proceeding from one. For there is one and the same Spirit which speaks by the ministers, and gives them sundry graces, 1 Corinthians 12. Briefly, by the blast of winds we understand the free preaching taken from the holy scriptures.

### Paragraph 517

Secondly, we must know that both good and evil angels are mentioned in the Scriptures. Angels, as appeared before, are called ministers. And there are good and evil ministers: the good inspired by God and the good angel; and the evil by the evil angel. And the enemy of the truth stirs up men in all places of the world, in the courts of kings, in the places of judgment, in schools, in colleges, in cities, towns, and villages, which may hinder the free course of God’s word. Therefore, the proclamations of kings and bishops fly to and fro, are proclaimed and set up, prohibiting the reading of the Bible, the preaching of the Gospel, and so forth. And to have some pretense for their evil doing, they forge that the Bible is corrupt in a thousand places, and that heresy is learned and taught out of the same. Therefore also they prohibit and condemn the Bible and the books of the Gospel, the unworthiness of which thing cannot be sufficiently spoken before the church. They do the same that in times past Antiochus, Epiphanes, Diocletian, and other men of the same sort are recorded to have done. The expositors of the Bible in times past deserved exceeding great praise; neither was there any faithful who said the holy book to be corrupted, for that all translations did not agree among themselves. We live therefore at this day in a time most corrupt and most unthankful.

### Paragraph 518

The suppression of reading holy Scripture is the foundation of corrupt doctrine, and of confusing the conscience, and of the despair that follows from it. And by the Earth he understands those men dwelling on Earth; by the Sea and Isles, those men of islands, and those who dwell on the Sea; by trees, men, everywhere shadowed in Scripture by trees. For unless the winds blow, the trees do not flourish, nor does the earth grow green. The Prophet says, “Send forth your spirit, and they shall be created, and you shall renew the face of the earth.” And except the word of God be preached, the minds of men do not grow green, nor are the fruits of good works brought forth by men. And therefore, the angels prohibiting the wind are said to do harm; for in truth, there is nothing more pestilent nor pernicious than the suppressing of the free preaching of God’s word. The Lord by his spirit renew all parts of the world. Amen.

## Sermon 35: ¶ The faithful are sealed to salvation, which they obtain by the grace of God in Christ Jesus.

### Paragraph 519

.

### Paragraph 520

And I saw another angel ascend from the rising of the sun, who had the seal of the living God, and he cried with a loud voice to the four angels (to whom power was given to hurt the earth and the sea), saying, “Hurt not the earth, neither the sea, neither the trees, until we have sealed the servants of our God in their foreheads.” And I heard the number of those who were sealed, and there were sealed one hundred forty-four thousand of all the tribes of the children of Israel. Of the tribe of Judah were sealed twelve thousand; of the tribe of Reuben were sealed twelve thousand; of the tribe of Gad were sealed twelve thousand; of the tribe of Asher were sealed twelve thousand; of the tribe of Naphtali were sealed twelve thousand; of the tribe of Manasseh were sealed twelve thousand; of the tribe of Simeon were sealed twelve thousand; of the tribe of Levi were sealed twelve thousand; of the tribe of Issachar were sealed twelve thousand; of the tribe of Zebulun were sealed twelve thousand; of the tribe of Joseph were sealed twelve thousand; of the tribe of Benjamin were sealed twelve thousand.

### Paragraph 521

After this I beheld, and lo, a great multitude (which no man could number) of all nations and people, and tongues, stood before the throne and before the Lamb, clothed in long white garments and holding palm branches in their hands, and cried with a loud voice saying: “Salvation be ascribed to him that sits upon the throne of our God, and unto the Lamb.” And all the angels stood around the throne and of the elders and of the four living creatures, and fell down before the throne on their faces, and worshipped God, saying: “Amen. Blessing, and glory, wisdom and thanksgiving, and honour and power, and might, be unto our God for ever and ever. Amen.”

### Paragraph 522

We have heard, brethren, concerning the opening of the sixth seal [note 35.1], that the sun was made black, the moon bloody, and the stars fell from heaven to the earth, and the rest that we have rehearsed. By all of these things was signified the corruption of doctrine. A sorrowful and fearful matter was shadowed with most sorrowful and most terrible parables. We have heard how there followed in the world a great turmoil of things, and with it, grievous despair; and that the winds also were restrained, that they should not blow. But we have experienced how great a grief it is, yea, and destruction also, to lack the air or wind, so much that without breathing and cooling, men must needs wither and be weakened and choked. But with so great an evil are those vexed who are destitute of the preaching of God’s word.

### Paragraph 523

Some person here might say that the entire world will perish in heresies, in the Alcoran, in Papistry, and other corruptions. In what state, then, were our forefathers? Do you think they are all damned? John prevents these things, and with a vision wholly Evangelical—that is to say, with a most profitable consolation—shows that God has an innumerable multitude of them, who even in the midst of those Antichristian times or difficulties are made safe: and that by the mere grace of God, through the intercession of Jesus Christ, of whom alone is salvation: that is to say, whom alone those who are saved may thank for their salvation.

### Paragraph 524

We now have to respond to those of opposing views, who constantly object,[note 35.3] arguing that if our ministers are condemned, it would be wicked to condemn all, and therefore they must be saved. They claim they have not heard of our new doctrine, but have maintained the old, and therefore will be saved by adhering to the old ways. To this, we reply that our ministers were saved, we gladly grant and believe it as well; but we add that it is by the free grace of God, as we shall understand more clearly presently, and not by popish superstition. Nor will you be saved by that means; but you must also be saved by Christ, if you desire to be saved.

Rather, considering the singular goodness of God today, the gospel is preached, and is preached even to you, to which you show yourself a rebel, declaring yourself not to be among God’s children, who hear the word of God with joy and keep it. You will have no excuse or pretense to cover your sin. If your ancestors had had the same opportunity that you neglect, good God, how far they would have run before you! Therefore, you knowingly and willingly speak against God and wilfully cast yourself into destruction. Die, then, through your own fault.

### Paragraph 525

Neither does this passage testify only that many will be saved by God’s grace from corruption and in the true faith, even when, in human judgment, there appear to be none or very few faithful; and even when very few or none are saved, because of the exceedingly great corruption of every time? We have also heard and read in the third book of Kings, chapter 19, that Elijah, complaining most grievously of the scarcity of the faithful, understood that God had reserved seven thousand men who had not bowed their knees before Baal. Therefore, the Lord always has his chosen, who, by grace through Christ, are saved in the midst of destruction and perdition.

### Paragraph 526

[note 35.6]And the Author of this salvation and preservation is first declared to be an angel ascending from the rising of the sun: namely, the Lord Christ, that sun of righteousness, rising up in those most intense Antichristian darkness, to those who seek God, and driving away the darkness for them. For Christ is the true light of all times, illuminating all those who are illuminated. He gives his people also preachers, who by the word may defend God’s people, that they be not destroyed with that common destruction.

### Paragraph 527

For it is diligently expressed that this angel had a seal, and not a seal only, but the seal of God, and even of the living God. For Christ, who is the image of God unseen, that is to say, the print or express image of his substance, in whom we know, as he himself says to Philip, the Father, has a seal, which is an instrument with which we seal such things as we will have sealed, saved, and confirmed, and discerned from that which is counterfeit, and kept safe against deception. But the Lord has no such seal as we have in this world; but so by a figure is called the spirit of God, with whom he inspires his faithful, by whom he gives also a lively faith, by the word of the living and eternal God. This seal therefore is the seal of the living God, the spirit of life, and lively faith. The Apostle Jerome, speaking: “We also trust in Christ, after the word of truth heard, the gospel of your salvation, wherein after the believed, they were sealed with the holy spirit of promise, &c.” These things are not divided. For faith is not without the word, nor both these without the holy ghost in the faithful. For Christ works with me by a lawful ministry, by whom he inspires certain that may teach and admonish men, unto whom he gives his faith and spirit, sealing their minds. Christ therefore doth prohibit the ministers of Satan, that they in restraining and letting the free preaching of God’s word, should not proceed to hurt men, before the minds of the chosen be sealed; that is to say, teaches how soever the truth is restrained, and the preaching of the Gospel obscured, yet that the minds of many shall so be furnished with God’s word, and with godly inspiration, which may so live, and be of such efficacy in them, that seducing can either have no place in them, or if it have any at all, can not abide or persevere to the end.

### Paragraph 528

There are also two other places in Scripture [note 35.9] testifying that seals were given to people, with which they were marked and delivered from present evil; and they do not contradict our seal of the living God. In Exodus 12, the doorposts of the Israelites were sprinkled with the blood of the lamb. The sign itself would have availed nothing, unless the virtue of God, instituting and consecrating the sign with his word, had turned away the destroying angel; nor has the sign lacked faith when used by God’s holy people. For the godly do not receive God’s ordinances without faith. Therefore, the same power of Christ preserved the Israelites from destruction, which now keeps the faithful from the infection of Antichrist. In Ezekiel 9, one seals the foreheads of the faithful, having the stylus of a scribe and priest. Truly, Christ has always defended his people. And he seals by impressing or writing this mark or letter Tau. That mark signifies, that is to say, the Law, or direction, or Rule. For whoever has the law of God, the word of God, and even the rule of faith engraved in his heart, is safe and secure from all evil. The ancients in old times called the rule of faith and direction, the very articles of the Christian faith—"I believe in God," and so on. You see, therefore, how all these signs in truth come to one point: those whom the Spirit of God has inspired and illumined with faith by the word are safe and secure from evil. This much concerning the seal.

### Paragraph 529

Now let us also consider who they are that are sealed. We read in Ezekiel, pass through the city of Jerusalem and mark Tau in the foreheads of those mourning and lamenting for all the abominations done in the midst thereof. And here it is said, till we seal the servants of our God. Therefore, the servants of God, and those who are sorry for abominable wickedness, are sealed. The contemners of God—hogs and dogs—are neglected.

### Paragraph 530

Furthermore, it is shown in what part of the body they are sealed. [Note 35.11] In times past, the blood of the lamb was anointed on the doorposts. In Ezekiel, Tau is marked on their foreheads. Here also, the seal of the living God is imprinted on the foreheads of the faithful. And the forehead represents the mind, the chiefest and most excellent part in man. Spirit and faith are put into the minds of the faithful. Nevertheless, the mark is aptly fixed to the forehead, not to the back, or shoulders. For those enlightened by the word and spirit, and who have faith, do confess the same, and dissemble nothing; and much less are they ashamed, but desire that their glory, which is their faith, might be known of all men. We display most notable things written on our foreheads: that is, most manifest things, whereof we are not ashamed.

### Paragraph 531

If we now apply these things to events that occurred in the past and continue to occur today, they will shed great light upon them. There were found good men, faithful and fearing God, mourning and serving God. And there are still people today, amidst Islam and Roman Catholicism, who explicitly condemn and have condemned this kind of life, openly confessing that it is not the true way of life, that there is no more wicked group of men than their priests, and that they would not entrust themselves or their salvation to them, but rather would consecrate themselves wholly to God’s mercy. And others, who have spent a significant portion of their lives with good zeal, though perhaps not according to knowledge, in those trifles and superstitions, when they reach the end of their lives, despise them altogether; and freely professing the truth, they condemn those trifles and commit themselves wholly to the Christian faith, esteeming nothing more excellent or surer than the rule of faith, which they also desire to hear recited to them as a true confession, and die in the same. All these have the mercy of God sealed with the seal of the living God, and delivers from all taint of the Antichrist and Satan, from corruption and destruction, through Jesus Christ our Lord.

### Paragraph 532

But lest we should gather in every age only here one and there one, the Lord himself now makes here a great reckoning, first of the Jews by every tribe he gathers twelve thousand, and after by multiplication, one hundred forty-four thousand; and of the Gentiles a multitude innumerable. Wherefore in every time and age innumerable obtain salvation: however much error, seduction, and destruction reign and rule in the world. These things do highly commend God’s mercy and comfort us exceedingly. And where certain gather hereof that there shall be yet in this world before the judgment a Saturnalian or golden age, wherein these things should be fulfilled, and that all men should come to the kingdom of God, it alludes too much to the gross error of the Millenniums, which is already expelled out of the church of God. These things were fulfilled in old time, and are at this day, and shall be fulfilled likewise, so long as the world shall endure. The kingdom of Satan and of Antichrist shall continue always to the last judgment, and shall still impugn the kingdom of Christ, and seem even to oppress the same; much less ought they to promise us so great security. When the Son of Man shall come, says the Son of Man himself in the Gospel, think you shall he find any faith upon earth? And again: it shall be as in the days of Noah and Lot, the words of the gospel are known, as also those of the blessed Apostles Peter and Paul, 2 Peter 3:1, 1 Thessalonians 4.

### Paragraph 533

[note 35.14]Those who do not like this explanation of ours briefly argue that the prophecies of the prophets concerning the restoration of Israel have not yet been fulfilled, and that, according to the truth of the eternal God, they must be fulfilled. They suppose, and indeed contend, that a certain or predetermined time must remain, in which all these things may be accomplished. To this I answer plainly, that we shall err shamefully with Papias, Augustine, Irenaeus, Tertullian, and Lactantius, and with those who are called Millenniums, unless we judge rightly here. I believe, therefore, that the same restoration of which the prophets speak must be divided into three times: the first to be historical, extending from King Cyrus to great Pompey, and which Ezra, Nehemiah, and the Author of the book of the Maccabees describe and teach to be fulfilled. The second to begin with the coming of our Savior and proceed until Antichrist and his destruction, which the Apostles and Evangelists have most diligently described, and in which they testify many things to be accomplished. And that the third time should begin from the restored gospel and the last judgment, and continue forevermore: which restitution truly seems to be of all others most perfect and complete, in which God will give to humankind fully what things ever he has promised by the mouths of the prophets and Apostles. Saint Peter has manifestly made mention of this in Acts 3:21, saying: “It behooves Christ to take heaven, until the time of the restoration of all things, which God has spoken by the mouth of all his saints from the time of the prophets.” And the Lord himself in the gospel, speaking of the last judgment, said: “Lift up your heads, because your redemption draws near.”

### Paragraph 534

Perhaps we may divide this matter in this way for greater clarity: the restoration of Israel, or of all the faithful, is certainly either physical or spiritual. The physical restoration may be called historical, and was accomplished by Cyrus, Zerubbabel, Joshua, Ezra, Nehemiah, and the Maccabees. The spiritual restoration is being fulfilled, or will yet be accomplished, by the coming of our wholesome Messiah, our Lord Jesus Christ. And the coming of the Lord is of two kinds: the first, indeed, is in the flesh, in which we believe many things, the Apostles bearing witness to their fulfillment in Christ. In the latter, he shall come again from Heaven for judgment. In that coming, he will most fully accomplish those things that we see as yet unfulfilled. And without doubt, all our hope is directed toward this coming and is comforted by it. Those things spoken by the Apostle in Romans 11 concerning the conversion of the Jews are being fulfilled partly, and partly are being fulfilled daily, and will continue to be fulfilled.

### Paragraph 535

Now we return to the abundance of those who shall be saved and are already saved from the midst of the kingdom of Antichrist, to be declared. John divides the universality of humankind into Jews and Gentiles. Of the Jews, one hundred and forty-four thousand are accounted. And according to our judgment, of a thousand Jews, scarcely one or two seem to be saved; but where, by the testimony of our Savior himself, so great a number is saved, there remains certainly, of this number, an infinite multitude of this stubborn people to be gathered, who shall be saved. And they are not saved by the Law, or by circumcision, or by their damnable obstinacy, but by the grace of God in Christ their Messiah, the only Redeemer, revealed to them mercifully, and received faithfully by them. For if the thief on the cross might be saved, now leaving his life, what shall prevent innumerable Jews from being saved by the same means? Nevertheless, I will determine no measure here. Neither will I by this means frustrate the ministry of the word and Sacraments. However, I know the things here spoken to be true; the measure or manner is known to God, and with him nothing is impossible. And to this serves the Apostle’s doctrine, in 11 to the Romans.

### Paragraph 536

This doctrine will make men neglect their own salvation; for already there are those who say, “If the end is well, then all is well.” As though they should say, “Live however you please in this world, immersed in pleasures and bloodshed, and devoted to gluttony, but believe only at the very end of your life, and you shall be saved.” I am not ignorant that there are many unclean pigs and filthy swine, abusing the true doctrine and comfort of the Gospel; but will the abuse of ungodly men take away the truth from us? The children of God, who know that there is no other propitiation or satisfaction for sins but the offering of Christ, therefore do not cease to renew their lives daily by repentance.

### Paragraph 537

Thus, although the faithful doubt nothing but that countless numbers are converted and saved by the Lord at the end of their lives, they do not abuse God’s mercy to license the flesh, but are afraid. For there are other passages that restrain them in order and duty. For the Lord says, “You have been forgiven much, go and sin no more, lest something worse befall you.” Also, “Let us do good while we have time; the time will come when we cannot work.” The parable of the ten virgins declares the same to us. Moreover, if the righteous scarcely be saved, where will the sinner and wicked appear? Furthermore, do not tempt the Lord your God. And countless others of like sort. And when the saints have spent their entire lives behaving blamelessly in God’s righteousness, yet in the last time of their lives they do not trust in that, but in God’s mere mercy through Christ. They always remember how gravely he was rebuked in the Gospel, namely, he who envied the good fortune of him who labored with him in the vineyard, for that he had received so much wages, coming into the vineyard about the last hour of the day, as he had received that who had labored all day long, and again, the thrifty son, for that he was sorry that his wasteful and prodigal brother was received again by his father, and a feast also made for him, and for him who was always obedient and took pains continually, no such thing was prepared.

### Paragraph 538

[note 35.18]But the Gentiles he does not arrange into any certain number, but says how he saw a great multitude, which no man could tell; no more than they could the stars, sand, herbs, or grass how many they were in number. He signifies therefore, that in all the world, at all times, innumerable people are saved by Christ. Nevertheless, lest any man should think that it would avail or hinder him to salvation to be born of this or that nation, tribe, or tongue, John adds immediately, “of all tribes, peoples, and tongues to be ordained to salvation indifferently.” Therefore, this difference hinders salvation nothing; but people are found in India, Ethiopia, Barbaria, and in the furthest part of Libya, in Scythia, Tartaria, and in the uttermost ends of the world, who are saved by the grace of Christ.

### Paragraph 539

And because it is difficult to reason about things to come, John here speaks most expressly of those who are not to be saved, but have already achieved salvation and are in Heaven, so that we should not doubt of their salvation. He also depicts the manner of salvation and everlasting blessedness. This treatise refutes those who suppose that souls sleep, nor have the enjoyment of Godhead before the judgment, nor are yet in Heaven. First, he says how they stand before the throne and in the sight of the Lamb. For the first blessedness in the blessed life is to see God as he is and to enjoy his glory, to be with Christ in glory. John 17:1. John 3. White stoles are the garments of triumphant and clean people. As will be declared more at large hereafter, and has been noted once or twice before. It signifies that the blessed souls are adorned with light, and so forth. And the palm also is a token of victory. Pliny treats much of the palm in chapter 4 of book 13. All writers say that the palm was the most ancient badge of a conqueror. And why this tree was chiefly chosen for this use, Aulus Gellius shows the cause in chapter 6 of book 3 of Noct. Att. Writing that in a palm tree there is a certain peculiar thing which agrees with the nature of stout and noble people. For if you lay great weights upon the wood thereof, the palm gives not place downward, but rises up against the weight and bears upward. And for this he cites the authority of Aristotle and Plutarch, to whom you may add also Pliny, book 16, chapter 24.

### Paragraph 540

All of these things are accompanied by an exceedingly great rejoicing, whereby not only do they give God thanks and praise his mercy, but also openly demonstrate and testify whom they may thank for their salvation. And they say, “Salvation to him,” and so forth. For that is the more accurate rendering, as Erasmus has also noted. They signify that God is blessed not only in himself, but also for having communicated this salvation to them and saved them. Of the throne or seat of God, we spoke before in chapter 4. God the Father himself sits on the throne. It is therefore a phrase of speech, which has this meaning: we owe this our salvation and blessedness to our God, who sits in his throne.

### Paragraph 541

Again they communicate this salvation to the Lamb, that is to Christ. For God by his grace through Christ saves the believers. And since Christ is called the Lamb, the whole mystery of the incarnation and redemption is remembered in the word Lamb; that being indeed reconciled to God by the blood of the host [note 35.22], we are now the heirs of God and the sons of God. &c. Therefore, the saints in heaven, and our fathers already saved and dwelling in heaven, do testify, and in testifying teach, that they are justified and saved not by Islam, or Roman Catholicism, or any other observances, but by the mere grace of God in Christ.

### Paragraph 542

Hereby are refuted two opinions, deeply harmful to the whole world. The first claims that Papists are saved because of their simplicity and strict discipline. For they say that because Papists are ignorant of better things, and perform their works with good intentions, they are saved by those same works. This is most vain and most ungodly. They add that unless we judge this way, surely not one Papist would be saved. I say plainly, no one is saved by Papistry, no more than by Mahometry. For it is called the way of perdition even by Saint Peter himself. 

However, I do not think that no one among the Papists is saved. But I believe that countless numbers, as I said before, have at length seen the corruption of Papistry through God’s illumination, have forsaken Papistry, and embraced the sincere Gospel, and so are saved by Christ alone.

### Paragraph 543

The latter assumes that every person, regardless of their religion, will be saved. Against this, the faithful here proclaim: those who are saved are saved by the grace of God through Christ. Therefore, no other religion offers salvation. There is no other name given to people, in which they must be saved, but that of Christ Jesus. No other way is open into heaven, nor any other door: whoever claims otherwise is called by the truth a thief and a murderer. Indeed, they utterly abolish Christ and the entire scripture, those who contend that every person is saved by their own religion. I cannot imagine any doctrine so harmful. Therefore, let us maintain what all the faithful in heaven have taught us: that salvation is from God through Christ.

### Paragraph 544

All the angels in heaven confirm these things, lest anything should lack that belongs to a sure and certain testimony; and also teach us by their example what we should do. They sing together “Amen,” thereby also testifying that salvation is by grace alone through Christ. Again, they fall down and worship God. But how much more ought we, mortal men, by worshipping to attribute this honor to him? And by singing a hymn, they exhibit to us a form of serving God, finally of judging rightly of God, that we attribute nothing to any creature to the reproach of the creator, which belongs to God alone, but ascribe all things to God wholly. The words of this hymn are explained in chapters four and five, so I need not linger on them here. They use “blessing” for “praise”; the rest of the words are plain.

### Paragraph 545

And now let us learn, being taught by so many testimonies and examples of all the saints, forsaking all vain and wicked opinions, to give all glory to God through Christ: to whom be praise and thanks giving. Amen.

## Sermon 36: ¶ Here is expounded who they be that are are clothed in white, from whence is salvation, and what is the true blessedness.

### Paragraph 546

.

### Paragraph 547

And one of the elders answered, saying to me, “What are these who are arrayed in long white garments, and from where do they come?” And I said to him, “Lord, you know.” And he said to me, “These are those who have come out of great tribulation, and have made their garments clean, and made them white in the blood of the Lamb. Therefore they are in the presence of the throne of God, and serve him day and night in his temple, and he who sits on the throne will dwell among them. They shall hunger no more, neither thirst, neither shall the sun light on them, neither any heat. For the Lamb who is in the midst of the throne shall feed them, and shall lead them to fountains of living water. And God shall wipe away all tears from their eyes.”

### Paragraph 548

[note 36.1]John saw the souls of martyrs resting under the altar, clothed in white garments. He also saw an immense multitude from every nation and people, saved from the idolatry of the Gentiles and the superstition of Antichrist, also clothed in white garments. On this occasion, three things shall now be explained to us: what are those who are clothed in white garments? Whence do they have that whiteness, purity, and salvation? Finally, what is the state or felicity of these, or what is true blessedness?

### Paragraph 549

[note 36.2]What time John had seen them, he marveled without doubt as to what they were; nevertheless, he seems to have inquired nothing. But of his own accord, one of the twenty-four elders offers himself as an expositor to him, doubtless an excellent teacher, a patriarch and prophet, finally a celestial master, to whom we may justly give credit. Here appears the ignorance of human wit. For just as the eunuch of Ethiopia in chapter 8 of the Acts of the Apostles acknowledges his ignorance, except an interpreter and skilled teacher were given to him, so at this present also blessed John himself, when asked whether he knew those who were clothed in white, confesses his lack of knowledge; yet he ascribes the knowledge to his teacher, by this means requiring a further declaration through a most humble modesty. Finally, here appears the immeasurable goodness of God, which undertakes to teach us who are rude and unworthy. We have many examples of this everywhere in the prophets and in the holy gospel of Christ.

### Paragraph 550

And in the beginning he declares plainly to John, and to all the faithful in the world,[note 36.3] what these are who are clothed in white: and he explains fully from whence they came. For with one and the same answer he dispatches both. He says briefly, that the clothed with white in Heaven are the godly people of all times and ages, who at length have escaped out of great tribulation. Tribulation is found to be sundry and diverse. For first, it is tribulation which comes from the plotting and persecution of tyrants. This pertains to martyrs alone. Whereof we have spoken in chapter 6. Which, for as much as in this world they were overwhelmed with unspeakable reproaches for the word of God, they have in another world received white garments. Then there is another tribulation, which arises from the fear of God, and is a care of obtaining salvation. This is sorrowful because of the unrighteousness and corruption of man. It is sorrowful because of the grievous abominations of Antichrist. And these also, albeit they be not made martyrs, yet are they in another life clothed with white. Finally, they have tribulation and are molested in the flesh, so many as mortify their flesh with its lusts. And because they mourn here in the world, they shall receive comfort and consolation in the world to come.

### Paragraph 551

Again, lest any man should ascribe life and salvation to martyrdom as to our work, and to repentance as to our desert, the Lord moreover declares expressly from whence that life and salvation proceeds, and how that whiteness and purity comes to us. And they have washed their garments, he says, and made them white by the blood of the Lamb. And here is found a diverse reading. Some suggest that they have enlarged their garments, so that it might appear he alluded to the families of great princes, who use, for the setting forth of their renown, to put on most wide and most sumptuous garments. But in my opinion, the Complutensian copy and Arethas seem to read more rightly and more simply, namely, they have washed, as also the old interpreter has translated it. For by this exposition follows, and have made white. Primasius reads, and have made their garments white in the blood of the Lamb. And hereby is signified that the salvation and purification of the faithful is of the blood of Christ, and of no other thing. Where verily blood sprinkled whitens not, but pollutes. Therefore we must understand these things spiritually, to wit that the very, natural and human blood of Christ, shed upon the cross, being sprinkled upon us spiritually (as Paul to the Hebrews 10 explains), and received with faith, although it touches us not naturally and corporally, purges us from all sins. And therefore we read in another place that Christ purges us with his blood. For because sanctification is the only work of God, therefore where the saints are said now to have washed and whited their garments by the blood of the Lamb, it signifies that they have received by faith the purification prepared by the blood. And this doctrine is catholic and of the right faith, which has so many and so great testimonies in the holy Scriptures. Finally, we perceive how those who are saved from the kingdom of Antichrist are saved by the merit of Christ alone, and by no other thing, as I have also shown you before. Moreover, it follows immediately: therefore they are in the sight of God’s seat. For what cause, I pray thee? Because they have washed and whited their garments in the blood of the Lamb: therefore for the merit of Christ have they entered into heaven, and there are surrounded with eternal light.

### Paragraph 552

Finally, the elder proclaims at length the condition of the saints and the true blessedness of the faithful. These things are certain glimpses, provided here for consolation. Or else, such things as the eye has not seen, or the ear heard, that God has prepared for those who love him. He recounts many things, from which we are to understand the excellence of eternal salvation and what good things we attain in it.

### Paragraph 553

First, the saints stand before the throne of God. In the throne is the majesty of God to be worshipped forever, and the blessed Trinity. And the saints stand before the seat, not as those who tarry before the gates. For as God’s most complete friends, they are always in God’s sight and have the enjoyment of his deity. Concerning this, the Lord speaks in the Gospel: “Pray that you may escape these things and stand before the Son of Man.” And also David: “The fulfillment of joys is in your sight, and pleasure at your right hand forever.” And there is added another thing, which may explain that standing: they serve God in his temple both day and night. That service has pleasure and no pain. And they serve God in the temple, as God is wont to be served in the temple. For they keep holy days, are glad, rejoice, are merry, praise, and so they offer up sacrifices and are refreshed with heavenly repast. And this joy shall be everlasting and perpetual, which is signified by day and night. Otherwise, in the blessedness everlasting there is no might at all, nor any changeable course of time. Hereunto is added that he who sits in the seat, that is, the divine majesty, [?will] dwell in them: that is to say, God will be all in all, or he will lean over them, and as it were a tent or tabernacle, will overshadow them, defend and keep them, and give himself whole to be enjoyed by them, as most familiar and friendly to them. Moreover, they shall hunger no more, neither shall they thirst. For all infirmity and misery is taken away from the blessed souls and glorified bodies. They are filled with all good things without any loathsomeness, with a most joyous fulfillment. Now the sun falls not upon them, nor the heat: which phrase of speech betokens that they are put to no labor nor pain, but are delivered at once from all displeasure and all pain, and to be at most pleasant rest.

### Paragraph 554

Again, this is established for so great a blessedness, Christ the Lamb—that is to say, Christ the mediator and redeemer—in the midst of the throne, that is, very God. For he, as both Ezekiel 34:23 and the Lord himself in John 10:11 testify, will feed them like a shepherd, and as a leader of life will lead them to the fountains of living water: that is to say, he will quicken them forever and preserve all his in that blessedness. He uses in this treatise words of the prophets most commonly used, so that, climbing to higher things, we might in some measure esteem heavenly gifts. Hereunto he joins yet a notable blessing: and the Lord will wipe away all tears from their eyes. Which words he has borrowed from Isaiah. For saints in this world, troubled by various evils, have shed most plentiful tears; but in the world to come, the Lord will comfort them, gladdening them with everlasting joy, never giving them any occasion of grief. And therefore he said in the Gospel, “Verily I say unto you, you shall weep and lament, but again the world shall rejoice; you shall mourn, but your mourning shall be turned into joy. And your heart shall rejoice, and your joy no one shall take from you.” We shall hear similar things in Revelation 21 as well.

### Paragraph 555

Hereby they perceive how shamefully they sin, those who constantly say, “If I should endure this present life for the sake of religion, who will tell me what that other life to come will be like?” Perhaps if I neglect this, I will gain nothing in another world. For here we have a most manifest testimony that, as most assured salvation is prepared by God in heaven for the faithful, so it is also most ample and great; in so much that the Apostle in another place says that the afflictions of this present time are not equal to the glory which shall be revealed to us. The Lord grant us that we may acknowledge these things.

## Sermon 37: ¶ Whilst the vii. Seal is opened, and the Angels with trumpets come forth, Christ the intercessor of his church offers up before his father the prayers of his faithful.

### Paragraph 556

.

### Paragraph 557

And when he had opened the seventh seal, there was silence in Heaven about the space of half an hour. And I saw seven Angels standing before God, and to them were given seven trumpets. And another Angel came and stood before the altar, having a golden censer, and much incense was given unto him, that he should offer the prayers of all saints upon the golden altar which was before the throne. And the smoke of the incense, which came from the prayers of all saints, ascended up before God from the hand of the Angel. And the Angel took the censer and filled it with fire from the altar, and cast it into the earth, and there were voices and thunderings and lightnings and an earthquake.

### Paragraph 558

Surely there are no books in the world, written by whomever or whenever, which can compare with the books of holy Scripture, concerning sincere truth, pure simplicity, and plain order. And perhaps this is no marvel to anyone who knows that these writings were indeed written by men, but inspired by the Holy Spirit. There are buildings most skillfully constructed by men, framed and arranged in a most beautiful order. But what beauty will you judge them to have if you compare them with the creation of the world, and with that most beautiful order which we see daily in all created things, and the changeable course of times? The most excellent works of men have nothing in them, or seem vile, if you compare them with the workmanship of God the creator.

But for the most brilliant order and most plain treatise, this book of the Apocalypse has among others a most notable, excellent, and wonderful praise. John promised some of the matter, signifying that he would speak of those things which should be done in the church from his time until the judgment. And the faithful know to what end they should take these things, not so that their curiosity might be maintained or satisfied, but that, being sufficiently warned beforehand, they should not fall, but take heed to themselves and hold fast to true salvation.

And inasmuch as there is much talk among men about why God does this or permits that, and why he does not prohibit these or those things, John has exhibited to us a most wholesome vision, by which we may learn not to talk against God, nor to contend with him, but to acknowledge all his judgments to be righteous and just. This thing, indeed, both all the saints in heaven and also angelic spirits acknowledge, and attribute all glory to God.

And thus, having prepared the minds of the listeners, he comes to the matter itself and declares the fatal destinies of the church. Under the sixth seal he touches generally the corruption of doctrine, which sin is more perilous and more pestilent than all dangers of the body or outward perils. He reasons more fully thereof, and now particularly under the opening of the seventh seal recites how far the same stretches. For he declares how many, how great, and what manner of sects, heresies, and troubles shall arise in the church, and how hurtful they shall be to the church.

And this place contains a history of the corrupt doctrine, of heresies or sects, and troubles from the time of John until the last judgment; it is extended throughout the eighth, ninth, tenth, eleventh chapters.

### Paragraph 559

Nevertheless, before the trumpets sound, for a consolation, as it were by a brief digression, a remedy is placed there, which the faithful in all ages may use in that pestiferous corruption to keep safe their souls and the integrity of the same. For many times in this book are brought in most strong consolations in matters of greatest difficulty. The tenth chapter will also serve this argument. And the remedy that he shows is this: that we must flee unto Christ, redeemer of mankind, intercessor and propitiator. And that we shall be safe under his defense, we must offer up to him our prayers continually. Verily, the Lord in the Gospel, reasoning of the greatest dangers of the Devil, and being at hand, yet adds by and by that which might comfort their sorrowful minds: “I have prayed for thee, Peter, that thy faith should not fail, &c.” Behold, we are saved in greatest distress through Christ’s protection, that we should not faint in faith. However, as the Evangelical and Apostolic letters indicate everywhere, our continual prayers, which we offer to God through Christ, must be joined to our trust in Christ. And in few words, the intercession of Christ at the right hand of God, and the effect and manner of the prayer of the faithful are here set forth to behold.

### Paragraph 560

But we shall declare everything in order. He spoke generally under the sixth seal of corrupt doctrine; in the seventh, he will declare the same particularly and most abundantly. And while the seventh seal was opened, there was silence in heaven almost half an hour. Concerning this silence, commentators write diversely. But as I think, this silence excites the hearers to a diligent and attentive hearing. For silence has an air of admiration and expectation of weighty matters. Solomon says in the ninth chapter of Ecclesiastes, “the words of wise men are heard in silence.” When weighty matters should be proclaimed and set forth, the herald is wont to proclaim silence. Indeed, the matters that follow are of great importance, and unless we observe them with great attentiveness, we shall perish in sects and seductions. Those spiritual wickednesses are more dangerous than corporal perils.

### Paragraph 561

And now, while they remain silent, eagerly awaiting what will come, the last seal is opened. Behold, seven angels with trumpets appear, of whom we shall speak later.

### Paragraph 562

Now a remedy is presented for these great evils, as I mentioned. To make it more vivid and impress it more deeply upon our hearts, it is shown through a divine vision. Before the throne, and almost encompassing it, appears a golden altar. And an angel comes and stands at this altar; he holds a golden censer into which the saints place their offerings. He offers them before the throne, and the fragrance of the incense ascends from the angel’s hand before God.

### Paragraph 563

We said in another place that the golden altar of incense was the Lord Christ himself, who is both altar and sacrifice and priest, as Paul witnesses to the Hebrews. He is called an angel, namely, the one of whom Isaiah makes mention in chapter 9, and also Malachi, saying, “Behold, I send my angel, who shall prepare the way before me, and suddenly the Lord shall come unto his temple, whom you seek for, and the angel of the covenant, whom you desire: behold, he comes,” says the Lord of hosts. The former angel, that is to say, messenger or ambassador, was John the Baptist, who prepared the way for the Lord. He, namely, the later angel, came immediately after the preaching of John and made complete that everlasting covenant. He now appears on the right hand of God in Heaven.

### Paragraph 564

And two things are said about him. First, that he stood before, or in, or upon the altar. We must not imagine anything physical here, but understand that this manner of speaking signifies the priesthood of Christ. He always appears before his Father on our behalf, as Paul taught in Romans 8 and Hebrews 9. Therefore, he pleads the cause of his church before God and is advocate for the faithful. The same one who stands before the altar also stands in the midst of the throne. For he is coequal with the Father in deity, according to which he sits on the throne; and in his humanity, he is of the same substance as we are, according to this dispensation, as both Bishop and true man, to stand before the altar. The latter, which is to be noted, is that Christ holds a golden censer in his hand. For he has taken our very nature without sin, so that he might intercede for us and offer our prayers to God the Father.

### Paragraph 565

And lest any man should doubt that he receives our prayers and offers them to God, finally that the true office of the Church might also appear, offering up all things by Christ, there is added, “to him are given many fragrances.” But to what end? That he might give them upon the golden altar, and that before the seat, as though to say, that he might bring them into the sight of God.

### Paragraph 566

And because of a further declaration, lest we should not know the true fragrances, which please God, and which the faithful offer unto God through Christ, one or twice he adds that those fragrances are the prayers of Saints. And he means by Saints, not those that dwell in heaven, but us in the earth, which are sanctified with the spirit of our God, with the blood of Christ, baptism, faith, and word. John 13. And the prayers are invocations and expressions of thanks. And he says expressly of all Saints, lest any should fear that he and his prayers offered by Christ were excluded: If thou believe, thou art holy, and thy prayer is of God accepted. What the prayers of Saints are, it appears in the Lord’s Prayer, which we offer up to the Father in the name and words of Christ: hallowed be thy name, thy kingdom come, and the rest, which all contend against those sects and corruptions of true doctrine.

### Paragraph 567

Irenaeus refers to this passage in chapters 31 and 32 of the fourth book. By this, he calls the Eucharist—the giving of thanks—the sacrifice of Christians. Those who uphold papistry corrupt this passage, twisting it to mean that the priest should sacrifice the actual body of Christ for the living and the dead. But the holy Bishop of Lyons knew this foul error. Away with them and their sophistry, to where they deserve to be. I have also spoken before, somewhat about the same matter.

### Paragraph 568

And that it might clearly appear unto all men, that the prayers of the faithful, offered to God through Christ, are pleasant and acceptable, there is added: and the smoke of the odours ascends, that is to say, the prayers of the faithful were accepted by God. Therefore let us offer diligently our prayers unto God through Christ. For he hears us, and delivers us from evil. And the scripture many times calls our prayers an acceptable sacrifice to God. The passages are in Hosea, 14, in the 50th Psalm. And in many other places. In Psalm 141, the prophet says, “Let my prayer be directed as incense in thy sight, the lifting up of my hands an evening sacrifice.” Primasius, expounding this place, said how Christ is said to have taken of the prayers of the saints. For because through him the prayers of all may come sweetly unto God. Hereof the Apostle: by him we offer up always a sacrifice of praise unto God, that is to say, the fruit of lips confessing his name.

### Paragraph 569

Hereby is refuted the opinion of those who suppose that the saints in heaven are intercessors of the faithful, who should commend their prayers to God and make the way open to God. For what need have they to procure for themselves other intercessors or advocates? What lack do they find in Christ? Or whom may they prefer or compare with Christ? What shall we say that even now the offerings are offered by the hand of the angel? The celestial saints were present with the Lord and were seen about the throne, but which of them taking the censer and gathering the prayers of the faithful, offered them to God? It displeased King Azariah [? Ozias or Asarias] that he took in hand the censer, intending to sacrifice and to perform the priests’ office; the same would be worse for those who dwell on earth; indeed, they would not remain in heaven if they took upon themselves the office of the only Bishop. &c.

### Paragraph 570

After this, we have heard that Christ filled the censer with fire taken from the altar, and sent it down into the earth. By this narration, he returns to finish the exposition of the trumpets. This fire is the grace of the Holy Spirit, which is put into the censer, taken from the altar, and sent down into the earth. For Christ took the fullness of the Spirit, as John shows in the 1st and 3rd chapters. Christ is altar and censer. Of the altar, fire is taken. For the Holy Spirit is the spirit of the Father and of the Son. “Him,” he says, “I will send you from my Father.” He sent him into the earth under the shape of fiery tongues; he sends him also at this day into the hearts of the faithful, that he may inflame them. This is the same fire which the Lord in the Gospel of Luke says he will send into the earth, and would that it should burn.

### Paragraph 571

Moreover, the effect of this fire follows immediately. For there were thunderings, and voices, and lightning, and earthquakes. By the voices of the Gospel, the wounds of sinners are healed, and the hearts of men enlightened by the illumination of the Holy Spirit, and so forth. Of these things we have also spoken in chapter 4, sermon 24. Concerning the preaching of the Gospel, as Haggai also prophesied it would come to pass, it sows a wonderful commotion of all nations, and so forth. Satan also was stirred, who raised up his ministers throughout the world against the wholesome preaching of the Gospel. For sects sprang up, whom the defenders of the truth resisted, fighting with them. Whereof now he will reason at length. The Lord grant grace that these things may both be spoken and heard with much fruit.

## Sermon 38: ¶ Of the seven Angles trompetters, and of the trumpets: and of the first ii. and iii. trompet.

### Paragraph 572

.

### Paragraph 573

And the seven angels who had the seven trumpets prepared themselves to blow. The first angel blew: and there was hail and fire mingled with blood, which were cast into the earth; and a third of the trees were burned, and all the green grass withered. And the second angel blew, and as it were a great mountain burning with fire was cast into the sea, and a third of the sea turned to blood; and a third of the creatures which lived in the sea died, and a third of the ships were destroyed. And the third angel blew, and a great star fell from heaven, burning as a torch, and it fell into a third of the rivers, and into the fountains of waters; and the name of the star is called Wormwood, and a third of the waters were turned to Wormwood. And many men died from the waters because they were made bitter.

### Paragraph 574

Our Lord Jesus Christ has kindled on earth a bright and wholesome fire, which the Apostles and those following in their footsteps have continually inflamed. But contrarily, Satan seeks to quench this wholesome fire, and not only to corrupt and deprave this doctrine of salvation, but also to abolish it and overwhelm it with lies. The manner of this is presently described, and even vividly portrayed, for no other end than that the faithful, being warned and fully taught, might be well aware of that pestilent infection. For the scope or end of this book is to preserve the church safe and sound from corruptions, or at least to repair it if it has been corrupted.

### Paragraph 575

Therefore, John saw seven Angels standing in the sight of God. To stand signifies to minister, and encompasses the faith and diligence of ministers. Servants stand before kings, ready to do service and to execute their commandments. We read in the first chapter of Job: “The sons of God came and stood before the Lord, and Satan also appeared among them.” The blessed Angels are called the children or sons of God. They come to do service before God; Satan appears among them, for he too is a minister of God, for the execution of those things that pertain to the wrath and indignation of God against the wicked. All elements are God’s ministers, and finally, all the creatures of God. For he is the Lord of the Sabbath, the God of hosts, who for the salvation and judgment of men uses all his creatures well and rightly, each according to its nature and disposition. For he uses the ministry of Angels like Angels, and so the service of devils as devils. But since the number seven is the number of fullness, containing within itself all times—for there are seven days of creation and rest, there are seven worlds or ages—certainly seven Angels appear before God, for they signify all battles that shall be fought to the end of the world.

### Paragraph 576

For to these seven angels are given seven trumpets, and the angels already had the trumpets, and even prepared themselves to blow the onset.[note 38.2] Where chiefly the use of trumpets is to be searched for. The same is most plentifully described of Moses, in the tenth chapter of Numbers. The use of trumpets was diverse, as it is also at this day. First, by the sound of the trumpet, the people of Israel were called together to consult about the common good. Again, at the sound of the trumpet, the senate of princes of the people did assemble. Moreover, they were warned by the trumpet when and who should remove their tents. Furthermore, the trumpets blew to battle, when they joined to fight, as may be seen in the twentieth chapter of Deuteronomy. The people moreover were called together with trumpets on the holy days for public and divine service. "Sound with the trumpet in Zion, call the congregation, gather the people," says Joel. There was moreover a feast of trumpets, and a Jubilee, having the name of the blowing and sound of trumpets, as appears in the twenty-fifth chapter of Leviticus. Finally, the preaching of the truth was figured by the sound of trumpets, and neither might any other blow it but priests. For it behooves one greatly to whom you commit or deliver the signs public.

### Paragraph 577

Of this varied use of trumpets, none shall agree better with our matter than the warlike. For this world has a shape of war. In it are the camps of good men, and the camps of evil: the tents of Catholics, and the tents of heretics. The chief ruler of these is Satan, and of those Christ: the Captain and Emperor of these is the Devil, of the other the Son of God. And now the angels sound their trumpets and blow the onset: not that the good angels and God himself is the author of heresies and of heretics, whose origin is referred to Satan and sin; but sounding their trumpets they give indeed warning to all men, and signify that most grievous wars shall arise in the world, and even in the church itself. But diverse men are diversely moved and work in war according to their natures. The true Catholics, being warned by the trumpet, take heed to themselves, pray, and finally taking in hand spiritual weapons, prepare themselves unto battle and manfully fight for Christ, and for maintaining and defending the truth. Heretics, sectaries, and men of corrupt minds, according to their malice, taking to them also armor, run forth and fight against Christ and the truth, defend lies, and those who are weak they take, spoil, beat down, and destroy. The good shepherds are the trumpets of God and of Christ: the Devil blows up arch-heretics and beginners of sects.

### Paragraph 578

Regarding the righteous and their struggle, we will hear in chapter 11 and the following chapter. Nevertheless, in every conflict, we must understand that the saints do not sleep, nor are they anywhere idle, but perform their duty everywhere. It was indeed sufficient for the Lord to show us the heretics and sectaries scheming, and to declare how much harm they can inflict, so that we might watch more diligently and beware of all corruption.

### Paragraph 579

The first angel sounding the first trumpet announces the first conflict. All battles share some likeness and difference. They are alike in that all heresies attack Christ and seek to extinguish or distort the truth of the Gospel. They are diverse in that, at different times, Satan, assailing different doctrines, has spread various heresies throughout the Church. Therefore, as the angel sounds the trumpet, that is to say proclaims war, he warns the saints to be watchful. And as he continues to blow, through God’s permission, according to his just judgment, and by the means and suggestions of Satan, hail and fire mixed with blood were created and sent, or fell, upon the earth. For Paul acknowledges spiritual powers in the heavenly spirits. And Scripture in a certain place aptly depicts wholesome doctrine as a heavy dew and shower that makes the earth fruitful; therefore John rightly compares false and heretical doctrine to hail, for it destroys the fertile places of the earth and utterly ruins the abundant produce. Wherefore, just as elsewhere perverse doctrine is called darnel, leaven, chaff, and so on, it is here called hail. But this hail is tempered and of a wondrous mixture. For it combines fire and blood. These things must be understood allegorically, not literally. Hail is water frozen by cold. And water has been called the wisdom of Scripture; therefore, hail signifies false wisdom. Yet fire is added to it, representing the pretense of Scripture and the inspiration of the Holy Spirit, to which is added blood, the evil affections of humankind—namely, the vices of ambition, wrath, contention, hatred, and similar affections. A hurtful and pestilent doctrine is compounded of these. For when false doctrine governs or corrupts Scripture, and the wicked affections of teachers are joined with it, a pestilent doctrine arises. Such was the doctrine from the beginning of the Nazarenes, or Mimeorites, and of the Hebionites, contending that justification did not come by the alone faith of Christ, but by the law. Our people fought sharply against this pernicious doctrine, namely Paul and the other apostles. From the beginning, some have corrupted the faith with philosophy, others have been blinded by human traditions, and have produced most corrupt opinions. Histories bear witness to this. And Tertullian not without cause called philosophers the patriarchs of heretics. For Paul diligently warned the godly to be wary of philosophy. Those who have not kept themselves from it, and who have valued philosophy and I know not what traditions more highly, have cast great heavy hailstones into the church instead of the heavenly dew and sweet showers.

### Paragraph 580

And have indeed harmed the church greatly. For a third of the trees was burned, and also all the green grass. This number is hinted at in four trumpets, and in six as well. And it seems to signify that a great portion of unstable and wavering people are deceived and abandoned, giving themselves to be destroyed by wicked men; again, the best part of the faithful are to be saved. The Lord himself knows the number exactly. It is enough for us to know these things which he has revealed to us, neither to search curiously any further.

### Paragraph 581

That men are signified by trees appears in the ninth chapter, where it is said, “They had been commanded not to harm the grass of the earth, nor any tree, except men.” After he had said, “Save only those trees which were not marked,” he preferred to say “men,” as if with this key he might unlock the mystery. It is not uncommon in Scripture to represent men by trees, flowers, and grass, as we may gather from Psalm 1, Isaiah 40, and Matthew 12. But that latter point, that all green grass was burned, must be favorably explained. For who could believe that all men were destroyed by those first heresies? We understand therefore that the minds of the faithful were diversely afflicted and tormented by those errors and troubles, but yet, like gold tried in the fire, not to be utterly consumed.

### Paragraph 582

The second angel sounds the trumpet, signifying that new wars are now brewing; and therefore exhorts that all the godly defend themselves with weapons. And there is cast into the sea not a mountain, but as it were a mountain burning with fire. The sea bears a figure of the world, than which there is nothing more unstable. It is a thing most frequently used by the prophets to call this our world, wherein we live, a sea. By mountains are signified kingdoms, as witnessed by Isaiah in chapter 2, Daniel in chapter 2, and Zechariah in chapter 4, and Christ himself in Matthew 7. By removing of hills or mountains, signifies any difficult thing, and by the opinion of many, an impossible thing. Now therefore springs up an heresy and a doctrine in the church, as it were a burning mountain, which was in deed most furnished, and as it seemed invincible. We read that such was the heresy of the Valentinians, whose sect the holy martyr Irenaeus taught to be divided into many. Such was the fury of the Manicheis and Montanists. They seemed to many to burn with the spirit of God, and to be wholly nothing else but the spirit, and all their oracles to be of the holy ghost. Manichaeus called himself the Apostle of Jesus Christ. The Montanists boasted of a new holy ghost. There was most great plenty of this darnel throughout the universal church. Neither was the success thereof small. For the third part of the sea was made blood. The Apostle signifies the wickedness of sects. For how vile and impudent were the heretics called Gnostics, the Valentinians, and Manicheis, as testified by Irenaeus, Augustine, and Epiphanius. And a great part of the creatures in the sea perished. And he speaks of those who have souls, not of fishes in deed, but men. Many ships moreover were lost, to wit, mariners and land men, being corrupted with these heresies.

### Paragraph 583

Those heresies arose from the writings of the authors I mentioned, but they are not yet completely gone; corrupt men persistently associate with one another, reviving the old errors. Consequently, a bitter conflict remains in the church to this day, and we are continually warned to be on guard against these corruptions.

### Paragraph 584

The third angel blows his trumpet, proclaiming new wars; and behold, a great star fell from heaven, burning like a torch, and infected the third part of the rivers and fountains of waters. That star is called wormwood. I told you in the first chapter that stars are called preachers, bishops, and notable men in the church. Therefore, it signifies that some notable man should fall away from the true faith into heresy, with which he should infect a great part of the world, corrupting the Scriptures and the sound doctrine of faith. And these things seem to be fulfilled in Paulus Samosatenus and Arius. This torch burned horribly and inflamed the whole world without recovery. That pestilence denied the deity of Christ and made the whole Gospel most bitter to us. For if Christ is not very God, how is he a Savior, King, bishop, intercessor, mediator, and salvation of the faithful? He quenched the light that denied the deity of Christ. Therefore, he is called by the name of wormwood. The prophet Jerome used the selfsame allegory, or metaphor, or allusion, in the ninth and twenty-third chapters. And Amos in the sixth chapter, where he says that the judges have turned judgment into wormwood.

### Paragraph 585

The Scripture and doctrine, beautifully depicted by rivers and fountains, was the occasion of death for many, having been corrupted by the Arians. The Scripture and doctrine of the Gospel is of itself mortal to no one, but rather lively to all; yet corruption makes it deadly. Poison put in wine makes the wine deadly: the wine itself kills no one, but rather glads and rejoices all people. Read the ecclesiastical histories of Eusebius, Theodoret, Sozomenus, Socrates, and others, and you shall perceive how aptly John has written all these things, and how rightly they are all fulfilled. No small part of that bitterness has flowed unto our time, while that old error is often renewed by the instigation of the devil. For what that unclean beast, Michael Servetus, a Spaniard, vomited against the Son of God, for his impenitent wickedness and continual blasphemy, the world knows; he was burned at Zurich. We must therefore pray to the Lord that in such dangerous conflicts he would keep us safe and sound. Amen.

## Sermon 39: ¶ The fourth and fifte trompet is expounded, of the opening of the botomleſſe pit, and of grass hoppers creping out into the Earth.

### Paragraph 586

.

### Paragraph 587

And the fourth Angel blew, and a third of the sun was struck, and a third of the moon, and a third of the stars, so that a third of them was darkened. And a third of the day was struck, so that it would not shine, and likewise the night. And I beheld, and heard an Angel flying through the midst of Heaven, saying with a loud voice: Woe, woe, woe, to the inhabitants of the Earth, because of the voices yet to come of the trumpet of the three Angels, which were still to blow.

### Paragraph 588

And the fifth angel blew, and I saw a star fall from Heaven to the Earth. And to him was given the key of the bottomless pit. And he opened the bottomless pit, and there arose a smoke from the pit, as it were the smoke of a great furnace. And the sun and the air were darkened by reason of the smoke of the pit. And there came out of the smoke locusts upon the earth, and to them was given power as the scorpions of the Earth have power. And it was said to them that they should not hurt the grass of the Earth, neither any green thing, neither any tree, but only those men which have not the seal in their foreheads, and to them was commanded that they should not kill them, but that they should be vexed five months, and their pain was as the pain that comes of a scorpion when he has stung a man. And in those days shall men seek death, and shall not find it, and shall desire to die, and death shall flee from them.

### Paragraph 589

The fourth trumpet declares a harmful and prolonged conflict,[note 39.2] which arose in the church concerning the doctrine of Pelagius. This Pelagius taught that the sin of Adam only harmed him, and not mankind, and therefore that in him all men do not die. He asserted that man has free will, so that he may do good. Neither did he believe that he would be free if he needed the help of God; if he had it, he could more easily do good; if he did not have it, he could nevertheless work it by his own virtue and deserve everlasting life. Therefore, he maintained that our victory is not from the help of God, but from free will, and that remission is not given to the penitent after the grace and mercy of God, but after the merit and work of those who, through repentance, are worthy of God’s mercy, as the residue which Saint Augustine recounts in the one hundred and sixteenth Epistle to Boniface, that Pelagius had renounced; which nevertheless in another place he shows that he had taught it and returned to his vomit, as in the register of heresy, heresy number 88, in the second book, chapter 2, to Boniface. The Manicheans, he says, deny that a good man had the beginning of evil from free will. The Pelagians also say that an evil man has free will sufficiently to fulfill a good precept. The Catholic doctrine reproves both of these, and says to them, “God made man upright,” etc. And to them, “If the Son has made you free, you are truly free.” And in the ninth chapter, the same author says that the will of man to evil is free, but that to do good, it must be made free by the grace of God, which contradicts the Pelagians. And where he says that the evil, which was not before, came from him, it is against the Manicheans. Moreover, in the eighth chapter, Pelagius says that the thing which is good may be accomplished sooner if grace helps thereto. By this addition, “more easily,” he truly signifies that he thinks that, although the help of grace is lacking, he can, albeit with more difficulty, perform what is good by free will. Again, the same in the forty-seventh Epistle to Valent.[note 39.5] That man, he says, falls into the error of the Pelagians who supposes the grace of God is given for any merit of man, which grace alone makes man free through Jesus Christ our Lord. But again, he who thinks, when the Lord shall come to judgment, that man is not judged after his works which might now, by reason of his age, use the free choice of will, is nevertheless in error. He says essentially the same thing in the second book, chapter 18, concerning the merits and remission of sins.

### Paragraph 590

With this doctrine of Pelagius—that is to say, darkened (for so John himself explains a little later, saying that a third of them was obscured, and so on)—the third part of the sun, namely Christ, the true sun of righteousness, was struck. For the Pelagians' doctrine denied the grace of Christ, and with human merit trampled underfoot the merit of Christ. By this, also the third part, that is to say a great part, of the Moon, namely the church, is shown to be struck and darkened; moreover, the third part of the stars, I mean preachers and ministers, have not taught with the light that became them. For histories witness that this heresy has severely infected diverse parts of the world, so that even bishops and learned men have followed this noisome error. A synod of bishops was assembled in Palestine in the East, which drove Pelagius to recant. They also disputed sharply against the Pelagian doctrine in Rome, and councils were assembled which condemned it. Synods were assembled in Africa, and after much reasoning, sentence was pronounced against Pelagius. For many were daily taken with this infection. For the doctrine is pleasing, which even today lacks neither maintainers nor defenders. It seems godly, and for the study of virtue needful, to affirm free will and human merit; again, it appears licentious to attribute all things to God’s grace.

### Paragraph 591

He adds that neither the day shone with a third of its light, nor the night with a third of its light. For just as grace could not be fully understood through the doctrine of Pelagius, neither could sin. And Saint Augustine in the second book of original sin, chapters 23 and 24, says that Christian faith properly consists in the account of two men. For by the one we were sold into sin, by the other redeemed from sin; by the one cast down into death, but by the other delivered to life, and so forth. And while all these things are spoken, they are spoken to this end, that we might beware of those heresies.

### Paragraph 592

And thus far we have spoken of the four trumpets and the greatest conflicts in the church. Three trumpets remain, and before them, a brief introduction is set, so that the minds of the listeners might be stirred.

### Paragraph 593

And John says how he saw an angel flying through the midst of heaven, and heard him crying: Woe, woe, woe to the inhabitants of the earth, because of the things that would happen to people when the other three trumpets should be blown. Therefore, to each trumpet is joined a woe. We express this very well in Dutch by “woe, woe, woe.” For the Greeks say, and John wrote in Greek, [? *a word lost here*]. And it signifies truly that the times of the former conflicts were sharp, but that those that follow shall be much sharper and crueler. For I told you in another place that this word “woe” encompasses the evils both of this present life and of the life to come, of body as well as of soul. Therefore, the times of Papistry, Mahometry, and of the last judgment shall be most dangerous.

### Paragraph 594

Note 39.9: The Complutensian copy has an eagle where we read an angel flying through the midst of heaven; perhaps because it found it so in Arethas. Indeed, the common translation, commonly called Saint Jerome’s, has an eagle for an angel. And therefore Primasius reads it so likewise, which seems to have followed the old translation in all things. But the eagle is swift and of most sharp sight, signifying the almighty knowledge of God and expedition unspeakable in doing of things.

### Paragraph 595

Note 39.10: The fifth trumpet describes a most terrible battle, which the Pope instigated by introducing errors into the world, indeed by promoting, establishing, and defending them through his ungracious Locusts that consume everything. He will endure until the end of the world. He will discuss him more fully and appropriately in chapters 13 and 14, and following.

### Paragraph 596

[note 39.11] The origin of this evil is attributed to the fall of a star. For a star has fallen from heaven to the earth. Stars, as I showed you at the beginning of this book, around the end of the first chapter, represent the state of ministers, or bishops. For just as stars shine in heaven, so bishops, illuminated with heavenly light, ought to shine in the church, both in doctrine and in an honorable life. And they remain in heaven so long as they perform their duty; they fall to the earth when, forgetting heavenly conversation and doctrine, they think upon earthly things, speak of, and pursue honors, pleasures, and such like corruptions. A little later, he will call him an angel, whom he now calls a star. [note 39.12] The Church of Rome was notable and pure, commended also by the praise of the Apostle. It had bishops, that is to say, ministers of the church, under Emperor Constantine, around 32, for the most part very well learned, most holy (yet men), and most glorious martyrs of Christ. Again, from Emperor Constantine to Gregory the Great, there are counted bishops or pastors of the church of Rome, around 32, among whom there were not a few diligent, learned, and godly. But there were also among them those who, blinded by the evil of ambition, began to incline more to seek honors and glorious titles than the doctrine of Christ concerning humility and simplicity, and the example of Christ and the apostles. Christ fled when the people would have chosen and made him king. He said that kings should reign, but that apostles and their successors should serve. If kings had offered them realms and riches, they should not have received them. The stories plainly declare what certain bishops of Rome practiced with the churches of Africa, and how they would have ruled over them. Nevertheless, among the later bishops were figures like Pelagius and Gregory, surnamed the Great, who gravely accused the bishops of Constantinople for attempting to establish the church of Constantinople as chief of all others in the world, and the bishop thereof universal. Neither was Gregory ashamed to say expressly that he is the forerunner of Antichrist, whoever would covet the name or title of universal bishop. But Boniface the third of that name, not long after Gregory’s death, moved nothing herewith, required and obtained from Emperor Phocas that the church of Rome might be called and taken as the chief and head of all churches. [note 39.13] By this, the bishops of Rome were plucked out of heaven and cast to the earth, utterly beginning to cleave to earthly things, to care for earthly things, even to aspire to the empire and chief rule and government. Here you have what star fell from heaven to the earth.

### Paragraph 597

And to this star (he calls him afterward the Angel of the bottomless pit) or Bishop (I name one, I understand all of that state and succession in that seat) was given the key of the bottomless pit.[note 39.14] Christ truly keeps the key of David: as I showed in the second chapter of this book. He gave to the Apostles the keys of the kingdom of heaven, power to open or to shut heaven—that is to say, the ministry of preaching the Gospel, whereby is shown and assuredly promised the forgiveness of sins and eternal life to believers; and the retaining of sins, and certain damnation is threatened to the unbelievers. No godly man doubts but that these keys were given also unto Bishops of Rome; yet every man knows that the later popes would not use them lawfully, but corrupting the Evangelical truth and infecting the lawful ministry, have obtained counterfeit keys. Therefore, to them of the Prince of darkness is given the key of the bottomless pit, namely corrupt and counterfeit doctrine, and not the apostolic, but apostate ministry, whereby, as it were from hell set open, they have brought forth outrageous errors and superstitions, and ungodliness of all sorts.[note 39.15] And I suppose it has happened not without God’s providence that Bishops of Rome are called *Clauigers* or key bearers, and wear keys in their arms. But you shall not understand them to be the keys of the kingdom of heaven, but of the bottomless pit rather: for he is a teacher of errors and of all abomination, author moreover of all wars and dissensions, leading them even unto Hell.

### Paragraph 598

God is indeed the source of unending goodness and all truth. This was revealed in Christ by the Apostles through the preaching of the Gospel, and it refreshes with wholesome water all who thirst for eternal salvation. Isaiah refers to this fountain in chapter 55. And Jeremiah in chapter 2. The Lord also speaks of it in the Gospel according to John, in chapters 4 and 7, and in various other passages.

### Paragraph 599

[note 39.17]Against this lively fountain of ever-flowing waters, is set the bottomless pit, unsearchable I say because of the malice of Satan, full of ungodliness, abomination, and a kind of lying. From hence blubbereth up into the world by false teachers and ministers of Antichrist whatever error and abomination there is in the world. For Satan, the father of lies, spreads abroad in the world by his instruments whatever darkness there is.

### Paragraph 600

Therefore the star or angel of the bottomless pit—that is, the Pope or Bishop of Rome—opens the bottomless pit with a key, and shortly thereafter ascends with the smoke of the pit. For I have spoken thus far of the beginning of evil; now I shall speak of its progression and manifestation.

### Paragraph 601

The Pope, through his corrupt ministry, opens the way to Hell, and not to Heaven. From Hell ascends a smoke. Smoke in some places of Scripture is a token of God’s presence, wrath, and vengeance, as when a smoke rose in the Temple of Solomon, 1 Kings 8, Isaiah 6:19. In Exodus 19:18, we read that smoke ascended from the mountain as from a furnace. You read in Psalm 18 that smoke went up in God’s wrath, and fire burned before his face. At present, smoke seems to signify hurtful and devilish opinions. Smoke hurts the eyes and prevents clear sight of the truth. So also does perverse doctrine; it darkens the eyes, takes away judgment, and blinds with error. And rightly do they suffer these things from the smoke of God’s wrath and from the lies of deceivable men, who have forsaken the light of the Gospel and the clarity of God’s truth. Under the name of this infernal smoke are contained the opinions and abominable doctrine that the Bishop of Rome, as he is the prelate of the chief church and Apostolic See, is also to be universal pastor and Apostolic, moreover the head of the militant church, the vicar of Christ on earth, whose voice must be heard as much as Christ’s himself; that he has full power in the church, the keys of the kingdom of heaven, and so forth. That he ordains and gives to all churches bishops or pastors, who should govern all other churches according to the prescription of the church of Rome, and so forth.

### Paragraph 602

But how great this smoke is, and how powerfully it is expressed: it ascends, he says, like the smoke of a great furnace. It signifies that the papist opinions and doctrine are thick, or gross, manifold, and apparent; where in truth they are nothing but smoke and vanity, puffed up and vain. But it is of such a power that it darkens the sun and the air. I have told you often that Christ is the sun of righteousness. And we call the air the wholesome doctrine, by which the souls of the faithful are refreshed. Therefore, by the papist doctrine, the sun and the air—that is, Christ and the Gospel—are obscured. Christ is the universal pastor, the high and only Bishop, the head and health of the faithful, who freely forgives sins; which is preached by the Gospel. This doctrine becomes corrupt when the Pope is admitted as head of the church, with the full power of granting indulgences for all sins. Thus is the sun darkened.

### Paragraph 603

However, the evil proceeds further, and establishes itself in the church with much greater effect. For out of the smoke came forth Locusts upon the Earth. For when, through the false persuasion of corrupt doctrine, the eyes of all men were blinded and did not rightly look upon Christ and his only gospel, and all men revered the Pope as the vicar of Christ, the head of the church, and a man apostolic, as it were the mouth of God, and he now made bishops and priests, and nourished, advanced, and established monks and friars: an infinite multitude of the clergy increased most abundantly, in number that could not be numbered. For he himself immediately in the following words, and with a fuller exposition, declares that he speaks nothing of those little worms, the Locusts. For he says, and it was commanded them, that they should not hurt the grass or hay of the earth (and truly, the clergy does not live on hay), neither any green thing, nor any tree, but men only. As though he should say, I speak nothing of grasshoppers such as in times past destroyed Egypt, but I speak of pestilent men, afflicting men with the poison of doctrine. But a little after, they are so described in every point that no man need to doubt that the false clergy is thereby signified. The which thing Primasius also saw, who in his commentaries upon this book said: he puts the authors of evil doctrine. For just as the Locust hurts with its mouth, so do they tear with their preaching, as we read, greedy wolves not sparing the flock, and so forth. Thus says he. There are also other causes why he likened the false clergy to Locusts. If the Locust is alone, it seems most contemptible; so there is nothing more vile than a solitary monk or friar, priest or sophist. But if they swarm together, they are a terror to men, and cannot be driven away with any force; they eat and destroy all. When the prophet Joel would show a great evil to come, he says that the Locusts will come. Sometimes they sing, leap, live at ease and pleasure, to the loss and hindrance of husbandmen. The same things may you see also in the clergy. I speak nothing here of holy priests, that is, lawful ministers of the church, of good men, honest and learned; I speak nothing of the ancient and holy monks, which were burdensome or grievous to no man, and were no preachers, but very laymen, getting their living with their hands, in the church subject with other faithful to the pastors of the church, and so forth. I speak of the unlawful, sluggers, idle bellies, devourers of victuals, but chiefly of false teachers.

### Paragraph 604

And doubtless the Popes’ clergy is most rightly compared to grasshoppers or caterpillars. For both are innumerable, and they occupy and consume all things. In times past, the ministers of the churches might be numbered, for the number was but small; neither were unprofitable or unnecessary persons nourished of the church goods. There remains a constitution of the emperor Justinian, where among other things, we ordain that there be not at any time in the sacred great church above sixty priests, men deacons, and one hundred subdeacons, ninety readers, and one hundred and ten, nor above twenty-five singers; that the whole number of the clergy of the greater church may consist in four hundred and twenty-five persons, and besides one hundred doorkeepers, as they term them. Therefore, in the most holy great church of this our noble City of Zurich, and in those three churches to the same united (to wit, in the church of our Lady, Saint Theodore, and Saint Irenaeus, let there be so great a multitude of the clergy. This some of the ministers of this imperial city and most large church established five hundred and twenty-five persons. But how many at this day may you find at Rome, or in another great city, priests, monks, friars, and nuns? They exceed this number four times and more. And to leave out many things that might here be brought in, Pope Pius Sabellicus shows in the ninth book of *Aeneidos*, the seventh chapter, that the sect of gray friars was so greatly multiplied throughout the world that then they held and possessed forty provinces, and under every one diverse cloisters and convents (wardens they call the rulers), and exceeded the number of three score thousand men; in so much that the master of the whole order, whom they call general, has been heard many times to offer the pope, preparing an army against the Turks, thirty thousand fighting men of the order of Saint Francis, which should be well able to serve in the wars, and yet be no hindrance or let to their religion or service. And now who is it that knows not how many orders there be of monks and friars? You may therefore account other orders after the rate of the order of Saint Francis, and though you attribute to every one but the one half of that number, to what a sum will it amount? To these if you add the colleges, more and less, throughout so many dioceses, persons, vicars, chaplains, and parish priests, you will grant that not without cause the popish clergy is compared to locusts.

### Paragraph 605

I need not explain with many words how they seize upon and consume everything. It is commonly said that whenever you see any place, fertile and wholesome, wherever you ride or go, you will find it full of clergy and possessed by religious men.

### Paragraph 606

He also speaks explicitly of the power of these Locusts. [note 39.26] He presents them as a parable: and he says that power was given to them as scorpions of the earth have. A scorpion is a flattering and, in appearance, a domestic worm, which suddenly strikes with its tail, or rather with the sting of its tail, and so poisons. Therefore, with flattering words, the clergy of Antichrist deceive and wield power through the poison of venomous doctrine. Thus the Apostle also speaks of false teachers in the sixteenth chapter to the Romans. Through fair speech and flattering, they deceive the hearts of the simple. Their power, therefore, is nothing other than evil doctrine, with which, as with the venom of scorpions, they infect simple Christians, but especially those who despise the doctrine of the Gospel.

### Paragraph 607

For it follows a declaration concerning whom these Locusts may harm. There are two kinds of men. One, willingly and knowingly, will perish, and are the open and professed enemies of the holy Gospel; whom, by God’s just judgment, these Scorpiolocusts destroy with their poison. The other is more simple, erring rather through ignorance than through obstinate malice; these are not stung by the Scorpiolocusts, for they have a seal in their foreheads (as spoken of in chapter 7). For the power of this evil is limited, and not without measure. Therefore, the locusts were given this power, that they should not kill—not those wicked who would rather die than live—but these simple ones. They harm them truly, but not to death. And they torment them for five months. And that torment is the trouble of the conscience, which they torment with threatenings, hypocrisy, and wonderful terrors.

### Paragraph 608

There is added for a comfort, five months. [Note 39.28] The locusts certainly come out in the month of April, and live until September, and when they have lived completely five months, immediately they die. It signifies therefore that those who are consecrated to godliness shall feel these torments for a little while: neither that the deceivers shall always prevail: but that there shall be spaces to rest and breathe in, wherein the godly through the truth may be recovered. For the locusts destroy not, and are seen all the year long. There seems therefore to be a comparison here in this determinate number, that the sense should be: like as the locusts live not longer than from April to September: so doubtless there is a time prefixed to those seducers, and false Papist clergy. Even thus has also the Apostle Paul himself comforted the church: which after he had prophesied that the church should be wonderfully vexed of hypocrites and false teachers, immediately he adds: and like as Jannes and Jambres resisted Moses [Note 39.29], so do these resist the truth, men of a mind corrupt, and lewd as concerning the faith: but they shall prevail no longer. For their madness shall be manifest to all men, like as that was of the other. And Primasius says that they are meant here, which although they were entangled with false doctrines, yet having remorse about the end of their life, they receive God’s verity. Again we see, as I warned you in Chapter 7, that all did not perish who were once entangled with the snares of Antichrist. For at the length through the mercy of God they escaped, and sought the grace of God to be given them through Christ, forsaking all superstitions. We see moreover, by reading of histories, how God has at certain times opened the truth by his faithful ministers, through whose preaching the lewdness of the Locusts is interrupted, that men began to smell them out, and to eschew the same; notwithstanding the regenerated, many times have returned, &c. And likewise other ministers have returned home, &c.

### Paragraph 609

And furthermore, he declares how great was or is the force of this evil. Their torment, he says, is like the torment of a scorpion when it has struck a man. At first, there is no great pain felt, but little by little it gathers strength, and at last it aches exceedingly. If remedy is had in time, the poison is not deadly; if it is not taken, he who is stung with it dies. To the declaration of this torment, which men feel in their consciences, pertains this that follows: in those days men shall seek death, and so on. And it is a like phrase of speech in a manner as is that same, “mountains fall upon us and cover us,” and so on, whereof I spoke in chapter 6. And it is the voice of one who is sore afflicted and brought in a manner to despair. Doubtless the popish doctrine of merits, of monastic perfection, and of other such like doctrines, have driven many headlong into despair. Hereunto is added that the times of the locusts were most full of sorrows, whereof all histories complain. The life was not pleasant; the locusts did so set men together by the ears among themselves, and so on. And to be brief, they brought men into such a case that they wished to die. The Lord Jesus deliver us from the poison of these locusts.

## Sermon 40: ¶ The Locusts are described by a maruelouſe Hypotipoſis, the Popiſh clergy: and is showed, of what sort the Antichristian war shall be.

### Paragraph 610

.

### Paragraph 611

And the likeness of the locusts was like unto horses prepared for battle, and on their heads were as it were crowns, like unto gold: and their faces were as the faces of lions. And they had hair as the hair of women. And their teeth were as the teeth of lions. And they had breastplates, as it were breastplates of iron. And the sound of their wings was as the sound of chariots when many horses run together in battle. And they had tails like unto scorpions, and there were stings in their tails. And their power was to hurt men for five months. And they had a king over them, which is the angel of the bottomless pit, whose name in the Hebrew tongue is Abaddon, but in the Greek Apollyon.

### Paragraph 612

We have already spoken of the origin and power of the locusts. [note 40.1] Nevertheless, lest anyone should be hindered by obscurity, so that he could not know the locusts, and beware (for this is the end of the whole prophecy to understand the mysteries of Antichrist, and beware), now also he describes the locusts with a wonderful figure, and their fight against Christ, and against the doctrine of godliness of all other most perilous.

### Paragraph 613

And there is no doubt but that the whole army of the Pope is here described, [note 40.2], especially the spirituality as they term it. For the soldiers of the Emperor, kings, and all princes serve him, whom they call secular. But in the Pope’s tents of the spiritual army are Cardinals, Patriarchs, Archbishops, Bishops, Abbots, Prelates, and there is no noble of priests, and religious persons of both sexes. Hereunto appertain many universities, Doctors, and Masters, great champions of the Pope: these are verily those locusts of whom the Lord Jesus speaks here. I know how displeasingly many will take this my exposition. And I would gladly (God is my witness) have spared them: but all the blame is in them, which in words and works betray and declare themselves to be locusts. For except the thing itself cry out that those things are done of them, which by the exposition are now brought to light, I will not desire that credit should be given to me. I speak nothing here in the favor of any man, neither for hatred. Let God himself be judge betwixt us, let the truth itself judge. Certainly all expositors with one consent understand by locusts false teachers.

### Paragraph 614

But let us see the description of the Apostle John by the revelation of Jesus Christ, which does no harm to anyone, which does not slander anyone. [Note 40.3] And he shows the similarities, that is to say, the likenesses of locusts, by which they may be depicted, and as it were set before our eyes, to be like the things which he brings forth. For unto every part he applies a parable or similitude, whereby he expresses most aptly the disposition and manners of the locusts.

### Paragraph 615

First, he says that the locusts are like unto horses prepared for battle. By this parable he signifies many things at once: that the clergy should not only be ambitious and proud (for a horse is an image of pride) but moreover rebellious and bold, and even cruelly fierce, and in their incredulity and in all their errors most obstinate. They are utterly ignorant of repentance. For John seems here to have alluded to these words of Jeremiah: “How does it happen that this people is not turned away from so froward an aversion? They cleave stiffly to deceit; they refuse to return. I have marked and heard, and they spoke not right; there was none that was sorry for his evil, and that would say, ‘What have I done?’ Every one of them did run his course, as it were a horse dislodged into battle; certainly, with this kind of people there is no amendment.” They think rather about how they may allure others into errors with them. He signifies moreover that the clergy shall be warlike and the authors of wars, and shall instigate wars against the saints and true worshippers of God. For they have the secular power, as they call it, ready. For a long time now there have been no wars that have not been raised by this kind of people. Histories bear witness hereof. Yea, and in this our time cardinals and bishops have led armies, and so forth. Finally, it is signified hereby that the clergy shall continually vex and weary with spiritual war also, the true church of Christ. Wherefore, in chapter 11, we shall hear how the beast comes out of the bottomless pit and makes war with the excellent prophets of God. They mix therefore and practice as well spiritual as corporal wars. Last of all, it is signified that the popes’ clergy shall be well fed, fair and well-looking, and given to voluptuousness, lusts, and pleasures of the body. For this kind of people represent not horses that are gaunt or lean, such as go to plow and cart, but such as are well kept and fed even to serve upon in the wars. For behold with me and consider of what sort the clergy is (for the most part), and you will say that they are here set forth in their colors.

### Paragraph 616

Secondly, upon their heads, he says, as it were crowns like unto gold. Rabanus Maurus in chapter 3 of the first book of the Institution of Clerks calls the shaving of the priests’ crown a kingdom, truly a token of the dignity of a king and priest. For priests and monks or friars boast themselves to be kings and priests, and yet in deed are neither of both. For the true faithful before God are kings and priests, 1 Peter 2. But by the ordaining or shaving of the Pope, they receive nothing either of kingdom or priesthood. John, therefore, says upon their heads, as it were crowns like unto gold, for he says not they were crowns, but like as they were crowns of gold. They were not crowns in deed, neither were they due to them. And yet notwithstanding in the end of the world now they have taken upon them diademes, or miters, and crowns of gold also, and the same most precious. Yet have they done this by no right. In times past, Bishops of Rome wore white miters, in token of purity and sincerity, finally of the knowledge of both Testaments; but none of the Apostles nor apostolic men wore them. Therefore they betray themselves like a rat with their own utterance, the which I suppose to be done by God’s providence, that they might be known, and eschewed of Christ's sheep as crowned wolves.

### Paragraph 617

Their faces were like the faces of men, not like the faces of locusts. Similarly, in Daniel, the Antichrist is attributed the eyes of a man, signifying industry and policy. These men pretend to great humanity; they are furnished with fair speech. One might think that if humanity were truly elevated, it would be found in them. But they feign these things, with the intent that by creeping thus into people's hearts, they may accomplish their purposes and deceive. In cunning, deception, shrewdness, and practice, as they term it, the Popes’ Legates, Ambassadors, Priests, and Religious persons excel all other wise men of the world. They press into all assemblies of all men, seeking to be made privy to all things. They consider everything as a means of accomplishing their purposes, they feign and dissemble everything, and they can easily outwit and beguile even those who are most astute. Moreover, they are learned, witty, eloquent, and wonderfully crafty in all things. The thing itself speaks and testifies that I write the truth.

### Paragraph 618

And they had here, like the appearance of women: by this likeness he notes their wantonness, idleness, whorish apparel, and effeminate minds. For they are combed and adorned, and very finely apparelled, delighting in women’s jewels, wearing costly garments, especially in the church, where they ought above all to show humility and frugality. Which of the Apostles ever went so decked (or rather disguised) in the temple or outside the temple? The excess and costliness of apparel of Priests and Monks gives no place to the costly array of the Persian Kings. Again, the thing itself speaks. Saint Augustine in a commentary on the seventh of the Apocalypse, in this appearance, says he would understand and show, not only an effeminate or womanly sex, but also either or both sexes. This is what he says, which I leave to be construed and examined by others.

### Paragraph 619

In attributing to them also the teeth of lions, he signifies their cruelty against the poor and faithful followers of Christ. They are most cruel in persecutions, and of blood most thirsty, and are not moved in this by any compassion. They destroy all things with the sword; many devise various tortures. They excel in tyranny Busirides and Phalarides; the thing itself speaks again. For if kings, princes, or magistrates would spare these wretched people, the priests and friars cry out that it is not lawful; finally, they incite the minds of all princes and magistrates against Gospel adherents by prescribing forms of inquisitions and oppressions. To this is added that some of them are hoarders, accumulating with insatiable covetousness and religious robberies, kings’ treasures. Again, some other wasters follow, who spread evil-gotten goods and waste them prodigally in riot, debauchery, whoring, or wars. Therefore, the teeth of lions are rightly attributed to them, in the same way as Amos is said to have attributed to the false prophets. They had also breastplates (thoraces), a defense for the breast, called a breastplate or a vanguard. Others interpret it as sureties, but they cover all the body; breastplates properly cover the breast. By this is signified that their hearts should be obstinate and inflexible. They are stiff-necked and rigidly bound, neither departing one hair’s breadth from their errors, but even assert that the sea itself cannot err; yes, and that the Pope cannot err. For they will not abide to be taught and admonished, but plainly say that the Church of Rome has never erred; therefore, there remains nothing but that you must subscribe to it, or else be condemned as a heretic and suffer death. It is signified moreover that these shall be through another man’s protection most safe. For they have their immunities, they have their privileges, they have the secular power always ready to fight at their request, they have their fraternities, fellowships, leagues, and affinities. What should we say that bishops and abbots are the sons, brethren, and cousins of princes? Whoever therefore touches them, he has touched the apple of the prince’s eye. For even for maintaining them and their state, all men fight as it were for life and lands.

### Paragraph 620

To the Locusts moreover are ascribed wings. For they are exalted above the common state of men, while they are taken and accounted the most fortunate and most excellent in the world, and so forth. Yea, and impudently they boast that herein they are worthier and greater than the Virgin Mary, for that she bore one in her womb, the Son of God, but they can call him daily to the altar? And while they fly, they make such a noise as horses do, drawing warlike chariots, and now ready to invade the ranks of enemies: that is to say, all their doings are most vehement, most warlike, to men horrible, and deadly. Hereunto pertain the clamorous disputations of the Sorbonne and other schools, excommunications, sentences given at Rome, the popes’ bulls and writings, the boastings of decrees, and they are in obstinacy invincible. All these things make a noise together, and thunder terribly.

### Paragraph 621

Furthermore, it is added that through their decrees and counsels they will tear apart or invade. Consequently, Daniel also attributes prosperity to Antichristians, saying, “He shall act, and he will prosper.” And they will invade in such a way that, as we have said before, people will desire to die, believing that there is no deliverance.

### Paragraph 622

[note 40.12]I have spoken before, in a previous sermon, of the tails of scorpions and of five months. Their venomous doctrine is noted, which nevertheless at certain times shall be reproved, that godly men may beware of it. And who does not see, who does not feel, how grievous or difficult this fight or battle is, waged by such locusts? Therefore, the Lord’s mouth has rightly joined a curse with the locusts. Men would rightly wish to die, so that they might be delivered from such great dangers. Let us weigh and consider these things at this day, and let us pray that we may overcome and escape the most pestilent poison of Antichrist.

### Paragraph 623

For now also the king of these locusts is brought forth, and is pointed out as it were with the finger of Christ. He sets him out by three titles, that he may be better known. The locusts, he says, have a king over them. This king is not lawfully given to them, but they themselves have chosen him. For who does not know that, through the policy of the spiritual fathers, the Pope has been exempted from the jurisdiction of princes, so as to rule over all spirituality? They acknowledge no other magistrate than the Pope of Rome, and they rail upon secular princes (as they call them), and will not obey them. All of them bind themselves and swear allegiance to the see of Rome, which they seek to preserve and maintain, even if all others perish. The form of swearing is known, made by bishops, abbots, and doctors to the Pope. And if kings and princes should so much as touch with their finger one who is anointed with the bishops’ oil, even if he be a church robber, a murderer, a thief, and a parricide, they are held accursed, and they and their realms are excommunicated. Thus I say, the locusts have the Pope as their king.

### Paragraph 624

It is also called the Angel of the bottomless pit [?]. And immediately in the eleventh chapter, he shall be called the beast which ascends out of the bottomless pit. Christ descended to us from heaven, the Angel of the Testament and great counsel. Whoever disregards him rightly hears the angel of the bottomless pit, that is to say, Antichrist sent by Satan himself from hell. For he is the adversary and enemy of Christ, in whom the Devil dwells corporally, as also thought Saint Jerome, that the Devil should wholly inhabit that great Antichrist.

### Paragraph 625

Therefore also a true name, and a true title most fitting is given him. For they lie who greet and call him “most blessed father,” “most holy Pope,” and so forth. Christ sets forth this figure with another style and gives him other titles. His name, he says, was Abaddon in Hebrew, and in Greek Apollyon. He publishes his name in either tongue for no other reason than that in either Testament—whereof the one is written in Hebrew, the other in Greek—this title is attributed to him. Abaddon or Abbaddon, or Apollyon signifies a destroyer. But Daniel in chapters 7:8 and 11, and Zechariah in chapter 11, attribute this virtue and property to Antichrist. Paul calls him the son of perdition, that is, most lost, most damnable, and the greatest author of perdition and damnation, which finally shall be to many an author of slaughter, by sundry wars. For through false doctrine he destroys souls, and with tyranny by fire and sword he wastes the land, and those who refuse to obey him, most cruelly. Let the Popes’ acts be considered, and the practices of spiritual fathers; let them be applied to these oracles of God, and then let a comparison and judgment be made. And this is as it were the key, opening to us the sense of this place, and that it should be expounded of Antichrist, whom Paul called the son of perdition. Habad in Hebrew signifies lost or destroyed, and from that comes Habbedon, perdition or destruction. So in Greek Apoleo and Apollume signifies to lose and destroy; from this is Apollyon. The Lord Jesus shall slay this destroyer with the breath of his mouth, and take him away utterly by his glorious coming.

## Sermon 41: ¶ The ſixte trompet is expounded, where is created of Saracenes and turkiſhe matters.

### Paragraph 626

.

### Paragraph 627

One woe is past, and behold two woes come yet after this. And the sixth angel blew, and I heard a voice from the four corners of the golden altar, which is before the eyes of God, saying to the sixth angel who had the trumpet: “Loose the four angels who are bound in the great river Euphrates.” And the four angels were loosed, those who were prepared for an hour, a day, a month, and a year, to kill a third of mankind. And the number of the horsemen were twenty thousand times ten thousand. And I heard the number of them. And thus I saw the horses in a vision, and those who sat on them having fiery armor of a yellowish and brimstone-like color, and the heads of the horses were like the heads of lions. And out of their mouths went forth fire, smoke, and brimstone. And of these three, a third of mankind was killed—that is to say, by fire, smoke, and brimstone, which proceeded out of the mouths of them. For their power was in their mouths, and in their tails, for their tails were like unto serpents, having heads, and with them they hurt.

### Paragraph 628

The sixth conflict is that of Mahometanism, by the Saracens, Turks, and Tartars, fought with great cruelty and much suffering. And would God it were fought. For we perceive daily through the events themselves the mystery of the prophecy, and see its fulfillment, and even experience it as well.

### Paragraph 629

At the sound of the trumpet of the sixth angel, John hears a voice from the four corners of the golden altar, that is to say from the midst of the altar: and there is no reason why we should seek a mystery in the fourth number. And he speaks of that altar which is before the eyes of God. That voice commands the angel trumpeter to release the four angels bound in the great river Euphrates. As soon as this was done, an innumerable army of horsemen marched forward, and slays and destroys a third of the earth, that is, a third of mankind. And the description of those horsemen, and their force or power, is most diligently given.

### Paragraph 630

We have recently learned that the golden altar signifies Christ, sitting at the right hand of the Father. He is purer and more precious than gold; he is priest and sacrifice for all the faithful, standing before God’s eyes, that is, pleasing God, in whom or through whom his soul is pleased with all faithful, whose virtue is sufficient for all. And God the Father desires that such a one be preached and believed by all the faithful in the world. The ancient church, through the Apostles, believed and taught this same Christ; until, by the work and instigation of the Devil, most corrupt men arose in the church, some denying Christ’s deity, others his humanity, others separating the person consisting of God and man, and others confusing the natures or properties of the natures. God’s goodness patiently endured this assault, many times sending faithful and open defenders of the truth, who might root out these blasphemous errors—which we have read to be done by various bishops or preachers of the church, or by ecclesiastical assemblies, which we call councils: such as were the councils of Nicaea, Constantinople, Ephesus, and Chalcedon, in which Arius, Macedonius, Nestorius, Eutychius, and other monsters of heretics who opposed Christ were condemned. Nevertheless, the incurable perverseness of men gaining the upper hand, there was no end to alterations and blasphemies. For two great bishops of no small churches, Peter, patriarch of Antioch, and Severus of Constantinople, arose during the reign of Emperor Justinian, and impudently and most wickedly affirmed (as the acts of the Fifth Council of Constantinople amply declare) that the body of Christ was utterly incorruptible and truly deified, not subject to any affections as mortals are. For they asserted that the Word became flesh so that immediately it began to be one nature, namely, divine, that Christ was made [illegible], that is, incorruptible. These things seem to proceed from the most wicked school of Valentinus, Marcion, and Manichaeus. James of Syria, surnamed Zalos, of whom the Jacobites are named in the eastern lands, took it upon himself to defend the doctrine of Severus. He taught that Christ, because he was incorruptible, neither suffered nor was crucified, but that some other person suffered for Christ, while Christ only stood by invisibly and looked on. This opinion was refuted by many testimonies of Scripture and finally repulsed and overthrown with the articles of our faith.

### Paragraph 631

For we profess in our belief that he who suffered under Pontius Pilate was crucified, dead, and buried. The prophets declared expressly beforehand that he should suffer and die; that he has suffered and died, the Apostles have witnessed, among whom John saw the death and suffering of Christ on the cross. Neither do we read that the Lord was ever so much offended with his disciples as he was with Peter, who sought to dissuade the suffering as unworthy of the Son of God. For he said, “Go behind Satan; you do not favor what is of God, but what is of me.” Therefore, they should not renew the error and madness that has been refuted. He seems to reason plausibly that God might have redeemed the world by another means than through the incarnation or suffering of the Son of God; to think it unworthy that the Son of God should be beaten with the blows of the wicked, and moreover slain. But this plausibility is from unclean flesh, not from God, indeed it is from Satan himself. Yet this absurd and most wicked opinion has found not a few followers. For the heresy of the Jacobites, contained in the Quran, is spread far and wide throughout the East. It is plain that the golden altar was by them most filthily polluted; and the merit of Christ’s suffering denied, the dignity and majesty of the priesthood and sacrifice of Christ trod under foot. There were, besides these, other most corrupt opinions in the West, and so forth. This thing worthily kindled the just wrath of God. For, in his just judgment, he permitted Muhammad to make new laws and to spread Jacobism far and wide throughout the world. For those who will not hear Christ worthily hear Antichrist, which thing the Apostle affirmed also in the first chapter to the Romans and the second chapter to the Thessalonians. Therefore, a voice is heard from the altar of him who sat on the right hand, commanding the four angels bound in the river Euphrates to be loosed, that is to say, to bring forth into the world destroyers who may overrun a great part of the world.

### Paragraph 632

The heresies of the Nestorians, Jacobites, Monothelites are gaining force, as are monks and friars, and the Benedictine order of Cas-sina. Moreover, idols and images are increasingly present in the church, gaining strength, and the pride and unfaithfulness of the bishops are growing. Mahomet, the destroyer of the world, was born in Mecca, a city of Arabia, of very obscure parents. He was raised by Sergius, a vile monk polluted with all kinds of heresies. After he reached the age of twenty-five years, he claimed to be the prophet of God. Through sedition, he was driven out of Mecca, where a great rabble of Jews, Jacobites, Christians, pagans, and heretics had gathered. He then went to neighboring cities and houses, and for ten years secretly instilled his doctrine into the people, so that throughout Arabia there was found a great multitude of Mahomet’s sect. Then Homar, a bold man, taking about seventy other men ready to fight, asked Mahomet what he would do. He answered: “Truly, my will is that, executing the commandments of the law, you cleave to it in riches and poverty, and cleave together with mutual and steadfast love; that you do not defile other men’s wives through adultery; that you abstain from evil and prohibit others; that you do good yourselves and persuade others; that you make war in the name of God, and that by fear and force you set forth the laws to the disobedient; for these things truly I promise you Paradise.” At this talk, they gave each other their faith. Homar, with his sword drawn, swore that he would not suffer the preaching and law of Mahomet to be kept any longer secret. Thus, by preaching and by the sword, Mahometism prevailed greatly in a short space. They broke into Mecca, put down other religions, and beheaded those who resisted. There, this new Solon, Mahomet, proclaimed a new law in the Temple at Mecca. A great multitude of servants and wastrels resorted to that wicked fellow, who sent ambassadors to the people around him and solicited them to receive his religion, persuading many to wickedness. These things were done under Emperor Heraclius about the year of our Lord 620. And that wicked and most absurd law of Mahomet yet remains, and is called the Alcoran, so that it needs no further declaration. Nicephorus, in his history, says that the Saracens began the desolation of the whole world. Saracens were called the followers of Mahomet. Certainly, they subdued Arabia. The Saracens and Persians invaded Syria and Egypt, Chaldea, and Armenia. Afterward arose the Turks and Tartarians, receiving the religion of Mahomet, who have subdued in a manner all the provinces of the Roman empire in the east and to the south.

### Paragraph 633

By the river Euphrates—famous throughout all of Asia—floated Babylon, the seat of the Eastern Monarchy. The mighty peoples of the East, the Assyrians, Babylonians, Medes, and Persians, who had ruled the world before the monarchies of Greece and Rome, seemed to be drowned, buried, and hidden, even bound within that same river. For the Macedonians of the West were governors of the world, and after them the Romans; and these mighty nations that we now name served them. But after the golden altar, as I said, was defiled, and countless people in the East and West revolted from the true Christian faith, God stirred up again the Eastern destroyers of the world, which had lain dormant as if asleep. For the prophets testify that these nations were the scourges of the world; therefore, God brought them forth again by his just judgment. Certainly, we read in the tenth chapter of Daniel that there was an Angel of Greece and an Angel of Persia, and that by them the entire people are understood. So now the nations of the East—Arabs, Saracens, Turks, and Tartarians—are raised up, who, for sin, might devastate the world, and the East might rule again, as Lactantius, from the Sibyl, prophesied would come to pass: let the West serve.

### Paragraph 634

Let us learn from this discussion that all evils, and especially the devastation and ruin of kingdoms, arise from turning away from true religion to false religion. The foolish people of today judge the matter completely differently, and for that very reason are miserably destroyed. Let us learn that fierce nations are restrained and held back by God, so that they should not harm; that God stirs them up to inflict just punishment on the unrepentant. Thus it was with Sennacherib, Shalmaneser, and Nebuchadnezzar, who were called the servants of God and executed his judgments. Therefore, let us fear God and persevere in the true religion.

### Paragraph 635

Moreover, the Saracens, Turks, and Tartarians are most diligently described. First, a wonderful expedition and celerity are commended in them, the principal virtue in wars. He says they are ever ready at every moment to execute the judgments of God. Therefore, he rehearses all parts of time, every hour of the day. And so there is no security from them; you can be in no surety. They are by and by in armor, and come unlooked for; they invade and speed their matters most luckily. He adds that through their most cruel and speedy armies, the third part of men in the world should be slain. Verily, Asia, Africa, and Europe have felt the most cruel slaughters and destructions of the Saracens, Turks, and Tartarians, ever since the time of Mahomet unto our days, about the space of nine hundred and twenty years. And also, the priests of Mahomet are very quick and diligent to allure men into their errors, and they want no lucky success.

### Paragraph 636

The number is also noted in an infinite manner, and the noble says he, of the army of horsemen, is twenty times ten thousand. And Mirias is the number of ten thousand, so that two myriads of myriads should make twenty times a thousand Myriades. And so the old translator has read or translated it: and Erasmus, twenty times ten thousand. The Dutch translation has, many thousand thousands. Laurence Valla in his annotations upon the New Testament does interpret, as has the Dutch translation, thousand thousands. But however it be, certain it is by the conference of other places, that a certain number is put for uncertain, that is to say for exceedingly great: and to be signified, that the horsemen of the Saracens, Turks, and Tartarians should be innumerable. For we read in the seventh of Daniel, thousand thousands served him, and ten thousand millions stood before him. And he speaks of angels (whom he signifies to be innumerable) and of their ministry. So also in the fifth of the Apocalypse, I heard, says he, the voice of many angels, and thousand thousands saying with a loud voice, etc. Certainly, the histories testify that the Saracens came out of Spain into France in number four hundred thousand. Paulus Aemilius in the second book of the acts of Frenchmen, recites that Charles Martel overcame three hundred three score and fifteen thousand Saracens. And Matthias a Michon in the first book, eighth chapter, of Sarmatia in Asia, Tamerlanes, says he, had an army of twelve hundred thousand. Moreover, it is plain that there were never in any age or memory greater armies of horsemen led out of any nation than of Turks, Saracens, and Tartarians. John adds that he heard their noble, either for that he would so confirm that he had said how their power should be greatest, or for that he would partly signify that their victories also were numbered, and should have an end. That in Daniel is most notable, Mene, Tekel Pheres: that is, has numbered, has weighed, has divided. He has numbered, says he, your kingdom, and has brought it to an end.

### Paragraph 637

And at the beginning and also in the times that followed, the matters of Muhammad increased exceedingly. After Muhammad himself, they had in order twenty-five emirs [? rulers], for so they called their kings or princes, who ruled with great power until the year 860. About this time, the fifteenth emir, called Muhammad, sought to drive out and oppress Imbrael, the governor of Babylon, who sent for Muchulet the Turk out of Scythia against the emir. And the Turk dispatched his forces, and drove many of the Saracens out of Asia, and the Turks began to reign in the East. And the Saracens expelled from those regions came into Africa; from thence sailing into Sicily and other islands, they possessed Spain also, and overthrew nations never before conquered, and invaded Italy, spoiling Rome and consuming with fire many fine buildings. Concerning this matter, one may read Volateran in the twelfth book of Geography, in the threefold Arabia. About the year of our Lord 1300, the Turkish emperors began the Ottoman dynasty, who possess at this day a great part of Asia, Africa, and Europe. Baptista Ignatius has written about this in the end of the second book of the Roman emperors, and Paulus Iovius. Many of the Tartarians received the religion of Muhammad, and have most grievously afflicted the world, as Mathias a Michon writes in Sarmatia of Asia. And there is no doubt that the people of Muhammad have been of very great power, and are so still even at this day.

### Paragraph 638

Now the horsemen and the horses on which they ride are also depicted [note 41.16], that is to say, the manners and power of the Mahometans are described. The horsemen had armor, not of iron, but fiery, yellow and brimstone-colored. Therefore, fire, hyacinth, and sulfur served as their breastplate, their armor. For the hyacinth immediately suggests smoke. The hyacinth, in color, resembles smoke next to fire and flame. And the horses had lion heads and serpent tails with heads. The horses breathed out of their mouths fire, smoke, and sulfur. With these plagues, he says, namely fire, smoke, and brimstone, a third part of men were slain. They also hurt me with the serpentine tails. He adds, their power was in their mouth, and the hurt in the tail. These appear to be understood and explained both spiritually and corporally. For the Mahometans, by their wicked doctrine, which is aptly compared to fire destroying, to smoke blinding, and to stinking sulfur, have destroyed countless people. Finally, with a lionish or tyrannical force, they have compelled many people to receive their Alcoran. And moreover, whenever their false prophets seem to flatter—for Isaiah says a false prophet is a tail—they act like serpents and infect people with the most corrupt poison of doctrine. Out of their mouths proceed not only blasphemous laws, but also marvelous praises, great boastings of victories, abominable blasphemies. Where do they say, is your Christian faith? Our religion of Mahomet overcomes all. All your things are miserable. Being vanquished, they serve like bound slaves everywhere. The thing itself declares that our religion is true and yours starkly false. Indeed, the Mahometans reign in a manner everywhere, in victories and riches they are fortunate and noble. That thing makes Christians afraid and causes many to revolt. For what is done among us is manifest to all. The gospel preachers have fought unluckily once or twice and endure great persecutions every hour. The papists overcome and rejoice. Therefore, there are many thousands who say, “How the thing itself speaks, whether religion is better.” Doubtless, this great felicity holds many still in error, who would otherwise be gentle and tractable. Therefore, it is no marvel that Turks or Mahometans prevail greatly with their mouth, since among Christians, victories and the felicity of this world are of such force in manner with all men. And yet they excel with their mouth and boastings, in the thing itself and in truth. For although the Turks are victorious, yet their religion is most false, most wicked, and most absurd.

### Paragraph 639

And physically, how those things may be understood, there is no man who does not see, who knows Turkish histories. The Muslims burn with fire and brimstone, for scarcely is there any other nation which has so wasted the world with fire as this. Whichever way they turn, all things burn with a fierce fire, all is full of smoke. Their princes are lions, and their government is like lions, all tyrannical. They command cruel things, and nothing but bloody words come from their mouths. Therefore, many of them have called themselves the wrath of God and the whip or scourge of God. And truly, this wrath of the Lord follows corrupt doctrine and swerving from the faith. With these three plagues—fire, smoke, and brimstone—one-third of the world is slain and destroyed.

### Paragraph 640

Moreover, the serpent’s tail admonishes chiefly [note 41.17] that they do harm very much. For if the Muslims or Turks have entered into league with Christian princes anywhere, they have not done so without craft and guile. Those who have believed their promises and flattering words, and have sought and received aid from them, have nurtured a serpent in their bosom.

### Paragraph 641

Hereof remain two notable examples. [note 41.18] A discord arising between the emperor of Constantinople and his princes. Whilst Marcus, lord of Bulgaria, joined himself with the princes, or lords of Greece, the emperor was compelled to require aid of Amurathes the first of that name, the third Turkish emperor after Ottoman. And he aided the emperor gently. For he sent into Greece twelve thousand chosen Turks, with whom the emperor being aided, he discomfited and put to flight Marcus himself and the rest of the rebels. But that same amity was the beginning of the destruction of the empire of Constantinople, and of all the calamities of Greece. For when Amurathes understood by the soldiers which returned home that Greece was both a most goodly country, and not strong, by reason of the discord and dissension of princes, he determined to transport thither immediately, under pretense of persecuting the emperor’s enemies. And so began to possess Greece itself, which both his sons and nephews within a hundred years brought wholly into their subjection. In our time arose a discord for the realm of Hungary between Ferdinand, who now is emperor, and John Vayuode, prince of Hungary: which being not able in strength to match Ferdinand, was driven to crave aid of Solyman, emperor of Turks. The Turk was by and by ready with great faith and diligence placing John in his kingdom; howbeit we see that immediately he being extinct, the Turk enjoyed the kingdom of Hungary. Would God therefore that Christian princes would not trust the Turkish navy and warfare. For whilst the Mahometan laughs upon the Christian with a friendly countenance, he intends to put a serpent into his bosom, and to destroy him. And we are also at this day in this sixteenth, as also in the fifteenth fight, in the Papistical and Mahometical corruption, wickedness, and tyranny. The Lord Jesus deliver us from all these evils by his glorious coming unto judgment. Amen, Amen.

## Sermon 42: ¶ What should be done to the reside we of impenitentes, in this mean while feeling none evil, of the Locusts and Horses.

### Paragraph 642

.

### Paragraph 643

And the remnant of the men who were not killed by these plagues did not repent of their actions, so that they would not worship devils and images of gold, silver, bronze, and stone, and of wood—images that cannot see, hear, or go. They also did not repent of their murders, their sorcery, their adultery, or their theft.

### Paragraph 644

It is spoken abundantly about how great calamity shall come upon the world of the locusts and horses under the fifteenth and sixteenth trumpets. And where it is sufficiently known that not all are subject to the locusts and horses, nor to be punished by them, even though some commit things worthy of punishment, one might marvel whether those who are free and exempted from these plagues may safely lead an impenitent life. He foresees and says, “And the remainder of men, who also commit shameful things against God, and yet are not slain with these plagues set forth, may not think to escape unpunished.” For even they shall be punished also by God, most just. For the speech is defective, and therefore to be made up, both by the tenor hereof, and also by the catholic sense of the whole Scripture, which is that all impenitent persons are punished by God, and that so much more grievously, the more carelessly they have abused God’s longsuffering, being nothing moved by any examples of God’s judgments. Yet he says not this by express words. It was enough for him to rehearse the wickedness wherein they were drowned. For from this, every man may gather what is due to such offenders. Arethas, a Greek expounder, expounding this place, says this speech shows an excellence of insensibility, that is, of the wantonness and lasciviousness of those who have spent the time granted them by God to repent, in vanity, that even for the worthiness of their slothfulness they might receive their reward; yea, even before the eyes of the ungodly, the very reward is put in effect. Yet these men, not only by the sight of these terrible things which they had present before their eyes, were made never a whit better, but also worse, and more and more wrapped in sin, have fulfilled their course, and so forth.

### Paragraph 645

Hereof we may gather that it is not sufficient for a godly and blessed life that a man be not a Romanist or a Muslim; but that of every one of us is required a true faith, which may make us to walk in all the commandments of God; and that we should know that all must be grievously punished by God, so many as transgress the law of God, of what religion, condition, age, state, or degree so ever they may be. For God, most just, has no respect of persons. Whoever has sinned without a law, says the Apostle, shall perish without law; and whoever has sinned in the law, by the law shall be judged. Certainly John seems here to bring forth both the tables of the law, and thereby to reprove the sins and wickedness of the ungodly men, of whom he will also that judgment be gathered. The first table sets forth the service of God, commanding to worship one God, not to worship idols, and so forth. The second gives precepts of living, and teaches the love of our neighbor, forbidding murder, adultery, theft, and like mischiefs. John brings forth two sins done against the first table, and three or four committed against the second. Neither is there any doubt but that he comprises under these all like or unlike, more or less offenses against God and his will. Whoever therefore you are, if you offend against the divine law, you shall be punished. If you seem in this world to escape freely and to flit from hence happily, the same may chance unto you that happened to the rich glutton, whose judgment is described in Luke 16. Briefly, he shall be punished, whoever offends God. God knows the manner, whether he shall punish here and in the world to come, or in the world to come only, and grant here a [illegible] voluptuous life.

### Paragraph 646

And we must chiefly observe in this treatise that sinners are not here condemned. For we are all sinners; otherwise no one would be saved. Damned are those who do not repent, those who truly die in their sins without repentance. The Apostle denies that idolaters, adulterers, thieves, covetous persons, extortioners, and so on, will possess the kingdom of God, but he adds: “But such you were, but you have been washed, but you have been sanctified, but you have been justified in the name of our Lord Jesus, and by the Spirit of our God.” And if you doubt whether you may again come into favor with God, if you, having once been enlightened and justified, fall again into sin, learn from the fall and sin of Saint Peter that you may be restored. And as it is written, “Seven times the righteous falls, and rises again,” and so on. Therefore let us learn from this how effective repentance is, and how harmful a lack of repentance is. If you are, or have been, an idolater, you ought not to despair; turn to the Lord and do penance. If you fall again, do not remain in your wickedness. I have spoken more about this in another place. But if you will not return to God, nor leave the evil custom of sin, never look for any grace of God. You shall perish in your sins.

### Paragraph 647

It remains that we declare in few words the forms of sins, set forth here by John, under which, as I said before, he has undoubtedly included similar offenses, so that the same judgment may be had of like things. First, he says, as it were generally, that they have not repented of the works of their hands. For although with this note or mark idolatry is condemned in the prophets, I extend it to all other deeds proceeding from the force of the flesh. For our work is truly sin; and the good work is of the grace of God and of regeneration. And this general thing once set forth, he adds diverse kinds and forms, two against the first table, and four or three against the second.

### Paragraph 648

It is against the first commandment to worship devils. For our very God will have himself alone taken for God, honored and worshipped. And who is so mad, say you, that will worship devils? Verily, there are certain people in the East, who are said to worship devils, for no other end but that they should not hurt them. This is a barbarous and foolish people; why do they not rather worship him who is only able to restrain the devil, that he can not hurt? However, this wickedness stretches far. For they indeed worship the devil, which will seem to worship God. For this matter is not esteemed of the opinion or intent of the worshipper, but of the law maker. For the gentiles would not seem that they sacrificed to devils, but would have taken it most displeasingly, if any should have said that they worshipped the devil. “You are a most vile and most impudent varlet and slanderer,” they would have said, “which dare so reproach the gods and us.” But Paul nevertheless, I say not, says he, that an idol, or that which is offered unto idols, is anything: but this I say, that the things which the gentiles offer up, they offer them to devils, and not to God. For where there is one only God, and he allows only these sacrifices, which are offered to him, he calls strange gods devils, and idol offerings sacrificed to the devil: of this judgment is the thing esteemed, and not of the fond intent of men. King Saul would have offered to God the burnt offering of Samuel: but Samuel told him, that he should commit idolatry and magic, and so forth. This is a hard saying, but yet true. Whereof I have spoken in another place more at large. The worshipping of images of God and of the saints is against the precept of the first table. For all idolatry is prohibited. John here with clarity defines idols, and taunts them also, alluding to the words of the prophet in Psalm 114. The idols of the gentiles are silver and gold, the work of men's hands; a mouth they have and speak not, and so forth. Therefore it appears of the matter that images have nothing of religion. For they are of earth, of gold, brass, stone, timber, and so forth. Again, of the form and shape it appears that images are vain. For the form resembles a most gross shape, and even a lie. For neither God nor the saints were of that shape which the idols represent. And now there is no virtue in them. They see not, they hear not, and so forth. How then do they represent God or the saints? I have spoken of idols elsewhere. They that think how there is a diversity between the idols of Christians and those of the gentiles, let them show that theirs are not of wood, or that those other do see, hear, and so forth.

### Paragraph 649

The sins that follow are against the second table, which commands, “You shall not murder, you shall not commit adultery, you shall not steal.” There are many kinds of murders. For they slay most cruelly those who have no sword, but a venomous tongue. Many kill with corrupt doctrine. There are schemers, parricides, and murderers, etc. Unless these repent, they shall not enter into the kingdom of God. And those who swell with envy and malice are homicides, as Saint John said in his canonical Epistle, 1 John 3:1.

### Paragraph 650

[note 42.9] Poisoning, witchcraft, or sorcery, or enchanting, pertains to murder. Love potions and enchantments were in the time of Saint John most frequented throughout the Roman Empire; at this day those wicked arts are renewed. But they shall be punished by God, so many as apply themselves to the same.

### Paragraph 651

Fornication also has various forms. This includes whoredom, incest, adultery, and anything else more abominable than these. The Gentiles believed that simple fornication, that is, between two unmarried persons, was no sin. But the Apostle defines the contrary in 1 Corinthians 6:15. This pestilent opinion is revived in many today. Nevertheless, it is certain that a fornicator will not enter the kingdom of God. Ephesians 5.

### Paragraph 652

Finally, this is completed with all its parts. [Note 42.11] Whereof I spoke once in the exposition of the ten commandments. The Lord Jesus preserve us from all defilement by sin, and so forth. Amen.

## Sermon 43: ¶ Christ a strong Angel, is set against Antichrist: and is excellently described, shining again in the darkness of the church with consolation.

### Paragraph 653

.

### Paragraph 654

And I saw another mighty angel come down from Heaven, clothed with a cloud, and the rainbow upon his head, and his face was as the sun, and his feet as it were pillars of fire. And he had in his hand a little book open. And he put his right foot upon the sea, and his left foot on the earth. And he cried with a loud voice, as when a lion roars. And when he had cried, seven thunders spoke their voices. And when the seven thunders had spoken their voices, I was about to write. And I heard a voice from Heaven saying unto me: Seal up those things which the seven thunders spoke, and write them not.

### Paragraph 655

Hitherto we have heard many things about the most dangerous conflicts against Christian piety and truth. But nothing has been spoken of the defenders and maintainers of the true religion; rather, the success and wonderful felicity of the wicked has been preached, especially in the fifteenth and sixteenth trumpets, namely, under Papistry and Mahometism. Therefore, it might seem to many that truth itself was not only oppressed and lost, but also that the truth of God’s promises began to faint and decay.[note 43.2] For the godly are oppressed; iniquity triumphs in all places; truth is trampled underfoot, and lawlessness reigns everywhere. Who, therefore, would not think that the things spoken everywhere about the rewards of the good and the punishment of the evil are in vain? Doubtless, the children of Israel doubted the faith of God’s promises when the term of their captivity was extended only to seventy years. What marvel is it, then, if the faithful at this day also, seeing the servitude or bondage of Mahomet and the tyranny of the Pope, or Antichrist, has continued for many years, begin also to be tempted, as also the saints were tempted in old times, witness Asaph, Psalm 73? And truly, there will be those at this day who will say, “Perhaps this world will always be, will never have an end; Papistry will reign forever; the Mahometans will conquer forever; the godly will be miserable forever; therefore it is better to conform to the world.” Our fathers, eight hundred and a thousand years ago, have long looked for your judgment, and thought the last day of the Lord was at hand, but no end appears anywhere, but all things are daily renewed; therefore, the same face of the world will always be, the courses of time duly returning. Who, therefore, knows whether a reward is prepared for the godly, or punishment for the wicked? For far other things chance to holy men than many look for. They look for blessing, life, and felicity, and behold they are overwhelmed with the curses of all men, carried to execution, and seem the most unfortunate of all men. He who cleaves to the Papists, Turks, and enemies of the Gospel, goes through luckily enough, and so forth.

### Paragraph 656

Therefore, as everywhere in this book, joyful things are mixed with sorrowful, so here also, after the most grievous battles of the fifth and sixth trumpets and the most strong temptations, he joins a most joyous gospel for the consolation and confirmation of the faithful, lest they should anywhere doubt of God’s promises or revolt from the true religion to the false. Therefore, against Antichrist, the black angel of the bottomless pit, is set the bright or shining angel of heaven, the Lord Christ. Here is gallantly described and is said to return unto his; the same swears solemnly that there shall be no other time, but that in the seventh trumpet the very mystery of God should be fulfilled.

Moreover, the Lord Christ commands John to eat the open book which Christ held in his hand and to prophesy again. By all these things, for the comfort of all the godly, is signified that Christ shall return into the church, out of which he seemed by his enemy and vicar to be cast out, with great glory and power; neither the hope and expectation of the faithful to be vain, however the last day of judgment be deferred into many ages, and the godly feel great adversity. Finally, that punishment and reward is prepared by God, and that this shall be given to the godly, and that inflicted to the wicked. For to the intent we might hereof be most assured, Christ takes a solemn oath and says it shall come to pass that the catholic and Christian truth shall again come into the field and fight valiantly against the Antichristian and Mahometan doctrine. Hereof, therefore, shall we learn not to despair in the long persecutions of Antichrist and Mahomet. We shall learn also how to fight against Antichrist and how he must be overcome not with warlike but spiritual weapons. These he is not able to match. He whets one sword on another. And hitherto, in deed, in these two last chapters, has been spoken of the wars of heretics and of the ungodly, and of Antichrist, the head of all evil; hereafter shall follow of the contrary fight of the godly and the maintenance of godliness.

### Paragraph 657

Before this is set a description of Christ most elegant, most wholesome, and most full of consolation, declaring his power in the ministry, by the ministers of the word, whom he has clothed with authority from above, and by the weak things of this world, overcomes and defeats the strongest things of this world, which seemed invincible. Christ, king and bishop, animates all his faithful with his spirit and word equally, always and everywhere works many things by his power, so that he is now felt by all to have come again, unto whom he seemed hitherto to have absent himself somewhat long. And I doubt whether in all the canonical books, in the prophecy of Isaiah, after the story of the Gospel, and especially after the Gospel of blessed Saint John, there be any other book which has more and more goodly descriptions of Christ than this book. They are deceived and much abused who suppose a rare gospel to be preached in this book. But let us see the description of Christ by parts.

### Paragraph 658

We have shown in the fifth and sixth verses that Antichrist, the Pope and Mahomet, are powerful. Now, a mighty angel, the Lord Christ himself, stands against them; an angel in deed, but not in nature or dignity. For he took the seed of Abraham, not the nature of an angel, and is much greater than angels, though he is Lord of angels. Which the Apostle shows in the first and second chapters to the Hebrews. But it is the angel of great counsel, namely, the ambassador of God the Father unto us, as Isaiah and Malachi called him, to teach the will of the heavenly Father, and to work our salvation, and now also appointed that from the right hand of the Father, as king and bishop, he should keep and defend his church. This Christ, I say, given to men, is strong, not weak. Strong to overcome and break through all the force of this world, of his enemies, of Antichrist, of the Devil, and of hell gates. No man therefore need to doubt that he may be defended by this strong giant; no man need to despair in any perils or matters, be they ever so desperate, since Christ almighty lives, who is able to help weary matters. He therefore must be called upon by us in all dangers; of him we must look for help patiently and with steadfast faith.

### Paragraph 659

A star has indeed fallen from heaven, but in the time of Antichrist, something has crept out of the bottomless pit, and Christ descends from heaven. The Lord will not come from heaven in a physical body until the final judgment, but he is said to return spiritually whenever he appears to have withdrawn or absented himself. For he is never absent; he is always present with his people, who are everywhere. He is said to be absent when his help is not felt, but people are troubled by adversity and seemingly broken by evil. Christ therefore descends to us whenever he gives us aid and counsel when we are in need. So, in the most dangerous conflicts of Antichrist and Mahomet, and finally of heretics and all other adversaries, Christ himself with his spirit is present with his people, helps them, encourages and strengthens them, and defends them.

### Paragraph 660

Antichrist, as it were clothed and wrapped with the smoke of the bottomless pit, is evil favored. [note 43.6] And Christ covered and as it were decked with a cloud, both shows himself to be God, who can convey him above the air and clouds into his kingdom, and can rain down dew or wholesome showers to water and make fruitful his chosen. A cloud many times in the holy history was a token of God present. A cloud took up Christ as it were girded from the eyes of his disciples. Clouds with Isaiah in the 45th and 60th chapters drop down grace. Finally, we shall be taken up in the clouds with all the chosen to meet the Lord in the air. Antichrist is crowned with a crown, which he obtained for himself by crafty means, by flattery, threatenings, and deceits, and keeps it by wicked practices, bloody weapons, [note 43.7] and all kinds of crafty dealings and ungodliness. But the rainbow is on the head of Christ. For he is king of peace, pacifying and reconciling all things unto God, reconciling also themselves together with the bond of charity. The rainbow is a token of the league and amity of God towards us, as we mentioned before in the matters of Noah. Of Antichrist is woe, desperation, anguish, and a most afflicted and troubled conscience. Christ is a consolation and peace of the conscience, so that no man need to say, “The mountains fall upon us and cover us.” For the faithful, being delivered, cry, “Abba, Father.” The same Lord Christ said also in the Gospel, “In the world you have affliction, but in me peace.” And again, “Come to me all that labor and are burdened, and I will refresh you.” Therefore, the rainbow sits on no man’s head more rightly than upon Christ’s head. For in the head of this angel appeared the rainbow, as though it had been his crown or diadem. And now we understand from whence consolation and pacification is to be looked for in most grievous perils and afflictions of Antichrist. [note 43.8]

### Paragraph 661

The sun, seen through the smoke of the bottomless pit, raised by the angel thereof and darkened most filthily, was a great grief to the world. But the face of this our angel is clear, and clarifies, and shines as bright as the sun, as in times past also he is said to have shone in Matthew 17. The bright face of Christ brings joy and unspeakable gladness to those who behold it, and pacifies the minds. And this is seen by us spiritually and by faith. Paul is the best expounder of this passage in 2 Corinthians 3 and 4. Certainly, full knowledge of Christ, whereby we believe him to be given to us by the Father, that by him all things should be accomplished, and he himself alone should be everything to us, arises as a most joyful light in the minds. For we are enlightened by his light, as is declared in the gospel of John. However, the terrible darkness of errors and calamities may be overcast in the world by the prince of darkness through the pope and Mahomet, and other corrupters of God’s doctrine, yet the light which is in the minds of the faithful through faith in Christ directs, lights, comforts, and preserves.

### Paragraph 662

Antichrist has the tails of scorpions, most venomous; but the feet of the Lord Christ are as pillars of fire. All things of Christ are firm, right, and most clean. In another place also, the truth of the Lord is figured by pillars. And God himself is called a consuming fire, burning up all uncleanness. Moreover, fire pierces, and cannot be easily quenched when it begins to catch and to burn. And who shall restrain the course of the fire of the Holy Spirit? With this is joined another thing, that the Angel set his right foot upon the sea, and his left on the earth. But to set foot is to challenge possession to himself. For as we say in Dutch, to kindle the mind to dwell in any place; so we read of the Hebrew phrase, whatever your foot shall tread upon, shall be yours: that is, whatever he shall win, shall be your own right and possession. And Christ seemed as it were to be cast out of his possession by Antichrist, Mahometans, and the rest of his enemies. For we have heard many times that they have won lands and islands. But this vision instructs us that Christ recovers again, reenters possession, and brings under his subjection such places which seemed to have been taken and lost both by sea and land, throughout all parts of the world. For by the preaching of the gospel, many are recovered who were plucked away from the true religion; that now I need not rehearse that Christ has in all parts of the world his chosen, who never bowed their knees to Baal. In them he has a most strong and most purified kingdom, figured by the fiery pillars. Certainly, the Apostle calls the church a pillar and foundation of truth. 1 Timothy 2. And Christ himself also says that his kingdom in the world is immovable, however it be assailed by Antichrist. The gates of hell shall not prevail against it, and his church.

### Paragraph 663

And just as Antichrist desires to have the book of the gospel shut, and closed, obscure and imperfect,[note 43.11] so the Lord Christ holds the book open in his hand. He opens, and no man shuts. You understand therefore from what it arises that Antichrist, although he be of most great power, cannot at this day shut the gospel book, which he seeks with all his force to do. Of Christ and his Spirit we have the word bright and clear. By the grace of Christ we have the bright preaching of the gospel, contrariwise a dark and an intricate sophistry of Antichrist: of this book shall follow hereafter more plentiful things. Hereunto appertains the worthy invention and godly benefit of Printing, never commended enough. This opens books,[note 43.12] and sends them abroad into the world in spite of all the enemies of God’s truth, and scatters them abroad in every corner of the world. So that they which cannot hear preachers, to them come godly books not without fruit.

### Paragraph 664

The sound of Antichrist’s wings is like the noise of chariots when many horses run into battle; therefore Christ also makes a noise, and cries indeed with a loud voice. John adds a parable. He roars—or loweth—which is as much as to say, he roars. For Erasmus, perhaps, says that the Greeks do not keep the difference between *brysō* and *goggysō* as the Latins do between *rugire* and *mugire*, that is, to roar and to low. He roars as a lion. We have heard before how Christ is called a lion of the tribe of Judah. Therefore, just as when a lion rears, as Amos notes in chapter 3, all are afraid of themselves; so when the Lord Christ cries by his word, all the wicked are assembled. This signifies that the gospel shall be preached against Antichrist constantly and with authority, to the terror of God’s enemies. And without doubt, although the princes, both spiritual and temporal, seem to despise and utterly reject the preaching of the gospel, it is certain by many tokens and conjectures that they are exceedingly afraid of that preaching, which they despise as vile. For they endeavor with all their strength of mind and apply all their counsels to abolish that same preaching. But if they suppose it to be of no force, why are they so afraid of it? Why are they at such great expense? Why can they never be at rest? To this pertains the common prophecy, which tells that while Antichrist reigns, Elijah shall come, who with a sharp, lively, and manly preaching shall refute the triflings of Antichrist. Therefore, the spirit of Elijah and his earnest preaching is that roaring of the lion, which roars out Christ’s truth.

### Paragraph 665

And as soon as that roaring was heard, the thunders also uttered their voices. By which voices are signified the sundry graces of the Holy Spirit, and chiefly the terrible preaching of the truth of the canonical scriptures, as appeared in chapter 4 of this book. For with the gospel in the latter ages, the scripture of the prophets shall be expounded again, which seems as it were to thunder against Antichrist, against sins, and wicked people. Verily James and John, brethren and Apostles of Christ, are called by the Lord Boanerges, that is, the sons of thunder, thunderers, that is to wit, excellently sharp in preaching, and to be feared.

### Paragraph 666

And St. John would by and by have written the voices of the thunders, but he is forbidden so to do; he is commanded only to seal them. For since the holy scripture, through the inspiration of the Holy Spirit, was written and brought forth already by the servants of God, the Prophets and Apostles, what need were there to write and set it forth again? Those things are sufficient for the godly that are written. From these, preachers may take thundering, lightning, and thunderbolts against Antichrist and all sects. And where he is commanded to seal up the things already set forth, it alludes to the last chapter of Daniel; and that sealing is referred as well to the godly as to the ungodly. Sealed letters are doubtless of most authority. St. John therefore, by his sealing, does not now make the Scriptures authentic, but in sealing them declares that they are authentic enough; so to the scriptures no godly man goes about either to add or diminish. Thus I say, the scriptures are sealed to the godly, as to those who are most persuaded that the scriptures are most perfect and authentic, which may most fully suffice in the church for true piety against all ungodliness. Where the wicked will not see this, and seek not for all things of godly religion in the scriptures, neither care greatly for the scriptures: what marvel is it, though the scriptures be sealed to them, that is shut up, which they neither greatly care for, nor understand, nor yet will understand? And on this wise is Christ set against Antichrist, and recovers again his church, discomfits and subdues Antichrist: to whom be malediction forever.

## Sermon 44: ¶ The Lord Christ performs an other, and confirms his elect, that they should not doubt of the faith of Gods promeſſes, &c.

### Paragraph 667

.

### Paragraph 668

And the angel I saw standing upon the sea and upon the earth, lifted up his hand to heaven and swore by him who lives forevermore, who created heaven and the things within it, and the sea and the things within it, that there shall be no more time; but in the days of the voice of the seventh angel, when he shall begin to blow, even the mystery of God shall be finished, as he preached through his servants the prophets.

### Paragraph 669

But while the wicked triumphed, and the enemies of God, Antichrist and Mahomet, overcame with great success, while all good men were oppressed, and deception and lying reigned everywhere, many men will think that there shall never be an end, neither of these evils, nor yet of the world. For the Apostle Saint Peter knew this, saying that in the latter days mockers will come, who will follow their own desires, and will say, “Where is the promise of his coming?” Of whom Malachi also reasons in chapters 3 and 4. But to the intent that the goodness of God might heal the wounds of the godly, and might advance them in the truth against lying and revolt, and establish them in the same, Christ appears and swears solemnly in the sight of all men. This thing must be explained by all circumstances. For it is a matter of great weight, full of comfort, and right wholesome and necessary for all men.

### Paragraph 670

There is no doubt that he alludes to the final chapter of Daniel, wherein the Angel of the Lord swears, confirming with a solemn oath that all the things previously told to the prophet by prophecy will be fulfilled in their times. Therefore, this mighty angel swears now also, indeed, Christ himself, who set his feet on the sea and land. For by the posture and bearing of his body, he demonstrates steadfastness, lest we should doubt anything of his faith and truth; since he is lord of all, he stands moreover on feet, not fleshly, but of fiery pillars. All things therefore of Christ are certain, sure, and immovable. He who rests on him stands surely; he who believes his words shall not be confounded. And it is no new thing that Christ swears. For we read very often in Scripture that God has sworn. We read in the Gospel that the same Lord Christ has most often repeated, “Verily I say unto you,” “Verily, verily I say unto you,” which is an oath of one swearing. When Caiaphas adjured the Lord in judgment, Christ did not conceal and dissemble by holding his peace, but with express words confessed the truth. From this you may learn that the Lord, when he forbade swearing at all, did not intend the sacrament of swearing. Which, where the boorish Anabaptists will not understand, they stir up wonderful trouble, worthy to be put to silence with more severity.

### Paragraph 671

But why, or to what end oaths are made or taken, the Apostle (drawing from the law in 22.2) has declared at length in the 6th chapter to the Hebrews, namely, that those who are wavering and doubtful might be confirmed, assured, and made quiet. If anyone doubts whether you are dealing faithfully with him, God commands that it be affirmed by a sacrament, so that all distrust may be removed. “Apostle,” he says, “swear by him who is greater, and to whom all controversy comes to an end, if it is confirmed by another.” In this consideration, God, intending to show more abundantly to the heirs of promise the unchangeable steadfastness of his counsel, expresses another. Even so, at this present, where divine providence foresaw that under the kingdom of Antichrist the hearts of the faithful would be most grievously tempted, and that many, because of the most prosperous fortune of Antichrist and all the wicked, would be hardened to believe God’s promises, and that many [?], which thing also Daniel in chapter 11 prophesied, would turn to Antichrist: it seemed good to God to confirm his promises by another, and that a solemn confirmation by his Son, so that those who wish to be wise may consider: if an honest and true person confirmed his promise to you by another, you would think it unworthy to doubt of his promises; how much less should it be lawful for you to doubt the promises of the Son of God, and all his words confirmed by a solemn confirmation? Believe the Son of God sworn, believe his Gospel most confirmed, even though the sky should fall and the earth gape wide. God cannot lie, which is truth, and that the eternal truth: which deceives nor is deceived: which is merciful, and loves people so that he adapts himself to their capacity. For even for us and for our infirmity he performs a sacrament, lest he should seem not to satisfy us in all things: and that all occasions of incredulity, and revolt to Antichrist, and to the filthy world might be cut away.

### Paragraph 672

Now we come also to consider the manner or form of the other. Two things are here recited: the manner of the swearer, and the solemn words of the swearer. For he says, how the angel lifted up his hand toward heaven: which indeed is the most ancient rite and holy ceremony of swearers. For we read the same of Abraham in Genesis 14. And in Daniel 12, it is written of an angel: which lifting up to heaven his right hand and his left swore. We verily hold up our right hand. But where we say, that giving of voices we will hold up both our hands: we signify that we will utterly be of that sentence that we hear there propounded. Therefore the holding up of both hands doth signify a most perfect fidelity, and most assured confirmation of the thing sworn. Certainly in the holy scriptures the lifting up of the hand is often put for another. Whereof perhaps we Germans have borrowed, where we say, that is to say, thou shalt confirm me this by another. And in matters most serious and grave we are wont to use some outward ceremony, whereby we may make the words and the thing itself as it were more notable and grave. Whereupon when we pray unto God, we lift up our hands. And verily another is as it were the calling upon the name of God. Whereupon it is commonly accustomed, with great fear to perform oaths. For all men arise, and put off their caps, as they were ready to fall on their knees before the sight of God himself. When bargain or contract is made with words, the right hands are joined together also, in token of fidelity. Therefore when we take a solemn oath, we lift up our hand toward heaven, where we believe that the Lord shows himself glorious to the faithful: from whom we feel that all good things come unto us: from whence we perceive also that vengeance doth fall upon the perjured, and contemners of God. Hitherto therefore Christ applies himself unto us: and after the manner of men, to the end that men may be made the quieter, he lifts up his hands unto heaven.

### Paragraph 673

And the solemn words of the swearer are these: he swears by him who lives forevermore, who made heaven and the things that are therein, and so on. Thus we read of Abraham in Genesis 14. “I lift up my hand to the Lord, God, possessor of heaven and earth.” And in Daniel 12, “He swears by him who lives forevermore.” Also in Jeremiah 4, “And you shall swear, ‘The Lord lives.’” We say so truly as God lives, and again, “So God help me.” And this is a true manner of swearing. God the creator is here most plentifully and most properly expressed, and here are all creatures severally expressed. He alone is the creator, he alone is living forevermore, as he who is life of himself and gives life to all. This creation and vivification are not communicated to others. As also he alone knows the hearts of men; that hereof we may learn to swear by the name of God alone, not to add to him any creatures, which know not the hearts, neither are life of themselves, but are less than he, and less than men, as they that are made for men. Next after God, there is nothing greater than man. Therefore let not man swear by any other than by God. For all the gentiles swear by a greater; if you swear by the saints or by the gods, you swear by men, equal verily, and not greater. God alone is greatest and best. Therefore we must swear by the name of God alone, as the scripture teaches elsewhere, in Deuteronomy 6 and 10, Exodus 23, Joshua 4 and 5, Jeremiah 45, Psalm 65, and elsewhere.

### Paragraph 674

But to say that this is indeed God himself, how does he swear, you say, by him who lives forevermore—that is, by God? He swears, without doubt, by himself, as in many other passages of Scripture. Or else he swears after the dispensation and assumption of the human nature; after which he said, “My Father is greater than I,” which, notwithstanding in his deity, was nevertheless coequal with the Father.

### Paragraph 675

And what I have just explained is the simplest and truest doctrine of oaths and the proper way of swearing. Yet, some understand this doctrine well enough, but for the sake of pleasing people would gladly swear by saints, and therefore ask whether they may join saints to God, especially with the reasoning, “Unless I do this, I will not be considered among the saints.” I answer that they may not, both because we have no explicit manner of swearing that we ought simply to follow in obedience to God’s honor, and also because those who require and prescribe this form would have us swear by the names of saints in heaven, and so acknowledge that we are helped and punished by their virtue and power. If you do this and acknowledge it, there is no doubt that you are gravely transgressing your sincere religion. Certainly, if you confess God before people, he will also confess you before his Father and his angels; if you deny him, he will also deny you, etc. Another [illegible] is like your confession, whereby you confess whom you acknowledge and believe to be your chief happiness, the refuge also from evil and the rewarder of good. If you therefore join saints to God himself, and equate them together, and say, “So help me God and his saints,” you are coupling them with God, granting them to be your gods, which can both help and hurt you. Therefore, take heed what you do. Read Augustine in the 145th Epistle to Publicola.

### Paragraph 676

However, we must also see what the angel swears by these customs and solemn words. For in this one thing consists the whole sum of the matter. The angel in the 12th chapter of Daniel swears that for a time, times, and half a time, and in the winding up, to scatter the hand of the holy people, all these things shall be fully done. So this angel here swears that there shall be no more time, but in the days of the voice of the seventh angel, when he shall begin to blow his trumpet, that the mystery of God shall be fulfilled. But here let no man understand that all time utterly, and everlastingness itself should be abolished and that there should be nothing more after the judgment; but there shall not always be such a time as now is, and such as the wicked enjoy in this world, supposing that the courses of times shall be always, that the world shall continue always, that they shall always flourish in honors and pleasures, oppressing the godly. This shall not be, says he, neither shall there be any more such a time that shall perish and be subject to changeable courses. For about the last judgment shall perish, or be renewed, all these things that shall perish, and life and glory everlasting shall succeed, I mean the time of eternity with all joy most replenished. Therefore he says not simply that there shall be no more time, but adds, in the days of the voice of the seventh angel, that is to wit, at the last judgment, that the mystery of God should be made consummate, perfect, and fully complete. What this secret, or mystery of God is, the apostle expounds and says, 1 Corinthians 15: Behold I tell you a mystery, we shall not all sleep: and the residue which are left, the mystery of God therefore is nothing else than that the end of all corruptible things is at hand, and the happy and everlasting world shall succeed; for that Christ shall then come to judgment: that Antichrist by him shall be abolished, that he with the whole body of the wicked shall be destroyed, the dead raised up again: the wicked to everlasting perdition, the godly to eternal life: for that death, sin and all corruption must be taken away from the godly, and be trodden under foot, and all misery imposed to the wicked, that they may be tormented world without end. And for as much as many times men doubt thereof, as I have said now often, Christ hath sworn that all these things shall assuredly come to pass, and that the godly shall be consummated with all glory, and that the wicked shall be consummated with all kinds of torments. And this is that great mystery of God, for which the very Son of God being incarnate, dead, and raised again from the dead ascended into Heaven, that he might convey us thither to him, having subdued Hell, Satan, Antichrist and all ungodliness. Therefore as in the 6th chapter was said to the martyrs, that they should rest for a little season, till the number of the chosen be fulfilled: so here we hear also, that the misery of God shall at length be fulfilled, &c. This is spoken to this end also that the godly should be of quiet minds, and patiently abide deliverance. If therefore this consummation be deferred, let us abide patiently and constantly, confirmed in Christ, and his Evangelical verity, as also the apostle of Christ, Paul, hath taught us out of the Prophets, in the 10th chapter to the Hebrews.

### Paragraph 677

Moreover, for a further declaration is added, [note 44.10] as God has evangelized, that is, has preached a good and fortunate message, namely by his servants the prophets concerning the end of the world, the last judgment, the everlasting punishment of Antichrist and all the wicked, and the glorifying of the godly, and so forth. Neither did he say these things for a declaration only, but for confirmation also. For by the oracles of the prophets the faithful are comforted, oracles which they have never failed in anything; neither shall they deceive in the end concerning those things which they had prophesied about the last judgment. And again we see how great is the authority of the ancient scripture, and that its use is excellent in the evangelical church; wherein we see both Christ and his Apostles confirm all their sayings with prophetical scriptures, and also illuminate, set forth, and declare. The testimonies of the prophets concerning the last judgment, the reward and punishment of the godly and ungodly, the abolishing of Antichrist, of death, and of all corruption, are in Psalm 110, in Psalms 24, 26, 27, and 46, also in Daniel 7, 11, and 12, in Zechariah 14, in Malachi 3 and 4, and also elsewhere. The Apostle has cited Hosea 1 and Corinthians 15.

### Paragraph 678

Therefore let us lift up our heads, brethren, let us watch and pray, for our redemption draws near. Deliver us, Christ, from all evil. Amen.

## Sermon 45: ¶ S. John devours the book received at the Angles hand, and prophecies again to the gentiles, nations and Kings.

### Paragraph 679

.

### Paragraph 680

And the voice which I heard from Heaven spoke to me again, and said: Go and take the little book, which is open in the hand of the Angel, who stands upon the sea, and upon the earth. And I went to the Angel and said to him, Give me the little book. And he said to me, Take it, and eat it up, and it shall make your belly bitter, but it shall be in your mouth as sweet as honey. And I took the little book out of the hand of the angel, and did eat it up, and it was in my mouth as sweet as honey. And as soon as I had eaten it, my belly was bitter. And he said to me: You must prophesy again concerning many nations, tongues, and peoples, and to many kings.

### Paragraph 681

This is the third comfort, which is contained in this tenth chapter.[note 45.1] For under the person of St. John, it is shown here that the apostolic and evangelical doctrine must be restored in the last times before the judgment against Antichrist and Mahomet. And he might have briefly said, “The apostolic doctrine, as it was preached by John, shall flourish again,” but he preferred to express the same by a beautiful vision, at the last adding a plain and brief exposition of the vision. Which is, “You must preach again, &c.”

### Paragraph 682

And all those who expound these things agreeably, first concerning the person of John, who, under Emperor Nero, returned from exile to Asia after about five years, and again preached the gospel. For he lived until the third or fourth year of the reign of Emperor Trajan. Secondly, concerning all preachers before the final judgment, endowed with the spirit and doctrine of Saint John, and constantly professing Christ against Antichrist. Primasius, expounding this passage, says that the certain meaning is directed to Saint John, who must yet be released from exile, not only to bring this revelation to the knowledge of Christ’s church, but also to preach the gospel more deeply to people and nations, to tongues and many kings. Nevertheless, no one doubts that this voice also agrees to the whole Church, which never ought to cease from preaching, and so forth. Thus he says. The ordinary gloss expounds these words: although this is understood of the very person of Saint John, yet even herein is understood that the Lord will have his church likewise instructed and taught by other preachers also. This pertains to the consolation of the faithful, who shall live in the days of Antichrist and the residue. Thomas of Aquin also: In Saint John himself, says he, other preachers are understood, whom the Lord in the time of Antichrist will have to preach instantly to great and small. So much says Thomas.

### Paragraph 683

Arethas, Bishop of Caesarea, an expositor of this book, recounts concerning this passage from Saint John that the opinion of the common people was that Saint John, along with Enoch and Elijah, would return into the world before the judgment, namely, corporally, and earnestly and constantly to preach against Antichrist. The same Arethas repeats with a more plentiful exposition, where in the eleventh chapter. Where in chapter 44 of Ecclesiastical History it is written that Enoch was translated that he might teach the Gentiles, many have interpreted it as though he should corporally return, that he might teach the Gentiles against Antichrist; whereas by the very translation made in times past he teaches rather the Gentiles that there is another life prepared for the servants of God, and that the same is also due for the bodies, since Enoch was translated both in body and soul: against the opinion of Epicurus and the mad world, supposing no other life to remain after this, and that the bodies do putrefy and never rise again. This Enoch seems to come spiritually to that last age, for that the Lord himself prophesied that a like thing should come unto it, as happened before the deluge or flood of Noah. For just as many then, being careless, scorned the judgments of God, neither feared they any peril, nor hoped for any better life, so it comes to pass also in the last age, in which Enoch constantly preaches through those who establish and maintain eternal life and the resurrection of bodies against the Epicureans.

Elijah appeared in glory with our Savior Christ to three chosen Apostles on Mount Tabor; neither is it to be thought that about the end of the world he must be thrust out of the heavenly palace, and again be subject to corruption, and subjected to the cruel hands of Antichristians, which might tear him in pieces. For just as in the time of our Savior Christ Elijah in virtue and spirit, I mean Saint John the Baptist, went before Christ the Lord, so also before the judgment Elijah shall preach in those who are endowed with the spirit and virtue of Elijah, who shall draw away the minds of all men from the worship of creatures to the adoration of the eternal and only God. Elijah cried out, “How long halt you on both sides?” If the Lord be God, follow him; if Baal be God, follow him. And now the Elijahs shall cry: if Christ is the perfection of the faithful, what need is there of human inventions and constitutions to work a perfection? If Christ is our justification, satisfaction, purification, our only mediator and redeemer, why are these things attributed to human merits? Why are saints accounted intercessors in heaven? Why is salvation ascribed to many other stinking things? Elijah cried out, “How long halt you on both sides?” As though he should say, it is not lawful to part your hearts between two gods, neither is it lawful to attribute all things of life and of salvation but unto God alone. The fellowship of the kingdom is in this case envious indeed. The Elijahs shall cry: if righteousness is of the law, Christ died in vain. No man can serve two masters. Christ shall profit you nothing, which seek salvation in the traditions of men. Come to Christ: he is the perfection of the faithful, and in him we are complete. And just as Elijah grievously accused Ahab, Jezebel, and the Balamites, so shall the Elijahs most sharply inveigh against kings, bishops, idolaters, and Antichristians. Thus I say, Elijah comes again, has come, and shall come before the judgment.

### Paragraph 684

Neither shall John the Apostle prophesy otherwise before the judgment. [note 45.6] He shall not return to the earth in his body from heaven; but preachers endowed with John’s doctrine shall renew all his doctrine, they shall expound those things which he has left to the church, written in his Gospel, in his Epistles, and in the Apocalypse. This book has for a time lain hidden, and has been disdained by good and learned men; yet preaching what is contained and set forth in this book will bring it to light and establish it, as is plainly seen to have happened with many others in our memory. And from all these things we clearly perceive how Antichrist must be opposed and slain, not with carnal armor, but with spiritual means: namely, by sincere doctrine, framed after the example of Enoch, Helie, and John, and taken from the holy scriptures. Of this we shall speak more fully in chapter 11. Briefly, John’s doctrine about the last judgment shall be renewed and become known to the world, despite and against their will. And by John’s doctrine we understand the entire evangelical and apostolic doctrine, in the writing and setting forth of which John also exerted a singular effort among the most excellent.

### Paragraph 685

[note 45.7]And in the meantime, within the same vision, the entire manner of faithful and lawful preachers being matched with Antichrist is depicted, showing what they ought to be and with what qualities they should be equipped. First, John is called by a voice spoken to him from heaven, with a commandment to go. Therefore, God’s vocation is chiefly necessary, lest any man should take upon himself this office with an evil affection. Moses was called, the prophets and Apostles were called: some indeed immediately from God, not of men, nor by men; some of God also, but yet by men and of men. The apostles of Christ were not called of men boasting of lawful succession, from Caiaphas, Annas, and the college of priests; nevertheless, they had their vocation from Christ, and approved their vocation in deed, by preaching the truth. Therefore, although we cannot at this day refer our vocation to the Pope and bishops, boasting of lawful succession, yet forasmuch as we are able to approve it in very deed, and by the testimonies of Christ, that our doctrine is Christ’s doctrine, and therefore that our ministry is lawful, we care not a whit for their opprobrious and railing words, which cry that we are not called, that we are not ordained by the Pope.

### Paragraph 686

But to those who are called, a sure commandment is given: to take the book—not every book, but the book open, and that from the hand of the angel, and again from the angel standing upon the sea and land. That angel is Christ the Lord, Lord of the whole earth, of the sea and all things contained therein. He, with his hand, offers to his ministers a book open—namely, the holy scripture, and especially his sacred and holy gospel, wrapped with no darkness, neither closed, but right manifest to those who will see. For although, because of the antiquity of the language, the propriety of speech, the figures, and the rites, places, things, and stories, some passages may appear somewhat difficult, what does this darken or obscure the mystery of faith and salvation, which is nevertheless most open and plain? Who does not understand what he should believe, what he should do, and how he should pray, even from the articles of the faith, the ten commandments, and the Lord’s Prayer? The sum of faith and of doctrine is certain and plain.

This book, therefore, opened, Christ offers to his ministers. Saint John speaks of a little book, not of a book. For if you compare the holy Bible, especially the gospel book, with other laws, books, and especially the decrees and decretals of the Pope, the little book of the holy gospel will seem very small. Primasius, expounding this passage, seems to understand the truth of the law and prophets manifested in Christ; therefore he says not now, as before, that he takes the sealed, but the open book. For Christ is the end of the law for righteousness to all who believe, and so forth. Therefore, the Lord Christ himself gives to sincere preachers no other preaching than his own, namely, the Evangelical. For he is the light and redeemer of the world, righteousness and life, and there is no other salvation. This preaching is not drawn from or taken from others than from the hands of the angel, not from the hands of the Pope or bishops. Christ says, “Go forth into the whole world and preach the gospel to every creature, teaching them to keep all things which I have commanded you.”

### Paragraph 687

Now also is required obedience of the ministers, that they obey the commandment of God: and that they seek and receive that which they are commanded to ask and receive. It is in vain that some look for a drawing and working of salvation outwardly, and without effort on their part to be finished, through the only invisible operation of God. If God will have me blessed and just, they say, let him work in me what he will. Moreover, they themselves are not careful how they should apply themselves to the grace of God working by grace. It is against their ungodliness that we hear how John applies himself to the commandments of God, not without grace. For he goes to the angel and says, “Give me the book.” For the Lord must be prayed to; we must read diligently, as Paul also commands; we must learn and obey the commandments of God, and not tarry until God without us draws us.

### Paragraph 688

And the Lord denies nothing to those who are willing, do, and are diligent, as he says in the Gospel: “I will give you a mouth and wisdom that your adversaries will not be able to contradict.” Moreover: “My good Father will give his Holy Spirit to those who desire it from him.” Therefore, the angel now says: “Take the book.” Nevertheless, he adds another commandment: “Eat it.” He alludes to the second and third chapter of Ezekiel, where the prophet is likewise commanded by God to eat a book offered to him. For John invents nothing new here. Jerome says that to eat a book is to lay up the understanding of the scriptures in the secret bowels or entrails. He seems by a trope both to intimate an earnest desire and to signify a singular diligence. For we devour with a greedy desire such things as we have long and much coveted to eat. They are said to have devoured books and authors, which they have perfectly learned and can. We say in Dutch, “Er hat den Galen, oder Priſciane gar freſſen,” that is to say, he has learned them perfectly. It is required therefore of the preachers that they learn the holy scriptures with a desire, and that they learn and remember them whole and exactly. Without a desire and fervor of mind, you shall profit little in the study of holy scriptures; and unless you learn the Gospel exactly, you shall preach it unprofitably. The ministers therefore may be ashamed of their ignorance, which is given to idleness, taverns, hunting, dressing [?], and other worse things, than to the study of holy scriptures. They being far unlike the apostle John, shall in this warfare against Antichrist win small renown, unless they do awake out of their profane sleep, and faithfully do their duty, without doubt most holy.

### Paragraph 689

Neither is the effect of the ministry and word preached concealed here. It is sweet in the mouth as honey. For David also sings: “The judgments of the Lord are to be desired above much gold and precious stones, and sweeter than honey or honeycombs.” This sweetness is always felt in the inward man, and the faithful, enlightened by the truth, have always continual joy. But we must not conceal what it seems to the flesh, and what the effect thereof is in the outward man. It truly makes one violently agitated—which is also a phrase of speech, to which others respond, signifying that what is proposed to us is both painful and grievous. The word of God therefore brings the mortification of the flesh, trials, pain, the cross, and innumerable adversities, which with a strong and constant patience we must overcome. For the Lord in the Gospel preached repentance or mortification, and among other things made much mention of persecutions with which his followers would always be exercised. Primasius: “When you have devoured the book,” he says, “you will indeed be delighted with the sweetness of the divine word, and with the hope of salvation promised, and the pleasant taste of God’s righteousness; but without doubt you will feel bitterness when you begin to preach both to the devout and the ungodly.” For the preaching of God’s judgment, once heard, doubtless through the bitterness of repentance, some being converted to better are changed; and others again, being offended, are more hardened, and bear great hatred and malice toward the preachers. The wise man says, “You shall rebuke a wise man, and he will love you; reprove a fool, and he will hate you,” and so forth.

### Paragraph 690

These things are not said only, but also done and felt; for John says, “And when I had devoured it, my belly was made bitter.” And we feel at this day the most grievous hatred from powerful men, namely, of spiritual fathers and temporal princes. Many are driven into exile, innumerable are shut up in prisons, and an infinite multitude are slain with sundry kinds of deaths. All these things the prophets foretold should come to pass, our savior himself in the Gospel gave us warning thereof; the Lord here tells us again the same tale. Therefore let us be strong and steadfast in the Lord, and fight against Antichrist until the end of our life. The Lord will not forsake us, lest we should be overcome by those adversities, having told us of them diligently beforehand. And thus must they be instructed who shall war against Antichrist before the last judgment.

### Paragraph 691

As I stated at the beginning of the sermon, here is a brief explanation of the vision. For the angel says to John, “You must prophesy again to the Gentiles, and so forth.” By this vision, he says, I would declare nothing else but that you must preach again to the world, first by yourself in Asia, after you shall return from exile; secondly by faithful ministers even to the world’s end, who shall spread abroad this doctrine now set forth by you, with sundry tongues through nations, and therewith shall defeat Antichrist. Those accustomed to reading the scriptures know that to prophesy is taken to mean to preach. For prophecy is preaching; they were in times past called prophets, which at this day are preachers, as we may gather from 1 Corinthians 11 and 14. And the doctrine of John is turned into the Syrian, Ethiopian, Egyptian, German, Spanish, French, English, Italian, to be short, in a manner into all languages: in all these, John preaches at this day by faithful ministers. The Gentiles, however barbarous and rude, hear John teaching, and so do the people of many nations. All these receive not a little comfort in these most dangerous days of Antichrist, and have received of them also before this time, which long sins renewed the apostolic doctrine against Antichrist. The same doctrine is brought at this day, and was brought in times past also unto kings and Popes, though they kicked and spurned against it. The thing I speak is not doubtful. For we both hear and see these things even at this day. Histories also report many things hereof. Land and glory be to God. Some copies in the Latin are corrupt, which have *Igitur* for *Iterum*. For John said, “You must prophesy [illegible],” which signifies *Iterum* again, not *Igitur*. For he signifies that he, being dead also, must preach to many nations in sundry tongues, by faithful ministers that shall fight against Antichrist. The Lord assist with his spirit all godly preachers of the Evangelical verity and apostolic doctrine. Amen.

## Sermon 46: ¶ S. John measures the temple, and shows that God hath a care of it: and the quire he excommunicateth.

### Paragraph 692

.

### Paragraph 693

And a reed like a rod was given to me, and I was told, “Rise and measure the temple of God, and the altar, and those who worship there, and the choir within the temple; cast them out and do not measure them, for they are given to the Gentiles, and the holy city they will trample underfoot for forty-two months.”

### Paragraph 694

The Lord continues comforting, and describing the hostile war against Antichrist, showing that the church will not be forsaken in those Antichristian and Turkish difficulties. The enemies will never so quietly enjoy all things, but the church will also have her champions or defenders, who will most valiantly resist Christ's adversaries.

### Paragraph 695

And those things are figurative, which are rehearsed in the beginning of the chapter, and seem to be taken from the fortieth chapter of Ezekiel. As are also those which are spoken of in the seventh chapter concerning the faithful sealed, from the ninth chapter of the same prophet. For he is commanded to measure the temple and to cast out the inner sanctuary; of which he shows the cause. And he does not mean the Temple of Jerusalem, which lay in ruins, nor should it be repaired after the prophecy of Daniel and Christ, but the very church of God, that is, the whole number of the elect. For Paul calls the faithful the Temple of God, living indeed, as also Peter in 1 Peter 2 and 1 Corinthians 3 and 2 Corinthians 6. We have said often that Christ is the only altar in the church, and sacrifice for sin, and priest and intercessor on the right hand of the Father. The worshippers are those who worship God through Christ in spirit and truth, and serve him lawfully or with fear. So many as are such, that is, whoever cleaves to Christ, the only peacemaker of the faithful, and serves God truly by faith, they are the very Temple of God and the true church. John has measured these, so that we should understand how the Lord sets his mind to build up the church, not to destroy it. For those who build measure the foundation, upon which the buildings should be set, as appears in Ezekiel 40. Then the temple was also destroyed by the Chaldeans, as the church is now devastated by the Papists and Turks. But the Lord promises by this measuring that he will repair the ruins of the church through the merit of Christ and faithful worshippers; moreover, he signifies that the faithful in these troubles are numbered (as we heard they were sealed) and sure, whom no hostile power can hurt in all these difficulties. For as the altar, Christ, is undefiled and cannot be polluted or destroyed by any power of the Devil, so are the sheep of Christ known to God and do not perish, as also the same Lord Jesus Christ testifies in John 10 and the apostle in 2 Timothy 2. Briefly, the faithful of Christ are in communion with God and with all his good things, in his care, building, number, and defense. This is a most assured consolation. However, where the Lord in the gospel prophesied that the true faithful would be excommunicated by false teachers, he also foresaw what would happen to the ungodly pastors of the false bishops. He says that they do not belong to the building of God, but are to be excommunicated by God, so that the godly should not fear their censure and cursing. And here is the distinction of two sorts, the first of which is more approved: namely, the hall or sanctuary that is within, cast out; that is to say, declare those who are in this sanctuary to be cast out from God. Indeed, the Antichristians will be within the temple, or the inner parts of the temple, and the chiefest part of the church, so that whoever does not acknowledge them and does not follow them in all things, and conform himself to the church of Rome, is judged to be a heretic. The inner sanctuary in the law was the station of priests, the place wherein they were when they should do sacrifice. And while he says the sanctuary must be cast out, he signifies figuratively that the Antichristian priests shall be thrown out. For the place is set for the thing contained therein. And where he says, cast out, he says it of those whom God has shut out; declare them to be cast out. For God does excommunicate; man pronounces and executes God’s judgment. The latter reading is of this sort: and the sanctuary which is without, cast out. So has the Spanish copy. But how shall you cast out that which was without before? Therefore I like, as I said, the former reading. But we reject this reading neither. For the hall that is without signifies the fellowship not communicating with the only altar, Christ, or with the true church of Christ, such as all this book shows the Popes to be with all their family. Moreover, the Pharisees and priests cast out the one who was born blind, John 9, that is to say, they excommunicated him for the confession of Christ, and the Lord says in John 15: “If anyone does not abide in me, he is cast out as a branch and withers.” Therefore, while John is commanded here to cast out the fellowship of priests, he is verily commanded to declare that those priests were excommunicated, which would be and seem the chief prelates of Christ’s church. He is also forbidden to meet this fellowship. For because God will not edify but destroy them, neither will have them honored among his. For he has rejected them. Who then will hereafter care much for the excommunication of those who are excommunicated? Wicked popes have excommunicated emperors, noblemen, and godly people; and, discharging their subjects of their fidelity, have set them in their princes’ places. The story of Gregory II against Leo Isauricus is known, and of Gregory VII against Henry IV, and of Innocent III against Frederick II, and of other bishops against right good princes. Doubtless, the chief string of the Popish tyranny has been excommunication, which the Lord here loosens.

### Paragraph 696

The Lord does not conceal why he pronounces the priests, or inner choir, excommunicated; for it is given to the heathen. This phrasing is as forceful as saying that within the choir they do not act as priests or faithful ministers, but as gentiles who have taken possession of this place. But gentiles are rightly excluded from the fellowship of God and the church, where the Lord himself said in the Gospel: “If he does not hear the church, let him be to you as a heathen and publican.” Undoubtedly, those who are not in the Temple or church, or are in the inner choir—that is to say, those who are considered among the prelates of the church—yet do not hold to Christ, but are more conformed to the heathen than to Christians, are justly accounted excluded among the gentiles.

### Paragraph 697

And now let us see why he identifies the Antichrist as the pope, with his followers and their images resembling the heathens. Those born of God hear the word of God and glorify it; those not yet born of God, but remain gentiles, not only do not hear God’s word, but also blaspheme it. So these men will not hear God’s word, and seek with all their endeavor how to frighten people away from the scriptures, which are God’s word. They say that they are obscure, doubtful, uncertain, and imperfect. Those who believe and cleave to them they call heretics, and the doctrine taken from them is heresy. Again, those who do not have Christ as their head, and as branches do not grow from the vine, have no communion with Christ and are gentiles. But such is the pope and his adherents, still persecuting Christ and all those who affirm Christ to be the only head of the church, Christ alone to be our righteousness and life, so that all the faithful are made fully complete by Christ. He who thus believes, they pronounce him a heretic.

Moreover, the gentiles worship idols, call upon creatures, suppose God to be honored with corruptible things such as gold, silver, and precious things, dedicated to the Temple and set up to beautify it. But what other thing do they do in the church at this day? You plainly see heathen temples when you see their churches. The life also of the gentiles is shameful and filthy; they are given to voluptuousness, full of surfeiting, addicted to filthy lust, they stink in whoredom, and exceed in gorgeous apparel and pampering of the body. See what things the Apostle writes of the life and conversation of the heathens in the fourth and fifth chapters to the Ephesians, and in the first chapter to the Romans, and in the first to the Corinthians, the fifth and sixth chapters.

Now what the life of the pope is and of his spirituality, the thing itself openly testifies, that even for this cause alone they might and ought to be accounted among the excommunicated, the Apostle himself pronouncing the sentence of excommunication in the place which we have now cited, the first to the Corinthians, the fifth. We may add hereunto their Epicureanism. For if they value any religion, if they have any fear of God in them, why do they sell all things in the church – forgiveness of sins, heaven, Christ, the oblation of Christ, matrimony, ministry – briefly, all things? Why do they call into doubt diverse articles of our belief? What mean these doubtful disputations of the immortality of souls and the resurrection of bodies? Why do they make a mockery of the life everlasting?

### Paragraph 698

Hereunto is added that they trample, yea spurn, the holy city; for this reason they may justly be declared excommunicated. This holy city is not that earthly Jerusalem, but the church of God, of which the holy city was a figure, as Paul explains in chapter 4 to the Galatians. For the earthly Jerusalem, according to the sayings of the prophets, having played its part, lies in ashes never to be restored. The Lord therefore signifies that the holy church of Christ should, through the tyranny of Antichrist and Antichristians, be trampled underfoot. And it signified more that he said to tread upon than if he had said to afflict and persecute. For treading is joined with the greatest contempt of him who is trodden on; and hereby is signified an extreme assault and wonderful cruelty of the enemies, which they practice on those they overcome, and have to use at their pleasure. We read in Daniel of the Romans: The beast had great iron teeth, eating and breaking small, and the rest treading under its feet. For wanton beasts are wont to tread with their feet such things as they cannot devour when they are full. And Solomon in the 27th of the Proverbs says, “A soul that is full treads the honeycomb.” Malachi in chapter 4, speaking of the joy of the godly, says, “They shall leap as calves of the herd, and they shall tread upon the wicked, which shall be as dust under the soles of your feet.” Briefly, John by treading signifies the oppression of the church joined with great tyranny, wantonness, and the exceeding great mockery and gladness of the wicked. And plainly seems to have alluded to these words of the godly prophet: “O God, the heathens have come into your inheritance, your holy temple they have defiled, and made Jerusalem an heap of stones. The dead bodies of your servants they have given to be meat unto the fowls of the air: and the flesh of your saints unto the beasts of the land. Their blood have they shed like water on every side of Jerusalem, and there was no man to bury them: and the rest that follows in the 78th Psalm.” And a little after in this chapter shall follow more things of the persecution of Antichrist. Neither shall these things be obscure, if you compare them with those which are done at this day in the church of Rome against the lovers of Christ’s gospel.

### Paragraph 699

Here is shown a certain time in which the persecution of Antichrist would be cruel against the church, namely the space of two and forty months. Some, in accepting this, torment themselves marvelously. I suppose plainly that a certain time was assigned, and that not without cause, yet notwithstanding an uncertain time to be understood. A certain time therefore is assigned, that we might understand that God has appointed an end to their furies; which, as he himself alone knows, he would signify to his faithful, not in years, but in months only, for consolation. For we suffer more easily that which we perceive shall continue but a few months. This sense has also Arethas, after a sort, touched, writing thus: we suppose that the time of forty-two months expresses a shortening of time about the coming of Antichrist; for the affliction to be executed upon the lovers of God, Christ, very God, says that those days should be abbreviated. And these forty-two months are three years and a half, wherein it shall come to pass that the faithful, and the very tried, shall be trodden and suffer persecution. Thus says he.

### Paragraph 700

Doubtless all commentators, in a manner being truly taught by this passage, have attributed to the kingdom of Antichrist and to his most cruel persecutions no more than three and a half years. For so many years make forty-two months, if one puts twelve months to a year. However, the Scripture and the thing itself speak that the kingdom of Antichrist should be a great deal longer. Whereupon I said that a certain time is assigned by the Apostle, and an uncertain time is understood: that is to say, all that same time that is reckoned from the fatal years 666, whereof is mentioned in the thirteenth chapter of the Apocalypse, until the last judgment. And why I expound a certain time by an uncertain, these are the causes. First, because the same number of months is put here in the thirteenth chapter and is ascribed to the old Roman Empire, so that in their tribulations the saints might understand and comfort themselves that there is an end appointed to their tyranny, which is known of God; and that the saints should no more be sorrowful than if they were constrained to abide their tyranny a few months only. Otherwise, if one should account from the first year of Julius Caesar and bring the course of time until that year wherein Odacer at Rome, all emperors of the west being taken away, was acknowledged as king, you shall not find only three years and a half, but about five hundred and seventeen years. If you shall bring the account from Julius to the empire taken away and given to the pope, you shall find about seven hundred and sixty-seven years. The later cause: for Daniel, the Lord Christ, and the Apostle Jerome agreeably do say that the persecution of Antichrist should last until the judgment. But who shall reckon to us the years and days of the last judgment? And therefore the number certain must be expounded by the uncertain, and must think that all things are numbered and prefixed in the counsel of God, which never neglects his faithful. To him be glory forevermore. Amen.

## Sermon 47: ¶ Of the two prophets fighting manfully against Antichrist, and of their power.

### Paragraph 701

.

### Paragraph 702

And I will give power to my two witnesses, and they shall prophesy one thousand two hundred sixty days clothed in sackcloth. These are two olive trees, and two candelabra standing before the God of the Earth. And if any man will hurt them, fire shall come forth from their mouth and devour their enemies. And if any man will hurt them, this way must he be killed; these have power to shut heaven, that it reign not in the days of their prophesying; and have power over waters to turn them to blood, and to smite the earth with all manner of plagues as often as they will.

### Paragraph 703

These things also pertain to the comfort of the faithful. [note 47.1] For the Lord promises that he will send prophets—that is, preachers—who will maintain and defend the truth of the Gospel and the glory of Christ, confront the Antichrist, and destroy his kingdom, and advance the salvation of the faithful. In the previous chapter, verses 8 and 9, the conflict of the Antichrist and heretics against God and his Christ, and against his church, was described; and now, in a few words, the opposing conflict is set forth, and the army of Christ is mustered.

### Paragraph 704

And he brings forth two prophets, that is, preachers: not that there shall be only two, but to signify that the power of Christ in the world would appear small to worldly men (as I shall tell you anon). Meanwhile, he understands all faithful preachers and pastors of all times who offer themselves to resist Antichrist and heretics. Some interpret these things as referring to Enoch and Elijah, who will appear physically before the judgment. However, Saint Jerome in his letter to Marcella refers that opinion to Jewish fables, signifying that these things must be understood spiritually, as are most things in this book. And, generally, all interpreters agree in interpreting these things as spiritual prophets, not physical ones according to the literal sense. I suppose there are only two prophets mentioned here for two reasons.

### Paragraph 705

First, for that he would allude to the old history or prophecy of Zechariah, which is in chapter 4. It was thought then, also, that the people of Israel, returned from Babylon, would find the rebuilding of the Temple impossible, for they had many and mighty adversaries, and they were weak and few, and their governors, Zerubbabel and Joshua, were contemned. But through the mighty hand of God and his faithful aid, it came to pass that the power of their adversaries vanished away as vain, and they, in defiance of hell gates, built up their Temple right, so the Lord says it shall be in that later age, that the ministers, most contemptuous and very few in number, shall build up Christ his temple and repair it, and shake the most mighty power of Antichrist. Hereunto I suppose belongs that saying of Daniel: and when they shall fall, they were helped with small aid, &c. Secondly, for this cause chiefly he accounts only for two witnesses, for it is written in the Law, “in the mouth of two or three witnesses every word shall stand.” It is judged therefore a full testimony, which shall be confirmed with the agreeable declaration of two. Wherefore, when the Lord says that he will give two prophets, it is as much to say as that he will give so many ministers as shall suffice, which shall both build up his church and also pluck down and restrain the kingdom of Antichrist. Some of the commentators think that by two witnesses are understood two testaments. However, we see that the Lord speaks here of witnesses, not of the thing testified or to be witnessed, which nevertheless we do not separate from the witnesses.

### Paragraph 706

The apostles and those considered apostolic are called witnesses everywhere in the Gospel and in the first chapter of the Acts of the Apostles. Witnesses are appointed in judgment to faithfully declare what they have seen or heard, to invent nothing of their own, and to neither add nor subtract anything from what is testified. Similarly, God places witnesses of God—that is, ministers—in the church of God, and it is required of them that they imagine nothing of their own, nor add to or take away anything from God’s word, but simply declare to the church of God the things they have seen in the narrative of the Gospel and heard from the prophets and apostles. Therefore, those who do not bring the Gospel are false witnesses, unworthy to be called witnesses of God and of Christ; they are rather the Popes’ witnesses, whose decrees and decretals they bring forth and testify to for the unlearned people. Therefore, those two prophets will be witnesses of Christ and will bring witness for Christ from the most true Scriptures.

### Paragraph 707

And the beginning of these things is here referred to God and to his Christ, as the origin of Antichrist is reduced to the devil himself. [note 47.4] “I will give,” says the Lord, “to my two witnesses, and they shall prophesy.” Christ sends preachers, and gives to them also the ability to preach. This is a wonderful comfort. For just as the devil often sends, instructs, and helps his false prophets, so Christ does not leave his church destitute, but gives to these ministers the ability to teach and do things successfully. For in the Gospel also he promised and said: “I will give you a mouth and wisdom, which they shall not resist, as many as are against you.” These things ought to comfort us in the grievous consultations, betrayals, and assaults of the enemies of the Gospel. Christ will not forsake his ministers, so long as they are faithful and depend upon Christ alone.

### Paragraph 708

Now is also declared the time of the preaching of the gospel against Antichrist, truly all that time wherein Antichrist shall tread the Temple and holy city. For a thousand two hundred and sixty days make forty-two months, if you put thirty days to every month. But we heard before that Antichrist should tread the church forty-two months. Again therefore, a certain number is put for an uncertain [?]. And here is signified, and that with a mystery is here defined the time of days, not of months or years. For though the function of the ministry be never so hard and dangerous, yet so shall God comfort and confirm them, that they may appear a few days only, not months or years to suffer persecution, and to travel in this laborious work of the Lord. And where I have said that those numbered days are put for an uncertainty of time, this has moved me, that by and by in the twelfth chapter, the same number of days shall be assigned; for which yet he has set before, for a time, and times, and half a time. Which appears plainly to be taken out of the seventh and twelfth chapters of Daniel. I know that the same is expounded by many for three years and a half; that the time should signify a year, times, two years, and half a time, half a year. But every man may perceive that the thing itself is repugnant to that number of years, if he be at the least any thing seen in histories. In the seventh chapter of Daniel, he says, the other beasts gave over their rule, and spaces of life were granted, for a time, and a time. But who will expound these things of two years only, since it is evident that the Babylonians, Persians, and Macedonians reigned many years? He signifies therefore that those kingdoms should reign so long as God would permit them, and give them power to reign. We say in Dutch, where yet we appoint no time prefixed. In the same chapter of Daniel is put the same phrase of speech, that the saints shall be delivered into the hand of Antichrist, for a time, times, and half a time. And in the twelfth chapter, he says that his prophecy shall be fulfilled in a time, times, and half a time. But who shall believe that within three years and an half all those things should be accomplished, which he declared in the whole work? Why then do they restrain the times of Antichrist to three years and an half, especially his persecution? Why see they not the destruction of Antichrist, and the peace of saints, and the day of judgment, to be the same day? For Daniel says, that the beast should be cast down headlong into Hell, when the seats be furnished. And Paul says, whom he shall destroy with his coming, and who shall show unto us the certain day of judgment? It is known to the Father alone. Let them leave therefore with their computations to strive with the Gospel. It appears therefore that the Lord by that kind of speaking, as it were by a riddle, has defined no time certain; but rather has admonished the godly of long suffering, of patience, and constancy; and has commanded that we should not over curiously search into the instant of this time, but should rather permit it to Christ himself, in another place saying: It belongs not to you to know times, and the moments of times, which the Father has reserved in his own power; but watch, that when the Lord shall come, he may find you watching. Therefore whether so ever the Lord shall defer his judgment a long, short, or mean time, be you constant. So at this present he says, how the ministers of Christ shall preach all that time, wherein Antichrist shall persecute. And truly, if thou read the histories, thou shalt find that the most virtuous and best learned men, have in all ages, now for the space of these seven hundred years and more, constantly resisted the Popes enterprises, their great abominations, and crafty jugglings and seductions of monks and Friars. Of the persecutions that they have suffered, I will speak hereafter.

### Paragraph 709

Furthermore, the appearance of these prophets is also shown, so that the manner of their doctrine may be understood from it. They will not be clothed in soft or precious garments, such as velvet, satin, or damask, or crimson embroidered, but in sackcloth. And sackcloth, as appears in the Prophets, is for a mourning garment and for those who are penitent. Therefore, just as Saint John was plainly clothed and preached repentance, so these also will urge repentance and amendment of life, and persuade people to frugality, while persecuting riot and all intemperance. Certainly, all good and learned men for these past seven hundred years have desired nothing else of the Pope and clergy, and of the people, but repentance and reformation; for which they have received little thanks from them. But what the apparel of the Antichristians is, no one is ignorant at this day. Some of it does not differ much from that of harlots. Consequently, he declares more fully and at length, of what sort they will be, and also their ministry, what the effect and virtue of their preaching will be. And he sets this forth and declares it with various figures taken from the scriptures.

### Paragraph 710

First, he alludes again to the fourth chapter of Zechariah. These are two olive trees, and lamps are nourished with oil; therefore, oil signifies the matter of preaching or of sermons. For candelabras bearing lights are preachers, showing abroad the light of Christ and his gospel throughout the world. And that preaching of light is taken from Scripture, just as the light of a candelabra is nourished with oil. Oil is a sign of the holy of holies. Therefore, John also calls the Holy Spirit “unction.” Certainly, Scripture is the inspiration of the Holy Spirit. Therefore, those preachers shall preach Christ from Scripture. And so, preaching the gospel of Christ through the inspiration of the Holy Spirit, they are said to stand before the sight of God of the earth: that is to say, they are in the protection, the care, and providence of that God by whose providence are governed whatsoever are in heaven or on earth.

For he appears to have alluded to these words of Zechariah: “The eyes of the Lord look over the whole earth,” and these are the two children of oil, which stand before the governor of the whole earth. And these things exceedingly comfort faithful preachers, who see that God has care of them, I mean God the Lord of all.

### Paragraph 711

Again, they are neither olive trees nor lampstands displaying the light of the Gospel, but those of the Antichrist’s party consider them the dregs and filth of men, placing them in the stead of the oil of the Holy Spirit and putting them into the lamp. Neither do they display any light, but darkness and the opinions of the most corrupt men. Against these, John reasons, saying, “These things I have written to you concerning those who deceive you.” And the unction which you have received from him remains in you, and you have no need that anyone teach you; but just as the very unction teaches you all things, so is it true, and no falsehood.

### Paragraph 712

Now also the weapons of these preachers are described, wherewith they may defend their cause and fight against their enemies. If any man will hurt them, fire issues out of their mouth and devours their enemies. This signifies, with pretended malice and against justice to hurt or injure: and first he said to hurt. If therefore any of the champions of Antichrist shall assail those preachers and shall blame their doctrine and ministry, straightways they shall bring forth of the holy scriptures God’s word, and so shall repress and overcome their enemies. For that these things may not be expounded after the letter, that same chiefly proves that by and by we shall hear that those prophets shall be vanquished and put to death by Antichrist: to wit, corporally. Who then can not gather hereof that the victory of preachers is spiritual, that their adversaries vanquished of the verity may live in deed bodily, but through the virtue of the verity they may seem to be ghostly slain? And therefore as it were by an interpretation is added: and if any will injure them, so must he be slain. So I say, by fire verily which goes out of their mouth. And who will say that material and natural fire should come forth of a man’s mouth? And Paul also expounding these things, taking the manner of speaking of Isaiah, reasoning of Christ and of Antichrist: whom he shall kill, says he, with the breath of his mouth. Behold, Paul calls it the breath of the mouth, which John named fire. We read also in the twenty-third chapter of Jeremiah, “Is not my word as fire, and as a mallet breaking the rock?” And again in the fifth chapter, “As much as you speak this word, behold I will make my words in your mouth fire, and this people wood, and it shall consume them.” Of Elias we read in the fourth chapter of Kings, first chapter, that calling down fire from heaven he bodily burnt the king’s servants. Which example where the disciples James and John alleged, the Lord forbad them, that he might admonish them of their function, to wit, that they must fight with long suffering and with the word of the verity. Which the Apostle in another place commands expressly, to wit, in Second Timothy, second chapter. Whereby we are plainly taught that Antichrist must not be vanquished with corporal weapons by the ministers, but with spiritual. For he must be slain with the gospel, with that most sharp sword, and fall down and die in the breasts of men, that he may be utterly contemned and known to be Antichrist. And where many confound the ministry of the word and the power of the magistrate, and for the same cause take the sword out of his hands, commanding that in this case he may not strike heretics and blasphemers, affirming that they ought not otherwise to be punished than by the word: let them learn to discern better betwixt offices, and not to give that liberty to blasphemers and to all manner of seducers, and to such as having been a thousand times convicted of heresy, cease not to infect innumerable and bring them into perdition, unless they be straightly punished by the magistrate. Let every one therefore apply their own office, and herein follow the rule of verity and equity, and then shall things be in better order.

### Paragraph 713

Furthermore, he adds more explicit statements concerning their power and ministry, even alluding here to various passages of Scripture. First, he says they have power to shut heaven so that it does not rain during their prophesying. He alludes to the story of Elijah, which is read in 1 Kings 17. And these things must be spiritually applied to our situation. For just as Elijah, through the power of God, prohibited rain, so the preachers of the gospel will, toward the disobedient, or those who will not hear the word but prefer to be seduced by papal abominations, shut up heaven itself, that is, assuredly testify that it is shut off from God, for as much as through Christ alone, as the only gate by which it is opened, they are restrained; and they will also tell them sharply that the grace of God is denied them, which is only granted by Christ. For the prophets are authorities that rain signifies the grace of God and the fruitful watering sent down from heaven. Therefore, during the time of their prophecy, they will steadfastly testify that those who prefer the Pope’s doctrines to the true bread from heaven are, through their own fault, deprived of that heavenly grace, light, and life. And again, we understand that they have power given to them to open heaven to believers. There is no place here to speak of that now. The things written in the gospel concerning the keys of the kingdom of heaven are more manifest, and chiefly belong to this, than that I should now rehearse them; I have both at other times and before in this same book spoken of them at length.

### Paragraph 714

Secondly, he alludes to the story of Moses, and says [note 47.13] that power is given to these prophets to turn waters into blood, which contradicts nothing with the former point. For the water of godly wisdom is a figure of the grace and relief of the Spirit. Blood signifies offense and punishment. For that saying of the law and of the Apostle is well known: “Your blood be upon your own head.” Therefore, these prophets will testify that God has truly sent his word of salvation to save all believers, but that it will be to the unbelievers, through their own fault, unto condemnation. For those who hear the preaching of God’s word and do not believe it, hear it to their own condemnation. And so the gospel is preached at this day to many without fruit, being corrupted by the popish doctrine, and by force will not be wise, etc.

### Paragraph 715

Finally, they have power to strike the earth with every plague, as often as they will. But they will not, except God’s word, by which they, being inspired and instructed, are governed, shall command them. For they will do nothing wilfully; they will not follow their affections, but the word of God. However, they are said to strike the earth with plagues when, out of God’s word, they threaten that God with plagues will punish the sins of men. Those plagues are recited in chapters 26 and 28 of Deuteronomy. Therefore, if they threaten impenitent persons with war, pestilence, famine, sicknesses, and other calamities, God will send them to those who are incurable, as the Lord himself says often in Jeremiah. Again, and on the contrary, they shall enrich with all blessings those who obey God’s word, when they show forth the Lord’s blessing.

### Paragraph 716

Thus much he has spoken hitherto concerning the preachers of the Gospel, who shall fight against Antichrist in that last age before the judgment, and shall build up the church and confirm the believers. You yourself shall observe in what preachers you perceive these marks, and you shall acknowledge them as the lawful prophets of God. And you shall acknowledge with all how great a benefit of God it is to have true and faithful preachers of God’s word. The Lord our God confirm all ministers of his word in the setting forth of his truth, to the world’s end.

## Sermon 48: ¶ Of the cruel fight of Antichrist against the Prophets of God, whom he overcomes and ſleyeth, and shamefully uses them.

### Paragraph 717

.

### Paragraph 718

And when they have finished their testimony, the beast that came out of the bottomless pit shall make war against them and overcome them, killing them. And their bodies shall lie in the streets of the great city, which spiritually is called Sodom and Egypt, where our Lord was crucified. And some of the people, kindreds, tongues, and nations shall see their bodies for three and a half days and shall not allow their bodies to be put in graves. And those who dwell upon the Earth shall rejoice over them and be glad, and shall send gifts one to another, for these two prophets vexed those who dwell on the Earth.

### Paragraph 719

We have heard of the constant preaching of those who will oppose Antichrist and his army for Christ’s truth and the church of the faithful, and that throughout the time Antichrist exercises tyranny against the church. Consequently, our Lord Jesus Christ will teach us through the Apostle and Evangelist Saint John, concerning the condition in which the saints will fight and how Antichrist will encounter them; which also pertains to consolation and a necessary admonition, lest any man should be discouraged by the prosperity of the Antichristians and the calamities of the faithful. He speaks therefore expressly of the grievous persecution of Antichrist, which has now continued these many years—I mean all that time wherein the Bishop of Rome has usurped authority over all churches—with some small spaces of respite permitted by the Lord. This persecution of Antichrist is more grievous and longer than any that occurred among the ancient people of God or in the primitive church. Certainly, for these fifty years, whoever, of whatever state or condition, began to speak even a little against the church of Rome, immediately felt hatred, imprisonment, banishment, and death. History testifies to this, and shows also that the persecution increased as the bishops themselves and their champions, monks and friars, increased in number and power.

### Paragraph 720

And the Lord declares most diligently when, who, of what status, where, when, and with how great cruelty the Antichrist shall act as a tyrant against the faithful servants of God. He adds immediately that all his endeavors shall be utterly vain; and how great shall be the rewards of steadfast ministers, and also the calamities of the Antichristians.

### Paragraph 721

First, indeed, he plainly admonishes when persecution must arise: not before the testimony of the prophets is finished. I showed you before that the testimony is the sincere preaching of the gospel. Arethas asks: what testimony? That the one who will be present is not Christ, but a deceiver and a pestilent seducer, and so on. And so great is God’s goodness, loving his church, that he will not allow the preachers to be taken away until they have finished their preaching. For the gospel must be openly preached to all men for salvation, and deliverance from anguish, schemes, and deceits, and from the seducers of Antichrist. And they shall finish their ministry with numerous writings and continual preaching. They shall finish, I say, when it pleases God. For some preach and abide safe and sound for many years, being safe and secure from persecutions. And others are immediately apprehended, cast in prison, and slain. Thus are these things done, as it seems good to God; which must ever be credited, whatever means he uses, to advance his glory and further the health of his church. Here comes to pass also, as we read often in the gospel, that the Lord was not taken, because his hour was not come. Therefore, a certain hour also will be appointed by God to the preachers. Before this hour they are safe and secure, though the devil be never so furious, tyrants rage, and bloodsuckers and enemies of faith lie in wait. We marvel sometimes how the preachers of the gospel could preach in so great a company of wolves for so long a time, and directly against wolves. Why were they not immediately torn to pieces?

The Lord God almighty has kept them, so that they might thoroughly finish the testimony of the truth. He allowed their enemies and gave strength to his servants to preach. To him we shall render thanks that many good preachers in times past, and of late days Doctor Luther, and Doctor Zwinglius, and other faithful witnesses of God, could in so wicked a world and in so great a power of Antichrist, execute their ministry for so many years, in defiance of the gates of Hell. Nevertheless, the princes and magistrates deserve also to be praised for the lawful defense shown them; yet this should have been none at all, unless the power of God would have had it so.

### Paragraph 722

And when the faithful in the church have been sufficiently warned, so that those who desire to be wise and are not of a steadfast purpose may all escape the snares of Antichrist and live in Christ, persecution will immediately follow. For as soon as the Pope hears, with his supporters, that he is being attacked, he will immediately begin to threaten and unleash his power, ultimately seeking to incite secular authority against heretics. He explicitly shows who this enemy of these prophets and their preaching will be, namely the beast, which is the Bishop of Rome, recognizable by his cruel, tyrannical, and beastly power. More will be said about the beast in chapters 13 and 17, where we will hear that she arises from the earth, from the bottomless pit, and from the deep pit of hell. For the origin of that wickedness is attributed to no other source than the devil, the prince of hell, a liar and murderer. And the thing itself proclaims this today: all persecutions and conflicts are instigated by the Pope and his bloody ministers of evil. From such authors arose all the calamities of former times.

### Paragraph 723

And he contends against the ministers and ministry of Christ with sophistry, with crafty and subtle practices, excommunication, death, and terrors. Jerome says: Antichrist will inflict various kinds of torments, and those he cannot overcome, he will attempt to vanquish with doctrine. He will give rewards and will promise sweet words, and will also display false miracles, and so forth. And saying it has pleased the Lord to call that seat the beast, why should we call it the holy See? If the Pope is that bloody beast, why should we address him as most holy father?

### Paragraph 724

He shows moreover, with what fortune Antichrist shall contend with the prophets. He shall overcome, says he, and kill them. The same the Lord said plainly in the Gospel, Matthew 10, and John 16. And before also Daniel in chapters 7 and 11. Some things are spoken also before in this book of the holy martyrs. The Lord gives this warning in time, lest if we should see the preachers of the evangelical truth slain, we should doubt of the truth of the preaching, or esteem the matter of religion after the felicity of this world. Which nevertheless many do at this day. For most men say: if this were the preaching of the truth, as it is said to be, the most true God would not forsake his own cause. But now since the preachers are oppressed and destroyed, why should we not gather that their matter is false, and theirs true that overcome? But if we might so reason, then the Prophets, Christ, and the Apostles defended a very evil cause. For all in a manner being oppressed of their enemies, in the end were slain also. Full good then was the quarrel of the Jews, Pharisees, and the most wicked enemies of God’s word? How be it, you will say, since the truth is invincible, how is Antichrist said to overcome? He shall not overcome doubtless, by sure testimonies, by holy Scriptures, or strong reasons: but by force, imprisonment, sword, and fire. For therefore by interpretation immediately follows: and he will kill them. Therefore by killing he shall seem a conqueror. For if in a combat Aeneas shall overcome and slay Turnus, Aeneas shall be called a victor. And hitherto in deed Antichrist overcomes: and although the martyrs be slain, yet do they before God receive the reward of victors: because their cause is just, and the truth overcomes in them. The enemies overcome with the multitude, pomp, authority, power, favor, riches and other like things: we in the goodness, and excellency of the cause, and finally by better testimonies of the Prophets and Apostles.

### Paragraph 725

Now we have the manner of the conflict and victory. He will fight and overcome with physical weapons, and will be subdued with spiritual armor. Added to this is what cruelty he will use against the prophets. He expresses this in two statements: “Their bodies will lie in the streets of the city, and they will not allow their bodies to be placed in graves.”

### Paragraph 726

The first sentence indicates extreme cruelty combined with utter contempt. For all foul things are cast out into the street, even the refuse of all streets is trampled underfoot. Antichrist will therefore treat the prophets with great shame, so that all will believe they have power over them, and will, as it were, spurn them with their feet, considering them outcasts and wicked persons, whose removal will make all things safe. Indeed, in some cities, it is customary to cast the bodies of those who are executed into the streets, so that all may tread upon them and drive carts over them, for the terror of others, and to signify that those executed were most detestable and put to death for serious crimes. And to this pertains that Antichrist, by secular power, will publicly hang up some ministers of the church on the gallows, and fasten others with chains to a post, and then burn them with a slow fire, and finally kill them, binding them so terribly to the pale in chains that he may hoist them up and let them down into the fire again, singeing them and lifting them up again, for the terror of all who look on. What will they say about him, judging them unworthy of the last honor? Burial is the last honor done to a person, but he will not allow the bodies of the faithful to be buried. Thus, perhaps, he explains what he said before, that their bodies will lie in the streets. At this day, not only burial is denied to those who suffer for the gospel, but also the bones of the dead are dug up and burned, [?those] who, while alive, would not receive the papal sacraments. For if anyone dies without having whispered a confession to a priest, confessing all his sins to him, without asking for absolution or receiving the Lord’s bread, or allowing himself to be anointed with extreme unction, although he died in true faith, because he has not used those papal ceremonies and subjected himself to the Pope, he will not receive Christian burial, but will be buried on a dung heap with dogs. The practice itself demonstrates this today. Moreover, these Antichrists will seek by this means to abolish all memory of the godly. For monuments are made to preserve the honorable memory of the dead. But the righteous will be in eternal memory, and so on. And they indeed think that they are acting like good Catholics, but the Lord Jesus exposes their work, saying that it is extreme cruelty. Then what shall you think of those who, blinded by hatred of the true religion, like wolves and ravens, swoop down on the bodies of dead martyrs, tear them apart, and treat them with great shame?

### Paragraph 727

But that cruelty is harsh and bloody, and it shall continue for three and a half days [note 48.12], a period that all commentators understand to be brief, certainly so, but yet uncertain, as I mentioned previously regarding months and years. Therefore, I suppose this shortness of time is brought about for consolation. We also say that the Lord grants the afflicted spaces to breathe, and shortens the sorrowful days, so that we might be able to endure it. If, therefore, our patience is tested in a severe persecution by Antichrist, let us think that our Lord God has accounted for all the days of our calamity, and that he has shortened them for the consolation of the weak.

### Paragraph 728

And the very place where this cruelty must be wrought against the prophets, he expresses plainly, and as it were points out with the finger. To wit, the great city. And it is the city of God, and is also the city of the devil: it is the city of Abel, the innocent, it is also the city of Cain, the parricide: it is the catholic city of saints, it is also the synagogue of Antichrist. These cities are open throughout the whole world, and are enclosed with no straight walls. You might call this city the lordship, dominion, kingdom, or empire, or fellowship of the wicked. Wherever therefore Antichrist or the Pope of Rome has jurisdiction, and even in the Roman church itself, throughout all nations and people, these things which we have heard shall be done against martyrs. For setting forth that city with more plain tokens, their bodies, he says, shall lie in the streets of the great city. And by explanation he adds: which is called spiritually Sodom and Egypt. Moreover, where also our Lord was crucified. And again: and they shall see of people, and kindred, and tongues and nations; therefore by this he understood not any straight, nor yet any large city enclosed with walls, but that city stretches throughout the world wherein dwell nations, kindred, people, and so on. Sodom and Egypt are far apart, neither can they be joined together with any walls. Again, our Lord was crucified in the city of Jerusalem, which is also called by the Prophets Sodom and Gomorrah, but he is crucified daily in his members throughout the world. And there is one and the same city and society of all the wicked in the world, as there is one body of the godly. Let us know therefore how that city wherein the bodies of the prophets lie in the streets, is the city of Cain, and the Roman church scattered over the world. The same is called Egypt and Sodom, but spiritually.

Where we see this vocable [illegible] spiritually, to be used in a sense far from the letter, for otherwise there is no spirit at all, either in Sodom, or in Egypt: for they are altogether flesh. Therefore, in sense of a parable, and by comparison, this city is called Sodom and Egypt. What Sodom was, appears in the 19th chapter of Genesis and the 16th of Ezekiel. Her sins went up to heaven. But what manner one is at this day, and has been for a long time, the church of Rome, all men know, except it be they that will not know. And the Apostle in 1 to the Romans has expounded. And Egypt robbed the children of God of their liberty, oppressed them with vile bondage, and prohibited them from the true worship of God. So likewise the Roman church has spoiled the church of Christ of the liberty gotten by Christ, has wrapped her in filthy servitude, that she might serve in the dirt of men's traditions. It prohibits moreover by all force and power, that she shall not, in returning to the gospel, serve God truly. Verily, our Lord Jesus Christ was crucified on Mount Calvary, at the city of Jerusalem; and also the articles of our faith expounding the same, say, that he suffered under Pontius Pilate. He was the Roman governor, it is manifest therefore, that Christ suffered under the Roman Empire. Under the same Empire, and under the judgment thereof, were executed the Apostles and ancient martyrs. Under the Empire of new Rome fall those prophets also at this day by sword and fire. Also, people, kindred, tongues and nations obey this Empire, now called the church, so that preachers in all places, the beast so willing and commanding, are apprehended, and slain with cruel deaths, people, kindred and nations looking on.

### Paragraph 729

To amplify and most aptly express the cruelty, this pertains chiefly [note 48.17], that these earthly men (possessing in these lands a church entirely carnal) shall rejoice and be glad over the calamities and miserable deaths of prophets. The Lord himself prophesied this before in John 16. Verily, verily, I say unto you, you shall weep and lament, but the world shall rejoice, and so forth. Yea, they shall send gifts, he says, and letters of rejoicing. That this was done in the council of Constance, when John Hus [note 48.18] and Jerome of Prague were burned, histories make mention. We have heard very lately, how after England fell again to the Roman religion, what joy and gladness, what banquets and triumphs the Papists made in all places. So often as the ministers, or other faithful are burned, they keep a festive air and pleasant feasting, singing “Te deum Laudamus.” And letters of rejoicing fly to and fro; in some other places, with solemn procession, they rejoice at the misery of the faithful, but the Lord sees these things, which foretold sins that the same things should now come to pass.

### Paragraph 730

And the cause of this great joy is no other than that those prophets have troubled those who live on earth. For those who love the earth and desire earthly things are greatly offended by the free preaching of the truth, which they hate more than dogs and snakes. For they either seek to attain honor, riches, and pleasure, or, if they have them, they desire to keep them. But they are greatly afraid that through preaching, these things will be severely shaken or entirely taken from them. Therefore, they desire nothing more than to be rid of and delivered from their clamors, and to have them removed immediately. For so they think they will be safe and enjoy their pleasures at will. With like affection and counsel as in Herod’s feast, they coveted not a kingdom nor a great sum of gold, but the head of John the Baptist. The Popes would rather at this day have the heads of certain ministers of the church than so many thousand crowns. Yes, moreover, ministers of churches are called plagues, disturbers, seditious, and injurious against God and his saints, and against all men. Therefore, they wish with all their hearts to be eased of this burden. The Lord Jesus forgive them this sin.

## Sermon 49: ¶ The enterprises of Antichrist in wedyng out the preachers to be vain: howe great shall be the rewards of Preachers, and of the punniſſhement of the wicked.

### Paragraph 731

.

### Paragraph 732

And after three days and a half, the spirit of life from God entered into them. They stood upon their feet, and great fear came upon those who saw them. They heard a great voice from heaven saying to them, “Come up here.” They ascended into heaven in a cloud, and their enemies saw them. In that same hour, there was a great earthquake, and a tenth part of the city fell. In the earthquake, seven thousand names were slain, and the remnant were terrified and gave glory to God of heaven.

### Paragraph 733

Hitherto he has spoken of the wicked joys and gladness of Antichrist and the ungodly men of the last age,[note 49.1], concerning the slaughter of the holy prophets of God. They will think how they shall reign forever in those their errors, superstitions, and pleasures; and suppose by their murdering to have silenced the preaching of the gospel, which is most unpleasant to them. But consequently, the Lord shows that their hope is most vain, their attempts to be frustrated, and their joys short; yea, and quickly to be turned into mourning and misery. For first, he declares that the prophecy or preaching shall be repaired by God through new prophets, and that to the greatest grief and terror of the Antichristians, who looked for no such thing. After this, he shows how great rewards are prepared and given to the preachers oppressed in this world and treated with great villainy. Finally, he signifies that the wicked shall not live in continual pleasure, but that God will disturb their joys, bringing misery upon them even in this world. Which, although he begins to punish at the last in this world, will in another world more abundantly augment their everlasting torments. And all these things shall need no great exposition, so that we mark diligently what things have been done a few ages past, and what are done also at this day. And all these things appertain to the consolation and comfort of the Saints.

### Paragraph 734

First that the free preaching of God’s word against Antichrist shall be restored, which seemed to him self to have overcome and oppressed all prophecy, he declares by these words: and after three days and an half, the spirit of life from God entered into them. He signifies by that number of days, as I told you before, a very short time, as though he should say: they shall not long enjoy their false and bloody pleasures. For God shall raise up other prophets in the place of those that are dead. And he speaks as though God should raise up the self same prophets, which Antichrist had slain, and that he would object them again to the wicked in their own bodies. How be it they shall be raised again in their bodies at the last day: but now shall other preachers succeed in the place of those that rest, unto whom God shall give that spirit of his, which he had given to the others that are dead. Therefore he calls this the spirit of life, for as much as those which were slain for the same doctrine, seem as it were to have lived again. Verily for likeness of doctrine, John the Baptist, Elijah, and the prophet Jeremiah seemed to have been revived in Christ, as is read in Matthew 14:12 and 16:14. And here is expressly said, that the same spirit did not proceed of the Devil, or of men (as it is said at this day of many), but of God. For he with his spirit (which is one) inspires his ministers, and directs the same by his word, that the latter wholly answer to the former in doctrine, and severe rebuking of sins, and so forth. For the lively effect of that spirit follows, and they stood upon their feet: that is to say, they lived again. Their doctrine seemed overthrown and trodden under foot, but God’s word stands again upon its feet, and runs most swiftly. We say in Dutch of such as be restored, to expound the effect that same also pertains, that the Antichristians saying other preachers succeed in the room of the that were slain: being stricken with fear, know not whether to turn them. By the way therefore is signified, that the course of the word shall be fortunate, and the which these men can not stop by any means, however they rage and murder. All these things shall be better understood by the Histories of later times, and of such things as are done yet at this day. And to the intent, that omitting the eldest things, I may touch those of latter time: the Bishops of Rome had thought they had won the field in the council of Constance, when they had burned John Huss and Jerome of Prague: but within a short time after many godly and well learned men sprang up in Bohemia and in other countries, in whom those slain appeared to have taken again the spirit of life. In Italy Laurence Valla taught to his great praise, and also Jerome Savonarola, and so forth. In Germany taught many godly men, as in France also in Englande, and other nations. Thirty years passed through the grace of God was brought a light into the world by Mirandola, Reuchlin, Erasmus, Luther, Zwinglius, Oecolampadius, Melanchthon, and innumerable others, in whom the spirit of life uttering itself after every man’s talent, set forth the Scriptures, detected the Romish wickedness, and rebuked the vices of all states, but especially of the clergy. The Romans are afraid of this spirit, and fill the ears of the emperor and kings with complaints and accusations, and cry out that we should all with our books be destroyed and burned. How be it the power of God nevertheless makes the prophets to stand on their feet, and their preaching to run a pace: however these rage in their fury, and persecute God’s verity preached throughout the whole world. To God be the praise and glory.

### Paragraph 735

In this consolation are mixed also rewards prepared for faithful ministers,[note 49.3] whom the Antichristians, after slaying, first excommunicate, so that they may seem bound and delivered to the Devil, to be tormented with everlasting punishment. And until now, all preachers who have spoken against the church of Rome and have suffered therefore at the Pope’s command, have been thought to have perished both in body and soul: their bodies, I say, consumed with fire, and their souls cast down into hell. For they were condemned as heretics and enemies of God and the church, even as plagues of mankind, and so taken out of this life. But contrariwise, the Lord here pronounces and declares everlasting rewards to be prepared for them. For their souls, released from their bodies, are straightway taken up into heaven, and their bodies raised at the last judgment, ascend into heaven also, that there they may rejoice with Christ forevermore. But to the intent that this godly promise of the everlasting and inestimable reward might be of more authority and credit with all men, the Lord propounds it not simply, but most gloriously decked and adorned, for he sets before that a voice was sent to the prophets, and that from heaven; moreover, great and loud. For great is the agreement of Patriarchs, Prophets, and Apostles with the very Son of God in assured doctrine, upon which we believe without doubt that those who suffer for the confession of Christ are saved both in body and soul. And that doctrine was brought from heaven, so that there is no place left for doubtfulness. There are testimonies in the scriptures, both manifest and many, as in Isaiah 26, Daniel 12, Matthew 10 and 16, John 14, and diverse others. What should we say is now brought as an express testimony hereof? For a voice sounds from heaven over the afflicted by the tyranny of Antichrist:[note 49.4] “Come up here.” That is to say, as much as, “I see the lewdness and cruelty of the Antichristians to be such that there is no place left for you on earth. They torment and persecute you as plagues, and you are unworthy to live on earth; come therefore here to me, into the heavenly palace, whither I myself came also after the cross and ignominious death.” We read in the Gospel that the judge shall say to the righteous, “Come, you blessed of my father, &c.”

### Paragraph 736

Furthermore, lest any man should think these words to be vain, the Lord adds through John, “and they ascended into heaven,” not because the resurrection has already occurred, but for the certain knowledge of the thing he speaks of, he speaks of the future as if it were past. Many similar phrases are found throughout the Prophets. Elijah in times past ascended into heaven, both soul and body, as we read in the fourth book of Kings, chapter two, by the same miracle he showed them also what reward the Lord has prepared for faithful preachers of God’s word, and nothing else is repeated here. He adds how they went up in a cloud. For a cloud took up Christ, our head, from the eyes of the disciples, and we also shall be taken up in a cloud to meet the Lord in the air, as the scripture states in the first book of Acts, and in first Thessalonians, chapter four. Although therefore preachers, and those who believe the preachers, are excommunicated by Antichrist, and through open and shameful punishments would seem to be sent to the Devil, Christ receives them, delivered from all evils, unto himself, into the palace of Heaven.

### Paragraph 737

He adds another point as well: their enemies saw them. [note 49.6] They saw, I say, with a terrifying fear, for as they witness those whom they condemned as God’s enemies being in glory as true and honorable friends of God, they will conclude that they themselves are destined for fellowship with the Devil. Read a full commentary on this in chapters 3 and 5 of the book of Wisdom. Although, therefore, preachers of the Gospel in this present world are judged and appear to the world as damned, in that day when all people will be assembled—all who ever have been, are now, or will be—it will be manifest to all that these are God’s dearest friends, and that their cause is just. With this, the Lord comforts those who are persecuted, condemned, despised, and scorned for preaching God’s word. By these things, he prepares and establishes the minds of the faithful, so that they are not discouraged by the rebukes, revilings, and oppressions of Antichrist and his followers.

### Paragraph 738

Finally, the Lord also adds certain things concerning the miseries of the Antichristians, [note 49.7] with which the righteous Lord begins to punish them and to interrupt their wicked joys, so that at last in another world he may put the same to torment that never shall have end. In that same hour, that is, doubtless in the time wherein they shall afflict the prophets, there shall be a great earthquake, and the tenth part of the city shall fall. And the tenth part we understand to be great, yet so that the greater part shall remain in error. As Saint Peter prophesied should come to pass, in 2 Peter 2, and the Lord himself also in Matthew 7. He seems to recite two evils which hang over them: calamities and revolts. For Saint John himself seems to add an exposition, saying: and there were slain in the earthquake the names of seven thousand men. And the remainder were afraid and gave glory to God of Heaven.

### Paragraph 739

Therefore I suppose the earthquake to signify exceedingly great alterations, commotions, seditions, wars, slaughters, and destructions. And he said the names of men after the Hebrew phrase, for a number of men. And he put 7,000, a certain number, for an uncertain one: as where it is said to Elijah, “I have left me seven thousand men, who have not bowed their knees to Baal.” For if it signifies a great multitude, likewise he signifies here that no small number of Antichristians shall be dispatched out of the way by slaughter and sundry or all kinds of calamities. Again, he signifies that the tenth part of the world, that is to say, the adherents and favorers of the Roman church, will revolt, not a few of them, from the same church, being frightened by the preaching of God’s word, and by plagues inflicted upon the enemies of God’s word, and so they shall forsake the Roman church, that they may give all glory to the God of Heaven.

### Paragraph 740

Having been misled up to this point by Roman trivialities and deceptive opinions, they have not given all glory wholly to the true God, creator of heaven and earth, and the inhabitant and ruler of heaven, while attributing more to creatures, human inventions, and errors than to truth. They have shared the glory that belongs to God alone with saints and the works of their hands. But now, being instructed by the preaching of the gospel, they will depend on God alone and will ascribe all glory to him through Christ.

### Paragraph 741

Now if you consult histories, not ancient (for I would not trouble you with a lengthy recounting), but recently written, within these last hundred years, you will find a remarkable clarity regarding this matter. When the preachers of Bohemia were burned at Constance, a great disturbance of the people immediately arose, the Bohemians waging a deadly war against the Romans. Aeneas Silvius himself wrote of that war, in which many thousands of men were slain, and many places destroyed and laid waste. Moreover, countless people deserted the authority of Rome. In our memory, a grievous persecution was stirred up against the faithful through the instigation of Rome, and thousands of faithful were slain, besides the expectation of all people [note 49.10]. Rome was taken in the year of our Lord 1527, and so devastated and plundered that this calamity could be compared with the most severe that ever occurred. Neither do the princes cease to wage war among themselves, and to weaken themselves with mutual destruction, which consistently shed the blood of the faithful. But we are glad and rejoice that a wonderful number at this day revolt from that Roman see, and give to God through Christ all glory. To him be glory and rule forever and ever. Amen.

## Sermon 50: ¶ The seventh Angel blows the trompe, and the elders singe a song of praise.

### Paragraph 742

.

### Paragraph 743

The second woe is past, and behold, the third woe will come shortly. And the seventh angel blew, and there were great voices in heaven saying: “The kingdoms of this world have become our lords and Christ’s, and he will reign forevermore.” And the twenty-four elders who sit before God on their thrones fell on their faces and worshiped God, saying: “We give thanks to you, Lord God Almighty, who are, who were, and who are to come; because you have taken hold of your great power and have reigned.”

### Paragraph 744

By the seven angelic trumpets, not only the destinies of the church are shown, but all the faithful are also spurred to watchfulness and to wage spiritual warfare.

### Paragraph 745

And to the last three trumpets, as they are the most dangerous, are joined three woes [note 50.2], signifying, as I said at the end of the eighth chapter, that all kinds of troubles and the most grievous afflictions will occur in these times, whereby people will be brought into the greatest distress. And the first, indeed, he has separated from the second and third, by these words: “One woe is past, and lo, two woes are yet to come after this.” This manner of speech does not interrupt the matter, but arranges the speech in order. For the papist woe does not cease when the Turkish woe comes on, but afflicts the churches together. That manner of speaking is given to establish the order; so now he distinguishes the third woe from the second, signifying indeed that Muhammad’s law will endure until the last judgment; and yet in the meantime, he does not deny that papistry will continue also, of which he has hitherto, in the eleventh chapter, discoursed many things, having finished the matters of Muhammad in the ninth chapter. Therefore, the sense of the Apostle’s words seems to be this: “You have heard of the first and second woe, hear furthermore also of the third and last woe.”

### Paragraph 746

And we must note—which greatly contributes to the comfort of the faithful—that the Apostle states plainly that the first and second woes have passed. Thus, he indicates that these two greatest tyrants will have an end, and that God has even appointed them specific limits and boundaries that they cannot exceed. Let us therefore rejoice, that God is watchful over us, and will not abandon us, nor permit the wicked to do more than is fitting.

### Paragraph 747

The third woe will affect not the godly, but the wicked, at the time they are oppressed by the final judgment and are led astray from their expectations by the Devil to eternal torments. No language, however eloquent, can express the unspeakable pains of this third woe. Therefore, Daniel also says in chapter 12: “The time will be such as has not been since the beginning of humankind.” But when this woe will occur is not expressed nor determined, as neither is the day of judgment, which is known only to the Father and therefore must not be investigated by us with excessive curiosity. It is sufficient for us that it will come shortly. For the Lord says in the gospel that he will shorten those difficult times for the sake of the elect. And again, when these things begin to happen, look up and lift up your heads, for your redemption is drawing near. But these things do not begin now; they are already accomplished. Therefore, it is certain that our redemption is at hand. So cast aside the worry and care with which many torment themselves, thinking that God delays too long, that he allows much to the wicked, that the godly are sorely vexed and, as it were, forsaken and neglected. For the truth says: “Behold, the third woe will come quickly, in due time.” For in chapter 10, he affirmed by a solemn oath that he will come for judgment. Now, as for the precise moment and opportunity of time, give glory to God and acknowledge him in the courses of time and in all things and creatures, using an opportunity most exquisitely. Therefore, when you confess in your creed, “I believe that the Lord will come from the right hand of the Father to judge the living and the dead,” confess also that he will come in due time. And just as from the beginning of the world, he has never forsaken or neglected those who served him, so he will not neglect them at the end of the world.

### Paragraph 748

For it follows that I may expound the things that go before. And the seventh Angel blew. For he declares that the judge is now at hand, raising from the dead the godly and ungodly: the godly unto joy, the wicked to everlasting pain. These shall be new battles, but to the wicked unfortunate, and altogether miserable. Of the trumpet of this Angel, you read in the gospel of Saint Matthew in chapter 24, and in Saint Paul in chapter 4 of the first to the Thessalonians. He should now adjoin the whole manner and discourse of that last judgment, but he will defer it to chapters 19 and 20. In the meantime, he will recite, as he has promised, the furies of Satan against the church, and how he will use those notable instruments, the old and new Roman Empire, to commit murder and, in manner, to destroy the church; wherein, nevertheless, the wicked shall in this world also be put to most grievous punishments. Now omitting, or rather reserving these things to their own place, he celebrates the gratulations, rejoicings, and praises of the Saints.

### Paragraph 749

The pride and arrogance of the wicked, and especially of the Antichristians, have seemed hitherto intolerable in the world. They have oppressed the godly, boasted of their victories, and bragged of their own felicity. As we shall hear in chapter 18 of this book, that beast has said, “I sit as queen, and am no widow, and shall never see any sorrow.” For voices are heard from Rome: all empires are ours. It is known what manner of things Augustine Steuchus, an Italian and chief champion of the Pope’s authority, has set forth in this cause against Lawrence Valla, about the donation of Constantine. And daily are heard the brags and rejoicings of the papists, concerning the everlasting continuance of the See of Rome, her victories, and the oppression of the preaching of the Gospel, and that she has her power stretched throughout the world, etc.

But in that day, when our Lord Jesus Christ shall abolish all power, rule, and authority, and shall have made all his enemies his footstool, according to the scripture in Psalm 110 and 1 Corinthians 15, there shall be heard again the voices of the glad and joyful, singing true and eternal triumphing songs in heaven. For angels and saints shall sing together; wherefore the voices shall be greater and more durable than the voices of Christ’s enemies, which last but a small season.

### Paragraph 750

Now he also recounts the triumphant song and rejoicing: the kingdoms of this world are made our lords, and Christ’s, and he shall reign forevermore. He shows two things: that all kingdoms are made the fathers’ and the sons’ and that he shall reign forevermore. All kingdoms were before also our Lord Jesus Christ’s; but the same did not appear so plainly to all men, when the bishop of Rome also usurped the same for himself and oppressed those who celebrated only the name of Christ. But when it shall truly appear, and to all flesh, that all kingdoms were ever, and yet remain of one and the eternal God, Christ therefore overcomes, the truth overcomes, the gospel overcomes, the church overcomes: those who are vanquished shall be led to hell, Mahomet with his followers, and the Bishop of Rome with his.

It is added that Christ shall reign forevermore. Antichrist in deed has reigned, and the wicked have rejoiced in this world for a very short time; but now the godly shall reign with Christ forevermore. He does not now divide the kingdom of the Father and the Son, but shows it to be common, where he says that the kingdoms are made: that is to say, it is openly declared that all kingdoms are of God the Father and the Son, and that he shall reign with his elect forevermore. So you may see that the place of Saint Paul may not be explained according to the letter, which is written in 1 Corinthians, chapter 15, concerning the Son being subjected and delivering the kingdom to his Father. For he shall deliver the kingdom, that is to say, shall bring and present it to the Father, and in his members shall be subject to the Father, with whom nevertheless he himself shall reign forever.

The affirmative word “Amen” is annexed, lest any man should doubt one whit of these celestial mysteries. However, he explains afterward more plainly what those voices are that were spoken in Heaven, while he annexes the narration of the twenty-four Elders and of such things wherewith they praised God.

### Paragraph 751

And here the most excellent and beautiful arrangement of this book seems to me worthy of observation. At the beginning of this vision, he introduces the same elders, teaching us by their example and hymns what we should do; therefore, he brings them again also at the end of this vision, that we might be instructed again by their words and actions, not only concerning the final judgment, of what sort it shall be—most righteous, no doubt, as all his judgments are (which the entire vision confirms)—but also that we should understand what is fitting for us, and what we should do: truly, that we should worship God, and wholly submit ourselves to him; and steadfastly believe that both the judgment shall assuredly come, and that it shall also be most just.

### Paragraph 752

The hymn or prayer they offer to God is a form of praise. It is a thanksgiving or rejoicing for victory. In this way, they give God thanks, yet also exalt God and rejoice with themselves and all the faithful in their salvation. They give God thanks for their salvation and commend his justice and truth, which he demonstrates in this judgment, rewarding the good with good things and the evil with evil. Therefore, just as they rise from their chairs and fall down before almighty God, so also we ought to do now and always. This is discussed further in chapter 4. Here we should learn humility, and that God alone is to be worshipped, and that all prayers, invocations, or thanksgivings must be offered to him alone—a matter entirely contrary to the papal doctrine.

### Paragraph 753

We see now the very thanksgiving, for none better can be found. They give thanks unto God. Let us therefore thank him also. And also commend and exalt him while they call him the Lord and God almighty, and also celebrate his majesty, where they say: “Which art, and which wast, and which art to come.” They allude to the words of God, spoken in old time to Moses in Exodus 3. By the diversity of times, the eternity of God is figured. But of this kind of speech I have spoken more in the first Chapter.

### Paragraph 754

And now they declare why they give thanks: for you have received your great power, and have reigned. God truly never relinquished his power, so that he needs to receive it again; but when he does not display it, and permits much to the ungodly, so that by their power they can infringe and prevail against God’s word, he seems to have laid it aside. Therefore now that he oppresses the wicked, and as a judge advances the godly, maintains the truth, and destroys falsehood, it is truly said that he has received his great power. Likewise now it is said that he reigns, not because he did not reign before, but because the Lord has reigned in the midst of his enemies, so that at times it was doubtful and uncertain whether Christ reigned or Antichrist; yes, that he has had the upper hand, and Christ has been oppressed: now that Christ has broken all the power of his adversaries, it is said most truly that he reigns. And Erasmus rightly admonishes in his annotations on the New Testament, that the translator would have turned more aptly *Ebacilenſas* if he had said, “You have obtained a kingdom.” For the Latin men say, *Regnauit*, “He has reigned,” which has left reigning: as they have lived, who live no more. But with the Greeks it is otherwise, at least in these words. To our judge, most just, most mighty, and most righteous, be praise and glory, forever and ever. Amen.

## Sermon 51: ¶ The thanks geuyng of the Elders is expounded, the Temple is opened in heaven, the arcke appears, and there were made lightenynges, &c.

### Paragraph 755

.

### Paragraph 756

And the Gentiles were angry, and your wrath has come, and the time for the dead to be judged, so that you may give reward to your servants the prophets and saints, and to those who fear your name, both small and great; and to destroy those who destroy the earth. And the temple of God was opened in heaven, and there was seen in his temple the Ark of his Testament. And there followed lightnings, and voices, and thunderings, and an earthquake, and a great hail.

### Paragraph 757

I showed you how the Elders gave thanks to God for their salvation, also praising God’s righteousness and excellent truth, which he demonstrates in his judgment as most righteous, wherein he rewards the godly with just rewards and punishes the wicked with deserved punishments. And by this figure of speech, they teach us that judgment will assuredly come, and that it will be entirely holy and just. Would that those who judge them to be speaking of trifles would diligently consider these things themselves, those who speak of that horrible and most dreadful day of judgment. For we look for things more terrible than any tongue, however eloquent, is able to express.

### Paragraph 758

He recounts the wrath or tyranny of non-believers, against the faithful, cruelly and continually executed, so truly that to many it seemed God was unresponsive, and neither could nor would be angry. But the judgment ones made, the elders extol God’s truth, and say the wrath has come. Doubtless the holy Prophets of God have always threatened punishments, testifying that God is angry, both with sinners and with sins. But where the wrath of God did not appear immediately, the Prophets seemed to frighten people with empty terrors, as if making them afraid of their shadows. But now, say the elders, the truth has appeared, and the wrath of God is come. And the wrath of God shows itself in God’s just vengeance.

### Paragraph 759

Moreover, they also praise this for demonstrating God’s truth and justice, for the time of the dead has arrived, that they be judged. Until now, while the world prospered, they seemed to tell fables and old women’s tales, which spoke of the resurrection of the dead and the life to come. For the resurrection of the dead was scorned by philosophers and worldly men. But the elders rejoice also, that this time has come, and that the dead are raised up, that is, that the bodies of the dead have risen again and come to judgment. Concerning this, the Apostle says: “We must all appear openly before the judgment seat of God, that every one may receive according to what the body has done, whether good or evil.” 2 Corinthians, chapter 5.

### Paragraph 760

Furthermore, they greatly praise God’s justice and truthfulness when they clearly explain how God, by his just judgment, renders to every one according to his deeds. He declares what he rewards and whom he rewards. First, he pays wages or hire. For reward is promised by God for good works. For in Jeremiah 13, the Lord says, “Restrain your voice from weeping, for there is a reward for your work.” And the Lord also says in the gospel, “Be glad and rejoice, for your reward is great in heaven.” And again, “The Son of Man shall come in the glory of his Father with his angels, and then shall he render to every one according to his doings.” So the Apostle said that every one must rise in his own body, that every one may receive such things as are done by the body, whether it be good or evil. While this world flourishes, and the wicked rejoice in their indulgence, and the godly are afflicted, and afflict themselves with continual mortification, the flesh judges that these lose both labor and expense, while the others are very happy. This too is declared in chapters 3 and 4 of Malachi. But at the final judgment it shall finally appear that the godly have not labored in vain, nor have the wicked scorned God unpunished and despised godliness. For God rewards every one according to the quality of his work, which he calls wages. Nevertheless, the godly do not abuse this saying in the meantime, acknowledging it to be of free mercy that they have believed and worked with good faith, and that their good work is therefore accepted by God because it is in Christ. Of this I have written in book 3, chapter 10, concerning the grace of God justifying, showing that reward cannot be proven by merit.

### Paragraph 761

Secondly, they declare to whom he gives reward; I say to two sorts of men: good, I mean, and evil. Again, he recounts many kinds of good men. First, he calls these the servants of God, as those who are subject to the empire of God alone and obey him in all things. Immediately, he names the prophets, teachers of churches. Concerning their state, more is spoken in chapter 11. These are sometimes more unfortunate than any others in this world and are accounted by many as great offenders, so that if they were removed, all clarity would seem to return. Therefore, they are justly recounted in the register of those who receive a reward from the Lord, namely, in recompense for their labor. Now into this account come also the saints, that is to say, all the godly who, being sanctified through faith with the spirit and blood of God, have lived a holy life, keeping themselves from all worldly pollution. Moreover, into the reward and honor of the holy saints are reckoned those who fear the name of the Lord: that is, they who are very holy and religious in deed. Finally, lest any man should think any of the faithful excluded, he adds, “to small and great”: that is to say, to men of all ages, status, and sex, and so forth.

### Paragraph 762

[note 51.6]After he comes to the evil, and adds: “and shouldest destroy those who destroyed the earth.” These things seem borrowed from the prophets, with whom there is much mention of those who destroy the earth, whom the Lord will ultimately destroy. And under the name of destroyers, John understands first Tyrants, Kings, and Princes, who are persecutors of the church. Also, men of war and soldiers, who by unjust wars destroy all things with sword and fire. Secondly, he understands unjust judges, moreover oppressors of the poor, who afflict widows and orphans; moreover, those who through usury, theft, deceitfulness, extortion, and evil means are harmful to all men, and by their insatiable covetousness bring about a scarcity of all things. Finally, those who by whoredom and adultery defile and break holy matrimony. Last, heretics destroy the earth, and those who infect people with corrupt doctrine, who dwell upon the earth; into this number also come seditious persons and traitors, and other wicked men.

### Paragraph 763

These the Lord will destroy with everlasting destruction, so that they do not cease to be those who perish, but become much more miserable while they are tormented with torments that will never end. Those who are wasteful and extravagant are said to be lost, yet in perishing thus they do not cease to be, but daily proceed to be more miserable, which is destruction itself.

### Paragraph 764

Furthermore, Saint John questions this doctrine of the reward of the godly; and what he previously presented as thankful praise and a joyous triumph, he now proposes as if it were to be seen with the eyes through a heavenly vision. He concludes this vision with the opening of the Temple, just as he began it with the opening of Heaven. For the loving Lord opens to his servants heaven itself to be seen by the mind, so that we should doubt nowhere the glory prepared for us in Heaven; neither should say, “Who has seen those heavenly things that are promised us?” For just as the blessed fathers, the Prophets and Apostles have had many visions of this sort, effective, true, and godly, so may each of us with the eyes of our mind, through true faith, look into Heaven itself. I know well that worldly men pass nothing upon such visions, as the Lord said in the Gospel: “The world cannot receive the Spirit of truth, for that it sees him not, neither knows him.” Let us not care for their contempt.

### Paragraph 765

Let us therefore,[note 51.9] see what is prepared for the servants of God in another world. First, John saw heaven open: and in heaven itself, he sees also the very temple of God, open, that is, to all the godly. By the Temple of God, he understands the secrets of God, the inward and privy parts of Heaven, whereinto he will receive to the fruition of himself all believers. But in that divine temple of heaven was seen the Ark of his Testament.[note 51.10] For God made a covenant or league with the faithful, that he would be their God, their fullness, and a most plentiful sea of all goodness, a most abundant, and most sufficient plenty of all things. The confirmation, testimony, and declaration whereof is the Ark of covenant, the very Son of God, in whom dwells all fullness of deity, and in whom we are made perfect. For he is the Ark, in whom are laid up all celestial treasures, full of grace and truth. This Ark of good things, and of eternal felicity, appears in heaven. For the Son of God is on the throne of God. The liberal and bountiful heavenly Father will pour out this Ark upon his children, granting to them through Christ his only Son all heavenly gifts, that we might be partakers of all Christ’s benefits, even to the deity, wherein he excels his brethren. Hereby it appears how Moses prepared the Ark, after the example of the same which he saw in heaven: and the figure whereof was the Ark of the covenant, &c. Otherwise we shall hear in the twenty-first chapter of this book, that there is no temple in heaven, &c.

### Paragraph 766

These most beautiful things to be seen, and most pleasant to be heard, the Son of God has set forth to be seen and heard by us. Consequently, he adds that punishments are prepared for the wicked, and explains them in various ways, presenting them to be seen. Hitherto, there had been in the world lightnings, [note 51.11] voices, and thunderings, etc. The Holy Spirit shone to the world, drawing through the doctrine of the truth, moving and inspiring fear. But the reckless world would not understand, nor even hear the manner and way of salvation; therefore, divine justice requires that they be addressed in another language, and therefore, by God’s just judgment, are now made lightnings, etc. And by this heap of words, he signifies the horrible punishment that God will inflict on the wicked. He appears to have alluded to the burning of Sodom, also to the words of the godly Prophet: “It shall reign upon sinners snares of fire, brimstone, and spirit of tempest,” in Psalm 11. Therefore, this vision is concluded, as in the story of Saint Matthew’s Gospel: and these shall go into everlasting punishment, and the righteous into life everlasting.

### Paragraph 767

In these last eight chapters, we have the third part of this book, and a significant summary of the church’s history from the time of John until the end of the world. Through this, we are instructed in the true faith and warned of all dangers and betrayals by which the true faith is attacked, so that, being watchful, we may avoid all corruption and cunning deception and be kept safe. To God be praise and glory.

## Sermon 52: ¶ The description of the church and of the red Dragon, fighting against the Church.

### Paragraph 768

.

### Paragraph 769

And there appeared a great sign in the heavens: a woman clothed with the sun, and the moon under her feet, and upon her head a crown of twelve stars. And she was with child, and cried out in labor, in pain ready to give birth. And another sign appeared in the heavens: behold, a great red dragon, having seven heads, and ten horns, and seven crowns upon his heads. And his tail drew a third of the stars of heaven and cast them to the earth. And the dragon stood before the woman, who was ready to give birth, to devour her child as soon as it was born. And she brought forth a male child, who would rule all nations with a rod of iron. And her son was taken up to God and to his throne. And the woman fled into the wilderness, where she had a place prepared by God, that they should nourish her there for one thousand two hundred and sixty days.

### Paragraph 770

The fourth section of this book presents to us the third vision, which others who divide the second vision into two consider to be the fourth. The Lord has frequently and extensively mentioned in the second vision the persecution and conflict of the faithful with Antichrist and with God’s wicked enemies, especially in chapters 6, 9, and 11. He now proceeds in the third vision, and at length, to discuss this conflict, setting forth the events as if they were to be seen with the eyes, in the three chapters immediately following chapters 12, 13, and 14. He repeats all things more deeply, and vividly and diligently describes the aspects of this conflict, and afterward the fight itself. Therefore, after describing the church, which endures the brunt of this war, he also describes the Dragon that instigates the war. He declares how busily he watches, and again, lest anyone should be discouraged, he adds that, despite his efforts, Christ certainly overcomes him, and that God ultimately hinders and defeats his endeavors, yielding him vanquished to the faithful. Now he describes the chief instruments that Satan uses in assaulting and persecuting the church, namely the old and new Roman Empire, and within this, the corrupt papistry, in which Antichrist is also vividly portrayed. Nevertheless, he adds to these unfortunate events, for the consolation and comfort of the godly, that the Lamb nevertheless stands on Mount Zion as a conqueror, having his church with him, however much this world rages and may be never so mad and cruel, and that the gospel is preached in defiance of Antichrist, and all men are warned to beware of Antichrist. Here he also begins to reason about God’s judgment against the wicked, so as to prepare him to speak in the fifth section about the pains or punishments of the Antichristians, a treatise which he begins in chapter 15. Hitherto, therefore, he treats of the fight or conflict of the church and of the wicked, namely of Antichrist, all of which the father of all murder and of all iniquity, the Devil, inspires.

### Paragraph 771

Therefore, just as this entire book is drawn from scripture and explains the older scriptures very well, these things also which were mentioned earlier seem to be taken from the third chapter of Genesis. There the Lord says, “I will put enmity between you (meaning the Serpent) and the woman, between your seed and her seed: her seed shall bruise your head, and you shall bruise his heel.” For you will read at the end of that chapter as well: “And the dragon was angry with the woman and departed to make war with the remnant of her seed.”

### Paragraph 772

And he describes above all things the parts of this conflict, her who was besieged by war, and the one who instigated the war, namely, the church and the Dragon. And he says how a sign of these things appeared in the heavens. For he would not only say or write, but also set them forth to be seen by the eyes, and in a manner to depict, so that all things might be seen more plainly. And where he says those signs were great, he admonishes that they were and are things of the greatest weight, and matters of the greatest importance.

### Paragraph 773

[note 52.5]First, he describes the church of God throughout all times using the image or figure of a woman. It is neither strange nor rare that the woman began to represent the image of Christ’s bride, the church, at the very beginnings of things, as can be seen in Genesis 2. And so the Apostle explained the image in chapter 5 of Ephesians. I need not now recount that Isaiah has repeatedly used the image of a woman to represent the church of God: “Rejoice, you barren one who does not bear,” he says, and so on. Finally, Paul in Galatians 4 uses Sarah as a figure of the church, which Solomon also explored at length in his Song of Songs, describing his bride. The church is that woman united with Christ, her spouse, in true faith and continual love. Afterward, he applies certain things specifically to the Virgin Mary, although the things that precede and follow do not entirely agree with this, a point shown by Methodius and Primasius, and other commentators as well.

### Paragraph 774

This woman is clothed with the sun. [Note 52.6] Scripture calls Christ the sun of righteousness and the light of life. Paul commands the church to put on Christ. Therefore, he is the light, the life, and the righteousness of the church; by Christ, the nakedness of the church is covered; Christ is the ornament and beauty of the church, through him it shines in the world.

### Paragraph 775

The Moon is subject to alterations, is variable, and receives sundry colors: she increases, and decreases: and although it shine, yet it always appears full of spots, and borrows her light of the Sun. Therefore all courses and alterations of times, and whatever is mutable and corruptible in this world, all affections also and infirmities, the church treads under her feet. All the light that she has, she has it of Christ; the light of her righteousness increases and decreases. Finally, she gathers always some spots of the nature of flesh, which she cannot leave but by death. Therefore she shines indeed, however the church feels some obscurity. As the Lord has said also, every branch bearing fruit he purges, that he may bring forth more fruit. And he that is washed, is all clean, and needs no more but to wash his feet.

### Paragraph 776

Moreover, a crown is the honor of the head, and a sign of a kingdom. Christ is the beauty, comeliness, and king of the church. In this crown there are no precious stones, but stars. For in Christ are the patriarchs, prophets, and the twelve apostles, who beautify and enlighten the church, and who have the light of the crown, and dispense that power into the church. Thus is signified the doctrine of the ministers, as in the first chapter of this book. Neither is the shining ministry the smallest portion, but it represents the most excellent things of the church.

### Paragraph 777

Moreover, that woman has in her belly, which in a certain phrase of speech is as much to say as that woman was with child; and had not only a great belly (as we say) but, after the manner of women in labor, cried out and, laboring with pain, was about to be delivered. This properly does not pertain to the Virgin Mary, but to the church. For the primitive church, concerning that first promise of the blessed seed, conceived in her mind a most assured hope that, at length, the Son of God should be born of a virgin, namely, the seed promised, which would crush the serpent’s head. Therefore, the church with an earnest desire and with most fervent prayers coveted and wished that Christ might once be engendered in and by the excellent womb of the same holy virgin. Moreover, Christ is begotten in his faithful when, through his virtue, they are regenerated. For Paul says, “My little children, for whom I labor again, until Christ is shaped in us.” The church therefore travails and brings forth in two ways: bodily, while she earnestly desires, with pain, that Christ might be born of the virgin; and spiritually by faith and regeneration, while she desires to be made conformable to Christ in her members. This therefore is the nature and disposition of this woman, that with a great desire embracing the incarnation of Christ and redemption, she would fain have it known to many; and that many times she wisheth to be regenerated and reformed after the image of Christ.

### Paragraph 778

This is truly a fine description of the church. Compare this to those who today claim the title and the guise of the church, and judge how well they align with this description. But this true church of Christ is brought into danger and conflict.

### Paragraph 779

Now let us hear, in the second place,[note 52.11] and as it were on the opposing side, what kind of person is the adversary or enemy of the church: namely, that old serpent, who was a liar and a murderer from the beginning, the sole author of all evil, all mischief, all errors, all iniquity, murder, and disturbance, and a most ungracious devil. Whom afterward he calls Satan, the seducer of the world, and adorns him with other titles, fitting for such a majesty.

### Paragraph 780

This is the Dragon, and that the great Dragon,[note 52.12] namely, of great power throughout the world in his followers. And a Dragon, because in ancient times he assumed the form of a serpent and deceived our parents. Pliny and other authors write many things about dragons. Scripture in some places,[note 52.13] calls the Devil a twisting serpent. For he is wonderfully subtle and can transform himself into countless forms, that he may deceive and keep the deceived in error.

### Paragraph 781

He is red, for he is full of fire, and the blood of saints and of the innocent. A true bloodhound, the parent and patron of all persecutors and bloody soldiers. In him remain the stains of the blood of Abel. He still smells of the shedding of the blood of the Prophets and Apostles.

### Paragraph 782

This has seven heads; upon each of these is seen a royal crown. He also has ten horns. For the Devil is called the prince of this world and has indeed been the governor of wicked rulers of all ages, and the ringleader of all horns and bloody realms. He was therefore the head of Ninus, the king and prince of Pharaoh, chief captain of Balthazar, king of Babylon, of Cambyses the Persian, of Antiochus the Macedonian, of Julius Caesar the Roman, and likewise of all other tyrants.

### Paragraph 783

The prophet Isaiah called a false prophet a tail, because of his soothing and flattering words; for with his honeyed tongue and sweet words, he creeps into favor with great men. Therefore, with flattering and deceptive words, and lying promises, with which (as in times past) he promises his worshippers godly things, he persuades them to all wickedness, stars—that is to say, preachers and notable men—whom he draws away from heavenly things and casts unto earthly things, so that, having forgotten celestial matters and their holy office and duty, they may now cling to earthly things, being wrapped in the earthly snares of the devil’s tail. And thus in deed he shall corrupt not a few. For he puts the third part of stars, meaning a great number of notable men, whose ministry he uses against the church. Thereof there are many and notable examples from all times in all histories.

### Paragraph 784

And after he has described this foul and filthy beast, and sworn enemy of all saints from the beginning of the world, he also reveals his attempts, treasons, and bitter poison against the church, and how he began to stir up war. This dragon, he says, stood before the woman who was ready to give birth, and he stood watching, diligent, attentive, and eagerly awaiting at all times, and he observed, taking advantage of every opportunity to harm the church. But the end of all his endeavors was to devour the son, born of the spouse of God. He has always, even from the beginning of the world, sought to intercept the glory of Christ; and if any faithful person of the church is spiritually regenerated and made conformed to Christ, he also attempts to lead them into error and destroy them. Therefore, Saint Peter did not say without cause that the Devil goes about like a hungry lion, seeking whom he may devour.

### Paragraph 785

He now, by the way, shows that Christ, as he was promised, is presented to the church, so that the dragon could do nothing against him. Whereupon he wants us to conclude absolutely that he shall have no power over us either, if we abide in Christ. For now he shifts from the universal church to the singular and most excellent member thereof, the virgin Mary, and in few words knits up the mystery of the incarnation: that excellent woman, of whom it is spoken in Genesis 3. The daughter of that matron, I mean the church, the holy Virgin, brought forth a man child, that is to say, her firstborn, king and priest, as Luke testifies in chapter 2. Immediately he declares what and of how great power he is, and why he called him a man child. He is the one of whom David prophesied in Psalm 2, that he should rule all nations with a rod or scepter, not of wood or lead that is pliable, but of iron—that is, strong and durable—namely, the word of God. But those who do not obey God’s word, he will beat down far and never, with an iron staff, that is, with power, which no man is able to resist. But for this mighty prince, Satan laid an ambush, that old dragon, who stirred up against the chief of the Jews and Gentiles. But he found nothing in him at all, as the Lord himself said in John 14; nor shall he find anything in the faithful of Christ at the last. Moreover, while the dragon attempted great things against Christ by the elders of the Jews, being risen from the dead, the Lord was taken up, as it were, out of the throat or hottest assaults of the dragon, to his heavenly Father, and sat on the right hand of God the Father, the old serpent’s attempts made frustrate. And thither also he will receive his faithful, though the serpent’s guts should burst. For through hope we sit together with our head in the supercelestial places, Ephesians 2. And this is the chief and greatest hope of the church in this conflict. For thus he gathers: the dragon most strongly and fiercely invades not only the ancient church but even the very head of the church, the redeemer Christ; yet with his fury, he could nothing prevail. Therefore, he shall no more prevail against his members.

### Paragraph 786

Now he returns again to the church, and says: after the dragon could bring nothing to pass against the son of God, he went and made war against the church, and the church fled into the wilderness. Certainly, Judea in the prophets is compared to a place most frequented; the Gentiles are called a desert or wilderness. Therefore, after Christ’s ascension, the Apostles departing out of Judea, repaired to the Gentiles; yea, and the Jews, inspired by the red dragon, cast out the church out of their limits, which was constrained, as appears in the Acts of the Apostles, to flee unto the Gentiles. And where the Lord has prepared a place for his church, and the church was greatly augmented among the Gentiles, certainly it was through his grace (and by no merit of man). Which prepared the place, which calls, directs, and keeps his sheep, the same has disposed, and yet does dispose for this church ministers or pastors, which may feed it, as the ravens did Elias, all the time that shall be, unto the world’s end. For as for the nobility of those days, I discoursed before. And by this exposition is signified that the dragon shall fight stoutly against the church, so that she shall be compelled to flee; but however much he shall rage against the church, the Lord God shall yet prepare a place on earth, wherein she may dwell safe, and will ever send pastors to feed. He shows moreover, that the flight shall not always be reproachable. The Lord save and keep us. Amen.

## Sermon 53: ¶ The description of the conflict of Christ and the Church with the Dragon: the drago is overcome, the heavenly dwellers sing praises.

### Paragraph 787

.

### Paragraph 788

And there was a great battle in heaven; Michael and his angels fought with the dragon, and the dragon fought with his angels, but they prevailed not, and was cast out of heaven. And the dragon, that old serpent, who is called the devil and Satan, was cast out, he who deceives the whole world. And he was cast to the earth, and his angels were cast out with him also. And I heard a loud voice in heaven, saying, “Now has salvation and power and the kingdom come to our God, and the power of the Lord is Christ’s, because he has been cast down, he who accuses them before God day and night.” And they overcame him by the blood of the Lamb and by the word of their testimony, and they did not love their lives to the point of death. Therefore, rejoice, heavens and those who dwell therein. Woe to the inhabitants of the earth and the sea, for the devil has come down to you, full of great wrath, because he knows that he has but a short time.

### Paragraph 789

The Apostle has spoken of the parts of the remarkable conflict and worthy battle; he has spoken also of attempts and purposes of the dragon, which indeed directs all his schemes to this intent, that he may devour all godliness, that is, destroy it utterly. He has shown how he began to move war against the church, which fled into the wilderness; and now, as it were leaving the woman in the wilderness, he seems to bring forth other soldiers, who give battle to the dragon, and most valiantly do challenge and also defeat him and all his power. John therefore describes the singular fight of one most excellent, namely Michael, who overcame the dragon; and describes the general fight joined with that particular. For he adds that all the angels of Michael fought against the dragon.

### Paragraph 790

First, it is shown to be the place of the fight or conflict. For he says, “A great battle was fought in heaven.” And it is evident that Satan was, at the beginning of all things, cast out of heaven into the earth, and therefore he does not wage war in heaven, nor raise any tumult there. For heaven is a place of rest and joy, not of debate and contention. Therefore, this must be attributed to the vision. For the Lord has, in heaven, represented this battle by signs to be seen, which in reality is fought on earth in the midst of the church.

### Paragraph 791

Here is set forth an image of a significant battle, showing what has been, and what is yet to be done on earth. I said even now that this vision was indeed specific, but to have a general battle appended. For Michael fights, as captain of this war, and Michael’s angels fight also. [Note 53.2] This must be carefully discerned, although Michael and his angels make up only one part. On the other side fights the dragon, as emperor of this war, and his angels fight also. And these truly make no other parties than we have heard before at the beginning of this chapter: that the parties in this fight are the church and the devil. Nevertheless, lest the victory should be attributed to the church, and not rather to Christ, the woman must now be omitted, and Michael brought in fighting. There is some difficulty in these things, but it will be easy enough for him who will mark everything in order.

### Paragraph 792

First we must see what Michael is, and there is no doubt that the angel Michael appeared in the vision, with an army of angels fighting. And on the contrary part, the Dragon fought with a host of devils. But since we heard at the beginning that these were signs, they must needs signify and betoken other things. I suppose, therefore, that Christ is signified here, the head of his church, king and protector, with his members, Apostles, Martyrs, and faithful. It is not rare that Christ should be figured to us by angels; indeed, it is most customary that angels are called the ambassadors of God and the faithful servants of Jesus Christ. Therefore, Christ, the head of the church, and the faithful members of Christ, fight against the Dragon, yet in a different way. For Christ overcame him alone in the combat without the help of any creature, while in temptations he discomfited him at last, and also by dying on the cross and rising again from the dead, he utterly broke his head. This is the only, true, and singular victory, whereby afterwards are obtained the victories of Christ's members, derived from that general fight, wherein Christ fights not now only hand to hand with the Devil, but all the members of Christ at all times, under Christ their Captain, fight against the Devil, and in the virtue or victory of Christ, fight and overcome, as we shall hear by and by in the song of praise.

### Paragraph 793

But for great and sundry causes we affirm that Christ is figured and signified to us under the type of Michael. We know, by the scriptures as many of us as are learned, that Michael, as also Gabriel, are the names of good angels of God. Michael signifies, “Who is as God?” And who, I pray you, is such as God, but in whom the express image of the Father’s substance, which is the invisible Image and word of the Father from the beginning, I mean the very Son of God Jesus Christ? Michael in the tenth and twelfth chapters of Daniel, is president, protector, and patron of the Jewish nation. And it is plain that the people of Israel had from the beginning no other tutor and patron but Messiah himself, the blessed Seed. This appears in the seventh chapter of Isaiah, where we read that the Lord spared the people of Judah and the princely city for Christ. In another place he says most openly, “I will defend that city for myself and for my servant David.” And David is called Christ in the thirty-fourth chapter of Ezekiel. Christ is therefore in very deed governor of his people, which nevertheless in defending and delivering his people, uses the ministry of angels; who also attribute nothing to themselves, but all glory to God alone. Moreover, that excellent victory cannot without offense of godliness be ascribed to Michael the archangel. For so, omitting our Messiah Christ, we should commend angels being made and worthy to be called angelical, rather than Christians. In the law was written, “The seed of the woman shall bruise the serpent’s head.” But the Lord never took the nature of an angel, but the seed of Abraham, and by sin condemned sin. There shall follow anon in the song, “Now is salvation and power, &c.” And it is added, “for the Devil is cast out.” And this salvation has Christ alone accomplished, wherefore it is necessary that Christ the conqueror of Satan be signified by Michael.

### Paragraph 794

And the dragon fought hand to hand against the Lord, [note 53.6] not only contending with him in the wilderness, but also never ceasing to tempt and assail him as long as he lived on earth. He stirred up against him the Pharisees and princes of the people, kings and the Roman governor, and thus at last broke the Lord’s heel. This was the greatest fight of the dragon. The same dragon now inspires kings and princes, wicked priests and cruel men, his agents who may wage war upon the church. And all these truly do persecute and vex the church in the power of the red dragon. Histories declare the same to have been done before Christ’s time; the same testify, and experience proves, that the like has been done from the ascension of Christ into heaven, unto this present day, and unto the world’s end.

### Paragraph 795

Now it is also declared with what fortune they fought on either side: namely, most fortunately concerning Christ, and most unfortunately as touching the Devil or red Dragon. And in this fight, as also in the song immediately following, is contained the whole fruit of this disputation. For from this, all the faithful may learn that Satan, our enemy, is unarmed; and that Christ in this conflict is on our side, as our Emperor and captain at all times, by whom all the faithful may easily overcome in all conflicts. Therefore, this matter of battle and victory is placed immediately after the beginning of the most dangerous battle with Antichrist and Antichristians, who are the brood or offspring and scales of the serpent, and champions of the Dragon, for comfort and consolation. And the natural order is here altered, which does not treat of the success of the battle until he has set forth all the conflict before. But this battle shall be continued hereafter in the rest of the twelfth and all of the thirteenth chapter.

### Paragraph 796

He states this in three phrases: first, the victory of Christ; secondly, the victory of all Christians. The first is, “they prevailed not, they had no strength.” Doubtless the devil’s power is great if God permits, and clearly greatest, in consideration of God’s just judgment, as also appears in Job, that he is able to quench and break the strongest things. But the Lord says in the gospel: “The prince of this world came, and against me he has nothing.” Again in the gospel: “The gates of Hell shall not prevail against her, the rock I mean, and secondly against the church.” Although therefore the devil makes a horrible uproar and cruelly rages against Christ and his church, he is without force. For the virtue of Christ prevails.

### Paragraph 797

The second part is, neither is their place any longer found in heaven, which manner of speech signifies nothing other than that the condemned angel is removed from all dignity, glory, and power; moreover, that he has no more place in the church, or among the elect of God; not that the devil should not return, or should not tempt, or renew war, but because he has no permanent place. This relates to what the Lord so often repeated in the Gospel, and now the prince of this world is cast out, in the 12th and 16th chapters of Saint John. Moreover, by other passages of Scripture it is manifest that the devil is shut out of heaven. And it shall be easy for us to shut him out, who, being cast out by the Son of God, has no place in us, unless we ourselves give place to him. Which we should not do, the Lord admonishes us diligently, that we should watch. The story is known in the 12th chapter of Matthew, of the devil, proposing to return, and therefore taking unto him seven worse spirits. But why do you hear him, why do you obey him, whom you see shut out of heaven? Nevertheless, hereby is also signified that the devil was so fully vanquished by Christ, that he was also driven to forsake the place of the battle.

### Paragraph 798

For the third member, as it were expounding the second, adds: and he was cast to the Earth. For those who are thrown to the ground are judged to be overcome. Therefore, a full victory and perfect conquest is signified. However, he was once most valiantly cast to the Earth. Of our Lord Jesus Christ, in the mystery of our redemption, and in the virtue of the same, is daily cast to the Earth, by the faithful. And just as the Devil has no permanent place in heaven nor in the elect, so verily he inhabits all earthly things, that is to say, men savoring the earth and contemning heavenly things. Yea, and we hear that his angels are cast out with him. For the Lord in the Gospel of Saint John, chapter 16, says: “In the world you have affliction, but be of good cheer, I have overcome the world.” And Saint John in his canonical epistle says, “You are of God, little children, and you have overcome them, for he is greater that is in you, than he is in the world.” And this is the victory that overcame the world, even your faith.

### Paragraph 799

And by the way he explains what we should understand by the dragon, of whom he has spoken until now, namely, the old enemy of humankind. He depicts him with his titles, attributing four names to him, so that we may also better understand his nature and beware of that wicked murderer. First, he calls him the old Serpent. For at the beginning, by the serpent, he infected our first parents with the poison of death and sin, and by that, the entire world: as is seen in Genesis 3 and Romans 5. Therefore, I said at the beginning of this chapter that he is called a dragon. Afterward, he calls him the Devil, that is to say, a slanderer or a false accuser. For it immediately follows, which may explain this word: “the accuser of our brothers has been cast out,” etc. A fitting example of this thing is declared in the first and second chapters of Job. [illegible] signifies to accuse or blame, and [illegible] is an accusation, and [illegible] a crime or complaint.

### Paragraph 800

Thirdly, he calls him Satanas, in the Hebrew word, namely an adversary, for he is in all things against God, and opposes himself and resists people in holy matters, if perhaps he might hinder or corrupt them. Lastly, he is called a deceiver, a betrayer, or one who subverts and betrays the entire world. For this the Lord attributes to him in John 8, for that he has been a liar from the beginning, and is the father, that is the fountain and origin of all lying, deception, errors, and seduction, and of all evil. For all errors and heresies, all deceptions, and all falsehoods, finally all kinds of evils, have flowed out of this most filthy wellspring. And who is he that hears these things, which will not abhor that vile beast? They must needs be stark mad who seek by all means to be in favor with that wicked spirit.

### Paragraph 801

He should now consequently add the remainder of this account, namely how the Dragon persecutes and attacks the woman, and she again resists by fleeing and overcomes through Christ. But he suspends this narration for a little while,[note 53.13] and places now a song of victory and triumph of the saints in heaven, of the angels and blessed souls. The sum of it is that Christ has overcome, and that the faithful overcome in Christ; and therefore even the heavens themselves, and all who dwell therein, must rejoice and sing. And I repeat that these things are interlaid in the dangerous Antichristian and Roman view, for a consolation, lest the saints should in those great dangers by reason of their natural infirmity be discouraged; but calling upon the name of Christ, should fight manfully, when they understand under whose banner they fight, and with whom they fight: truly, with one overcome under Christ’s standard. And when we hear that the Dragon’s force is broken, we shall think that the furies of either beast, as well the ten-horned as the two-horned, are weakened in the faith of Christ. This gives also no small courage in this conflict, that we see that the Dragon has no power over those who are sprinkled and purified with the blood of Christ, but over earthly and worldly men. And this triumph is heavenly. For voices are heard out of heaven, singing a merry note, to the intent that the rejoicing of the blessed spirits might have more authority, grace, and efficacy among the poor afflicted.

### Paragraph 802

They all with one voice maintain simply that salvation and power are now made perfect, for by the Lord’s death and resurrection, God has wrought power and made perfect the salvation promised to the fathers, namely while he trod down the serpent’s head, abolished sin and death, and restored life. Thus is the kingdom of God in this world established in the elect, while even by the power of Christ, the Prince of this world is cast out and overcome. For this is the reason we must rejoice, and to see the virtue and power of Christ revealed, and how salvation is made perfect: because, he says, through Christ, the Devil is cast down, that is to say, overcome and vanquished, so that he can no more accuse humankind before the judgment seat of God. To this belongs what Saint Paul wrote: Who shall accuse the elect of God? It is God that justifies; who is he that condemns? It is Christ, who died, yes, who rose again, who is also on the right hand of God, making intercession for us.

### Paragraph 803

Moreover, the heavenly beings do not only preach the victory of Christ, but of all the faithful, which they obtain against Satan in the faith of Jesus Christ; that it may hereof at least appear what we should understand before by Michael and by his angels. And he emphasizes that Christians overcome Satan not by their own merits, force, or strength, but by the merit and grace of Christ. And they say, that is, the angels of Michael, overcame the Dragon by the blood of the Lamb. For inasmuch as the faithful are purified by the blood of Christ, Satan has nothing against them; but if they have the spirit and faith of Christ, they overcome the Devil also. So in times past, the destroyer had no power over those houses which were marked with the blood of the Lamb, Exodus 12. And he adds another thing for which the faithful overcame: for the word of the testimony of Christ, which is the gospel. Which because it is invincible and eternal, they overcome all things of this world, whoever abides in the lively and eternal word of truth. And even in the most true gospel, the Lord himself has promised that he will not forsake his own, and will fight for them. Therefore, the faithful must needs overcome. To these things is added the effect of Christ’s purifying: they loved not their life more than Christ; and therefore they have given it for Christ unto death, and so have overcome. For many are vanquished by this one thing, that they will not hazard their life for Christ.

### Paragraph 804

For these great benefits of God, they now exhort the heavens themselves, and all the inhabitants of heaven, that is to say, they exhort one another to sing a joyful song. And that which the blessed saints say they do here, they teach the saints on earth to do also, and instruct them about what manner and sort they ought to be, who shall overcome Satan in battle—namely, purified by the blood of Christ, cleaving to the testimony of Jesus Christ, and contemptuous of their own lives, to whom it seems not grievous to die for Christ’s sake.

### Paragraph 805

Finally, concerning the song’s ending, they declare that the Devil will reign and hold sway in earthly and fleshly men—those who truly deride godly things and value only worldly possessions, which are destined to perish. For the acquisition and preservation of these things, they will be willing to do anything, no matter how difficult. For Christ’s cause, they will endure nothing, neither action nor suffering. To these, they pronounce a terrible woe, namely, the curse of this present age and of the life to come.

### Paragraph 806

But in him whom the Devil possesses in his kingdom, he also utters his malice against the elect, and that his great malice. For he rages most cruelly against the godly, and against godliness. He rages also most extremely against those his worshippers, whom he pollutes with all kinds of filthiness, and with all shame and reproach defiles.

### Paragraph 807

Again I suppose that same pertains to the comfort of the godly, that is spoken of a short time. For Satan, indeed, through Antichrist shall most cruelly rage against the church, but those days shall be shortened for the elect’s sake. By the way, it is noted also the wicked nature of Satan, which, knowing that the last judgment is at hand, wherein he must be thrown headlong into hell, thinks to requite and recompense the shortness of time with the cruelty of his wrath and devilish fury.

### Paragraph 808

And thus far concerning the victory of Christ and his saints. Now follow, with less terror, yet horrible things concerning the war which the dragon most greedily and fiercely instigates against the mother of God. The Lord Jesus brings him under our feet. Amen. Amen.

## Sermon 54: ¶ The Dragon perſecuteth the woman: she is defended and preserved of the Lord. The Dragon stands on the sand, &c.

### Paragraph 809

.

### Paragraph 810

And when the dragon saw that he was cast down to the earth, he persecuted the woman who bore the male child. And to the woman were given two wings of a great eagle, that she might fly into the wilderness, to her place where she is nourished for a time, times, and half a time, from the face of the serpent. And the dragon cast water out of his mouth, as a river, after the woman, that he might carry her away with the flood. And the earth helped the woman, and the earth opened her mouth and swallowed up the river which the dragon cast out of his mouth. And the dragon was angry with the woman, and went and made war with the remnant of her offspring, who keep the commandments of God and hold the testimony of Jesus Christ, and he stood on the sand of the sea.

### Paragraph 811

[note 54.1]Having begun to speak of the persecution of the Dragon and the flight of the church, and having deferred it briefly to declare Christ’s victory, he now resumes and completes that discussion, vividly and most clearly describing the battle, and repeatedly emphasizing the help of God, which is given to the church through God’s grace.

### Paragraph 812

When Satan, therefore, saw himself overcome by Christ and completely cast out, he began to rage against the church redeemed with the blood of God’s Son and vexed her with grievous persecution. For immediately after Christ’s ascension, a great persecution was stirred up against the Apostles and the Apostolic church. The Apostles, put in prisons, were grievously rebuked with words, and also scourged with rods and whipped. Stephen was stoned, James beheaded with the sword [note 54.2], finally, by the means of Paul (who then was called Saul), innumerable were cast into prison and put to cruel torments.

### Paragraph 813

On the contrary, he recounts the present help of God, which he expresses through a figurative kind of speech, [note 54.3] following the manner of the vision for greater effectiveness. For he says that to the woman, meaning the church, were given two wings of a great eagle, by whose aid she fled into the wilderness, where she concealed herself for a time, safe from the dragon’s sight. And here is signified that a great power is granted to the church to flee and escape the fury of Christ’s enemies, and to proclaim the gospel among the Gentiles.

### Paragraph 814

Whereof you may read in the eleventh chapter of the Acts of the Apostles. And he does not mention the eagle’s wings without cause, and that of a great eagle. For Moses in Deuteronomy makes mention under this figure of the defense and aid of Almighty God, just as the eagle flies over her young, and stretches abroad her wings, and carries them on her shoulders, so the Lord has also kept and advanced them.

### Paragraph 815

Furthermore, he says there is a place given to the church in the wilderness, namely, provided by Christ, which turns to him whom he will, and prepares for himself a spouse. And he nourishes the church among the Gentiles with his evangelical word, as he nourished his people in old time in the desert with manna. And the time of the church he does not prescribe. For he uses again a kind of speech as it were a riddle, borrowed out of Daniel, which God employs when he wishes the time to remain unknown to us; which, since we know that good and just things consist in him, we should not curiously inquire after. Of this I have spoken before. Doubtless it is plain that the church among the Gentiles shall continue and remain until the last judgment. But the day of judgment no man can define.

### Paragraph 816

Again he declares, [note 54.6] with what fury the Devil inflamed, shall make a new and continual war against the church. When he saw the church among the gentiles to be daily increased and established, he vomited after the woman water: and that we might know the figure, he adds, as it were a river. For he signifies that the Devil has unleashed a sea of evils into the church, sects I mean, dissensions, tumults, seditions, persecutions, wherewith the whole world has overflowed. Verily he raised up everywhere all magistrates and priests against the Apostles and Apostolic doctrine. Read the Acts of the Apostles chapters 13, 14, 15, and the chapter following. Neither is it rare in the Psalms, by waters, floods, and rivers, to understand all kinds of afflictions. And to this end he raised up those great evils, and unleashed them upon the godly, [illegible], that he might cause the church to be carried away by the flood: that is, that he might take away the godly and the doctrine of piety. And this is the continual endeavor of Satan; hereunto he applies all his consultations and doings. So in the empires of Nero and Domitian, he studied to wash away the church by the blood of saints, but yet in vain. For therefore I suppose it is said, he vomited a flood of evils after the woman, not upon the woman.

### Paragraph 817

For God never failed his afflicted church. It is a thing to marvel at, that the earth opened her mouth and swallowed the flood poured out from the serpent’s mouth. This earth drank up in old time and covered the blood of Abel. Here is signified that the godly, enduring persecution, have help from an unexpected source, as David is recorded to have been delivered by the help of the Philistines, expecting nothing less than to deliver David from the hands of King Saul. But while they intended another thing, they brought to pass that which seemed good to the Lord, who can turn the evil counsels of evil men to the profit of the godly. And we see many times in the Acts of the Apostles that the earth has swallowed a flood of evils; that is to say, that worldly men, doing something else in the meantime, have procured peace for the church. So did the town clerk or recorder of Ephesus pacify the multitude of Ephesians, who were all in an uproar and worse than mad. Lysias, the captain of the guard, takes Paul away from the bloody hands of the Jews; so did the centurion defend Paul, that he should not be slain by the soldiers in the shipwreck. The civil wars that began immediately after the death of Nero gave peace to the church until the reign of Domitian. But the old serpent, who can never rest, attempts new wars. For now, being mad with rage at the church, he goes to make war against the remnant of the woman’s seed, which is truly to be borne until the judgment of the church by the word. And so he makes way to the Roman persecutions, which followed immediately after the time of Saint John, in the Roman Empire, and the Antichristian persecutions raised after the Empire was subverted. This will be spoken of in chapter 13 and following verses.

### Paragraph 818

Nevertheless, from this it chiefly appears what Saint John understands by the woman, namely that which engenders the seed of God. The church is called both Mother and daughter. The daughter, because she is engendered by the word preached in the church; the Mother, for that by the word, she brings forth spiritual children to Christ. For the seed of God, and the seed of the woman, are all those who keep the commandments of God and have the testimony of Jesus Christ. They keep the commandments of God, which esteem God’s law and frame all parts of their life according to the same. They do not keep God’s commandments, which set nothing by the law or word of God, nor frame their life after the same. Of this matter, there is spoken at large in the fourteenth [?chapter]. The testimony of Jesus Christ is nothing else but the gospel of Jesus Christ, preaching to us the free remission of sins. They have this, which possess it by faith.

### Paragraph 819

And where he says,[note 54.8] that the dragon stood on the sea shore, it is a preparation for what follows: for presently he says how the beast, the principal instrument of the dragon, came out of the sea by the devil’s means. And it is a comfort that the dragon is said to stand on the sand, and not on the rock. For it signifies that the fury of Satan shall not long endure against the church, and that the kingdom of the devil shall be shaken and fall to decay, whose foundations are laid upon the sand.

## Sermon 55: ¶ He exhibiteth a noble instrument of the Dragon to be scene, the old Roman Empire, which describes what manner a one it is, &c.

### Paragraph 820

.

### Paragraph 821

And I saw a beast rise out of the sea, having seven heads and ten horns: and upon his horns ten crowns, and upon his head the names of blasphemy. And the beast which I saw was like a leopard, and his feet were as the feet of a bear, and his mouth as a lion. And the dragon gave him his power, and his throne, and great authority: and I saw one of his heads as it were wounded to death, and his deadly wound was healed, and all the world wondered at the beast, and they worshipped the dragon which gave power to the beast.

### Paragraph 822

John proceeds to describe, through the revelation of Jesus Christ, the prominent agents of the devil, by whom he has afflicted the church of Christ with continual and most grievous persecution. And he speaks of the old and the new Roman Empire. John could not, without exceeding great danger, utter, much less describe, these things without some human assistance, and while banished and driven into exile. For the Roman Empire was considered godly, invincible, most sacred, and everlasting. Nevertheless, the Apostle speaks and writes of this in such a way that it seems he cannot avoid being considered a seditious person, and offending against the holy majesty both of the emperor and the Empire. But what, I pray, would you do, if God commanded you to speak and write so?

### Paragraph 823

The world also rages at this day, when they hear kingdoms and governments rebuked by God’s word for sin and wickedness committed. Some princes, acting imperiously, issue proclamations commanding that such things be heard no more. But the Lord says in the gospel: “If these remain silent, the stones will cry out,” signifying completely that the truth must be preached, and that it cannot be suppressed or extinguished by any decrees, threats, force of arms, or punishments. Therefore, if those to whom the office of preaching is committed were to remain silent today, the Lord will stir up other preachers, who, even if the entire world says no, will bear witness to the truth. Therefore, I would counsel princes not to waste their efforts in vain attempts against God’s truth, for they shall not prevail. The truth will triumph. For he who then equipped John against the Roman Empire, at a time when it was most flourishing and powerful, the same also at this day reveals his truth to the world, now broken and aged, and will surely overcome. Woe to those stubborn natures that love to deceive. Let all preachers learn, by the example of the Apostle John, to freely utter the things they have received in commandment, and to fear no man. He who is in us (as the same John said in 1 John 4) is greater than he who is in the world.

### Paragraph 824

And the beast he calls the Roman Empire of great authority, and as it were godly, not without weighty considerations. For the Lord still keeps the language of Scripture, imitating Daniel, who in chapter 7 attributes the name of beast to the Roman Empire. And Saint Jerome, explaining the prophecy of Daniel, understands the beast to be the Roman Empire, and supposes that is why it is not called a lion, or a bear, or a leopard, but a beast: so that whatever cruelty anyone can imagine in beasts, by that they may understand the Romans; for surely in their manners they have shown themselves beasts. Mithridates, the renowned king of Pontus, speaking of the Romans in the thirty-eighth book of Justin, as they themselves report, says that as their founders were nurtured by sucking a wolf, so all that people have wolfish minds, never satisfied with blood, hungry and empty for rule and riches. And now how filthy beasts many Roman princes have been, their own writers testify, chiefly Suetonius, and others who have written of the lives of the emperors. And that the people of Rome were also of beastly manners, chapter 1 of the Epistle to the Romans proves.

### Paragraph 825

You will say, I know well, sins that John includes under this image refer to the entire Roman Empire. Shall we call Constantine, Theodosius, and other godly emperors beasts? I say that the Scriptures use this manner of speaking, and by beasts, in truth, understand empires, even though they do not call all those who dwell in those empires beasts without distinction. Therefore, we understand that those living a life acceptable to God are exempted within all empires, and know assuredly that neither Daniel nor John would have defiled such innocent and praiseworthy men with words. Indeed, in all this treatise concerning the empire and Antichrist, we always except those who are innocent and excel in virtue. Of this we shall happily speak more hereafter.

### Paragraph 826

And first he shows the beginning of this Empire. A beast comes out of the Sea, on the sand whereof stands the Dragon: and in the 17th Chapter it is said how the beast came out of the bottomless pit. Therefore the beginning of it is referred to Satan. Nevertheless, we must here take diligent heed that we do not take anything away from the Lord our God, which he claims for himself. Scripture in sundry places, but chiefly by two most excellent witnesses, by Daniel in the 3rd Chapter and Paul in the 13th Chapter to the Romans, has set forth that kingdoms and empires are of the Lord, and that he sets up and deposes kings. There is no power, says the Apostle, but of God. And hitherto the apostles command obedience to princes and magistrates. How then is it that we hear that the Roman Empire came out of the bottomless pit, since the Apostle speaks of the same? Doubtless the Roman Empire is not absolutely of the Devil. For God is the author of monarchies, and preserves realms and policies, giving thereunto certain faithful servants. But Satan meddles with men's matters and corrupts both kings and kingdoms; and so long they are of the Devil. The Christians in all political matters obeyed emperors, but commanding idolatry they obeyed them not. Certain it is that God did institute the kingdom of Israel or of ten tribes by the prophet Ahab; yet nevertheless the Lord cries out in another prophet, they have reigned indeed, but not by me. For the Lord would have had those kings to have framed all things after his word, and to reign in the fear of God: and where they did not so, but following the instigation of Satan ordered all things after their own lust, they are rightly said to reign, not of God, but of the devil. Therefore have the godly obeyed kings: but they obeyed them not commanding wicked things, although they took them for their kings. God had instituted the order of priests; nevertheless Christ calls the doings of the same priests the works of darkness. And Peter says: we must rather obey God than men. So verily the Roman Empire, which was of God, came also out of the Sea (as Daniel says also) out of the troublesome world, and even out of Hell, being made great through slaughter, murder, sedition and treason. For the people of Rome with the most part of emperors regarded the Devil and the world, and not God.

### Paragraph 827

And what the Roman empire is today, he now depicts: it has seven heads and ten horns, and each horn has a crown, signifying truly that horns represent kingdoms. We are not here introducing any new or strained interpretation. In chapter 17, the angel himself explains this, saying that the seven heads signify seven mountains or hills, and also kings. Rome is considered to have many hills, but there are seven notable ones: the Palatine, Capitoline, Aventine, Caelian, Esquiline, Viminal, and Quirinal. Propertius explains the same in one verse (which I have expressed in two): “Seven hills the city rises above, overlooking the entire world.”

### Paragraph 828

And therefore is called by the Greeks, “Apocalypse,” meaning “revelation” or “discovery.” And indeed, the city is taken to represent the entire Empire. So there have been many kings and emperors encompassed within the seventh number. But it is certain that the seventh number of kings is exactly found in the history. For at the beginning, when Rome was first built, there reigned seven kings in succession: Romulus, Numa, Tullus Hostilius, Ancus Marcius, Tarquinius Priscus, Servius Tullius, Tarquius Superbus. They were expelled because Lucretia was ravished by the king’s son, and they were ruled by consuls, by ten men, and by dictators, until the time of Julius Caesar, who first usurped a kingly crown for himself. After him reigned Antony and Octavian, called Augustus, Tiberius, Caligula, Claudius, and Nero, again seven. In Nero, the Empire receives a plague. From thence again are accounted seven: Otho, Galba, Vitellius, Vespasian, Titus, Domitian, Nerva. From him, the Empire devolved to Ulpius Trajan, a Spaniard. Therefore, the Roman Empire could not be expressed by plainer marks. To this Empire also Daniel attributed ten horns, both because it was composed of many kingdoms and also because it was dispersed again into many. This will be spoken of in chapter 17. And it is a common thing in the Scriptures to signify kingdoms and power by horns.

### Paragraph 829

And to this kingdom the Lord Jesus ascribes open wickedness, which he calls blasphemy. For he adds: and upon his heads the name of blasphemy, that is to say, whatever blasphemy may at any time be devised anywhere will be found manifest in this Empire, and chiefly in its leaders. For if one beholds the hills of Rome, chiefly the Capitoline Hill, one will find it called by Cicero the dwelling place of the gods, truly because it contained, in a manner, the images of all the gods. For on those hills were seen the temples of Jupiter with all his attributes, and so forth; the temples of Saturn, Juno, Minerva, Mars the avenger, Hercules, Janus, Venus, Apollo; also the temples of Fortune, Health, Victory, Concord, and such others. But if one looks upon the emperors themselves, Gaius wanted his images set up in temples, and the people to swear by his name. Nero blasphemed the name of Christ, and by shedding innocent blood sought to abolish the Gospel. Domitian commanded himself to be called God and Lord. And others also have demanded divine honors, men steeped in blasphemies, and reeking in all wickedness.

### Paragraph 830

Furthermore, by an image composed of diverse beasts, he shows how the Roman Empire increased and obtained such power, and what its manners were. In the seventh chapter of Daniel, the eating of the mountain signifies the monarchy of Greece or Macedon, the bear signifies the Persian monarchy, and the lion signifies the monarchy of the Chaldeans or Babylonians. And it is plain that the Romans, overcoming those nations and subjugating those monarchies to themselves, came to the supreme height of governance. For they subdued the eastern regions chiefly by Lucullus, Pompey, and Crassus; Macedon and all Greece by Paulus Aemilius; a good part of Africa by Scipio and Marius; Egypt by Octavius Augustus; and so forth. And just as they were in religion ungodly, so in other manners they were not unlike wild beasts. For as the Libyan or panther is spotted and of diverse colors, so are the Romans, a collection of many nations, prone to sedition and slaughter. The bear only goes upon his feet, but with them also strikes and catches his prey; so the Romans did nothing else but strike, fight, and take spoils. And just as the strength of a lion is most excellent among four-footed beasts, and the lion’s mouth is insatiable and stinking, so was the Roman Empire most strong, covetous, never contented, and the very essence and corruption of mischief.

### Paragraph 831

And Saint John declares indeed more expressly,[note 55.14] that the Romans possess from the Devil all wickedness, cruelty, and mischief. The Dragon says, “He gave to that beast his power, and that great authority.” This is, in effect, as if he had said: The Devil reigned wholly in the Romans, and the Romans wrought by the Devil all that they did. For the Devil is the origin of murders and lies. Of the Devil’s seat I have spoken in the second chapter of this book. However, we must know that all power is of God; but he, by his just judgment, permits many things to the Devil over the children of unbelief. For Saint Paul in 2 Thessalonians 2, when he had spoken of the mighty working of Satan, by signs and false wonders, with which they should be deceived, those who would not receive the truth, he adds immediately: “Therefore God shall send them strong delusions, that they may believe lies, and be judged all those who believed not the truth,” and so forth. For (as I have often admonished) we must take good heed that we mix not the works of God and the Devil together. Good works are of God, evil are of the Devil. Now, lest any man should marvel why God permits so much to the Romans and the Devil their head, and does not infringe their force for the elect’s sake, Saint John interweaves a heavy fate of the people of Rome and of the whole Empire, which befell them immediately after the first persecution raised against the church of Christ, and after the most shining Apostles were executed, verily to avenge that innocent blood. For he sees one of those heads, as it were, wounded to death. For Nero, who first of the Emperors, stirred up the first persecution against the church, killed himself. And he was the last Emperor of that family. And he left the Empire so afflicted that it was likely to fall to decay. Certain provinces revolted. Galba, Otho, and Vitellius fought among themselves and made civil wars. This Vitellius moreover, drove Sabinus, Vespasian’s brother, suspecting no evil, with others, into the Capital house, and setting the temple on fire, destroyed both Temple and men together, and made all one heap. Neither does Orosius conceal why these things happened, saying: “Rome fell swiftly by the murder of princes and civil wars, for the injuries done to the Christian religion.”

### Paragraph 832

Nevertheless, the Apostle adds that the wound was healed again. For Sextus Aurelius Victor says that Vespasian, in a short time, refreshed the weary world that had been wounded. Here, he says, is the deadly wound that was healed again. For other writers, discussing this matter more fully, describe how Vespasian, returning to Rome, considered nothing more noble or better than to establish and beautify the commonwealth that was sorely afflicted and decayed, to bring order to the provinces and cities that were disordered by tumults and seditious uprisings, to reform the military discipline that had become lax, and to punish offenders. He repaired the city, damaged by old fires and ruins, with new buildings; he rebuilt the Capitol that had been burned; and he erected the Theater in the midst of the city, the most ancient monument of the Empire, and so on.

### Paragraph 833

Moreover, he now touches sharply upon the foolishness and wickedness of the world. And there was great wonder throughout the earth, and so forth. For the world follows present fortune, and judges all things by their good or evil outcomes. For that religion, they say, is most noble, stable, and true, which is celebrated by victories and shines with the ornaments of this world. Therefore, because of the majesty of the Roman Empire, which they held in great admiration, most people received the Roman religion and defended it as sincere. But John, declaring the enormity of this sin, says, “And they worshipped the dragon,” and so forth. He does not say that they worshipped gods, or wood and stones, but that they worshipped the Devil. Idolaters will say that they worship and honor gods, and are not ignorant that images are made of corruptible material? And that the worship they give to them returns, not to those dead signs, but to those of whom they are signs. Thus, truly, all idolaters say. And if you say to them, “You worship wood and stones,” they will quickly answer that great injury is done to them. For they are not so foolish, they will say, as to worship that which they made with their own hands, and so forth. But the Apostle, who knew well these civil explanations and crafty shifts of idolaters, speaking frankly against them, did not respect what they alleged for themselves, but rather what God judges, and the truth of the thing declares, and says, “And they worshipped the Devil or the Dragon.” So Paul in 1 Corinthians, chapter 10. “The things that the Gentiles offer up,” he says, “they offer to demons.” But the Gentiles deny this. But God in this case does not pass judgment on the judgments, intentions, and denials of men, but pronounces according to his own judgment. In Leviticus, chapter 17, he says, “If anyone offers oblations to me other than those I have prescribed, they shall defile yourselves with blood.” Let now the Mass-singing priests cry out until they are hoarse again, “We offer to the Lord God, not to strange gods!” Yet the Lord’s sentence shall stand most true forever, that they transgress with unlawful worship no less than if they committed parricide. As also Isaiah bears witness in chapter 66. The Lord God allows the sincere obedience that we show to his laws; he cares nothing for our inventions and good intentions. Thus, at this present, he shows in few words what the thing is in truth, that all idolaters worship the Devil. If we would at this day esteem these things rightly, we should not contend as if for life and lands, about maintaining images in the church. The Lord Jesus enlighten our hearts and minds to see his truth.

## Sermon 56: ¶ The beast is worshipped, and he blasphemes the name of God, and the Saints of God, and finally makes war with the Saints.

### Paragraph 834

.

### Paragraph 835

And they worshipped the beast, saying, “Who is like unto the beast? Who is able to war with him?” And there was given to him a mouth to speak great things, and blasphemies: and power was given unto him, to do forty-two months. And he opened his mouth unto blasphemy against God, to blaspheme his name, and his tabernacle, and them that dwell in heaven. And it was given unto him to make war with the saints and to overcome them.

### Paragraph 836

He states that the world worships the dragon; now he adds that it worships the beast. If it is said that the beast represents the empire, some might wonder how the empire could be worshipped. But we will explain briefly how people worship the empire: by accepting its decrees, rituals, and superstitious ordinances, and by being entirely dependent on it. And there were not a few at that time who, in order to gain favor with the Roman Empire, denied the faith of Christ, revolted from the church, and joined themselves in religious practices and sacrifices to the fellowship of the empire. They truly worshipped the beast. Furthermore, that which is due only to one God, the Romans attributed to their empire. But whoever ascribes divine properties to anything, does in truth deify and worship it. And the properties of God are these: to be without equal or end, that he alone is greatest and best, immortal, eternal, most powerful, most invincible. For so the Prophets say: Who is like unto you, O God, in heaven and in earth? Who is as you? Who can resist God? But the Romans attributed all these things to their emperors and to their empire, saying, as Saint John also recounts: Who is like unto Rome? Who is able to war with it? They called their emperors gods, best, greatest, most powerful, and most invincible. They called the empire itself eternal. You can still see these things in ancient authors and coins. Therefore, those who were not ashamed to attribute these things to the Roman rulers and kingdom rightly worshipped the beast. And what other thing is done today, while for the favor of emperors, kings, popes, and their realms, the truth is denied or twisted according to the affections of people? These also worship the beast.

### Paragraph 837

Now the beast is also given a mouth, speaking great things and blasphemies. [note 56.2] Of blasphemies we shall speak more fully later. But since the Roman Empire obtained the greatest victories and held the most magnificent triumphs, it seems to have been given opportunity to boast proudly of those victories, and to claim as their own things that were in truth accomplished through the power of God. And without doubt, the greatest and most unrestrained boasts of the Romans remain, that they are conquerors and lords of the world. But such pride was severely punished in Nebuchadnezzar the king, as you may see in chapter 4 of Daniel. Saint Peter affirms that God resists the proud and gives grace to the humble. God hates the arrogant and takes away their names from the earth.

### Paragraph 838

And where some person might demand,[note 56.3] what will be the end of injuries, pride, and finally, intolerable authority and blasphemies? John anticipates this and says, “And power was given him to do, that is to work violence, for forty-two months,” that is to say, for such a time as seems good to the Lord. Although he wishes the time to be unknown to us, it is known to him, so that the faithful may promise themselves that this evil will endure only a few months. I have reasoned about this number in Chapter 11 and Sermon 46, and have shown in previous passages that those numbers are equivalent: namely, the thousand two hundred and three score days, the forty-two months, the time, two times, and half a time. God, therefore, admonishing us as it were by a riddle, does not want us to curiously inquire after times which he has kept in his own power; it is sufficient for us that he has assigned all things within their proper limits.

### Paragraph 839

A plentiful treatise of Roman blasphemies. [note 56.4] First, he says by a trope, he has opened his mouth: whereby he has signified his boldness, and liberty, yea licentiousness of speaking. For we say he would not once open his mouth, when we signify any man that will not speak frankly. But the Romans, and companions of the Roman superstition, blaspheme God in three manners. For first, they blaspheme the holy name of God in this, that they do prefer their false gods and their superstitions to the true God, to the true and most holy religion. For where they did admit in the city of Rome the gods of all nations, and their religions, the religion of the only God of Israel they utterly refused: for that they understood how he would be worshipped alone, and by no other rite than that which he himself had prescribed. But they had rather retain wickedly those their many gods, and their religion, although most absurd, than to commit themselves into the tuition of one, and to receive a moderate and simple religion. Augustine, *Aurel.* I recount not now the blasphemous words of them, uttered against the true God, about that time chiefly, when Vespasian and Titus triumphed, after the Jewish war finished, both of the City destroyed, and the people of God overcome. There were carried about in the triumph the holy vessels of the Temple, and even the God of the Jews, as vanquished and bounden, was seen led into the Capitol house, to make his supplication to their great god Jupiter, as it pleased them. Whereupon we understand that the name of God was no whit less outrageously blasphemed at that time than it was in old time of the Palestinians or Philistines, what time they set the Ark in the temple of their god Dagon; likewise of Rehoboam [?] and Sennacherib, moreover of Belshazzar, King of Babylon, in the 5th chapter of Daniel. But the offenders are found out at the last.

### Paragraph 840

Secondly, the Romans blasphemed the Tabernacle of God. That same ancient Tabernacle of the people of Israel was not only the office, or place of religion and worship, but also a token of God’s presence. For God is now present in the midst of his Church, a figure of whom the Tabernacle of witness represented. But the Romans called the Christian church wicked, foolish, seditious, harlike, and detestable; and they most grievously persecuted it, seeking to destroy it by all means, and bent their whole power to this end.

### Paragraph 841

Finally, they also blasphemed the heavenly inhabitants, the blessed souls of saints, prophets, and apostles, whom they called wicked, seducers, peace breakers, blasphemers, heretics, and sinful persons. For at this time, while John wrote these things, several apostles under the Roman Empire had already been executed and slain, considered plagues of the world, and their memory and doctrine condemned. But from this, you perceive how displeased God is when anyone reviles godly preachers and holy ministers of the church. For the Lord takes the reproach spoken as if against himself. There remain even today blasphemies of this sort, as found in Cornelius Tacitus in the twenty-first book of Augustine, written against Moses and the people of God.

### Paragraph 842

Moreover, God permits the beast to wage war against the saints and to overcome them.[note 56.6] For the Roman Empire, until the time of Constantine the Great, instigated ten most severe persecutions against the church. You may read of this in Eusebius, bishop of Caesarea, and Orosius in the history he wrote to Jerome. And this passage especially pertains to the instruction and comfort of the church. For the Lord also in the Gospel prophesies of the church’s destinies, for the consolation and information of the faithful, as appears in chapters 15 and 16 of John. And how the saints are overcome I explained in chapter 11. The Lord Jesus preserve his church. Amen.

## Sermon 57: ¶ Of the power of the Roman Empire, and who worship the beast: and of the destruction of Rome, and the Roman Empire.

### Paragraph 843

.

### Paragraph 844

And power was given him over all relatives, language, and nation, and all who dwell upon the earth worshipped him, whose names are not written in the Book of Life of the Lamb who was slain from the foundation of the world. If anyone has an ear, let him hear. He who leads into captivity shall go into captivity; he who kills with the sword must be killed with the sword. Here is the perseverance, and the faith of the saints.

### Paragraph 845

The Apostle, through the revelation of Christ, also speaks of the power and majesty of the Roman Empire. The Roman Empire was indeed of greatest power in the time of Octavian Augustus, also in the time of Domitian, and in the reign of Trajan, also under Hadrian, Aurelian, Diocletian, and Constantine. The greater part of the inhabited world obeyed, namely all of Europe in a manner, Asia, and Africa, as both Latin and Greek histories testify. However, the Lord warns us not to curiously search the counsels of God, being inquisitive, why God gave so great power to the Romans, whom he knew would abuse it to the oppression of Christ’s Church. For where he says that the power was given to Rome, he stills and appeases all murmurings. For empires are of God. But he is most wise, righteous, and holy. Therefore, where he made the kingdoms of the world subject to Rome, he did it wisely, justly, and holily. In that the Romans corrupt God’s ordinance and commit themselves to be governed by the Devil, it comes of evil.

### Paragraph 846

Let our arguments here cease, for the wise man also says that wicked men and hypocrites rule because of the people’s sins. And that he recounts lineages, languages, and nations, he does in imitation of the Prophet Daniel, who by such phrasing is accustomed to signify a very extensive and powerful empire.

### Paragraph 847

But what does this have to do with us, or what benefit (do you say) comes to us from it, that the Roman Empire is so widely extended throughout the world? Truly, we see how this prophecy has accurately predicted everything that has come before; therefore, there is no room left for doubt about the things that follow. Let us consider moreover that the most powerful kingdoms, which appear invincible to people, may be dissolved by God without any difficulty. Let us therefore learn to fear God, and to walk in his commandments, and to despise these earthly things.

### Paragraph 848

Now he declares more expressly who will worship the beast:[note 57.2] for he said that people in the world will be captivated by the beast and will worship it. He now clarifies this, placing the word "worship" so that it may be understood to apply both to those who are present and to those who are to come. For he speaks not only of people of his time, but of all who, captivated by the admiration of the empire and its majesty, will either deny or despise the faith of Christ. And he says that all who dwell on the earth will worship the beast—and lest anyone refer this absolutely to all, as though none of the true worshippers of God were included, he adds, "whose names are not written in the book of life of the Lamb," namely, the reprobates, not the elect: the unbelievers. I say, those who despise the word of the gospel, disdain to hear it, and are rebels to Christ. Arethas, the expositor, says they dwell on the earth, those who are unconcerned about heavenly things or the glory that is there, or who devote themselves to earthly pursuits and live a beastly life, according to the same. Thomas Aquinas also brings a testimony from Jeremiah 17: They who depart from me shall be written in the earth. For they have forsaken the fountain of living waters, even the Lord himself. Of the book of life I have spoken in chapters 3 and 5, and will speak of it again in chapters 19 and 20 of the Apocalypse.

### Paragraph 849

Hereunto he adds a noteworthy observation, in the manner of the Apostles, who always, whenever they have occasion to celebrate and explain Christ and the mystery of his redemption, do so. John says that the Lamb has been slain and offered from the beginning of the world. And it is without controversy that by the Lamb is understood Christ.

### Paragraph 850

It is therefore demanded, how he was slain from the beginning of the world. Many torment themselves at length as they expound, that Christ was slain in Abel, and in all saints, by participation, not by passion. Certainly we may not expound this place after the letter. For Christ could not be slain before he was born. Moreover, the Apostle affirms that Christ, who was from the beginning of the world, has not been slain oftener than once. Read what he says in chapter 9 to the Hebrews. And yet the most true word of God cannot be contrary or repugnant to itself. Therefore, we say, after the common rule of expounding the Scriptures, that the signs have the names of the things signified. For the Lamb was called a Passover, whereof it was a sign. Circumcision was called the league or covenant itself; sacrifices are named sins. So, verily, from the beginning of the world, sacrifices were slain, which were symbols or signs of Christ to be incarnated and offered up once for the cleansing of sins. We understand therefore by this testimony of Christ, that all the sacrifices of the ancient fathers were sacraments of Christ, and that the redemption of Christ has from the beginning of the world been of efficacy to all the faithful. Therefore, this place is notable and worthy to be observed. Hitherto pertains the Apostle’s testimony in 1 Corinthians 10:1: that all our forefathers have eaten of the same spiritual meat with us, and drunk of the same drink, and that they drank of the rock following them, which was Christ.

### Paragraph 851

And up to this point, he has spoken of the majesty of the Roman Empire, its blasphemies, and its sins. Now follows a discussion of the destruction of so great an empire and the punishment of sins. However, this will be addressed again in chapter 17.

### Paragraph 852

And with an acclamation, most commonly used in the gospel, and as it were peculiar to Christ, he stirs up all his listeners, and cries out, “He that has an ear let him hear.” Truly it was to men a wonder, and seemed incredible, that so great a majesty could fall; but yet it falls. The faithful also marveled what should be the end of blasphemies, slaughters, injuries, abominations. Moreover, the doctrine that follows is notable, excellent, and worthy to be kept in memory. Therefore he stirs up all men to attentiveness, and then he says: “Whosoever shall lead into captivity, shall go into captivity; whosoever strikes with the sword, &c.” For in such sort he declares the destruction of Rome and the Roman empire, that he confirms with all the justice of God’s judgment. And also with a marvelous brevity of God’s sentence given or pronounced against Rome, he partakes of that immeasurable power. And this is both by the law of God, by the law of nature, and by the law of all nations received as a thing most just, that every man should look to have the same done to him that he doth to another. For to this belongs the sentence rehearsed of Noah in Genesis 9: “He that sheds blood, his blood shall be shed.” The same is repeated in Leviticus 33: “Woe to him that spoils, shall not thou be spoiled?” A testimony whereof is Nineveh with the prophet Nahum, and Babylon with all the prophets. Therefore has the Lord taught in the gospel, “Whatsoever you would that men should do to you, do you the same unto them also.” With what measure you mete unto others, with the same shall others mete unto you again. Whosoever strikes with the sword, with the sword shall perish. Therefore it is most reasonable, that since Rome has spoiled the whole world, and injured all nations, and made cruel war upon all men, it should be again of all nations invaded, spoiled, torn, and trodden under foot. Let us mark this judgment of God, and let us fear God, and do good unto men. For here is sentence given against all men that do injury to their neighbors, but especially those which invade innocents with unjust wars, and which they are hired to make, &c.

### Paragraph 853

And here we must repeat something from Histories,[note 57.5] whereby the truth of this prophecy may be better known and understood. When the most excellent Prince Constantine received the government of the empire, as if abhorring Rome, he built Constantinople and made it the seat of the empire. And from that time the majesty of Rome began to fall into ruin. Under the emperor Gratian, a prince most witty, the Barbarians were a great terror to the Romans, whereupon Gratian made a league with them. Stilicho, father-in-law to Honorius, a Vandal by birth, diminished the wages of the Gothians and other league fellows of the people of Rome; for which cause they took armor. Yet being pacified again, they were stirred up afterward through the malice of Stilicho and of Duke Saule, and under the conduct of Athalaric their king, they marched to Rome, laid siege to it, and besieged it for the space of two years, at length taking and spoiling it. This siege and spoil Saint Jerome in his Epistle bewails much. Orosius writes much and Christianly hereof in the 29th chapter of the 7th book of Histories. It is reported that Rome was taken the first day of April, in the year 412. Yet the Gothians immediately leaving the city, removed into other places thereby; nevertheless, being again inflamed with fury, they returned, and under their captain Athaulphus, they plagued and spoiled Rome worse than they did before. The king had determined, extinguishing the name of Romans, to have called the city Gothia, if he had not been dissuaded by Galla Placidia, daughter to Honorius. A few years after, Rome was taken again by Gensericus, king of the Vandals; and that which was enriched and replenished with the robberies of all nations was by fourteen days together emptied clean. After came Odacer with the Germans; and putting down the name of emperor, reigned over the city himself as king for the space of 15 years. Whom Theodoric of Verona expelled and slew. And there reigned with his Ostrogoths about 50 years. Then it was recovered by Belisarius for Justinian, emperor of Greece, but to the utter destruction of Rome. For Totila, king of Gothia, discomfited both the Greek and Roman army at Placentia; after he besieged Rome, scaled, took, sacked, overthrew, and set it on fire. The city burned thirteen days. Neither was there any man in it for the space of forty days. Read the 4th book of Sabellicus the 8th Aeneade. Perhaps I shall discourse more at large of the destruction of Rome in the 17th chapter. Wherefore within the space of 136 years, Rome came seven times into strangers' hands, and was sacked most cruelly, and fell by the edge of the sword, and was led into captivity:[note 57.6] which hath long stricken with the sword, and led away all nations prisoners. This was the just judgment of God.

### Paragraph 854

And Saint John adds a doctrine, how the faithful should behave in such great troubles and adversities. That is to say, while the Romans rule and rage, also in those bloody and cruel changes and destruction of the Roman Empire, the saints will need patience, or perseverance, and faith. These two virtues will preserve the faithful, so that they do not perish as well. Concerning patience, the Lord speaks in Saint Luke, chapter 21: “In your patience you will possess your souls.” Blessed John speaks of faith: “And this is the victory that overcomes the world, even your faith.” Impatience and incredulity have led many into the denying of the faith, to idolatry and to all ungodliness. Thus, we also learn how to arm ourselves in our days against all ungodliness. The Lord deliver us from evil. Amen.

## Sermon 58: ¶ Of an other beast, which comes up out of the Earth: that is to say, of Antichrist.

### Paragraph 855

.

### Paragraph 856

And I saw another beast rise up out of the earth, and he had two horns like a lamb, but he spoke as the dragon.

### Paragraph 857

The Apostle Paul plainly testifies that the things written were written for our learning, so that through the patience and consolation of the scriptures we may have hope; therefore, we must also apply these things to our present circumstances. For Christ, the Lord of all, when he foresaw how greatly Satan would, through his chosen members, the old and new Roman Empire, afflict the church, would have us diligently admonished of every thing, so that all afflicted persons should learn patience from this, and conceive comfort and hope, and not be discouraged by the heavy burden of evils. Just as he has therefore diligently described the old Roman Empire, and showed as it were pointing with the finger what mischief it should work to the church, and admonished all to have faith and patience; so will he from hence forth describe papistry or Antichristianism, in which description he sets forth before our eyes whatever the saints shall suffer; that being warned beforehand, they may endure persecution more manfully and yield less to misfortunes.

### Paragraph 858

And in goodly order he begins to set forth Antichrist after the Roman Empire is torn and taken away. For Daniel says that a little horn should arise among the ten horns, and that three of those horns should be pulled down, plucked off, and cast away, so as to attain to great power. For he signifies that the Roman Empire, being divided and brought now to decay, Antichrist shall arise, who will procure to himself a new and counterfeit empire. And Paul says also that Christ shall not come to judgment until Antichrist has gone before; and that he shall not come neither unless this be first taken away, which hinders and prevents his coming. The which Jerome and other holy expositors understand of the Roman Empire, which must be plucked up and taken away, and then Antichrist shall arise. But the majesty of the empire was destroyed about the year of our Lord 480, when Odacer invaded Rome. For from that time, for the space of 300 years and more, there was no emperor of the West after Augustulus. And besides this, under the emperor Justinian, Rome was burned and laid waste by Totila. Since which time the bishops of Rome have begun to look aloft and to think upon a new kingdom.

### Paragraph 859

And therefore the Lord says that this beast arises from the very earth. The kingdom of our Lord Jesus Christ comes from heaven and brings things to heaven. Papistry comes neither from Christ nor from his doctrine, but arises from the earth: that is to say, from evil means, ambition, avarice, treason, and cruelty. What ministers of the church Christ ordained is easily perceived by the Gospel of Jesus Christ. That he forbade them government, supremacy, authority, and majority (as they term it) appears in the 18th and 20th chapters of Matthew and the 22nd chapter of Luke. Therefore, the Acts of the Apostles and the doctrine of Peter testify that Peter was a minister, and not a lord of the Apostles, much less a prince of the city or empire of Rome. For they lie boldly that say how Rome and Italy are the patrimony of Saint Peter, given him by the Lord. At first, the Apostles and apostolic men, ministers of churches, governed the churches equally, neither did one take upon himself more preeminence than another. Which thing I am able to prove by many testimonies of ancient writers, if need required. About the time of the Council of Nice, and a little before that time, when churches were greatly multiplied, Metropolitans were ordained and customarily received, instituted indeed by a laudable [?but yet man] ordinance: that is to say, in a certain province or head city a bishop or pastor was ordained, who should have as it were the oversight of the rest, and should serve for the calling of Synods or assemblies. Yet it was then diligently provided that he should not be called Primate, lest any man should think himself preferred before others in power, but in order. Neither was the bishop of Rome at that time exalted above all others, but there were diverse Metropolitans, whereof the bishop of Rome was one. The Council of Nice confirmed that same custom and would have it ratified. Socrates in his Ecclesiastical History, book 5, chapter 8, recounts many Metropolitan churches in Asia. Saint Jerome to Evagrius and in an epistle to Titus says plainly that in old time churches were governed by the common counsel of priests or elders, and that then bishops and priests were all one: After, by the custom of the church, not of the verity of the Lord’s ordinance (I rehearse Saint Jerome’s words), bishops were preferred before priests, yet must they govern churches together.

### Paragraph 860

And of that same custom, yea rather of the abuse of the custom, Antichrist had his beginning. For Boniface, Bishop of Rome, began first to assume dominion over the churches of Africa. But he was immediately restrained by the six African councils, where Saint Augustine is also said to have been present. After that, the Bishop of Constantinople also began to claim supremacy for this reason chiefly, that Constantinople was then the imperial palace and chief city of the empire. However, certain bishops resisted him; among whom was Leo, Bishop of old Rome. There remain certain letters of his to the Emperor of Constantinople, to the bishops of the East, and to others. So was this trouble appeased for that time also. But straightway another Bishop of Constantinople, blinded with ambition, demanded anew to have the supremacy given him. Whom Pelagius and Gregory, Bishops of Rome, withstood. And this latter so impugned the supremacy of the Patriarch of Constantinople that he did not hesitate to call him the vaunter of Antichrist, who would usurp the title of general bishop. There remain not a few letters written of this matter in his register.

### Paragraph 861

Nevertheless, a few years later, when the bishops of Rome were greatly afraid that that dignity would be given to the bishops of Constantinople, Boniface III obtained from the emperor Phocas—a parricide—that the bishop of old Rome might be considered the universal bishop, and Rome the head of all churches. This constitution established the Pope in authority [note 58.6], so that he was now regarded by many western bishops as Apostolic, and many matters were brought before him for judgment. Through this, he gained the favor of many princes, chiefly of France, by whose aid he drove out of Italy both the emperor of Greece and the kings of Lombardy, and brought Rome and the best and most flourishing parts of Italy under his own dominion. Thus I say, out of the earth comes up the second beast.

### Paragraph 862

Furthermore, Christ calls the Roman papistry a beast, because in avarice, covetousness, tyranny, cruelty, and even in bestial behavior, it differs nothing from the old beast, of whom I have spoken before.

### Paragraph 863

Up to this point, we have discussed the origin of the Antichrist or Pope, and of the new empire. Furthermore, Saint John proceeds to describe that second beast vividly, so that we all may know and avoid it. And first, he reasons about the power of the Antichrist.

### Paragraph 864

[note 58.8]That other beast, he says, had two horns: and he adds, like a lamb. For of them is spoken in the fifteenth chapter of this book. And the Lord signifies the priesthood and kingdom, which the Popes usurp for themselves, affirming that power is given them in Heaven and on Earth, in spiritual matters and temporal. For therefore they give in their arms two keys, that is to say two horns. They boast that they have two swords. Of these blasphemies, he who will be fully instructed, let him read the words of the beasts of Boniface VIII, in the sixth Decree of Majorities and Obedience. One, in Clement V, second book of oaths; finally, Gregory IX, or rather the first book of Innocent III, title on Majorities and Obedience. All histories make mention that Boniface VIII died in the year of our Lord 1300, instituted the first Jubilee, and in the same openly before the people showed in the way of ostentation the Pontifical and Imperial majesty, while on one day he appeared in the apparel of a Bishop, on the other having put on purple robes showed himself to the people like an Emperor. They carried before him two swords. And he himself cried, “Lo, here are two swords,” as though he should point with his finger to the whole world, that he and certain of his predecessors and all his successors were that two-horned beast. What shall we say that all bishops by him consecrated wear upon their heads miters or two-horned caps? Unless therefore we be blinder than was Tiresias, we see with our eyes who is that great Antichrist.

### Paragraph 865

And here we must observe that he does not [note 58.9] say that those are the horns of a lamb. For Christ still keeps both the priesthood and kingdom with the faithful in the church; neither does he resign them to any other, he has appointed no Vicar. For he continually executes, at the right hand of the Father, the offices of both King and Bishop, and this all faithful feel with joy. He says therefore, like a lamb. For the Pope will make all men believe that he has received from Christ priesthood and empire, that he is Christ’s Vicar; where he is nothing less. He boasts everywhere that he is the great shepherd and has received the keys of the Kingdom of Heaven; and that from the very Lamb of God, through the Apostle Saint Peter; and therefore that all bishops are subject to him, finally, all kings, princes, and people.

### Paragraph 866

He proceeds to show moreover, what the talk of Antichrist is, what is his doctrine, and what is his speech. He spoke, says he, as did the Dragon. The Dragon is the Devil, as was shown plainly before. Therefore, he ascribes to Antichrist, or Popery, a diabolical doctrine, or a devilish mouth or tongue. We must see, therefore, how the devil speaks, that we may understand rightly how Antichrist speaks. In Paradise, he so tempers his talk that he calls into doubt the certainty and truth of God’s word, and by that occasion places his own word instead of the word of God. Is it so, says the Devil, under peril of your life, has God forbidden you to eat of the fruit of the tree of knowledge of good and evil? Yea rather, if you eat thereof, you shall be made like unto God. And after the same sort, Antichrist in his Popery brings the truth of the Scripture into doubt, which by all means possible he discredits as imperfect, maimed, obscure, and doubtful. And immediately upon that occasion, he brings in his traditions and decrees, with which he may patch up what he contends lacks in the Scriptures. But in his traditions, he affirms things contrary to God’s word and so deceives men. And all men know, who have any skill of Popish matters, that the first and chief principle and foundation of Papistry is that the Scriptures are imperfect and obscure, and therefore need traditions. Moreover, the Dragon speaks openly against the laws of God, and so does the Pope manifestly. God will be worshipped alone; the Pope adds saints to him. God forbids idols and idolatry; the Pope commands them plainly. God will have his name to be sanctified, and to be sworn by only; the Pope, by dispensing with oaths, pollutes the name of the Lord and commands us to swear by the names of God. God commands us to keep holy the Sabbath day; the Pope brings this into contempt, sets forth his own holy days, and makes double feasts. God commands us to honor our parents; this the Pope abrogates and commands greater regard for abbots and abbesses. God commands, “You shall not kill, you shall not commit adultery, or steal.” The Pope grants most ample indulgences and pardons to his soldiers for rash wars made at his will and pleasure; he spoils with his sacrileges all churches; and he, with his maiden priests, fills all the world with adulteries and whoredoms, to speak at this time of nothing else more filthy. And where God forbids lies and false witness, the Pope and his whole doctrine (which he sets forth besides the Scripture) are sown with lies. He not only dispenses with false witnesses but permits also to break safe conduits and public faith given, and if he hates the Prince, absolves the subjects from their fidelity and obedience; gives liberty to all lusts, and makes laws which nourish the desires of the flesh. And which shall be the Dragon’s mouth, if this is not it? The Dragon moreover is said to have spoken and said to the Lord: “All these things I will give thee (for he showed him the kingdoms of the world) if thou wilt fall down and worship me.” What other thing speaks the Pope? Does he not enrich his obedient children with the riches of this world, especially those who will fall down and kiss his feet? I suppose the Devil would never be so shameless as to offer to the Lord his foot to kiss; but that beast, in the sight of God and his angels, and of all the world, dares put out his foot, marked (not without great mockery) with the sign of the Cross, and proffer it to be kissed by all the children of God. I cannot bring forth the horrible and innumerable blasphemies out of the decrees and decretals. For I am ashamed of such ungodliness. Who therefore will not acknowledge that Satan himself does in this beast reign and rage? God shortly confound the same. Amen. Amen.

## Sermon 59: ¶ Again of the power of Antichrist, and howe the former beast is worshipped.

### Paragraph 867

.

### Paragraph 868

And he does all that the first beast could do in his presence. And he causes the earth and those who dwell therein to worship the first beast whose deadly wound was healed.

### Paragraph 869

Again he reasons about the power of the second beast, or Antichrist, and of papal authority. He does or executes, he says, the power of the former beast; that is to say, he exercises the same authority that the old Roman Empire exercised. Where he adds, in his presence: Arethas expounds, following immediately after, and even imitating the same. But what power and authority they exercised, I declared before in his place, about the beginning of this chapter. Therefore, just as the Roman Emperors considered all kingdoms and provinces to be theirs, and to belong to them, so the Bishops of Rome claim that all realms are theirs. I give nothing here to affection or hatred. There recently came forth a book printed at Lyons, by Augustine Steuchus, keeper of the Pope’s library, in which he recites from the register of one Gregory (I suppose Gregory VII) all the kingdoms of Europe, Spain, England, France, Denmark, Hungary, etc. The ownership of which belongs to that seat of Rome, the use to the Princes, clients of that see. Often the Popes have attempted to bring the kingdoms of the East under their control also, and that under the pretext of the holy war and recovering the Lord’s sepulchre. And just as the old Romans troubled with continual war the nations that did not acknowledge or obey the old Roman Eagle, so the See of Rome in our time, and in the memory of our forefathers, has put to business and trouble those kingdoms, and nations, and people, who sought to revolt and would not acknowledge those double keys, that is to say, two horns. For who does not know with what cruel wars he vexed in times past the land of Bohemia? Who does not know what Germany and England have suffered in former years? So truly the second beast exercises boldly the tyranny of the old beast. The old beast issued proclamations concerning religion, and the payment of tributes and customs, and in this way impoverished almost all realms, their riches being brought to Rome. And what other thing does that seat do today? What has it done now, to reckon the least, these five hundred years? Who therefore does not see that the second beast exercises most abundantly the power of the first beast? A certain man composed verses in Latin mocking the covetousness and deceits of Rome; and where Rome magnifies herself to be the head of the world, which in Latin is *Caput*, thus he says: If *Caput* comes from *Capio*, which signifies to take, then Rome may well be called so, which does nothing forsake. If you decline *Capio*, *Capis*, and come to the roots, her nets are large and cannot miss, to catch both all and some.

### Paragraph 870

He adds another point, that this second beast arranges for those who live on Earth to worship the first beast. This, without doubt, we see fulfilled in the Papist kingdom in two ways. First, the Papists have secured such authority and reverence for the Roman Empire, which they call both sacred and holy, that as many as live today, when they hear the name of the Roman Empire spoken, imagine a certain divine thing, brought to them from Heaven. I grant that there have been many truly noble princes, godly, and altogether worthy of praise, within that same Empire: such as Constantine, Augustine, Gratian, Valentinian, Theodosius, and many others. I grant that under these and similar rulers, the Empire was holy, and in truth an empire of Christ. For Christ was acknowledged with true faith; yet we see how the Lord Jesus has nevertheless, as Daniel has done also, called that Empire a beast, doubtless figuratively and for the tyrants. Therefore, we must wisely and justly attribute to each ruler what is his, and not without consideration embrace and reverence that bloody Empire as sacred and holy. And we have declared before, in what manner kingdoms are of God, and how far their works are to be approved that are in kingdoms. And of this there shall be spoken a little later yet more fully.

### Paragraph 871

Secondly, the second beast causes men to worship the first, in this chiefly,[note 59.3] that Papistry has brought back the pagan manner, the names only changed. For I told you before, that the first beast was worshipped, in this that simple people received the Roman religion and worshipped idols. The pagans did truly confess the almighty high God, but they joined to him many gods, to whom they submitted elements, diseases, arts, countries, cities, the members and parts of man, and such other like things. Aeolus was god of the winds, and Neptune of the sea, Pluto ruled in the earth, Mars was god of war, Minerva and Apollo of arts, Aesculapius over diseases, Hercules and many more. Venus was lady of love, and the goddess Juno of marriage. Neither was there any member in the body that did not have his god: so had all countries and cities their saving gods, and every house its domestic gods. To them afterward they framed idols, that is tokens and memorials, which might bring those heavenly gods into the memory of the earthly dwellers. They built chapels and churches, they instituted priests, holy days, rites, and ceremonies. These things are found in the books of the gentiles, and in our histories, and also in the writings which have refuted the pagan idolaters. But in the papal kingdom at this day, the names being only changed, who can deny that the same cult, the same religion, indeed very superstition is not renewed? Of these things I have treated at length in my book *De origin erroris*. The Papists teach that the saints in heaven reign with God, and that to them are subject sicknesses, arts, limbs or members, cities, and all things, and must therefore be called upon and worshipped. Saints are expressed and represented by images, to these images are erected altars and churches: briefly, it is done to them that was done to the gods and idols of the pagans. Who therefore understands now that Antichrist has procured that the first beast might be worshipped, that is to wit, might be powerful again, and that the old idolatry and superstitious worship might be renewed and frequented?[note 59.4]

### Paragraph 872

And as we have read before, it is said that all who dwell on Earth worship him, those whose names are not written in the book of life from the Lamb. He says this plainly here as well, and he causes the Earth and the inhabitants of the Earth—that is, those who seek and regard only earthly things—to worship the first beast. Not all are polluted with papal idolatry. To this pertains the noble history of Leo the third, Emperor, and Gregory the second, and of other Popes, through whose wickedness idolatry was again brought into the church; of which I wrote at length in my work *On the Origin of Errors*.

### Paragraph 873

Neither is this annexed here without a mystery, whose deadly plague was healed. For he seems to compare together the first and second beast, and to show their likeness. And I told you how many men at the first were kept still in the Roman errors and idolatry, because the gods, by Vespasian’s means, were said to have preserved the commonwealth, which else with civil wars was as it were brought to ruin. Finally, we read in histories that the Empire of Rome has many times received deadly wounds; but yet, immediately, through the wisdom and valor of some noble men, the gods (as they speak) so willing, have been healed again. In that number are Lucius Septimius Severus, Valerius Aurelianus, C. Aurel. Val. Diocletian, and so forth. By whose lucky success, triumphs, and victories restoring the Empire, many have been moved to say, who does not see that Rome shall be eternal, and that the Roman religion is to the gods most acceptable, and that the emperors also and public weal is endowed with a certain deity, and is to be honored? After the same sort, the kingdom of the Pope or Antichrist, having tried diverse chances, has very often escaped out of desperate dangers. Force and policy have afflicted it, and also the religion of Henry III, Emperor, and of his son Henry IV, Frederick I and II vexed the popes. There were also other mighty princes who inflicted mortal wounds to the See of Rome.

### Paragraph 874

Again there were bishops of Rome who with great cunning have healed their wounds.[note 59.6] Among these were Gregory VII, Urban II, Paschal II, Calixtus II, Alexander III, Innocent III, Honorius III, Gregory IX, Clement IV and V, Boniface VIII, John XXII, and many others. But was that situation not in greatest danger in times past, when three popes were created at once: one residing in Rome, the second going to Avignon in France, and the third living in Spain? But all three were deposed by the power, diligence, authority, and policy of Emperor Sigismund and the council of Constance. That deadly wound was fairly healed in Martin V. And this felicity, and the restoration of the papal kingdom, persuades many effectively that popery is of God, and the popish religion to be most certain and true: as that which has often been attacked by mighty princes might indeed be shaken, but never yet overthrown. The acclaim of all the Romanists is known: the ship of St. Peter is indeed tossed with storms, but cannot be drowned. But Daniel himself also prophesied that this would come to pass, saying: “He shall prosper and shall do what he will, and shall kill the strong and holy people at his pleasure, and guile shall be directed in his hand.” Those who are offended at this day by the felicity of that chair of pestilence and the beast thereof do not mark these things. Therefore, just as the days of mourning and sudden destruction came upon old Rome and utterly destroyed both the city and empire, even so we shall hear in the 17th and 18th chapters that Babylon shall have her fatal destinies. The Lord Jesus confirm us in the faith of Jesus Christ, and deliver us from the guiles, lucky success, and felicity of that Roman Antichrist. Amen.

## Sermon 60: ¶ Of the signs of Antichrist, and Image of the beast of him raised.

### Paragraph 875

.

### Paragraph 876

And he did great wonders, so that he made fire come down from heaven into the earth in the sight of men, and deceived those who dwell on the earth by the means of those signs which he had power to do in the sight of the beast, saying to those who dwell on the earth: that they should make an image unto the beast, which had the wound of a sword, and did live. And he had power to give a spirit unto the image of the beast, and that the image of the beast should speak. And should cause that as many as would not worship the image of the beast should be killed.

### Paragraph 877

He proceeds most diligently to describe Antichrist and his kingdom, which so greatly challenges the faith of Christ and afflicts his church, so that he might be known and avoided by all people.

### Paragraph 878

He says now he shall do great wonders,[note 60.1] by which he understands miracles. Some of these are true, and some are false. I call those true miracles which are done in reality and are not by any crafty juggling counterfeited, and which draw me to the truth and set forth the truth. Of this sort, without doubt, were the miracles of the Prophets and Apostles, holy Martyrs, and chiefly of Moses and Christ. These do good to people, do not harm, nor empty people's purses; rather, they glorify God and make the truth to be believed, by drawing people only to God as to the fountain of all goodness. So John testifies of the Lord’s first miracle done in Cana of Galilee, and says: “This is the first sign that Jesus did at Cana in Galilee, and showed his glory, and his disciples believed on him.” This sign was true, and free from any suspicion of trickery; it was a benefit bestowed upon poor people newly married, by which God was glorified. His disciples were moved by this, and the Spirit of God working also inwardly, believed on Christ. So do all true miracles testify Christ to be healthful and beneficial, and therefore alone to be called upon and worshipped. So do John and Peter interpret the signs or miracles that they themselves wrought, as recorded in Acts 3. And of such miracles we have great abundance in the Evangelical, Apostolical, and Ecclesiastical history: and they have no other end than that we must believe in the Son of God, as which alone gives life and all good things.

### Paragraph 879

And I call false those signs that are done through devilish craft or enchantment,[note 60.2], or by the crafty juggling and subtle trickery of wicked men: such as are the deeds of witches and enchanters; such as were the wise men of Pharaoh, and Simon Magus; and those of which mention is made in Deuteronomy 13. Finally, such were the miracles of the Friars, burned at Bern in Zurich: and it is greatly to be feared that most of the miracles of all Monks and Hermits have been of this sort. Likewise, those are called false also, which although they are done in deed, yet bear witness to a lie against the truth, confirming the Pope to be head of the church, that images are to be worshipped, that we must pray unto Saints, and go on pilgrimages for religion, yea rather for superstition’s sake, that we must worship a new God lying hidden under the form of Bread and Wine, that God must be honored with vows and monastic practices, and if there be any other thing of like sort. With such deceptive signs are filled at this day all temples, churches, and chapels. These have persuaded many even wise men, and do also at this day. Which thing the Lord prophesied should come to pass, saying: there shall arise false Christs, and false Prophets, and shall show great signs and wonders, that even, if it were possible, the elect should be brought into error. And Paul also: the coming of Antichrist, says he, shall be after the working of Satan, with all power, and signs, and lying wonders, and the rest, which is read in 2 Thessalonians 2. And we know that many Bishops of Rome have wrought signs: but that same is not so excellent, but that the Bishop of Rome has confirmed whatever miracles have been wrought in all Christendom, and augmented them with his bulls and indulgences. Doubtless all had a contrary end to the miracles of Christ, and yet have, for they do not profit men, but empty their purses, put men to sundry charges, and lead them away from the faith of Christ, to the faith of Antichrist, confirming his religion, superstition, and doctrine. Neither does he place these his miracles among the last of his arguments, when the simplicity of the gospel is challenged. But if we be wise, we will beware of them, as of a most present pestilence.

### Paragraph 880

And among his miracles, the Lord, through John, recounts that above all things, he causes fire to come down to the earth, and that in his presence. He seems to allude to the story of Elijah, as we read in the fourth book of Kings, chapter one, and as we mentioned in chapter eleven. It was no small miracle, as recorded in the Acts of the Apostles, that the laying on of the apostles’ hands gave the Holy Spirit. Simon Magus also desired the same grace, but he was severely rebuked by Peter the Apostle, as we read in the eighth chapter of the Acts of the Apostles. And it is to be observed (as Augustine has also admonished in the fifteenth book of *De Trinitate*, chapter twenty-six) that the apostles did not give the Holy Spirit, for it is God alone who gives the Holy Spirit. Therefore, at the prayers of the apostles and at the imposition of hands, the Holy Spirit was given from heaven. Wherefore John the Baptist said that he baptized with water, but that Christ should baptize with fire and with the Holy Spirit. And by fire is signified the Holy Spirit.

But Antichrist, the Pope, will boast that he, having power given him from heaven, grants the grace of the Holy Spirit. Doubtless, during consecrations, he says that he gives the Holy Spirit. Likewise, in auricular confessions and absolutions, they brag that they give full absolution of sins, which indeed is a great miracle. They lay their hands upon the confessing sinner and say that they absolve him from the punishment and guilt, by the power received from the most holy See of Rome. Primasius, explaining this passage, says: “It is no marvel if that beast, which falsely assumes the name of the Lamb slain and yet living, also fraudulently claims this gift of the Holy Spirit, colorably by imitation, and feigns a donation to his ministers, as we remember that Simon Magus coveted but could not obtain.”

### Paragraph 881

There is another fire also, which Antichrist calls down from heaven, and casts and throws at his enemies, to avenge himself on them: namely, the dart and thunderbolt of cursing. This was terrible to kings, princes, and people. And these have so much feared the thunderbolt of excommunication that they have done and granted many things, which otherwise no one could have obtained from them. The story is known of Emperor Henry IV. For Platina in the life of Gregory VII shows that this emperor was excommunicated by the Pope. After he adds these things about the emperor, he came speedily to Canossa, where the bishop was with Mathilda, and immediately laying aside his royal robes, went barefoot to the gates of the city, and humbly requested to be let in. His entry was denied, which he took in good part, even though the winter was sharp and everything was frozen hard. Remaining three days in the suburbs of the town, and continually asking for pardon, at length, at the request of Mathilda and the Earl of Savoy and the Abbot of Cluny, he was absolved. Frederick Barbarossa, that he might be reconciled to the Pope, laid his neck under his feet to be trodden on, which clearly demonstrates the injuries done to other kings and people by this beast. I will yet tell of one. The Venetians besieged Ferrara, which pays tribute to the Church of Rome, for which cause they were excommunicated by Clement V. Therefore, France Dandalo, who was afterward created Duke, went into France, where the bishop was at that time, to seek pardon for that offense. It was long before he was admitted to come into the Pope’s sight. At last, he was led with an iron chain about his neck to the bishop’s table, like a dog, and there forced to lie under the table among the dogs so long, until the wrath of Clement, without any clemency, had passed, and he removed that ignominy from his country. And he was therefore always afterward called “dog” by his own countrymen, because he had lain like a dog at the Pope’s table to obtain absolution. This is written in Sabellicus in the end of the ninth book, chapter 7. The Pope, when excommunicating, uses tapers of wax burning, which he throws down to the ground from on high, so that even thereby we might perceive that it is he who calls down the fearful fire from heaven upon men on earth. And these things the beast does in the presence of men, namely, with great confidence and boldness, finally to make me afraid and to keep them in awe. For after the same kind of speaking, Paul also commands to rebuke a great man offending before all men, that others may be afraid. 1 Timothy 5.

### Paragraph 882

Nevertheless, the Lord adds the use and effect of Antichrist’s wonders, so that the church may rightly judge them. For he deceives the inhabitants of the earth by reason of miracles, and so forth. He will seduce by these signs or miracles, namely, by that grace of the spirit which he feigns to give, and with those his excommunications by which he would seem to cast men down into hell, those who are more inclined to earthly than heavenly things. And he will lead them from the faith of Christ to his deceptions. Therefore, we must judge the tokens and doctrine of Antichrist, because they seduce people. Cease therefore to marvel at how it has happened that the Pope has allured so many men of wisdom and learning to his side. You have already heard by what means this is done. Therefore, do not always be foolish; learn, take heed, and believe Christ and his Gospel, and cleave to it.

### Paragraph 883

Again he says,[note 60.7] that power is given to the beast (truly by God’s just judgment, that according to the apostles’ saying, those who prefer to believe lies rather than the truth may be judged) that he should work those miracles in the sight of the beast. What is it to work miracles in the sight of the beast, but to do them in the presence of all men, boldly and without fear, even to frighten and deceive the beast itself? Here therefore now we hear of two beasts: the beast that does the wonders, and that later beast in whose sight the other, former beast does those signs.[note 60.8] Yea, it shall follow hereafter that both the beast and the false prophet, who works these miracles before the beast, by which he also deceives the beast, shall be cast both into everlasting fire. Who therefore is the former and the two-horned beast, but the Pope? The very same is the false prophet also. And who is the beast, in whose sight the Pope works wonders, but the Image of the beast, and therefore a beast also, inasmuch as empire is raised from the beast and governed by the spirit of the beast.

### Paragraph 884

For it follows that the beast sets up an Image of the beast, [note 60.9] and that of the same beast which had the wound of a sword and lives: that is to wit, of that old Roman Empire. Now therefore a new Roman Empire is erected, which nevertheless is not plainly called a beast, but the Image of the beast: that is to say, an empire in deed, but which does not come so near the old one by as far as an image differs from the true example. For the old Roman Empire is as it were an example, of which the empire set up by the Roman Antichrist is only an image, representation, shadow, and as it were a dream, having nevertheless some similitude of the same.

### Paragraph 885

I have explained before, and demonstrated through historical accounts, how the old Roman Empire was torn apart and dissolved. In past times, one emperor ruled in the East at Constantinople, and another in the West at Rome or Ravenna. But from the time of Augustulus, for a period of three hundred years or more, there was no emperor of the West. And the lands that had belonged to the emperors were now held by others, and the empire was utterly lost. Therefore, around the year of our Lord eight hundred, when Charles the Great, King of France, came to Rome on Christmas Day, Leo III, Bishop of Rome, placed the crown upon Charles’s head, and the people proclaimed with a loud voice, “To Charles, the Emperor crowned by God, long life and victory!” These events are recorded in all histories, especially in the fourth book of Aventinus’s Chronicles of the Bourbonois.

### Paragraph 886

[note 60.11]Again, when this empire seemed to waver and to decline, as if it were about to fall, the Bishop of Rome instituted seven Princes Electors. Some refer this ordinance to Gregory V, who was Pope when Otto was Emperor. And some to Gregory X, who called Ralph of Abpurga to the empire. More will be said about this anon. But the Lord by John says expressly how the beast said to the inhabitants of the earth that they should make an image of the beast. For the Popes have, by speaking, and not by fighting (as appears in the histories of Bishops of Rome, especially of Leo III), erected a new empire. For by preaching, persuading, and practicing, they brought the empire to King Charles. Certainly Platina, in the life of Leo III, says that the Bishop, intending to gratify King Charles, who had deserved well of the church, after solemn service done in the church of Saint Peter, by a loud voice declares Charles Emperor, and crowns him, with the voices and prayers of the people of Rome. &c.

### Paragraph 887

But now we must examine more diligently wherefore the newly established empire is called by the Pope, the Image of the old beast. And here, indeed, many things could be alleged, but I shall recount only a few. Above all things, it is called the Image, both because it is named the empire itself, and would be taken for the old empire, where in truth it is a name without the thing, and a vain title, lacking that ancient power, majesty, and glory. For unless the emperor possesses the kingdom by his own inheritance, what kingdom shall he have by the name of emperor? Shall he have Rome? Shall he have Italy, the old seat of the empire? Shall he have France, Spain, Hungary, Germany? For although Germany is now considered the seat of the empire, she has her own princes, her own free cities, and which enjoy their privileges, even though they are called imperial. Theodoric of Niem, a German, and a familiar friend of certain popes, who also wrote the lives of certain bishops of Rome, who were last before the council of Constance, in the third book, chapter 43, of his Histories, writes: Of what magnificence the Roman empire is, at least visibly seen in Germany. For you shall have there an archbishop or a bishop, who has of yearly revenues twice as much as the King of Romans receives in all his dominions. And again, a temporal prince, that has more lands than has the emperor. And so forth.

### Paragraph 888

Moreover, in the old Empire, there was a mighty monarch who wielded full authority and was honored by all men as a god on earth. Such as Caius, Domitian, Diocletian, and others. His image represents the Pope, bishop, and king, and is like a certain earthly god, the greatest monarch, with fullness of power. Furthermore, Rome, or the old beast, had a most honorable Senate. So too does the Bishop of Rome have a princely Senate of proud cardinals. For they are, in effect, all princes. The book of the Roman governments recounts the vicar, or lieutenant of the Diocese of Asia (a diocese, in Greek, signifies a disposition, administration, dispensation, government, or jurisdiction), the vicar of the Diocese of Thrace, and of Pontus. Likewise, there was a noble man, president of the governments in Italy, who had many dioceses under him. And no fewer had the lieutenant of France. Just as the Count of Strasbourg, the captain general of the soldiers at Speyer, and the general of the soldiers at Worms acknowledged the Duke of Mainz as a proconsul, so at this day, the bishops of those cities are subjects to the Archbishop of Mainz. The bishops, therefore, seem by the Pope’s ordinance to succeed in the place of the Roman governments. Certainly, you shall see most of these bishops called not only most reverend fathers in Christ, but also most noble and mighty Dukes and Princes of the Empire. And this is also manifest: that the Emperor of the old beast had his legions, the Roman eagles or standards, and most expert and powerful captains of war. But the high bishop and king of Rome has, in that imagery, obedient children, kings and princes in Europe, not to be despised, whom he may command if need require, to stretch forth the secular power. For so thunders Boniface VIII in the first book *De Major*. And obedient, he says, is he who denies the temporal sword to be in the power of Peter, he misunderstands the word of the Lord, saying: “Put up thy sword into thy sheath.” Therefore, both swords are in the power of the church, namely, both the spiritual and the material sword; but this must indeed be exercised for the church, the other by the church. The spiritual by the priest, the material by the hand of kings and soldiers, but at the will and patience of the high priest.

The old beast had his laws written and published daily in a manner new. The Popes, therefore, after the imitation of the imperial laws, have written decretals and often make new laws. Indeed, they say that the voice and precepts or commands of the Pope are to be received and taken as the words of our Lord Jesus Christ, the Son of God, and Apostle Paul. They add moreover that we must stand by the Pope’s determination; that where the Pope is, there is the general council; where the Pope is, there is our common country. He is compelled or reproved by no man, even though he be called an heretic. He has all laws in his breast, or in the scroll of his breast; he may interpret or expound all things. The same ratifies no sentence; and it is in him alone to take away one man’s right and give it to another. He may take away privileges and, at his will and pleasure, not only change bishops, but also depose the emperor himself and declare no sentence of the emperor. All the world is the Pope’s diocese; and the Pope is the ordinary of all, having fullness of power as well in spiritual matters as temporal. For he is Lord of Lords, and has the right of the King of Kings over all subjects. For he has no peer; he is all things, and above all, and it is necessary for salvation to be under the Bishop of Rome. For there is one consistory or judgment seat of God and of the Pope. These things I have taken from their own books, namely, from their Decretals and glosses. There is a book by Antony Russell of Arezzo, on the power of the Pope and the emperor, where you may read innumerable things of the same sort. But of these things which I have noted hitherto, I suppose it is made plain enough, how the Pope, who is here also called the false prophet, has set up the image of the beast.

### Paragraph 889

Hereunto John adds another thing [note 60.16] that the empire, once established, and all things set in order, the beast or false prophet moves all that weight, and puts life into the image, so that it can speak: namely, that the false prophet has given it the ability to speak. For unless the pope confirms the election of the King of Romans, he shall not be considered worthy of the name of Emperor. 22. quest. 5. de forma, in the gloss, the emperor swears to the pope, as a client to his lord. You may read the same in the first book, title 9, de jure iurando, in Clementines. Moreover, who does not see how both the emperor and other princes are surrounded by a company of bishops, who inspire them as to what they should speak or do, and how they should behave in all things? For this reason, legates are also sent, called *legati a latere*. And it is not unknown that in all princes’ councils, the spiritual leaders generally have the chief authority. They are, for the most part, chancellors, secretaries, ambassadors, and so on? And the pope and king openly say how he ought to judge all men, but to be judged by no one. Indeed, his officials also usurp this authority for themselves. If there is any assembly, the Bishop of Rome commonly rules by his spirit and governs the chiefest matters, especially matters of religion. For unless the decrees please the fathers, they threaten to abolish whatever the states have decreed. But if a general or national council is called, it is wholly ruled by the pope’s spirit. It speaks and determines as it pleases the pope. For unless it decrees according to the pope’s pleasure, he will seek to abolish everything altogether. For we recently heard that the synod or council is where the pope is. Innocentius the ninth, in the third question [note 60.17], says, “The judge shall be judged neither by the emperor, nor by the whole clergy, nor by kings, nor by the people.” And the gloss on that passage notes that the council cannot judge the pope, and so on. Therefore, if the entire world were to give a sentence in any matter against the pope, it appears that we ought to stand by the pope’s sentence against them all. Indeed, the same glossier in another place says: “The pope, if he will, may dispense against the council, for he is more than the council.” Most truly, therefore, did the Lord say at this present how the beast had power to give a spirit to the beast, and that the image of the beast should speak. For whoever do not show themselves obedient and willing instruments unto this beast in all his affairs are accounted for dead and rotten members, and therefore to be cut off from this vital body. Indeed.

### Paragraph 890

However, in the meantime, lest I should blame any worthy person, or seem to harshly criticize those who have done no wrong, or should be said to fail to acknowledge the goodness of God working in empires, but rather to find fault with it, and to confound and mix together all things both good and evil without discernment, certain things are here by a lengthy, yet necessary digression, to be admonished and better declared. I admonish therefore and repeat, that the Lord our God is the author of empires, and ordains them for the benefit of people; but that the Devil joins himself with God’s good institutions, and according to his evil nature corrupts those good institutions of God, by influencing people’s affections diversely and applying them to evil matters. Whereupon in governments very many things arise which are to be disliked by the godly: such as tyranny, alteration of the state, and similar things. Nevertheless, although God hates all wickedness and cannot allow any evil, we see that he, by his infinite goodness, uses the evil governments of men to the good or profit of his people. For God loves his church exceedingly, and seeks to relieve and comfort all mankind by empires, although not altogether, or in all things commendable.

### Paragraph 891

I will therefore not deny that sins within the Western Empire were renewed; that is to say, sins the Image of the beast was set up. These seven hundred years, they have many times governed so, that it has easily appeared that God has wrought the health of his people in the governments. Daniel figured by beasts the four Monarchies of the world, which nevertheless did not suppose that all their Princes were beasts, nor did he condemn all Princes, nor did he think that there had been or should be no good thing in them, although the most part were most corrupted. There were found in the old Roman beast (to speak nothing in the meantime of the Princes of Assyria, Babylon, Medes, Persia, or Macedonia) who have set forth profitable laws, set in the books of Justinian. There have been found under that most cruel old beast, who have advanced the true religion of Christ, and defended earnestly the church of God, such as before we said were Constance, Constantine, Theodosius, and diverse others; which come all under the number of the Empire, but not of the beast. But inasmuch as the beast signifies the Empire, so may there be found Princes under the Image of the beast not a few, who have both set forth wholesome laws, and have employed great benefits upon mankind: as have done Charles, Lewis, and Lothaire of Saxony and others. Notwithstanding that they themselves in many things cannot be allowed of the godly. There are found among the later kings of the new Empire, which in power and majesty were not much unlike the old, in virtues not much behind them, but in certain things [illegible]. There are found who have assayed to purge the empire from Popish corruptions, and to bring the Popes under [?Corum]. But with no great or good success. For what the Otthones, Henries, Lodouicks, Fridericks, briefly many French Princes, Saxons, Swedes, Bavarians, and of Austria have been, many notable testimonies of histories do report; which testify that certain Kings both of France and of other realms also, have not bowed their knees to this Baal; or if they have done at any time, yet have they repented, and have showed some token at the least wise, whereby the wise might perceive that they set not much by that beast.

### Paragraph 892

Therefore, all holy and excellent men who have lived throughout time, during the period when the Image of the Beast has reigned, are to be excused. I mean emperors, kings, princes, bishops, states, cities, and the people of the empire and other realms who lived, but were not under the unhappy image of the beast; for they did not offer themselves to the spirit of the beast to be moved and governed by it, nor did they speak explicitly what the beast commanded. Rather, they spoke against the beast and greatly disapproved of his actions. Therefore, just as I have not included the godly princes and good government within the old monarchies, specifically the old Roman Beast, nor condemned them of beastiality [?], so now, when criticizing the Image of the Beast, I do not confuse the good and godly princes and people, and their government which is not evil, with the corrupt actions of Antichrist. For I always except moderate and profitable empires, honest and godly men, however they live under the Image of the Beast, yet not according to the inspiration of the beast or false prophet.

### Paragraph 893

Hereunto I add this also, that the empire was not suddenly established after the will and pleasure of the Bishop, but by diverse spaces of time, sundry attempts, and treasons innumerable: therefore at the length it devolved to an extremity of corruption, and as I may say, bestiality. Whereby it appears that the prophecy of Saint John is to be applied to the things themselves, and to the times, after the state, maliciousness, and corruption of every thing and time.

It is most certain, and by common consent of all historians plainly testified, that in Charles the Great, through the means of Pope Leo the Third, the empire in the west decayed was renewed: and that thus the image of the beast, that is to wit, of the Roman Empire, was erected. And albeit that at this time the empire decayed in the west was restored by the Pope, yet it is evident that the Popes in the beginning of this empire, by certain donations and gifts, did not use so great power as they usurped to themselves afterward, when they had overthrown and deposed certain emperors. For although the donation seemed to be made by King Pippin, and the Pope was then to have received the beginning of his kingdom, yet that he was subject to emperors and kings with the City of Rome also, this same image and other things prove. In the French Chronicles of the Acts of King Charles in the year of our Lord eight hundred and one, it is found written: afterward having set in order the matters of the city and Bishop of Rome, and of all Italy (therefore Italy then also obeyed the Emperor), not only public, but also (mark) ecclesiastical and private (for all the winter the emperor did nothing else), departing from Rome with his son Philip, he came to Spolet. The same author in the Acts of the year eight hundred and sixteen, says he, Stephen, elected in the place of Leo the Third, takes as great journeys as he could to come to the Emperor, sending in the mean time two ambassadors which might treat with the emperor (Louis the Pious) for his consecration. So likewise in the Acts of the year eight hundred and seventeen, it is shown how Paschal being chosen sent an embassy to Emperor Louis. In the Acts of the year 823, the same Bishop stood at the examination and judgment of the emperor. You may find in the Acts of the next year that Emperor Lothair established the matters of Italy and Rome. Yet doth the same author again make mention of the donation of King Pippin, which gave to Saint Peter Ravenna, and Pentapolis, and all the government. Yet doth he make no mention of the donation either of Charlemagne or of Louis the Pious. The 43rd distinct makes mention thereof. I Louis, &c., in the gloss is written thus: There Louis gives Rome and diverse other things to Saint Peter and to Paschal the Pope. All historians in manner make mention of the donation of the Kings of France. An abridgment of all gathers out of the library Dolaterane in the third book of Geography, in the Acts of Pippin and Charles. Whereby one may easily conjecture what manner of canon is set forth in the 96th distinct, in these words: Constantine the Emperor hath given and granted to the Apostolic See the Crown and all the Imperial dignity in the City of Rome and in Italy, and in the west parts. Which by and by after he discourses with a lengthy exposition out of the life of Saint Sylvester, written (as they say) by Gelasius, in the chapter following. But Antony, Bishop of Florence, denies in his History that this donation remains in any old books. Cusanus and Laurence Valla have impugned the same: neither hath Otto, Bishop of Friesing, in the 3rd chapter of the 4th book of his story, nor Marsilius of Padua in the defense of peace, nor Raphael Volaterane allowed the same, nor many more that I could rehearse. Moreover, in the Chronicles of the Kings of France, set before the story of Paulus Aemilius of the Acts of Kings of France in the year 755, thus you may read: Pippin again entered into Italy, and Aiulf subdued; he gave gifts to Maximus, Bishop of Rome, also the Dukedom of Ravenna of very great lands, lest any man should unthankfully and unjustly take away this largesse from the French Kings, ascribing to the emperor Constantine which Pippin gave to the church of Rome, against the will of the Greek Emperor, affirming the same possessions to be the right of the Empire. From thence Pippin first received and brought into France the ecclesiastical rites of the Romans and ceremonies of songs.

### Paragraph 894

However, the rule of Charles’s empire and his descendants was not very stable or permanent. From the first year that Charles was made Emperor, to the seventh year of Conrad, who was Lewis III’s nephew, by his brother, the last of the house of Charles, there are reckoned about one hundred and nineteen years. For Charlemagne reigned for fourteen years, Lewis for twenty-six, Lotharius for fifteen, Lewis the second for twenty-one, Charles for two years, Charles, surnamed the second, for three, Crassus for twelve, Arnulf for twelve, Lewis the third for ten, and Conrad for seven. Conrad, lying on his deathbed, nominated King Henry, Duke of Saxony, surnamed Falconer. And thus the empire was transferred to the Germans. This Henry, called the first, never went to Italy, nor was he consecrated or crowned by the Pope. His son, Otto, the first of that name, was summoned to Italy and is said to have gone there with a great army, being received at Rome and hailed by the people as Emperor and Augustus. Otto, writing in the sixth book of Histories, chapter seventeen, affirms from the decrees that Pope Leo VIII of that name did consecrate this Otto, the first King of the Germans. For his father, Henry, refused it. Albert Krantz, in the tenth and eleventh chapters of the fourth book of Saxon matters, affirms that Pope Leo made a surrender of all things that the Popes had received from the kings of France and the author defends this surrender as true. However, the keeper of the library testifies that Otto confirmed the donation of the kings of France, Pippin, Charles, and his sons. Moreover, there remains in the decrees a copy of another document, the forty-third distinct, whereby King Otto binds himself to the Pope, promising that he shall not interfere with anything concerning the Pope and the Romans; secondly, that he shall restore all the lands of St. Peter that come into his hands. Let the reader judge what these lands are.

### Paragraph 895

Shortly after this time, around the year of our Lord 996 [note 60.23], it is said that by the decree of Pope Gregory V and with the consent of Otto III, the seven prince-electors were assigned, to whom the defense of the church (as says Wimpelingius) and the Roman Empire was committed. All historians and writers agree on this, including the Italians Blondus, Platina, Sabellicus, Volateranus, Egnatius and others; Germans, Albertus, Nauclerus, Carrion, Functius, and several others. However, some have made no mention of this ordinance. Therefore, Aventinus in the fifth book of Chronicles, folio 510.707, says that he knows [?I cannot tell how certainly] that after the death of Frederick II, the electors were instituted and confirmed by Gregory X. But however that matter stands, it is certain that there have been many among the seven prince-electors who were fervent and earnest in true religion, and excellent in all kinds of virtues, especially the seculars, as they term them, who have greatly disliked the tyranny and impiety of the Popes of Rome, so much so that they have stoutly and often withstood them. Our age is undoubtedly indebted to this order or state, because a good part of the preaching of the holy gospel has been reformed, which both they and other princes of Germany, most worthy of praise, valiantly defend and maintain [?by God’s inspiration]. The Lord increase in them, and in other godly princes throughout the world, his gifts, and mercifully keep and preserve them. But to return to the progress and order of the history, it is certain that immediately after Gregory V, the Devil invaded the see of Rome. Neither could Platina dissemble this thing, a writer of Popes’ lives known to all men [note 60.24], who has very favorably spared his lords and masters, and many times has covered their abominable acts; yet, writing of the successor of Gregory V, Sylvester II, before called Gilbert, a monk of Florey, he says that he deserted his monastery and followed the Devil, giving himself wholly to him. And immediately he adds: Gilbert, moved by ambition and a devilish desire to rule, through bribery first obtained the archbishopric of Reims, then of Ravenna, and afterward, with greater entreaties, the Devil furthering him, he obtained to be pope; yet under this condition that after his death he should be the Devil’s wholly. &c. He who would know the full story, and abridgment taken from Antoninus, Nauclerus and others, let him read the ninth book of Funccius’ Chronicles, under the year 998. Benno, a Cardinal, supposes at this time to be fulfilled those thousand years, after which the Devil, breaking loose, began again to rage in the world, of which certain things will follow in the twentieth chapter of this book. Nevertheless, I shall not refuse to gather here certain things from this Benno Cardinal, and briefly to recite them here for the declaration of our matter.

### Paragraph 896

Therefore, in the life and acts of Hildebrand, called Gregory the Seventh, one Gerbertus, who had infected the city with sorcery, (he says) after the thousand years fulfilled, coming up out of the pit of God’s permission, was Pope for four years, and changing his name, was called Sylvester the Second. And after Gilbert, for twenty-five years (I suppose it should read thirty-two years), and how they reigned, histories testify, and that very evil Theophilactus’s scholar violently achieved the seat, called Benedict the Ninth. He had a dear friend and privy to all his doings, one Gratiane, Archpriest of St. John, near the Latin gate. To whom Hildebrand, a monk of Cluny, forsaking his abbey, did familiarly cleave and became a familiar friend of his. But Benedict, fearing himself, sold his seat to Gratiane, the master of Hildebrand, receiving from him five hundred thousand pounds, which promoted to the office was called Gregory the Sixth. Nevertheless, they shortly had a third Pope, Sabinus, and he was called Sylvester the Third.

The emperor, therefore, Henry the Second, a godly man, valiant, wise, and stout, going to Rome to purge the church (for as yet the bishops used not full authority), compelled Benedict or Theophilacte the Magician to flee, cast Gregory in prison, and sent Sylvester away to his old bishopric. And he holding a council, placed the bishop of Bamberg, whom he called Clement, in the seat, of whom also he received the crown. And he brought Gregory with his disciple Hildebrand with him into Germany. In the meantime, Benedict returning to Rome from flight, vexes Clement, and with much enchantment infects the city; and by letters received from Hildebrand out of Germany, he learns what is done in the emperor’s court. Gregory dies there in prison, and left Hildebrand his heir both of his false packing and of his money. Clement dies also, whom Damasus the Second succeeds immediately, but straightways poisoned; by reason of the tumult that was in the city, the emperor sends Bruno (bishop of Toul, come of the noble house of the earls of Holstein), a worthy man. Here Beno annexes: in whose train, through the overmuch sufferance of the emperor, Hildebrand was permitted to return; by this permission to subvert both the bishopric and empire under pretence of religion. And this Beno herein was a true prophet, which says also in the story of Hildebrand: and telling Bruno many things, by the way crept into his favor; and as soon as he came to Rome, obtained of him that he was made one of the keepers of St. Peter’s altar. And in a short time he filled his coffers. And he also reconciled his old lord and master Benedict, feigning repentance deceitfully to Leo the Ninth (for so Bruno being made Pope was called), and through the counsel of Benedict, otherwise called Theophilacte, he armed Leo against the Normans and betrayed him unto them. The Germans therefore slain by treason, scarcely the Pope escaped all desolate. This says Beno. And certainly it is that this monk Hildebrand, from that time forward aspired to get the seat; and in the meantime, whilst it was governed by others, he incensed and ruled the Popes, as Leo the Ninth, Victor the Second, Stephen the Ninth, Benedict the Tenth, Nicholas the Second, and Alexander the Second. But they smell of Hildebrand’s style, that are set forth in the name of Leo, Nicholas, and Alexander. But at length, he himself climbed into the chair, in which he so used himself that no man unless he were stark blind but might see that his devilish government had requited most abundantly Henry the Fourth, the son of Henry the Third, his father’s carrying of him into Germany. And he began openly and impudently to take upon him the power of the emperor. Neither can it be told in few words, in what detestable wise this beast did afflict both the emperor and empire, all the while he was Pope, for the space of twelve whole years. An abridgement of that story has John Functius compiled in the tenth book, under the year of our Lord 1074.

### Paragraph 897

I know that Platina, and many Italian writers, yes, and some Germans also, do highly commend the religion and virtues of this Gregory the 7th. By this, the Papal tyranny, under the guise of religion, is wonderfully augmented and confirmed, and many are blinded. Yet, through God’s grace, it has come to pass that men of grave authority, religion, and virtues have fairly and well stripped off the disguise from this beast. Therefore, Synods and Councils are not to be condemned; first, the Council of Mainz, where 19 famous Bishops were present. Then, a Synod was assembled at Brixia, with 30 Bishops and most of the nobles of Germany and Italy. There was also a Council assembled at Worms, where King Henry was present, and all the German Bishops (except those of Saxony) deposed the Pope from his office. The Epistles and fragments of these Councils are found in the Chronicles of Verspergen, chiefly. He is openly accused by these of all wickedness and ungodliness, of hypocrisy and cruelty. We have rehearsed a little before what Cardinal Benno, a writer of his time, has committed to writing. There remain also testimonies of Sigisbert, an old writer, concerning this Pope. Anyone who wishes may read the 5th book of Aventinus, from page 162 and onward, and also the preface of the 6th book. The same author, in the 7th book, reporting the words of Eberhard, Bishop of Salisbury, who was present at the Council of Regensberg, says that Hildebrand, 170 years under the guise of religion, laid first the foundation of the Antichrist’s kingdom. This wicked war he himself first began, which his successors have continued to this day. First, they excluded the Emperor from the Pope’s election and referred it to the people and priests of Rome. Afterward, they mocked and thrust out [?]; they go about now also to bring us into subjection and bondage, so that they might reign alone. And the things that follow declare that there have not lived many Popes more bold and impudent than this, who have advanced more highly the majesty of the seat. He excommunicated Emperor Henry the 4th and deprived him of the imperial dignity; moreover, he stirred up his subjects against him and absolved the rebels and traitors from their oaths of fidelity. And he himself, like a monarch, gave the Crown of the Empire to others at his pleasure. The power and treasure of the Empire have been so worn and wasted, both by civil and foreign wars, that for many years now the kings of Germany have neither been able to recover their strength nor yet to resist the most arrogant tyranny of the Popes. And thus, at last, the Pope has become a monarch, and Emperors, Kings, and Princes are made their clients and wards.

### Paragraph 898

When Gregory VII was dead, Mokes of Hildebrand’s sect and faction succeeded him—of his manners and corrupted nature, as if heirs and sons who remain entirely within his character—Victor II, Urban II, Paschal II, and Gelasius II [note 60.29]. Paschal II caused the son, Henry V (oh, wicked and detestable parricide), to war against the father, the miserable Henry IV. Shortly after, Gelasius II and Callistus II excommunicated Henry V as well, and did not cease to harass this prince until they had wrested from him the gift or election of the Bishopric [note 60.30]. And that to the great and inestimable profit of the See of Rome, and to the unrecoverable loss of Germany, and so forth. These things are described more fully by Urspengen [?] in the Chronicle of the year 1122.

### Paragraph 899

In the times following, the audacity and power of the Popes has increased hourly, and the German kings have resisted them stoutly enough, but yet with small success. Where in the mean season we must remember the words of the Lord, uttered by Daniel, saying: and there shall arise a King of a shameless face, and understanding propositions, and his strength shall be fortified, but not in his own force: and it cannot be believed how he will destroy all things, and he shall prosper, and do &c.

### Paragraph 900

I will therefore touch briefly what things in the times following the Popes have attempted against kings, and boldly done for the establishing of their empire and monarchy. Pope Alexander III excommunicated Frederick I, called Barbarossa, and trampled him under his feet. And when the good prince said how he showed this obedience to St. Peter, the beast exclaimed, setting himself also before Peter, and said, both to me and to Peter, and stamped on him. Pope Innocentius III could not abide, much less allow Philip, the son of Frederick, created Emperor; but commanded the electors to choose another, I mean Otto, Duke of Saxony, whom notwithstanding shortly after he excommunicated also. That proud beast said that he would take from Philip the imperial crown, or lose his Apostolic mitre. Unto Innocent are ascribed those most proud words, which are read in the decretal of Gregory IX, de Elect., in title 34, chapter 6. on this wise: that the princes have right and authority to choose a king, and afterward to advance him to be Emperor, we acknowledge, as we ought, as to whom of right and ancient custom it is known to appertain; especially since that such right and authority came unto them from the Apostolic Seat, which translated the Roman Empire from the Greeks to the Germans in the person of great Charles. However, the princes must know again that the right and authority to examine the person chosen king, and to be promoted to the empire, belongs unto us, which do anoint, consecrate and crown him, &c. The same is written in the first book, title 33, de maior. and obedient., writing to the Emperor Constantine. So much diversity, says he, as there is between the sun and the moon, so great a difference is there between Popes and Kings, in God’s name.

### Paragraph 901

But Emperor Frederick II, nephew to Barbarossa, an excellent prince, was excommunicated by many Popes: Honorius III, Gregory IX, and Innocentius IV. Indeed, Gregory IX, while Frederick, that excellent prince, waged war in Syria with the Sultan, invaded and occupied Frederick’s provinces. [note 60.35] There were most cruel wars and discords between the Popes and this Frederick. The same Innocentius IV excommunicated Conrad IV of that name, son of Frederick II, and stirred up the Prince of Thuringe against him. And when Emperor Conrad was dead, the Pope obtained the goodwill of the Neapolitans to yield themselves to the See of Rome. Conrad had left a son and heir, Conradine, and Manfred his bastard brother, who would be called king of Sicily. Therefore, Pope Urban IV (some say Clement IV) against Manfred summoned Charles, brother to King Louis the French king, Earl of Provence and of Saintonge, to come with an army into Italy,[note 60.36] and called him King of both Sicilies. He overcame and slew Manfred at Benevento and received the kingdoms of Sicily from the Pope to do him homage. But Conradine, Duke of Swabia, accompanied by Frederick, Duke of Austria, led out of Germany a well-equipped army into Italy against Charles for the recovery of his old and paternal kingdom. But vanquished by Charles at Lake Fucino, he was taken with Duke Frederick. [note 60.37] It is said that 12,000 were slain. The occasion of so great an evil were the Popes, chiefly Clement IV, who, when questioned by Charles, the worthy prince, about what he should do with his prisoners, answered in a way that the Frenchmen understood meant they must suffer. Therefore, he put them both to the sword. In this, the house and posterity of the most noble Dukes of Austria and Swabia is said to have failed. Paulus Aemilius discusses this matter more at length in the 7th book of French Acts, and Aventinus in the 7th book. But even so, the ire and fury of those most holy fathers could not be pacified, so that the most noble Dukes of Swabia had, for God’s glory and the common good, most godly and most constantly resisted the Roman bishops, wolves, I would have said.

### Paragraph 902

But these parricides and bloody wars displeased all good people everywhere, and especially wise and godly princes, so that they understood how they must avoid that empire and flee from it as from the plague. For it was not only a shadow, but would also utterly consume the yearly revenues and treasure of whoever received the office. Now it was known throughout the world what the most valiant and excellent princes of Germany had endured for 119 years, from Henry IV to the sons of Frederick II, because of the bold ambition and incredible malice of the Popes; and that many of them had lost both their lives and their ancient kingdoms, and most excellent liberty.

### Paragraph 903

And there was the empire without any emperor for several years, which I am accustomed to call a desolation of the kingdom or empire. For the popes with their invincible and intolerable pride and tyranny had so weakened the force of the emperors, that the empire seemed subverted and destroyed; neither could there easily be found anyone who would take it up or thought it worthy to be desired. At last, at the commandment of Gregory X, who held a council at Lyons, Count Rudolph of Habsburg was chosen. Although he did not refuse the offered office, yet when often requested to come to Rome, he is said to have answered: “The wayward steps of feet do fear me sore,” meaning by this saying that he did not trust the popes, who by their schemes had destroyed both many princes of Germany and also innumerable people coming to Rome. And this Rudolph is said to have been crowned king in the year of our Lord 1273, the 200th year after Gregory VII. And so long a time lasted the conflict of popes and emperors. A little while after, while Albert, the son of Rudolph, was chosen emperor, and the election was referred to Boniface VIII of that name, he stoutly rejected it, and showed immediately in word and deed that he was both pope and emperor, who rightfully held both powers. Which I explained in sermon 58, and the same Albert Krantz declares exceedingly well in book 8, chapter 36, of Saxon matters. In the place of King Albert, was substituted Henry, prince of Luxembourg. But what authority over him and the empire challenged Clement V, he who wishes, may know from the Clementines. For there is a long treatise thereof in book 2, title 9. I could also rehearse many other similar things of Pope John XXII and others, if I did not think it superfluous.

### Paragraph 904

For these things which I have discussed up to this point, it appears sufficiently clear that the popes themselves, with a mischievous boldness, have seized imperial power, boast of being monarchs, and abuse the service and ministry of kings, treating them as wards and clients, yet pretending the name of sons, so that they may have them the more obedient. For so Gregory VII wrote to Géza, king of Hungary, as can be read in chapter 17 of this book, sermon 75. Yet understand in the meantime that the great part of princes and nobles have not known this beast, but have rather opposed him, and therefore have not come into the number of the beast. But inasmuch as they lived under the Empire, they were far estranged from the beast.

### Paragraph 905

By this I would have them addressed, who will exclaim and say: who can take it well to have the holy Empire called the Image of the beast, and so many noble Kings and Princes, Cities and worthy people? But I neither ought nor will change the manner of speaking which the scripture uses. These are the Lord’s words all, which Daniel in old time, and now Ibon, have revealed to us; but I may except and excuse those who are excused by the testimony of scripture. The way is ready and brief: whosoever will be free from the beast, let him take heed that he be not influenced by the Pope’s spirit; and that he speak not and do not that which the Pope commands against godliness. Let him rather be ruled by the spirit of Christ; and so it shall come to pass that, dwelling in the midst of Babylon, he shall not live after the iniquities of Babylon, but in the Kingdom of Christ.

### Paragraph 906

It follows: and the beast shall cause that whosoever shall not worship the Image of the beast, shall be slain. And it is all one offense to worship that old beast and to worship the Image of the new beast. Of the worshipping of him, I have spoken a little before. Therefore they worship the Image of the beast, which admit the decrees and those ordinances of the seat and Empire, speaking the inspiration of the beast: which allow the Roman religion, which fall to the kissing of the feet, and show themselves in all things obedient children of the seat, and are faithful to the popish Empire. Now if any will not be such a one, and would be content with Christianity, would abhor Rome the seat of the beast, and detest the Image of the beast, he like a church robber and traitor is judged unworthy of life. There is a canon in the fifth book of Decretals, the seventh title of heretics, wherein without any circumstance of words, Lucius the third of that name, determines plainly that heretics are struck with an everlasting curse, whosoever believe and teach otherwise of the Sacraments than the church of Rome believes and teaches. He commands moreover that such being deprived of all dignity, shall be committed to the judgment of the secular powers to be punished with due correction. But if the temporal magistrate will not punish and so defend the church, then he shall be also deprived of all honor, etc. But why do I tarry in rehearsing these things? All men at this day know and see what things are done daily. They are condemned, exiled, excommunicated, shut up in prisons, vexed with sundry torments, at the length also cruelly slain, whosoever shall refuse to worship both the beast and his Image. The Lord Jesus, the true King and Bishop of his church succor us, and restrain the cruelty of the ungracious beast. Amen.

## Sermon 61: ¶ Of the mark and number of the name of the beast.

### Paragraph 907

.

### Paragraph 908

And he made all people, both the humble and the powerful, the wealthy and the impoverished, the free and the enslaved, to receive a mark on their right hands or on their foreheads. And that no one might buy or sell, except he who has the mark, or the name of the beast, or the number of his name. Here is wisdom. Let him who has understanding calculate the number of the beast, for it is a man’s number, and his number is six hundred threescore and six.

### Paragraph 909

He adds the rest, by which Antichrist may be recognized and avoided. Indeed, he may chiefly be known by these things that follow.

### Paragraph 910

And he speaks of the subjects of Antichrist and of this new king and bishop. He will procure to himself, says he, an infinite multitude of all kinds of men, of all states and degrees. For his kingdom shall be ample and large. Therefore the Lord resists here certain kinds and states of men: and under the same he understands whatsoever is of the same state in the whole world. The Roman Antichrist brought under his subjection small and great, rich and poor: free, namely nobles, and bond. For we see that emperors, kings, dukes, marquesses, earls and baronesses, realms, countries, cities, patriarchs, archbishops and bishops, prelates, doctors, clerks, and laymen obey him: also men of greatest power, riches and wisdom, together with the poor people. There is no such kingdom, and so diversely composed in the world, no not among the Mahometists. And all these verily willingly are subject to the seat: yea they have persuaded themselves that they cannot well live, that they cannot be saved, unless they be subject to the See of Rome.

### Paragraph 911

And just as princes distinguish their subjects and servants by colors and signs, and common people also mark their livestock with different brands and marks, so that they may be known to belong to whom or serve whom, for every man has his colors: one white and black, another red and blue, another white and red, some black and yellow, which they give to their soldiers and servants to wear, and by this they declare themselves to be retained by him or him. And as they mark their horses with their brand and set their marks upon household vessels, so Antichrist will undoubtedly have his sign, namely, his mark, by which he may both bind people to him, and so bind them that they may be distinguished from others, and by this means wear the badge, as it were, the colors of their lord and master. And he will give his mark on the right hand, or on their foreheads.

### Paragraph 912

According to Arethas and Primasius, and indeed all commentators agree, it explains this as the confession of the mouth, and the study and work of a good deed. We have heard how Christ, in the seventh chapter, imprinted on the foreheads of his servants the faith that works through effective charity. Indeed, the sign of God’s children is faith, and the love that arises from that faith. So do the writings of the Evangelists and Apostles testify. Nevertheless, Christ also has the external mark of his servants, those wholesome sacraments of the church, Baptism and the Lord’s Supper.

### Paragraph 913

How be it if any are baptized at this day, and partake of the Lord’s Supper, call upon God the Father with the Lord’s Prayer, and utter their faith by a sincere confession of the Apostles’ Creed; moreover, confess those to be good works which are done in faith after the rule of the ten commandments, and besides this, shine in good works: shall he be taken for a true Catholic and a right Christian man? In old times, doubtless, all men would have embraced him as a brother. But what should he be at this day in the Pope’s kingdom? You shall seem, by all these things, to have confessed nothing at all of the true faith, except you plainly profess that you believe after the faith and tradition of the church of Rome; and that you do acknowledge those for good works which the church of Rome has approved. Unless you believe and profess on this wise, in vain shall you confess all the former matters. No, though you say moreover that you believe the law, the prophets, the gospel, and the Apostles. They will like you a great deal better if you say you are an obedient child of the Apostolic See and church of Rome, than if you should say that you are the child of God, a Christian man, that you put your whole trust in the Son of God, which is the only salvation and righteousness. Yea, those shall find it, who will by and by at these words cry out that they smell of heresy and a mind infected with poison. I say nothing, experience itself will witness, that I say truth. And thus does the Pope mark his men both in the forehead and right hand. Thus are the Roman whelps discerned from other faithful, as it were by marks.

### Paragraph 914

Besides this, there is another thing. All Papists plainly testify that unless a man be marked in the forehead with chrism by the bishops’ hand, he is no Christian, however he be baptized and believe in Christ Jesus. Whereof it follows that they attribute more to their confirmation and anointing by the bishop than to Christian faith. Read the book called *Summa angelica* in the title of confirmation. This therefore is a sealing of the Papist religion; the Christian marks of Christ are sufficient. The Pope also by another way imprints his mark on the right hand, by making of vows and performing of oaths, as they term it. For those who make a vow in entering into any religion, as they call it, do bind themselves to the Pope and See of Rome by a stipulation. Furthermore, Antichrist, the Pope, by oaths also to be performed by the holding up of the right hand, doth bind and bring into danger emperors, kings, archbishops, princes, bishops, doctors, universities, and all states. They promise that they will attempt nothing against the church of Rome, nor against the high bishop thereof, nor yet against the privileges and statutes of the See. There remain the manner of oaths in the decrees and decretals. I touch these things briefly. They see more who do not shut their eyes. And all men behold how the Pope has set his mark on the right hand and forehead of men.

### Paragraph 915

There follows again the cruelty and bloody tyranny which Antichrist practices against the Christians, that is to say, against those who will not receive the mark of the beast: that is, who will not prostitute themselves to the lust of the Pope and the seat of pestilence. Antichrist, says the Lord, by his power shall bring to pass that none may buy or sell, save he that hath the mark of the beast, and so forth. And these come all to one effect: the mark of the beast, the name of the beast, and the number of the name of the beast. For he hath the mark of the beast, which acknowledges the seat and professes the faith of Rome: and even to whom the Christian faith is not enough. He hath the name of the beast, whoever he be that will be named an obedient child of the holy See of Rome, and acknowledges the Pope to be head of the universal church. He hath the number of the name of the beast, which has a society with the beast, and that society that number betrays or shows. Therefore, except thou acknowledge the Pope to be supreme head of the church on Earth, with the fullness of power; unless thou profess to follow the faith of the holy church of Rome, and to detest all things whatever that See has condemned, thou art forbidden fire and water. That same has the Lord called to prohibit, that thou may neither buy nor sell. We say in Dutch, signifying one that is banished out of all men's company. He understands therefore excommunication, that horrible thunderbolt of the Pope, wherewith are stricken all those that have set more by Christ than by the Pope, or which have loathed the Pope’s decrees in comparison of the Gospel. Let him read, who desires, the Sixth Decree of Boniface the Eighth, in the Fifth Book, the Second Title, *de hereticis*. Also, Clement the Fifth, Book Three, Title *de hereticis*. But he that will know exactly a compendious treatise of tyranny, and a glass of butchery, let him read the Bull of Martin the Fifth, which is subject to the Sessions of the Council of Constance, and is written to Bishops and inquisitors of heretical depravity. Among other things, there is one which gives a wonderful light to this place which we now expound: where it is commanded that they do not permit those that despise the communion of the church of Rome, to keep or dwell in any house or lodging, to make any bargains, or occupy any traffic or trade of merchandise, or to have any comfort of humanity with the faithful of Christ. Read thou the rest, leaf 134. Hereunto may be added, that in Popish churches is the greatest buying and selling of all. But unless his crown be shaven, and his hands imbrewed with oil, that is, except he hath received in the forehead or head and in the right hand the mark indelible (for so they term it, that cannot be put out), he hath no merchandise left him in the house, nor so much as a little corner. But Christ whipped these merchants, or buyers and sellers, once or twice out of the temple: Antichrist hath brought them in again. And this is verily a wonder: they show more favor at this day to Jews, Turks, and heathens, than to Christians. For unto the only gospelers is no place permitted: verily for that they ascribe all to Christ, preach Christ only, leave nothing to the Pope, but rather accuse him most constantly, and bitterly.

### Paragraph 916

But what shall we say to those whose hands and foreheads have been defiled with the mark? I urge them to wash themselves in the blood of Christ, forsake Antichrist, and turn to Christ, abandoning their errors and repenting of them. If you have bound yourself to Antichrist through another, do not fulfill that rash and wicked oath by ungodly speaking against the gospel. Do penance, cleanse yourself, return to Christ, and you shall be saved.

### Paragraph 917

Now lest anyone here should prattle, so that we may be Christians and abundantly instructed in heavenly wisdom, even though we hear or speak nothing of the pope and papal matters, and consider those disputations unprofitable, yea odious, and pertaining to the stirring up of troubles, therefore hurtful and foolish: the Lord anticipates this, and says expressly, “Here is wisdom; in the knowledge and right judgment of these things consists the true, heavenly, and godly wisdom.” Unless we are wise in this thing, we shall be fools, and not wise. The Lord therefore excites the hearers to the study of inquiring after Antichrist, and to beware of him when he is found. For in the 14th chapter, we shall hear that they shall drink of the Wine of God’s wrath, as many as have received the mark of the beast and worshipped his image. Wherefore they shall drink at the same table with Christ of the cup of life and of the grace of God, so many as have despised papal authority. And who shall deny it to be the true wisdom, by which we may come from the wrath of God to the grace fellowship and participation of the same? Moreover, the Lord adds that those endowed with understanding, not witless and full of hurtful folly, should reckon the name of the beast; that is to say, should be diligently occupied in this matter, that those things should be diligently searched for, which worldly men affirm to be curiously sought and inquired after, [?not only without any profit, but with loss also]. Moreover, the Lord commands us to account the number of the name of the beast. He adds that the same is not hard to do, for this number to be the number of a man, to wit which a diligent man may easily by faith and industry attain to. For so does Jerome expound it also, saying that number is common and known to men. Let them leave then to trouble our godly studies, which blame our sermons made against the Pope, and laugh at our diligence such as it is in expounding papal abominations, finally which suppose we spend our time in vain in the account of times. They hear here, except they will hear nothing, that we have received commandment of the Lord so to do; moreover that the Lord testifies that wisdom is herein.

### Paragraph 918

And here I give warning that the manner of speaking is to be observed, [note 61.8] lest we and our listeners weary in vain through the insistent examination of a certain name in the numbers. For it is said to be the number of his name, as though a certain name should be gathered and composed of these characters: as most often it is gathered from these three letters, Christ. Nor do they err who think that by these three characters nothing else is signified than the name of Christ, which the Lord himself in Matthew 24 prophesied that Antichrist should use. Surely, he calls himself Christ’s vicar. I know very well that the proper names of great men have been sealed by prophecies and signified before: as Josiah, Cyrus, Jesus. But here you can gather no such thing, but forcibly and as it were against the text. I understand therefore by the number of the name of Antichrist or beast, the very acceptance whereby we come unto his name. And a name is a brief definition or description of any thing, whereby it is known of what sort and manner it is. Which thing in this our cause, the nobility supplies, which brings us unto those times which give him his name, from which he takes his name, that is which times reveal unto us Antichrist spoken of before in the prophets, and show us who and what he is, or who we should take for Antichrist, even him verily, which having brought low three kings, he himself starts up of naught, and to the destruction of the true religion begins to reign.

### Paragraph 919

And now he shows us expressly this number, which I may call nominal, and vocal, which may lead us to Antichrist, that we may know who it is; and when we know him, we may beware of him, and commands us to number the years six hundred sixty and six. For so many signify these Greek letters. In explaining this number, commentators have varied greatly. I prefer the explanation of the blessed Martyr Irenaeus, who, perhaps a hundred years after the publication of the Apocalypse, wrote his book against heresies, and encountered some who had heard John preach. Irenaeus also includes Andreas, the good bishop of Caesarea, who, with Arethas, speaks thus: the perfect reckoning and just account of the noble [?], as likewise other things which are written of the same Antichrist, the opportunity of time will reveal, and experience will confirm, to those who watch diligently. For if it were necessary, as some of the Doctors suppose, that this name should be manifestly known, he who saw him would surely have revealed it. But divine grace did not allow the name of this pestiferous beast to be rehearsed in this godly book. Thus far Andreas.

### Paragraph 920

After the same manner, the holy Martyr of Christ Irenaeus, before Andreas, wrote in the fifth book against Heresies. For at the end of the book, he says it is safer and without danger to wait for the fulfillment of the prophecy than to suspect and guess at every name, since many names may be found having the aforementioned number, whereby the question is not answered. Yet, he adds immediately, “the name contains the number 666.” And it is very likely to be true. For this term has a kingdom. For the Latins now reign. This is what he says.

And without doubt, this good doctor did not err in the slightest, being endowed with the Holy Spirit of God. For we see that the church of Rome is called the Latin church, and the Pope the high bishop of the Latin church. We see all services in churches performed in the Latin tongue. In courts and all judgments of bishops, the Latin tongue is used alone. Moreover, no one may serve in this church unless he is a Latinist. What will you say that these Latinists call the Hebrew, that is to say, the holy tongue, by a contemptuous name, “Jewish”? They call the Greek church and tongue heretical. The Bibles in Greek and Hebrew are suspected by them. For they insist that only the Latin Bibles are authentic and should be read by all as authentic. But these things are better known than that I need to admonish and recite them here with many words.

Nevertheless, this holy man Irenaeus does not wholly affirm this conjecture as most certain, although what he said was most probable and likely to be true. For he adds, “we will not be in danger herein” (for he also recited the name [illegible], the royal or tyrannical name of Nimrod), “neither will we definitively pronounce that he shall have this name,” knowing that if it were necessary for his name to be manifestly preached at this present time, it would certainly have been uttered by him who saw the Revelation. But he showed the number of the name that we might beware of him when he comes, knowing who he is. And he concealed his name, for it is not worthy to be preached by the Holy Spirit and so forth.

### Paragraph 921

Nevertheless, this same matter shows us how to understand those 666 years. For he says: knowing the certain number, which Scripture reveals to us—that is, the number 666—let the faithful remain steadfast or look for this: first, the division of the kingdom into ten, then the same power reigning and beginning to reform its affairs and to expand its kingdom. After this, one who comes suddenly will challenge a kingdom for himself and will put the aforementioned kings into fear, having a name containing that same number, so that he may be recognized certainly as the abomination of desolation. He says this again.

### Paragraph 922

But who does not see that the holy Martyr [Jerome] sends us to the prophecy of Daniel, which in chapter 7 says how the Roman Empire shall be divided into many kingdoms; and how in the midst of those kings a little horn should arise, which would overthrow and abase three horns; and that the same would begin to reign proudly, tyrannically, and wickedly, against both God and men, but chiefly to the faithful, intolerably?

### Paragraph 923

Let us then see how and when these things are fulfilled. [note 61.14] Where the Roman Empire had godly emperors, even then wicked Rome would not bow her stiff neck to Christ, but always most obstinately aspired to her old and wonted idolatry, which it coveted to have restored. And finally, when the fatal time was at hand, wherein the Lord, most righteous, thought to requite bloody Rome, he armed against her the Goths, Vandals, and Germans, which subdued and destroyed the lady of the whole world, and destroyed the whole empire. On this matter, seek more in sermon 57, and in the sermons following. [note 61.15]

### Paragraph 924

And it is evident by Histories that the Roman Empire, the Goths beginning to invade it, did decline, with provinces revolting in every place, and was divided into many kingdoms. For to speak nothing of Asia and Africa, the Persians ravaged the former, and the Vandals the latter, all Greece followed the Emperor of Constantinople, and likewise other nations never. The Westgoths possessed all of Spain, and the Franks subdued Gaul, Germany, and the nations adjoining to the same. The Eastgoths and Lombards obtained Italy. Thus, truly, many kingdoms were established, and in place of Rome reigned many kings. However, while these kings considered how they might best enlarge their kingdoms and put down and expel others, full craftily the Bishop of Rome played his part also. For he obtained supremacy over all bishops and so gained great authority with kings and realms, yea, and linked himself in league and amity with kings and princes. Whereupon quickly and suddenly, or as the Martyr of Christ prophesied, upon the sudden, he appeared, and at last usurped a kingdom, to wit, of Rome. For by his judgments falsely taken for apostolic, he put down King Childeric, of the lineage of Merovinges, the lawful king of France, and advanced Pipin, then Captain of the French guard, to the crown. And so he overthrew or plucked down one horn and bound unto him a most mighty king, by whose power afterward he was a terror to the kings of Greece and Lombardy.

### Paragraph 925

[note 61.16] Concerning the year of our Lord 269, the Emperor of Constantinople, expelling the East Goths, instituted a new government in Italy. But since this kind of rule and governance is not known to all men, I will briefly recite what and how great it was, by the words of Nauclerus the Historian, from his *Generations*, book 20. “Then began the City of Rome and Italy to have a new manner of governance, by which they lost more of the dignity, glory, and fear over all the world than of all the calamities which these 160 years have afflicted them, and at the last had left Rome to be inhabited of wild beasts.” For Longinus brought in a new name of dignity, the *exarchate* of Italy: that is, the high magistrate. Which keeping still at Ravenna, went never to the City of Rome. And in the government of Italy and of Cities, he kept first this order, that the president should not govern the province or region, but every City had their magistrate to govern them, whom he called Dukes. Wherefore making Rome equal with other Cities and Towns, in this thing only he honored the same, that he called the magistrate’s place in Rome, *praesides*. But they that did succeed him were called Dukes, as they were afterward for many years, so that it was called the Dukedom of Rome, as the Dukedom of Narnia and Spolet. Neither after Narses and Basil had it any more either Consuls, or Senate lawfully assembled: but by a Duke of Greece, whom the high magistrate sent from Ravenna, the commonwealth of Rome was governed a long time.” Thus much he.

### Paragraph 926

Any man may easily perceive here that the prophecies are fully accomplished [note 61.17], and the Roman Empire has fallen into ashes. For she who had been the most powerful ruler in the world is now seen to be a vile government, no more excellent than that of Spoleto and Narnia. And it is to be known that this Exarchate in Italy was the third lordship instituted since Augustulus was slain, in whom the histories say that the Empire of the west was finished and ended. For first, when Augustulus was slain, the Germans under their king Odacer possessed Rome. Afterward, the Ostrogoths, under the leadership of their Duke Theodoric of Verona, expelled and killed Odacer and reigned at Rome and in Italy. Last of all, the Ostrogoths were expelled and killed by the Lombards, and this Exarchate was instituted. And the Lombards, being called into Italy by the Greeks against the Goths, would no longer leave, for they saw the land fertile and rich, pleasant and abundant with pleasures [note 61.18]. Possessing great power in Italy, they subdued many cities and people of Italy, establishing now the fourth dominion, which they called the kingdom of the Lombards. They had powerful kings. However, that Exarchate of Ravenna, although they waited diligently for it and attempted to invade it, could never extinguish it, until the Bishop of Rome put his hand to it, pretending the sincerity of religion.

### Paragraph 927

Historians record sixteen Exarchs in succession, who ruled for approximately one hundred and eighty years. The fifteenth was named Paulus. Nauclerus states that in January of the twenty-fifth year, Leo the Third, Emperor of Constantinople, ordered that all those subject to the Roman Empire should tear down their images, break them, and burn them. Conversely, the Pope—some say Gregory the Second, others Gregory the Third—wrote to the entire world, instructing them not to obey these wicked commands of the Emperor. Platina writes more in the life of Gregory the Third: Gregory, with the consent of the clergy of Rome, deposed Leo the Third, Emperor of Constantinople, from both the Empire and the communion of the faithful, because he had torn down images. Nauclerus adds that the papal decrees held such great authority at that time that first the people of Ravenna, and then the people and soldiers of Venice, openly rebelled against the Emperor and the Exarch in Italy. And the treason increased daily. For Marinus Spatarius, Duke of Rome, and his son Adriane, while traveling through Campania, were slain by the Romans. In their place, they created one Peter as Duke of Rome. The people of Ravenna also, while some sided with the Emperor and others with the Pope, in a riot, killed Paulus the Exarch and his son. Thus writes Nauclerus.

### Paragraph 928

In these disturbances, the Lombards, believing the opportunity they had long desired had arrived, invade the lands of the empire through the agency of their king, Luitprand, and besiege Rome itself. But Pope Gregory, the instigator of all the unrest in Italy, the soldier and practitioner of it, and resembling neither a priest nor a preacher, summons Charles Martel, father of King Pippin, with his French warriors into Italy against the Lombards. However, this Charles persuades the king of the Lombards amicably to withdraw from the city. But not long afterward, Aistulf, king of the Lombards, plunders again the lands of Ravenna, renews the Italian war, and seizes Ravenna itself, demanding tribute from the city of Rome. But Pope Stephen II, who aspired to the governance of Ravenna and wished for the Lombards to be destroyed, requests aid from King Pippin of France, to whom, not long since, Pope Zacharias by his unjust judgment [? a false decree] (as many suppose) had given the kingdom, essentially offering him the kingdom. Therefore, the Frenchmen, armored and eager to conquer Italy, are present. While King Pippin entered Italy, he encountered the ambassadors of the Emperor of Constantinople, who requested that he restore Ravenna, the exarchate, and its lands to the empire, which rightfully belonged to it, and not to the Pope or the Romans. Pippin replied that he was warring for Saint Peter and the Pope, and to prevent the Lombards from harassing the church, and that he would take the exarchate and other territories of Italy from them and deliver them to the Pope, which he did. For he overcame King Aistulf, took from him the governance of Ravenna, and delivered it to the Bishop of Rome.

### Paragraph 929

Herein may all men see, unless it be those who will see nothing, how this contemptuous bishop, and very small horn, has at one push overthrown two horns. For he has put the emperor of Constantinople from the government of Italy and has put down the King of Lombardy, causing his people to be driven out of Italy. For a few years after, the Pope, by the force of Charlemagne, put down Desiderius, the last king of Lombardy, and destroyed withal the whole people of the Lombards. And thus the Pope arose and became as it were king of old Rome and of the chief part of Italy. And now are the beginnings of the kingdom laid, but as yet he reigned not with full authority, as is declared before. Eberardus, Bishop of Salisbury, whose words I recited in the preface of this book, extends these things further. But I suppose this our exposition to accord with the prophet, with the things and times. And the Pope gave to King Pippin for so great a donation a title, as Platina shows in the life of Stephen the second, that all kings of France should be called most Christian. Afterward was the Image of the Empire bestowed upon Charles, whereof is spoken before.

### Paragraph 930

And lest the Pope should seem to have received nothing, while King Pippin gave him the exarchate, the stories report thus: [note 61.22] the exarchate was divided into two regions, in Pentapolis and Aemilia. Pentapolis had five cities: Ravenna, Cesena, Classe, Forum Livii, and Forum Popilii. In Aemilia were Bologna, Reggio, Parma, Piacenza, and all the lands that lie from the borders of the Piacentines and Ticini to Adria, and from Adria to Ariminum, &c. But he who desires may read the Donation of Louis the Pious, in Volaterranus’s Geography, where he names the kings of France. We say nothing yet of how he afterward usurped power to himself over kings and realms, finally over all churches and souls, so that we must confess that a more marvelous prince never lived.

### Paragraph 931

You have here a brief and concise account, declaring how the Pope, having humbled and overthrown three kings, began to be made a king himself. But let us now apply the number of the name of the beast [note 61.23], so that it may be known to the whole world, that there is no other Antichrist to be looked for than the bishop of Rome, who has indeed laid the foundation of his kingdom under the emperor Phocas, built it under the kings of France, and enlarged it under the emperors Henry and Frederick, finally having established it under the emperors following; reigns in our time, and has done so for certain ages already past, &c.

### Paragraph 932

The calculation of 666 years must be reckoned from the time when John saw the Apocalypse. Irenaeus says it was seen not long since, but in a manner in our days, about the end of the reign of Domitian. And Eusebius in his chronicles says that it was in the year of our Lord 97. Therefore, there remain yet three years to complete a hundred years from the birth of our Lord. Add to a hundred years these years of the number of the beast, 666, and subtract those three years of the first hundred, and you shall have the year of our Lord 763, which was the 13th year, or thereabouts, of King Pepin’s reign, and the 7th of Pope Paul. Notwithstanding that there be writers of histories and times who attribute to Paul but one year, etc.

### Paragraph 933

Now must we not look only at what things happened in the very year 763, but what followed in the years before and after. Of this, I will recite a few things from the writers of histories and chronicles.

### Paragraph 934

Nauclerus, in the sixteenth generation, says that in the year of our Lord 1150, under Pope Zacharias and Emperor Constantine V, the twenty-sixth generation began, in which there was an alteration of the Kingdom of France, an abolishment of the Kings of Lombardy, and a translation of the Roman Empire from the Greeks. These great alterations were happily foreshadowed by the wonders that occurred at this time. In Mesopotamia, the earth split open for two miles, and a mule was said to have spoken with a man’s voice. Ashes fell down from heaven. There were wonderful earthquakes. Crosses appeared upon men’s garments. These things Nauclerus wrote. Similar events are recorded in the history of Eutropius, in the twenty-second book, under the year of Constantine VI; moreover, in Vincent’s historical chronicle, and in the Fasciculus Temporum.

### Paragraph 935

In the year of our Lord 751, through the counsel of Zachary the Pope, Pipin, master of the king’s household, oppressed his lord Hilderych, king of France, began to reign, and reigns. Eighteen years. This writes Aemilius in the second book of kings of France. And in the year 755, Pipin enters Italy with an army, vanquishes the king of Lombards, and gives the entire government of Ravenna to Saint Peter, [note 61.27] against the will of the Emperor of Constantinople. Vespergenensis in chronicles. You see how, instead of the Emperor, the Pope begins, in a way, to reign at Rome and in Italy, the horns are shaken off, according to the prophecy. Matthew Palmer in his chronicles, under the year 756, says that the Roman Empire revolted at a pace in the East, and the Emperor was persecuting the Christians (Idolaters he should have termed them). Pope Stephen gave to the kings of France the imperial titles and dignities, and confirmed Pipin and the successors of his lineage as kings only, all others utterly excluded, and in the name of the people of Rome, called him Patrician. Hitherto Palmer.

### Paragraph 936

John Functius in his Chronicle writes: In the year of our Lord 756, the rites and ceremonies of the church of Rome were carried into France and first received. In the year of our Lord 757, Paul is made Pope, and immediately follows that fatal year of our Lord 763, which falls as the midpoint between the year 750 and 770, or 773. Within this time, all those things occurred that both give the name to Antichrist, and whereby, just as everything else is known by his name, so too he receives his name and is known.

### Paragraph 937

In the year of our Lord 768, Stephen III held a council at Rome in the Lateran church with the bishops of France and Italy: [note 61.30] and decreed that no one should be ordained bishop of Rome except a Cardinal. He also condemned the Greek council of the Emperor Constantine against images, which he commanded to be possessed and worshipped. These things writes Antoninus in Chronicles, title 14, chapter 1 and 5.

### Paragraph 938

After this, that great Charles, son of Pipin, was summoned to Italy by Pope Adrian [note 61.31], took Desiderius, King of the Lombards, and extinguished the Kingdom of the Lombards. This occurred in the year of our Lord 773, and also the two hundred and fourth year after the Lombards arrived in Italy. He confirmed and augmented the donation of Pipin his father, as many historians relate. John Functius in his Chronicle adds that throughout the entire realm of France [note 61.32], at the command of Charles, the ceremonies of the Roman church were instituted. Now we have the name of Antichrist, of the number 666. We know who he is, and whom we should beware of.

### Paragraph 939

I cannot here omit, but must note a few words concerning the Sibyl’s prediction regarding the origin of the Antichrist, which, in my judgment, is very consistent with what has been said before. For the eight books of the Sibyl’s oracles, taken from the library of the honorable commonwealth of Augsburg, were published by the most godly and learned Doctor Sixtus Betuleius in the year of our Lord 1545, and that in Greek. And this Sibyl of Erythrea, or whatever she was, prophesies in the eighth book that Rome shall fall and be burned with fire. The words of the Sibyl in Greek are to this effect.

### Paragraph 940

And this event occurred at the time Totila, King of the Goths, burned the city, as we have recounted previously. Shortly thereafter, these words are added in the same prophetic writings: When emperors have oppressed the world with great bondage from east to west, the number fifteen will be fulfilled. Then a king will appear wearing a white hat, who will never add his name to Pontius [?], as if he were to be called Pontifex. He will live for worldly pleasures and will reward with his wicked foot. And the rest that are read there.

### Paragraph 941

She bids us count fifteen kings from the burning of Rome. After whom a new king will come, whom she describes. And it is manifest that Rome was taken, plundered, and burned under Emperor Justinian. After Justinian the Younger, fifteen kings are counted to Emperor Theodosius. After Theodosius, Leo III succeeds, whom she calls “delicatos,” that is, given to pleasures. Because most of them were not very valiant, but under Leo III, Italy revolts from the emperor. And shortly also the government called the Exarchate was given to the pope by King Pippin, against the emperor’s mind. We see therefore that the reckonings agree. For we have also brought the years to King Pippin, 666. And so a new king arises, whom the Sibyl names [illegible] notable because of his white hat or mitre. For so she notes the bishop (who in old time wore white miters on their heads) that should be a king. She gives him a name also. For she says how he has a name never unto Ponti. For add to the word Ponti, *fex*, and you have Pontifex. She annexes certain notes or marks also: that he shall regard earthly things, and not heavenly; and that he shall also provide [illegible] and give rewards with his ungracious foot. And that is rightly spoken, since after Domitian and Diocletian, none of all the kings, save the Pope, has offered his foot to be kissed: whereby fools think they receive great rewards. But omitting these things, let us return into the way.

### Paragraph 942

The blessed martyr Irenaeus, speaking of this king, in the same fifth book, says that in the beast’s coming, there is a recapitulation of all iniquity and all deceit, so that all apostatical power concurring and concluded in him [note 61.33] might be thrown into a furnace of fire. And that he has spoken this thing by the spirit of prophecy, all men will confess who have read the lives of the Bishops of Rome: but especially of Sylvester II, Benedict IX, Gregory VI, Gregory VII, Urbanus II, Paschalis II, Alexander III, Innocentius III, Gregory IX, Boniface VIII, Clement V, John XXII, to speak nothing of various others. What has been done in our days by Julius, Clement, Leo, and Paul? Spain, France, England, Hungary, and Germany, and other realms speak, which have been set together by the ears and entangled among themselves with most cruel wars. The blood of martyrs shed speaks, which cries unto the Lord. What remains therefore, but that we should take heed to ourselves and beware of this man of sin, and cleave to our Redeemer Christ our Lord, beseeching him that he would come shortly and deliver us from all evil. Amen, Amen.

## Sermon 62: ¶ Christ stands upon Mounth Zion, having his church: and is deſerbed by notes, which and what shall be the sheep of Christ.

### Paragraph 943

.

### Paragraph 944

And I looked, and behold, a Lamb stood on Mount Zion, with him one hundred forty-four thousand, having the name of his Father written on their foreheads. And I heard a voice from heaven, like the sound of many waters, and like the voice of a great thunder. And the voice that I heard was like the sound of harpers playing on their harps. And they sang as it were a new song before the throne, and before the four living creatures, and the elders, and no one could learn that song except the one hundred forty-four thousand who were redeemed from the earth. These are they who were not defiled with women, for they are virgins. These follow the Lamb wherever he goes. These were redeemed from among people, the first fruits to God and to the Lamb, and no deceit was found in their mouths. For they are without blemish before the throne of God.

### Paragraph 945

Like as he has hitherto mixed joyful things with sorrowful,[note 62.2] and annexed a consolation to most hard and cruel chances: so now he adioyneth also to the tyranny of the Roman Empire an exposition having both a consolation and an exhortation most grave and weighty. Undoubtedly, by the description of the Romish tyranny and reign of Antichrist, it might have seemed that the Church and the preaching of the Gospel had been utterly lost, and that ungodliness should have triumphed forever; he declares therefore by a most excellent vision how Christ shall reign notwithstanding in his chosen, and shall overcome, and shall have his church continually, and that right famous. He describes what the elect shall be. He adds that the preaching of the Gospel cannot be so oppressed, but that it shall rather be preached with great constancy throughout all the world. And that Rome also shall fall, and all the ungodly be punished. He exhorts therefore most earnestly that we have nothing to do with Antichrist, lest also we be made partakers of his damnation. And to the intent there might want nothing that concerned a full comfort, he adds that thing which may chiefly confirm the minds of all the godly even in the greatest dangers, how they that die in Christ do flit straightways from corporal death unto life everlasting. Which finished, he turns to the description of the punishment to be taken assuredly of the Antichristians. Wherefore if the Books of the Gospel and New Testament are to be esteemed for the manifold description of Christ, and of salvation by him obtained for the faithful, if they are to be esteemed of the comfort and preaching of the gospel: this is doubtless a book most gospel-like, as that which by a continual tenor to perilous things annexeth consolation.

### Paragraph 946

John therefore sees the Lamb standing upon Mount Zion. Christ therefore sleeps not; he is not ignorant of the perils and conflicts of his church, but he stands prepared to aid and succor his people. He stands as an invincible king, whom neither the Dragon, nor the old beast, nor the new beast has overthrown. For I have told you often, especially in chapter 5, that by the Lamb is understood Christ. For he is the Lamb and price of our redemption until the judgment, but laying aside the office of an intercessor, he shall be a most severe and also a most holy judge. And Christ stands, not in the sand, as did the Dragon, but on a Mount, and that upon Mount Zion. Mount Zion was a figure of Christ's kingdom, as appears plainly in Psalm 2 and Isaiah 2. And the kingdom of Christ is the church, both triumphant and militant; therefore, in the fellowship of saints, stands Christ, the joy and glory of those in heaven, and the life and helper of those who fight yet on Earth. Let us believe therefore that in the Antichristian persecutions Christ will never fail his faithful, as he is known never to have failed the old saints under the old Roman Empire afflicted. For this consolation serves chiefly for us, who are vexed by Antichrist, and served for them also, who suffered martyrdom under the old Roman Empire. Neither is there any doubt but that they confirmed themselves with this in the greatest persecutions.

### Paragraph 947

It is most comforting to know that the Lamb is not alone, but has with him one hundred forty-four thousand—a vast and complete church. Therefore, even though the beast rages and slays the confessors of Christ, there will always be a church that cannot be uprooted from the earth. He sets a specific number for something that appears uncertain, yet is certain and determinate: for the number of those who will be saved will seem small compared to those who worship the beast and perish. Nevertheless, we understand that the number of those who will comprise the body of the Church, under their head Christ, will be greatest, even at the time when the Pope and all the agents of Antichrist have unleashed all their fury. I have spoken of this number of the elect in Chapter 7, where the same number is given.

### Paragraph 948

And just as those who oppose Christ bear the mark of Antichrist on the right hand and foreheads, so truly the sheep of Christ, which will be the church, the bride of Christ, under their head Christ, will have their mark also on their foreheads, namely, the name of the Father of the Lamb. For "Eius" is to be referred to the Lamb. And he speaks not of an external mark, which should be printed on their foreheads, but of the mark of their minds. This is faith, the sign of all God’s children. And faith in the Father and the Son, which is not without the Holy Spirit. And how could you believe that almighty God is your father, unless you understand the same to be obtained through the Son? This faith, therefore, is here understood to be a Christian faith, not a Jewish or Turkish faith, which yet confesses God to be the Father. But since they do not have the Son, as Saint John says in his epistle, they neither have the Father. Therefore, the true members of the church of Christ, the true sheep, do believe that they have a merciful Father through the Son, by whom they know that the Father, being pacified, has given all things of life and of salvation in his Son. Those who do not seek salvation and all goodness in the only mediator, the Son of God, do not, without doubt, have the true mark of the children of God on their foreheads. At this day all will claim to be Christian, but neglecting Christ, they depend wholly on saints. Therefore, their faith is not the true mark of the children of God. No, they neither know the Father nor the Son. And therefore they persecute those who cleave wholly to the Father through the Son. And saying that Christ is with his church, what need has the church of a vicar? Certainly, it cannot be the true church, which has a vicar of Christ, for then it lacks Christ, whom the true church cannot lack.

### Paragraph 949

It was not enough for the Apostle to have simply said that the church is united with Christ; unless he had also added, with many words, how he has seen the church afflicted and how she has behaved when the beasts torment her, then truly, so that we may learn from it, what the hope of the saints is in the greatest dangers, and of what sort we ought to be in persecutions and temptations.

### Paragraph 950

First he hears a voice from heaven, as the voice of many waters. Waters in the Scriptures many times signify people. We understand therefore hereby that the church shall be populous and speaking, intending to dissemble nothing, but freely to profess Christ. And therefore he hears also the sound of a great thunder. For the church receives from heaven power to preach and show forth the Gospel gravely, though the world’s bowels burst. And verily, concerning the fervent and constant preaching of the gospel, John and James are called with Mark, sons of thunder. And concerning the preaching of the gospel, more shall follow afterward. He hears moreover a melodious harmony of men singing to their harps, and singing as it were a new song. This is chiefly referred to the saints in heaven, singing eternal praises to God; secondly, to the saints living here yet in earth, who also offer unto God continually praises and thanksgivings. Therefore, however their hearts be made sorrowful in perils and adversities, yet their spirit rejoices in the Lord. For no man could learn that same song, except the elect. For just as none of the heavenly dwellers can express or understand the excellency of the joys of the life to come and the praises of God, except he dwell among the heavenly inhabitants and be partakers of the most godly life; so except any man living yet here in earth be regenerated, he neither sees how great is the felicity of the faithful, nor can he justly esteem the praises which they offer unto God. Touching the new song, I have spoken in chapter 5. And certainly, to worldly men, the things seem as they were new, which the faithful bring forth of God’s word.

### Paragraph 951

Now John describes what kind of people the sheep of Christ will be, those who will remain in the church of Christ, disregarding the rage of the beasts. To them also the mark of their fathers’ name in their foreheads is explained.

We shall also perceive what the true marks of the faithful are. First, they are redeemed from the earth. Doubtless, all of us bearing the earthly image of the earthly man were sold into sin, for which cause we are also subject to condemnation. But the Lord has bought us with the price of redemption, paid upon the cross, so that now we are reshaped into the image of the heavenly man, and are adopted as the children of God. Concerning this redemption, the Apostle has spoken in 1 Corinthians 7 and to the Romans 3, and in other places. Peter also, in 1 Peter 1. And because the faithful know themselves to be bought and adopted by Christ for the heavenly inheritance, they are eager to serve their Redeemer only, and inseparable cleave to him.

### Paragraph 952

Moreover, they are virgins [note 62.9], not defiled by association with women. Some who discuss these matters torment themselves, fearing that something here might seem to diminish the value of holy matrimony, by which, without doubt, the Apostle testifies in 1 Corinthians 7 and to the Hebrews—no one is defiled. I am ashamed to bring forth the trivialities of the Papists. For who could hear the most corrupt reasoning about purity? They will use this to justify and defend their single lives; but all men see, except they be more blind than moles, what filthiness has been committed and is committed daily under the guise of this ungracious and most unclean singleness. But the Lord speaks nothing at this present of physical marriage, but rather spiritual. For it is evident that the apostles, as the bride leaders of our Savior, have brought the church to our Savior as a chaste virgin, which has not had dealings with any strange or foreign woman—that is to say, which is not defiled by the participation of evil doctrine. Read Solomon reasoning about that woman seriously in chapter 4 of Proverbs. Read moreover the Apostle in 2 Corinthians 11, teaching exceedingly well that the faithful are an undefiled virgin, the spouse of Christ. The faithful therefore who lived under the tyranny of the beasts received no strange doctrine of idols and of other profane cults, neither do they at this day admit the Popish infection, but keep their maidenly minds for their husband, Christ, being betrothed to him by sincere faith.

### Paragraph 953

These follow the Lord, wherever he goes. That is to say, they care for no person but Christ; they desire no person but Christ. In him they place all their help, all their comfort, all their joy, all their salvation; to him alone they always give respect, in him they know themselves to be complete, which one and alone is to them all things. Moreover, wherever Christ calls the faithful by doctrine and example, even to death and most cruel slaughter, they follow willingly and cheerfully. By this means, they can never be separated from him in the world to come. For wherever Christ is, there is also Christ’s minister, as he himself witnessed in the twelfth and fourteenth chapters of John.

### Paragraph 954

They are also redeemed from humankind,[note 62.11] truly delivered through the grace of Christ, that they should not follow this corrupt and unclean world, by all manner of pollution. For Christ by his Spirit and word calls his out of this world, that although in body we are conversant in the world, yet we should with all our mind abhor the world and the things that are therein. Furthermore, for this intent he has chosen and redeemed his from the bondage of humankind or of the world,[note 62.12], that they should be first fruits to God the Father and to his Son. Which place the most godly and excellent learned man Dr. Francis Lambert expounding in his commentaries upon the Apocalypse: it is manifest, he says, by Leviticus 23:15, Numbers 18:15, and Deuteronomy 18:18, what first fruits are, and that they were gathered for the Lord and went to the high priest. But Christ is that high priest, to whom the spiritual first fruits appertain, namely, the godly and sanctified to God. These things are confirmed by the Apostle, who said that Christ gave himself for us, to the end he might redeem us from all iniquity and might purify us to himself, a people for good works. Therefore, the true faithful singularly apply this to godliness, that they may be the first fruits and a most excellent offering to the Lord, since they know themselves to be redeemed for this end, that all the rest of the time of their life they might serve God.

### Paragraph 955

In their speech, there is no deceit. He does not say that no desire or evil impulse is found in the hearts of the faithful, but that there is no deceit in their speech. For although the faithful are troubled and vexed by the affections of the flesh, they love the truth so much that, to their knowledge, they will deceive no one. And they chiefly dissemble nothing that pertains to the confession of truth, nor use any deceit in the doctrine of the Gospel.

### Paragraph 956

Moreover, they stand without blemish before the throne of God,[note 62.14] not by their own merit, but by the sanctification of Christ, as Paul also affirms in chapter 5 to the Ephesians. He spoke appropriately of them standing before the throne. For Augustine said that our sanctification will ultimately be perfected in the world to come.

### Paragraph 957

These are what I say are the true marks of the true faithful, and of the true church of Christ. Let every man search the secret corners of his heart and consider diligently in his mind whether he is marked with these signs. Let him earnestly pray to God that, if he feels them, the Lord would confirm them; if he does not feel them, that the Lord would impress them deeply in their minds.

## Sermon 63: ¶ The Angel preaches the eternal gospel of Christ.

### Paragraph 958

.

### Paragraph 959

And I saw an angel flying in the midst of heaven, having an everlasting gospel to preach to those who sit and dwell on the earth, and to every nation, kindred, language, and people, saying with a loud voice: Fear God, and give honor to him, for the hour of his judgment has come. Worship him who made heaven and earth, and the sea and the fountains of water.

### Paragraph 960

Antichrist desires nothing so much as to be oppressed by the preaching of the Gospel. For even therefore has he instituted inquisitors of heretical depravity to dare call the Gospel heresy. Therefore he burns Gospel books and preachers of the Gospel, and everywhere restrains the reading of the Gospel and Evangelical books. Wherefore the simple suppose that it can not be but that the Gospel with all its adherents should perish utterly. Now therefore, in the Lord’s consolation, is brought in a vision of an Angel (for he is still in the vision), flying in the midst of heaven, having the everlasting gospel, and preaching to the world. Whereby is signified that the Gospel shall be preached unto men, in spite of all the enemies thereof. And he gathers a brief summary of such things as by the Gospel are preached to the world. Those same appertain also to the comfort of the church, which under the old beast suffered persecutions for the Gospel. We will briefly consider every thing.

### Paragraph 961

First, it is evident even from what has been said that “angel” signifies the ministers of the word, and the ministry of the gospel itself. Scripture calls preachers “angels.” John the Baptist is so named from the prophet Malachi, as was explained earlier. By this honorable title, ministers are reminded of purity and most sincere faith. For angels are God’s ministers, who regard, love, and honor only Him, and execute His commands most faithfully, sincerely, and diligently. Preachers should be like this in their kind and office. Just as angels cannot be harmed by the treasons and injuries of men, so God defends His ministers until the appointed hour. Peter was delivered from prison in Acts 12. Paul was protected in the shipwreck, and so on. And another angel is mentioned, because he has already introduced various visions of different angels. Nevertheless, “another” seems to refer to the first, for he adds two more angels, the first of whom he calls “another,” and the later, “the third.”

### Paragraph 962

And this angel flies in the midst of heaven. By this is signified the fortunate course and progression of the preaching of the gospel. It is also written in the prophets, that his word runs swiftly. Psalm 19. David compares the running of the gospel’s preaching to the course of the sun, joyful as a giant he runs his way: he arises to the farthest part of the heavens and runs again to the same, and neither can any man stop him, nor hide himself from his heat. The sun shines in all places. Therefore, the preaching shall be free. For as we cannot pluck back or hinder the things that are above us in air and sky, so shall we neither pluck down nor hinder him that flies in the midst of heaven. The words and writings fly, they fly far and wide. Neither can the truth be oppressed. God has given to the world Printing, whereby the gospel is preached and runs far, wide, and most swiftly.

### Paragraph 963

And this angel has the everlasting gospel. Wherein is the greatest comfort. For it signifies that the truth will be in the world invincible. And for many reasons the gospel is called everlasting. First, because the truth is immortal, which cannot be bound, however the ministers are fettered and slain. 2 Timothy 2. Secondly, the gospel is eternal, for because it was shown to our first fathers, prophesied in the law and prophets, fulfilled by Christ, declared by the apostles, and by the grace of God brought unto us. Indeed, it was predestined before all times. Read 1 Ephesians. For this cause it is called everlasting, as much as it pertains to us and to our posterity until the end of the world, and not only to our elders. And because it is everlasting, those lie who at this day call it a new doctrine or learning. Papistry is new, which has its origin when everything was ordained. Moreover, the Apostle says: If I or an angel from heaven should preach any other gospel, or besides the one that you have received, let him be accursed.

### Paragraph 964

And we hear expressly that the angel had not only the gospel, but that he had preached the gospel. Many indeed have the gospel, but it is silent and written in books. The gospel must be shown forth and proclaimed. He declares also to whom the gospel must be uttered and preached: to the inhabitants of the earth; for it must be cried out to those who are drowned in earthly matters, and they must be raised out of their sleep. And after his manner and imitation of blessed Daniel in the seventh chapter, he lists nations, kindreds, tongues, and peoples, thus signifying that the gospel shall be preached throughout the whole world. Which thing the Lord said also should come to pass in Matthew 24, and then that the end should come. And we see at this day that the gospel has in a manner thundered throughout the whole world. And here I give warning lest any deceive himself. The Apostle in 1 Timothy 3 and 1 Colossians says that the gospel was preached throughout the whole world in his time. However, not all men had then received it, but a few. Do not therefore imagine with yourself that the gospel is not preached unless all receive it. They are deceived who promise themselves before the judgment a concord of all nations, for it is written that there should be one shepherd and one flock. For the same was accomplished when the Lord prepared to himself one church from the Jewish synagogue and the dispersion of the gentiles, of which Christ is head and pastor, and Antichrist shall at the length by his last coming be abolished. Therefore he shall always resist Christ.

### Paragraph 965

Furthermore, where he sees and hears this Angel preach the gospel with a loud voice, he means that the preachers shall with great constancy and frankness, also with shrill voices and utmost earnestness preach the gospel against Antichrist. And we see at this day that the more cruelly the faithful are grieved and oppressed, the more fervently and louder they cry, and that they also are called clamorous criers.

### Paragraph 966

Moreover, he includes in a brief summary what things are to be presented in the preaching of the gospel, especially in the last times. First, he says: fear God. The fear of God is the beginning of wisdom; therefore, not fearing God is the beginning of foolishness and of all errors. The fear of God has nothing in common with the fear of the world. The godly man is not afraid of God as a guilty servant fears his master, and that punishment more than his master, whom he hates rather. For the fear of God has the reverence and love of God. It attributes to God supreme majesty, embraces faith, and has a faithful care, whereby it waits upon God, worships, praises, and professes him. Doubtless, because we fear men more than God, we fear more the Pope and the malice and hatred of him and his, therefore we do not execute justice uprightly, nor profess the truth frankly, neither yet set forth the gospel. But the Lord in the Gospel says: fear them not, which may kill the body and have no power over the soul; rather fear him, which condemns both body and soul to hellfire. Certainly, the fear of God is not only the beginning but also the bond of all virtue. Hereafter we shall hear that the fearful shall be cast into hell, with the beast and with the false prophet. Therefore let God be our fear, as Isaiah teaches in chapter 8. Let us fear God for our sins committed. Here few are afraid; but many are afraid to speak the truth, to maintain godliness, and to rebuke wickedness.

### Paragraph 967

Secondly, the preaching of the Gospel comprehends the honor of God. For he says, "And give him honor." [note 63.9] And you do not separate the Son from the Father. For he, in John 5, says thus: "The Father has given all judgment to the Son, that all should honor the Son, as they honor the Father. He who does not honor the Son does not honor the Father who sent him." And indeed, the Father cannot be honored but by the Son. For we honor him when we believe him to be true and receive Christ, the Son of God, as the only righteousness and perfection of all the faithful. By faith, therefore, we chiefly honor God, then revering him by faithful obedience and walking in his commandments. John, in his Canonical writings, says that he who does not believe in the Son makes God a liar—see how greatly one may dishonor God—who does not believe the testimony which God has testified about his Son. And this is the testimony, that God has given to us eternal life, and this life is in his Son. He who has the Son has life; he who does not have the Son does not have life. We are therefore forbidden, elsewhere than in Christ alone, to seek life and all goodness. But the papists honor the Pope, and his authority, his saints also, and honor God alone. They carve into their tombs, *Soli deo gloria*—to God alone be glory—but yet in the meantime they persecute those who will not ascribe the glory due to God alone to their foolish trifles. But the Gospel will cry out that to God alone all glory is due.

### Paragraph 968

Hereunto is added a spur, which may prick them to fear and glorify God: for the hour of his judgment is come. The Gospel therefore in the latter days shall impress upon men the final judgment. This has a wonderful effect in obtaining a reformation of life. And it is purposely said, it is come: and not, the hour of his judgment shall come. For so is the certainty of his judgment expressed, and we are warned to look for that same day every moment. The Apostle used the same argument in the 17th chapter of the Acts, addressed to the Athenians, and to the Corinthians in the 2nd Epistle, 5th chapter. Let us remember, O brethren, that strict judgment; let us amend our faith and manners, and all things that do not agree with the gospel. For certainly we shall die, certainly we shall be judged. But then, when we shall promise ourselves peace, sudden destruction shall come. Watch.

### Paragraph 969

Finally, the Gospel teaches us to worship God alone. Do not the faithful then worship idols, wherever they may be erected? They do not worship the Pope, who is overwhelmed with wickedness; much less do they kiss, and by kissing worship his ungracious and foul feet. They do not worship the god Mauzim of the wafer makers, the god in the box, which is worshipped in palaces and churches as shut up in the pixe. They do not worship saints, but God alone. Therefore, lift up your hearts to heaven and worship. We have here on Earth wonderful works which may move us to worship this God alone. He is the maker of heaven and Earth, and of the sea. Who is greater? Who is mightier? Therefore, worship him, the true God, as in Matthew 4. He adds here fountains of water, for the miracle and benefit of water is great. For if we consider the original spring, substance, pleasantness, and usefulness of fountains, we shall be compelled to wonder. God be praised.

## Sermon 64: ¶ An other Angel preaches, that Babylon shall fall: and an other diſwadeth all men from the felloweſhippe of the religion of the beast.

### Paragraph 970

.

### Paragraph 971

And there followed another angel, saying: “She has fallen, she has fallen—Babylon the great city—because she has made all nations drink of the wine of her fornication.” And the third angel followed them, saying with a loud voice: “If anyone worships the beast and his image, and receives his mark on the forehead or on the hand, that person shall drink of the wine of the wrath of God, which is poured out in the cup of his wrath. And that person shall be punished with fire and brimstone before the holy angels and before the Lamb, and the smoke of their torment ascends forever. They have no rest day or night, those who worship the beast and his image, and whoever receives the imprint of his name.” Here is the perseverance of the saints. Here are those who keep the commandments and the faith of Jesus.

### Paragraph 972

For the comfort of the faithful flock of Christ, [note 64.1] another angel appears, a type of all godly preachers, who proclaims with great constancy that the kingdom of Antichrist will fall, however it promises itself everlastingness. From this, the saints gather that persecutions will certainly end, along with all other abominations throughout the world. For, considering the continual persecutions of the wicked, the saints cannot but be marvelously saddened; they must needs receive no small joy and comfort from this understanding, that they here will not endure forever.

### Paragraph 973

and here it is said that Babylon shall fall. And in deed it were foolishness to expound these things of the old Babylon in Asia, which was fallen long since, scarcely any token thereof being left; we must therefore understand it of another, which is in her flower, and even by a figurative speech, we must understand it of Rome. For there is a great correspondence as it were, betwixt both: Babylon was the first monarchy, Rome is the last. Babylon sore afflicted the people of God, so doth Rome grievously vex the church of God. Babylon burdened Israel with a grievous captivity: so Rome vexes the church, with more than a long captivity. Babylon overcame the people of God, and burning the City of Jerusalem, and destroying the temple, led away Israel captive: so Rome also having razed the city of Jerusalem, and subverting the temple, triumphed of Israel. Babylon planted idolatry, superstition, and all abomination, advanced, maintained, and set forth the same unto all men: but at the length when she would have thought least of it, the people of God being suddenly delivered, she was utterly subverted. So is Rome also, the mother and nurse, and reviver of all abominations in the church of the last time, wherein she shall perish at the last, all those that believe truly in Christ being delivered. And especially it is called great. For how great and mighty the church of Rome is, all we see and by experience know at this day. Neither am I the first that understand by Babylon, Rome. For many expositors reading the first Epistle of Saint Peter, in the end of the epistle, do understand by Babylon, Rome. Certainly Oecumenius says: And here he calls Babylon Rome, for the excellency and brightness of the Empire: the which Rome obtained a long time since. But this the more ancient writers expound more plainly, as Tertullian in his book against the Jews, which says: so Babylon with Saint John bears the figure of the City of Rome, therefore also great, and proud in her kingdom, and a murderer of the saints. The same words in a manner, he repeats in the third book against Marcion. And no less plainly Saint Jerome calls Rome Babylon: and that same Babylon whereof S. John speaks in the Apocalypse. Read the epistle of Paula and Eustochium written to Marcella, by the help of Saint Jerome. Read himself in the 11th question to Algasia. Again in the preface to the book of Didymus of the Holy Ghost, to Pauliniane. Also in the end of the 2nd book against Jovinian. The same in the life of Saint Mark: Peter, saith he, in the first epistle, under the name of Babylon, doth figuratively signify Rome. But Saint John will expound himself in the 17th chapter. And we understand that the City of Rome shall fall chiefly, with all her ungodliness: and with the same also, the Roman superstition and abomination, throughout the world. And the Angel in deed, says she is fallen, which is yet to fall: and that by the prophetical manner of speaking, wherein that which shall assuredly come to pass is uttered, as though it were now done. To signify the certainty thereof, that reduplication or iterating of the word also appertains, she is fallen, she is fallen. This is also repeated in the 18th chapter, where it shall be showed how it is taken out of the Prophets, &c. Notwithstanding both a desire and joy also, might seem here to be signified. For such things as we have long, and with a desire looked for, we receive now coming and say, thou art come, thou art come at the last long looked for, and now make me glad. For the saints with a great desire, look and long for the destruction of that most wicked, most vile, and most troublesome kingdom of Antichrist.

### Paragraph 974

The cause of the destruction of the commonwealth and the church of Rome is also shown, for she has caused all nations to drink, and has made them drunken with the wine of wrath of her fornication. Indeed, the effect of wine on people is greatest. Therefore, doctrine is compared to it in the Prophets. Thus, Rome, with her unclean and corrupt opinions, has made all people drunken. And it is called the wine of wrath. For, as God is angry with those who err in the way of the Roman church, he allows them to do so. Because God has revealed the sincere doctrine of life through his only Son and most chosen Apostles, and people do not receive it, God is justly offended with them and gives them over to a reprobate mind, that they may follow shameful errors, as Saint Paul prophesied would come to pass in 2 Thessalonians, chapter 2. This wine is also called the wine of her fornication, by which she herself, having first played the harlot, has now become the master of fornication, and as it were a bawd to all others. This manner of speaking is well known even from the Prophets. Rome did not persist in the doctrine of the Gospel and of the Apostles, but invented a new one, and that contrary to the Gospel, concerning the vicar of Christ on earth, of the power of keys, of indulgences and pardons, of justification by works and merits, of satisfaction and confessions, of the worship of images, and praying to saints, of celebrating masses, and worshipping the sacrament of the altar, as they term it, of monasticism and vows, and such other innumerable things. This doctrine, as she presents it as Apostolic, ancient, and Christian, she offers to all people, and so plucks them from Christ, withdraws them from the Gospel, seduces them from the old Christianity, and destroys innumerable souls. Therefore, God pours out to her also from the cup of his wrath, and brings her also to destruction forever.

### Paragraph 975

And on this occasion he dissuades all men from fellowship with the Roman church or papistry, that we have nothing to do with the Roman religion unless we will be partakers also of the everlasting punishment. He reasons therefore of the loss and punishments, and describes grievous and horrible pains, if perhaps men might so be frightened from that ungodliness. The angel therefore cries, and that with a loud voice. Wherefore let all ecclesiastical preachers learn that they must earnestly and terribly cry out in this case, that all flee the communion of the Roman or popish church. I know doubtless what the common people believe and say, that all shall be saved at the last day, whatever religion they be of; and especially if any remain an open papist. But we can neither condemn nor absolve any man, set them in Heaven, or cast them to Hell. God is a righteous judge. He alone knows who shall be saved or damned. We ought therefore of right to credit his judgments. But whereas he pronounces openly that the favorers of the Roman church shall be damned, who am I to say the contrary, or what men will pronounce otherwise? Let us hear therefore the sentence of the just judge, and let us believe the word of the Son of God, and let us beware of the popish religion.

### Paragraph 976

I have already explained in chapter 13 what it means to worship the beast and his image, and what it means to receive the mark on the forehead and on the right hand. Briefly, they worship and receive the mark of the beast, those who participate with the Roman church or religion; finally, those who obey the wicked decrees of the empire and persevere in obedience to the See without repentance. Arethas, interpreting this passage, says that to worship the beast and to receive his seal is to consider Antichrist to be God, and in word and deed to promote the things he desires.

### Paragraph 977

And here in a horrifying manner, and with prophetic words, is described everlasting damnation, prepared for those who, forsaking Christ the Savior, cleave unto Antichrist the destroyer. Like as they have drunk of the corrupt doctrine infused by the Pope, so again shall they drink, that the just Lord shall pour out of the cup of wrath. And the wine that is poured into the cup of God’s wrath is the straight, exquisite, and most grievous judgment of God, wherein, being angry, he inflicts upon the Antichristians horrible and unspeakable punishment. A like manner of speech is read in Jeremiah, chapter 25. And like as pure wine, not delayed, is of most efficacy and pierces: so the judgment of God, wherein he will proceed against the Antichristians, shall be most grievous, such as no tongue, be it never so eloquent, can express.

### Paragraph 978

And for a further declaration shortly after follows what they must drink: truly fire and brimstone. Perhaps the Lord alluded to these words of David in the eleventh Psalm: Upon the ungodly he shall rain snares, fire, brimstone, storm, and tempest; this reward they shall have to drink. He seems moreover to have alluded to the burning of Sodom, and to the thirtieth chapter of Isaiah, in the end of which it is shown that hell shall be wide enough to receive all the ungodly, and that matter shall never want to nourish the fire, nor the fire be quenched. He expresses moreover a grievous pain, where he says that they shall be tormented, and that in the sight of the Lamb and holy angels, that so they may receive deserved punishment forever for their contempt, whereby they have despised the Lamb and messages of angels. Likewise in Luke 13, the Lord says: There shall be weeping and gnashing of teeth, when they shall see Abraham, Isaac, and Jacob, and all the prophets in the kingdom of God, and you to be shut out, and so forth.

### Paragraph 979

And that same pertains also unto evidence, and to stir up a terror in the minds of all men, where he adds by a figurative speech: and the smoke of their torment ascends up evermore. Therefore, the burning and punishment of the ungodly will be everlasting, and never to be finished world without end. And we seem here at this description, as it were before our eyes to see the flames of eternal damnation carried up on high, and cast up with them great heaps of smoke, to roll up and disperse them far and wide. I remember here that of Virgil.

### Paragraph 980

And that no kind of terror might be lacking, he expresses most aptly and most abundantly the perpetuity of everlasting punishment, saying: neither have they rest day nor night. So says the Lord in Mark 9. Their fire is never quenched, and their worm shall never die. They err, therefore, who promise to the damned deliverance from their torments after many worlds.

### Paragraph 981

And it is not in vain that he repeats what he had said before, how those who worship the beast shall suffer these things. And therefore he repeats it, lest, as it happened, we should consider it a trivial matter. They shall be condemned, says the truth, who receive the Papist cult and religion, and persevere in the same. [note 64.13] To all this is annexed an acclamation, or double sentence, notable and wholesome. For inasmuch as the wisdom of God did foresee what adversity remained for the godly in this world, which they might surely expect at Antichrist’s hand, who professed the truth, therefore for a comfort and consolation he adds: here is the patience of the saints, which is as much as if he had said: and here shall patience take place, whereby the saints may overcome all evils. Here we have need of stout courage, and a sure and constant mind. In Luke 12, the Lord likewise requires patience in persecutions. Here therefore is counsel given, how the saints should behave themselves, namely that they should suffer patiently those evils that Antichrist shall work against them. And there follows another sentence, which clarifies this: here are they that keep the commandments of God, and the faith of Jesus. They shall overcome through patience in so great evils and dangers, which keep the commandments of God, the foundation of which is the faith of Jesus Christ; which truly put all their trust in Christ, hear the word of the gospel, and keep the commandments of God, not of men. The like unto these are read in Matthew 24 and in Hebrews 10. Arethas says in this same time of Antichrist, the patience of the saints is shown. Then the speech is figured, as it were by a question moved. And who are they whom he calls patient? After, as though he should answer: they that keep the commandments of God, and the faith of Jesus. For they when perils approach will value God more than death and temporal evils. This he says. I pray God these things be as faithfully performed by us, as they are easily understood. The Lord grant us his spirit.

## Sermon 65: ¶ The faithful assuredly and straight way flitte from the corporal death unto life everlasting.

### Paragraph 982

.

### Paragraph 983

And I heard a voice from heaven, saying to me: Write, “Blessed are the dead who hereafter die in the Lord.” Indeed, the spirit says that they rest from their labors. But their works follow them.

### Paragraph 984

Although he has spoken on many occasions about the state of souls in the world beyond,[note 65.1] and about the blessedness of the faithful who are killed for the sake of religion, it was chiefly necessary to address this matter here. For I asked, how many must be killed by the beast. Now, lest they, fearing death, should choose rather to worship the beast than to be slain, and lest, perhaps, having lost this life, there be no other life to be hoped for in the world to come, he diligently and certainly treats of the state of souls, and of the felicity and blessedness of souls, which as soon as they die, they attain, assuredly and immediately departing from this world into everlasting life. But those who know these things, and have conceived them through true faith, and understand that they will undoubtedly pass from bodily death into blessed life, cannot but more boldly despise this present life.

### Paragraph 985

And this wholesome doctrine is comprised in three points. First, he shows its certainty; secondly, he declares what it is; last, he sets forth and illuminates it by circumstances. At the first, verily, he seems to allude to the manner, customarily received by all nations, that such things as they would have thought to be certain and undoubted, they would also commit to writing to leave them unto posterity. But the certainty and verity or authority of the thing is esteemed by authors, who first have dispatched any matters among themselves, and after have caused the same to be put in writing. At this present, therefore, is God shown to be author. For John says: and I heard a voice from heaven. And by and by adds: the spirit says. Therefore, there is no doubt but that the Son of God himself has spoken and revealed these things.

For him he saw at the beginning of this revelation; after he sees diverse kinds of angels, but he sees not Christ speaking to him. But he hears now his voice from heaven; he hears the spirit speaking, by whom the Lord said, whilst he was yet conversant in earth with his disciples, that he would treat and speak all things in the church. Let us believe therefore that the words which are here recited, by Christ’s doing, are a celestial oracle, certain and true, whereof we ought not to doubt. And John the Apostle and Evangelist is commanded to write the sayings of Christ from the heavenly seat. Which thing he doth; and so at Christ’s commandment sends them unto all posterity, unto us also and to our offspring even to the world’s end. But if writings written by the chancellors or secretaries of kings and princes, being notable men, deserve credit, we may much more justly and rightly believe this writing, which the Son of God indites from heaven; and that beloved disciple of Christ, the apostle and Evangelist John writes. You had once a confidence in the Popes’ bulls (they may well be called bulls, since they are more vain than bulls or blabbers in the water), sent from the See of Rome, wherein you, as one assured, did put full trust to have remission of sins and blessed life. And shall you not now be accounted mad and out of your wit, in case you will not believe this heavenly writing? That other was indicted by the spirit of Antichrist, by the Pope, the man of sin, and child of perdition; and written of some deceiver infected with Simony and sacrilege, which in life and manners was filthiness itself. But in John is nothing but cleanness, purity, and integrity; and the very Son of God which prescribes these things to John, is the very verity and life, the light of the world and lord of heaven and earth, of life and death. See then how safely you may lay yourself to this heavenly writing, which here is offered and given to you freely. You need not to disburse for the same one farthing. The Pope instituted in the church buying and selling and devilish bargaining about indulgences and other things, which were plain deceipts and illusions, plain mockeries, and open blasphemies, and therefore accursed forever; as Peter also pronounces in the eighth chapter of the Acts. God himself dissuades all men from such tromperies, and bargains wicked and vain, in the fifty-fifth chapter of Isaiah, where he promises again that he will give to the godly all plenty of all good things.

### Paragraph 986

And now let us hear what the writing is, and what Saint John is commanded from heaven to put in writing. It is a short sentence, as also in many places, the wisdom of God comprehends in few words the true sum of blessedness: so providing for our infirmity, that we need not to complain that the doctrine were too long, which we with our slender understanding are not able to attain to. The Lord therefore pronounces them to be blessed who die in the Lord; then we must see what he understands by blessedness, and who they are that die in the Lord. Blessedness is that high felicity which happens to the faithful in another world, in which we shall see God himself as he is, and have the fruition of him unto a joyful and never-wearisome fullness. We shall live in the same with all the Saints forever, and shall have joys that cannot be expressed with the tongues of men. Of which more shall follow afterward. They shall rest from their labors, and more plentifully in the twenty-first chapter. And they die in the Lord who, by faith grafted in Christ, are joined to him alone, depend wholly upon him, only regard him, and desire nothing else but him alone. For they are said to live in Christ, in whom Christ lives by faith; they that live in Christ do frame their whole life after the will of Christ. And they die in the Lord chiefly and before all, who for the confession of the Lord’s faith, suffer death, and offer themselves to torments. And not they alone, but those also, who although they die by the sword of the persecutors, yet die when the Lord calls them in the true Christian faith. For these are also blessed, as the Lord in Saint John says, “Verily, verily, I say unto you, if any man keep my word, he shall not see death forever.” However, they do not die in the Lord who either deny God that they might not be slain, or trust to their own merits and intercessions of Saints, or to other men’s works, be they monks, friars, or mass-serving priests, and so depart out of this life, thinking that they shall be helped by other men’s works. To be brief, the truth of the Lord pronounces them all blessed and fortunate who depart out of this world in true faith.

### Paragraph 987

Finally, the Lord himself adds a notable declaration of this brief statement. For he sets forth the circumstances of the time and the manner of the blessedness. For it is wont to be demanded, what time salvation and felicity happens to the dead? whether immediately, or after a time? That is, do our souls depart immediately after the death of the body to the blessed places? Or are they intercepted for a certain time, so that they might be purged in purgatory before they enter into heaven? Or are they held in sleep, awaiting the resurrection of the bodies, so that they might then awaken and together with their bodies enter into heaven? To all of which things the heavenly oracle answers immediately, saying that the same felicity comes to souls immediately. In the Latin copies, this place is pointed thus: “Blessed are the dead who die in the Lord.” Immediately now says the Spirit, that they may rest from their labors. In like manner reads the Spanish or Complutensian copy. But Arethas and the Greek copies, and also the exemplar of Paris are pointed so that [illegible] should be the end of the sentence, as Erasmus notes. After follows [illegible], which is, “verily, certainly,” says the Spirit. The sense is therefore, that the faithful, being dead, shall straightway and immediately achieve salvation. For the word which Saint John uses signifies from the very instant, from that hour immediately, incontinently. This allows no space, but expresses that which we are wont to note by the Dutch phrase. Being admonished therefore by a divine oracle and confirmed by a writ brought from heaven, let us all be assured that the souls of all the faithful depart from bodily death into everlasting life. These things are confirmed and made plain also by other passages of Scripture innumerable. I will choose out only a certain few, and those also the testimonies of our Savior, who is the light of the world and the word of life. In chapter 3 of Saint John, he says expressly that the faithful are delivered from death by his cross, as in times past, by the sight of the brazen serpent, the Israelites were delivered from the deadly sting of venomous poison. And it is plain that they were delivered immediately and most fully. In John 5, the same says, “I have passed from death to life.” Let this passage be weighed diligently, and it shall appear that it alone suffices in this matter. In John 6, he says openly, “I will raise him on the last day.” But he raises not the bodies only at the last judgment, but in every man’s last day, that is, in the death of every one, he preserves the souls, that they should not perish or be tormented, and so forth. We have in the gospel examples most clear: namely, of Lazarus the beggar, who was immediately after his death carried up by angels into the bosom of Abraham; and of the thief, who heard the Lord say, “Today you will be with me in Paradise”; and of Stephen saying, “Lord Jesus, receive my spirit”; but especially of our Savior, saying on the cross, “Father, into your hands I commend my spirit,” and so forth.

### Paragraph 988

By these things, whatever the monastic and Antichristian doctrine has built—regarding purgatory, indulgences, and the miserable state of souls in the world to come—is completely overturned. From this, they made a most shameful profit. Those who believe that souls are mortal are also refuted, as well as those who believe that souls sleep in another world, where they cannot even sleep here in this state of weakness. Therefore, you will say it is madness to think that souls sleep, being freed from the burden of the body.

### Paragraph 989

Concerning the manner of the blessedness of the saints, they rest from their labors. [note 65.7] Therefore, salvation is a most joyful tranquility. All diseases, sicknesses, griefs, anxieties, sorrow, famine, thirst, cold—briefly, all things that vex or trouble people—are removed. Rest and tranquility, joy and blessing take their place. And since the dead rest from their labors, who can believe that they are afflicted with torments? But lest any man should doubt even a little, he adds a confirmation, [illegible], yes, or certainly, verily, the spirit says, the dead shall be quiet from all their griefs. Let no man therefore doubt.

### Paragraph 990

And he adds another thing, that the works of the saints follow, that is to say, after the saints have departed from this world, they are rewarded in another world, if they have done anything well, if they have suffered hard things. For a reward is prepared for virtues. The saints hope for and receive this reward without boasting of their own merit, and not in contempt of the merit of Christ. For they acknowledge that God in his saints crowns his own gifts. And this is spoken of the reward of works for the consolation of those who suffer many things in this world. So the Lord said in the Gospel: “Your reward is plentiful in heaven.” And the Apostle affirms everywhere that rewards are prepared for those who are crucified here with Christ. And here let us mark diligently that these things are spoken also of the spirit of Christ under the guise of another. For the world despises religious persons and those who suffer for religion, and objects that they lose their labor and cost. Contrarywise, the spirit by another [?vouchsafes] avouches that a reward is prepared for virtue.

### Paragraph 991

Let us also note this: their works, and not those of other men, are what we should follow, and they are not sent after by others. Let no one deceive himself, let no one think that after his death, priests will send him into purgatory a portion of other men’s merits. Those are not good works which are done by priests and friars apart from and against God’s word, but provocations of God’s wrath. Those are not excluded from the kingdom of God in the Gospel who go to others to buy oil. Scripture says elsewhere: Let us do good while we have time, for a time will come when no one can work. Let us therefore watch, and in deed do good works through faith.

## Sermon 66: ¶ The Judgment of the Lord is described under the paraboles of harvest and vintage.

### Paragraph 992

.

### Paragraph 993

And I looked, and behold, a white cloud, and upon the cloud one sitting, like unto the Son of man, having on his head a golden crown, and in his hand a sharp sickle. And another angel came out of the Temple, crying with a loud voice to him that sat on the cloud. Thrust in the sickle and reap, for the corn of the Earth is ripe. And he that sat on the cloud thrust in his sickle on the Earth, and the Earth was reaped. And another angel came out of the Temple, which is in Heaven, having also a sharp sickle. And another angel came out from the Altar, which had power over fire, and cried with a loud voice unto him that had the sharp sickle, and said: Thrust in thy sharp sickle, and gather the clusters of the Earth, for her grapes are ripe. And the angel thrust in his sickle on the earth, and cut down the grapes of the vineyard of the Earth, and cast them into the great wine vat of the wrath of God: and the wine vat was trodden without the city. And the blood came out of the vat, even unto the horse bridles, by the space of a thousand and six hundred furlongs.

### Paragraph 994

Now he proceeds to describe God’s judgment, especially against the Antichristians and all the ungodly. This section could be joined with the following material, as it deals with the same argument. It pertains to the consolation and confirmation of the faithful, persecuted by Antichrist. Some people think there will never be any judgment. However, because they oppress their neighbors, they believe they will never feel God’s displeasure. Moreover, the faithful are tempted when they see the wicked prosper and themselves decline daily. Therefore, they think the Lord is taking too long. They expostulate with the Lord, saying, “When will the injustices end?” If Christ will come in judgment, why does he delay it, and cause such great suffering to his people? The Lord now shows that the judgment will certainly come, and will occur when all things are ripe—that is, when the iniquities of the Amorites are complete and the measure of iniquity is full. When wicked people are ripe, the Lord will come to judge. In the meantime, we must persevere in constancy and patience, like farmers waiting for the harvest and the grape harvest. If anyone rebels through impatience, he will not be approved by the Lord, as the Apostle states from the prophet in Hebrews 10. And just as we may desire and long for the harvest and the grape harvest, so we ought not to expostulate with God because he tarries longer than we wish; likewise, we ought not to contend with him as to why he comes so late in judgment. And just as the harvest and the grape harvest are certainly expected and arrive, so without a doubt God will punish the wicked and save the godly. These are truly tastes of what is to follow abundantly and are more expressly declared, and are appended to the former matters, because they pertain to the consolation of the godly.

### Paragraph 995

And to the intent that all things might be more clear, using parables to present them before our eyes. He employs two parables drawn from the prophets and the doctrine of the Gospel. For the prophets frequently depict God’s judgment through harvest and vintage. Certainly in Joel 3, the Lord says: “I will sit in the valley of Josaphat to judge all nations. Put in the sickle, for the harvest is ripe,” and so forth. It is also well known what is said on the same matter in the story of the Gospel. We must therefore write these things more deeply into our hearts, and fear God, and persevere in his redemption with patience.

### Paragraph 996

First is treated the parable of harvest, then the parable of vintage:[note 66.2] both showing that the Lord will be judge and that in his appointed time, against all those who either think there shall be no judgment, or question the Lord’s timing, saying he comes too slowly and late, &c. And first is described the owner of the harvest, the Lord himself and judge Jesus Christ. He is said to be like unto the son of man: not because he is not now the very son of man, nor because he shall not come to judgment in the very human nature which he once took from us and never laid it aside (for he is truly the son of man,[note 66.3] and remains on the right hand of the Father, and shall truly come in human nature to judge the living and the dead), but he seems to have alluded to Daniel, and to have expressed his speech, saying: “I looked in a mighty vision, and lo, there came one in the clouds as it were the son of man,” &c. Where we read also the description of the judgment against the beast. And therefore he has here made mention also of a cloud: “and I saw a white cloud, and one sitting on the cloud,” &c. Moreover, the angels in the Acts say, “so shall he come, as they have seen him go up into heaven.” And they saw him taken up, and a cloud to receive him and convey him out of their sight. Therefore shall he come again in a cloud unto judgment. We read oftentimes in the Psalms that God sits on a white cloud. By the way, therefore, is signified the deity of the judge. Therefore this judge is very God and very man, the Savior of the faithful, the revenger and judge of the unbelievers. We are sent therefore by Jerome to the seventh chapter of Daniel.

### Paragraph 997

Then he wears a golden crown on his head: not that there is any corruptible gold in Heaven, but for corruptible men so he speaks, that they may understand their judge to be the high king; and may gather thereof, that no man is able to resist the power of this king. For otherwise our Lord hath no need of any corruptible gold.

Finally, our Lord here has a sickle, and that right sharp. Whereby is signified his judgment exceeding straight, and destruction of the wicked. In the third chapter of Matthew, the judgment of the Lord is compared to a fan, by blessed John. He adds that the axe is laid at the root of the tree: whereby he signified that certain judgment was at hand, or rather destruction.

### Paragraph 998

Now follows an exposition of the proceeding of the judgment, and he perseveres in the parable. For he speaks as if a servant returning home out of the fields, did show unto his Master who looked for the hour of harvest, that the corn was now ripe (the hardness of the grain is a token of ripeness) and that it is time to be reaped. For else it is no need to admonish him that knows all things of any thing, that he remembers not: much less of the hour of judgment which no angel knows, but the Father alone. Therefore we ascribe this wholly to the parable: and we understand that a certain hour of judgment is appointed, which when it shall come, the godly will be delivered without delay, and the ungodly condemned. Another angel, says he, came forth. For before we heard how diverse came forth. This cries with a loud voice, as one that will tell of a matter most great and certain, and to be declared in the church with exceeding great outcry, to the comfort of the faithful, who ought nothing to doubt of the judgment, and to the terror of the wicked, who seem to contain the same. And this crying angel comes out of the temple. For we heard before that John saw a temple in heaven. And where the crier of the judgment comes out of the temple, it signifies that no unrighteousness of the judge is to be imagined. For the temple is consecrated to holiness and righteousness, and is called the house of God. Justly therefore he judges, and in just time he judges, and justly executes all things. The angel bids the judge do that which he of himself was about to do. Thrust in the sickle, says he, and reap. Two causes are alleged. First, for the hour is come that you should reap. Therefore a certain hour of judgment is appointed, which when it comes, the judgment shall be most certainly. And it is come for you, says he, for all judgment is given to the Son. Then, for the corn of the earth is ripe. As though he should say: the iniquity of earthly men is grown up to the highest, therefore it is reasonable that it should be cut down. And God alone knows when the iniquity of the earth is fulfilled. But when it shall come thereunto, there shall need no great preparation, deciding, or pondering of causes. At one word he finishes the judgment, and the execution of the same, and as it were swallows up and devours the whole earth in a moment, saying: herewith he thrust in his sickle, which sat upon the cloud, on the earth, and the earth was reaped. The rest of the things which seem to belong hereunto, take out of chapter 13 of Matthew. And that which he has said hitherto, he repeats, and beats in by another parable. For by this he shadows the same which the other parable did commend. That plenty makes for the plainer evidence, and beats in most diligently the certainty and verity of the judgment, lest herein we should doubt anything, and waver with the unfaithful world. The parable is taken from vintage. The same is used very often by the prophets, speaking of the destruction of any nation. And the Lord also in the gospel compares his people to a vine. And the angel holds in his hand a sharp sickle. He represents a figure of Christ, who has all power of judgment alone. A sharp sickle is the straight judgment, as was spoken of the sickle before. This angel comes out of the temple also, to wit a judge most righteous. Unto him cries another angel, which had power over fire, which comes out from the altar. For before we heard that there is an altar in the temple, and that under this altar do rest the souls of the blessed martyrs. Here therefore is figured that God does now remember the bloodshed of his servants, who for the profession of the only altar (that is Christ the priest and only sacrifice) were slain, and now to proceed to take vengeance hitherto long delayed. Therefore this angel is said to have power over fire. Fire many times in the Psalms signifies God’s vengeance. This angel therefore is here, as it were master of execution, and captain of vengeance. For angels in Daniel also, as God’s ministers, are said to have rule over things: not that we should worship and honor these ministers, but the Lord that works by them. The sun and moon are the lights of the world: but therefore no wise man will worship them. Here is signified plainly that vengeance is certainly prepared for them which shed innocent blood on the earth, and that this vengeance shall chiefly be executed in the end of this world. Albeit that he punishes nevertheless grievously before the end also here on earth, namely parricides: in so much that the Psalmograph says, men of blood shall not live half their time.

### Paragraph 999

And as in the parable of harvest, harvest was finished with a short sentence: So here also, the vintage is ended with few words. For as soon as the ungodly shall see Christ in the clouds, with the prints of his wounds, and his Saints with him, whom they have scorned, hated, persecuted, and slain, they will immediately gather, that they by their just desert must be allotted with devils, whom they have followed and served. Therefore, there shall need no long discussion of the matter. Every man’s conscience shall accuse him, and the sins of every man shall be manifest to all creatures. The ungodly shall stand before the judge with great confusion, in utter contempt, in pain and fear, and sorrows not to be expressed, and shall go straight ways into pains and torments that shall never have end. [note 66.10] Hereof I say, it behooves often to make mention, hereof it becomes often to warn all men, that they may beware in time, and take heed to themselves.

### Paragraph 1000

However, Saint John himself describes the everlasting damnation and vengeance that God executes upon his enemies. [note 66.11] And he uses the image of a wine press or wine vat, so that he may linger in the allegory, and it is made outside the city. By explanation, he calls it the great wine vat of God’s wrath. For it is hell, or the place of punishment and condemnation. Into this wine vat shall be gathered the clusters of the earth, or the grapes of the earth, that is, the earthly and ungodly men. And the city of God is heaven itself, the seat of the blessed, which shall afterward be described most abundantly in chapter 21. But that wine press is set outside the city. For in another place in the Gospel, the Lord also says that the wicked must be cast out into utter darkness, where there is weeping and gnashing of teeth.

### Paragraph 1001

But this wine vat is rightly called the wine vat of God’s wrath. For the wrath of God is executed therein: and those with whom God is angry for their sins are shut up therein, that there they may, according to their demerits, be tormented and vexed forever, and without end. And he calls it great, for that the place is wide enough to receive all the ungodly. As also Isaiah has admonished at the end of the thirtieth chapter. Others read of the great wrath of God.

### Paragraph 1002

It is added that from the fat or wine press no wine flows, but blood, and that in great abundance. This he illustrates with a marvelous and horrible hyperbole. The blood flowed far and wide, for the distance of a thousand six hundred furlongs. Again, it was very deep; for it rose to the bridles of the horses, those I mean, which went and wrestled in the blood, namely, in their own blood. By this hyperbolic speech is signified that the multitude of the ungodly will be greatest, and that God will most abundantly avenge that immeasurable blood which the wicked have spilled on earth. They delighted themselves while living on earth with wars, slaughter, persecutions, and martyrdoms; therefore, God, most just, will give them in another world enough blood, so that being drowned in their own blood up to their chins, they may seem to bathe themselves in their own blood. And here we must remember that the horses prepared for battle of which we spoke in chapter 9 will be drowned in everlasting torments. Thus, thus at the last will the Lord avenge himself upon his enemies. Let us call upon him, and abide patiently and valiantly. The Lord grant us his grace.

## Sermon 67: ¶ The Angels of seven plagues are brought forth. Moreover the triumph and praise of Christ's holy Maritrs is described.

### Paragraph 1003

.

### Paragraph 1004

And I saw another sign in heaven, great and wonderful. Seven angels, having the seven final plagues. For in them is fulfilled the wrath of God. And I saw as it were a glassy sea mingled with fire, and those who had conquered the beast, and his image, and his mark, and the number of his name, standing on the glassy sea, having the harps of God: and they sang the song of Moses, the servant of God, and the song of the Lamb, saying: Great and marvelous are your works, Lord God Almighty, just and true are your ways, O King of Saints. Who shall not fear you and glorify your name? For you alone are holy; for all nations shall come and worship before you, for your judgments are made manifest.

### Paragraph 1005

[note 67.2] Concerning the harvest and vintage, explained in the last part of the previous vision, is appended the fifth part of this godly work, which presents to us the fourth vision of this work, which some consider to be the fifth. This concerns God’s judgments and has two parts; therefore, it might also be divided into more visions, but we prefer to use fewer. First, he discusses at length the pains or torments prepared by God and to be executed upon Antichrist’s followers, and all the ungodly. Here is treated the judgment of the whore of Babylon, the destinies and ruin of Rome and the Roman church, the rejoicing and song of the saints, the coming of the judge for judgment, and the pain and everlasting destruction of all the wicked. These matters are addressed in chapters 15, 16, 17, 18, 19, and 20. Then also he reasons excellently of the reward of the saints and of everlasting blessedness throughout chapter 21. Everywhere, hell itself and heaven itself are set open; and we are given a glimpse, while in this mortal flesh, even into hell itself and into the very palace of Heaven. You will not find anywhere in all the Scriptures a continuous and abundant discussion of the judgments of God, the torments of the wicked, and the felicity and joys of the godly, as in this present work.

### Paragraph 1006

And this treatise is especially necessary in this our last and ungracious world,[note 67.3] wherein men, neglecting the spirit of God, have become like brute beasts, altogether carnal, regarding the flesh and wholly depending on it. Happy are all the victorious, wealthy, honorable, and glorious Antichristians; miserable are the poor and despised true Christians, and subject to the injuries and persecutions of all men. Therefore, carnal men esteem all things of the present fortune, and proclaim that their religion and conversation pleases God, and that the Christian way displeases. The godly are here also grievously tempted, as they were also in times past, read Psalm 73 and the first chapter of Habakkuk. The ungodly promise themselves that they shall reign forever; at length also they contain the judgments of God, neither thinking that it will ever come to pass that they shall be punished. The talk of punishments is devised by melancholic persons, and uttered of malice; and therefore they say and think them not to be regarded, but to be merry in this world. Therefore it behooved the place of God’s judgments to be most largely and diligently decided, and to be set as it were before the eyes of the hearers: to the end all might rightly understand what should be assuredly the end of good and evil. But the punishments of the ungodly are diverse, to wit, of this life present, and to come. And the punishments of this present life are almost innumerable; and the torments of the life to come are eternal and unspeakable; and as there is no comparison between the pained and true fire, so is there none between the punishments of this present life and that to come. But if men would earnestly believe that unspeakable joys and everlasting torments are prepared by God for good and evil, doubtless all would sin less and serve God more diligently. But let us now see what is the treatise of Saint John concerning the same.

### Paragraph 1007

First, he demonstrates that all that follows is not earthly, but heavenly. For he sees another token in heaven. He says “another,” because in chapter 12 we heard that mention was made of another certain sign. And he calls that a sign or token, which signifies another thing, and therefore is not to be considered of itself, but in as much as it brings into knowledge another certain thing, and that much greater than it shows at the first sight. He calls this sign, that is to wit, that same vision, great and marvelous. For the judgments of God are greatest and most wonderful. While they are executed, the ungodly marvel, which had thought such things should never have come to pass; the godly also marvel at the great power of God, his most just righteousness, and his ripeness and faithfulness in delivering and saving his people. Then he declares what sign was shown him in heaven, and by that celestial vision: he saw seven Angels, having in seven cups, plagues. That is, he perceived God prepared and furnished with divine power, with which he both might and would send plagues and deserved punishments, as well upon Antichrist himself, as upon his members, and all the ungodly men on Earth, for their wickedness committed against God. And as we have many times warned you in this book, the seventh number is the measure of fullness. Wherefore God has ministers enough and enough, by whose service he may plague and destroy the ungodly. And therefore seven plagues are all manner of plagues. Temporal plagues are abundantly recited in Leviticus 26 and Deuteronomy 28. The Lord is rich, and in everlasting plagues of most diverse kinds also. For Scripture in certain places rehearses a gnawing worm, a fire unquenchable, weeping and gnashing of teeth, outward darkness, and many others of like sort. But these seven plagues he calls the last, and immediately shows the reason, for in them is the wrath of God fulfilled. For on those last and most corrupt ages the Lord will pour out his plague, and that the plagues of his just wrath, and shall pour them out most fully to the end, and shall execute his full wrath against the ungodly, forevermore.

### Paragraph 1008

Yet now he suspends a while that narration concerning the Anglo-Saxon rulers of the plagues [note 67.6] and places or sends before the great joys of the blessed Martyrs, triumphs, songs of praise, rejoicing, and thanksgiving. And this joy is interlaced here in the treatise of punishments, for the consolation of the faithful, that they should know themselves delivered from punishments. And if it happens, while the wicked are punished, that some displeasure touches them also (as it cannot be avoided, but that when the wicked are plagued, some detriment must also arise for the faithful), that they may understand yet that the dangers and detriments must be recompensed with the excellent abundance of joys. For hereby is signified how the godly rejoice while the Lord executes his justice. To be also the changeable course of things, that those who have once wept in the world, should now be glad and joyful, according to the saying of our Savior in John 16. Moreover, it behooved by the testimony of all Saints to be declared to the Saints who dwell on Earth, that the judgments of God are righteous and true; which thing, understood, questions and sundry murmurings against God cease.

### Paragraph 1009

First he sees those who overcame Antichrist, and have had nothing to do with him; as we say in Dutch, for this I suppose is signified by that plentiful rehearsal of certain members (the declaration of which is set forth before) in heaven, not in some torturous place, or nowhere, as some men gather. He saw, I say, in heaven the blessed souls standing on a glassy sea [note 67.7], mixed with fire. And in another place I have told you that the sea figures the world, because of its rage and instability. Certainly Daniel takes it so in chapter 7. And it is called glassy because of its frailty and brittleness. For worldly things shine, but they are soon broken. Whereupon it is said that worldly things are as brittle as glass: which, while they shine, break. And not without cause is fire mixed with worldly things. For the saints, while they are conversant in earth, always feel in a manner the fire of affliction, as spoke of by Saint Peter, 1 Peter 4. And they stand on a glassy sea mingled with fire. For conquerors tread upon the world, and upon all the torments and mockeries of the world, triumphing over all worldly things. The prophet in Psalm 66 brings in the saints singing a joyful song unto God, and among other things saying, “You have brought us into snares, you have laid tribulations upon our backs, you have set men in our necks. We have passed through fire and water, and you have brought us out into a place of relief.” Therefore, alterations follow in another world. Wherefore Arethas, expounding this place, says that the glassy sea seems to indicate nothing else than that by the sea is meant the multitude, by the glass the brightness, by fire the purity of those who are worthy of that blessed life. And certainly the same words in diverse respects may signify diverse things, and make the sense agreeable.

### Paragraph 1010

Hitherto we have heard that the saints are in heaven, where they triumph over the world, which is vanquished. But now we shall hear more clearly what they do in heaven, and how they sing to the Lord a song of thanks and praises, which fully agrees with Psalm 66 [note 67.8]. And he attributes harps to the blessed martyrs, as he did to the Elders. These he calls of God, as you would say divine and celestial, fitting to set forth the praises of God. For a celestial jubilee is signified, of which it is spoken in the fifth chapter. He adds moreover, to express the music: and they sing. And he declares also the manner of their singing [note 67.9], the song of Moses, the servant of God, and the song of the Lamb. Therefore, this song of the saints is a rejoicing ditty, triumphant and of thanksgiving. For just as in times past Mary with the company of Israelite virgins, at the appointment of Moses, sang a song when the Israelites were delivered out of the bondage of Egypt, and Pharaoh was drowned in the Red Sea with his whole army—of which you may read more in the fourteenth and fifteenth chapters of Exodus—so the blessed souls in heaven praise God, who has delivered them from Satan, Antichrist, and the world. And the song of the Lamb is the Christian thanksgiving, by which the virtue of Christ and his redemption is praised by the saints. For just as the old fathers after eating the Paschal lamb made a jubilee and gave God thanks, so the blessed saints now enfranchised with the full liberty of the children of God, give thanks unto Christ their deliverer.

### Paragraph 1011

Finally, they recite the order and form of their song. [note 67.10] God is highly commended here, called the Lord, God, almighty, King of Saints—the one for whom the saints fight, by whom they are governed, and who defends, maintains, and keeps the saints. He is called holy, in whom there is no blemish, no iniquity. And above all, they praise his works, which they call great and marvelous. These are manifest in heaven and on earth. They declare the power, wisdom, and justice of God. Therefore, they immediately conclude that God’s ways—that is, the considerations he follows in governing and doing things—are true and just. For he deceives not, he does no man wrong. Therefore, God is just in punishing the Antichrist and delivering his own. For although he may seem to neglect them, he keeps faith with the godly, just as a king never neglects his subjects.

### Paragraph 1012

Now they claim that it is fitting for all people on Earth to do likewise: it is right that all men fear you and glorify you in all things, neither to accuse nor murmur at your judgments. There is added another reason: he alone is holy, without sin, and without spot. None of all creation possesses this. Although many Gentiles now despise God, yet they will come and worship; they will know their own filthiness and the holiness and righteousness of God. For the justice and judgment of God, which are not yet revealed and therefore are despised, shall be revealed, that all the godly of all nations may attribute glory to the righteous God. These things also prepare the reader and hearer for the treatise now following concerning the judgments of God and the punishments of the ungodly. The Lord open the eyes of our minds, that we may see these things with fruitful understanding.

## Sermon 68: ¶ The seven Angles are described, coming forth to execute the seven plagues.

### Paragraph 1013

.

### Paragraph 1014

And after that I looked, and behold, the Temple of the Tabernacle of testimony was open in Heaven, and the seven Angels came out of the Temple, which had the seven plagues, clothed in pure and bright linen, and having their breasts girded with golden girdles. And one of the four beasts gave unto the seven Angels seven golden vials full of the wrath of God, who lives forevermore. And the temple was full of smoke, from the glory of God, and from his power: and no man was able to enter into the Temple, until the seven plagues of the seven Angels were fulfilled.

### Paragraph 1015

He now returns to the description of God’s judgments, from which he had made a brief digression. This treatise is very fruitful. For God’s judgments are the punishments or pains of the wicked, and testimonies of God’s righteousness and truth. Again, the godly are strengthened in their hope. For they see that not one jot falls from the words and threatenings of God, even though he may be long-suffering, tolerates them for a time, and even seems to favor and spare the ungodly. Therefore, the godly perceive that their hope is not vain. They also learn to fear God and to pray continually, lest, being intoxicated by the pleasures and felicities of this world, they turn away from God to ungodliness. Finally, the wicked are terrified by these pains, are provoked to repentance, which, if they refuse, they will undoubtedly feel plagues, as Pharaoh did.

### Paragraph 1016

But before the Anglo-Saxons’ power brings forth the cups of plagues received, they are most fiercely and diligently described. And it is shown from whence they came, that is, what is the origin of God’s judgments. They come out of the temple set open, and that out of the temple of the Tabernacle of Witness, which is in heaven. For Moses saw a temple on the mount, and that also in heaven, after the similitude whereof he was commanded of God to make the tabernacle of Witness. Therefore, the tabernacle of Witness was fashioned and built after the shape exhibited and seen in heaven, which the blessed Apostle to the Hebrews calls the pattern, that is, the very example or patron. For it was said to Moses, “See that you make everything according to the pattern which was shown you on the Mount.” Which thing Moses did accordingly. But such things as came forth of the Tabernacle of Witness made on earth, seemed to the Israelites just and holy. Hereof were delivered the oracles and answers of God, which it was not lawful to speak against. Therefore, when we hear now that the very judgments of God against the wicked world, pains and punishments, come out of the true temple itself, the pattern I mean and that celestial, who should hereafter doubt that all the judgments of God, with which he plagues the ungodly, be sacred and holy? And whilst the ungodly are plagued, that we must think nothing else, but that a sentence, as it were an oracle, is come or pronounced from heaven, which it is unlawful to gainsay? To conclude, the divine judgments do proceed out of the throne of God, wherefore they can not but be most holy. Otherwise, we shall hear in the twenty-first chapter that there is no temple in heaven. These are therefore types and figures, not matters true and permanent; but after they have signified this, for which they were instituted, they are passing and fading away.

### Paragraph 1017

Hereunto also appertains the apparel of Angels,[note 68.4] that hereof men may also esteem the judgments of God. They are said to be clothed in pure linen, or clean and white, or bright, whereby is signified that the judgments of God are unspotted and bright. For we have heard that these things which John saw were signs. Therefore we may not imagine carnal things in heavenly matters, but spiritually to expound such things as in the sign seem to be as it were corporal. The garment in this world is changed with the state of things. For they use white garments in victories and triumphs, black at burials and mournings, red in battle. Here is signified therefore that the judgments of God are most pure, and that God overcomes and triumphs over the ungodly. At the resurrection and ascension of our Lord Angels appeared in white garments, and shining bright, to signify the glory of Christ. Now is the very breast girded with a girdle, and that in deed with a golden girdle. Gold is a token of purity. In the breast is the seat of affections. The girdle binds, moreover prepares for the journey. Therefore it betokens that the judgments of God are prepared, and in a readiness: the same to restrain affections, that is to say, not to be pronounced or done of envy or malice, love or favour, but to be just, moderate and upright.

### Paragraph 1018

And one of the beasts gave to the seven angels, revengers and punishers, seven bowls, and these were full of God’s wrath. Now, although God does not need the help of creatures, nor receives anything from them as if lacking anything, yet he did not make his creatures in vain, and he acts in order. All creatures, without doubt (for I said in chapter four how the beasts signify the universality of creatures), contribute their labor against the wicked, and whatever they have from God (and they have all things) at his will and commandment, they willingly and freely employ to execute God’s judgments. Thus, fire falling from heaven upon Sodom and the cities around it, ministered the plague or cup of God’s wrath to the angel revenger. So the water overwhelmed Pharaoh and his host. So the earth opened and swallowed the company of Dathan and Abiram, and so forth. Thus, the armies of the Gentiles employ themselves to punish the ungodly. The walls of Jericho fall, the hail destroys the Canaanites. Thus God, without any difficulty, punishes his enemies, saying that all creatures are ready to aid and assist. And the bowl or cup is of gold, for again is signified the justice and equity of God’s judgments. And where God is called a ruler living forever, his eternity and majesty are signified, which neither the transitory things of this world nor human infirmities shall overcome. In the sight of the living God, all the wicked shall fall and perish eternally.

### Paragraph 1019

After this, the Apostle sees the temple filled with smoke for the majesty of God and for his power. That smoke is a sign of God’s presence, it appears by many places in Scripture, but chiefly in chapter 8 of the 3rd book of Kings. It is also a token of God’s wrath. For Aretha says, “smoke is a token of God’s wrath,” according as it is said, “smoke ascended in his wrath.” And neither is smoke without fire, nor fire without smoke. Moreover, smoke hurts the eyes and makes them blind. So in Isaiah 6, the temple of God, which Isaiah sees, is filled with smoke. And at this present, not only the presence of God and of his wrath appear to be signified, but also to be figured, that the judgments of God are unsearchable, so that the things which he himself reveals not to us, we cannot attain to. For his majesty is infinite, and his power passes all things. Primasius, Bishop of Utica in Africa, expounding this place, says, “Think that smoke signifies that all men cannot penetrate the secrets of God’s judgments, and that the eyes and minds of mortal men shall, at the contemplation of the plagues inflicted, grope in darkness.” Which now he determines to utter, and unto the final end of the same, he affirms the smoke to abide still in the temple. Thus says he.

### Paragraph 1020

It now appears to explain the same [note 68.7], and no man could enter into the Temple, and so forth. But it is certainly true, by the authority of the evangelical and apostolic doctrine, that the souls departing from the body before the end and final judgment go directly into the blessed places and there experience the joys promised by God, which is true. Therefore, another thing is signified, namely, that before the end of all things, the saints cannot clearly see all of God’s judgments. For here we see through a glass, but there face to face, and we shall know God himself, and the truth and manner of his judgments. Primasius, neither could any man enter into the temple: that is, could penetrate the secret, until the seven plagues of the seven angels were finished. Wherefore, the Psalmist: “This,” he says, “is labor before me, until I may enter into the sanctuary of God and understand the conclusion of matters, and so forth.” Here is signified therefore, that saints before the judgment will not know the secret mysteries of God’s judgments. Let it then suffice us that he himself has deigned to reveal to us; for the rest, let us believe that the Lord is just in all his ways and holy in all his works. To him be glory.

## Sermon 69: ¶ The three former Angles power out their vialles upon the Antichristians, and all the ungodly.

### Paragraph 1021

.

### Paragraph 1022

And I heard a great voice out of the temple saying to the angels: Go your ways, pour out your vials of wrath upon the Earth. And the first went, and poured out his vial on the earth, and there fell a noisome sore upon the men which had the mark of the beast, and upon them that worshipped his image. And the second angel poured out his vial on the sea, and it turned as it were into the blood of a dead man: and every living thing died in the sea. And the third angel poured out his vial upon the rivers and fountains of waters, and they turned to blood, and I heard an angel of the waters saying: Lord which art and wast, thou art righteous and holy, because thou hast given such judgements: for they shed the blood of saints, and prophets, and therefore thou hast given them the blood to drink: for they are worthy. And I heard another angel out of the altar saying: Even so, Lord God almighty, true and righteous are thy judgements.

### Paragraph 1023

After speaking generally about God’s righteous judgments, he now proceeds particularly through the number seven and declares at length the plagues of God, [note 69.2] which he also inflicts in this world upon the wicked, but especially upon those who oppose Christ. This passage corresponds to the same, or at least has many similarities to, the passages in Moses’s book of Exodus from chapter 7 through chapter 12. For in those chapters are described the ten plagues of God, with which, for sin, he afflicted King Pharaoh and the entire realm of Egypt. These plagues are summarized in fine verses by Dr. Musculus, our worthy Godfather.

### Paragraph 1024

These plagues are also explained in Psalm 150. In chapter 15 of Exodus, the Lord says: “If you diligently hear the voice of your God, and do what is right in his sight, and keep all his statutes, I will not send upon you any disease that I sent upon the Egyptians, for I am the Lord, healing you.” We learn, therefore, from the account of God’s plagues, to fear God and to walk in his commandments. It is not contrary to this saying of God that we read how Job and other holy men, walking in the commandments of God, were afflicted with grievous diseases. For these are private [illegible] are not chiefly inflicted for sin, but for the exercise of faith and the increase of virtues.

### Paragraph 1025

Men for the most part ascribe the causes of plagues to the stars and to other matters [note 69.4], and therefore do not turn to the Lord, who is striking them in response to a most wicked life. But we are taught by the treatise of Moses, which we cited from Exodus, and by this present discussion of John, that God himself punishes the sins and wickedness of men, although he uses the service of men and elements. To these, as to the proximate causes, men attribute the evils they receive, which they justly suffer from God for their sins [note 69.5]. For this cause, a voice is now heard, not from the air or from the earth, but from the Temple of the Lord, true, just, and holy, commanding the angels to pour out their vials upon the heads of men. The wicked are therefore plagued by God himself. But a vial is nothing other than the just judgment of God, or vengeance deserved by men. Angels pour out their vials whenever men are punished with plagues through means appointed by God. And that voice which is heard from the temple is great, for no man can resist God, nor infringe his decree. When he commands, all creatures obey.

### Paragraph 1026

But while this first angel, executor of God’s judgment, unleashes his plague upon men, a terrible sore botch fell upon them. This plague corresponds to the sixth plague of Egypt. And that botch signifies a canker, a fistula, and swelling sores or boils, but chiefly the pox of Jude, which others call the disease of Naples, some the French pox, and some the Spanish; truly, for that in the war of Naples (which was waged by the French and Spaniards in the year of our Lord 1494) it first appeared in the camp of prostitutes, infecting the army. Mainardus the physician discusses this at length. But however diverse and venomous sores do infect many grievously, yet the French pox chiefly corrupts the abbeys of monks and nuns, and colleges of priests, above others. For they giving themselves to most filthy fornication, do abhor and detest in others holy matrimony, and therefore receive thereof the reward of their iniquity. Therefore it is said here expressly that the Antichrist should be vexed with this disease, or rather punished. You shall find some whose face is eaten with this disease. All whoremongers and adulterers for the most part are troubled with this plague. Job also, the excellent servant of God, was covered with sores and boils, but by the singular counsel of God, as I touched upon before. Therefore it is no marvel, though sometimes very good men, free from the uncleanness of whoredom, are also infected with this disease.

### Paragraph 1027

The second angel poured out his vial on the sea, and immediately the sea became like the blood of a dead man, corrupt and lifeless. All that lived in the sea died. The sea is always in motion and changeable; therefore, it rightly signifies the world, or unconstant people in the world, who, because of their sins, are afflicted with the plague and die in great numbers. In the words there is the figure of synecdoche, where every living soul is said to die. This second plague corresponds to the fifth plague of Egypt. Under this plague we understand all kinds of pestilences and plagues. Hezekiah also suffered from the plague, as do many godly men, but by God’s unique counsel.

### Paragraph 1028

The third angel poured his vial upon the rivers and fountains of waters, which were immediately turned into blood. [Note 69.9] This corresponds to the first plague of Egypt. The Egyptians had drowned in the Nile the newborn infants and had oppressed the innocent Israelites; therefore, they were worthy to drink of the Nile, for it became blood.

### Paragraph 1029

Water otherwise in the Scripture signifies doctrine, as in Ezekiel and Zechariah. Therefore, the rivers and fountains of water signify ecclesiastical preachers and rulers, whom God has given to the people for defense and relief. Certainly, Saint Peter calls false prophets dry wells. In 2 Peter 2:17, we shall hear that by waters are understood people. This, then, is God’s plague: rulers of the people and preachers of peace have become the instigators and leaders of rebellion and war, in which they fall and kill one another, shedding the blood of the saints. And although in wars the godly are also afflicted, the Lord knows how to requite their suffering and ease their sorrows. Saint Augustine settles this matter at length in the first book of Christian doctrine. But if we look upon the tumultuous history of Italy, France, Germany, and Hungary, and other realms that glory in being called Christian, we shall find them to have been blazing firebrands of war, which by right ought to have been the rulers of peace.

As the Lord says in the Gospel, a prophet must not die anywhere except in Jerusalem, so no war must be instigated except by the Popes of Rome, bishops, and prelates. I will only recount a few examples. Pope Gregory II, through sedition, expelled Emperor Leo Isauricus from Italy. Pope Stephen brought King Philip of France into Italy against the Lombards. The same did Charlemagne at the urging of Pope Leo III, driving him clean out of Italy, having slain many of them with the sword. Pope Gregory VII, a most wicked man, stirred up Peter, King of Hungary, to war with Emperor Henry IV, entangled all of Italy and Germany with wars, and drove Henry to fight many battles, which were not light. Urban II, of that name, stirred up war in both the East and West and all other parts of the world he called holy, attempting to recover Jerusalem. This war was long, cruel, great, and bloody, such as cannot be found in all the world. What Alexander III accomplished against Frederick Barbarossa, and how he raised up all of Italy against him, the histories tell. And while Frederick II warred in the Holy Land, Gregory IX took Naples from him. Here, the Abbot of Ursperge observes that such wickedness should be committed by a Pope. Great factions arose in Italy, of the Guelphs and Ghibellines, through the means and instigations of the Popes. Clement IV brought in the French army, under the leadership of King Charles, into the kingdom of Naples, and deposed Conrad, Duke of Swabia, from his inheritance, and caused Frederick, Duke of Austria, and Conrad to be slain together with many thousands of Germans. Pope John XXII armed Frederick, Duke of Austria, and Leopold against Emperor Louis IV of the house of Bavaria. Boniface VIII commanded King Albert, Duke of Austria, to bear hostile banners against King Philip of France. As Martin V stirred up a grievous war against the Bohemians. Eugenius IV betrayed Ladislaus, King of Poland and Hungary, to Amurath, the great Turk, to be vanquished and slain through treason, sending his legate Julian Caesarini, a cardinal, about the practice, who also perished in that disastrous defeat. This recalls the saying in Virgil: “Even the soothsayer himself is slain.” Pope Sixtus IV sent to the powerful Swiss nation a red scarf with a bull attached, granting large indulgences to those who would fight for the church of Rome. Julius II, through much bloodshed of the Swiss, began to expel the French king from Italy, which at length Leo X accomplished, receiving Emperor Charles V, whose son still rules in Italy. Clement VII began to oppress him again, but death thwarted his endeavors. Paul III joined the forces of Italy with Charles V, and warred against the Germans for obedience denied to the See of Rome and the Gospel received. In this war, Philip, the Landgrave of Hesse, and John Frederick, Duke of Saxony, an Elector Prince, were taken. Great villainy and cruelty were wrought by the soldiers in Germany. Pope Julius III began to parley with the Frenchmen, and stirring up the war of Parma and Mirandula, brought the Frenchmen to Siena. There arose a most grievous war by sea and land, both in France and Italy, and also in Germany, which continues to this day: the princes and people tear one another asunder, they drink their blood abundantly, while nevertheless, in the meantime, they persecute Christ’s church most grievously. The Lord send peace.

### Paragraph 1030

And now where the godly might marvel why God so suffers the world to be shaken and troubled with mutual wars, the Angels prevent the marveling and complaint, and show not only the cause, but also praise the justice of God in these judgments. And he brings in two Angels, as meet and sufficient witnesses of this business. The one he makes ruler of waters, the other speaking out of the altar. He seems here to follow Daniel, who also in the tenth chapter, says that Angels as governors were set to rule over provinces. Not that God does not work and govern all things in waters and in all elements and regions; but for that he uses the travel of Angels as his ministers. But where the Papists gather hereof that Saints rule over elements, diseases, limbs, cities, and every part in man, it is foolish and superstitious, and smelling of idolatry. For the manner of Angels and of blessed souls is clean diverse, moreover, the Scripture attributeth unto them far other things than it does to these. You shall read nothing of the blessed souls as having anything to do with men here on Earth in the whole Scripture. But in sundry places of the Scriptures, you shall read that Angels are set to be men's keepers, and to serve them with divers ministries. Again, you read not that the godly have for this cause given any godly honor to the Angels; no, we shall hear in this book how Saint John would have worshipped an Angel, but was prohibited of the Angel once or twice. Moreover, here the Angel renders a reason why the water is turned into blood, and commends here in God’s justice. For turning his talk unto God, “Thou art indeed, O Lord which art, and which wast,” and so forth. He pronounces him righteous, as he that will do no man any wrong, and therefore calls him also holy. In the meantime he signifies his everlastingness, and that he gives being unto all things, where he says, “which art, and which wast,” and so forth. Of this phrase of speech is spoken in the first chapter. And the true righteousness gives to every one his. Therefore the Angel says, “Therefore Lord thou art righteous, and declarest thy righteousness to the world, in that thou hast given them blood to drink, which have shed the blood of the Prophets,” that is, of preachers, for preaching of the truth. And not their blood only, but have shed also the blood of thy holy faithful—I mean, whom for the true professing of the faith they have vexed, and at last slain. Therefore are they worthy that they themselves should again drink the blood of them and theirs; that is, should fall by mutual wars, tumults, and slaughters, verily before recited.

### Paragraph 1031

These things are confirmed by another angel who speaks from the altar, and not without cause from the altar. For we heard before in chapter 6 that under the altar the souls of those who have been killed cry out and say, “How long will you not avenge our blood?” Therefore, now a voice is uttered from the altar so that we should understand that God does not forget the blood of his saints, but avenges it in just and due season. Now, here, by the way, the omnipotence of God is commended, so that the ungodly may understand that in the time of affliction and vengeance, there will be no power able to resist the Almighty. To him alone be glory. Amen.

## Sermon 70: The iiii. and v. Angles shed their vialles.

### Paragraph 1032

.

### Paragraph 1033

And the fourth angel poured out his vial upon the sun, and it was given to him to afflict men with heat of fire. And men raged with great heat, and blasphemed the name of God who had power over the plagues, and they did not repent to give him glory. And the fifth angel poured out his vial upon the seat of the beast, and his kingdom grew dark, and they gnawed their tongues for sorrow, and blasphemed the God of heaven because of the pain and anguish of their sores, and they did not repent of their deeds.

### Paragraph 1034

The faithful do not regard their afflictions, sent by God’s just judgment, as the punishments of sinners [note 70.1], but as trials of their faith; although they acknowledge that they are justly afflicted for their sins. Yet, they commend God’s grace, which transforms the punishments of sinners into trials of faith. To the ungodly, punishments are plagues, which they cannot bear patiently, nor glorify God, but rather blaspheme him, and believe they suffer undeservedly. Therefore, God’s plagues are most grievous to the ungodly, although far more cruel things are prepared for them, namely, everlasting damnation in the next world. Therefore, the plagues of this world inflicted upon the ungodly are like certain preparations and introductions to more grievous torments.

### Paragraph 1035

The fourth Angel pours out his vial on the sun, and to him was given power to plague people with heat or fire. Many interpret this plague allegorically, understanding by the sun Christ enlivening the consciences of the faithful, and the latter to be darkened in the minds of people, choosing rather the darkness of Antichrist than the light of Christ. Therefore, consciences, erring and seduced with error, burn with various lusts and despairs, by which they are driven at length to various blasphemies. While I do not entirely reject this interpretation, in my judgment the sense will be clearer if we understand the fourth plague to be heat and great drought, a barrenness of the earth, a scarcity of grain, finally an intolerable thirst afflicting both people and beasts, and also breeding and intensifying hot diseases. For so we have read in the threatenings of the law: I will give a heaven of brass and an earth of iron. In the time of Elijah, for despising and rejecting the word of the Lord, God plagued Israel with a severe drought, as you may see in the third book of Kings, chapters 17 and 18. Jeremiah also describes a similar drought and heat in chapter 14. Again, the Lord defended Israel with a pillar of a cloud by day and a pillar of fire by night. Moreover, we have heard previously in the Apocalypse: the sun shall not shine upon them, neither any heat. And justly is this world plagued with burning heat, as it offends grievously, burns with various lusts, and also by wicked proclamations prohibits the spreading and refreshing of God’s word.

### Paragraph 1036

For this is the effect of the plague. And he says to me, burned with great heat. At first, he says, being inflamed with an exceedingly great heat, they were even raging mad. For we read in histories that those afflicted with too much heat have felt grievous displeasures and torments both of body and mind. Then he adds what follows from the former passage: the impatience of the heat provoked them to blaspheme God, and even him who has power over these plagues; namely, that because he has full power to do so, he will not deliver them so vexed with burning heat. Conversely, the children of Israel in their tents being stung by serpents, inflaming the whole body with the sting, did repent and did not blaspheme God. But coming to Moses, they said, “We have sinned, for we have spoken against the Lord and against you. Pray the Lord that he will take away from us these serpents.” Therefore, those who, through impatience, murmur against the judgments of God blaspheme the name of the Lord, and will not acknowledge themselves to be rightly and justly punished, seeking pardon. Finally, it is added that they did not repent so that they might give glory to God, and so on. For the Lord plagues us to the end that, being afflicted, we should repent and give God the glory, confessing, as I said before, that we are justly punished, and ought with weeping and wailing to turn to the Lord striking us. But these, like Pharaoh, do not acknowledge their sin, nor pray to God, nor are they amended, but many times overcome themselves in malice. Hereof we learn the difference between the godly and ungodly, and how both use themselves in afflictions. For they give glory to God and amend their life; these do not give God the glory, but become worse than themselves. To give God the glory is to give place to God, not to resist, but to acknowledge their sin and God’s righteousness; and not only this, but also the mercy of God and clemency toward the penitent, and the same to require humbly.

### Paragraph 1037

The fifth angel pours his cup upon the seat of the beast. It is more evident than it needs to be proven by testimonies that a seat or throne is used for a kingdom, as John himself immediately uses “seat” to mean a kingdom. Also, in times past, the masters, or rather ministers of churches taught while seated, and had their stoles and chairs in holy assemblies. The saying in the gospel is well known: “In the chair of Moses sit the scribes and Pharisees,” etc. It is known that in ancient times there were seats of patriarchs—Jerusalem, Antioch, Rome, Alexandria, Constantinople, and others—and that these are called Apostolic seats, because the Apostles taught there. And so the Apostolic seat is used to refer to the Apostolic doctrine itself. That seat erected and established at Rome by the Apostles and Apostolic men, the beast—that is, the Pope—has subverted, and in its place has erected the seat of pestilence, which he dares nevertheless to call the seat of Christ [note 70.7] and the seat of Peter. Christ has no more any seat on earth, except that he dwells in the hearts of the faithful church. Otherwise, the true seat of Christ is the right hand of the Father. The true seat of Peter is heaven itself. Rome is no longer his seat, for the Apostolic doctrine and patriarchal chair have been destroyed and trodden under foot, and in their stead an earthly empire or kingdom has been set up by the Pope. Indeed, he overthrows the Apostolic seats by force of arms. Now, therefore, God, having compassion on his people, pours out his wrath and plague on the see of Rome, illuminating men with the light of the Gospel, so that they may know and see the wickedness and abomination of the Roman see. This is a wonderful benefit to those who are enlightened, and a great grief and torment to the Roman people. For the effect of the plague follows: and his kingdom was made dark. This plague corresponds to the ninth of Egypt. For just as thick darkness plagued the Egyptians, bright light rejoiced the Israelites, so the Papists will be tormented by shameful errors, and it will grieve them to have their errors detected and their glory obscured; the faithful will rejoice in the light of Christ. For now the majesty of the seat, and of him who sits therein, begins—and has already begun—to be obscured. What was once called a holy seat is now, by the godly and learned, called wicked Rome, the whore of Babylon, the mother of all fornications, the den of thieves, Sodom, Egypt, the red harlot by reason of the purple senate of Cardinals, who wear red and purple. It is commonly said and truly, the nearer Rome, the further from Christ. They rightly call the Cardinals, bishops, and spiritual fathers the family and limbs of Antichrist, men deceived and deceivers, most corrupt with simony and filthy lust. Therefore, the kingdom of the beast—so he explains the seat—was made dark. Furthermore, it is added how the worshippers of the seat of the beast behave themselves. First, for pain and sorrow, indignation, wrath, and envy, they gnaw or bite their tongues, which is the gesture of angry men, and impotently angry, that is, burning with furious rage [note 70.8]. This is a phrase of speech signifying how they will rage with great fury against the truth that has been opened, which they would have utterly hidden and suppressed. Again, they blaspheme the Lord of heaven and maker of all, both because he afflicts them with sores and sundry plagues, and also because he casts a darkness upon their kingdom. For even now, the Romans call the preachers of the gospel deceivers and heretics, and the very doctrine of the gospel, heresy. But this reproach redounds to him who is the author of that doctrine. Finally, they do not repent of their doings, of their simony, of their crafty juggling, sacrileges, idolatry, and all ungodliness. And the apostle says that evil men and deceivers will grow worse and worse, deceiving and being deceived. Therefore, it is no marvel that you see the Papists at this day, with a stiff neck, proceeding obstinately in their errors. But the greatest plague is to be forsaken by God, and stubbornly to maintain their errors and ungodliness, and therein to persevere. The Lord deliver us from evil. Amen.

## Sermon 71: ¶ The ſixte Angel ſhedeth his vialle.

### Paragraph 1038

.

### Paragraph 1039

And the sixth angel pours out his vial upon the great river Euphrates, and the water dried up, that the way of these kings of the East should be prepared. And I saw three unclean spirits, like frogs, come out of the mouth of the Dragon, and out of the mouth of the beast, and out of the mouth of the false prophet. For they are the spirits of the devils working miracles, to go out to the kings of the earth, and of the whole world, to gather them to the battle of the great day of God almighty. Behold, I come as a thief. Happy is he that watches, and keeps his garments, lest he be found naked, and men see his shame. And he gathered them together in a place called in the Hebrew tongue Armageddon.

### Paragraph 1040

The sixth angel pours out his vial on the great river Euphrates; the end of this shedding is that the way might be opened for the kings of the East, that is, that Babylon might be taken. This plague chiefly pertains to Rome and the Roman church. The speech has an allegory, or a secret comparison and allusion to old Babylon. We read in the fifth chapter of Daniel that Babylon was taken the same night wherein Balthazar, king thereof, had prepared a sumptuous banquet and looked for nothing less than destruction. Herodotus and Xenophon report how the kings of the East, Darius Priscus, who is also called Medes, and Cyrus of Persia besieged the city round about; but where there was no hope to win it, Cyrus let out the Euphrates by ditches, so that the army might wade over the river; and so was the city laid open and taken on the side where it was fenced with the river. Euphrates therefore fortified Babylon and brought to it many other commodities and pleasures. Here are signified, by Euphrates, riches, defense, pleasures, commodities, tributes, and customs, which the Roman churches call sacred or of the holy church. These commodities and pleasures are diminished by the kings of the East, by true Christians, whom the Scripture calls the Kings and Priests, and derived and put to another use. Wherefore the power of the Roman church begins to decay, to the intent that at length it may be taken and abolished by the Lord Christ himself. Doubtless the true Christians understand, believe, and profess that Christ alone is the Savior, neither that there is salvation in any other. And that this is given freely; that they be mad and commit Simonry and sacrilege, which in this case practice and make merchandise. Read the fifty-fifth chapter of Isaiah, and the eighth of the Acts, finally the first and second chapter of Paul to the Colossians, wherein most diligently is declared, illumined, and set forth that by Christ alone we are absolved and in him alone have all fullness. And what time the common people do hear this, to wit that by those Roman trifles, fairs of pardons, and other crafty jugglings they are deceived and robbed of their substance, they shut by and by and make fast their chests, their purses, their cellars, and garners. And so dries up the river of wealth and pleasure; it dries up also when the godly deny to give other customs, as tithes, palles [?], first fruits, and such other like things. So is the way prepared for the kings of the East, so begins Rome, the second Babylon, to be taken and come to naught.

### Paragraph 1041

It follows furthermore,[note 71.3] how Antichrist will fight against the faithful and godly laboring to dry up Euphrates, for the maintenance and increase of his kingdom: and where he might briefly have said, he shall send forth ambassadors unto all kings and princes, to stir them up against the gospelers, for the defense of the privileges, rights and revenues of the See of Rome, he had rather most diligently describe those ambassadors and show their destruction. It forces very much to have known the Popes legates. For they are marvelous pestilent to the church of God: for we have not only experience of it at this day, but also by the reading of all histories, that great evils[note 71.4] and all calamities in a manner have been brought into the church, and are also at this day, through the instigations of those legates. I touched a little before what mischief Cardinal Julian Cesarine, the legate of Pope Eugenius, wrought in Germany, Bohemia, Poland and Hungary. What is done in our time, and has been done in our fathers’ memory, it were too long to rehearse. If our elders had by the doctrine of Jesus Christ revealed to the church by John understand and known the nature of the Popes legates, they might easily have eschewed wherewith they have undiscreetly entangled themselves, and suffered great loss and hindrance. I speak nothing here of ambassadors and embassies of kings and commonwealthes uncorrupted.

### Paragraph 1042

First, he diligently demonstrates the origin of legates, so that we may understand that they are led by a wicked spirit, and that their vocation is not godly, but devilish. He shows a threefold origin, within which, in truth, they may all be reduced to one devilish unity. He says that Primasius, explaining this passage, saw one spirit, and because of the number of the parts of one body, he says three, so that the number of the wicked might be expressed as being led by one devilish spirit. Therefore, the first den from which the legates break out, he calls the dragon’s mouth. The dragon is spoken of in chapter 12; and there is no one who does not understand that it signifies the devil himself. They therefore come forth from the devil. For all the affairs of their embassy consist in lies, schemes, practices, finally, in corrupting the truth and sincerity of the gospel, and also in factions and dissensions, in slaughter and bloodshed. And the devil was from the beginning a liar and a murderer, as the Lord himself says in John 8. And so far, they are from the dragon’s mouth. The same arise also from the mouth of the beast. For they come furnished with the Pope’s authority, legates lateral with full power. Of the beast, I have spoken in chapter 13. Finally, they come out of the mouth of the false prophet. The true prophet and pastor, supreme and only of the universal church, is Christ, the Son of God. Antichrist is that false prophet and chief seducer of the whole world, as is said in chapter 13. Therefore, the legates come, sent from the Pope, who has put into their mouths words, instructions, or commissions that they should speak those things which are of false prophecy. However, explaining himself more plainly, he declares of what sort the legates shall be: namely, three unclean spirits. An unclean spirit is everywhere in Scripture called the devil or Satan, truly in nature and effect. For as the Spirit of God is called holy, so this is, contrarily, unclean. For he himself is, by nature, or rather by his own corruption and revolt from God, impure, and the author of all impurity and uncleanness. He signifies therefore that those legates shall be men of a devilish uncleanness. And indeed, if you add to this the lives, manners, and conversations of those lateral legates, and of their families, you will find, in effect, nothing else but extreme uncleanness, filthiness, and beastliness, monstrous lust, whoredom and adultery, and detestable fornications, wonderful surfeiting, bloody schemes and counsels. Therefore, the thing itself speaks: and the things that the legates do everywhere are a commentary on this passage. And where three unclean spirits are reckoned, some explain it of divines, lawyers, and religious persons, such as monks and friars, from which three sorts the Pope’s ambassadors are mostly chosen. I understand simply by the third number that those legates shall be most furnished with all hostile authority, and that they shall all agree well among themselves, and all help one another: so that whatever one seems to lack, another may supply. Solomon in Ecclesiasticus says, “A threefold rope or line will not easily break.”

### Paragraph 1043

But now, so that no one should lack understanding, he presents, through a parable, what kind of men these legates will be—truly frogs of the marsh or fen, and insistent, tedious criers, foul and filthy. And he does not say that they are frogs in reality, but like frogs. For just as frogs, by their insistent crying, are most tedious and troublesome, and those of the marsh are also filthy, so these legates love earthly things and filthiness. And by their complaints, accusations, provocations, writings, and arguments, they are altogether frog-like and fen-like, hateful both to God and men. They are nothing ashamed; if they are interrupted a little, they immediately return to their old song, [?persisting]. For there is no other tune with them but [?that]. Primasius reasons greatly about frogs. Among other things, it is fitting for false prophets, like frogs crying in the night, to make a damnable noise by barking out errors. For frogs, because of their location, appearance, and troublesome noise, are so hateful that the Devil, with his works, is known to be abominable to the truth, and with just fire rightly condemned, and so forth. Thus he says. And just as the frogs of Egypt were raised from the dust by the devilish art of the magicians, and cried out against God’s truth, calling back the people of God through Moses and Aaron to true liberty and worship of God, so too the Pope’s legates harass with talk the preaching of the gospel, the free deliverance, the Christian liberty, and true service of God. And just as the frogs double and repeat even to the point of making one weary to hear, so these fen-like beasts of Rome always have in their mouths, “the most holy See,” “the most holy father,” “the holy church of Rome.” “The holy church of Rome does not err,” “the holy church of Rome must be obeyed.” “He who will not obey her is a heretic and a schismatic.” These things they will reiterate and repeat many times and often, to all people, and in all and every cause, their one and the same song [?sounding].

### Paragraph 1044

The Lord adds through John, and thus declares it even more clearly: for they are spirits of devils working miracles, &c. “Daemon” (which is here used in Greek for devil) has its name from sundry knowledge and skillfulness of things, and seems to be in a manner indifferent, although it is commonly put for the Devil. Nevertheless, for a difference, they are called Eudaimones and Cacodaimones, as if good and evil workers. For the Greeks say that “daemon” is called from “daina” that is knowing or skillful. For “demiurgos” is called an expert Artificer. The Lord therefore signifies that the Popes’ legates shall be spirits of devils, that is to say, spiritual fathers (but endowed with the spirit of Satan), wise men or skillful, crafty workers to bring their matters to pass. And therefore he adds, working wonders. Whereby he seems to allude to the magicians of Egypt, who also wrought miracles and detained King Pharaoh in lies against the truth. Paul moreover in 2 Timothy 3 compares the wise men and ministers of Antichrist to the magicies of Egypt. And it is well known that the legates do everywhere boast of miracles which have been done in their church and religion, and so keep still the hearts of kings and princes in papist errors. Of miracles speaks Paul in 2 Thessalonians 2. And I have said something hereof in Chapter 13.

### Paragraph 1045

Moreover, here is shown the end of all the treatises and counsels of the Popes’ legates: that they might go forth to the kings of the whole earth, to assemble them to battle.[note 71.10] And so they shall creep into the courts of all kings and princes. You will surely find the Popes’ legates in a manner in all the kings’ courts. And what do they do? They surround kings and princes. They see that no faithful man be admitted to the kings’ speech, they learn to know all the kings’ counsel, which they write and signify to Rome: and if they dislike anything, that they may invalidate and subvert the same; and that they always repeat that song of theirs, namely [illegible], that is truly obedience, which all men owe to the Holy See. Finally, that they arm kings and princes to defend the church of Rome and destroy heresies. This, I say, is the battle of that same day of the great God almighty, that is to say, which shall be marked by the coming of the Son of God unto judgment, and which shall endure to the coming of Christ unto judgment, when he shall then avenge the blood of his, from the hands of that hideous beast. And he calls the day of judgment, the day of the great God, as does also Jerome in the second letter to Titus. And the day of God almighty: as he who shall then show his omnipotence, and even his divine power, which seems now to the ungodly by reason of his long forbearance to sleep. This necessary and most profitable description Saint John has set here, by the revealing of Jesus Christ, so that we should watch and beware of them.

### Paragraph 1046

Hereafter follows a faithful admonition and exhortation to watchfulness, lest we fall asleep and perish with the Antichristians in the cares and pleasures of this world. And he says how that day of the Lord will come suddenly, and when we shall least look for it. For the Lord here repeats that thing which he said also in the Gospel: behold, I come as a thief. These things are read in Matthew 24 and are repeated by the Apostle in 1 Thessalonians 5. And verily, that same sudden coming of the Lord excites the minds of us all and provokes to watchfulness, lest we should at unawares be oppressed. He adds also immediately a profit prepared for them that watch. Happy, says he, is that man that watches. He adds moreover, how the godly should demean themselves in watching. How they must keep their garments, that they be not defiled; and take heed moreover that they walk not naked, lest their filthiness be espied. Touching garments I have spoken most largely in another place of this book. He keeps his garments, that keeps his life and conversation unspotted of worldly filthiness. He does not walk naked, which puts on Christ. But his shame is seen, that sins impudently. But chiefly is their shame seen, whose whoredoms, adulteries and fleshly lusts are known and open to the eyes of all men. And here is the state of them to be lamented who are called spiritual, and rather in deed to be detested than to be described. Blessed are they whose sins are covered, and happy are they that have learned to be ashamed. Unhappy are as many as cannot blush, but set such a face upon the matter that they glory in their sins and wickedness.

### Paragraph 1047

After this, he touches briefly on the destruction of both the legates and those deceived by the legates, as well as those who fight against God and true religion to maintain Roman authority. The legates indeed assemble people of their faction for battle against the faithful, but the Lord has gathered them into a place called in Hebrew [illegible], which some interpret as the destruction of the River, and others as the army of desolation. But however that may be, the sense seems clear: they are assembled by the legates that they might withstand, or prohibit the destruction of the River and the ruin of Rome. But the Lord will also assemble the very same people, that in the same place and through the same means they may be destroyed by the Lord. This, we believe, will be accomplished finally at the last judgment. To the Lord Christ, our redeemer and revenger, be praise and glory. Amen.

## Sermon 72: ¶ The seventh Angel powers out his vialle.

### Paragraph 1048

.

### Paragraph 1049

And the seventh angel poured out his vial into the air. And there came a great voice out of Heaven from the throne, saying, “It is done.” And there followed voices, thunderings, and lightnings, and there was a great earthquake, such as was not since men were upon the earth, so mighty an earthquake and so great. And the great city was divided into three parts. And the cities of nations fell. And great Babylon came into remembrance before God, to give unto her the cup of wine of the fierceness of his wrath. And every island fled away, and mountains were not found. And there fell a great hail, as it had been talents, out of Heaven upon the men, and the men blasphemed God because of the hail, for it is great, and the plague of it was sore.

### Paragraph 1050

The seventh and last cup poured out into the air signifies the disturbance and alteration of all elements, and the horrible, but just judgment of God, and finally the end of all things and eternal punishment. The things are expressed with figurative language, taken mostly from the Prophets, and by a secret comparison brought from holy history. This is done for this consideration, that all things might be more majestic, and that every person should more diligently search for the meaning of an excellent matter, which once found, he might keep and retain in perfect memory.

### Paragraph 1051

And when the air is disturbed, diverse and horrible tempests arise in the air. And the Lord Jesus in the Gospel according to Matthew testifies that about the final coming of Christ, the powers of heaven shall be moved. And as soon as the cup was poured out into the air, and a great tempest arose, a voice sounded, it is done. By which voice is signified how all things are at an end, even of the whole world, much more of wicked papistry. And this voice is heard out of the very temple of heaven and the throne of God, lest we should doubt anything of the truth and certainty of the sentence given, and again of the virtue and power of him that pronounces it. Therefore are they shamefully deceived, so many as affirm the world to be everlasting, and that they shall reign always upon earth, and enjoy the pleasures thereof. A voice from heaven, out of the most holy temple of God, and even out of the most sincere throne of the Almighty, speaks, that it is done. For he speaks of the time to come as though it were past, that we might as certainly know that all worldly and popish things should have an end, as we undoubtedly know the things to be done which are already past. Let us therefore watch, and put no confidence in the things of this world, which are most deceptive. All things shall fall to decay and come to naught, men only, and the blessed spirits, remaining through the grace of God, the unhappy also remaining perpetually, appointed to perpetual punishment by the justice of God.

### Paragraph 1052

And just as the holy Prophets depicted God’s judgment through figures for men’s eyes to see, so now the Lord Jesus, through Saint John, uses figurative language to shadow the terror of that horrible judgment. For he says that there will be thunderings, voices, lightning, and thunderbolts, and an earthquake so terrible that the world has never felt anything like it at any time. For Saint Peter also, at the end of his second letter, recounts terrible things about the last day and the burning of all worldly things. But the quaking and terror of men’s minds will be yet a great deal more terrible than all these.

### Paragraph 1053

The Lord in the Gospel according to Matthew. Then he says, “all the kindreds of the earth.” For the ungodly, whose consciences are wicked and corrupt, shall feel those terrors and torments unspeakable. The godly, as according to the saying of our Savior, do not come into judgment; so although that they also, by reason of the infirmity of the flesh, be somewhat astonished at the sudden alteration of things, and the terrible tearing and crashing of all elements, yet since they have known before that the same should come to pass, and believe the Savior’s saying, “your redemption draws near,” they gather up their spirits and comfort themselves in Christ, and rejoice in him, coming to judge or condemn the ungodly, but to save the godly. And herein is alluded to sundry stories of the holy scriptures, but chiefly to the burning of Sodom, to the drowning of Pharaoh in the Red Sea, and the ruin of Jericho, &c. Those were verily but several destructions, and yet terrible above measure; therefore what think we that last destruction will be, which shall be general?

### Paragraph 1054

Then shall that great city be cut asunder, [note 72.6] the universality of men in the great church, divided into three parts: that is to say, in the end there shall be found three kinds of men in the church. There are true Christians, who attribute to Christ his true glory, that is, all things of true salvation, and cleave to him alone by sincere faith. There are Papists, who after the letter ascribe many things to Christ, but not as is fitting: for they ascribe to Antichrist those things which belong to Christ alone; and in associating with him, things that ought not to be associated, they deny Christ. For if the Pope is head of the universal church, if he is king and priest, and so forth, why is Christ preached to have those things alone? There are moreover Moderates, who will not seem to deny Christ, and yet attribute a little to Antichrist, whom nevertheless in many things they despise utterly. These have no certain religion, but one established and conceived at their pleasure, as it pleases them to believe this or that. There is a great number of these men at this day, deriding and mocking whatsoever is not tuned after their own light, and wanton Lucianicall wits. You may find also in the gospel a field sown with sundry seed, to bring forth most diverse fruits, yea even cockle and darnel, which at length in the end of the world shall be gathered. &c. Matthew 13.

### Paragraph 1055

Moreover, the cities of the Gentiles (he says) will fall [note 72.7], by which I understand the Jewish, Turkish, and foreign religions, torn into various sects or heresies. But each of these has its societies, rites, and laws, which they commend as the best, and as those that will endure forever; but they too will fall. The sole religion or faith of Christ shall prevail and overcome. Arethas expounding this passage in the same manner: The cities of the heathen, he says, falling down, are diverse opinions about faith, and so forth. They (I say) have all fallen.

### Paragraph 1056

[note 72.8]But especially, he affirms and demonstrates diligently, it was fitting and necessary that the city and church of Rome should be destroyed and subjected to perpetual torments. I have declared sufficiently before that Babylon is Rome, which in truth is great, not only in Italy but throughout all France, Spain, Germany, and other realms. The city and church of Rome have seemed to many to be everlasting and triumphant forever. Herein the Epicureans cry out that God does not care for these lesser things, but that every man lives here, either happily or unhappily, according as he has discreetly and skillfully framed his life, that [illegible] knows nothing of our pleasures and displeasures, and our conversation. But contrariwise, Saint John affirms that the Lord has remembered Babylon, and to have remembered her so that he has determined to commit her to torments. He expresses this through a prophetic phrase of speech, that he might give unto her the cup of wine of indignation, or fierceness of his wrath, that is to say, that he might punish her accordingly, as the great indignation and wrath of God requires. Therefore she shall have no small punishment, for the wrath of God is not light, but most grievous and hot. For he requites and recompenses the slackness of punishment with the extremity of pain and torment. You may read similar things in Malachi 3. Concerning the cup also, of the wine of God’s fury, has been spoken of before from the Prophets.

### Paragraph 1057

Now also, among other things, by a figurative speech he shows that the ungodly have no refuge, nor way to escape. Otherwise, the wealthier would hide themselves far away on islands, to be out of range of danger; many would flee into the mountains, that they might lurk safely. But now he says that even the islands flee, and therefore, in fleeing, they cannot be overtaken. He adds that the mountains—that is to say, no places of refuge or hiding—can be found. Therefore, nothing remains but that all the ungodly, in general, being taken should be put to torment.

### Paragraph 1058

Furthermore, he adds that hailstones as big as talents should be cast down from heaven upon wicked men, and that those who have not been remembered to have fallen in no memory of him [?]. And he seems to have alluded to the story of the Canaanites, which is in the tenth chapter of Joshua. To be short, here is signified that the grievous and inevitable judgment of God pronounced against all the ungodly shall, at the general judgment, torment the wicked with such an extremity that no eloquence of men, no sense nor understanding can attain unto, for it is always more grievous. Primasius, expounding this place, says: he sets the wrath of revenge in hail, whereof we read: “The wrath of the Lord falls down like hail.” Neither does he without cause mention a talent’s weight, for with equity will he inflict judgment, and so forth.

### Paragraph 1059

Here is shown the stubborn and incurable rebellion and impatience of the wicked,[note 72.11], by which they are provoked against God’s judgments, and utter blasphemies against the Judge himself and his judgment. I have treated these things more briefly, for we have heard much the same before concerning the end of the 11th Chapter. To the Lord be praise and glory.

## Sermon 73: ¶ The Judgment or punishment of the pourple whore is described: and also the sin, and vngodlines of the same.

### Paragraph 1060

.

### Paragraph 1061

And there came one of the seven angels, who held the seven vials, and spoke with me, saying to me: Come, I will show you the judgment of the great prostitute, who sits upon many waters, with whom the kings of the earth have committed adultery, and the inhabitants of the earth are drunken with the wine of her fornication. And he carried me away into the wilderness in the spirit. And I saw a woman sitting upon a rose-colored beast, full of names of blasphemy, which had seven heads and ten horns.

### Paragraph 1062

He has hitherto discoursed generally about God’s just judgments; albeit, in the meantime, he has touched on some particular matters concerning Babylon or Rome, rather than reasoned. And now, consequently, individually and evidently, he handles the destruction or end of the beast, and of its image, of old Rome and new—I mean both the empire and papistry—which he seems to point out as it were with the finger. In chapters 13, 14, and 16, he signified something of this; now he pours forth everything with a notable abundance and evidence. And the same thing that I said at the beginning of this passage, I repeat here again: that hereby are affirmed how the justice of God is shown to be the end of good and evil, that the godly are confirmed, and the judgment to come established, as we expressly confess in the Apostles’ Creed. The sum of all is this: old and new Rome, the empire and the Popish kingdom, which is the kingdom of Antichrist, shall perish for sins and great enormities. For he joins together the beast and the image of the beast, the beast and one sitting on the beast, a proud harlot, so that they cannot be separated. Therefore, the passage must be expounded of both empires.

### Paragraph 1063

And lest anyone should think me to be led by a misguided sentiment in understanding these things concerning Rome, I will demonstrate by the testimony of both God and man that this interpretation is ancient, not new, true, and not contrived. For immediately the angel himself, as we shall hear, explains these things concerning Rome. Arethas, also an ancient interpreter of this book, says that diverse interpreters understand by the whore old Rome. And he continues immediately: While she speaks of the mother of harlotry, whether you will call her old Rome or new, or the time of the coming of Antichrist (behold, he says, the time of Antichrist), you cannot err from the truth, for both cities [Rome and Constantinople] have held the empire, and each has been satiated with the blood of holy martyrs, and so forth. Thus far he. And what the most ancient writer Tertullian and Saint Jerome have spoken repeatedly concerning Babylon and the purple whore, I have reported previously in chapter 14.

### Paragraph 1064

And in this order he proceeds. First, he shows the author of this revelation, after he gathers the sum of the revelation or vision. For again he treats by visions, to the end all things might be more lively and evident. Indeed, some divide this work from chapter seven to chapter twenty-one as the sixth vision, as I admonished in the beginning of this work. Then he notes the place and manner of the vision. Finally, he propounds the vision itself, and immediately adjoins the exposition thereof. And in the process of this matter, he uses a judicial kind of pleading, and that after a prophetic manner. For the prophets most often, and in the beginning, set forth the sins and wickedness of the people before the eyes of all men, and then annex unto it the judgment, pain, or punishment. For so does John also at this present.

### Paragraph 1065

The author of this horrible vision is the Lord Christ himself, but he uses the ministry of an angel, and that of one who, coming out of the temple of the divine majesty, was appointed with six others to pour out plagues and vials. This is the head minister. It seemed fitting that the judgment of Babylon should be uttered by an angel that had rule over thunderstorms. The Lord Jesus himself will take vengeance on the beast for whom this triumph is reserved. And we understand that such things as are set forth here have proceeded from the high Bishop himself, Jesus Christ, and that the manner of speaking is angelic, heavenly, and godly. Who then shall blame us, if we, using the words of angels and of Christ himself, say that the Bishop of Rome and all Popery is that purple, great, and most common harlot? It appears also to many who seem godly that moderation is neglected when these things are repeated by preachers. It seems that they would shut and stop the mouth of Christ himself. However, they attempt that in vain. For if the preachers hold their peace, the stones will cry. For it behooves that like as the glory of Christ, so the shame of Antichrist should be known to the whole world. But they offend most grievously who, in the sermons made against Antichrist, require I know not what modesty, as though he ought to be spared, which spares no good man. As though that doctrine were not modest, which is taken and received from the mouth and words of Christ. Afterward, in chapter 18, we shall hear the Lord command: “Render unto her as she has rendered unto you,” and so forth.

### Paragraph 1066

Secondly, he summarizes the essence of everything in a few words and indicates to what we should refer all things. “Come,” says the angel to John, “and I will show you the judgment, condemnation, and punishment of the great whore.” And when he says, “of the great whore,” he intimates what the crime or cause of punishment is: fornication, infidelity, or ungodliness. This vision also pertains to this, so that we might understand how Rome should be punished or destroyed, that is to say, the Roman Empire, or the kingdom of the Pope or of Antichrist, and why or how it deserves to be destroyed. She is a whore, and a great and wandering whore. And who does not know that a marriage is contracted between God and all the faithful? That God is the bridegroom, and the church his spouse? She is then bound and coupled to her husband alone in faith and troth. If she breaks this faith and loves others, gives herself to them, calls upon them, and honors them, she is a whore. Of this I have spoken many times, both in this book and elsewhere.

### Paragraph 1067

And it is a very common thing in the Scriptures to call unfaithfulness, impiety, superstition, and idolatry, fornication or whoredom. If anyone desires testimonies of this, he will find them in Judges 8, Isaiah 1, Jeremiah 2 and 3, Ezekiel 16, Hosea 1, 2, and 3, and other places. Rome, therefore, was a great prostitute, and is even now a most stinking harlot, for which she is full of idolatry, the worship of creatures, and abominable superstitions. Neither is she herself only polluted with all filthiness, but compels moreover the whole world to serve, and that to serve in idolatry and superstitions. What will you say that through the wonderful providence of God it came to pass that a woman feigning herself a man did climb up to the See of Rome, was created bishop, and called John 8, which was one Gylberta, a great whore, born at Mentz? For thus would God declare that the Bishop of Rome sits a whore upon the beast. And herein I follow the constant consent of all historians. Nevertheless, I am not ignorant that there are some who have thought how this John was intruded into the seat of an harlot, and for that cause was called an harlot.

### Paragraph 1068

Furthermore, old Rome had power to do these things, for she sat upon many waters, that is, had dominion and rule over many people and sundry nations. All the kings of the Earth have committed whoredom with her, lest they have submitted themselves to the Romans, bound themselves in league, and received of them superstitions and idolatry. For the children of Israel were also said to have committed whoredom with the Egyptians, for that they had joined amity with them and were become fellows in profane religions. And so new Rome, the Pope’s kingdom, stretches far and wide, and the kings and princes of the Earth commit whoredom with her. Therefore doth the word of the Lord call it filthy whoredom, which the Romans name a holy bond and obedience. It is added, and they that dwell upon Earth are made drunk. For he signifies that, being infected with errors, yea rather besotted and clean out of their wits, they have been mad in idolatry, and yet rage in their superstitions, like drunkards, and cannot, for fury, receive the preaching of the gospel. Touching this wine of fornication and whoredom, and of that drunkenness, I have spoken in chapter 14. And it is aptly spoken that dwellers upon Earth are made drunk, not so much for that men dwelling upon Earth are made drunken, as for that earthly minds and choked with earthly desires shall become faithful worshippers of the Roman See.

### Paragraph 1069

Thirdly, he explains the manner of the vision [note 73.11] thus: “I was carried away in spirit.” Therefore, while his body remained in Patmos, in spirit he saw a woman sitting on a beast, and consumed with fire. Such are many visions and sights in the prophets. And he notes also the place where he saw the beast, not in heaven, nor in the temple or tabernacle, nor in a fruitful place [note 73.12], but in the wilderness. Isaiah calls the Gentiles and heathens a wilderness. And truly, the Romans and their new superstitions would have had no place in the church, but are outside the church. God forbid that we should acknowledge the church of Rome to be the head of all faithful churches. And at this day, many of those who are called most holy and most reverent differ nothing from the Gentiles, their titles and hypocrisy only excepted, as was spoken of before in chapter 11.

### Paragraph 1070

Fourthly and last, he exhibits this vision, or a glimpse of old and new Rome, and the ruin and destruction of them both, and with all describes most diligently the wickedness of either. First, the beast must be considered, after the woman sitting on the beast. The beast represents the figure of old Rome; the woman, of the new and of Papistry. And the woman sits upon the beast, for the image of the beast has succeeded and has placed her seat in old Rome. For Daniel also affirms that Antichrist shall pitch his seat or palace between two seas [note 73.13], namely, the Adriatic Sea, commonly called the gulf of Venice, and the Tyrrhenian or Tuscan Sea. And the beast is rose-colored, of a red and bright color like crimson, for Rome has been most cruel and bloody, and swimming altogether in the blood of all men, but especially of Christians. How much blood shed Marius, Sylla, Pompey, Julius, and others, as histories Pliny has reported. Rome has with sword and fire destroyed the whole world. The ten persecutions of Christians before the reign of Constantine are most commonly known.

### Paragraph 1071

As I showed in the thirteenth chapter, the beast was full of blasphemous names. Rome was filled with chapels and idols, and daily blasphemed God, Christ, and the gospel, tearing the church asunder. Concerning the seven heads and ten horns, there is also speaking in the thirteenth chapter. And certain things will follow in this same chapter, plain enough. And thus much hitherto concerning the old beast: now follows the woman sitting upon the beast.

## Sermon 74: ¶ The same matter is yet still treated of, and the vision is expounded.

### Paragraph 1072

.

### Paragraph 1073

And the woman was arrayed in purple and scarlet, and adorned with gold, precious stones, and pearls; and she held a golden cup filled with the abominations and filthiness of her fornication. And on her forehead was written a name, a mystery: Great Babylon, the mother of fornication and of the abominations of the earth. And I saw that the woman was drunk with the blood of saints, and with the blood of those who bear witness to Jesus. And when I saw her, I marveled greatly. And the angel said to me, “Why do you marvel? I will explain to you the mystery of the woman and of the beast with which she is carried, which has seven heads and ten horns.” The beast that you saw was, is not, and will ascend out of the bottomless pit, and go into destruction; and those who dwell on the earth will wonder—those whose names were not written in the book of life from the beginning of the world—when they behold the beast that was, is not, and will be. And here is a mind that has wisdom.

### Paragraph 1074

He describes exceedingly well and vividly, setting forth for all men to see that very woman previously named the great whore. Toward the end of the chapter, he depicts death itself, and says: “And the woman which thou sawest is the great city.” [Note 74.1] Indeed, she is great, as one who holds dominion over the kings of the Earth. He means therefore the very city of Rome, and even the papal and Roman church, and the pope himself with all his officials and chaplains, who constitute a great and powerful city, ruling over all kings and princes of the Earth. For who does not know that Rome and the prelates of the church rule even above magistrates and princes? Consider what is done in the courts of kings and princes, and by whose advice and counsel the princes of the Earth are chiefly governed. And in calling papistry a woman, he alludes to the fifth and seventh chapters of the Proverbs of Solomon, who also likens crafty and deceitful philosophy, and worldly wisdom to a beautiful woman full of various schemes and deceptions.

### Paragraph 1075

And now also he clearly and vividly depicts the apparel and exceedingly immoral behavior of this woman. She is furnished with no good qualities inwardly, whereby she might commend herself to the world and to her lovers; therefore she excels in outward adornment, while inwardly she is full of all abominations. For she is like to all the world, like the Pharisees and hypocrites, whom the Lord says are outwardly shining tombs of marble, but within are filled with all corruption and filthiness, and even stench. And by this is chiefly signified that the church of Rome and the kingdom of Antichrist presents itself altogether with worldly trappings, namely with gold, silver, precious stones, and all costly array. These things are found among the people of old, under the discipline of the law; they are also found in the images of the gentiles, who suppose that God is not rightly honored, but with the precious things of this world. But we know that the Levitical priesthood is abrogated with all outward adornment, and that now the church decks herself with virtues, hates, and abhors outward embellishment. Lactantius refutes at length the external adornment in religious worship, in his book *Institutes on True Worship*. Moreover, all ancient writers show that God is not worshipped by Christians with gold and silver, but with faith, charity, and righteousness. What will they say who claim that Daniel in chapter 11 shows that Antichrist will worship God with gold, silver, and precious things? This is condemned and rejected.

### Paragraph 1076

And does it not appear here plainly that the Lord Jesus himself has set forth to us the Pope or Antichrist, painted as it were in a table? For he appears altogether such, and in such like apparel he shows himself to be seen of all men, as the whore of Babylon is decked at this present. And he challenges to himself this apparel by a certain right. For the Papists bring forth a false, feigned donation of Constantine, among other things, pronouncing thus, in Distinct. 96. We give and dismiss to blessed Sylvester and to all his successors the Lateran palace of our Empire; moreover, the diadem, to wit, the royal crown of our head (which the Pope calls a kingdom, and has made it triple) and also our purple robe and coat of crimson, and all our imperial array, &c. I cannot here omit, but must needs write out a few things of Platina, the Pope’s secretary, *De Pontificibus*. For describing the life of Clement V at the Pope’s coronation, he says that Philip, King of France, and Charles his brother, John, Duke of Britain, were present, who, overwhelmed by the falling of a wall, dies (see the divine and just judgment of God) with many others, whilst the pomp of the Coronation, as the manner is, was led through the city. King Philip also by the same ruin was sore hurt and lamed; the Pope, stricken from his horse, lost a ruby out of his mitre that cost six thousand ducats.

### Paragraph 1077

This harlot moreover drinks to all nations from Circe’s cup, which the Lord calls of Gold.[note 74.5] And it signifies doctrine. For to give drink is to teach, as in Ezekiel. Gold betokens sincerity and purity of doctrine. Doubtless under pretense of sincerity and divine truth, Rome has easily persuaded all people to receive the doctrine of the Roman See. For the Pope has named himself Apostolic, and the church of Rome also Apostolic. And in the Canons has left written: so are all the laws of the See to be taken as though they were confirmed by the godly mouth of Saint Peter himself. Distinct. 19. Read chapters 20, 21, and 22. Distinct. Therefore the more simple sort of the world have supposed that they receive the very word and laws of God, when they received the decrees and doctrine of the church of Rome. But our Lord Jesus Christ declares here what has been, and what is even at this day their doctrine, and says: full of abominations and uncleanness of her filthy lust. And the scripture calls abomination, idolatry, as in Deuteronomy, chapter 7. Moreover, the false worship of God, superstition, and such things like. The uncleanness of whoredom in the Prophets is nothing else but perverse doctrine and perverse religion, not attributing all good things to God alone by his Son, but rather dividing the heart and applying them both to creatures and to wicked worships. But such is the doctrine and religion of the Roman See. Therefore her great sin is here recited, that she has with her evil and venomous doctrine seduced and infected all nations, and even now retains them in superstition and idolatry. A like place is in chapter 16 of Ezekiel. And I suppose that our Lord Jesus Christ used here now words very filthy, to the intent he might pluck away from the Roman decrees and Canons their authority and visibility, and that their filthiness might appear unto all men, and be known, and eschewed.

### Paragraph 1078

Furthermore, lest anyone should be ignorant of what the woman is, who is here portrayed for view, and that all might flee that great witch Circe, he writes her own name on her forehead, so that all men might read it, and that she might by no means be unknown. For he calls a mystery the understanding or signification of a secret; for by a figure of speech Rome is called Babylon. Whereof I have spoken before. And after the true signification of the word, Babel signifies confusion. And Rome has brought an exceeding great confusion into the church. For just as the primitive church of Rome set forth the gospel in the western lands, so now, the first simplicity and purity being extinguished, later bishops, regarding more ambition and covetousness than humility, liberality, and godliness, have brought into the whole world all manner of idolatry and superstition. Certainly, she is called by the Lord expressly the mother of the harlotries and abominations of the earth. For we may thank the church of Rome for all the corrupt doctrine and ungodliness that is in the church. She is the origin of idols, the Mass, and other abominations. Therefore, she is most worthy to be punished with most grievous torments. And this verily is the most worthy title of the Roman church. Others call her Apostolic, divine, chief, and of all most holy; the Lord Christ calls her Babylon, and that great harlot, and even the mother of the abominations and harlotries of the whole world, fighting against God and his Anointed. Therefore, let all the holy and obedient children of God the Father flee from her. I will not here make rehearsal of how she is also the mother of all abominations and harlotries, even after the flesh. For whilst the church of Rome has prohibited lawful marriages and, permitted by God to ecclesiastical persons, it has opened the gates to fornications, adulteries, harlotries, and lusts abominable. There need no words; the thing itself speaks.

### Paragraph 1079

To these, he adds a crime of all others the most grievous. And here also he employs an amplification. For he says that great Circe, the most venomous witch and sorcerer, is not sprinkled, or immersed, or wet, but drunk with the blood of saints, that is, of holy martyrs, who have borne witness to Jesus Christ by preaching the gospel and ascribing all things of salvation to Christ alone. But how many thousands, nay millions of martyrs, through the instigation and means of the Bishop and church of Rome, have been executed with most extreme and horrible punishments within these six or five hundred years, histories make mention. What has been done, and what abundance of human blood has been shed even within these thirty years, which our memory reaches, my heart grudges to recite. Great is this crime also, for which Babylon is plagued by God with just and most grievous torments. And aptly here is made mention of martyrs, that is, of the witnesses of Jesus. For those who confess the Evangelical doctrine of Christ to be the true and absolute doctrine, that Christ is the only head of the church, the only priest and bishop, mediator and sacrifice, and will not join with all, while at the same time that the doctrine of the church of Rome is also most perfect, to be held in like reverence with the doctrine of the gospel, that the Pope is head of the church militant, and on earth the true vicar of Christ and general pastor, and that the saints in heaven pray for us, and that the mass is a true and real sacrifice for the sins of the living and the dead—they are condemned as heretics and schismatics, with sword and fire to be rooted out of the earth.

### Paragraph 1080

And thus far extends the mystery of the vision, with terrifying sights presented before our eyes, wherein is vividly described both the ancient Roman Empire and, principally, the Papacy of Rome with its sins and crimes, heinous and full of enormity. Afterward shall follow an exposition of the vision, to which at the last shall be annexed the punishment to be taken of Antichrist and the entire Antichristian City. But if anyone applies all these things which are spoken of the woman to ancient Rome, I will not object. For there was also a religion at Rome, which consisted in gold and precious things. Ancient Rome had a cup of false wisdom, with which she made drunken and infected all nations. She was therefore the mother of abominations and whoredoms, from whom the provinces learned superstitions, and so forth. Nevertheless, these things particularly concern the Pope. Notwithstanding that ancient Rome also was drunken with the blood of Saints, and so forth.

### Paragraph 1081

John was greatly astonished when he saw the woman. Daniel also was astonished, so much so that his heart almost failed him, when he saw that Roman beast, as appears in the seventh chapter of Daniel. Godly people are also astonished today when they see so great things tolerated or permitted by God to the church of Rome against the pure and sincere. For the prelates of the church are fortunate, victorious, powerful, and in favor with all princes, and they bring to pass whatever they imagine or desire. Happy is he who is not offended by them. Read the seventy-third Psalm. How good God is to Israel, to those who are upright in heart, &c.

### Paragraph 1082

The Angel, as chief minister, reveals the mystery to John and the whole Church: that is to say, he opens the secret and true meaning of the vision by parts most diligently. And he speaks in deed of the whole body of the beast, notwithstanding that the beast has certain things peculiar, and likewise the whore. Yet the Angel himself says, “I will show thee the mystery of the woman, and of the beast that bears her.” Nevertheless, this is also a dark speaking, wherein in the beginning of the exposition he says, “The beast which thou sawest was, and is not.” The Roman Empire was still present while Domitian ruled, but it was no more such as it had been before. For from the first monarch Julius, it was as it were by inheritance in the house of Caesars, until Nero. For in him the beast received a deadly wound, but it was healed, and diverse emperors reigned not of any one family. The Empire therefore had been in the power of one house before, but after Nero it was not so. Again, the Romans possessed the Empire after Nero. From Nerva, which is the seventh after Nero, the Empire was devolved to Trajan, under whom it was puissant and strong. Therefore it was and it was not. Whereof John himself will speak a little after. Moreover, the histories testify that the Empire of Rome was extinguished, and in its place sprang up another, which is also called the Roman Empire; whereof you may say most truly it was, and is not. For that old Roman Empire was the most ample and noblest Empire in the world, but this new, now erected by the Pope, is none such, but rather an image of the beast (as I said in the thirteenth chapter), a shadow and a dream. Therefore we doubt nothing but that in this vision is exhibited to us a type, both of the old and new Empire, but chiefly of Papistry.

### Paragraph 1083

And now he demonstrates no other origin of the beast [note 74.10] but a hellish and devilish one. For he says how it shall arise from the bottomless pit, as was spoken of before. All empires, certainly, as Daniel shows in chapter 2, are of God. But if the rulers are corrupted, the origin is attributed to the Devil, not of the empire itself, but of the corruption. Moreover, it is added here what end that unfortunate empire is likely to have at the last: and it goes into perdition. For it is cut up by the roots on earth, and in the next world is condemned to everlasting punishment.

### Paragraph 1084

But just as he demonstrated in the thirteenth chapter, concerning those who should wonder at—that is, honor and worship—the beast, so here he repeats the same, not the elect children of God, but earthly men and those condemned, whose names are not written in the book of life. Of this we have spoken before. He adds an exclamation: and here is a mind that possesses wisdom. The Lord urges all hearers to the diligent consideration of these things, lest we be deceived and perish. Those are fools who marvel at the beast’s prosperity, victories, the Pope’s majesty, riches, and pleasures, and submit themselves to him. Truly wise are those who understand that the empire is taken away, and that under the guise of an empire lurks Antichrist, the child of perdition and man of sin, to be shunned by all the faithful. For these have turned to Christ, in whom they know they have all things pertaining to life and salvation, and to live in him. To him be praise and glory.

## Sermon 75: ¶ The godly vision is yet more plainly declared.

### Paragraph 1085

.

### Paragraph 1086

The seven heads are seven mountains on which the woman sits; they are also seven kings. Five have fallen, one is present, and the other has not yet come. When he comes, he will continue for a time. And the beast that was, and is not, is also the eighth, and one of the seven, and will go into destruction. The ten horns which you saw are ten kings who have not yet received the kingdom, but will receive power as kings for one hour with the beast. These have one mind, and they will give their power and strength to the beast. These will fight against the Lamb, and the Lamb will overcome them. For he is Lord of all lords and King of all kings, and those who are on his side are called chosen and faithful.

### Paragraph 1087

The angelic interpreter proceeds to declare to John the mystery of the beast displayed, and of her judgment, in stages. And at this time, he explains three things: what the seven heads signify; why he said of the beast, “he was and is not”; and what the ten horns represent.

### Paragraph 1088

He explains two approaches. First, by seven mountains, upon which the woman sits [note 75.1], whom at the end of the chapter he calls the great City: namely, the great city, which is commonly called Rome, because it stands upon seven mountains. Moreover, the beast has seven heads, because it has had many times seven kings. Of this I spoke also in Chapter 13. At this point, he counts the seven kings, and there is no doubt that he is speaking of Rome. This, I suppose, is the Lord’s chief intent in these matters. For he could not speak more explicitly unless he had expressed the name of Rome itself; but the name of Babylon we heard expressed before. Five, he says, have fallen, namely, since the deadly wound was inflicted, in the death of Nero, within fourteen years. For immediately after Nero, Galba began to reign; who being slain, Otto reigned; which after killing himself, Vitellius succeeded, and he also was killed by the Flavians. For after him, Flavius Vespasian was emperor; after whom Titus, the best prince of all. And these five fell within fourteen years. He adds, and one of those is reigning now, namely, the sixth in order, Domitian, the son of Vespasian and brother to Titus, a most ungracious man, who persecuted the faithful and had condemned John the Apostle into exile. Another, says John, has not yet come: namely, Cocceius Nerva. For after he came to the empire, and lived most virtuously and righteously governed the empire, he did not remain long. For when he had reigned one year, three months, and nine days, he died. And thus much hitherto concerning the seven kings and the seven heads of the beast. These things, so certain, pertain not so much to the exposition of this passage as to the consolation of the faithful: who here may clearly perceive how empires consist in the hand and providence of God almighty, who knows his own and has care of the godly, although they may seem, by reason of their grievous persecutions and cruel torments, to be neglected by God.

### Paragraph 1089

Consequently, he explains why he said the beast was, and is not: namely, because he was the eighth Roman emperor, Ulpius Trajan. For he was the eighth from the empire wounded by Nero. Trajan was of the seventh, that is to say, was adopted by Nerva the seventh emperor. And until this point, the Roman Empire had been governed, first in deed by Caesars, then by the noblest citizens of Rome. But concerning this Trajan, who succeeded Nerva, the writers of histories say that he was the first foreigner to rule the empire, for he was a Spaniard. Therefore, the empire was or has been in the hands of the Romans, now it is so no more. For a Spaniard succeeds, so that the empire now seems as though it might be called Roman Spanish. And inasmuch as Trajan persecuted Christ and his members, he also went into perdition. And let no man think that this was the only and sole cause wherefore John said, concerning Trajan, that he was, and is not. For he has pronounced expressly, and he is the eighth, as though he should signify that there be other causes also, for the which it was said that the Roman Empire was, and now is not, whereof is spoken before.

### Paragraph 1090

Hereafter follows also the exposition of the ten horns. And the same horns are here recited, which are spoken of in the seventh of Daniel, and in the thirteenth of the Apocalypse. Neither is there any cause why one should superstitiously cling to the tenth number. For in the fourteenth of Numbers, the Lord says how he has been tempted ten times by the Israelites, for many times. Here is signified therefore, how the Roman Empire shall be dispersed into many kingdoms. For whether you say kings or kingdoms, the matter is all one. Doubtless the Roman Empire beginning to fall to decay, there sprang up kings in the East and West, which invaded the Roman Empire, Persians, Goths, Vandals, Lombards, and I know not what others; at the last in Spain, France, Hungary, I speak not of Africa and Asia, were found diverse kings, and the Roman monarchy ceased. Of these kings the Angel warns us for diverse causes. These, says he, have not yet received the kingdom. For whilst John wrote the Apocalypse, Domitian ruled, and the Roman Empire was yet mighty and strong, and so remained for certain ages. When therefore did they receive their kingdom? They receive, says he, power as kings at one hour with the beast, namely the second. For these things cannot be understood of the first and old Roman Empire. And Primasius, expounding this place, admonishes that an hour here is taken for a time present. Therefore at the same time, the beast, that new Empire grows up and increases, and the kings receive might and power. For the decay of the old Empire was the strength of kings, and of the new Papal Empire. And in deed the emperor Phocas commanded the church of Rome, and the Bishop thereof to be head of churches. Which gave a certain beginning to the Popes dominion, as also in the thirteenth chapter. I have recited: which he obtained at the length more fully under king Pippin and other Princes of France and Germany, but Nüchelus speaking of the Empire of Phocas in the twenty-first Generation, says the enemies of the Roman Empire, by the slothfulness and cowardice of Emperors, had taken away in the West country with the Islands of Germany, France, Spain, Hungary, Slavonia, and a good part of Italy, and thereto a great part of Africa; and in the East parts, Cacannus of Thrace, King of Huns invaded the Iberians, Armenians, Arabians, Dardanes, and the middle parts of Macedonie and Greece. And the Persians in a manner possessed all Assyria, the Saracens destroyed Egypt; fie for shame, our strength has so failed us through riot, covetousness, and voluptuousness, that the Roman Empire stood then only in name. Hitherto he. The same things have we discoursed more at large in the thirteenth chapter of this work. And verily Daniel shows how among those ten horns, one other little horn should grow up, which should strike off three, and take their place, and reign wantonly, cruelly, and wickedly. Wherefore the Papal Empire, and those sundry kingdoms grew up in a manner about one and the same time.

### Paragraph 1091

He also shows what kind of kingdoms those will be, and how they will behave toward that final beast, namely, the church of Rome: he says they all have one opinion; they believe one thing, and are of the same religion. He speaks chiefly of the western kings. For they all receive the decrees of the Bishop of Rome, and honor them, as most obedient children of the most sacred and holy church of Rome. They will deliver to the beast their power, their authority, or kingdom, for they submit themselves to the See of Rome. If the church of Rome has need of an army or force of arms, the kings send their power gladly to him, which the noble kingdom of Bohemia felt about a hundred years since, though it were to no great advantage and beautiful triumphs of the invaders. Yes, moreover, they acknowledge themselves to owe homage and fealty to the most holy and supreme Bishop in all the world. Hereunto chiefly pertains what Augustine, as recorded by Steuchus in his book against Laurence Valla, concerning the Donation of Constantine, in section 94, has written in this way: Gregory the Seventh to Géza, king of Hungary: “We suppose it is not unknown to you that the kingdom of Hungary, like as other most noble realms also, ought to be in the state of its own liberty, neither that it ought to be subject to any king of another realm, save to the holy and universal mother church of Rome, which has her subjects, not as servants, but as children.” Hereunto Steuchus adds: “Hear with what government the church rules, that she may sustain her subjects, not as servants, but as children. She does not displace kings from their possessions, but permits them to reign as her sons; while they are reigning, she reigns herself also.” Nevertheless, she will be known as Queen and Lady.

You hear how all the most noble realms are subject to the Apostolic See. Even there he shows that the most noble kingdoms of Spain, France, England, Denmark, Russia, Croatia, Dalmatia, Aragon, Sardinia, Portugal, Bohemia, Swabia, and Norway are subject and tributaries to the church of Rome. In section 97, he adds moreover: although the kings reigned and continued in possession, they are wont to acknowledge her as Queen, true Lady, and giver of their kingdoms. And in section 105, the old records of all Popes are full of high authority, whereby they have, with their empires, governed the whole world, having the rule and order of all lands, which power and authority that impudent praiser of the Roman See is not ashamed to call omnipotent or almighty.

### Paragraph 1092

And undoubtedly we see today, great embassies sent to Rome by the newly elected and crowned western kings, intending to pay homage to the Pope, whom they call the head of Antichrist, and to offer what they term due obedience. Therefore, he called them before not kings absolutely, but as kings. For they acknowledge a superior, and are even as it were servants or wards of the servant of servants. Of whom he has made verses: “The common people brought from the far ends of the world, The servant of servants, O Rome, you are now their lord.”

### Paragraph 1093

Hereunto the Apostle adds a thing yet more grievous.[note 75.9] These kings, I mean the confederates of the Pope, and obedient children of the Church of Rome, endowed with the spirit of the beast, shall fight with the Lamb. Whereby is signified the tyranny which kings, and princes, and certain other states of the Roman Empire do practice, and long have practiced against Christ and his gospel. Concerning the Lamb we have already spoken enough before. John the Baptist, pointing with his finger to Christ, says: behold the Lamb of God, which takes away the sins of the world. Therefore shall the Roman princes fight, not against Christ himself, for they will be Christians, but against the Lamb, that is, the sanctification, justification and satisfaction of Christ. For if any man say at this day, that the son of God is most holy, by whom alone sins are forgiven, and we are sanctified: and say not also, that the Bishop of Rome is most holy also, which purges by pardons granted, but shall say rather, that pardons are plain deceitfulness, and the Pope most unclean of all: he shall doubtless neither be taken for a right Catholic, neither shall he be spared for confessing the Lamb of God. If any man shall confess that justification is only in the son of God alone, and that men are justified by faith only, and not also by our works and merits, he shall be carried to death or to prison, neither shall the confession of the Lamb of God avail him anything. If any man shall say, that he is fully purged through the only oblation of Christ on the cross, as of a lamb without spot, and sacrificed from the beginning, neither that he needs any popish Masses, whereby the priests boast that they make a daily offering for the sins of the quick and dead, which in deed is both false and blasphemous, he is straightways hurried to prison, and from thence drawn to the stake and burned. We cannot deny but that this is true, saying there are at this day innumerable examples of Romish kings and princes in this behalf. We shall not need therefore to fetch our exposition far off, how these kings, which wholly depend on the Pope, shall fight with the Lamb. I speak here nothing of others, which cleave whole unto Christ.

### Paragraph 1094

And therefore, for comfort, what follows is joined to this: the Lamb will overcome them. [note 75.10] For although Papist kings and princes seem to overcome the saints, whom they burn, murder, and destroy, yet Christ lives forever; the redemption of Christ flourishes. As a very pious poet has sung: Christ lives still, and always will; his truth also will remain, while all that fulfills this world shall perish and be worthless.

### Paragraph 1095

Kings perish, kingdoms perish or are changed; but the truth is never changed, Christ perishes never. He adds a most strong reason: for he is Lord of Lords, and king of kings. Therefore shall they be made a footstool for the feet of the Lamb, as many as shall strive against him. You see again why Saint John said before: they receive power as kings. For all kings are under Christ, who excels all lords in the world. For to him is given power in Heaven and in earth. Let us therefore be of bold courage. For the Lord is Emperor, and our king almighty, immortal, and invincible. He will come shortly in the clouds of fire, to judge the quick and the dead, &c.

### Paragraph 1096

Moreover, victory is assuredly promised to us who are servants of Christ.[note 75.11] And those who are with him or on his side are called, chosen, and faithful. We are chosen in Christ before the foundations of the world were laid, so that we should believe in him and be saved, first, to the Ephesians. To this we are called by the preaching of the gospel. Read 2 to the Thessalonians, chapter 2. And we ought to give thanks to God forever, and so forth. Let us hold fast to these things and be steadfast in the troubles of this world, and without fear. To God be glory.

## Sermon 76: ¶ Again this vision is more fully declared, and the punishment of the beast is showed.

### Paragraph 1097

.

### Paragraph 1098

And he said to me: The waters which you saw, where the prostitute sits, are peoples and nations, and tribes, and languages. And the horns which you saw on the beast are those who will hate the prostitute and will make her desolate and naked, and will devour her flesh, and burn her with fire. For God has put in their hearts to fulfill his will, and to do with one accord, to give her kingdom to the beast, until the words of God are fulfilled. And the woman which you saw is that great city which reigns over the kings of the earth.

### Paragraph 1099

The angel sent by the Lord Christ reveals to John and the entire world the mystery of the beast, but especially her judgment or punishment for her heinous crimes. He will pursue this also in the following chapter.

### Paragraph 1100

And now he explains the meaning of the waters, over which the prostitute holds sway, namely the Roman power. Waters signify kingdoms dispersed throughout the world. He explains this in his own prophetic manner, as was noted before, by three words. For in naming people, folk, nations, and tongues, he encompasses as it were countless nations, distinct with diverse languages and customs. But where nothing is more changeable or unstable than waters, and when they are stirred up, they become more furious and outrageous: the common people are rightly compared to waters, which are also called changeable or unstable because of their fickleness, and furious and mad because of their rage.

### Paragraph 1101

Therefore, not without cause, all wise men have severely condemned seditions, which we are accustomed to call tumults or uprisings. For through such disturbances, many wicked individuals gather, and opportunity is given them to act on their desires and to inflict harm. But if many nations were subject to the Roman Empire and nevertheless strayed from the true faith, what will it avail to count up many kingdoms that might agree on a religion? As though the sincerity and truth of religion depended on a multitude of people agreeing in the same.

### Paragraph 1102

Now follows God’s judgment against bloody Rome, which is the chiefest thing in this vision: the sum of it is, Rome shall be rent in pieces and burnt with fire, as we heard also in chapter 13. Just as she has done to others, so shall be done unto her. And these things are to be explained first of old Rome, and afterward of new; and in the same, the words must first be considered, than a conference of histories must be had, out of which the truth of the prophecy may appear.

### Paragraph 1103

Ten horns signify kings who arose from the division of the Roman Empire, such as the kings of the West Goths, East Goths, Germans, French, Lombards, Huns, Vandals, and so on. These nations served the Romans and received payment from them, favored them, and to their own detriment carried out their business, no differently than an earnest lover serves a harlot, from whom he cannot be withdrawn, whom he loves most fervently, but at last, perceiving her deceitful behavior, he begins to hate her most intensely. So these nations and others began to persecute the name of the Romans, desiring that no monuments or traces of them remain anywhere. All the provinces of the Romans were filled with Roman statues, images, pillars, inscriptions, and writings; but in those places, especially in Germany and its borders, how many of such great abundance remain? The cities where the Romans had their garrisons are utterly destroyed, so that scarcely any traces of them appear today.

### Paragraph 1104

And just as an honest man, having a prostitute as his wife, a shameless strumpet, does not only hate but also forsakes her, troubles her, and strips her bare, having plucked from her all her wifely apparel and ornaments [note 76.5] (for so God in his prophets threatens to do to his people for their unfaithfulness), so nations revolted from the Roman Empire, destroyed and impoverished it, spoiling the riches thereof, which the Romans had heaped together by the robbery of all nations [note 76.6]; they spoiled everywhere also the Roman provinces. And where it is said that those kings shall devour the flesh of the beast, it is to be understood in the manner of speaking. For so we are wont to say, when we signify extreme cruelty and malice without mercy; therefore, just as Rome has been most cruel towards all nations, even so shall all nations most cruelly tear her, and finally shall burn her with fire.

### Paragraph 1105

Now let us compare these things with history, and see how they were fulfilled in old Rome, and may be yet fulfilled in the new. And first we will speak of old Rome, and after of new. And truly, old Rome grew many years, and practiced robberies throughout the whole world, and destroyed the saints of the Most High; wherefore it was worthy that the punishment thereof should extend and endure many years, and so as it were by degrees to descend to the last burning and destruction thereof. There are gathered the years of her punishments about 136, in which she, being impenitent, was vexed and tormented with continual calamities, slaughters, and vexations. And hereof I compiled an abridgement in the 57th Sermon of this work, the 13th Chapter. And here I will repeat a few things, and will rehearse certain other things more plainly and at large. As the Lord in punishing the Ninevites and Jerusalemites, declared his long suffering and clemency, and also his straight justice; right so in proceeding slowly to destroy Rome, he left them mercifully space to repent. But since they refused to do so, he wasted and destroyed them terribly as impenitent. He gave therefore to Rome excellent good rulers, Constantine, Iovian, Valentinian, Gratian, Theodosius, and so forth, by whose diligent labor and godliness he disclosed the furies and raging idolatry of the heathens, and also restored and established the true religion. But as in the time of Josiah the old, ingrained error and abominable idolatry could not be rooted out of their hearts, but that the greater part would rather keep the abominations of the Amorrites; so the Romans both in the city and in the provinces eagerly aspired to the restoration of old idolatry. Therefore, just as he tamed at length with grievous wars the invincible ungodliness of the Jews, and destroyed the city of Jerusalem, so by the war of the Goths and Vandals, and invasions of barbarian nations (as the histories term them), he destroyed proud and wicked Rome, with her provinces, and finally consumed the city with the sword and fire of the Goths. The very name of the enemy cried out that the vengeance was not executed by men, but by God himself. For the Germanic word for Goths signifies the people of God, or gods’ people. For God in Old High German is called *Got*; from that comes *Gotthes*, *Die Gotther*, the people of God. Therefore God, and not man, did chasten, torment, and at last destroy Rome. Which thing John the Apostle speaks most expressly at this present.

### Paragraph 1106

First, during the rule of Honorius and Arcadius, the Visigoths, through Alaric, besieged, attacked, captured, and plundered the city. Jerome to Principia greatly laments this misfortune of Rome in the Epitaph of Marcella; but Orosius, as I also mentioned in sermon 57, in my judgment more rightly commends the just judgment of God in the affliction of Rome. It is clear that Rome was then, because of the enormity of her sins, chastised with mercy, but where the Romans would not acknowledge the hand of the one who struck her, it came to pass that Alaric, being dead, the victorious army, now led by Adolphe, returned from Lucania and plundered the riches of Rome even more greedily than before.

### Paragraph 1107

From that time, Rome was again granted a period of repentance, about forty-two years. [note 76.10] In the meantime, through destruction and devastation brought by the Huns in their provinces, and that great and wondrous event, they were admonished to be wise. What will they say who observe that Attila himself, with his Huns, now invades Italy itself, and now hangs over the neck of Rome? Then there occurred an event that could have turned the Romans to the service of the true God, had there remained in them even a spark of gratitude. For the minister of the church of Rome, Bishop Leo (the ambitious pride of Popes was not yet known), a preacher of the Christian faith and a steward of Christ's mysteries, making supplication to Attila, obtained peace for Rome, and by a manifest oration turned away the bloody enemy from the necks of the Romans. This was an exceedingly great benefit, which God by his servant would have shown to the Romans, if they would yet cease to hate the religion of Christ and to slander Christ, as though he poured out evils into the world, and that no good nor quietness came from the preaching of the gospel. For even now (not to speak of others innumerable), he has implored upon them a benefit inestimable, and that by the preacher of the gospel. This was done in the year of our Lord 454.

### Paragraph 1108

However, while the Romans proceeded according to their accustomed manner, [note 76.11] and now also Valentinian, a prince not evil, was murdered, and by a riot many dishonorable things were done, nor did any sign of thankfulness toward Christ appear, or sign of true conversion. Through the influence of one woman, Eudoxia, the wife of Valentinian, who herself also suffered many dishonorable things in that riot, it came to pass that Genseric, king of the Vandals, sailed out of Africa with three hundred thousand men to Rome, and captured it, and for the space of forty days, he gathered up the treasures brought there from all parts of the inhabited world. Then the intercession of Leo could do nothing, except that the Vandals abstained from killing and burning: which was also a benefit of God not to be scorned. The first king of the West Goths who broke into Rome was called Alaric; others call him Atalaric; but this king of the Vandals is named Genseric, and so Rome, a harlot, is made desolate and naked, plundered, I say, which, having been enriched with the spoils of all nations, was hitherto proud. However, it was not now altogether destroyed and burned: which was no small benefit, which Christ again showed to Rome for an amendment.

### Paragraph 1109

And yet moreover, about twenty years have passed, during which time, nevertheless, as with the ten tribes of Israel before the destruction of Samaria, continual murders were committed while princes reigned in Rome. Yet, for all that, none of these rulers escaped being either slain, murdered, or expelled. Augustulus was the last. For just as Augustulus succeeded Julius and initiated the Roman monarchy, so Augustulus brought it to an end.

### Paragraph 1110

[note 76.13]For the Roman legions being extinguished, and the imperial title taken by the Germanic people, Odacer, (who took his name from destroying lands, Oedacher, and was called as if a destroyer,) captured Rome, and reigned there as king for about fifteen years. Yet he was expelled again, and slain (at the instigation of Zenon, Emperor of Constantinople) by Theodoric, Prince of the Ostrogoths. [note 76.14] And the Ostrogoths reigned at Rome for about fifty years, until the Emperor Justinian sent Belisarius into Italy with a Greek army to recover it. The Ostrogoths, being aided by a force of Germans sent to them by Theodeper, king of France, valiantly resisted. They warred in Italy for eighteen years, with fluctuating fortune.

### Paragraph 1111

[note 76.15]At the last, Totila overcame Baldeuille. He seized and burned the city of Rome, but not suddenly. For he gave a time to deliberate. But where he could not so prevail, he destroyed Rome, and as John has prophesied, burned her with fire. All histories mention this destruction. John Augustine in the third book of Chronicles has written of the same matter: Totila besieged Rome and took it, on the sixteenth day of January (the seventeenth of December) in the year of Christian salvation 548. Totila gave all the goods to the soldiers, but he commanded by proclamation that their bodies should be free. From there, he sent ambassadors with his wishes to New Rome (Constantinople) to Justinian. He demanded of the emperor Italy and the league as it had been under the emperors Anastasius and Theoderic, king. Which if he could not obtain, Totila threatened that he would raze the city, which he could not keep, and abolish the Roman name. Justinian answered that Belisarius was in Italy, to whom he had committed Italian matters. Totila, therefore, where the emperor would not grant his requests, determined to raze the city of Rome. He leveled the greatest part of the walls in most places and set the Capitol building on fire. He commanded all citizens with their wives and children to depart from the city. The common people of Rome were dispersed in the towns of Campania. The senators and nobility Totila kept with him as hostages. Then fire was put into every house. Thus, Rome being fired in all places, Totila left it vacant. For thirteen days, the fire burned freely. The city of Rome was forty days in that solitude, that there was neither man nor woman in it. After destroying the city, he moved his camp toward Lucania and Calabria. Belisarius came to the city left vacant, and sooner than one would have thought, fortified a part of the city with a ditch, wall, rampart, and wooden towers, for all could not be restored. Totila was with him, but repulsed, departed to Tibur. Belisarius was sent for into Greece by the emperor. Totila besieged Rome and took it. So in one year, Rome, the head of the world, the lady of all nations, was taken three times, thus writes Augustine. Leonard Aretine, writing of the Italian war against the Goths, in the end of the second book, says: After this, Totila departing from Rome with his whole army, left it utterly desolate and vacant, and so on. Who will say now that John has not in few words comprehended the destruction of old Rome, which the histories afterward have plentifully described: and finally, how, in the same manner as it was prophesied, it has followed the prophecy, after 451 years? And that so evidently, that you would think presently to behold Rome both falling and burning.

### Paragraph 1112

And just as in the story of the gospel the Lord intermixes a prophecy of the destruction of Jerusalem and of the end of the world, so that every man might, seeing the city of Jerusalem perish as he had prophesied, with no stone remaining upon another, gather with like truth and certainty that this world shall fall: so may we, seeing the once great city of Rome fallen and so great an empire, which was thought should have lasted forever, brought to naught, gather also that new Rome with her shadow or image of the empire shall as sure as day fall and be brought to naught. And first indeed the Saracens and Turks, who ruled and yet reign in the provinces subject to the Roman Empire, as in Asia, Greece, Egypt, Africa, Slavonia, and Hungary, and therefore are rightly accounted among the ten horns, do hate worse than dog or snake both Popery itself and Rome, and all that imagery of empire. Yea, and histories also testify that they have often made invasions and spoiled Rome itself. What is done at this day, experience itself teaches. But whether the Turk or the Christian princes themselves, converted to Christ by the gospel, shall spoil this new Rome, destroy it utterly and burn it with fire, the Lord knows, who seems here to intimate some such thing hereof. This is certain that Christ alone with his hand shall bring down Antichrist and abolish him with his coming. Certain it is that the Earth and all the works that are therein shall be broken. For thus is the apostolic doctrine: and that all these things shall be in the end of the world. Read Paul, 2 to the Thessalonians, chapter 2, and Peter, 2 Epistle, chapter 3. Moreover, there arise in sundry kingdoms of the world learned men, who, once bounden to the See of Rome, have defended her and her stinking idols; but after being converted to Christ, begin to hate both Rome and the Roman church, which also they assail and burn with the fire of God’s word. Therefore all the glory, dignity, and wealth of the Pope and Popery has perished and perishes daily in the godly. All that are godly and wise hate Rome and Roman wares. All cry out that this Sodom is worthy to be burned with fire falling from heaven. Neither is there any doubt but that a grievous vengeance is prepared against her.

### Paragraph 1113

And briefly is shown a reason wherefore the kings should rage so cruelly against the beast, and why these things are done in such sort and manner as we have heard. For God says he has given into the hearts of the kings, that is, that they should work his will and should do with one mind and consent. For where some refer this to the beast, that seems far off and strange. It is referred rather to the next, to God I mean, who put into the hearts of the kings to do his will, I say of God. For it is the mind and will of God that the beast should perish, that punishment be taken from her for shedding of innocent blood. The same God will procure that kings shall not be at discord, but at concord, that being of one mind and accord, they may execute God’s judgment. So we read in the Prophets that God put into the hearts of kings, Shalmaneser, Sennacherib, Nebuchadnezzar, Cyrus and others, that they should do as they are said to have done, namely in punishing the wicked and defending the godly. And there is also mention made in Histories how Alaric, king of the West Goths, was indeed dissuaded by a servant of God that he should not make such haste to destroy Rome, but that he answered, “There is one that continually troubles me and says, ‘Go, destroy Rome.’” And he that put that mind and will into the heart of Alaric, Adolphe, Genseric, Odacer, Theodoric, and Totila, the same if he will, and when he will, and into what princes he will, shall put, that they also shall do their duty against this new city and church of Rome.

### Paragraph 1114

The angel adds that God has also put into the hearts of kings, so that they should give their kingdoms to the beast, until the words of God are fulfilled. The interpreters explain and say that God has permitted them to conceive this counsel in their minds, to deliver the kingdom to the beast. But I believe it is clearer if we simply confess that God is the author of no sin, and that men sin not as compelled by any fatal necessity, but through their own fault and vice. Therefore, God would, as he has also expressed and taught in his word, that kings should deliver their kingdoms to Christ, the supreme King. But since it did not please them, and they preferred, for various worldly and fleshly reasons, to deliver their kingdoms to the Pope, and submit themselves to the so-called Apostolic See, God, in his just judgment, has forsaken them and given them over (as Saint Paul wrote to the Romans) into a reprobate mind, to do those things which God does not allow. And so the words of God, prophesied by the Prophets and Apostles, are fulfilled in this way. Surely these are the words of God, and not of men, which are read concerning this matter in Daniel and throughout this book of Revelation.

### Paragraph 1115

Finally, the angel explains what is signified by the woman sitting on the beast: namely, that great city of Rome, the head and leading power of the world, and the Roman church, Papacy, and authority, extending herself and her dominion over the kings of the earth. Of this we have already spoken sufficiently. To God be glory.

## Sermon 77: ¶ He shows that Rome shall assuredly fall: and adds the causes of her fall.

### Paragraph 1116

.

### Paragraph 1117

And after that, I saw an angel come down from Heaven having great power, and the Earth was lightened with his brightness. And he cried mightily with a strong voice, saying, “She has fallen, she has fallen, even great Babylon, and has become the habitation of devils, and the hold of all unclean spirits, and a cage of unclean and hateful birds. For all nations have drunk of the wine of the wrath of her whoredom. And the kings of the earth have committed fornication with her, and her merchants have become rich from the abundance of her pleasures.”

### Paragraph 1118

He pours throughout the 18th chapter the destruction of old and new Rome, also of paganism and Antichristianism, and that with a marvelous abundance and clarity of speech, even so that one would think that you saw all things presently. And he uses also a most godly order. For first, the angel declares the destruction of Rome with most fitting words. Secondly, counsel is given to the godly, how to behave themselves in so great dangers. Then is added the manner of the desolation, that just as Rome has greedily and cruelly spoiled and destroyed other nations, even so it shall befall her also. After this, a lamentation is made, wherein the princes and merchants mourn for the ruin of Rome, where they also recount the riches and pleasures of Rome. Finally, the Apostles and Prophets rejoice at the most just judgment of God. Again, the angel of the Lord cast a millstone into the bottom of the sea, that so the most certain, unrecoverable, and most weighty destruction of Rome might be signified. Whereunto again are annexed the causes of so great evils, and the same finished with the praise and congratulation of all the heavenly dwellers.

### Paragraph 1119

And it is fortunate that he imitates the holy prophets of God, of whom two in a manner follow the same pattern,[note 77.3] describing the destruction of old Babylon. Isaiah in chapters 13 and 21, and Jeremiah in chapters 50 and 51. And Ezekiel depicts the overthrow of Tyre in chapters 26, 27, and 28. For just as the fate and end of all the ungodly are alike, so the canonical Scripture, in portraying their destruction, agrees well with itself. The Apostles moreover,[note 77.4] although they spoke and wrote to the Gentiles in Greek, altered nothing of their natural way of speaking, and even compelled foreign languages to serve the holy, and not the Hebrew to serve pagan languages. For when speaking Greek, they observed the natural phrasing of the Hebrew speech, considering it first, divine, and holy. And though they could speak all languages, they never spoke or wrote any foreign language so that the Hebrew phrasing could not be perceived.[note 77.5] Let some therefore beware at this day, that they are not overly fastidious, and follow the purity of the Latin speech so that, in expressing it, they do not fall away from the simplicity of the holy tongue and lose some mysteries. Those who are not stubborn would rather conform themselves to the holy language and learn its phrases than force it against the ear to foreign languages and compel it to serve our delicate ears. Moreover, we have already admonished often what the end and purpose of this treatise is, concerning God’s judgments or punishments. For the truth and justice of God are confirmed, the afflicted receive comfort, and the wicked, and all God’s enemies are made afraid, and so forth.

### Paragraph 1120

[note 77.6]But when John published these things and prophesied of the destruction of Babylon, all men at that time clearly understood it to signify Rome, because of the recent destruction of Jerusalem and the grievous captivity of the Jews, which had recently occurred under Vespasian. For just as Babylon in times past had vexed the holy city and nation, so now Vespasian, the Roman, did. The godly truly believed these things to be true and that they would undoubtedly come to pass. The ungodly, however, mocked them as foolishness. Their elders had done the same. For when the Prophets also prophesied the destruction of Nineveh, Babylon, and mighty monarchies, they seemed to them to be mad. Nevertheless, just as they had said, so it came to pass. Therefore, the faithful believe God’s oracles, however long they may be deferred, which are prophesied to come, no matter how impossible they may appear to the world. For to God who speaks and wills, nothing is hard.

### Paragraph 1121

[note 77.7]And intending to demonstrate the downfall of Rome, he prepares his audience and gains credibility for the prophecy, while first of all he shows the author of the oracle or prophecy, the very Angel of God. Indeed, he highly commends this Angel to us, so that we should doubt nothing of the truth of those things which he speaks. For he says how he came from Heaven. From this we gather that those things he brings are divine and celestial, and are said to have great power, lest we should think impossible those things which he says will come to pass. For if God’s Angel is a minister of such power, what may we think the Lord to be, who sent the Angel? One Angel before the walls of Jerusalem killed one hundred eighty-five thousand men of war. One Angel in a night slew all the firstborn of Egypt. Therefore, when the most mighty Angel prophesies the destruction of old and new Rome, we need not doubt that it will utterly perish. Moreover, the Earth was lightened with the glory, that is to say, with the brightness or light of this Angel. For this prophecy is neither obscure nor will it be hidden, but chiefly and most clearly preached throughout the world.

### Paragraph 1122

Wherefore the same angel cries with all his force. For it behooves these oracles of God, wherein is treated of the glory of God and the salvation of souls, to be preached with loud voices, however the world prohibits and persecutes them. And let those who think that men may be restrained by proclamations, fire, and sword observe that they shall not, with the clearest voice, preach against Antichrist. These fools are deceived. They have fought and contended in this matter for six hundred years and more, and neither could any man, though he raged ever so fiercely, bring this preaching to silence. It breaks out repeatedly and pierces far even at this day throughout the whole world; therefore the glory of this angel is yet, and ever shall be shining and bright, and his voice and preaching most strong, though the Popes [? guts burst].

### Paragraph 1123

Now follows the prophecy of Thaungel, [note 77.9], some whereof is: Rome shall perish, and no trace of her will remain. He utters this prophetically, as he did also in the 14th Chapter. She is fallen, she is fallen, great Babylon. He said, “She is fallen,” for she shall fall: putting the time past, for the certainty of the thing, for the time to come, whereunto the doubling also appertains. Likewise spake the Prophets. Macrobius marvels at the wonderful brevity of Virgil. And among other things in the first Chapter of the 5th book of Saturnales, will you hear Virgil, says he, speaking with so much brevity, that brevity itself can be no more straightly hampered and drawn together? And fields where Troy was behold how in very few words he has consumed a mighty great city: and has left no ruin at all. Hitherto Macrobius. These things we shall more truly and more rightly apply unto our Prophets, most eloquent in their tongue, and chiefly to Jerome. For what could be thought more brief than that which he said, “She is fallen, she is fallen, great Babylon?” For Jerome both expressed the greatness and majesty of the city, and swallowed it up whole, no ruin at all left, for he signified that both old and new Rome, although it seem stout, invincible, and eternal, yet shall it fall: and so fall, that nothing thereof shall be left. Which shortly after he sets before our eyes more expressly by a certain sign, whilst the Angel taking up a millstone, and casting it into the bottom of the Sea, adds: thus or with such a violence shall Babylon that great city be overthrown, and shall be found no more. Therefore was there never any thing, is, or shall be in the world so mighty or impregnable, which the invincible power of God cannot bring to naught, when he will, and when the fatal hour is come. Old Rome is lost, and that mighty monarchy decayed: there is fall also the superstition and idolatry of the heathens, that has reigned many years: new Rome shall perish also with her Imaginary Empire: the kingdom also of the Pope or Antichrist which has long sotted and plagued the world shall fall, and fade with smoke.

### Paragraph 1124

Moreover, by a figurative speech taken from the prophets,[note 77.10] he shows the manner of the destruction by consequences: and is become the habitation of devils, and so forth. For so he signifies that it shall be destroyed, that the place which was before much frequented by men, shall now be the habitation of wild beasts and devils, delighting in wilderness, as our Lord also testifies in Matthew 12. And he alluded to the words of the Prophets. Isaiah in chapter 13: Babylon, the beauty of realms, shall be overthrown, as the Lord subverted Sodom and Gomorrah; it shall not be inhabited, but beasts shall there take their rest, and satyrs or harts shall there leap. The same things are repeated also in Jeremiah 50. And in chapter 51, he says: Babylon shall be in heaps, and an habitation for dragons, a water and a hissing, that no man may dwell there. Not unlike things are read in Ezekiel 26 concerning the subversion of Tyre. And that old Rome was destroyed, I showed before: and for the space of forty days and more, inhabited by no man. And that we see it inhabited again, it does not invalidate Christ’s prophecy. For France, Petrarch, an Italian and among the best learned Italians, in a certain epistle to a friend, expounding these words of the Apostle John, among other things, says, “You are verily become such already, for how much better is a wicked man, and of desperate doings, than a devil? Verily you are become the habitation, or rather kingdom of devils: which by their crafts, albeit in human shape, reign in you, and so forth.” Petrarch lived and wrote these things about two hundred years ago. And in another certain epistle, speaking of old and new Babylon, he says, “She was, of all others, worst, and at that time most filthy; and this now is no city, but an house of fiends and sprites, and to be short, the sink of all sin and shame, and that hell of the living, signified long before by the mouth of David, than it was founded or known.” And the same again: whatever you have read of Babylon in Assyria or Egypt, whatever you have read of the four Labyrinths or Mazes, finally whatever you have read of the way to hell, of the dark words there and lakes of fire and brimstone, compared to this hell, it is a fable: here is that proud and terrible Nimrod; here is Semiramis with her quiver; here is unmerciful Minos; here is Rhadamanthus; here is Cerberus devouring all things; here is Pasiphae put to the bull, among a cruel kind, as Virgil says, a young of double shape, Minotaurus by name, a monstrous progeny of unlawful lust. Finally, here you may see whatever is confused, whatever is black, whatever is or may be feigned horrible and ugly, and so forth. These things he has written, and many other more like these, in other epistles. But what do you think he would write now, if he saw the court of Rome at this day? Which is doubtless in many ways more corrupt than it was then. Briefly, John signifies after the sentence of Christ, king and judge, that Rome, both old and new, together with paganism and Antichristianism, shall perish utterly, and never be restored again.

### Paragraph 1125

The causes I have rehearsed once or twice, he repeats and emphasizes: ungodliness, idolatry, and the deception of all people and nations, whom they have compelled by torments to receive impiety. Where cruelty, tyranny, and bloodshed also hold sway. I spoke of the wine of whoredom before in chapter 14 and elsewhere, so it needs not to be repeated again with tediousness. And hereunto is added another new cause, and the merchants of the Earth who grew rich from the power or influence of her pleasures. And he said, of the power of pleasure: for through immeasurable, mighty, and insatiable lust they grew rich. For as Rome abounded with spoils, which it had greedily taken from all nations and brought to Rome, they were given to all kinds of riot and wantonness. Therefore the masters of luxury, the dividers of delicate pleasures, and the merchants of the most precious wares, repairing thither, always found that they would be entertained and honored, and were made rich by the luxurious and riotous life of the Romans. Therefore the Apostle notes an incredible pursuit of most sumptuous riot, in food, drink, apparel, building, in pampering and cherishing of the body. The Roman church also of our time, struck with the same rage both in Italy and elsewhere, accumulates excessive riches through riotous living. This is seen chiefly in those spiritual fathers, bishops, and abbots, and in the entire Roman clergy. But God never allowed riot and tyranny to go unpunished in any nation. Therefore Babylon has fallen, and therefore the church of Rome will fall too. Therefore let private men also love temperance, and abstain from riot and pride. To the Lord be glory.

## Sermon 78: ¶ Counsel is given to the godly, which are commanded to go out of Babylon. Enemies are stirred up against Babylon, and they are commanded not to spare her.

### Paragraph 1126

.

### Paragraph 1127

And I heard another voice from Heaven saying, “Come away from her, my people, that you be not partakers of her sins, lest you receive of her plagues. For her sins have gone up to Heaven, and the Lord has remembered her wickedness. Give her reward even as she rewarded you, and give her double according to her works. And pour into her double in the same cup which she filled unto you. And as much as she glorified herself and lived wantonly, so much power she will have in punishment, and sorrow, for she says in her heart, ‘I sit being a queen, and am no widow, and shall see no sorrow.’ Therefore her plagues shall come in one day, death and sorrow, and hunger, and she shall be burned with fire, for strong is the Lord God who shall judge her.”

### Paragraph 1128

The second point of this chapter is the faithful counsel of the Lord, given to the godly [note 78.1], concerning how they should behave during the prosperity and destruction of the city. Rome has indeed for a long time been mistress of the world; the riches and pleasures of the entire world have been seen at Rome. If anyone at Rome or in the provinces showed himself compliant and obedient to the Romans, and loved the Roman religion greatly, and conformed to the corrupt manners of the Romans, he was highly esteemed and could, as it were, by degrees, attain to high promotion and dignity, to the greatest riches, and the most desirable pleasures. If any man resisted the Roman religion and would not assent to the Romans, he was afflicted with persecution, plundered and driven into exile, imprisoned, or led to execution. Therefore, the godly were grievously tempted and did not know whither to turn. As we see the like done today in New Rome and the papist kingdom throughout the world. Wherefore, God, who does not will that man should perish but be saved, gives here the best counsel of true felicity and salvation; those who obey it are blessed.

### Paragraph 1129

And directly he shows the author of this counsel from the beginning, so that he might establish its authority, and that we might boldly receive it. “I heard,” he says, “another voice from heaven; therefore this counsel proceeds from God.” Those who follow it obey God; those who do not obey it despise and reject God’s counsel. And what is this counsel? Briefly, it is plain, possible, honest, and wholesome, without doubt. “Come away,” says the Lord, “from her—that is, Babylon, both old and new—my people: you who will be called the people of God and be written in the number of the citizens of God.” This is his counsel, and no other. God gave this same counsel to his ancient people through his prophets when they were in captivity in Babylon. For Isaiah says in chapters 48 and 52, “Depart, depart, come away from there; do not touch anything unclean.” “Come away from her, be cleansed, you who bear the vessels of the Lord.” And Jeremiah in chapter 51, “Flee from the midst of Babylon, and let every man save his soul, so that he will not be rooted out in her wickedness.” For the time of God’s vengeance is at hand; he will reward her. Therefore, the Lord counsels us to flee, so that our souls may be saved. For otherwise, unless we flee, we shall perish. However, the prophets did not teach the Israelites to flee from Babylon physically, by local motion, as they term it. For Jeremiah in chapter 29 exhorts the captive people to dwell in Babylon and to make their provision there until the time of deliverance arrives. For then they must come out of Babylon. In the meantime, he would have them depart not by physical motion, but by difference of manners. For although they shall dwell in the midst of the superstitious, ungodly, and idolaters, the Lord does not want them to be like them. Therefore, this fleeing consists in abstaining and refraining from ungodliness, idolatry, sins—namely, bloodshed, usury, pride, lechery, and other like vices—but persevering in true godliness and innocence.

### Paragraph 1130

In like manner, if the godly were to flee under the old Roman Empire, they would everywhere fall again into the hands of the Romans; as we also at this day, although we change our place, have papistry either ever-present or imminent. Therefore, the Apostle speaks well: we must get ourselves out of the world if we are to be conversant with sinners. This, therefore, is the true and godly flight: if remaining in this world bodily, in mind and manners we depart furthest out of the world, so that we abstain from all idolatry and profane worship, if we do not allow it, if it does not please us; if we neither assent nor conform ourselves to the manners of the ungodly; if we shall not betray our religion, either for people or for worldly gain. So therefore, the Christians who lived under the Roman Empire fled Rome, utterly abstaining from the worship of idols and the corrupt manners of the gentiles, although they lived among the heathens. For that the ancient churches in Asia were such, we have heard in the second and third chapters of this book. Albeit therefore that we also dwell under the Popish kingdom and empire that persecutes the gospel, yet must we flee papistry, that is to say, Popish churches; none of the godly ought, for worship or obedience’ sake, to enter in, to acknowledge, allow, or use any Popish rites or ceremonies, but from their vices and corruptions to flee as far as is possible. For so the Apostolic scripture teaches us in Romans 12, 2 Corinthians 6, Ephesians 5, and 1 Peter 4. And John, as it were, expounding himself, says, “Be not partakers of her sins, lest you share in her plagues.” And sins are not only those which are done against the second table, but also those that are committed, and that much more against the first table; of which sort are idolatry, impiety, the abuse of God’s holy name, strange worship, against the second and fourth precepts of the first table. Those were then, and so are at this day taken for very good works, where they are abominations. Partaking is chiefly in the communion of sacred things, again if they are given to the same dissolute riot with filthy men. If therefore we beware of those things, we flee out of Babylon and follow the good counsel of God.

### Paragraph 1131

But in this we commonly offend today, those who are called Gospel advocates. For many think it sufficient if they observe some religion in their heart privately, and openly associate with them, which may either help or hurt. You will find those who kneel before idols, who hear Mass and popish services. There are some who know many abominations of the popish priesthood, but nevertheless make their sons priests, because the promotions, and the clerical life—that is to say, the wealthy and pleasant life—appeals to them. There are some who, through the bond of marriage, force their children into the midst of Papistry; and these regard nothing else but riches, worldly honors, and friendships. Against all these, the prophets and the Apostles, and at this present Christ the Son of God from the right hand of the Father, thunder and cry aloud, “Come away from her, my people, and be not partaker with her sins.” These words do not admit of any witty or civil reasoning, nor carnal or crafty qualifying. For it follows, lest they receive of her plagues. For if the like Rome, if the like the Roman religion, if the Roman prelacy please you, riches and promotions, if the Roman corruption content you, let the judgment, pain, and damnation due to Antichristianism content you also.

### Paragraph 1132

We also have, at this time, what answer we may give to Romanists, objecting and claiming that we are rebelling or committing apostasy, and for the same reason, the crime of schism. “You have fallen from the holy church of Rome,” they say, “and by that very abandonment, you declare openly that you are sectaries and schismatics.” To this we answer that we make a distinction within the church of Rome. For we acknowledge an old church of Rome, venerable and apostolic. Of which Paul wrote: “Your faith is proclaimed throughout the world.” Whoever departs from that church, without doubt, will be both a schismatic and will perish forever.

There is again another church of Rome, new, and entirely contrary to the old, no longer apostolic but rather Papistic, wherein are not ministers of the word and sacraments, but either princes, not unlike the Gentiles, or merchants, from whom the sacraments, the remission of sins, heaven itself, and all things in the church are sold for a little money. They teach a doctrine that deviates entirely from the doctrine of the gospel. These are openly, not infected, but swimming and stinking of the most shameful vices, even the filthiness of whoredom—to say nothing now of Christian bloodshed. And there is no sign of repentance in them. To persevere with these, to commune with these, is to perish eternally. Therefore, from the company of these men, the Lord commands us here to depart, yes, and to flee from. That is what we have done, at the Lord’s commandment, which openly here commands us to come away, depart, and flee from the scarlet woman and from this Babylon.

There are also other notable passages commanding this departure, which whoever wishes to know and consider should read Deuteronomy 13, Jeremiah 23, the words of the Lord in the gospel of Luke 6:23 and 24, and Matthew 7:23. Read also both epistles of Paul to Timothy, especially 6:3-5 of the first, and 1:3 and 4, and 2:4 of the second. In Romans 16, he says: “I beseech you, brethren, mark those who cause divisions and put up to evil, contrary to the doctrine which you have learned, and avoid them. For those who are such do not serve our Lord Jesus Christ, but their own appetites, and through flattering words.”

### Paragraph 1133

And explaining the reason why we should flee from Babylon, he speaks of profit and loss. Lest we receive her plagues. For whoever aligns himself with the ungodly, idolaters, and those who are filthy and unclean, receives the same reward as they do: and the reward of this present life is a curse, a reprobate mind, and various calamities, as described in chapter 16 and elsewhere, and after this life, everlasting torments. Therefore, he is not speaking of a trivial matter when he warns of fleeing from Babylon or avoiding the Roman religion. Many believe these things, for they do not consider how great is the abomination of the church of Rome before God; and therefore they hear these things as a fable and persevere in the same kind of life in which they are and have lived until now. But he does not lie who says that those who do not prepare to flee from Babylon will shortly perish with Babylon and with the entire fellowship of the wicked. Woe to them.

### Paragraph 1134

However, since the wicked in this world are commonly fortunate (and many from this infer that God does not know our affairs, or at least if he knows them, does not care greatly), the Apostle, or an oracle brought from heaven, adds: “For her sins have risen up to heaven, and the Lord has remembered her wickedness.” God truly never forgets iniquities. For all things are evermore present before him. Yet he seems not to remember when he does not punish. For so men suppose: but when he punishes and visits sinners, he seems utterly to have had consideration of our affairs, and to have remembered wickedness and wicked men. Therefore God is righteous, and mindful of evil, and of good also: and when he sees that the time will recompense all men's works, and chiefly the evil. In the meantime, he signifies also that the sins of old and new Rome are great and full of enormity. For in Genesis 19, the sins of Sodom are said to have risen up to heaven, as if to cry out against those who committed them and demand vengeance. So we read in Jeremiah 51 that the sins of Babylon ascended up to the clouds. For John, in a manner everywhere, uses the places of Scripture, to the intent he might give his book more authority, although otherwise inspired by the Holy Spirit. Indeed, the old satirical poets, such as Horace, Juvenal, and others, wrote severely against the sins and vices of old Rome. There remain also at this day many sharp writings against Rome, and the Cardinals and Prelates of the Roman church, and Pasquils innumerable (Pasquil at this day is a satirical writer, one in place of many), so that as well at this day as in times past the sins of Rome cry up to heaven itself.

### Paragraph 1135

He proceeds after this to recount again the plagues and most certain destruction of Rome, which is the third section of this chapter. Here, the most horrible and cruel manner of destruction and subversion is exceedingly well described. For God is portrayed as calling on and exhorting the soldiers, the commissaries, and the executors of his judgment to vengeance: and that they should punish her most extremely, and spare her not, but reward her most abundantly, and mete unto her by the same measure wherewith Rome has measured to others. For here takes place that saying of the Lord, a common saying with all nations: with the same measure wherewith you mete, others shall mete unto you again, and there shall be given good measure, pressed, shaken, and running over. Therefore, saying that Rome has robbed the whole world and seduced the whole world, rightly and by the just wrath of God was she spoiled and utterly subverted. The Goths accomplished these things with great faith and diligence, so that we cannot doubt but that New Rome, and that Apostate See, must of her enemies, whom the Lord has prepared, and of the Angels gathering the tares, be plucked all to pieces. And what shall become of her in another world, we may gather hereof, that he beats in so often that her evils shall be doubled without mercy, her pain also, mourning, and grievous torments. These things doubtless are grievous and horrible. Would God they might be perceived by faithful minds. And again, this passage is written as it were word for word from the 50th chapter of Jeremiah: where you read to this effect: be avenged on Babylon, and as she did, do the unto her. Spoil and destroy says the Lord, and accomplish all that I have commanded thee. Destroy her, that nothing remain. Entrench round about, that no man escape. Reward her after her work: and according to all things that she hath done, do the unto her. For she has been proud against the Lord, and against the holy one of Israel. Thus said the Lord in Jeremiah. You see therefore where the Lord borrows his own at this present. You see what every city, or commonwealth, or man may promise himself, if being enriched by the loss of others, he live voluptuously and proudly in this world. For God is the same always, and his judgments are keen against all the ungodly.

### Paragraph 1136

And he has been involved in the causes of ruin,[note 78.10] cruelty, covetousness, extortion, slaughter, burning, with which Rome has laid waste the entire world. But he proceeds more expressly to recite other causes: namely, pride, glorying, and boasting, security, riot, pleasures, and indulgence. For it follows: as much as she has glorified herself and lived luxuriously, and so forth. And again: for in her heart she says, “I sit as a queen,” and so forth. He has borrowed these things also from Isaiah 47. Where Babylon glories thus also, and with so many words. Rome in times past gloried herself to be Lady of the world, and that she should be everlasting. For they stamped on silver coins the words “Rome eternal.” They had thought that the kingdoms would never be taken from her. She thought therefore that she would never be a widow. And I doubt not but the Germans borrowed from the Romans that Germanic word *Römer*, by which they mean to boast or brag stoutly, which seems to have been peculiar and proper to the Romans. She was careless and secure. She had not thought to be overthrown. She said, “I shall see no morning; I will have no morning cheer; I will always sing, ‘Gaudeamus.’” The Romanists at this day also boldly make their boast that no emperors, no kings, no people, no heretics, or schismatics (for so they term the enemies of Roman wickedness, godly and learned men) have yet successfully assailed Rome. That the enemies of the church of Rome have always been oppressed, that she has always triumphed over her enemies, these seven or eight hundred years and more. That the ship of St. Peter may be sorely troubled, tossed, and overwhelmed with waves and billows, but cannot be drowned; and therefore that the See of Rome shall be a perpetual queen and lady of all realms and churches. and so forth.

### Paragraph 1137

But hear now God’s judgment: for as much as she is proud, vain, glorious, careless, and wicked, her plagues will come in one day. Arethas notes that by “one day” is signified a sudden destruction, and that she will perish when she least expects it. He recounts her plagues in order: death, mourning, famine, and fire. Historical accounts testify that these things were fulfilled in old Rome, as I have spoken of before. Therefore, we doubt nothing at all, but that new Rome also will be torn apart by humans and by God’s angels and uprooted. And lest any man should think this impossible (for great is the power and majesty of Rome, so that anyone who had said in John’s time that Rome would fall would have seemed to speak something as impossible as saying the sky would fall), he adds immediately: for the Lord God who will judge her is mighty. Therefore, let us not doubt the fall of Papistry. For the Lord is true, just, and almighty. To whom be glory forever and ever. Amen.

## Sermon 79: ¶ A doleful song or mourning, and lamentation of Rome, which the Princes and Merchant make for her.

### Paragraph 1138

.

### Paragraph 1139

And the kings of the earth shall weep and wail over her, those who committed adultery with her and lived sensually, when they see the smoke of her burning, and shall stand afar off, for fear of her punishment, saying, “Alas, alas, that great city Babylon, that mighty city! For in one hour your judgment has come.” And the merchants of the earth shall weep and wail in themselves, because no one will buy their wares anymore—the wares of gold and silver, and precious stones, of pearls, and silk, and purple, and scarlet, and all manner of wood, and all manner of vessels of ivory, and all manner of vessels of most precious wood, and of brass, and of iron, and cinnamon, and perfumes, and ointments, and frankincense, and wine, and oil, and fine flour, and livestock, and sheep, and horses, and chariots, and bodies and souls of men. And the delights that your soul desired have departed from you. And all things which were luxurious and were of great price have departed from you, and you shall find them no more. The merchants of these things, who became rich from her, shall stand afar off for fear of her punishment, weeping and wailing, saying, “Alas, alas, that great city, that was clothed in silk and purple, and scarlet, and decked with gold and precious stones and pearls! For in one hour so great riches have come to nothing.” And every ship’s captain, and all who occupy ships, and the sailors who work on the sea, stood afar off and cried when they saw the smoke of her burning, and said, “What city is like this great city?” And they cast dust on their heads, and cried weeping and wailing, and said, “Alas, alas, the great city, in which all who had ships on the sea were made rich by reason of her wares! For in one hour she is made desolate.”

### Paragraph 1140

Fourthly, this chapter addresses the weeping, lamentation, or mourning of Rome, which will be struck down and destroyed. The text is abundant and marvelous, employing a clear hypostasis that presents all things before our eyes. Our Lord God consistently uses a familiar method: whenever he intends to clearly demonstrate and impress upon all people the downfall or destruction of a nation, kingdom, or city, he commands his prophets to compose an elegy, or a lamentable song. Such lamentations reveal not only the subversion but also the causes of destruction, and the manner of desolation is recounted. The purpose is also declared, so that others may not become like that nation and share in its destruction. We have clear examples in the writings of the prophets, especially the lamentations of Jerome, and the mournful song of Tyre in Ezekiel, chapters 27 and 28, which aligns more closely with this passage. Indeed, it appears that John has borrowed many things from that source.

### Paragraph 1141

There is nothing here to greatly concern ourselves with. The sum of it all is this: Rome will fall and perish completely, so that nothing remains, neither of the empire nor of that See, much less of the riches and pleasures. This was partly fulfilled in ancient Rome, and partly will be fulfilled in the new Rome at the day of judgment.

### Paragraph 1142

However, neither Christ himself, nor the Apostle is recorded lamenting the downfall of Babylon, but rather wicked persons are induced [?]. These are Kings and Princes of the Earth, merchants, or governors of ships, or mariners, who have all committed fornication with this harlot and, through her company, have become rich. Indeed, old Rome was supplied with the friendship of Kings; and again, the officials sent by them to govern provinces appeared to be Kings and Princes. And because the riches of Rome were great, and all nations were plunged into riot, the merchants there gained exceedingly. Moreover, ships sailed to Rome from the East, South, and West—that is, from Syria, Egypt or Africa, and from Spain itself, and the furthest parts of the world. But when Rome was destroyed and lay in ruins, and the Empire was torn apart, those whose profit and pleasure were lost could do nothing but lament.

### Paragraph 1143

New Rome has also, besides those temporal [note 79.3], even peculiar merchants and princes of her own. For the prelates of the church are princes. And in the church of Rome, all the saints of God know, how to occupy the traffic of merchandise. For what holy thing is not to be bought in that seat? Merchandise is practiced in forgiveness of sins, in pardons and satisfactions, in ecclesiastical benefices, in the worship of images and saints, in masses, in burials, in saying diriges for the dead, and almost in all spiritual matters. Hereof comes an immeasurable gain, and the greatest occasion of pleasures. Other merchants buy their wares very dearly; the Roman Canaanites pay not one denier or farthing for their wares, but sell them for an unreasonable price. Neither suppose I that ever there was any merchandise like unto this in all the world, nor yet a more gainful profit from a thing of naught. Erasmus has also touched these things, in the proverb to ask tribute of a dead man. And whereas before the day of judgment, the Lord Christ shall destroy Antichrist with the spirit of his mouth, and that gain begins to be diminished, we see how everywhere among these spiritual merchants, complaints and grudgings arise. Then what manner of lamentation and wailing think you will be, where the same Lord by his coming shall utterly abolish the same Antichrist, and they must go into fire everlasting?

Again, we must somewhat also consider the mourning. To mourn of itself, is no sin. For the best and holiest men have lamented their dead, and their calamities, and destruction of cities and realms. For Abraham mourned. The lamentations also of Jeremiah remain, over the city of Jerusalem. The faithful mourned with a great mourning for Stephen in the Acts. However, in their lamentation they kept a mean, and referred all things to the glory of God, and salvation of souls. The ungodly and worldly men do not mourn after this sort. They never remember the signs of men, for which the righteous Lord punishes the world, neither do they refer the evils of them and theirs to the glory, verity, and justice of God, or amendment of manners: therefore are they not sorry that God is offended, nor require forgiveness of sins; but it grieves them that occasion of sinning is taken from them, that their pleasures and profit is past.

And now, while princes, merchants, and mariners—not for the favor of God lost, not of true compassion or love of their neighbor, but for love of themselves, for the loss of earthly things, for the destruction of goodly, ancient, strong, and precious things—but chiefly for their profit lost and pleasures taken away. The Apostle makes mention of such grief in the second to the Corinthians, chapter 7. And surely this sorrow and mourning is nothing else but a description and a shadowing of a most certain and greatest destruction, and that of men ungodly. And full well and purposefully doth he set forth the wailing both in the behavior of the mourners, and also by their words. To their gesture appertains that they weep, wail, cry out, and cast dust on their heads. To their words are referred these things: “Woe, woe, alas, alas, that great city!” &c., which is repeated of the merchants and sailors.

### Paragraph 1144

Moreover, here are also touched upon the causes of destruction, the riot and luxury in which Rome indulged. And likewise are rehearsed the wealth, riches, majesty, pride, and pleasures of Rome. And here, by the way, are warned all worldly men regarding what they may expect if they devote themselves to the pleasures and luxury of this world, which was present in Rome and is immeasurable. We have not read in any histories that nations have long endured, whether they have been blessed or conquered by worldly pleasures. To build, to eat, to drink, to be clothed, and to have servants, men and women, is lawful; but a measure must be kept in these things, as in all others. The benefits of God must be acknowledged, and they may not be valued more highly than virtue. But at Rome, and in the world, godliness and moderation are passed over, and these things alone are regarded, desired, and beloved. In buildings and household goods, all things were sumptuous and immeasurable. They were made of gold, which could have been well of earth or tin; of silver, where wood or iron might have served. And when wood was chosen, it was not every wood, but the most excellent was chosen. [illegible] appears to be named of Thyia, a tree, to which Theophrastus attributes great honor, reporting that the famous buildings of old temples were made thereof, and a certain immortality of matter, incorrupt, enduring on houses against all weathers, etc. Pliny has this in Book 13, Chapter 16. In service also, they use men like beasts; neither do they have beasts for their own use, but mostly chosen ones. They have horses and mules exceeding fine. They have their horse litters, coaches, and chariots, truly notable; all things glisten with gold, precious stones, and purple; and all things are wrought and devised for pride and sumptuousness. What shall we say that the whole bands of their men go all in silks and velvet, wearing their masters’ colors? The lord himself of all, sitting on the shoulders of his Belphougers, is borne on high and is carried on men’s bodies as the most noble chariot. In the meat and drink of these men, all things are most delicate, exquisite, and varied. Their drink is costly, strange, and immoderate. The apparel of their body is also over-sumptuous. Their garments glisten with gold and are stiff with pearls. Their common garment is of crimson satin. They use also ointments and apples of desire, which may both be understood of the fruits of trees, and also of pomanders containing musk and smelling sweet, and of odoriferous savors.

### Paragraph 1145

Finally, in all things, we must consider what the ultimate fate is of riot, pride, and sensuality, and how unstable the favor and friendship of men are. Here, everything perishes; nothing remains secure. And those who have prepared for many years perish truly in one hour. They flee from us in danger, having received great gain from our hands. Indeed, they stand far off, out of danger, and lament the sorrowful circumstance; no one comes to help or deliver us. Every man is afraid of his own life. Let us therefore learn to trust in God, to despise pleasures, to place no confidence in flesh or the friendship of men. For while you are fortunate, you will have many friends; if the world begins to frown upon you, they will all forsake you, in whom you placed your trust, and leave you in the thorns. And this is the chief end of all these things, as I showed at the beginning: Rome shall fall and be made desolate forever. The Lord our God restrain all evil. Amen.

## Sermon 80: ¶ The reioycyng of Saints for the overthrow of Babylon, the drowning of the fame, and the causes of drowning or destruction are rehearsed.

### Paragraph 1146

.

### Paragraph 1147

Rejoice over her, heaven, and the holy Apostles and prophets, for God has given you judgment upon her. And a mighty angel lifted up a great millstone and cast it into the sea, saying, “With such violence will that great city Babylon be cast down, and will be found no more.” And the sound of harps and musicians, and of pipers and trumpets, will be heard no more in you; and no craftsman, whatever his craft may be, will be found no more in you; and the sound of a mill will be heard no more in you; and the light of a candle will shine no more in you; and the voice of the bridegroom and of the bride will be heard no more in you, for your merchants were princes of the earth, and with your enchantments all nations were deceived; and in her was found the blood of the prophets and of the saints, and of all those who were slain upon the earth.

### Paragraph 1148

In the fifteenth place of this chapter, the Angel of the Lord exhorts all the saints in Heaven to rejoice, [note 80.1] and that for the overthrow of Babylon. And this rejoicing of the saints is set against the wailing of the wicked. For as they lament for the loss of pleasures, so the saints rejoice over ungodliness oppressed, and the glory of God avenged.

We are certainly forbidden in the Proverbs of Solomon, and in the doctrine of Christ and his Apostles, that we should not be glad of the calamities of our enemies, nor that we should say evil or do evil to our enemies. This is a perpetual command, given to all men, never to be altered by any dispensation. But we must observe that men rejoice in diverse ways. Men are glad many times at the destruction of their enemies, and that out of hatred and malice, which is not done without sin. Others are glad again of the calamities and plagues of the ungodly, yet bearing no malice toward them, and would doubtless have wished them a better state if they might have been persuaded to turn. But they rejoice rather over justice revenged, and the godly delivered from the tyranny of the wicked. Of this we read that the prophet said in Psalm 58: “The righteous shall rejoice when he shall see vengeance; he shall wash his feet in the blood of the ungodly” (that is, he shall purge his affections and evil manners, [note 80.2] when he shall see the blood of the ungodly spilled, which he believes to be done as a warning, lest we should follow our evil affections, and that our blood should be shed also by the most just God through his ministers). And a man will say, “Verily, there is a reward for the righteous; verily, God judges the earth.” Therefore, the righteous are glad and rejoice when they see vengeance. And it is not said that they covet or wish for vengeance. “Vengeance is mine,” says the Lord, “I will reward.” When the Lord therefore rewards, they are glad for the deliverance, and for the truth established and confirmed, and rejoice not out of hatred they have toward the oppressors, whom they have wished lost and destroyed.

The godly always wish the wicked to be converted and to return into favor with God. But when they see them unmoved by repentance, but obstinately proceeding and falling into their own destruction, and that God intervenes for the salvation of the faithful and the deliverance of the godly, the godly rejoice at this deliverance and praise the justice of God. Nevertheless, they would always rather, if it could be, that the lost had led their lives otherwise. But now, since it cannot be otherwise through their own obstinate malice, they do not speak against the judgments of God, but rather commend them. These things, certainly, the saints on Earth do. And the saints in Heaven, since they are now purified from all affections, their rejoicing is altogether most pure, so that it would be superfluous to reason curiously about it.

But where the heavenly saints rejoice at the destruction of the wicked, we may easily judge how much they err who trust to the help or prayers of saints, and yet alter nothing at all of their wicked life. It shall be easy also to discuss their doubt and carefulness, who fear lest their brethren, sisters, friends, and kinsfolk be condemned. For the saints plainly consent to the will of God, and extol the judgments of God, and rejoice thereat, and can be sorry no more.

### Paragraph 1149

And he bids them rejoice, as many times in the Psalms we read a like phrase; unless you would rather understand by “Heaven,” heavenly dwellers, such as we believe the apostles and prophets to be. For at the same time when John wrote these things, all the apostles were in a manner slain. And here it is to be known that the Roman beast had devoured, that is to say, afflicted and slain, not only the Son of God, our Lord Jesus Christ, but also John the Baptist, all the Apostles of God, and all the martyrs of Christ. By the prophets we understand, not only those of old, but all the faithful preachers of the gospel. For we have heard more than once that faithful preachers of the word are called prophets. He adds moreover a reason why they ought to rejoice: for God has given you judgment concerning her. For in chapter 6, the souls of the martyrs cry under the altar: “Lord, how long before you avenge our blood on those who dwell on Earth?” Now therefore they praise God’s justice, which, as he promised he would avenge, so has he now avenged in deed.

### Paragraph 1150

And by this passage we learn that all judgment is given to the Son, and that no saint in heaven can judge or punish an evil man on Earth. For it is most false that saints are said to punish their enemies: Saint Anthony with the holy fire, Valentine with the falling sickness, and others with other diseases; God alone, as is declared at large in chapter 16, punishes and sends and takes away sickness. And most certain it is, both by this and by many other passages in this book, that God does not sleep, but will, when he sees fit, certainly avenge and punish. The martyrs, when they were about to die, had committed all their judgment to the Lord their God. He now judges the judgment of the saints of Rome: that is, after his just judgment takes punishment of Rome, for that she had with wrongful judgment oppressed the saints.

### Paragraph 1151

In the sixth place of this chapter he returns to the description of the subversion of Babylon. And it is a most clear, and even a certain visible and evident demonstration by a similitude. For taking up a great stone, in quantity like a millstone, he casts the same into the sea, and making a declaration of his doing, says, thus suddenly, and with such a violence, shall Babylon be cast down, &c. This place is taken out of the end of the 51st chapter of Jeremiah, where you read in a manner the like things word for word. And here is now brought in a strong angel, lest we should think that the force of Rome were happily stronger than that it could be broken. But it shall be broken by a strong angel. And the things that be suddenly drowned appear no more. Here is signified therefore, that with a sudden destruction Rome shall fall, that there shall no token thereof be left, and that it shall fall without any difficulty, it shall be made to plunge, and never more be seen. And the Lord in the gospel affirms, that the crime or slander must be punished with a millstone hanged about the neck: yea and that same not to be punishment grievous enough, although among the Syrians it was accounted for vile and shameful, since the crime deserves to be punished with a much more grievous or crueler pain. Wherefore Primasius supposed, that here by the way is signified how Babylon, for offenses given to the world, should be drowned in the sea, as it were with a millstone tied fast to her neck. Doubtless if ever any city, if ever any kingdom were hateful by reason of greatest offenses, and given to the Christians innumerable slanders: Rome and the Roman Empire, and even the papists of the church have hurt most by slander, and yet hurt. Wherefore it is no doubt, but that it has been plagued most grievously, and shall be yet more punished by the Lord.

### Paragraph 1152

Again, through prophetic and figurative language, he signifies a significant desolation, and that the place would never afterward be inhabited forever. You will find similar expressions in Isaiah 24 and Ezekiel 26, and in various other passages. All pleasure, he says, will perish, especially that which was formerly derived from music. All trades will cease. Briefly, there will be no more any dwelling place for people.

### Paragraph 1153

In the seventh and final place, I set forth again the causes of this downfall, and these are particularly noteworthy. The first: “Your merchants were princes of the earth.” For those who have engaged in commerce within the church of Rome, and still do, are in a manner princes. Of whom I have spoken before. Here is noted therefore their pride, avarice, and lavishness. Arethas says he calls them merchants, who trouble and vex the entire world, as if they were certain fairs, and so forth. The second: “By your enchantment all nations have been seduced.” There is no doubt that enchantment and magic reign in Babylon, and that there are plenty of fortune-tellers, necromancers, and enchanters to be found there; yet here chiefly is signified seduction, idolatry, and impiety, or error of doctrine. Such an enchanter was Jezebel, as appears in the fourth book of Kings, chapter 9, who practiced enchantment indeed, and bewitched people with corrupt religion. And even so Rome has seduced the entire world, and yet seduces. For this cause she deserves most grievous punishment. The last cause of subversion: “For in you is found the blood.” Blood shed cannot be wiped away nor cleansed from those who shed innocent blood. And although it is not immediately required, there will come a time when it shall be required of God, and then it will be found. And he makes mention of three kinds of blood. First, the blood of prophets, meaning those who have preached the Gospel and have been the fathers of the faithful. Secondly, the blood of saints, namely, holy martyrs. Finally, the blood of all men who have been slain on earth, namely, dwelling here and there throughout the world, whom we understand to have been dispatched and taken out of the way by the wars, seditions, and tyranny of Rome. So we also read in the first oration of Jeremiah, that God straightly requires the blood of his servants spilled. Doubtless all shedding of blood is grievous (the same excepted which is justly done by the magistrate), yet one is more heinous than another. For he who kills a preacher of the Gospel sins more grievously than he who dispatches a private person, and he who for religion’s sake slays a man, making him a martyr, sins more heinously than he who kills a man in war. Therefore, all the blood shed by Rome in any sort shall be required of Rome, and is required. Thus the Lord also spoke of the city of Jerusalem, Matthew 23. The Lord Jesus have mercy on us, and look upon us with eyes of his mercy. Amen.

## Sermon 81: ¶ The reioycinges and Himnes of saints are recited for Rome destroyed, and all vngodlines taken away.

### Paragraph 1154

.

### Paragraph 1155

And after that, I heard the voice of a great multitude in Heaven, saying: Alleluia. Salvation and glory, and honour, and power be ascribed to the Lord our God. For true and righteous are his judgments, because he has judged the great whore, who did corrupt the earth with her fornication, and has avenged the blood of his servants of her hand. And again they said, Alleluia. And the smoke of her ascended forevermore. And the twenty-four elders, and the four beasts fell down, and worshipped God that sat on the throne, saying: Amen. Alleluia. And a voice came out of the throne, saying: Praise our Lord God all those that are his servants, and those that fear him both small and great. And I heard the voice of a great multitude, even as the voice of many waters, and as the voice of great thunderings, saying: Alleluia. For our Lord God omnipotent reigns. Let us be glad and rejoice, and give honour unto him: for the marriage of the Lamb is come.

### Paragraph 1156

As the Apostle in this book has abundantly described the oppression of the saints, and the cruel and proud assaults of those who persecute the Gospel—those who mock God and torment his saints—it is said that the complaints of the godly have always arisen, as though God, through his suffering and great patience, were neglecting the oppressed. He also now discourses at length about the rejoicing and praises of the saints, by which they extol the truth and justice of God, never failing to punish the ungodly persecutors. They rejoice here chiefly, and praise God for the removal of Antichrist and all ungodliness with him. This is indeed the first point of this chapter. The second confirms all saints, lest they doubt anything about the salvation of the faithful, which he shows to be most certain. The third recounts the sin of blessed John and the faithful doctrine of the holy Angel, that we should worship no creatures, however holy. Finally, the judge, or avenger, Jesus Christ, coming in judgment is described; moreover, the perdition or punishment of all the ungodly, which the just and holy Lord will take of them, is described. This section, which began in chapter 11 of this book and has been suspended until now, is now finally completed.

### Paragraph 1157

And indeed the Jubilee of the Saints is diverse, plentiful, and manifold, over the lost and condemned enemies of the godly. First, he hears a voice, and that a great voice of much people in heaven. He shows therefore in general that all heavenly beings—the angels not excepted—sing praises to God in heaven. This we understand shall be at the last judgment, when all the ungodly are trodden under foot. And before these things are done, they are rehearsed and described, so that hereby the godly may comfort themselves in dangers and torments, and may abide steadfast in the true faith, believing that they also, though now oppressed, shall sing praises of thanks to God.

### Paragraph 1158

And indeed, he here compiles the entire hymn, spoken in praise of God, the avenger. He places “Alleluia” first, after which he adds the praises, “Salvation and glory,” etc. And “Alleluia” signifies “praise the Lord.” He uses a very common term, well known among all in the primitive church. For certain Psalms have this title, “Halleluyah.” For the chapter thus exhorted and stirred up the people to praise God. So, likewise, now the saints, as if embracing the argument of their song, say, “Alleluia.” These terms have more grace in our and foreign languages than in translation. Thus, they have remained in the church: “O[s]anna,” “Amen,” “Selah,” “Maranatha,” and diverse others. Concerning these, Saint Jerome also writes to Marcella and Damasus.

### Paragraph 1159

Now follows the hymn: salvation and glory, and honor, and so forth. They praise these things in God, and ascribe them wholly to him. This is what I explained in interpreting chapters 4 and 5 of this book. Moreover, they praise God for what is principal in this matter: for his judgments are just and true. This saying seems worthy to be imprinted deeply in the hearts of all men, as it may greatly encourage them in temptations. And because the judgments of God are just and true, he adds, because he has judged the great whore: that is to say, he has taken worthy and fitting punishment of the great whore. Hitherto, the Lord has seemed to many to be too slow and too favorable to Rome and the Roman church, but then they shall see that God is most just. The whore has been spoken of before.

### Paragraph 1160

Yet does he repeat here again her most heinous and greatest sins. First, corruption through whoredom and enchantment, whereby is signified seduction by corrupt and wicked doctrine. The later, the shedding of the blood of holy martyrs, whereof we have already spoken many times. Therefore God punishes the corruption of doctrine and the cruelty of the Roman church practiced against the saints of God.

### Paragraph 1161

And just as in the beginning they sang “Alleluia,” so also in the end they repeat the same. By this repetition, they declare that the praises which we pour out to God on Earth are most acceptable to God. Immediately following is a sentence that might seem to be from either John or the heavenly dwellers themselves. It signifies that the burning of the ungodly will be perpetual and will never end, as Isaiah said in chapters 30 and 46, and the Lord himself in Matthew 25 and Mark 9. For when he speaks of smoke, he understands that there is fire beneath. Let us earnestly consider these things whenever the pleasures and commodities of Antichrist flatter us. For this perpetual fire is prepared for all the ungodly, especially Antichristians. Then he brings in the twenty-four elders and four beasts praising God, by which the universality of creation is understood. Concerning this, see what is said in chapters 4 and 5 of this book. First, they not only kneel but also fall down, so that we should understand what we ought to do on Earth. They worship God who sits on the throne, neither angels, nor spirits, nor any creatures. Furthermore, with two words he shadows their hymn. For they sing, “Amen” and “Alleluia.” They confirm God to be just, and his judgments to be righteous, and that justly he punishes the whore, and therefore that he is to be praised.

### Paragraph 1162

Now also a voice comes out of the throne, [note 81.7] namely from God himself, but by the ministry of an angel. For it follows: “Sing praise to our God.” Behold, he says, “our God.” Therefore, he accounts himself here among those who share God with humans. Therefore, he was an angel who recited those things of God. Therefore, what the saints are doing now, they are commanded to do. For in the midst of the praises, this voice is heard from God by the angel. And he commands that they praise, and that the true and only God. He shows moreover who should praise him: all the saints, that is, all who fear God, whether they be great or small. By this commandment, therefore, is signified that God is delighted with the praises of holy men, and approves of them. From this, we now who dwell on Earth, learn to praise the Lord without ceasing, and with a sincere heart. We learn that no one is excepted, of whatever degree or age, sex or condition he be.

### Paragraph 1163

Again another hymn is appended,[note 81.8] as it were an example of obedience. For God by the angel commands the saints to praise. Now therefore they obey God, and offer to him praises. And how great these praises were, he shows by a double comparison, and by a marvelous brevity, and evident manner. For he says, how the voices of the singers were shrill, as the gushing and noise of many waters: also like the clapping or cracking of great thunders. If such brevity and perspicuity were found in Homer or Virgil, it should have many marvelers thereof, who would extol and commend the elegance. But no man marvels, no man sets forth or commends the elegance and efficacy of the holy Scriptures, lacking example. And again is appended a hymn, the beginning whereof, as of the former, is also Alleluya. And like as in the former hymn the saints have celebrated that God does justly punish the wicked: so in this they preach that God reigns, and shall seem even freely to save the saints. They command therefore to praise the Lord. The reason, for because he is omnipotent, he reigns. He has verily reigned evermore: but since so many things have been permitted by him to the ungodly, many have thought that the ungodly, and chiefly Antichrist, has reigned: but now he has oppressed him, and avenged his glory and his servants, it is made manifest to all men that God alone reigns forevermore. They allege also another cause, why God should be praised, the rather why the godly should be glad and rejoice: for the marriage of the Lamb is come. For as much as that time is now come, wherein the Lamb himself will bring in the children of God, his well-beloved spouse, those I mean whom by his bloodshed he has redeemed, to joys everlasting. Of the marriage shall be spoken a little after at large. Praise and glory be to our Redeemer Christ Jesus the Lord. Amen.

## Sermon 82: ¶ Of the marriage of the Lamb, and of the making ready of the Lambs wife.

### Paragraph 1164

.

### Paragraph 1165

And his wife prepared herself. And it was decreed that she should be clothed with pure and beautiful silk, for silk is the righteousness of the saints.

### Paragraph 1166

The saints celebrate the Lord with praises, rejoicing, and hymns. There are countless reasons for this, yet two stand out above the others. The first is that the Lord has judged the harlot and avenged the blood of the saints. The second is that the marriage of the Lamb has come. They rejoice therefore at the justice of God, by which he has punished the ungodly, and at his mercy or grace, by which he gives to the godly a blessed life. But here we must speak of the marriage.

### Paragraph 1167

There is much mention made in the holy Scriptures, both of the Old and New Testaments, of matrimony and marriage. This cannot be explained according to the literal meaning, but by an allegory, lest we fall into shameful and most outrageous absurdities alongside the Turks and Mahometans. For spiritual things are represented by physical matters. The spiritual essence of this is: God the Father, the lover of mankind, will save men by his Son. This is declared by a parable of wedlock and marriage. And in matrimony there is a contract or a binding agreement, there is coupling or a handfasting of each party, and finally, marriage.

### Paragraph 1168

In the contract not only the young man and the maid are affianced, but also the whole manner of the marriage to come is appointed, and an order taken. For the lawyers say that affiancing is a promise of the marriage to come. This contract was made at the beginning of the world, where God promises that he will deliver mankind by his son, and receive him into glory. Hereunto appertain all the promises of Christ, of the remission of sins, and everlasting life. Moreover, the duties of the spouse are prescribed. She promises to be obedient, and other things, &c. Christ, the son of God the Father, as bridegroom, affiances himself to all the chosen through his free grace: he promises them his righteousness, all heavenly gifts, and eternal life. He takes upon himself moreover all the infirmities of the bride, and purges her filthiness. And the bride is affianced to him by faith, as it is with Hosea, and binds herself wholly to him: after whose will and law she frames herself wholly. For she is the body of a lively head. As Paul says in 5.1 to the Ephesians. The bride’s leaders are the prophets, patriarchs, Apostles. So John the Baptist in 3.29 of John, calls himself the friend of the bridegroom. He adds, to be the spouse of Christ. Paul, 2 Corinthians 11, “I have espoused you to one husband, a chaste virgin,” &c. Hereunto the 16th chapter of Ezekiel seems to appertain.

### Paragraph 1169

The joining together of either party is performed after they are betrothed, with certain ceremonies: namely, by taking each other by the hands, and certain words spoken, a token or a ring is given, and so forth. But immediately after the beginning, a league or bond was made between God and men, which is often read of, not without ceremonies, certain words and sacrifices being offered, as by Abraham, Moses, and others. God binds himself to men, and men to him, and that not without sacraments. Hereunto belong all those things that God would be in league with man, and have men bounden to him, and all his things communicated to us. And this marriage, of all others, is most strictly joined while the Son of God has united our flesh into one and the same person with him, and has commanded the Apostles to preach unto all that he will have a communion with the faithful. Of this communion many things are read everywhere in the scriptures. And he has given a pledge of faith and perpetual amity, not a ring of gold, but rather the sacraments: yea, even the Holy Spirit, as Paul says in 1 Corinthians 1 and to the Ephesians 1.

### Paragraph 1170

And the marriage [note 82.4] shall be celebrated in the resurrection of the dead. The souls truly pass from bodily death into everlasting life; but the full restoration and salvation of humankind is not made perfect unless the body also comes. Therefore, at the resurrection comes the marriage of the Lamb, that is, of Christ our Redeemer. Then we are carried to meet Christ in the air, and he brings his wife into the bridal chamber of eternal glory and blessing; then shall be held that feast and delightful supper, then shall the bride enjoy forever the love of the bridegroom. This shall truly be the marriage of the Lamb. And the marriage shall be the merrier, for that the harlot being cast out and condemned, the wife and virtuous matron shall have the full and perfect joy alone. At this joy and at this marriage, the holy inhabitants of heaven rejoice.

### Paragraph 1171

Moreover, the saints experience here also a certain preparation of the Bride, so that by the way the godly may understand what is best for them and to what they should dedicate themselves in the last age. Let us prepare ourselves to meet the Bridegroom. For we look for the judge every hour. And we prepare ourselves not in one hour or day, but throughout our lifetime. And how we should be prepared, the Lord himself shows by the parable of the ten virgins. Let us adorn ourselves with true faith against Antichrist in the later days. Let us beautify ourselves with the works of charity, the works also of righteousness, chastity, and temperance. Let us not be corrupted and defiled with drunkenness, bloodshed, and the cares of this world.

### Paragraph 1172

Furthermore, lest anyone should ascribe this preparation to his own merit, strength, and virtue,[note 82.7] and that we should see also that this preparation chiefly consists in providing the garment, John adds immediately, “And to her was granted that she should array herself.” If it is given, then it is not prepared by our force or means, 1 Corinthians 4. If it is given, then it is not bought by popish dealings. Read Acts 8. He also expresses the kind of garment, of clean or pure silk, and shining or bright. For we read also in the gospel of the wedding garment. The Apostle frequently exhorts us that we should put on the Lord Jesus. These things are in allegory. But he immediately explaining this kind of garment, says that silk is the righteousness of saints.[note 82.8] Saints he calls the faithful. But where there is only one justification through faith in Christ, John speaks of justifications in the plural number. For those who are freely justified through faith alone, immediately perform sundry and many works of righteousness. For he who is righteous, as John says,[note 82.9] does works righteousness. Therefore there are justifications, namely, the righteousness of faith justifying, and the righteousness of works justifying: that is to say, declaring us to be justified by faith alone. For we are purified by the blood of Christ freely, which we receive by faith, and are fully justified, as Paul witnesses in Romans 3. Again, those who are righteous do sundry works of righteousness and commend themselves to God. Thus they do not appear naked, but clothed with their wedding garment, as we touched upon in the third chapter of this book.

### Paragraph 1173

And it is very fitting that the bride’s garment is called pure or clean, not because of her own person, whom we know to be always hindered and weakened by the flesh, but because of the Spirit sanctifying and the blood of the Son of God. As Paul testifies in Ephesians 5:1, and in 1 John 1:1. The garment is also said to be shining and bright, and that for the glorifying of the saints to come. This is mentioned in Daniel 12 and Matthew 13. For righteousness leads to glory. For whom he has justified, he has also glorified. To him be praise, honor, and glory.

## Sermon 83: ¶ Of the certetie of the salvation of Saints, and what bless or Salvation is.

### Paragraph 1174

.

### Paragraph 1175

And he said to me: Write, “Blessed are those who are called to the marriage supper of the Lamb.” And he said to me: “These are the true sayings of God.”

### Paragraph 1176

The second part of this chapter concerns the certainty of the salvation of the faithful. It explains, in the meantime, what kind of blessing the faithful receive. Enough has already been said about the marriage of the Lamb, which refers to the glory and blessing of the elect. However, many things in this life suggest doubt and seem to undermine salvation, so the wavering minds of believers are now strengthened. This doctrine is beneficial for afflicted and troubled consciences, and it refutes the teaching of sophists, who claim that humans are never assured of their salvation, because in another place the wise man says: “Man does not know whether he is worthy of love or hatred.” He spoke that on another occasion and for another purpose, as I have explained in my book on the grace of God, and so forth.

### Paragraph 1177

At this present, therefore, it is shown that the salvation of the faithful is most certain. For first, the angel commands the Evangelist to write. This is taken from the manner of men, who put in writing their testaments, covenants, and bargains, and then seal them, for the cause of credibility and for a perpetual memorial of the thing. And those who have such writings are of a quiet mind, and think themselves safe and assured against all deceptions and subtle practices. Therefore, to the intent that the mind of man might be quieted in the matter of salvation, he causes as it were an instrument to be written, whereby all the godly might be assured of certain salvation. The same manner of writing our Lord follows in other places in weighty matters, as we may see in the eighth and thirtieth chapters of Isaiah, and in the second chapter of Habakkuk. Wherefore it is less to be marveled why the Apostle Paul so often alleged that same out of Habakkuk: that the righteous shall live by faith. For this only testimony of God, as that which is cited out of the godly instrument, might suffice as all.

### Paragraph 1178

And where God openly commanded Moses and Jerome to write (concerning which we may certainly judge, and undoubtedly gather that other prophets also, apostles, and evangelists wrote not without commandment), we see the authority of the books of the Old and New Testaments among all the faithful. For they are divine, authentic; they are God’s instrument and testament, the books of God himself, which are rightly believed without any other help or confirmation. We believe the testaments and sealed writings of men; how much more ought we believe the canonical books of Scripture?

### Paragraph 1179

Again, it is plainly declared to John, what he should write:[note 83.4] Blessed are those who are called to the marriage of the Lamb. Therefore, it is evident that now it is confirmed both by the divine oracle and lawful instrument, that those who are called to the Lamb’s supper will be and are blessed. This was proclaimed by divine oracle and written authoritatively. What place of doubtfulness is left? Undoubtedly, the faithful are blessed, grafted in Christ. For they are now called, to whom the gospel is preached, by which they are called to the participation of the gifts of God, but chiefly to eternal life through Christ; and those who not only hear the word of the gospel, but also receive it and believe it with their heart. For many are called, and few are chosen. For the gospel is preached to many, and the grace of God is offered in Christ, but they do not receive it. But those who, through the grace of God, do receive it with true faith, are blessed. For they are not only called to the Marriage, but also come to the marriage and enjoy that wedding supper. These things seem to be taken from the doctrine of our Savior, which he taught in Luke 14, of those who were bidden to the marriage. Read that passage.

### Paragraph 1180

Nevertheless, it is declared along the way what that blessedness of the faithful is: nothing else truly than the enjoyment of the marriage supper of the Lamb. A supper is prepared when the day draws to a close. So also is full salvation given to the godly toward the end of the world, at the resurrection of the dead, as was explained in the previous sermon. And truly, all those things are altogether allegorical, representing to us a certain signification of eternal life and glory. Otherwise, we have learned from the doctrine of the Prophets and Apostles—which the ear has not heard, nor the eye has seen, nor has entered into the heart of man—that the same God himself has prepared for those who love him.

### Paragraph 1181

Finally, a most weighty assertion or confirmation of this is added. For he hears it uttered by an oracle from heaven. These words or sayings of God are true. They are truly true, and are from God. Or else, they are true because they are from God. Erasmus has translated: these words of God are true. And so has the vernacular translation: these sayings of God are true. By a double reason therefore are these things confirmed, which are here proposed: both because they are true, and because they are from God. Although they come to one point. For since they are from God, who is truth, they cannot but be true. Therefore let us believe these things, and leave no place to doubtfulness.

### Paragraph 1182

Here is an explanation of why the scriptures and preaching delivered by men from the scriptures are not the word of God: for they are written on paper with ink, and pronounced with a human voice, and with a sound that passes away, whereas the word of God is not human, nor corruptible, nor passing away. For the heavenly oracle declares plainly that the sayings were written into the book of Saint John, and pronounced by the angel, and are true, and are God’s word. Paul also affirms this in 1 Thessalonians, chapter 2, that the word he preached about him was the very word of God. Likewise, Saint Peter in 1 Peter, chapter 1. Therefore, let inquisitive people cease raising these paradoxes and stop arguing that the word of God written and preached is not the word of God. Whatever is written or spoken that does not agree with the holy scripture of God is not truly the word of God. The minds of the faithful should rather be drawn to this point: that they believe and adhere to all words of scripture declared in their proper sense, as the most certain words of God. For otherwise, what can we trust? What after this can we have as undoubted and certain? To God be glory.

## Sermon 84: ¶ The fact of S. John is declared, which would have worshipped the Angel, and of the Angel prohibityng.

### Paragraph 1183

.

### Paragraph 1184

And I fell at his feet to worship him. And he said to me, “Do not do that. For I am your fellow servant, and one of your brothers, and of those who have the testimony of Jesus. Worship God. For the testimony of Jesus is the spirit of prophecy.”

### Paragraph 1185

Here is the third point of this chapter, namely the actions of the Apostle John and the Angel of God. John wished to worship the Angel; but the Angel forbade him, instructing him to worship God. And before this act and undertaking of John seems chiefly to be considered. Angels are certainly noble creatures and of great power, by whom the Lord executes his greatest affairs. They often take on the form of men, and frequently appear to men, serve, protect, and do good to them, according to how God uses their ministry. For the Apostle, speaking of Angels (as I told you in Sermon 29), says they are not all ministering spirits, sent forth to serve those who will be heirs of salvation? And Scripture makes these things plain through various examples. Three appeared to Abraham in human likeness, who were Angels, instructing him; two delivered Lot himself out of the hands of the Sodomites and brought him out of the fire; whole armies of Angels surrounded Jacob, defending him against the force and violence of his brother Esau. The Lord sent his Angel before Moses and the children of Israel to lead them through the wilderness into the land of promise. Fiery chariots surrounded Elisha. An Angel lifted the siege of Jerusalem, slaying one hundred eighty-five thousand Assyrians. Daniel had Angels familiar with him. Likewise, the fathers and other Prophets. An Angel delivers Joseph out of all care; the same delivers the wise men from Herod’s treachery; immediately he commands that Christ be conveyed away to Egypt. Angels minister to Christ; in white garments they testified that the Lord was risen and ascended into heaven. The same brought the Apostles out of prison; one of them delivered Peter out of Herod’s prison. An Angel was sent to Cornelius, an Italian captain. Angels often spoke with Paul. They frequently bestow great benefits upon men. They declare themselves through God to be of great power. And while people observe these things, they would worship Angels—as is even the case now, where the Apostle John understood that Christ himself, by his Angel, had opened to him such great mysteries for the profit of the churches. While he marveled at his brightness and godly gifts, he was about to worship this Angel, the bringer of mysteries. Not that he intended or purposed to revolt from God, and to worship an Angel instead of God—for neither is it lawful to imagine such wickedness of so great an Apostle. He would therefore have worshipped and honored the Angel with *Dulia*, as they term it (and as Augustine explains it), not with *Latria*: that is to say, to worship and honor God as God, but the Angel somewhat less, as an excellent messenger of God. However, in this he offended, so that all men should understand that they sin, however many worship and honor Angels or excellent creatures with godly worship. As all the worshippers of Saints do at this day in Papistry. Neither do they have any other way to excuse their error than that same distinction, that God is worshipped and honored with worship *latrical*, and Angels and Saints with worship *dulical*, and the Virgin Mary with honor *hiperdulical*, and I know not what other things, which I am both ashamed and loath to rehearse.

### Paragraph 1186

And it appears that St. John here was entangled with the same error [note 84.11], whom otherwise we must needs confess to have sinned by Apostasy, and would have worshipped the Angel for God, or with God. Which are both two wicked, and unworthy such a man. But if he worshipped God, and nevertheless would have worshipped the Angel also, what thing else did he do than offend in the worship itself? And verily God hath permitted so worthy a man to err, as he did also Peter and Thomas; to the intent he might heal their infirmities: that is to wit, that by their errors we might learn to believe more rightly, and to honor God more purely. For this present place teaches openly, and other like examples of errors, that all the sayings and doings of Saints are not to be allowed without examination.

### Paragraph 1187

Now, consider the matter of a most excellent angel, that is to say, a godly refutation of error. [note 84.12] First, he does not lightly say, “Do not as you have planned,” but severely condemning his action, he says with a certain vehemence, “See that you do it not.” We have a similar phrase in Switzerland, when (signifying that one must be very careful) we say, “Long und thu das nitte.” See that you do it not. Therefore, we have learned by the testimony of the angel that now neither angels nor saints are to be worshipped. For the Lord himself says of saints, “They shall be as the angels of God.” I do not see why they should not be equated with angels. And we have certainly learned that they may be worshipped neither with liturgical nor secular worship. And to worship is to honor with the mind, to fall at the feet, to bow down and kneel, as I have said elsewhere.

### Paragraph 1188

After the angel shows reasons why he ought not to worship, for I am your fellow servant. He does not say “servant,” but “fellow servant,” meaning of the same office as you, under the same lord and master. For angels serve God after their manner, and so do men serve God after their manner; yet all are servants, and that of one master. And it is against reason that one servant should honor and worship another of his fellows, being of the same state and creation. It is therefore an unworthy matter that the faithful should worship the Apostles, Prophets, or Martyrs; much less does it become them to honor their dead bones. And lest any man should say, how the angel, indeed in respect of the most excellent Apostle John, confesses himself to be his fellow servant; but that there is another consideration to be had of other men, who do not come near the dignity of blessed John: and therefore we are much inferior, we may worship angels and apostles our superiors, he prevents and says, “and of your brethren.” And who are the brethren of the Apostle John? The angel himself answers, and says, “which have the testimony of Jesus.” The testimony of Jesus is the gospel, and the very faith fixed on the gospel, comprehending with a faithful mind Jesus. Wherefore all the faithful of Christ are John’s brethren; therefore the angel is their fellow servant also. And therefore none of the faithful ought to worship any angel or apostle; the Lord himself also in Matthew 12 calls all that obey his word or preaching brethren. And here is diligently to be noted, that by faith we are made the brethren of Christ, of angels and apostles. This the monks and friars should have beaten in and set forth, and not the fraternity of our Lady and fraternities of Saints, unless they had been the Apostles of that great and abominable Antichrist.

### Paragraph 1189

Moreover, the angel himself, explaining his own words again, shows what the testimony of Jesus Christ is. For the testimony of Jesus is the spirit of prophecy. And the spirit signifies revelation or understanding, and prophecy the prophetical and apostolic doctrine. Therefore, the meaning is: the testimony of Jesus Christ is nothing other than the revealing of the doctrine of prophets and apostles in the mind of the faithful through the Holy Spirit and faith. And therefore, the apostles are called witnesses in the gospel, and the gospel is a testimony. To testify is to preach. From this explanation, the following argument may be drawn: the reason for your worshipping me, John, is undoubtedly that excellent revelation and prophecy that I have revealed to you. But if I should therefore seem worthy of worship because there is in me an excellent spirit of prophecy, then by the same reasoning you should worship all your brethren, in whom is the same spirit of prophecy, namely, the testimony of Jesus, the true faith. But when you see this, and are compelled to grant it, you will find it very absurd. I perceive it to be absurd if you should worship an angel.

### Paragraph 1190

The last and strongest reason [note 84.16] why he would not be worshipped is this: worship God. This command is derived from God’s eternal and unchanging authority and law, revealed in Deuteronomy 6 and repeated by our Savior Christ in Matthew 4. If we would obey God’s law, all cults, worship, and invocations of saints would long ago have been banished and exiled from the church.

### Paragraph 1191

Furthermore, there are other passages [note 84.17] that praise the ministries and virtues of angels, yet teach that we should honor and call upon God himself. Read the excellent Psalms 34 and 91. And if anyone desires the agreement of the church fathers, let him read Augustine, saying that angels must neither be worshipped nor called upon, nor should sacrifices be made to them, nor churches erected. The key passages are found in the writings on the true religion, chapter 55, against Maximinus, an Arian bishop, first book, page 77; *The City of God*, book 8, last chapter; chapters 10, 16, 19, and 20. To God be the glory.

## Sermon 85: ¶ The description of Christ the Judge coming to the last judgment.

### Paragraph 1192

.

### Paragraph 1193

And I saw heaven open and behold a white horse, and he that sat upon him was called Faithful and True, and in righteousness he judges and makes war. His eyes were as a flame of fire, and on his head were many crowns. And he had a name written that no man knew but himself. And he was clothed with a vesture dipped in blood, and his name is called the Word of God. And the warriors which were in heaven followed him upon white horses, clothed with white and pure silk. And out of his mouth went a sharp sword, that with it he should smite the heathen. And he shall rule them with a rod of iron, and he shall tread the winepress of the fierceness and the wrath of almighty God. And hath on his vesture, and on his thigh a name written: King of Kings, and Lord of Lords.

### Paragraph 1194

Hitherto we have heard many things of the various punishments of the ungodly.[note 85.1] And because it is manifest that God inflicts punishment on the mischievous and wicked at sundry times and in diverse ways, but most fully and most severely in that very last judgment, and from then on forever, and John has once, twice, thrice begun to treat of the last judgment, as in the end of chapters 11 and 14. And yet he has always deferred, suspended, and reserved it for another place; at last thinking it time to set before all men's eyes a description chiefly necessary, he takes it in hand and now finishes it as a matter of the greatest importance. He annexes therefore to a plentiful treatise of the torments of the ungodly, a most full and evident description of the judge most righteous and greatest, and of that last judgment, most strict of all others, wherein most fully and severely the pains shall be executed upon all Antichristians and ungodly forevermore. This place (which is the fourth of this chapter) and this treatise stretches unto chapter 21. The language is grand, smelling of prophetic majesty, apostolic clarity, and efficacy. You shall find not a few of this sort in the prophets, especially in chapters 24, 25, 26, and 27 of Isaiah.

### Paragraph 1195

And truly this doctrine is very profitable and necessary to be learned and understood most diligently by all faithful people [note 85.2]. It was set forth with great diligence and abundant clarity to this end by the Prophets and Apostles, but especially by the Lord Jesus Christ himself, both in the Gospel and also in this most godly revelation. For unless you are kept in your duty by fear of the coming judgment and Judge, it is no surprise that you run mad and perish with this foolish and wicked world. In the treatise of the last judgment, one sees the end of all people, life and death, happiness and misery, pain and unspeakable, weighty reward. Whoever remembers these things well, abhors wickedness and walks in holy fear before God.

### Paragraph 1196

And we have learned of the doctrine of the Gospel, that the same day of the restoration of all, and oppression of the ungodly, and also of all ungodliness, is known to no mortal man, but to the Father alone: and therefore to inquire of the hour and moment thereof is most foolishly done, much more wickedly. Notwithstanding, the good Lord has shown and signified tokens, which when we shall see to be fulfilled and accomplished, we might lift up our heads, knowing that our redemption draws near. Behold your redemption, he says, not your torment. For he speaks of the godly, looking for their redemption from heaven, at the return of our Savior and Redeemer, our Lord Christ: which shall also inflict vengeance on his enemies, as Paul says in 2 Thessalonians 1. Therefore, let us not be curious here, searching for things unsearchable: but rather let us watch and pray, according to the wholesome precept of our Savior, judge, and revenger. Let us have our loins girded, and let lights burn in our hands. Let us look for him steadfastly in faith, and with holy hope. Let us rather take heed that the care of this world does not possess our hearts, and beware of drunkenness and surfeiting, and that we are not of the number or conversation of those who, in the days of Noah and Lot, regarded worldly things only, despised heavenly things, and laughed to scorn those who gave them good counsel, until the wrath of God was kindled and fell upon them when they least looked for it. We see all the tokens that are said should come before the day of our Lord to be fulfilled. Let us watch therefore: and these things considered in this wise, let us see and hear with great and diligent attention, what manner of judge of all shall come, and what that judgment shall be—most wished for by the godly, horrible and with trembling to be feared by the ungodly.

### Paragraph 1197

[note 85.4]First, John in the vision sees heaven open. For by a vision, to the end all things might be more evident, he not only tells so great a matter, but sets it before their eyes to behold: and that he says, he says of the revelation of Jesus Christ, lest any should object and say, “Are you not a mad fellow to talk thus of matters unknown?” For what is he that knows who or what that judge shall be? Or else what that judgment shall be? Therefore he tells these things from the judge Christ himself, and by a heavenly revelation. For other places of the Scripture show that the Lord shall come in glory and majesty: therefore with a great and most shining brightness of light, with fire and exceeding great clearness. For so it is said in Matthew 24:25 and 25:1, in Daniel 7, and in 2 Thessalonians 1. Therefore, by the opening of heaven is signified that the whole world shall be lightened with glory and brightness, and that the same day shall be most shining and clear. Others understand that the judgment cannot be fully perceived but through celestial revelation. Which as I confess to be most true, so I think I hear some greater matter to be signified.

### Paragraph 1198

[note 85.5]Now follows the description of the judge, as of a noble and valiant warrior, composed of many parts. The faithful understand by this that the guardian, watchman, and avenger of the church does not sleep, whom the wicked do not perceive, failing to recognize the wrongs they commit against the godly, nor caring for those superstitious Christians, as they term them. They see moreover that they err if they think Christ at any time is too lenient, or overlooks for too long the calamities of his servants. For now he comes forth as a judge and revenger. There are many excellent descriptions of Christ in this book, as in any other; but this one is most elegant and vivid, which I have, according to my limited ability, expounded by parts. You shall always consider matters of greater import, until it is given to behold them directly with our eyes.

### Paragraph 1199

Our judge comes on horseback, and that on a white horse: [note 85.6] not that he needs the help of corruptible horses in heaven, but thus he speaks after the manner of men, that we might imagine greater things. Conquerors ride on white horses. Here is signified therefore, that our judge shall be a conqueror and a triumpher. Others suppose the white horse to be a sign of his most pure humanity. I understand rather the white cloud. For the same token he took up from the eyes of his disciples, when he ascended into heaven from Mount Olivet. In the same way he shall come again to judge. And just as kings are carried on horses and chariots, so the Psalmist ascribes to God clouds as horses and chariots.

### Paragraph 1200

Our judge is faithful and true. Faithful to his faithful. True in all his promises toward the godly and ungodly. They are deceived, and shall see themselves to be deceived at the judgment, so many as have contemned the promises and threatenings of God as vain, and esteeming things after the success of this world, judge the wicked to be happy and fortunate, and the godly to be wretched and miserable. Hereof hath the Prophet Malachi reasoned in the third and fourth chapters. And saying the judge is faithful and true, he judges and fights in righteousness, to wit, giving every man his own, rewards to the good, and punishments to the evil. This king doth not judge and fight as the kings of this world are wont, following vanity and corrupt affections. And Christ is said to fight when he rewards the ungodly after their demerits, the Apostle: we must all, says he, be manifest before the judgment seat of Christ, that every man may receive such things as he hath done by his body, according to that he hath wrought be it good or evil. 2 Corinthians 5.

### Paragraph 1201

3. The eyes of the Judge are like a flame of fire. For just as no one can escape or hide from a judge or judgment, as he searches the secrets of all, and nothing can be hidden from his sight, so are his eyes terrible and fearful against the ungodly. The godly, however, are filled with all pleasure, joy, and gladness by the sight of the Lord. Flaming and fiery eyes are also attributed to Christ in the first vision; see there for more. And Scripture testifies everywhere that the judge knows all things, even the secrets of hearts. Therefore, you act foolishly if you think you have won the battle and have sinned unpunished, having escaped the knowledge and judgment of men. Another judgment remains, in which all the deeds of the wicked will be revealed to their utter shame and confusion before all the world. The sins of the godly are covered by him through whose benefit they are justified and absolved from pain and guilt.

### Paragraph 1202

4. Our Judge has very many crowns upon his head, for he alone governs all realms and nations. As Daniel signified in chapter 7. He alone might truly be called Africanus, Europeanus, and Asiaticus, Parthicus, Persicus, Germanicus, Gothicus, and others. Which our kings have foolishly claimed for themselves, attempting to assume the monarchy, when Christ alone is the true Monarch forever. This Judge and mighty Prince shall strike off the triple crown from the head of the Bishop of Rome. Moreover, there shall be no king so mighty in the whole world that shall be able to resist him, and make war against him.

### Paragraph 1203

5. Our judge has a name written, which no man knows save himself. This will be explained more clearly in due course. Christ has a name unspeakable, for he is the true God, eternal, incomprehensible, and Almighty. This name knows no one but himself. For first, the majesty of God is greater than can be comprehended by any creature; again, the name of God is agreeable to no one but him alone, for the name of God, in this signification, may not be communicated. For he is very God, and besides him none, which thing Isaiah repeats often. He is the Savior, King, Monarch, and Judge: these things all belong properly to him alone and are not common to others. Moreover, the Lord himself says in the gospel: no man has known the Son but the Father; neither has any man known the Father, save the Son, and to whom the Son has pleased to reveal. Besides this, we see it imperfectly: and the glory of the divine majesty is so great, as I said before, that human capacity is unable to conceive such a glory. No man therefore save God alone knows his name.

### Paragraph 1204

6. The vesture of our Judge was sprinkled with blood.[note 85.12] This signifies victory and the slaughter of his enemies, which will be explained further toward the end of the chapter. He took this observation about our Judge from chapter 63 of *Explanations*. He alludes to conquerors returning from battle, whose garments and armor are soaked with the blood of the slain. And it betokens the just severity of the Judge and great slaughter of the enemies.

### Paragraph 1205

7. The name of the judge is now expressed, [?which] is utterly unknown to the ungodly. And the Judge is called the word of God. For the Son is the word and speech of God, the express mark of the divine substance: in whom the Father himself is expressed: and of whom, as of the word, we understand the will and mind of the Father, the true messenger of the heart. These holy words of the gospel are known: In the beginning was the word, [?and] the word was with God, &c. Therefore Christ, the word, was made flesh, the Lord God and Judge of all.

### Paragraph 1206

8. To the Judge is added an Army, not of angels only,[note 85.15] with whom he often repeated in the gospel that he would come in judgment: but of all the faithful, or saints, who at no time, not even here, are separated from their head. For at the sound of the trumpet, the archangel arises, and the saints arise, and the living along with the dead are changed and are taken up to meet Christ in the air. Here, here in the clouds and bright air, appear with Christ the happy and blessed victors. Immediately, the ungodly also rise, and those who lived at that time are changed with those who rise again, to pain and confusion. But they see the saints with Christ in heaven, in glory, and immediately feel unspeakable torments. These things, without doubt, come to pass and are fulfilled, as described in the 3rd and 5th chapters of Wisdom. Saint John therefore says that this army is in heaven, not on earth. He says how they follow Christ. For the same thing the Apostle also said in the first letter to the Thessalonians, the 4th. Moreover, he adds that they were clothed, and did not appear naked, and he expresses the kind of garment. They were clothed (he says) in silk, white and clean. For saints in Christ obtain righteousness and glory, are made clean and are glorified.[note 85.16] And this meaning Saint John himself has opened to us a little before, saying: silk is the righteousness of saints.

### Paragraph 1207

9. From the Judge’s mouth proceeds a two-edged and sharp sword, which cuts on either side. It is not sharp on one side and blunt on the other; it cuts on both sides equally. This signifies a just sentence pronounced by God’s mouth against the wicked. For against them, God’s sentence is a sword, piercing even to their hearts. Therefore it is also called sharp. The judgment of our Judge is straight and severe, but yet just and righteous. What that sword is, is declared in the Gospel: namely, that heavy and immutable sentence, “Get you hence into everlasting fire.” Matthew 25. Upon this it follows in the words of the Evangelist: that with the same he may strike the heathen, that is, to condemn and put all unbelievers to perpetual torments.

### Paragraph 1208

10. And he shall rule them with a rod of iron. In the same way, he says the same thing he said before. For those who would not receive or acknowledge with repentance the staff of instruction and discipline afterward, shall find judgment and feel the iron scepter, with which he will break them all to pieces, like a potter’s vessel. Neither shall any power resist or prevail against him. And this manner of speaking is taken from Psalm 110. For Saint John gladly uses the words of Scripture to the end to make his book more commendable or more pleasant and acceptable.

### Paragraph 1209

11. He treads the winepress of the wine of wrath. And again he says the same, that he did before [note 85.19] but now utters another parable, taken from the scriptures, namely from the sixty-third chapter of Isaiah. The effect of it, or some part of it, is that he will pour out his wrath upon the ungodly and punish them most severely with his almighty hand, to which all things give way, bowing their heads. See what is said here in the fourteenth chapter of this book.

### Paragraph 1210

12. Again, the name of this judge [note 85.20] is shown, and in the name is the majesty and power of all others greatest. He has the name written on his garment and on his thigh. By these is declared the true humanity of Christ, after which he is exalted, as the Apostle says in the second chapter of Philippians. And to him is given a name which is above all names. Here he is called King of Kings and Lord of Lords, very God, Lord, monarch, and judge of all men. For so do the other Apostles speak also in the second chapter and seventeenth verse of Acts. And there might seem in this name of the Judge, as it were a cause to be shown, wherefore he is here appointed judge over all, because he is King and Lord of all. To whom be glory forever. Amen.

## Sermon 86: ¶ The description of the judgment, wherein punishment is taken of Antichristians and ungodly.

### Paragraph 1211

.

### Paragraph 1212

And I saw an angel standing in the sun, and he cried with a loud voice, saying to all the birds that fly in the middle of the sky: “Come and gather yourselves together to the supper of the great God, that you may eat the flesh of kings, and of high captains, and the flesh of mighty men, and the flesh of horses, and of those who sit on them, and the flesh of all free people and bond people, both small and great.” And I saw the beast and the kings of the earth, and their armies gathered together, to wage battle against him who sat on the horse, and against his soldiers. And the beast was seized, along with that false prophet who had worked miracles before him, with which he had deceived those who received the beast’s mark and those who worshiped his image. These two were cast alive into a pool of fire burning with brimstone, and the remnant were slain with the sword of him who sat on the horse, a sword that proceeded from his mouth, and all the birds were filled with their flesh.

### Paragraph 1213

Immediately after the description of the Judge, and a vivid depiction, a description no less evident follows of the judgment: that is to say, how Christ, having vanquished his enemies, commits them to perpetual torments. And the Apostle uses a prophetic phrase and eloquence. For by a figurative speech, all fowls are called to the slaughter and feast, that they might be filled with the flesh of the slain. And first an allusion is made of such as slay wares and prepare a feast, whereunto they may call their friends, and make them cheer. Again an allusion is made to the murder and slaughter of enemies, whereof wild beasts and ravening fowls are filled. Neither is there any thing hereby signified, but Christ shall overthrow all the ungodly, and take punishment of the same. Before was set forth a supper for the godly, wherewith they are refreshed and fulfilled. Now is prepared a feast of the solemn slaughter, whereby the ungodly receive no commodity, neither are they satisfied, but rather are slain and devoured, that is to say, perish. For no man will imagine that the wicked shall be overthrown at once, and after wearied by wild beasts, and gnawed by fowls, and so all punishment gathered together. For so should their pain seem to be none at all. But by temporal parables, eternal things are figured. These are taken out of the Prophets, namely out of the sixty-sixth chapter of Isaiah, and the thirty-ninth of Ezekiel, where are read in a manner the same words when he covers and wraps the ungodly with calamities, I mean wars, and destructions, and with other torments as it were kills them: but chiefly, when at the last judgment he commits them to pains everlasting.

### Paragraph 1214

And this slaughter is declared by the angel, [note 86.2], standing in the sun and crying with a loud voice. This signifies that the day of judgment will be solemnly proclaimed with great fanfare and will be a remarkable day, so that no one can be ignorant of it, but all things concerning it will be heard by all people. Therefore, he cries with a loud and audible voice. And he gathers, as in Ezekiel, the birds and fowls to devour kings, and people of all ranks, ages, and genders: that is to say, that all these must be called together to suffer eternal punishment and destruction. Therefore, he lists diligently kings, chief captains, strong men, and people of all sorts, namely, persecutors of Christ, Antichristians, the ungodly, scoffers, and impenitent persons.

### Paragraph 1215

Here is shown the cause of the damnation and destruction of the ungodly, while their enterprise, endeavor, and attempt is shown. They are now assembled to fight a battle against Christ and his elect, that is, against the church, by the beast, the Kings of the Earth, and their armies, captains, and soldiers, fighting both spiritually and physically. And here no lengthy commentary is needed. Read the histories of the church for the span of these five hundred years. And see what is done today by Popes, Bishops, and Princes, and by their counselors and ministers. Parliaments are called, in them bloody decrees and laws are made against the evangelists. And a grievous persecution is attempted against the Gospel and church. The spiritual fathers cluster together, they consult, and cry out that a council must be called. Wherefore, I pray you? to the intent that the new doctrine (as they call it) of the Gospel may be cut up, rooted out, and plucked up by the roots. And therefore they stand always ready to fight, are many times assembled against God and his anointed, to fight with the Lord Christ and with his chosen. Now days, if at any time peace is concluded, and the most bloody wars of Princes are taken up, or in the composition of peace, or immediately after, consultations are had, how to oppress the godly. But there is no other kind of sin more heinous than to impugn the truth of the Gospel. And therefore this is the chief cause of the condemnation of the wicked at the same day. It is accounted at this day among the chiefest virtues of Princes, if a Prince will give no place to the preaching of the gospel, but will shut it out, defend and maintain the church of Rome, with those doctrines, rites, and ceremonies falsely called old. Such are called right and good Catholics, most Christian, and defenders of the faith.

### Paragraph 1216

But now their condemnation is certain, and the manner of it is described. And the beast is captured. He speaks of these well-known offenders, as if they were seized by the manner, and taken in fact beyond their expectation. For in the midst of their plans they are intercepted, while they are yet in great hope, and truly believe they will accomplish many things against Christ and his church; then, in the heat and wickedness of their attacks and persecution, they will be apprehended. Consequently, it is also clear that persecution of the truth will continue until the end of the world. But who will be taken? The beast and the false prophet, who performed miracles. Concerning these matters, see what is said at the end of chapter 13. The cause of eternal damnation is again presented here. For he has deceived the world by his enchantments and crafty trickery, by his decrees and commandments. Of this I have spoken before, more than once, as also in chapter 17 and following verses. The fellowship of his condemnation is also joined: all those who have received the mark of the beast and have worshipped the beast. Concerning this, we will now repeat nothing. These things are declared in chapter 13. And although nothing can or ought to be proven by the pictures, it is nevertheless certain that painters have borrowed from this source for their old paintings of the last judgment.

### Paragraph 1217

We see the old depictions of the final judgment, painted a hundred years and more, intended to represent and display to us a large crowd of priests, monks, and friars, and all kinds of spiritual leaders. Most notably, the spirits of kings and popes are shown being driven to hell, where they burn in its bottom with everlasting fire. Consequently, it is said that more priests than farmers go to the devil, and so forth.

### Paragraph 1218

Furthermore, the damnation itself, and the manner of torment, or the torment itself, shall be fire. For John says: these two are cast quickly into a pool burning with brimstone into a lake or standing water as is found in fen countries, for such is the description of hell and of extreme punishment, as is also described by Isaiah in chapter 30. Tophet was long ago prepared, and the same is prepared for the king, which he has made both wide and deep; the inner chamber of it is fire, a great store of wood, which the blast of the Lord, or stream of sulfur, sets on fire. Not dissimilar things are read in chapter 66 and in the gospel of Jesus Christ, Mark 9, Matthew 25, and in other places without number. And the plague of Sodom is known to all men, Genesis 19. The laughings therefore of the ungodly scoffers, the godly would rather believe these things than prove them. For here is hell set open as it were for us to look into it. Let us fear.

### Paragraph 1219

And two here are singularly named. For sins, they have been authors of all evils, and rightly must be chief in pains or torments. For the wise man also has said, terribly and surely he will appear to you, for it shall be a most severe judgment to those who bear rule, and the mighty shall suffer mighty torments. For Scripture also shows in another place that there are degrees of punishment, according to the quality of the crime. And let us not think here that the head is punished without the members. For the whole body of Antichrist shall be condemned to torments. All the ungodly shall be punished, as hereafter shall be more plainly declared, and in the end of the twentieth chapter.

### Paragraph 1220

But this also is especially to be noted, that it is said how they shall be cast quick into hell. For so is signified the resurrection of the dead. Here is signified that in the judgment the world yet remaining shall be taken in the flesh, not as yet dead but living, which Saint Paul expresses vividly. 1 Thessalonians 4. And we pronounce openly in the creed, saying: from thence he shall come to judge the quick and the dead: not only the just and unjust, but the dead, to wit remaining in the flesh, and living. Antichrist therefore shall live at the day of judgment, and shall not be extinguished before. The persecution of Antichrist shall endure, with all ungodliness, even to the last day. And just as Korah, Dathan, and Abiron, and the rest of the conspirators, were taken in the very crime of rebellion, and swallowed up quick with their tabernacles, and all their things of the earth opening, so at the day of judgment hell gaping wide shall receive and swallow up all the ungodly, but chiefly the Antichristians. The which many now do not believe, but in that day shall find it with unspeakable pains, and horror incredible, and all we shall see it with these our eyes. Every one puts trust in his own sect, and hopes to obtain salvation in his superstition. But the things that we here at this present are told us of the judge himself, Christ, as most certain and undoubted, and after a sort are set forth to behold.

### Paragraph 1221

And what shall be done at that judgment with the remainder of the ungodly and impenitent? Shall only Antichristians be damned for Antichristianism? John adds: and the remnant are slain with the sword, and so forth. For in Matthew 25, the judge declares: “Go into everlasting fire, prepared for the devil and his angels.” For I was hungry, and they gave me no food, and so forth. For if they are to be damned by the just judgment of God, who, when they might have done good to men, have not done it, what, I pray thee, shall become of those who not only have shown no generosity to the needy, but have moreover plundered those who had honestly and were generous of their goods, and afterward spent them in riotous living? And so they have brought those of honest substance into extreme misery, and by this means have robbed the poor also, who were wont to be helped through their generosity, of their help and succour? Here are also comprehended heretics, Jews, Gentiles, Mahometans, and all others like them.

### Paragraph 1222

At the end, it is repeated that all birds are filled with the flesh of the damned. We understand this to be expressed through figurative language, and it should not be interpreted literally, but rather as signifying that all ungodly and impenitent people will be punished most abundantly. Primasius, interpreting this passage, says that we should not understand it carnally, believing that the saints—which he explains as being represented by birds—are satiated with the flesh of the wicked. Instead, the equity of God’s judgment, revealed to the saints, shows how, in redeeming the full number of the elect, the remainder is decreed to be damned. They are said to be filled with this knowledge of righteousness, which a person may hunger or thirst for in this life but cannot fully comprehend. Isaiah also speaks of the ungodly: “And they shall be for the fullness of sight unto all flesh.” Here, I suppose, is the fullness of souls mentioned earlier. Immediately following, “and the fowls may be taken in the evil part,” referring to the angels who transgressed. After leading their followers to destruction and accomplishing their evil desires, they are said to be filled with the flesh of the condemned, taking satisfaction from their damnation, to which they were authors of errors. This is what Primasius says. However, while I do not disagree with these things, which are undoubtedly spoken truthfully, I believe they should not be examined so closely. What is spoken figuratively and in the prophetic manner seems to indicate nothing other than that all the ungodly will be destroyed by the great power of God and extreme torments. Therefore, let us fear God, to whom alone be glory.

## Sermon 87: ¶ Of the bright verity of the gospel, which by the ministry of the Apostles was spread abroad through out the whole world, and by a thousand years.

### Paragraph 1223

.

### Paragraph 1224

And I saw an angel come down from heaven, having the key of the bottomless pit, and a great chain in his hand. And he seized the dragon, that old serpent (who is the devil and Satan), and bound him for a thousand years, and cast him into the bottomless pit. And he bound him and set a seal on him, so that he should deceive the nations no more until the thousand years were fulfilled. And after that, he must be released for a little season.

### Paragraph 1225

[note 87.2]Saying that the greatest aspects of religion and true godliness consist in the true knowledge and understanding of the final judgment, as I often advise: Saint John diligently discusses the treatise on the final judgment to our great benefit. And following his customary method, to ensure that all that he proposes is clear, he not only declares the matter in words but also presents it through visions as if they are events to be seen with the eyes—and this is for the faithful. For to the unfaithful, all these things, although most godly and divine, seem mere trifles and fables. But the wisdom of God will laugh at them also when she sees her time, as she threatens in the Proverbs of Solomon. And he also resolves certain questions that are commonly raised about this matter.

### Paragraph 1226

And he said, how the beast with the false prophet and all his adherents should be cast at the last day into hell; but where the first part is neither Antichristian, nor yet Christian, but rather of their own sense and judgment, to be a rule and law to themselves—such as are the Nestorians, Jacobites, Georgians, and so forth—or those who are still heathens or gentiles, moreover Jews and Turks—some man might marvel and demand, what shall be done with them, or what will come of them? John answers: “And the remnant were slain with the sword of him that sat on the horse,” and so forth. Again, where a godly man might wonder how they should be condemned, who, being born among Turks, heretics, Jews, and gentiles, never heard the Christian truth? John anticipates this imagination and, at the beginning of chapter 20, shows with what majesty, clarity, and evidence the truth of Christ’s gospel was notified to the world; how also all force and power was taken away from the devil, and that for the space of a thousand years, wherein the preaching of the gospel thundered continually, so that those who have not received the gospel of Christ are utterly inexcusable. For the preaching of the gospel was not obscure, but most clear and manifest, nor short and contracted, but published for the space of a thousand years; it was not received by a few little ones, but by all people and nations under the sun. Therefore, it is a gross ignorance of the Turks, heretics, Jews, and gentiles. For although in times past the truth seemed to have been notably known, now it is not so; yet it is certain that the majesty of the gospel has been so great in the world that there is now mention of it with all men, and by their own malice they hide their eyes, which understand nothing of Christ. Therefore, that saying of the Apostle is even now of force: “If our gospel be hid, in those that perish it is hid, unto whom the God of this world hath blinded the minds of them, which believe not, that the light of the gospel should not shine unto them,” 2 Corinthians 4. Whereupon we now gather that none of those who are condemned in the world are condemned without deserving. Which thing the Apostle Paul also touched upon in the Epistle to the Romans in chapters 1 and 2. Here, therefore, is a profitable and necessary place treated, of the famous preaching of the gospel throughout the world, the course of which endured a thousand years.

### Paragraph 1227

And this treatise proceeds in this order. First, the angel is described, after which his work or effect is declared. And last is the sealing of the time. Concerning the description of the angel, first he is named an angel,[note 87.6] and comes forth abroad; however, the entire apostolic state is understood here, in which Saint Paul, the doctor of the Gentiles, shines exceedingly. It is not surprising that the order of the apostles is signified by an angel, for an angel signifies a messenger, ambassador, or an apostle. And therefore, the prophet Malachi called John the Baptist, the herald of our Lord, an angel: “Behold, I send my angel before thee,” and so forth. Ministers of the church are often called angels in this book. But if the worthiness and nobility of the name pleases the ministers, let the angelic purity and excellent faith please them also. An ambassador does and says nothing except that which he has received in commission from him who sent him; so also let the ministers set forth nothing except that which they have received from the Lord, in the Scriptures.

### Paragraph 1228

Secondly, this excellent angel is said to come down from heaven, not that the bodies of the Apostles came from heaven, but because their vocation and office were given to them from heaven. For the Son of God, who came down from heaven, chose the Apostles and sent them forth into the world. This is declared in Matthew 10:20 and John 20:21, Mark 16:6, and Luke 24:48. And Paul says to the Galatians that he was called and ordained an apostle neither by men nor from men, but by God through Christ. Thus it appears how great is the authority of the Apostles. For they do not speak, but the Spirit of Christ and of the Father speaks in them. Therefore, he who despises their doctrine despises God the Father and the Son. They lie moreover, who say that the gospel is a new doctrine forged by clever men. Read the first and second epistles of Peter.

### Paragraph 1229

[note 87.8]After this, the Angel is said to hold in his hand two excellent instruments, the key and chain. Let us see what is meant by them. Doubtless, by these two instruments John understands nothing else but the free, true, holy, and lively preaching of the gospel: by which it came to pass that both hell was locked from the faithful, and the devil was held and kept fast bound in chains, that he could not hurt the godly so much as he would, and seduce whom he listed. For so John will explain himself later.

### Paragraph 1230

And the keys of binding and loosing the apostles received from the Lord, in Matthew 16 and John 20. They open the bottomless pit and hell itself to the ungodly by preaching the gospel, when they show them their damnation in hell for their ungodliness. They shut up hell from the godly, while by preaching the gospel they open heaven and bring the faithful to celestial joys. I have spoken of the keys at length in another place. A chain is the sign of captivity. By the preaching of God’s word, the Devil is taken and bound, which is why common painters have depicted the Devil bound with chains to certain notable preachers.

### Paragraph 1231

The things that follow this key, and the significance of the chain, are better explained, as is the effect of the apostolic preaching. For he adds: “And he seized the dragon,” referring back to chapter 12, where everything is explained. You may also look there for the same information. And the angel bound Satan, which is the purpose of the chain: to prevent him from stirring up, invading, and destroying the faithful. Furthermore, there follows something even more forceful: “And he cast him into the bottomless pit,” that is, he hurled him headlong into the depths of hell so that he could not murder the faithful. Another, even more serious thing follows: “And he shut him up so that he could not come out again.” This is the purpose and effect of the key. Furthermore, he placed a seal upon him. Letters, prisoners, and graves are typically sealed for security and credibility, to prevent tampering and to ensure they remain shut, sealed, and safe. All these things signify a complete and perfect victory that we have obtained through Christ, by the word of the gospel preached to us and communicated through faith. For our sake, he overcame, defeated, bound, imprisoned, and sealed the enemy, so that we might be safe and secure from him. Finally, it follows—this may explain all the preceding points—that he should deceive the people no more, by the same means he had previously used to seduce them before Christ’s victory and before the gospel was preached throughout the world. For then, everything was full of ungodliness and error. Temples of God or idols were everywhere, idols were worshipped, and they gave oracles. Altars smoked with the blood of men and beasts. All wickedness reigned. Arts, magic, witchcraft, and parricide were practiced without punishment. It is difficult to express, even in a lengthy speech, how shamefully Satan had deceived the world, how assuredly he reigned, and with what complexity he had bound humanity to himself like a slave. Let anyone who wishes examine Greece, Italy, and Asia; Corinth in Greece, Rome in Italy, and Ephesus, the leading city of Asia. They will find abominable abominations and say that the devil had reigned completely there, constantly bewitching the simple and wretched with new deceptions. But after Paul alone (I will say nothing now of the other apostles) came to Corinth, Ephesus, and Rome, and preached Christ, who cannot see that John here truly saw the devil bound and securely shut up? I am brief in this matter, which is abundant, because I think I have done enough if I have shown only a few stepping stones, by which one may come to a much more thorough consideration of these things. Relevant are the divine words of Paul, which are read in Acts 26, spoken before King Agrippa, the princes of Syria, and Festus, the proconsul. “For this purpose I have appeared to you,” the Lord says to Paul, “that I might appoint you a minister and a witness both of the things which you have seen, and of the things which I shall appear to you after this, delivering you from the people and nations to whom I now send you, that you may open their eyes, that they may turn from darkness to light, and from the power of Satan to God, that they may receive forgiveness of sins.” In 1 Colossians, as well as in other places, Paul shows that Christ has overcome Satan, and that the same Christ has redeemed us and brought us out of the kingdom of darkness into the kingdom and light of the Son of God. Therefore, when the apostles and ministers are said here to bind and shut up Satan, it is to be understood through the means of their ministry. Each person can also judge whether he has profited from the doctrine of the gospel that he has heard for a long time in the church. For if you yourself are still bound by the devil’s chain, you have not yet heard the gospel as it pertains. But if you feel that the devil is bound by the chain, and that you rule the devil, and the devil does not rule you, then things are going well. Pray to God: “Lord, confirm and increase what you have worked in us.”

### Paragraph 1232

And concerning the time of this most brilliant truth of the gospel, [note 87.14] it is said how it shall endure in the world a thousand years. For he says expressly: he bound him for a thousand years. And again: that he should no more deceive the people, till the thousand years were fulfilled. I know that the opinions of the expositors, concerning these thousand years, are diverse. I know how the heresy of the Chiliasts or Millenaries, by Papias, the author of this view, as Eusebius recounts in the third book of the Ecclesiastical History, was derived from this. I will not here linger to refute the opinions of others, which would also be overly long and tedious, and not of such great profit. I will only express my own view to be weighed by godly readers, and then I will leave it free for every man to follow that thing which he shall think most agreeable to the truth and profitable for the faithful. And I understand plainly and simply that Saint John speaks of a thousand years, which ran on by continual course from the time of Christ, until the last corrupting of the Evangelical preaching and church of Christ.

### Paragraph 1233

I am not particularly concerned with precisely determining the length of these thousand years.[note 87.15] I simply set the beginning of the reckoning to the time of the open preaching of the Gospel, when the word began to be received by the Gentiles. Therefore, I suppose that there may be three periods appointed, which nevertheless will all fall within the same reckoning, differing little or nothing among themselves, or having only a small diversity, not exceeding half a year more or less. You may, if you wish, begin the calculation of the thousand years from the thirty-fourth year of Christ’s birth, when Christ also ascended into heaven, and Paul, being called to the ministry, began to restrain Satan by drawing the Gentiles into the fellowship of God’s people through the preaching of God’s word. And you will arrive at the year of our Lord 1034, and the papacy of Pope Benedict the ninth, who, after unlawfully seizing the chair of Saint Peter, as they call it, practiced magic and entered into league with the devil; of whom he was also carried away, as he had sold his bishopric before Pope Gregory the sixth. Read the story of Cardinal Benon, as mentioned previously in chapter 13, and let other stories be read. It is certain that the Devil occupied the apostolic seat at those times, as they term it. Read the stories from Sylvester the second and onward.

You will then say that about that time the Devil broke loose again and seduced the people, especially through Popes. Or begin the calculation of the thousand years from the time when Paul, being bound for the Gospel at Rome, testified that the Gospel had been preached throughout the world. That was about the year of our Lord 60. From that point, reckoning a thousand years, you will arrive at the year of our Lord 1060, when Nicholas the second was Pope, under whom it is written that the truth was diversely tempted and corrupted, and that Gregory the seventh also, by his schemes and endeavors, troubled the whole world. Or begin the calculation from the destruction of Jerusalem, when the Jews cast out the Gentiles, and a large number of Gentiles entered and were received into the place of the rejected Jews, which was the year of our Lord 73, and you will arrive at the year of our Lord 1073, even to Pope Gregory the seventh, in which time not a few historians write that the Devil himself reigned. Doubtless, no man has harmed godliness or more boldly advanced impiety than this Gregory, otherwise called Hildebrand. I have spoken of him before in chapter 13. There I also admonished you that Cardinal Benon calculated those thousand years from the birth of the Lord and concluded them in Sylvester the second.

It is evident, therefore, that the Gospel has held a notable place in the world and has not been extinguished for the space of a thousand years, that is to say, from the time of the apostles until the year of our Lord was reckoned 1073 or thereabouts. What was done at that time and afterward, we shall hear when we come to that saying: “And when the thousand years shall be fulfilled, &c.”

### Paragraph 1234

Some man will say, “I cannot see that the preaching of the Gospel has continued in the world for so long a time, namely a thousand years, for it appears by histories that the doctrine of merits, satisfactions, and justification by works, immediately after the time of the Apostles, laid their first foundations. We know that the intercessions of saints and the worship of relics were defended by Jerome, who departed this world in the year of our Lord’s incarnation 422. We know that the Bishop of Rome immediately after the death of Gregory the first, claimed to be head and catholic pastor of the universal church. We know that about the same time, namely about the year of our Lord’s incarnation 630, Muhammad seduced a great part of the world. We know that shortly after, arose that detestable contention about the having of images in the churches of Christians. We have heard that John has assigned years to Antichrist, 666. Finally, it is manifest that the Devil has by murder, parricide, and all kinds of mischief reigned in the children of unbelief. Wherefore you say, “I see not how the Devil has been bounden a thousand years and locked in chains.” [note 87.16] I answer that the things which are alleged hitherto are true; yet nevertheless, to be true and so to remain always, which John by the revelation of Jesus Christ has affirmed, that the Devil should be shut up for a thousand years and remain bounden until a thousand years were at an end. And the same we declare on this wise. The Lord said in the Gospel, “Now is the judgment of the world; now shall the Prince of this world be cast out.” And whereas it is not lawful to doubt of the verity of Christ's words, yet nevertheless, he is not to be understood as being so cast out, but that he has been of great force in the world and has been called by the Apostles themselves, the Prince of this world. How then is he said to be cast out, and to tempt the godly, to reign, and to be cast out of his kingdom?

### Paragraph 1235

He is cast out of the church and of the faithful,[note 87.17] not that he does not come again, for always he returns and seeks to pluck back the redeemed (but that he possesses no more the full dominion. For Christ now lives and reigns in the church and saints. These, as Saint Augustine says, he assails from without). He is cast out of his ancient possession, but he labors to recover his old habitation. And thus was Satan bound and shut up for a thousand years, as one who did not possess the faithful of Christ throughout the world, nor rule them at his pleasure and according to his malice, although he has tempted and vexed them. So was the Holy Spirit denied to be given, not that he was not in the world and in the Prophets, but because he was never so plentifully poured out upon all flesh as after the glorifying of our Lord Christ. In the same sense we say that death and sin are taken away from the faithful and trodden under foot. As Paul, who in the first chapter to the Colossians said that we are translated out of the kingdom of darkness into the kingdom of light, nevertheless says to the Corinthians that the god of this world has blinded the minds of the unfaithful. So John at this present says how the Devil is bound and sealed for the space of a thousand years, and the very same says afterward that the rest of the dead did not revive until the thousand years should be fulfilled: that is to say, in all those thousand years did not believe those who valued the beast more than they did Christ. And they, truly through their own fault and the instigation of the Devil, did not believe and perished. Therefore did Satan exercise his power in them. Which to the faithful in deed is bounden and tied fast, but to the unfaithful free and over familiar. Likewise, Hell is shut to the godly, to the wicked open. Wherefore also we confess in the creed, life everlasting and not death or everlasting damnation. For the faithful have no Hell, or there is no Hell prepared for them, but for the ungodly. For Christ has broken Hell, but for his faithful; to the unfaithful all things of Hell are yet most strong, and these have Hell.

### Paragraph 1236

[note 87.18]Again, the Devil is said to be bound, shut up, and sealed: for since the redemption of Christ, his power has not been so great in the world as it was before. Therefore John explains himself, and says that he should deceive the people no more. What is this “more”? But that he shall not so seduce them from now on as he has done hitherto. Therefore, although in the meantime he shall deceive some, in those thousand years he has not reigned so fully, safely, and freely as he did before, and as it is permitted him after those thousand years to rage. Therefore, these things are spoken by comparison, and not absolutely. And the thing itself, or experience, teaches that they are not to be understood absolutely, and after the bare letter. Although therefore Satan has in these thousand years also blown his poison upon many, and has troubled the world, yet this is nothing in comparison of those things that have followed after the thousand years, even until this day, and shall follow hereafter unto the world’s end. In old time also he reigned fully among the Gentiles through idolatry. But a thousand years fell down their temples and idols, with all other instruments of ungodliness.

### Paragraph 1237

We read truly, how there were in the time of the Apostles who affirmed that people are justified by the law and works. Whereupon sprang up the doctrine of satisfaction and merits. But the same doctrine was refuted by the Apostle Paul, above other Apostles. Saint Augustine also, and after him Bede, most consistently have defended the doctrine of grace and redemption by Christ. This doctrine remained secure for the space of a thousand whole years. But afterward, monks gaining the upper hand, the doctrine of satisfaction and human merits did prevail; whereupon the doctrine of Jesus Christ concerning the free remission of sins and the imputation of righteousness was utterly obscured. There grew up an opinion with certain of the saints, making intercession or praying in heaven for their worshippers. Relics began to be worshipped too soon. Nevertheless, those who were illumined clung fast to the only intercessor, Christ, and did not honor relics. But after those thousand fatal years, many attributed more to saints than to the very holy one of saints. We see what is done at this day. The writings of monks and friars testify how much the worship of creatures has increased within these four hundred years, or thereabouts. Who will deny that exceeding many have been deceived by heretics? But who can gather thereby that the Gospel has been utterly lost, and that Satan has reigned fully?

### Paragraph 1238

The Bishop of Rome has ascended into the top of Mount Zion, and will be called the head and general pastor of the whole Catholic Church. However, the East consistently resisted, and so did other parts of the world as well. At length, after a thousand years, he made his boast most impudently, that the fullness of power was given him, which he obtained by trickery and clamor, and afterward usurped it. Muhammad seduced many, yet nevertheless the patriarchal churches persisted, and the East honored Christ, likewise the South and North, so that the thousand years again had their light; neither has Satan in these raged so much as he has since those years were complete. Doubtless, since the Turks began to rule and reign, all matters of religion grew every day worse and worse. And the war into the Holy Land did very much hurt to religion, and gave great courage to the Saracens and ungodly people; whereof I shall speak afterwards. And images began to be set up in temples, and to be defended. But the histories testify that the same was done with great difficulty, and hardly could the use of them be obtained, all good men most constantly resisting. And when they were now admitted, yet were not the idolaters so stark mad as we see they are now, and have been for certain years past. Wherefore it is rightly said how after a thousand years, Satan should be loosed from his chains, which before also moved the unbelievers, yet finally rages more furiously.

### Paragraph 1239

John assigned to the Antichrist a certain number of years, namely 666. From this, we might understand the name of the Antichrist. But it does not follow that the Devil was then completely loosed, or the light of the Gospel utterly extinguished. For the Apostle, speaking of him in his time, says, “The mystery of iniquity now works.” Therefore, the Antichrist has his seeds, his beginning, his rising, his growth, and his increase. But after a thousand years, he began to operate most impudently and boldly, which he had also done before, but now he spews out his poison most venomously, oppressing kings and all who speak against him, even if only a little.

We know moreover that in these past thousand years, the Devil has reigned in many through murder, perjury, and innumerable and unspeakable evils. But if you consider what has been done since those thousand years, and what is done today, you will say that those ages of the thousand years were like golden and silver worlds; and now for these five hundred years they are like brass, iron, lead, and clay.

Lactantius, in the seventh book of *Institutions*, chapter 15, toward the end of the world, says that the state of worldly matters must necessarily be altered, and iniquity prevailing, to incline to the worse, so that these our times, in which iniquity and mischief have grown to the highest degree, yet in comparison of that incurable evil, may be accounted fortunate and, in a manner, golden ages. For justice will wear so thin, ungodliness, covetousness, stubbornness, and lust will be so common, that if there are happily any good men then, they will be a prey to the wicked and everywhere vexed by the unrighteous. And evil men only will be wealthy, and the good tormented in all vexation and misery. All right will be confounded, and laws will perish. Then no man will have anything save that which is either evil gotten or evil kept. Gold and violence will have all. There will be no faith in men, no peace, no humanity, no shamefastness, no truth. And the remnant which are read there. By all of which our days now seem to be painted vividly.

### Paragraph 1240

But what does it mean that the Devil must be loosed for a little season? Does this “little season” last now five hundred years? We interpret this passage similarly to the saying in the gospel: “Unless those days were shortened, no flesh should be saved.” For it is clear from history that neither the Devil nor Antichrist has enjoyed his kingdom without disturbance. In every age, holy and learned men, illuminated and comforted by God, like Enoch and Elijah, have resisted the ungodly and maintained the true religion. Through them, consciences afflicted by Antichrist have received comfort, God, in his mercy, tempering matters so that the elect should not despair in great temptations, errors, and darkness.

### Paragraph 1241

Therefore both Satan and the Pope could enjoy these matters for only a short time. For immediately after the thousand years sprang up the Waldensians, who constantly opposed the Pope and his ungodliness. The Lord has raised up certain kings, who are images—namely, the German Emperors Frederick, Lewis of the house of Bavaria, and many others. The Popes themselves have been at disagreement, while many have been chosen, and each of them will be the vicar of Christ, thus tearing apart that entire ecclesiastical body of theirs with schisms. Against these preachers arise Wycliffe, Huss, Jerome of Prague, and diverse others, with earnest and vehement opposition. What is done at this day, and has been happening for these thirty years and more, against superstitions and idolatry, against the Pope and all his clergy, the Papists themselves cry out, and all parts of the world can testify. Therefore the Devil has been cast down for a little season. The Lord Jesus will soon tread him under our feet.

## Sermon 88: ¶ What those thousand years shall be, and of the certain felicity of souls after the death corporal, and of the first resurrection, and second death.

### Paragraph 1242

.

### Paragraph 1243

And I saw thrones, and those who sat on them, and judgment was given to them. And I saw the souls of those who had been beheaded for the testimony of Jesus and for the word of God, who had not worshipped the beast, nor his image, nor received his mark on their foreheads or in their hands. They lived and reigned with Christ for a thousand years. But the rest of the dead did not live again until the thousand years were finished. This is the first resurrection. Blessed and holy is he who has part in the first resurrection. On such, the second death has no power, but they shall be priests of God and of Christ, and shall reign with him for a thousand years.

### Paragraph 1244

By these words John declares himself, [note 88.1] explaining what those thousand years will be. Not such, without doubt, as many imagine (among whom are counted the Millenaries or Chiliasts), who say that there will be tranquility on earth, and that during those years the saints here on earth will reign corporally with Christ in exquisite pleasures and joys. For John himself refutes this opinion, while he shows how the saints will be beheaded by the beast and his image, and that those who remain in death will not live again or receive the gospel of Christ. It is manifest, therefore, that the beast and his image will be present during those thousand years. It is evident that the gospel of Christ will shine so brightly during those thousand years that Satan will be so strictly bound that nevertheless not all will receive the gospel, nor will there be quiet tranquility; but that the saints, for Christ’s truth, will suffer persecution by the beast, and that many will not believe the gospel but rather resist it and perish. Yet that the Devil in the meantime will not have so great power as he obtained after the thousand years were finished, nor that the gospel will in those thousand years be so darkened as it was after it was corrupted and depraved. And he touches upon certain opinions that are right notable and necessary, and opens them, namely, what the state will be of those who either are killed for Christ or reject Antichrist; for their souls do not sleep until the judgment, but live with Christ in heaven. He treats moreover of the first resurrection and the second death. Thus, to those who marvel where the souls of the dead will become and what they will do immediately after corporal death, he answers, and declares as much as is necessary to know.

### Paragraph 1245

Therefore John sees thrones, and those who sit on them. And who are those who sit? He adds an explanation and says, “the souls of those who were beheaded.” For it is taken as an explanation, as if to say, those who sit on the heavenly thrones are the souls of those who were beheaded. Souls are not beheaded, but bodies: the souls remain in their state and life. Therefore, he says the souls of those whose bodies were beheaded or slain. And here let us note that John does not speak of the bodies [illegible] at the last judgment, changed, or raised again, but of the souls released from the bodies of the martyrs. For he speaks of souls loosed from the bodies, before the Judgment, according as every one in his time lives here in this world and is called from hence by death. For Arethas, Bishop of Caesarea, also explains this as the souls of martyrs; yet he does not think that no one should be saved unless he die by the tyrants’ sword. For he adds this moreover: or truly he names to be beheaded, tropically, those who have mortified their members on earth. Thus far he. And we also have shown before that first and chiefly the holy martyrs are rewarded with eternal life, secondly all those who have honored God truly and have done penance, and have crucified their flesh with all its concupiscences.

### Paragraph 1246

And he states explicitly that the saints were beheaded, not for theft, murder, or wrongdoing, as Saint Peter also teaches, 1 Peter 4. But for the word of God and the testimony of Jesus Christ. The word of God is the very Son of God our Savior, and the testimony is that wholesome gospel, and the very preaching and professing of the same, as we have explained previously through the examination of Scripture. They are also numbered among the saints those who have not worshipped the beast, and so on. And these are the martyrs beheaded or slain, because they worshipped God, but they would not worship the beast or its image. However, not all who reject Antichrist are slain, and therefore he particularly mentioned them as a distinct member. But what it means to worship the beast and its image, and to receive its mark, and so on, I have explained at length in chapter 13.

Now let us see what their condition is, those who shed their blood for Christ and abhor Antichrist with all his enchantments: they live, he says, that is, by faith in this present world. As Paul also said: I no longer live, but Christ lives in me. And from that same life follows everlasting life in another world. Therefore John added, “and they reigned with Christ for a thousand years,” meaning that entire span of time. Not that they did not reign and live with Christ afterward, but that their souls until the judgment have not slept, but rather have lived a blessed life in heaven. Which from the beginning he indicates by another vision. For he sees a seat set, and the souls sitting on them. And by a figurative expression he signifies that certain seats and honorable places are prepared in heaven for the blessed souls, as the Lord himself says in the gospel: “In my Father’s house are many mansions, and I go to prepare a place for you.” He calls the seats thrones, alluding to the royal thrones of kings. But of these heavenly seats, we must conceive of greater, divine, and spiritual matters. They sit not because they do nothing but sit on a cushion, but they reign, triumph, rest, live, and have enjoyment of the comfort, joy, and everlasting glory. This is the manner in which the souls and spirits sit.

He adds moreover that judgment was given to those souls, truly because they are exempted from judgment and do not come into judgment (even as our Savior says), but have passed from death to life. It is also declared in another place, in what sense the saints are said to sit upon the seats and judge the world: where it is evident that all judgment is given to the Son. It is evident therefore by this infallible passage of Scripture that the souls of the saints do not sleep after the death of the body until the last judgment, but live in heaven with Christ. But at the judgment they shall return to their bodies raised again, and together with their bodies shall be received into blessed seats. And this is the condition of the faithful. From this hope let us never allow ourselves to be withdrawn. In my Decades I have discussed more at length about the souls separated from their bodies and have shown that they do not sleep.

### Paragraph 1247

And here I cannot refrain, but must set forth and recite that which Dr. John Funceius, a learned and diligent man, who has read much in the tenth book of his Chronology, under the year of our Lord 1332, wrote in these words: About this time, the most holy father, Pope John, the twenty-second of that name, fell into this heresy, which he also openly professed and taught, that souls do not see God before the last day. For so his father had taught him, deceived by the visions of Tantalus, which were commonly carried abroad in writing. And Pope John sent two preachers to Paris, namely, a couple of Friars, one of the order of preachers, and another Minorite, to profess his error there. But one Thomas, a preacher of England, resisted the Pope stoutly, whom the Pope committed to prison. And the King of France called a synod in his palace, in the forest Vitinian, where all those assembled subscribed against the Pope. Then the king sent ambassadors to the Pope, exhorting him to recant his error and that he would deliver Thomas out of prison. Which enlarged the prisoner; and also (as it is said), following the admonitions of his friends, at the hour of death repented. So much Funceius. It is a shame therefore for some, who at this day, in so great a light of the gospel, dare renew that most foolish error, affirming that souls separated from their bodies lie snoring, I know not in what dormitory, neither to feel anything, till at the day of judgment they be joined again to their bodies, and rise again.

### Paragraph 1248

John adds: and the remnant of the dead did not live again until the thousand years were accomplished. Not that they lived afterward, but that they revived never at all. As Scripture says in another place, Michal, David’s wife, remained barren until the day of her death: not that she had a child after her death. But whom does he mean by the remnant of the dead? Surely all those who descend from Adam are dead. As Paul declares well in chapter 5 to the Romans. But we have heard how some through faith have received Christ, and so being quickened, have shed their blood for Christ, and would not worship the beast or his image. Now is added to this number: but the remnant of the dead, who are neither regenerated through faith nor would bestow their lives for Christ, but would rather worship the beast and his image, these, I say, lived not because of their unbelief. For without faith there is no true life in this world.

We speak nothing here of the vital or natural life. And we say that life is double, or of two sorts: the one spiritual, which is of faith and of the Spirit of God, and of Christ, which is received by faith and lives in the hearts of his people, and his life in them. For the Lord himself says: “He who eats me shall also live for me.” The other life is everlasting, that is, of another world, in which we shall see God as he is, and shall be as he is, living in God and with God forevermore. Conversely, death is of two sorts: spiritual, whereby lacking Christ and his Spirit, and void of faith, we live in sin. The Apostle, speaking of this death, says that a widow living wantonly is dead while alive. And the Lord also to the disciple who wanted to return home and bury his parents, says: “Let the dead bury their dead.” There is also everlasting death, that is, everlasting wretchedness and misery, which follows the spiritual. Yet see what we have said of double death in chapter 3 of this book, in expounding the Epistle to those of Sardis. Therefore, John here signifies that there will be many in these thousand years who will not receive the gospel with a lively faith, and therefore will remain in death, as the Lord said in John 8. Therefore, they err shamefully who suppose that all nations in the whole universal world will come one day to an unity of faith and most assured peace in this life.

### Paragraph 1249

And Saint John himself again explaining his own words says, “This is that first resurrection.” Which I ask you, by which men receive Christ through true faith [note 88.8] and rise from sin into a new life. The Apostle speaks much of this in Romans 6. He also speaks, in Ephesians, from Isaiah: “Awake, you who sleep, and rise from the dead, and Christ shall shine upon you.” Therefore, those who do not acknowledge their sins, nor are regenerated, nor are quickened by faith in Christ, nor rise again with Christ in a new life, are not partakers of the first resurrection. The second resurrection is that universal resurrection of all flesh, in which all men will indeed arise, but in different states: the faithful to everlasting life, the unfaithful to everlasting death. The Lord himself also repeated this from Daniel chapter 12, in John chapter 5.

### Paragraph 1250

And he demonstrates, by example and in an apostolic manner, a threefold fruit or effect of the first resurrection. First, he says, “Blessed and holy is he who is partaker of the first resurrection.” He is blessed, meaning happy and heir of celestial and eternal life. Holy: that is to say, purified, sanctified, and justified. For faith in Christ does sanctify and make blessed. Therefore, in those who are thus sanctified, the second death has no place nor power. And the first death is the death of sin; therefore, the second death is eternal damnation. See what I have spoken hereof before in chapter 2 of this book, in the Epistle to the church of Smyrna. Finally, the faithful are made priests of God and of Christ, the elect, I mean, segregated, notable, excellent, both of God and Christ most dearly beloved, who in eternal life might offer eternal praises to God. It is repeated again, and they shall reign with him for a thousand years. And this signifies that all saints shall reign with Christ forever, but chiefly the souls, even also before the judgment.

### Paragraph 1251

Primasius, Bishop of Utica, explaining this passage, says that it is not spoken only of bishops and priests, but just as we call all Christians by reason of the mystical chrism or anointing, so are all priests, for they are members of a Priest, of whom the Apostle, Saint Peter, says, “an holy people, a royal priesthood.” Thus he says. But this whole passage concerning the binding and loosing of the Devil, the thousand years, and the first resurrection and second death, Saint Augustine has well and diligently, for his time and for so much as he could see, discussed at length in Book 20 of *The City of God*. I propose these things of mine to be diligently considered by the faithful. Let every man hold that which he shall think most consonant to the truth. To the Lord our God be praise and glory, now and evermore. Amen.

## Sermon 89: ¶ What shall be done when the thousand years are expired, of the world deceived, of war and grievous persecution of the godly, and of the everlasting payne of the wicked.

### Paragraph 1252

.

### Paragraph 1253

And when the thousand years are expired, Satan shall be loosed out of his prison, and shall go out to deceive the nations that are in the four quarters of the earth, Gog and Magog to gather them together to battle, whose number is as the sand of the sea. And they went upon the plain of the earth, and compassed the tents of the saints about, and the beloved city. And fire came down from God out of heaven and devoured them. And the Devil that deceived them was cast into a lake of fire and brimstone, where the beast and the false prophet are and shall be tormented day and night forevermore.

### Paragraph 1254

He declares here what will happen after those thousand years. And he states primarily two things: that the devil will be released from his prison, so that he may deceive the people in the world, and may assemble Gog and Magog for battle. To this he adds two more: a most severe persecution of the church, and punishment of the wicked, and everlasting damnation of the devil and his followers.

### Paragraph 1255

And the seduction of the world must again be explained by the figure of synecdoche. [note 89.2] For the meaning of the scripture will not allow us to understand that there should be no godly people at that time. For we believe that there is a church, and an holy church, and that it shall always exist in the world until the judgment. And we have heard moreover in this book how many thousands are sealed so that they should not perish. And also that the dragon must be loosed for a little season. Just as we read in the gospel that Satan is cast out and his kingdom taken from him, nevertheless Saint Peter warns and says that the devil goes about like a roaring lion, and seeks whom he may devour. Verily, for the greatest force of Satan is to infringe the faithful, by Christ that mighty champion and noble conqueror, the Devil notwithstanding going about and aspiring again to the empire, and to be restored to his former place. So at this present we understand that Satan, loosed after those thousand years, ranges now abroad more freely, exercises greater authority, seduces more people in the world, and rules further than he has reigned these thousand years. Yet so that there shall nevertheless in the world be a fellowship of Saints dispersed and vexed miserably. For immediately Saint John says that the beloved city of God is besieged by the enemies. Therefore the church shall be in the midst of the enemies. Wherefore all that same place must be explained not of the truth and religion wholly extinguished, but of the more large and ample power and seduction of Satan, the old serpent.

### Paragraph 1256

Therefore he says that when the thousand years shall be expired, the Devil will be released from that prison into which, through the power and might of Christ or the preaching of the Apostles, he had been shut. For when the bonds are broken—namely, the sincere doctrine and preaching of the gospel corrupted and depraved—he comes out. And for this purpose he comes out, that he might deceive the Gentiles, that is to say, all people and nations dwelling in the four quarters of the earth, meaning the whole universal world. And to the end he might allure Gog and Magog, namely fierce men, barbarians, worldly, mocking and despising true religion, addicted to robbery and given to evil things, and regarding only corruption and wickedness, that I say, he might draw such men to unrighteousness and keep them still in error. For such does Ezekiel signify Gog and Magog to be. But those who, through divine grace, are not such, shall not be deceived by Satan; but grounded in Christ, shall persevere in the doctrine of the prophets and Apostles, and shall rightly worship Christ, shall abhor Antichrist, and all wickedness in the world.

### Paragraph 1257

But a deceptive falsehood has passed through the world far and wide, since the thousand years expired,[note 89.4] as experience teaches, and histories of times testify. For it is plain that during those thousand years, there were famous churches of Christ in the East, which nevertheless we lament to have been destroyed within these five hundred years. Therefore, the wicked and abominable sect of Mahomet began six hundred years after the birth of Christ, and from that time forth was advanced by the Saracens, but prevailed at last after those thousand fatal years. For how great is the power of the Turks now in Africa, Asia, and Europe, no man is ignorant. And Papistry had its beginning and progression rather soon; but after a thousand years it was of full force. For the Bishops of Rome through the abuse of excommunication have oppressed even the most mighty Emperors and Kings. For who knows not with what shameless boldness the popes have withstood Kings and Emperors, Henrys, Fredericks, Louis, and many other Princes, whom their lewdness has vexed, vanquished, and overcome? After much and grievous contention, the Popes extorted to themselves the consecrating of Bishops. They usurped moreover the church goods also, by which (such a force has money) they might do in the world what they listed. For by this means Papistry received strongest sinews. Moreover, after those thousand years was raised up and established that God Maxim,[note 89.5] of whom also Daniel makes mention, which brought also a great strength unto Papistry. I mean transubstantiation, and the horrible polluting of the Lord’s Supper, and manifold abuse of the holy mysteries. And of the force hereof increased an infinite number of priests and filthy Friars. For after those thousand years at the length came up the sect or order of Jacobines, Celestines, Gilbertines, of Gray friars, black friars, white friars, and many other friars, and monstrous Monks, which have craftily crept into the favor of all princes, to the intent they might know all their secrets by auricular confession. Then began all things more impudently to be set forth and sold in the church than ever before. Superstitions and unprofitable and hurtful ceremonies overflowed. For we have seen thirty years since and more, how much increased daily idols and idolatry, worshiping of creatures, and abuses innumerable about the same, pilgrimages to dumb idols, and an infinite number of the same sort. I recite not that holy matrimony waxed now vile after those thousand years, in so much that ministers of churches were prohibited to marry. Then waxed whoredom rife, rape, and adultery, and yet more filthy things than all these, &c. I pass over here very many things; this only I rehearse, if you compare the rites, ceremonies, and superstitions of Papistry with the heathen gentility (as I have partly showed here and there in my works), you will say that Papistry passes far all gentility. For in case the false opinion and persuasion once taken away, you do way what Papistry is in itself: you will grant that there was never such a corrupt thing in the world. Full rightly therefore says Jerome, that Satan is broken loose out of prison.[note 89.6] By which proverb he signifies matters extremely corrupted, nothing to be done in his place or decent order, but all things confused, all turned up side down, at the will and lust of the evil spirit.

### Paragraph 1258

Hereunto is added another thing, that when the thousand years expire, Satan should gather Gog and Magog to battle. By these words, John has undoubtedly alluded the prophecy of Ezekiel, which we read in chapters 38 and 39. Ezekiel seems to have prophesied of the wars of Macedon and of Antiochus, speaking of them with a prophetic phrase and a hyperbolic amplification. The prophet says that Gog is the land of Magog. And it is evident that Magog was Japheth’s son, who dwelt at Mount Caucacus and extended his empire to Ethiopia and Egypt. And afterward, out of Asia and from the eastern parts, Antiochus Epiphanes made war on the people of God. This was a figure of Antichrist, as all commentators confess. Therefore, it appears that John brings forth these things by way of comparison, as though he should say: just as in past times the people of Gog and Magog severely molested and afflicted the people of God, so in the times of Antichrist, most grievous wars shall arise, with which the church of God shall be shaken and laid waste. He says truly that the host of these destroyers shall be innumerable. He adds, in the manner of Scripture, a parable for clarity: as the sand of the sea. And also, by another phrase of speaking, he signifies that the enemies of God’s people shall be bold and ready to overrun the whole world and to turmoil all things with wars. For he says: And they went upon the plain of the land. This means that, being swift and bold, they shall run over the whole world; everywhere and throughout the wide world, there shall be cruel wars.

### Paragraph 1259

For most purposely he adds:[note 89.8] and they surrounded the tents of the Saints and the beloved City. And means that the church of God shall be most grievously plagued with those Gogicall and barbarous wars. For in times past Jerusalem was called the chosen and beloved City; but after she rejected the word of the Lord, she was no more beloved of God, but rather rejected and hated. Therefore Saint John speaks of the Catholic church, which Saint Paul also in another place out of Isaiah names, Jerusalem that is above. The same is also called the tents of Saints. For the faithful are in the church as it were in tents, fighting against Satan, the world, sin, and flesh. And where he says, they surrounded the tents of Saints, he says something more than if he had written, they assailed or besieged or assaulted the tents of Saints. For they surround them, which give the assault round about, and vex them most grievously, as though they were already taken, that no hope can appear to any man, no refuge or way to escape.

### Paragraph 1260

Without a doubt, if we compare these things with histories, we will find that the church has been many times assailed with cruel wars; but never yet with crueler than after those thousand fatal years. I mean the holy war, as they term it. William Archbishop of Tyre has written at large about it, as has the Abbot of Uršpurge in Chronicles. Also Benedictus Cottes and Paulus Aemilius in the fourth book on the events of the Franks. Finally, Volaterranus in the eleventh book of Geography, concerning Coele-Syria and Palestine.

### Paragraph 1261

Historians report many things of the battle of Troy. Others suppose that those of Assyria and Babylon were greater. Many extol the wars of the Persians and Macedonians, as in truth they were terrible. The Romans also have their Punic, Mithridatic, Civil, Cimbric, and Germanic wars: but I believe that the war they call Holy was more cruel than all these, more bloody and grievous, and of longer duration. In this war, great battles were joined, with countless numbers of men, involving almost all nations and peoples of the world. Wonderful and monstrous slaughters were made. More than a hundred thousand men died, a number beyond belief. It continued for many years, more than any previous war in the world. Furthermore, it was waged with the most hostile intentions. And what is most relevant to this discussion is that the Oriental Saracens, Turks, Egyptians, Babylonians, and other barbarian nations were inflamed in this war, harboring an unquenchable hatred against the Christian religion, seeking to uproot it, and succeeding in uprooting a large portion of it, and they continue to do so to this day. This war, therefore, was the most grievous of all, and was the cause of the persecution of the faithful in the East and West. To note something of this, and to recount for those ignorant of history, it is clear that under Pope Gregory VII, there were many and famous churches in the East, and Patriarchal churches still secure. But while this Pope, above all others, dealt wickedly against Christ, the Son of God, and his holy church, just as we read in the time of Solomon, who, after he had revolted, brought forth many enemies against him, and most cruel ones: so in the wicked and tyrannical reign of Gregory VII, Solyman the Turk invaded Antioch, at which time the Emperors of Greece were said to have been dispatched from the East. And the Turks, advancing, are said to have invaded and harassed first the straits or ports of the Caspian hills and the country of Armenia, around the year of our Lord 764. There is now no time to speak of this in detail. After Solyman, Belchiaroke, the Turkish prince, whom others call Belzet, succeeded him, and he also invaded Greece itself, despising the Emperors of Constantinople. Alexius, who then was Emperor, is said to have sought aid from the Westerners against the Turks. And also one Peter, an Hermit (who certain historians blame most grievously, and not without cause), coming out of the East and traveling through the West, cried out an alarm. Urban II, whom some call Turbane, and a disciple of Gregory VII, called a great council at Clermont in France, where he proposed the question of recovering the Holy Land and delivering the Lord’s sepulcher from the hands of the Infidels. That council reminds me of what is described in the eighth book of Kings, chapter 22, under Ahab and Jehoshaphat, concerning the recovery of Ramoth-Gilead from the hands of the Syrians. For there was also a deceiving spirit there, there were Ahabs, there were Jehoshaphats, and many other similar things. To avoid making many words, a journey was decreed against the barbarian infidels of the East. This was done in the year of our Lord 1095. In the meantime, Peter the Hermit stirred himself and gathered certain thousands, whom he led through Hungary into Asia. Immediately after, followed the unfortunate captains Folkemar and Gottschalk, priests, who, by the way, destroyed everything with fire and sword, and were slain. At last, Godfrey and Baldwin, most noble princes, with certain excellent captains and noble warriors, with an innumerable multitude of men, were transported into Asia: which they say was done in the year of our Lord 1096. And within four years, or at most three, they had taken by assault or surrender the cities of Nice, Heraclea, Tarsus, Antioch, and Jerusalem. The Abbot of Urspurg reports that so much blood was shed in the city of Jerusalem that in the very temple itself, the horses stood up to the knees in the blood of the slain. The same man tells of a notable battle fought at Ascalon, in which about fifteen thousand foot soldiers and five thousand horsemen of Christians overthrew and routed Solyman of Babylon, equipped with one hundred thousand horsemen and four hundred thousand foot soldiers, and that more than one hundred thousand men were slain in that battle. And this journey of Godfrey was the first among the worthy voyages of Syria or Asia.

### Paragraph 1262

Following this expedition came others, better equipped. For while the victory and good fortune of those who first went to the East were greatly praised and commended throughout the West, William, Prince and Duke of Poitiers, filled with great hope, led approximately one hundred thousand soldiers into the Eastern lands. The year of our Lord was reckoned 1101. But of so great a number, scarcely one thousand are recorded to have returned home safely.

### Paragraph 1263

3. After, in the year of our Lord 1147, through the urging of Bernard of Clairvaux, Louis King of France, and Conrad King of Germany, and Frederick Prince of Swabia, took their journey into the East, leading with them an army almost innumerable; but the same army perished almost entirely, with scarcely the princes surviving.

### Paragraph 1264

In the year of our Lord 1189, when the city of Jerusalem was taken by the Sultan, the King of Persia—a city the Christians had held for only about 89 years—Emperor Frederick, nicknamed Barbarossa, King Philip of France, King Richard of England, and other powerful princes raised a very large army of Christian people to recover the city and the Holy Land. They were fortunate to transport their army into Asia, but afterward suffered great misfortune. For Emperor Frederick was drowned, and the entire army, as Urſpurgenſ testifies, died of the plague.

### Paragraph 1265

5. The fifteenth (and that famous indeed) voyage in Syria involved the mighty Kings Philip of France and Richard of England, nicknamed “Lionheart.” This occurred in the year of our Lord 1191. However, they returned without accomplishing anything of significance, and with a considerable loss of men.

### Paragraph 1266

6. And Palmerius, a chronicler, says that Henry reports the son of Emperor Barbarossa sent an army into Syria, which returned the following year. Consequently, the Christians, lacking assistance in Syria, lost all the dominion they had retained. He speaks of these things in the year of our Lord 1198.

### Paragraph 1267

7 Again, in the year of our Lord 1213, Pope Innocent III of that name sent public letters to all the faithful of Christ, urging them to take up arms against the infidels who possessed the Holy Land. [Note 89.14] If any man has leisure and desires to read the letters, he will find them in the Chronicles of Ursperg. And then, in the year of our Lord 1215, he held a general council in Lateran, where war was decreed against the Easterlings. Also, Honorius III, about the year of our Lord 1217, addressed and confirmed the same thing. Consequently, many Christian princes met at Acre, formerly called Ptolemais, and waged a deadly war upon the Easterlings. In this conflict, they captured the noble city of Damietta. Yet neither the outcome nor the fruit justified such great enterprises, costs, perils, and losses.

### Paragraph 1268

8 Therefore Frederick II, the emperor, hoping to do some good, marches also with an eager and well-equipped army into the East; this they say occurred in the year of our Lord 1234. In the meantime, while he valiantly fights in the East, the Bishop of Rome, Gregory IX of that name, taking an occasion (I use the words of Ursperger) of the emperor’s absence, sent a great army into Apulia and seized the emperor’s lands. This was done while the emperor was absent serving Christ (which is most wicked to say), and kept them subdued for his own use, and by no means would allow those who had taken the holy cross (that is to say, those who were going to war for the emperor) to take shipping or passage, but kept them under his power, both in Apulia and in Lombardy. And more such matters, which those with leisure may read in the same [? chronicle]. Therefore, the emperor, constrained, leaving his affairs there unfinished, was forced to negotiate with the enemy and returned, that he might recover what the Pope had taken from him.

### Paragraph 1269

And not long afterward, namely in the year of our Lord 1248, Louis, King of France, with his brothers Robert and Charles, and a very powerful army, sailed into Syria. There Robert was slain, and Charles was captured by the Sultan, and was scarcely released at the last, with great difficulty.

### Paragraph 1270

King Louis of France, in the year of our Lord 1270, embarked with his three sons from Marseille to sail into Africa. A plague struck his army in enemy territory, by which both the father and the sons died, and the entire army suffered an exceedingly great calamity.

### Paragraph 1271

And again, although they experienced misfortune in the wars against the Barbarians, it was nevertheless considered again in the council of Lyons under Gregory X, about the year of our Lord 1273, concerning the recovery of the holy land. But Palmerius in the year of our Lord 1291 states that many thousands of Christians were slain in Syria by the Saracens, and all the rest fled the country in fear. And the Chronicle of the Kings of France states that Aemilius ended the holy war (namely, in the year of our Lord 1291) when Ptolemais in the East was destroyed by the Sultan. It is manifest therefore that this barbaric and gogical war has lasted about 195 years—such a long time as I know no other war in the world that was ever waged with such obstinate minds, with so great armies, and so much shedding of human blood. We see in the meantime the tents of saints, and the city of God beloved, namely, the faithful church throughout the world, especially in the East, and also in the West, to be most grievously afflicted, and more than oppressed and destroyed, with only a few small remnants remaining; so that not without cause we may perceive what the Lord said in the gospel: but when the Son of Man shall come, will he find any faith on Earth?

### Paragraph 1272

The most holy and wise Prophet of God Daniel seems to have foreseen and prophesied all those things, as he did all the rest concerning Antichrist. After speaking at length about the power of Antichrist and the worship of God against the Apostles' institution, he adds in the 11th Chapter. And in the time of the end, to wit the end of the world and last judgment approaching, shall set upon him, namely upon Antichrist, the king of the South, and the king of the North shall fall upon him like a whirlwind, with chariots and horsemen, with a strong and great navy, and shall invade his realms. He shall overflow with armies, to wit innumerable, and he shall pass through, that is to say, he shall overcome all like a conqueror, doing what he lists. For we have perceived that the armies sent into the East by the counsels and motion of the Bishop of Rome have molested by sea and land the Turks and also the Soldan of Babylon and Egypt. What will they say that Daniel, pointing as it were with his finger, adds? He shall come also into the chosen land and invade the land of desire, namely Judea, which some time was called the chosen, delectable, and pleasant land. And many shall fall in the war verily, that shall be made for the recovery of the holy land. It follows in Daniel, these shall be delivered out of his hand, Aedom and Moab, and the Princes of the children of Ammon. For those nations are not known to have been so destroyed as the rest were, by the Saracens and after by the Turks, for that they framed themselves to them in time. Daniel annexes, and he shall lay his hands upon realms, neither shall the land of Egypt escape. For it is evident that the same also was possessed of the Soldans, princes of Babylon, and of the emperors of the Turks. It follows, and he shall have the rule of the treasures of gold and silver and all the precious things of the Egyptians. By which the prophet has signified the inestimable treasures and riches, and excellent majesty of the Soldans and Turkish Emperors. All which things, even so as the Prophet has said, experience proves to have been and as yet to be fulfilled. The Prophet adds, finally, the Libyans and Aethiopians shall be in his journeys. This translator has turned, “He shall pass also through Lybia and Aethiopia,” or as others have translated it, “they shall be in his way.” And he means that those regions shall be open to those barbarous Soldans and emperors of the Turks, by league, vicinity, and amity. Jerome expounding this place says, “when Egypt was taken, those lands were also afraid.” Wherefore he says not that he took them, but passed through Lybia and Aethiopia. Whether sense of these soever thou choosest, thou shalt not err, as I think, from the truth. And Daniel adds the disturbance from the East and from the North shall trouble him, so much that he shall go forth in a great fury to destroy and kill many. Jerome shows must be understood of Antichrist. The Pope of Rome affirms that the Patriarchal seats are subject to him, as Jerusalem, Antioch and Alexandria, and the holy land to be his right. And he hears from the East and from the North that all those parts are possessed of the Soldans and Emperors of Turks; he calls therefore great councils and decrees war against them. He hears moreover that Constantinople is taken, that the Rhodes is won, Dalmatia subdued, Bulgaria and Hungary vanquished, and so on. Again therefore he summoneth councils, he arms kings, he leads forth soldiers, he moves war, and decrees that war shall be made for the recovering of the holy land and to root out the Turks. So verily this Gogmagog war is not yet ended or appeased at this day. Whereby it comes to pass that an infinite multitude of men are slain on either side. Furthermore at the end of this Prophecy, the prophet shows, and as it were with his finger, the palace or seat of Antichrist, by Antiochus figured before, lest any man should not know where Antichrist were to be found. And he shall plant, or set the tabernacle of his palace betwixt two Seas: to wit the Adriatic Sea, called now the gulf of Venice, and the Tyrrhenian or Tuscan Sea, in the mount of desire of holiness, that is to say in the pleasant and holy hill. We have heard certainly that the palace of Saint Peter is preferred both before Mount Zion and also Sinai. There sits the most holy, in the seat of holiness. There is most full remission of all sins. There is the mother and supreme head of all churches. There is the high court and judgment, from whence may no man appeal. There sits the king of kings and high Bishop, which so far excels in brightness and majesty the Emperor and other kings, as the sun doth the moon and stars. There is thought to be perfect holiness, and all the treasures of Christ and of his Saints. Therefore said Daniel rightly, that Antichrist shall dwell in the noble and holy hill, namely in the seven hilly Rome, as we heard also in the 17th Chapter. Finally, he prophesies also of the end of this most puissant prince, Antichrist, and says: and what time he shall come to his end, no man shall help him. For Christ coming to judgment, shall thrust him out of his seat. And Daniel in the 12th chapter following describes the judgment. To Christ alone be glory.

### Paragraph 1273

Let us therefore proceed to add a few things concerning the torments of the ungodly, and the everlasting condemnation of the Devil and his followers. John says, "And fire came down from heaven and devoured them." And the prophet Amos, in chapter one, calls God’s vengeance fire, as others do also. Therefore, John signifies that the vengeance of God will fall upon all the enemies of the church. In times past, fire coming down from heaven burned up Sodom and Gomorrah, and also consumed the enemies of Elijah. And although fire does not always fall from heaven in a physical sense, the persecutors of the church will never escape punishment, because they have afflicted Christ’s saints. Surely, if we observe what happened in that holy war, and what happens daily, we will say that God’s vengeance is most present both against the Turks and the Papists. But if anyone understands that about the end of the world fire will rage and consume the wicked, as Peter also mentions of fire and burning, 2 Peter 3, I will not contradict it.

### Paragraph 1274

Finally, he also addresses the everlasting damnation of Satan and all his followers. For where the Lord said in the Gospel, "If the blind lead the blind, both shall fall into the ditch," it follows that both Satan, the deceiver, and the people he has seduced, will be carried together to hell. There, John places, as it were, joining together the devil, Gog and Magog, the Saracens, Turks, briefly all deceived nations, the Beast, the false prophet, and all Antichristians. We see therefore that God’s judgment is righteous, which he now returns to describe. And we admonished before, by this speech, that they shall be tormented day and night, and so forth, to signify the perpetuity of damnation. From this, may the Lord our God deliver us: to whom be glory forevermore. Amen.

## Sermon 90: ¶ The Judge, and last judgment is described, with the resurrection of the dead.

### Paragraph 1275

.

### Paragraph 1276

And I saw a great white throne, and he who sat on it had a face from which the earth and heaven fled away, and their place was no more found. And I saw the dead, great and small, standing before God, and the books were opened. Another book was opened, which is the book of life, and the dead were judged according to what was written in the books, according to their deeds. And the sea gave up the dead who were in it, and death and hell delivered up the dead who were in them, and they were judged every man according to his deeds. And death and hell were cast into the lake of fire. This is the second death, and whoever was not found written in the book of life was cast into the lake of fire.

### Paragraph 1277

John had begun to speak of the universal and final judgment, around the end of the eleventh chapter. And he resumed the same to be finished in the nineteenth chapter, where we heard that Antichrist should be cast down from his throne and glory into hell. A question then arose concerning those who, although they do not cleave to Antichrist, are also not joined with Christ: what will become of them at the final judgment? Having addressed that question and demonstrated the equity of God’s judgments, he returns, as it were, with an after-song to the description of the general and final judgment, describing it more generally now than he did in the nineteenth chapter, where he seemed chiefly to have treated of the destruction of Antichrist; yet so that he also showed, in a manner of speaking, what would happen to the other ungodly.

Now he handles the same judgment more generally, showing that all shall be judged herein, and sets it forth wholly as if it were painted to be seen by our eyes. For in his customary manner, he explains all this matter through a heavenly vision, so that he might not seem only to tell the thing to our ears, but also to show it forth to be seen by our eyes, so that it might be more deeply imprinted in our minds. And all these things are most certain and undoubted—as I also admonished you before—revealed by the judge Christ himself. But the judge and Lord himself can be ignorant of nothing concerning this matter. Nor can we perceive that John has hitherto been deceived or misled in anything that he has set forth to us, but has hit rightly all and singular points, as we see that testimonies confirm his prophecies to be fulfilled. Why, then, should we so much as doubt one of such things as are spoken of the judgment?

Therefore let us credit these things, and not be among the mockers whom the Apostle Peter prophesied would come and say: “Where is the promise of his coming?” Doubtless this matter is of greatest importance, the foundation and root of our faith. Here are explained to us not a few articles of our sincere and catholic faith, chiefly these: “I believe that Christ shall come to judge the living and the dead; I believe the communion of saints, the resurrection of the flesh, and life everlasting.” Let us therefore be diligent in hearing and marking these things, lest we be accounted among those who hear without any fruit the mysteries of the kingdom of God; but let us rather prepare ourselves to meet the judge, so that we may, with the wise virgins, enter with the bridegroom to the marriage and everlasting joys.

### Paragraph 1278

And the description or demonstration of this vision chiefly concerns: what the judge will be; who will be judged; how they will be judged; of what sort the resurrection of the dead will be; and of everlasting damnation; finally, who will be properly damned. These things I shall, in order and according to the grace that God has given me, declare as plainly as I can.

### Paragraph 1279

We have previously understood what kind of judge there will be; at present, he is shadowed by certain marks. These things align with the vision described by Daniel in chapter 7. By the way, we see again how this book contains testimonies of the prophets, from whom it is commended to us, just as John explains the prophets to us. John sees a throne, and it is white and great. For the judge himself said that he would come in glory and majesty, namely with great light. And we also believe that his judgments are righteous, just, and white. Arethas, an interpreter, says that the throne is great because he who sits upon it is great, of whom the prophet said, “Great is the Lord, and great is his power,” etc. And in the throne, as judge of all and supremely righteous, he sits, furnished with all power and virtue. For all of this signifies the word “sitting.” Those who are to be judged stand; he sits. Therefore, he calls him who sits, as one would say, “judge.” No other name does he give. But we believe that all judgment is given to the Son, and that he is appointed judge over all. John therefore sees, and also shows us to behold the Lord Jesus Christ coming in the clouds of the air, a righteous and mighty judge. Paul also in 2 Timothy calls him a great God: not that there is one great God and another lesser God, but that the majesty of our Lord Jesus Christ will at that day be most evidently seen, and the Lord himself will then reveal himself to the world with greater glory and power than ever before.

### Paragraph 1280

The same shall appear also most severe and most just. Whereupon John says figuratively, from whose face both heaven and earth fled away. For if those things which have not sinned dare not come into the judge’s sight, but seek as it were to save themselves by flight: where, I pray thee, shall the ungodly and summer [?]? And doubtless the prophet Malachi also: who, he says, shall abide the day of his coming? Or who is able to stand when he shall appear? So in the sixteenth chapter we heard that heaven fled back and was folded up like a scroll, that the mountains also and islands flitted, and that kings and princes and other men hid themselves in caves, saying to the hills and rocks, “Fall upon us, hide us from the face of him that sits on the seat, and from the wrath of the Lamb,” and so forth. By which words although a desperate conscience out of corrupt doctrine is described, yet the same shall appear chiefly in this judgment, when the severe and most righteous judge shall appear. A similar figure is read in the eighteenth Psalm. Where it is added, “and their place was no more found.” It is annexed to amplify the matter, not that heaven and earth shall be nowhere, but so much as they dare not (which is spoken by a figure) appear in the judgment of God. By all these things therefore is signified that the ungodly, being destitute of all counsel, shall not at that day know whither to turn themselves, or what to do, but trembling and despairing to be vexed with unspeakable torments before the seat. It might be thought in the meantime that John signifies this also, how heaven and earth should at the coming of the judge be renewed. Which also the Apostle Peter more plainly expresses in the third chapter of the second Epistle, which nevertheless refers and applies all those his sayings to the same sense that we have touched before. For he says: “Seeing then that all these things shall be dissolved, what ought you to be in holy conversation, looking for and hasting the coming of the day of God?” Arethas of Caesarea says that the flight of heaven and earth signifies no changing of place (for whither should they flee?), but flight and flitting from corruption to incorruption, and the last coming of the Lord, under which this mortal body of ours shall put on immortality, and the face of the earth shall be renewed. This, he says, is a like phrase of speech as is had in the twelfth of the Apocalypse, of the angels cast down out of heaven: neither was their place found any more in heaven.

### Paragraph 1281

Now he also touches those who will be judged, truly the dead. For he says, “And I saw the dead.” And shortly after, we shall hear that the dead will be raised up. Therefore, those who rise from the dead will be judged. Nevertheless, the living are not exempted, as the Apostle most manifestly states in 4.1 of 1 Thessalonians. But at this present, he names not the living, but the dead, for the resurrection of the dead is more hardly believed; and it is more easily believed that those remaining in flesh will be judged at that day. And truly, souls never die, bodies die. Therefore, where it is said here that the dead will be judged, we mean that all those who are dead at that day will come in their own bodies to the judgment of Christ. And all men must be judged. Wherefore, John sees great and small: that is to say, men of all sorts, status, sex, and age. Kings and princes are not exempted, the common people will not escape, neither children nor old folks, men nor women. He sees them all standing before the face, or judgment seat, of God. The guilty, or those accused, or to be accused, will be set before the judgment seat of God. And Paul also testifies expressly of this matter: “We must all appear before the judgment seat of Christ, that every one may receive in his body according to what he has done, whether it be good or evil.” 2 Corinthians 5. But they will appear in a diverse manner, both good and evil. For the wicked, as guilty, are brought to be judged and punished, and that their guiltiness may be openly known to all creatures. The good, for as much as they are justified and acquitted, and have no more guilt or crime by reason of Christ’s satisfaction, appear in judgment with glory, ready to judge the ungodly after their fashion and manner, and not to be judged by any. And this thing is singular: he says that we shall be judged in the sight of God. For who can appear in the sight of the terrible God, and fire consuming all things, save he that is purged with the blood of Christ? And what shall we think can be hid or escape the sight of God, knowing all things?

### Paragraph 1282

John also declares how the dead will be judged: he says that books are opened, and another book is opened, and so on. Therefore, by the books, and then by the book of life: that is to say, by the things written in those books, the dead are judged. Scripture ascribes to God the ways of humans, by which humans are accustomed to write down reminders for themselves, lest they forget things; but with God all things are one and always present; he neither forgets nor remembers. Nevertheless, Scripture attributes both to him. God is said to forget when he does not help or punish; again, he is said to remember when he does help or punish. In Malachi, the ungodly say that God has no care for human matters, neither does he care for the godly, nor yet punishes the wicked. But immediately an answer is made: then those who feared the Lord spoke to every one of his neighbor: the Lord gave ear and heard, and a book of remembrance was made in his presence, and so on. Therefore, their books opened, that is to say, the secrets of all people brought to light or made manifest, the Lord will judge whatever has been thought, said, done, or left undone. The books also of consciences (for the conscience stands in place of a thousand witnesses) will be opened in judgment, God revealing and judging all things. For Paul, speaking of the Gentiles, says that they show the work of the law written in their hearts, their conscience also bearing witness, and their thoughts accusing one another, or also excusing, in that day when the Lord will judge the secrets of men, according to my gospel, through Jesus Christ. And these are indeed the books which will be unclosed in the judgment. From this it appears that the judgment will be done with great speed; neither will every man be reasoned with by all, by books written to make the judge weary, as the ignorant might imagine.

### Paragraph 1283

But what is that unique book of life,[note 90.11] which will also be opened at the judgment? Reference is made to the book of life in chapter 3. Briefly, the book of life contains only one entry: whoever believes in the Son of God has eternal life. Therefore, we are judged according to what is written in the book of life. Those who believe are saved; those who do not believe are already judged, that is to say, are certainly damned.

### Paragraph 1284

And for as much as faith shows itself by works, incredulity also, hidden in the heart, betrays itself by works: therefore John adds immediately, “according to their works.” For in Scripture, man is likened to a tree. And the tree is judged by its fruit, whether it be good or evil. A tree has a growing life, which in Latin is called *anima vegetativa*, and a nature or disposition, bringing forth fruit according to its nature and kind. But that soul, that *anima vegetativa*, and that good disposition, bringing forth in us good fruit, that is to say good works, is a lively faith in Christ, where it is, there the man is regenerated and has a good disposition: therefore he cannot, because of his good disposition, but bring forth good fruit. Therefore, after our works we shall all be judged. For the judgment must be open and manifest: but faith does not appear, but in works. For it is the gift of God, and is of itself invisible, namely, a sure trust in the promises of God. And it is seen in works. However, it does not follow from this that men are justified by works also, and not by faith alone: but that by works faith is declared, which purifies and justifies, that afterward we may be able to bring forth the works of righteousness. It follows that in judgment no pretense, no hypocrisy shall be allowed. For many say they believe, which declare their faith by no good works. We learn hereof that no book shall be of force at the last judgment, save the books of God, or the books of conscience, wherein God writes with his finger: finally, the book of life written by God before the worlds were made, through his divine predestination, whereby he has predestined us, that he might adopt us for his children by Christ Jesus. And the rest, which Paul recites in 1 Ephesians. Therefore, the hurtful books of Jews, Christians in name only, and Turks, such as the Talmud, decretals, and Alcoran, shall perish. These shall be of no force at all in the judgment.

### Paragraph 1285

Now he returns to the dead, of whom he had made mention before, and lest any man should say: how shall the dead be judged, who were drowned in the sea, who were swallowed up by fishes, and devoured by wild beasts, who were consumed with fire, or in the earth, were brought into dust? he anticipates and declares that the bodies of the dead will rise again, and being so restored will come to judgment, and says: and the sea gave up the dead that were therein; that is to say, those who had perished in the sea. And by these words also he has touched upon the manner and means of the resurrection of the dead, and has sent us also to Genesis 1. The manner of the resurrection is God’s omnipotence, as Saint Paul also witnesses in 3 to the Philippians. For God by his omnipotence raises up, and calls those things that are not, that they may be. If this thing seems to you new or impossible, behold the beginning of things, and from that estimate the small restitution. Was not the sea or water from the beginning? But is it written to have had any fishes from the beginning? None at all. But God commanded that the water should be replenished with fish. And did not immediately, at God’s command, all manner of fishes appear, where before there was not one? What marvel is it then, if God in the end of things, commands the sea, and other elements also, to yield again their dead, and they obey their maker? Verily the Lord in the gospel says that those who are in their graves also [note 90.14] shall hear the voice or command of the Son of God, and shall rise again. The bodies moreover of those who die are turned for the most part into the same elements from whence they were taken out. There is that putrefies in the earth and is converted into earth. There are some consumed with fire. There are some that perish in water. Some hang in the air and are there consumed. But at the Lord’s command, by whatever kind of death they perish, they shall rise again to judgment whole. Arethas, also Bishop of Caesarea, perceived this and said: he recites these things to the intent he might declare what the final and universal resurrection shall be. For where many, not believing that the same shall be, say that it is by no means possible to be in those bodies which have been long corrupted and brought to that point that they are not at all, this sermon now corrects this, saying: Like as the bodies, when they were not, began to be, not by a certain chance or of themselves, but of the four elements, namely of Water, Fire, Air, and Earth: so also being reasonably returned again into the same, may be of the same composed again.

### Paragraph 1286

And for a further declaration he adds again: [note 90.15] and death and hell gave up those who were in them, dead. For he understands by death, any kind of death, as though he should say: death itself restores to the Judge and judgment, whomever, after what sort soever he has dispatched. Death therefore is feigned to be as it were a person, which holds the dead in himself, or in a prison. And hell has yet but a few bodies (for some we read to have gone down to hell quick) but the souls of the wicked. The same return to their bodies, that the whole man may be judged, body and soul. Others by hell, after the Hebrew phrase, understand a sepulture or grave. Again is repeated, that the whole man shall be judged body and soul, after every man works.

### Paragraph 1287

Thus much hitherto concerning the resurrection of the dead,[note 90.16] of which in our books elsewhere we have treated more at length. Finally, there follows a discussion of everlasting damnation, and who are properly condemned. And Hell, he says, and death are cast into the lake of fire. Of this, it has been spoken before. And Hell here signifies not the place of punishment, but those who are inhabitants of Hell, namely, whose souls are yet detained in hell, or appointed thither. Death also signifies those who are dead in sin, and those who, from spiritual or temporal death, go straightway to everlasting death. Whereupon is immediately annexed: This is the second death, by which truly those who are dead to Christ are devoted to perpetual fire, and those who live to Antichrist and the world. Others interpret these things to mean that after the judgment the Saints shall neither be buried again nor die. Paul affirms this also from Hosea in 1 Corinthians 15. Arethas and Primasius agree with us. For Arethas says: and he calls death and hell those who have committed things worthy of punishment, as fulfilling the number of the second death. And Primasius, by these names, says he, he signifies the Devil (because he is the author of death and pains in Hell) and also the whole fellowship of devils. For this is the same that he spoke more plainly before, by way of preventing: and the Devil, who deceived them, was cast into the lake of fire and brimstone. And that which he added there more obscurely, saying, and the beast and the false prophet, here is stated more plainly. So much Primasius. And who knows not that the members must follow the head, that all the ungodly follow the Devil, the head of all ungodliness?

### Paragraph 1288

And he expresses this most clearly, indicating who will be condemned to eternal fire at the final judgment: those whose names are not found in the book of life. Therefore, only those who are faithful in Christ, those predestined to eternal life, will be saved. All others, regardless of their religion or the kind of life they have lived, however righteous it may seem, will perish. Some interpret these words as referring to those who remain alive at that day. For we believe that the Son of God will judge both the living and the dead. Whether living or dead, it is certain that no one will be saved except through faith in Jesus Christ; all others will be damned. This is the ultimate fate of the good and the evil. To Jesus Christ, judge of all and redeemer of the faithful, be praise and glory forevermore. Amen.

## Sermon 91: ¶ That the world shall be renewed, the Saints glorified and made blessed: and what that felicity shall be, and howe certain.

### Paragraph 1289

.

### Paragraph 1290

And I saw a new heaven and a new earth. [note 91.1] For the first heaven and the first earth had vanished, and there was no more sea. And I, John, saw that holy city, New Jerusalem, coming down from God out of heaven, prepared as a bride adorned for her husband. And I heard a great voice from the throne, saying, “Behold, the tabernacle of God is with men, and he will dwell with them. They shall be his people, and God himself will be with them, and will be their God. And God will wipe away every tear from their eyes. There shall be no more death, neither sorrow, neither pain. For the former things have passed away.” And he who was seated on the throne said, “Behold, I am making all things new.” And he said to me, “Write this, for these words are true and faithful.” And he said to me, “It is done.”

### Paragraph 1291

I warned you concerning the beginning of chapter 15 of this book, that a fifteenth part of this work begins at chapter 15 and deals with God’s righteous and just judgments. And since God’s judgments are of two kinds—when he punishes evil according to their wickedness, and rewards the good with rewards—I said that this passage consists of two parts. First, I said that John abundantly describes the torments to be inflicted on Antichrist and all the ungodly. Secondly, he describes rewards, especially at the end of the world, to be given to all the saints. For we have often heard in this book that the souls separated from the body are, immediately after bodily death, taken up into everlasting life, but that the greatest felicity of all awaits the faithful at the end of the world, when the bodies are raised again and receive the rewards of everlasting glory. This passage is treated throughout chapter 21 and the beginning of chapter 22. Just as in the earlier part he has shown hell as if wide open, and the everlasting torments as if they could be seen presently, so in this later part he unlocks, as it were, or opens heaven itself, so that with the eyes of faith, we may see what hope and glory awaits the saints. And with all, the article of our faith is most clearly explained: “I believe in everlasting life.” And again, for greater clarity, he declares these things through a vision, which others note as the seventh and last. Therefore, all things are figured spiritually, not carnally, to be understood and taken. Doubtless, the matters are excellent even when understood literally; however, we must think of spiritual matters, and always consider them greater than human speech can express. For we know, as taught by the doctrine of the Prophets and Apostles, that it is always true that what no eye has seen, nor ear has heard, nor has entered the heart of man, is what God has prepared for those who love him. 1 Corinthians 2

### Paragraph 1292

And the chief points of this passage are these. First, he shows that the world shall be renewed. Secondly, he signifies that the saints shall be glorified and blessed. And he declares in general what that blessedness shall be. Immediately thereafter, he confirms these things with many reasons; moreover, he describes the place, the court, and palace of the blessed, and likewise the glory and blessedness of the saints. Which in the beginning of chapter 20, he finishes exceedingly well, under the figure of a river and tree of life. And just as he has for the most part borrowed all his things out of the books of prophets, which John also with his Apocalypse illuminates, so has he at this present borrowed these from chapters 65 and 66 of Isaiah, and chapter 37 of Ezekiel, and the last chapters of the same.

### Paragraph 1293

Of the renewing of the world he speaks plainly,[note 91.5], as does also the Apostle Peter in his second Epistle, chapter 3, that all things truly should be purged by fire, and not wholly abolished and annihilated, but should certainly be purified from all corruption: for Arethas says that he does not signify the extinguishing of the creation, but a renewing for the better. Therefore, John says expressly that he saw a new heaven and a new earth, and he adds by explanation that the first heaven and the first earth have vanished away: that is, they are changed in their qualities, so that the corruptible things are now gone, created for corruptible uses. For even so the sea is no more, also certainly subject to corruption, but changed into something better. Augustine and his colleague Primasius suppose that the troubled state of the world (signified not seldom in the scriptures by the sea) about the end of the world shall cease. Reading chapter 17 of the 20th book *De civitate dei*, he reasons likewise at length about this innovation of the world, in the same 20th book *De civitate dei*, and chapter 18 and other places. I think it best in this matter to put away all curiosity: and if any hidden thing appear therein, that it be reserved until that day, in which we shall see all things evidently. And I suppose that these things concerning the renewing of heaven and earth are not spoken so that there should any place be prepared for us, which we should inhabit again in these lower parts under heaven (for we believe that we shall fly up into heaven and go meet the Lord in the clouds, according to the doctrine of the Apostle, 1 Thessalonians 4), but so that our minds are thus confirmed, that the faithful shall undoubtedly be renewed and glorified. For if heaven and earth, made for man, be renewed and purified, who will doubt now that men themselves shall be most chiefly clarified?

### Paragraph 1294

Consequently, now John declares [note 91.6] that the saints will both be renewed and glorified and placed in blessed seats, and he signifies generally what the glory of the saints will be. Afterward, he will declare more fully and individually all those things most diligently. For he hears an angel saying: “Come, I will show you the bride, the wife of the Lamb,” and so on. He now names this figuratively a city, and that city is truly holy, the new Jerusalem. And a city signifies both the place and habitation, as well as those who dwell in the place, that is, the citizens themselves. This city is therefore not only the place of the blessed, but also the very communion of saints, previously prefigured in the city of Jerusalem. But he puts a great difference between this new one and that visible, physical Jerusalem [note 91.7]. For he calls ours holy; that other in the land of Palestine was profane, polluted with the blood of Christ, prophets, and apostles, and for the same reason utterly destroyed. Ours is also called new, for the communion of saints will be renewed at the same time. Therefore, it follows in interpretation, “coming down from heaven,” not that the habitation of the saints after judgment will again be on earth, but that the glory and renewal will be granted from heaven by divine majesty and power. As James is also reported to have said, every good gift and every perfect gift is from above, coming down from the Father of lights. And Paul also, in Galatians 4, said that the free church is the heavenly Jerusalem. The same in 1 Corinthians 15: The first man is of the earth, earthly; the second man is the Lord himself from heaven. Those who are earthly are like that earthly one; those who are heavenly are like that heavenly one. And just as we have borne the image of the earthly man, so we will also bear the image of the heavenly man. Therefore, John rightly said that the church of the saints comes down from heaven, that is, from heaven receiving its glory. For again, by demonstration, he says, “prepared by God, as a bride adorned for her husband.” The apostle, in 2 Corinthians 5, says, “We know that if our earthly house, this tabernacle, is destroyed, we have a building from God, a house not made with hands, eternal in heaven.” And immediately afterward: “He who has prepared us for this is God.” He removes all corruption from his saints, but gives and teaches them to be purified with all gifts of the body, so that they may be worthily adorned and dwell in the everlasting bridal chamber with their bridegroom, Christ. Therefore, this adorning consists in the abolishing of all corruption and mortality, and in the gift of incorruption, immortality, and glory. Paul, the apostle, also speaks of the purifying and adorning of the bride in Ephesians 5. In this world, the purging and trimming begin, and finally, at the end, it is finished most perfectly. For then the church will have neither spot nor wrinkle, all corruption truly wiped away, and all glory received. And here learn by the way that the saints are prepared by God; therefore, salvation is of mere grace.

### Paragraph 1295

And he proceeds to declare yet more plainly,[note 91.8] what the glory shall be: a subject on which he has been occasioned to speak more often than once. Blessedness chiefly consists in two things. For God will give unto his Saints all that is good, and will take from them all evil: and so shall these for ever enjoy the ultimate good, and the most perfect felicity, and shall lack all pain and misery. Saint Augustine in the end of his book *City of God*: “How great,” he says, “shall that felicity be, where no evil shall be, no good shall lack?” And this declaration of eternal felicity has its parts, by which it is made manifest. For first, a voice, and that a great one, cried from the Throne: “Behold, the tabernacle of God is with men.” The conjunction of God with holy men was in time past prefigured by the Tabernacle of Witness, whereby God testified that he would be in the midst of his people. And the same he will at the end, after the judgment, fulfill most abundantly. And therefore that voice adds:[note 91.9] “and he will dwell with them, and they shall be his people, and God himself with them, and will be their God.” The which Saint Paul seems to have uttered more succinctly and briefly, “and God shall be all in all.” For whatever is good, whatever is fair, whatever is pleasant and delectable, whatever the mind of man can imagine to be wished for, briefly, whatever pertains to the true and perfect felicity and blessed life, that same great God almighty will be wholly, and will show in himself most fully. And just as all and singular men enjoy to a pleasant satiety the amiable brightness and some heat of the sun, yet the sun nevertheless loses nothing by it: and although all men use the sun in common, each nevertheless enjoys it as proper and peculiar; right so in another world we shall use that eternal light, and joy everlasting and unspeakable. Whereof more plentiful things shall follow shortly.

### Paragraph 1296

And then, just as God in himself gives to the glorified all goodness, so will he remove all evil from them; so that they will not only be delivered from calamities, but those calamities shall never return, nor be feared again. This he declares most abundantly through words taken from the oracles of the prophets. “God will wipe away all tears from their eyes,” he says. He used this kind of speech also in chapter 7, truly taken from chapters 25 and 65 of Isaiah. And David also in Psalm 126: “They who sow in tears shall reap in gladness.” He seems to have alluded to mothers, who wipe the eyes of their tender children, comfort the sorrowful, and cherish them hurt or bruised. Therefore, if the saints have suffered any pain or grief in this world, it will be requited to them, provided that they shall feel no more adversity. The Lord also said in the gospel: “Verily, verily, I say unto you, you will weep and lament, but the world will rejoice; and you will be sorrowful, but your sorrow will be turned into joy,” and so on, in John 16.

### Paragraph 1297

Consequently, he declares even more fully by naming the calamities, that the saints in the next life will be delivered at once from all evil: and death will be no more. For they will be rewarded with eternal life. Therefore, there will be no more fear of death, which is in a manner more bitter than death itself. The same affirms the Apostle in 1 Corinthians 15, citing the testimony of the prophet Hosea. There will be no mourning nor sorrow, which dries up the bones, even though they are full of vitality. For the joy of the saints will be perpetual. There will be no clamor, no complaints, no questioning or grumbling. For there will be no injustice, no malice, or envy. This world sounds and reverberates full of the clamor and cries of poor wretches. But in the blessed realms there will be no misery. There will be no pain, labor, sickness, or weariness. The cause of this is that the old things are gone. There is now another life, indeed, a very different manner of living from that which we live now. Therefore, whatever is of sin and subject to corruption will be taken away, as the Lord said in the Gospel,[note 91.11] “The children of this world marry and are married: but those who are accounted worthy to attain to that other world and resurrection from the dead will neither marry nor be married. For they can die no more; for they are made like angels, and are the children of God, because they are the children of the resurrection.” Luke 20. But of eternal life we have spoken more in our commentaries on Matthew 12. And the Lord himself in John gathers the sum of all things and says how he makes all things new. Therefore, in the world to come we shall think of no carnal or corrupt thing, but a heavenly one.

### Paragraph 1298

But the minds of the faithful are grievously tempted in this matter, the Devil suggesting that the hope of the faithful is vain: and that is a foolish thing, to despise good things present and certain for an uncertain glory. There are innumerable others of the same sort, which come to the mind of man, and trouble and shake the faith of eternal life. The Lord, therefore, the faithful Physician of his people, lest they should feel any hindrance in this behalf, confirms these things strongly and in many ways: declaring the hope of the faithful to be most certain, and all things to be undoubted, which are or shall be taught of eternal life, of the felicity and glory of saints. And he places this assertion as it were in a bright ray, after he has certainly collected the sum of felicity, whereunto by and by he will add fuller things after the vision exhibited.

### Paragraph 1299

It is to be understood that the certainty of the blessed life is shown most clearly of all by these words: “And he said to me, write,” and so forth. Nevertheless, no weak reasons for the truth can be gathered from the preceding words either.

### Paragraph 1300

First, indeed, he says: “I, John, saw.” And we know John to be an apostle and witness of the truth, whose testimony it is unlawful to distrust. Saying therefore this godly man saw the things himself which he recounts; to doubt of them would be a wickedness.

### Paragraph 1301

Secondly, he hears a voice, and a great, resounding pronouncement coming from the throne, namely, from the twenty-four elders and angelic beings, and the entire heavenly host. And who can doubt their testimony, who already exist in everlasting blessedness? They know, and have experience of what felicity is; therefore they speak and testify that which is tried and known.

### Paragraph 1302

Moreover, he himself who sits on the throne speaks and testifies, saying: “Behold, I make all things new.” God is true, and in him there is no deceit. And saying this so plainly, he testifies that eternal life will be: and we see him declare also of what sort it will be. No place for doubtfulness is left hereafter.

### Paragraph 1303

And the things that he has shown and declared concerning the blessed life, he commands immediately to write. Things are written for a perpetual memorial of the thing, which we know to be true and substantial. For writings or testimonies which are written or made and sealed, by the law of nations and common custom of men, have the force of an undoubted testimony. But such letters or testimonies are made and sealed at the commandment of God. For God commands John to write those same things which are taught of the blessed life, and therefore they are true, undoubted, and infallible. As he himself immediately adds and says: “For these words are faithful and true, steadfast I say, and immutable.” What could be spoken more evidently than these? Here is also the authority of holy Scripture established. But he adds another thing almost more forceful: and he said to me, “It is done.” By this manner of speaking is signified either that the end has come, and all things are accomplished, like as it is used in chapter 16, or else that the thing which is spoken and believed to come, is so certain as though it were done already. We Germans, when we wish to signify that the thing which we have proposed, or promised and said, is sure, we are wont to say, *Es ist gemacht*, “it is done.” Let us therefore believe assuredly these and all God’s words. Moreover, let us give our Lord God most hearty thanks, which with so great faith and diligence sustains and confirms our hope, and has commanded these mysteries of our salvation to be put in writing and published to the whole world in all ages. To him be glory forevermore. Amen.

## Sermon 92: ¶ It is furthermore declared, that the hope of the everlasting and blessed felicity and glory to be certain and undoubted.

### Paragraph 1304

.

### Paragraph 1305

I am Alpha and Omega, the beginning and the end. I will give to the thirsty person from the well of the water of life freely. He who overcomes shall inherit all things. I will be his God, and he shall be my sun. But the fearful and unbelieving, and the abominable, and murderers, and adulterers, and sorcerers, and idolaters and liars shall have their part in the lake that burns with fire and brimstone, which is the second death.

### Paragraph 1306

To all the former comes now the sixth testimony of the certainty of the true felicity of the faithful,[note 92.1] taken from the very nature of God. For he proclaims of himself and says, “I am Alpha and Omega.” And immediately by explanation, “the beginning and the end.” He took this from Isaiah, with whom the Lord says repeatedly, “I am first and last.” And here let no man imagine that God is first in order, referring the beginning to consequences, as though he had a beginning; or that he is called the last or end, as though he should ever have an end; but the contrary rather is to be understood, namely that God has no beginning, no end, but is everlasting, from whom all things have their being, and by whose decree all things have an end; where he himself endures forever, and his years never fade, just as the prophet says in another place, and the Apostle also. And since he is eternal, without beginning and without end, who always lives, and quickens and preserves all things that live, how, I pray you, should he not quicken the faithful? So certain therefore is the life, salvation, and felicity of the faithful, as it is certain that God is life, and indeed everlasting life. For he is everlasting, and the life of the faithful. I have spoken about the phrase of speech, “I am Alpha and Omega,” in the first Chapter and third Sermon.

### Paragraph 1307

The seventh witness of our assured salvation comes from the truth of God and his promises, and it shares a certain connection with the previous ones. For that which God has promised, he can certainly fulfill without difficulty. He has promised a blessed life, and therefore he will assuredly grant it to the faithful. He cites God’s promise at this time, bringing in God speaking to John and to us with these words: “To him who thirsts I will give of the well of living water.” That is to say, I, who am life and eternal life itself, will give the faithful to drink the water of life, that is to say, I will quicken him, preserve him in life, deliver him from death and all evils, and reward him with all heavenly gifts. Who can doubt the truthfulness of the one who promises, especially since this place or this promise is found in more than one location? David sings plainly in Psalm 36: “Your mercy, O Lord, reaches to the heavens; and your faithfulness to the clouds; your righteousness is like the strong mountains, your judgments like the great deep. You, O Lord, save both man and beast. How excellent is your mercy, O God! And the children of men shall put their trust under the shadow of your wings. They shall be satisfied with the fullness of your house; and you shall give them drink out of the river of your pleasures. For with you is the well of life, and in your light shall we see light.” Many of these things are found in the Prophets, and are explained by our Savior himself in chapters 4 and 7 of the Gospel of John, where he shows that he gives water and wholesome drink to the faithful, which at length will spring up into eternal life. It is therefore certain that the faithful are quickened by Christ; and therefore, the blessed life of the faithful is, and will be, most assured and certain, as promised by so many explicit promises of God. We discussed this water of life somewhat in chapter 7 of this book toward the end, and we will have certain clear matters in the beginning of chapter 22.

### Paragraph 1308

But in the meantime and by the way, he shows and declares unto us, after the apostolic manner, who willingly and often declare unto us the manner of our salvation, how eternal life is communicated to us, to wit, freely, graciously, which notwithstanding, for the doubtfulness of speech or understanding of words, we do not properly express the force of the Greek word. They are justified, says the Apostle in Romans 3:24, freely through his grace: that is to say, by the mere mercy of God, by no merit of man. For the same Apostle in the same Epistle to Romans, chapter 6, “The reward of sin is death,” he says, “and where on the contrary side he should have set, and the merit of righteousness, eternal life, for this member he places rather, and the gift of God is life everlasting.” And he adds immediately, through Christ Jesus our Lord. Therefore John rightly says that eternal life happens to the faithful freely: that is, by the very grace of God, through the merit of Christ, and by no desert of man. For if we could by our works and righteousness deserve eternal life, then Christ had died in vain, to no purpose. There was no cause why he should die, saying we might of ourselves have been saved. There is no effect, nor merit of Christ’s passion, such an effect as it is in very deed, that by the blood of Christ alone we are purified. For if there were or had been another means of salvation, Christ needed not to have been incarnated and have suffered. And that this term ought after this way and manner to be explained, many other places of Scripture prove. In Matthew 10:8, the Lord says, “Freely you have received, freely give.” The Lord will not have his Apostles to receive any recompense for the gift of healing. But speaking of the ministry, he says, “The workman is worthy his hire.” In John 15, the Lord says, “They have hated me without cause, doubtless without my desert, or undeserved of my part.” In 2 Corinthians 11, the Apostle says that he preached the gospel freely, without taking reward or recompense therefore. And in 2 Thessalonians 3, he says, “Neither have I taken bread of any man for naught, in short, where John says that life is given to the faithful, he attributes all things of our salvation to the grace of God, and the merit of Christ’s passion, and plucks it from man’s merits. And the same affirms Isaiah also in chapter 55, rebuking foolish men, spending their money about things of naught. Here ought therefore to cease the fairs of indulgences and pardons, and holy things in the church. Let the Pelagians keep silence.

### Paragraph 1309

However, lest anyone, by the free preaching of the grace and merit of Christ, against the deservingness of humans, should gather that the blessed life falls to idle people, sleepers, and those ceasing from all good works: and that God alone works, and we work nothing: but only suffer the operation of God in us, and for that cause nothing is required of us, he anticipates this. First, the Lord says that he will give to those who thirst to drink of the water of life. Therefore, faith and a fervent desire for godly things are required of us; not that faith is prior [?], but that it is given by God. For by thirst, to signify the faithful desire of a godly man, the Lord himself is the author in Matthew 5, pronouncing blessed those who hunger and thirst for righteousness. And also in John 6, the Lord himself understands by drinking to believe. Therefore, faith is required of us, that is, that we should thirst for the water of life. This very thing also the Lord grants by his Spirit and word, as we have declared elsewhere. And he says how he who is freely justified must also fight; not only fight, but must overcome. Therefore, the duties of charity are required, of which there is frequent mention in chapters 2 and 3 of this book, wherein is made most frequent mention of this fight and victory. And God will acknowledge those who labor thus valiantly for his children; to them he will show himself a father, and take them as the heirs of all their fathers’ possessions. They are bastardly children, who, being idle, boast of faith, praise God with their mouth and words, and deceive him with their deeds. You see therefore that both must be preached in the church: that we are justified and beautified freely; and so, being justified, must work good works, to which, notwithstanding, as to their merits, they ascribe not salvation, but to the mere grace of God through Christ.

### Paragraph 1310

Consequently, and on the contrary, he lists those excluded from the fellowship of blessed life and the blessed, compiling a register of sins and wicked men, as he has done at the end of chapter 9, verses 21 and 22. And as the Apostle has, in a manner, recited to the Corinthians.

### Paragraph 1311

And we suppose that in John’s time these sins were most common, nor sufficiently known as was fitting. Many also at this day judge them more lightly than true godliness permits. And we doubt not but that in this record, which is comprised in eight kinds or members, are contained all other like sins and wickednesses. But we understand that hell fire is assuredly due to them for their sins committed, who have no faith at all, neither can by any means be persuaded to repent and turn unto God. For in the first Epistle to the Corinthians, chapter 6, he says, “You were such, but you have been cleansed by the blood of Christ and by the Spirit of our God.” Therefore, if we have been such at any time, let us repent; or, if we have fallen into these sins again, let us rise up and turn to the Lord, who calls to him sinners and promises pardon and grace. But woe to the incurable, walking always and without repentance in the way of iniquity.

### Paragraph 1312

And we shall consider separately eight aspects of this register.[note 92.7] First are listed the fearful. But the Lord himself was afraid, and even quaked for fear of death: the saints of God have feared also, and often fled for fear: yet are they not for this cause condemned in the Scriptures. Therefore another fear is meant, namely that immoderate fear, by which, compelled, we do for fear of men, that thing which God has prohibited: and we ourselves convict in our own consciences, understand that we sin in so doing: or when, through carnal fear, we leave undone that thing which God has commanded us: briefly, when we more fear men, as princes or leige fellows, or enemies, or any other men whatever they be, than our Lord God himself. And therefore the Lord himself in the gospel said: fear not them which kill the body, and cannot kill the soul, &c. Matthew 10. The same in another place says: he that denies me before that aduouterous generation, I will deny him also before my father in heaven. Doubtless it is a foul shame to fear more a most wicked man than most holy God. But men offend in this regard at these days most grievously. For some attribute so much to wicked and cruel persecutors that even for them they will command to pervert the preaching of the Gospel, or to keep silence altogether. There is one who will set more by the King, Prince, Earl, Baron, Citizen, or plowman, Bishop or Abbot, or some flattering friar, or vile massmongering priest, and will fain and dissemble for his favor, rather than he will freely confess the truth, and fear and glorify God, who is to be feared only. Unto them says Isaiah: say not conspiracy, and be not afraid of terror of the enemies, neither be you discouraged. But rather sanctify the Lord of hosts: let him be your terror, let him be your fear. He shall be the sanctuary, and stumbling stone, and the rest in the eighth chapter of Isaiah. For unless we put away this vain and wicked fear, and go about to furnish up the Lord’s work valiantly, constantly, and without fear, we shall surely be cast down to hell. Let fearful men think hereof, and call upon the Lord, and take unto them the spirit of strength, and of wise and godly boldness: and do the work of the Lord not negligently, but diligently, valiantly, and constantly. He is greater that is in us, says Jerome in his canonical writings, than is he that is in the world.

### Paragraph 1313

Unbelievers are not weak in the faith, modest, and fearing God; but such as do not believe God’s word, its promises, commands, and threats, neither follow God nor his Christ. Rather, they follow strange gods, prefer to believe fables, and have withdrawn their hearts from God. And of these there is a great multitude at this day, who nevertheless have all in their mouth that they believe God and his word, but they do not believe the preachers, thinking truly that their incredulity is thus sufficiently excused. But where the preachers show nothing else but the word of God, they cannot but restrain God’s word while they despise the sermons of the preachers.

### Paragraph 1314

In the third place follows that the torments of Hell are due to the abominable and detestable. For “abomination” signifies abomination and stench. He notes therefore abominable and detestable men, to whom all religion is a mockery, who deride God and his word, and blaspheme all holy things, the children of Belial, uncurable, and scornful. These, although they know the truth, yet they know it to their own condemnation, seeing they contemn it, and with dogs and hogs return to their vomit and wallowing in the mire. Whom also the Apostles have noted: Peter in the second epistle, chapter 2, verses 1-3; Paul in 3 Timothy; and 12 Hebrews; Jude Thaddeus throughout the chiefest part of his epistle; John himself about the end of chapter 22, reciting in a manner the same register, calls them dogs. And would God we lacked examples at this day of abominable men and such kind of dogs. But there is no cause why we should marvel hereat, considering that we live in the time of all others, the most corrupt of Noah and Lot, as Matthew 24.

### Paragraph 1315

Of homicides there are sundry kinds. [note 92.10] For we kill with the heart, mouth, and works. You may see the expositors of the ten commandments on this, chiefly Dr. Musculus. But I think the world has never known a more notable, more cruel, and more shameless murderer, indeed a parricide most certainly, according to the word of Christ in John 8:44, the firstborn son of the Devil, than the Bishop of Rome. For he, in a manner, at all times, for these five hundred years and more, has blown the trumpet to all the grievous wars of Europe or Christendom; and again has granted to murderers, especially those warring for the See of Rome, most large and ample pardons, and promised heaven to them that die in that warfare. All the which, were it not for the great mercy of God, would have destroyed both body and soul.

### Paragraph 1316

Then Saint John recounts fornicators. [note 92.11] And he names the lowest kind, so that we should understand the higher and more vile ones, such as rape, adultery, incest, and Sodomitical acts; nor should we exclude gluttony, drunkenness, and all kinds of riotousness and the fostering of voluptuousness. Wherein we undoubtedly see that Saint Paul, under the term fornicators, comprehends all filthiness and riot. But in our days, whoring is so common that every most shameful fornicator is admitted to the altar; a married priest who keeps holy matrimony is expelled from it. For this we may thank Syricius and other Popes, whom the Apostle has severely noted in the first letter to Timothy, 4.

### Paragraph 1317

Note 92.12: Sorcerers are discussed in chapter 9 of this book. John [? describes] them, and by this he means magic, enchanters, diviners, witches, and those who use devilish arts to influence affections. The Latin authors also understand them to be those who administer poison.

### Paragraph 1318

Idolaters are those who worship idols. And it is astonishing that the papists of today deny that they are idolaters. For what else is an idol but a shape or image made of some visible material, representing the form of God or a saint, but without spirit? An idol, therefore, is an image of wood, stone, or metal, representing the shape of God the Father, of God the Son, or of Saint Peter, and so on. David describes an idol, and says: “The idols of the heathen are silver and gold, the work of human hands. They have mouths but do not speak; they have eyes but do not see.” Psalm 113. And I would fain know what difference the idols of the papists have from these? Concerning worshipping them, they cannot deny that they worship those idols of wood and clay. For they attribute to them holy names, and even the sacred name of God, which is communicated to no other, saying: “This [pointing to stone or wood, that is, an idol of wood] is God the Father, this is God the Son, this is Saint Peter.” I tremble in my mind while reporting these things, especially since the Lord himself has said, “Whom will you liken me to?” Isaiah 40. And Saint Paul calls this plainly counterfeiting foolishness, and expressly denies that the godhead is like a stone artificially polished. Acts 17:29. Again, these images, which they call their gods and saints, made with human hands, they bring into the churches, namely a place of worship, and set them upon the altars. To these they go on pilgrimage, fall down before them and worship, kiss them, offer oblations to them, and hang jewels on them. Moreover, they attribute to them also a part of the heavenly doctrine and instruction, saying that the unlearned are taught and admonished by these. And what is worship if this is not? Let them see, therefore, whether they can excuse themselves in this before God and men, and provide rather to save their souls. However, they sweep away all these things as if with one word, and say: “We do not worship the signs, but the things signified.” Then, if the signs were taken away, would they not return to the idols on pilgrimage? Do they not think it done in a manner to God himself that which is done to the idols? Do they not punish an image breaker as a traitor against the divine majesty? For he shall not seem to have cut asunder wood, but to have defiled God himself. Therefore, they acknowledge something more in this wood than wood alone. For you think that some divine thing is hidden therein, and therefore this wood is accounted by you not common wood. Which thing you declare also by sundry tokens otherwise. Moreover, the gentiles excused themselves in the same manner, saying that they worshipped the things, and not the signs. But this seemed not a sufficient excuse unto godly men, as it is to be read in Lactantius and Athanasius in their books against the gentiles. But God has at one word refuted you and said, “Who has required these things at your hands? If any will exhibit to me worship, let him worship according to the prescription of my most holy law. They worship me in vain, teaching the doctrines of men.” These things I have declared somewhat more at large, to the intent that those who will yet hear any reason, and in whom the word and law of God have any place, might know and avoid that gross and mortal sin of idolatry.

### Paragraph 1319

And liars include those who are quick of speech, slanderers, gossips, whisperers, deceivers, covetous people, thieves, usurers, bribe-takers, and all kinds of hypocrites and unreliable people. For just as God is truth, so he loves truth, simplicity, steadfastness, and integrity. This vice of lying prevails widely today. For there is the least, or rather no faith at all in the world. The Lord be merciful to us.

### Paragraph 1320

Regarding the lake or pond, burning with fire and brimstone, and of the second death, I have spoken before in the nineteenth and twentieth chapters. And elsewhere. He signifies that all these and the like shall be cast down by the Lord into the everlasting fire of Hell. For he puts here a portion as inheritance, as also in Psalm 11: he shall rain fire and brimstone upon the ungodly, and this is part of their cup. And in Matthew 24: he shall put his portion with hypocrites. And we say also, he has not obtained his right; or he is punished as he is worthy. Just as the saints obtain the Kingdom of Heaven by inheritance, so are everlasting torments in place of inheritance to the ungodly. To the Lord, the righteous Judge, be praise and glory. Amen.

## Sermon 93: ¶ Here is set forth a goodly picture, a description or figure of the blessed seat, and of the heave lie life and glory everlasting.

### Paragraph 1321

.

### Paragraph 1322

And one of the seven angels, who had the seven vials full of the seven last plagues, came to me and spoke, saying, “Come here, and I will show you the bride, the Lamb’s wife.” And he carried me away in spirit to a great and high mountain, and he showed me the great city, holy Jerusalem, descending out of heaven from God, having the glory of God. And her shining was like to a stone most precious, even a jasper clear as crystal; and had great and high walls, and twelve gates, and at the gates twelve angels; and names written, which are the names of the twelve tribes of the children of Israel. Three gates on the east side, three gates on the north side, three gates toward the south, and three gates on the west side. And the wall of the city had twelve foundations, and in them the names of the Lamb’s twelve apostles.

### Paragraph 1323

John returns to the description of the heavenly city, which at the beginning of this chapter he had begun. He has included certain matters that are fitting, concerning the certain hope of the faithful; after which he seems afterward to unlock and open Heaven, that the godly, with the eyes of faith, might as it were look therein, and clearly see what is the hope and glory of saints to come. For under the image of a most beautiful city, he sets forth a picture or description most evident of the blessed dwelling place, or palace and city of God, or of the everlasting realm and triumphant church. We shall not here invent and fabricate things earthly and physical, but spiritual and celestial. For the Spirit of God intends that, by means of temporal things, our minds ascend to the eternal, and that by temporal things, we grasp what is more excellent. Therefore all things are figured, with amplifications, hyperboles, and full of other figures. We shall therefore imagine far greater things: as we are accustomed to do, when we read or hear such things as our Lord taught under the parables of weddings and feasts.

### Paragraph 1324

And first, it is declared to us [note 93.2] who is the revealer of this godly and wonderful vision: that is to say, who is the opener of the mysteries, truly an Angel of God, and the very same who previously in chapter 17 to the same John said, “Come, I will show you the condemnation of the great whore,” and so forth. For it is the same God who punishes the ungodly, and gives rewards to the godly, and announces to men his righteous judgments by his ministers. Moreover, we see the sins to be most certain, and partly also accomplished, which he showed before concerning the judgment of Rome. Would not anyone conclude that what he now utters and shows concerning the everlasting glory of the faithful is also most certain? Gathering some of the things which he will show him, he sets them before and exhorts him to follow, saying: “Come, I will show you the bride, the wife of the Lamb.” Of her there has been spoken before on many occasions. He signifies the congregation of Saints, coupled by faith to our Savior Christ. And not only does he show to John (and in the same, to us all) the spouse, but also reveals her glory from God. The meaning therefore is this: “Come, I will show you what shall be the glory of the church of Christ in the life to come, what shall be the state of the everlasting life.” Certainly, he speaks also very many things of the church, but chiefly of her glory in the world to come.

### Paragraph 1325

[note 93.3] Then he briefly touches on the manner of revelation. For he adds, “And he took me up in spirit onto a great and high mountain.” Therefore, just as in the former visions he was carried away in spirit, his body remaining in Patmos, and as we have read and pointed out before, that such visions and raptures happened to Ezekiel, so now he says that he is carried away in spirit, and in mind to have seen the things the angel showed. Wherefore, if we will read or hear these things to any profit, we must lift up our minds, and be carried up in our spirit, and think that all these things must be understood spiritually. Arethas: rightly, he says, the heavenly life and conversation of the saints was shown to him on the mountain. For with them there is nothing earthly, low, or abject, but all things lofty and high. This he. Certainly, when the Lord Christ wished to exhibit to his disciples a certain taste and foretaste of the glory to come, he carried them up onto a mountain and was transfigured before them, which Matthew affirms in chapter 17, to have happened to Peter, James, and John.

### Paragraph 1326

And now he turns to the vision itself, and generally and briefly describes or foreshadows the blessed seat and glory of the life to come. Afterward, he amplifies this more largely, particularly, and as it were by parts, and so enlarged and beautified he sets it forth as if it were to be seen by the godly. [note 93.4]

He calls the heavenly country, and the habitation of saints, the great City. For it is the city of the great king, and in it shall dwell an innumerable number of the blessed, and of angels thousands without end, and shall have the fruition of great glory. Neither is there any fear that the place should not suffice so great a host of men and spirits, or that it shall be too crowded. Great is the city of God, which is truly able to receive all good men abundantly. In the Gospel of John the Lord says: “In my Father’s house are many mansions,” and so forth, John 14:2. The same place is called holy Jerusalem. For just as no defilement shall be seen there, so shall no unclean person appear. Of the heavenly Jerusalem is spoken before. Thomas Aquinas says: She is said to have descended from heaven, for that whatever goodness the holy church has, she acknowledges that she has received it from the grace of God. But of this matter I have spoken in the last Sermon. [note 93.5]

And the city of God, that is to say heaven, has the seats of God and the blessed, the glory of God, that is to say the divine majesty and brightness, and whatever great thing so ever the mind of man can think or imagine, or in all things the unspeakable excellence of God, such as neither the eye has seen, nor the ear has heard, nor has ascended into the heart of man, 1 Corinthians 2:9. These things he has touched upon briefly and generally hitherto.

### Paragraph 1327

And consequently he declares by particulars and at large that celestial glory and blessed seats are present [note 93.6]. For whatever things are simple, whatever are in cities commendable, the same are plainly found in this our city most excellent, such as its size, strength, majesty, security, excellence, beauty, pleasantness, and abundance of things. These things I say, and all other like, wonderfully excel in the city of our God and in our fathers’ house. And whereas these things are set forth and amplified in this way most liberally, yet it seems nothing at all is said, if a man considers the unspeakable majesty of the celestial glory [note 93.7]. But all these things are alleged by the Lord through John to this end, truly, that we should be taken with the desire of so worthy a life, and should think in our tribulations and troubles that the afflictions of this present world are nothing, being compared with so excellent and sovereign glory; finally, that all are mad, who begin to doubt of the eternal hope of the faithful. Very many things of this sort are also read in Ezekiel in chapter 40 and following. We will touch every part of this treatise, using nevertheless a succinct brevity, lest we should be tedious to any man. And truly he touches the principal and most commendable things of cities, and in them shows that the city of God excels.

### Paragraph 1328

In cities and houses, the greatest commendation is when all things are bright and clear, for darkness is horrible and unpleasant. Therefore, an excellent light is declared to be in the city or house of the Lord. A parable is added, showing the excellence of this light. It is like a most precious stone, such as a jasper, as it is commonly called, or a chrysolite, or some other stone most bright. And John himself adds more, as it were a jasper stone like a crystal. This is a new way of speaking, but it has a marvelous grace if we understand it rightly. For a jasper is green, and a crystal is bright. He seems therefore to say that this celestial brightness is continually green and never withers; that is to say, the heavenly light is everlasting, and in itself, after a sort, wearing green, and in growing green, it shines brightly and rejoices all heavenly dwellers. For afterward follows: “For the glory of God has lightened her, and the Lamb is her light.” This brightness and most joyful light the Lord promises in sundry places in the gospel of John, and the whole blessed life, of this not the least part, is commonly called blessed light, or everlasting light, or light of heaven. It seems to have been prefigured in the golden candlestick of the tabernacle, and so forth. For if it were not difficult for our Lord God to give to precious stones wonderful colors and brightness; if he illuminates this world full of wicked men with most goodly lights, the Sun, Moon, and Stars, what light may we think to have in heaven, where no man shall dwell but the best, and of God most dearly beloved? Of this light much mention is made with Isaiah and in the psalmist.

### Paragraph 1329

Walls in cities are most notable and excellent when they are high, thick, and strong, able to withstand all force from enemies, and defend the citizens from all harm, keeping them in peace and security. Therefore, strong walls are both great and strong, and also high and impregnable. This signifies that the protection of the saints in heaven will be most safe and sure, so that the saints will be in perfect security, and exempt from all fear. No one will trouble or take away their joys, as the Lord affirmed in John 16. For there will be perpetual security and gladness in heaven, most perfect and everlasting.

### Paragraph 1330

Moreover, gates are placed in the walls, by which people may enter the city. Therefore, in the wall of the great country there shall be twelve gates, that is to say, a most ample entrance into eternal life shall be open on every side. And we believe that there is no other way to heaven, no other port or gate, or any other door or postern to remain, than the only and sole Christ Jesus our Lord, as he himself has taught in John 10 and 14. But because he has appointed angels, prophets, and apostles also, as porters of heaven, to whom he has committed the keys of the kingdom of heaven, and these do bring the chosen and let them into the heavenly country, many gates are truly said to have been and to be. And for a further declaration, it is added that in every gate was an angel, namely twelve. And we have heard in the beginning of this book that angels are God’s ministers and pastors of churches, sent by God for the salvation of people, that is, that they might bring them by the word of truth and holy ministry, through faith, into life everlasting. Moreover, we read how the soul of poor Lazarus, dying, was carried by angels into the bosom of Abraham. Why, then, should we marvel that angels stand at the gates? For by the true and only gate, Christ, they bring the faithful into the heavenly country.

### Paragraph 1331

And again, for a further declaration is appended,[note 93.12] and in the gates were names written, which are the names of the twelve tribes of the children of Israel. For the Lord would signify that he used the industry of the patriarchs and prophets of all tribes in opening heaven to me; and again, that all the elect of all tribes belong to the fellowship of blessedness. We shall see, therefore, in heaven, the patriarchs and prophets, and all the saints who, before the coming of Christ, are written in the registers of the heavenly realm, just as the apostles also saw Moses and Elijah talking with Christ on the mount. Wherefore, not without great cause, the Apostle to the Hebrews wrote: “You have come to Mount Zion, and to the city of the living God, the heavenly Jerusalem, and to the multitude of angels, and to the assembly of the firstborn who are written in heaven.” And the rest which is read in the twelfth chapter.

### Paragraph 1332

And he also addresses the matter of the gates. He assigns three to each part of the sky, and does not do this without reason. For our Savior himself says in the Gospel that they will come from the east and from the west, and will rest with Abraham, Isaac, and Jacob, in the kingdom of heaven. Arethas investigates this mystery more closely, and believes that no one will enter through these gates unless they acknowledge the eternal Trinity of God and understand the mystery of the cross of Christ. He says how the twelve tribes are divided by the Trinity, according to the fourfold figure of the world, &c. Aquinas also says: whoever are saved, they are justified by the faith of the holy Trinity, proclaimed throughout the four quadrants of the world by the preaching of the Apostles.

### Paragraph 1333

Now he also shows that the foundations of this city are most sure and immovable. For the wall of the City, he says, has twelve foundations. Concerning the foundation of the church and our salvation, David spoke expressly in the Psalms. Isaiah also in chapter 28. Our Lord and Savior spoke in sundry places in the Gospel. Peter moreover in the Acts and his first epistle: likewise the Apostle Paul, who said, “No other foundation can be laid than that which is laid, which is Christ Jesus,” 1 Corinthians 3. How then are there laid here twelve foundations? Doubtless Christ remains one and a sure foundation. However, inasmuch as in placing and revealing him, the Lord used the ministry of the twelve Apostles, for this cause the city is said to have twelve foundations. Not that the Apostles are in deed the foundations of the church and our salvation, but in this that Christ, that true foundation, was by the twelve Apostles made known to the faithful, and as it were laid under, whereupon the believers have built themselves by the Apostles’ faith. Whereupon he says purposefully, “and in those twelve are the names of the Lamb’s twelve Apostles.” For the gospel also, which is both in very deed and unchangeably Jesus Christ’s alone, is called of John, Matthew, Mark, and Luke, of Peter and Paul, because it has been preached by them. And we understand hereby not only the church which was before the coming of Christ, of Patriarchs and Prophets, being received into heaven to rejoice in God, but also the Apostolic church, I mean that all men in the whole world who have believed the Apostolic doctrine shall live with all the Saints in that heavenly country; all of which we shall both see, and with them also shall glorify God forevermore.

### Paragraph 1334

Primasius, Bishop of Utica, does not greatly disagree with this our explanation, explaining how the Apostles are called foundations. For thus he has written: since we know that the church has one only foundation, that is to say Christ, we ought not to be disturbed that here he says she has twelve. For in Christ the Apostles have deserved to be the foundations of the church; of whom the Apostle says, “another foundation can not be laid, besides that which is laid, Christ Jesus.” In him also the Apostles are said to be light, as he says to them, “You are the light of the world,” where Christ alone is the true light, which illuminates every man coming into this world. Christ therefore is the illuminating light, and they are the lights illuminated. And after a few words, the same author: “Here it behooves to acknowledge the twelve Apostles to be foundations, yet called in the only foundation, Christ Jesus.” Hereunto pertains also that he has not concealed the name of the Lamb. The Apostles therefore are foundations, but in one foundation, Jesus Christ. And Christ alone, without the Apostles, is rightly called the foundation; but the Apostles without Christ, could by no means be called the foundations of the church. These are the words of Primasius. Which Arethas, Bishop of Caesarea, declares more briefly and plainly, and says: they are indeed called foundations, for that they have laid the foundations of the Christian faith; and gates, for that by them—that is to say, by their preaching—there may be found a way to bring people to the Christian faith. Thus much he says. Doubtless the Apostle Paul in 2 Ephesians calls Christ the foundation of Apostles and Prophets, which verily in preaching they have laid, and to which they have leaned, and by which also they are saved. To him be glory.

## Sermon 94: ¶ Yet again is described the seat of the everlasting country in heaven.

### Paragraph 1335

.

### Paragraph 1336

And he who spoke with me had a golden reed to measure the city, and its gates, and its walls. And the city was built squarely, and its length was equal to its breadth. And he measured the city with the reed twelve thousand six hundred furlongs, and the length, breadth, and height of it were equal. And he measured the wall of it one hundred forty-four cubits, according to the measure of a man which the angel had. And the construction of the wall was of jasper. And the city was of pure gold, like clear glass, and the foundations of the walls and of the city were adorned with all manner of precious stones. The first foundation was jasper, the second sapphire, the third chalcedony, the fourth emerald, the fifth sardonyx, the sixth sardius, the seventh chrysolite, the eighth beryl, the ninth topaz, the tenth chrysoprase, the eleventh jacinth, the twelfth amethyst. And the twelve gates were twelve pearls, and each gate was one pearl, and the street of the city was pure gold, as a thoroughfare of shining glass.

### Paragraph 1337

He proceeds in describing the blessed dwellings and that life of the world to come, under the image of a most beautiful and most excellent city. We must understand all things not according to the literal meaning, but according to the spirit. All things are said for our comfort, and so that we should boldly despise this world and its pleasures, and the fury of persecutors; and should always desire those great and everlasting good things promised to us, as we have heard in the description. Indeed, we have seen four particular things of this heavenly city as if in a lively picture: what light it has, what walls, what gates, and foundations. Now, in the fifteenth place, follows what the width, capacity, or largeness of this city is. For cities are commended because of their size. It is necessary that the greatest number of citizens should have the largest city.

### Paragraph 1338

Therefore comes forth a messenger of this city, an angel sent to John from heaven, holding in his hand a reed, that is a long pole or measuring rod, not of wood or lead, but of gold. By the measuring, he would have us understand the quantity of the blessed seat. In the meter and in the measure, we shall not need to seek any great mysteries. For the eternal wisdom and providence of God has prepared seats for his chosen; and that in a golden order, that is to wit, most purified, which is signified by the golden reed or measure. For the judge in Matthew provokes the sheep to take the inheritance, prepared from the beginning of the world. He alone knows also, who are his.

### Paragraph 1339

The situation of the city is declared, planted in a square, whereby is signified the strength and stability of the blessed in heaven. For the place is no ball, bowl, or globe, rolling and easy to turn. Neither need we to doubt of the certainty thereof. For hope shames no man, and he that believes in Christ shall never be confounded.

### Paragraph 1340

Moreover, the length, breadth, and height of the city are equal. Every side, in its square, has twelve thousand furlongs, which make in all forty-eight thousand in the whole circuit. Concerning the furlong, what and how much it contains, I see that learned men vary. Pliny in book 2, chapter 23, attributes to a furlong one hundred and twenty-five paces, that is to say, six hundred and twenty-five feet. If you now calculate these things and divide them into miles, you will find that the city is most ample and large. Some reckon it one hundred and fifty German miles. Hereby I suppose to be signified that the place and space is great enough, whatever innumerable multitudes of angels, of blessed spirits, and men shall flit into the blessed seat and dwell therein. As also the Lord in the Gospel said: “In my Father’s house are many mansions.” In chapter 30 of Isaiah toward the end, it is shown that there shall be space and place enough in hell also for the wicked. And the equality on every side declares that men of people or nations shall have no prerogative. For whether you be of the East or of the West, whether you be Greek or Barbarian, so that you be faithful, you shall be received by the Lord. Moreover, in the gospel equality is declared, while the penny is paid not only to him who wrought in the vineyard all day long or half the day, but unto him also who came into the vineyard in the evening.

### Paragraph 1341

Note 94.6: The height of the wall is undoubtedly immeasurable. From this we gather that the blessedness is most certain, and that no one can enter it except through the gates. For no one can climb over such a height, nor scale those walls, whether he be an enemy who would trouble them, or a hypocrite who attempts to gain heaven by stealth.

### Paragraph 1342

Where he says, “and he measured the wall thereof, an hundred and forty-four cubits,” it cannot certainly agree with furlongs; therefore, we must understand it of the thickness of the wall. By this again is figured the strength and security of the blessed. It is added how the Angel measured with the measure of a man, which the Angel had: that is to say, that the Angel measured the customary cubits and furlongs used by men. Therefore, this Angel had the same measure in this measuring as is commonly used by men. For so would he signify that the place of eternal felicity should be determinate and certain. For there shall be after the resurrection true and determinate bodies. If there be any other mystery herein, perhaps it is the same which the Lord spoke of in Luke, namely of the blessedness of the faithful in another world: they are like angels, and are the children of God, since they are the children of the resurrection. If any man will account these numbers more exactly and show higher mysteries, I will gladly give place. I suppose here rather celestial things to be figured than either arithmetical numbers or geometrical proportions to be taught. Nevertheless, I can willingly grant that those arts help to the understanding of the Scriptures.

### Paragraph 1343

In the sixth place, he discusses the matter of this heavenly City. For cities are praised for their substance and material. [note 94.10] The saying of Caesar Augustus is well known, who is said to have spoken of Rome: “I found it of brick, I leave it of marble.” And cities built of stone are justly preferred to those that are of timber, and those built of squared free stone to those made of rough stone. But what is the building or material of the celestial City? He declares that by five parts or members. [note 94.11] First, the walls are of jasper. Let no man here imagine carnal things. The jasper is green. The celestial City always flourishes; God’s protection never fails.

### Paragraph 1344

2. The city itself, that is to say, the buildings in the city, the palaces and houses, are pure gold. For all things are purified in the eternal realm. There is no uncleanness, no evil affections, there shall be no trouble or pain. As the Lord said also in the 19th Chapter of Matthew, disputing against the Sadducees. Therefore, just as gold is most tried and pure, so shall the heavenly dwelling be most clean. Therefore, the bodies that shall dwell in heaven must also be clarified or glorified. He adds that this most pure gold is not glass, but in brightness represents most pure and shining glass. For in heaven all things are clear. There we shall be seen face to face. There we shall most perfectly know all things.

### Paragraph 1345

3. And first he says generally that the foundations of the city are beautified with all manner of precious stones; after which he recites particularly the stones that are most excellent. Doubtless nothing is more precious, nothing more excellent, than Christ, the foundation of our salvation, than the apostolic doctrine, whereby we are induced to the knowledge of Christ and of our salvation. And he sets in order twelve stones, to the intent we should understand that there is not one precious stone alone placed for the foundation, but a row of one sort in such a length as the side is square, and so consequently likewise in all parts of the square. For the first order therefore is placed a jasper stone, that is to say, in the first place of the foundation, jasper stones are set in their rank; again in the next row upon the jaspers are laid sapphires, throughout the whole space, in such length as the foundation was, and so consequently the other stones were couched and laid in order. By all which is signified that the foundation of our salvation is both most excellent and sure. Which we ought of right to set more by than by the price of all the jewels in the earth. And there are found men godly and beneficial, who, bestowing or selling these earthly jewels, according to the Apostle’s doctrine in 1 Timothy 6, prepare for themselves a good foundation in another world. There are foolish people who are overly in love with jewels, and many times instead of precious stones that cost very much when polished, they obtain glass. Truly worthy doubtless to be deceived. Verily, precious stones have their use and virtues, neither were they made by God in vain. But we must always remember that saying of the wise man: all things are not suitable for all men.

### Paragraph 1346

By the listing of precious stones, he appears to have alluded to the precious stones set in the attire of the high priest, as described in Exodus 28. I do not doubt that John took these things partly from Isaiah 54, which Jerome expounding, directs those who desire to know more about stones to Epiphanius and to Pliny’s 37th book on Natural History. Arethas in his commentaries applies the twelve precious stones to the twelve Apostles of Christ. Furthermore, the writings of Bede on this passage remain; from whom Thomas Aquinas took what he has in his commentaries on the Apocalypse. I see no way that I can profitably linger longer in this treatise. Therefore, I refer the curious reader to these authors; it is enough for me to have shown that these costly jewels signify the excellence of the foundation of our health and salvation.

### Paragraph 1347

Moreover, in the fourth place, the matter of the gates is declared. They were all of one whole pearl, each of them, the price of which is exceedingly great. The gate of heaven is Christ, and the gatekeepers of heaven are the Apostles, as has been declared before. Therefore, the gates are most precious and most strong. In the thirteenth chapter of the Gospel according to Matthew, Christ himself and the salvation that comes from him are compared to a pearl, for which the merchant sells all that he has, thinking himself rich enough if he may obtain this pearl.

### Paragraph 1348

5. In the fifth place, the street is also described, what it is. [Note 94.16] In earthly cities, the streets are often muddy, even if the cities themselves are otherwise famous and noble. Where they are noteworthy, they are paved with stone or brick; but the street of our city is paved with pure and bright gold. For in heaven there is no unpleasantness, no obscure darkness. All these things are undoubtedly spoken beautifully; yet far greater things must be understood and imagined. We must strive with all our might that the thing which the tongue of man cannot utter, nor our mind here conceive due to its greatness and excellence, we may at length behold in heaven presently, and experience it in our glorified bodies through Jesus Christ our Lord.

## Sermon 95: ¶ Furthermore yet is described the everlasting country in Heaven.

### Paragraph 1349

.

### Paragraph 1350

And I saw no temple in it. For the Lord God Almighty and the Lamb are its temple, and the city has no need of the sun, neither of the moon, to light it. For the brightness of God illuminates it, and the Lamb is its light. And those who are saved shall walk in its light, and the kings of the earth shall bring their glory and honor into it. And its gates shall not be shut by day, for there shall be no night there. And they shall bring the glory and honor of the Gentiles into it. And there shall enter into it none unclean, neither whatever works abomination, or makes lies, but those who are written in the Lamb’s book of life.

### Paragraph 1351

The Apocalypse proceeds in the description of the divine or celestial City, to comfort and keep the faithful in all temptations and afflictions. Therefore, in the seventh place, he discusses the temple. [Note 95.1] For in famous cities, there is no small consideration and praise of churches. This is evident by all writers of histories, places, and times. What temple then is in heaven? None at all. For Saint John says, “And I saw in the city of God no temple.” This passage does not contradict those things which are in the 11th and 15th chapters concerning the temple in heaven. For the temple is there exhibited in a figure and vision, not that there is in deed any temple in heaven, but that thus God’s justice and certain salvation promised in the Scriptures might be signified, as we have declared in those places.

### Paragraph 1352

And what is the cause that no temple appears in heaven? The divine revelation answers: for the Lord God Almighty, and the Lamb are the temple in that our heavenly country. The purpose of temples is this: the Lord first instituting the tabernacle, and then the temple, wanted it to be known that he would be present in the midst of his people, as a father, Lord, and defender. And therefore it is said in the scriptures that people come unto the Lord, whether they came to the tabernacle or the temple of the Lord. The temple moreover was erected for preaching and prayer and the external service of God, for receiving the sacraments, or offering up sacrifices. But the Saints in the heavenly country have no need of all these things. Therefore they need no temple. Therefore no temple is seen in heaven. For the Lord God now shows himself to them to be enjoyed; the saints are now with him, wherefore they need no token of his presence. We are taught by doctrine what God is, and what is his will, and that we are saved by the Lamb: but now that we see God himself face to face, and that salvation has come by the Lamb of God, what need is there of a temple in heaven? By prayer we require life and everlasting joys: now since these have happened to the elect, what need is there of a house of prayer? The Saints now, without any temple, offer up eternal praises unto God. And seeing that sacrifices and sacraments have no further place in the everlasting country, I see not why there should be any temple in Heaven. We rest and keep in Heaven an everlasting Sabbath. This place moreover proves that Christ is very God, coequal with the Father, as to whom he is joined inseparably in all glory. Neither is the Holy Spirit separated from the Father and the Son, which elsewhere is said to dwell in us: for this cause we are called the temples both of God and of the Holy Spirit, by the Apostle, in 1 Corinthians 3 and 2 Corinthians 6.

### Paragraph 1353

The eighth section of this description is repeated concerning the heavenly light, and that not without great cause; indeed, the same light is again commended in chapter 22. For in buildings, there is nothing more excellent than light. Otherwise, without light, all things are blind. Furthermore, he does not say that the Sun and Moon should no longer exist, but that the City of God should not need those lights. He shows the reason: for the glory of God has illuminated it. And the glory of God is the divine, celestial, and unspeakable brightness of his inapproachable light, which he inhabits, and according to his good pleasure communicates to the elect. The Lord Christ (who here is called the Lamb, for the mystery of redemption) illuminates the blessed. For by him we are clarified, and enjoy that eternal, most beautiful, and celestial light. John has borrowed this passage from chapter 60 of *Confessions*, where we read: “The sun shall not be there, for the light of the day. And the brightness of the moon shall not shine there: but the Lord shall be to you a perpetual light, and your God shall be your brightness.” Furthermore, the seats of the blessed are thought to be fixed above the sphere of the Sun and Moon, and also the brightness of the saints to surpass far the light of the Sun and stars. Augustine has testified this also in chapters 24 and 30. To God almighty and eternal light, be praise and thanksgiving, who has prepared such great things for us, and gives us gifts such as no tongue can express.

### Paragraph 1354

He shows in places more than one [note 95.4] who are partakers of that light, or who are citizens of this celestial city, and what is the state of the citizens. All nations and people saved are citizens of the eternal country. Here are two things to be noted. First, that the Gentiles are made inheritors of glory, and that without any choice. For here excels not the Jew, nor the Greek, neither Roman, nor Barbarian. Again, yet not all without respect, and confusedly obtain everlasting light, but the saved only: that is to say, whom Christ has saved and redeemed from sin, the Devil, Antichrist, and from the curse and the world. And Christ saves the elect and faithful. They therefore shall in deed be partakers of the light: these are the citizens of the country everlasting. But what is their state and inheritance? They shall walk in the light of God the Father and the Lamb: that is to say, they shall have the fruition of the light and of God himself, to their joyful sweetness and fulfillment. For it is a figurative speech, to walk in the light, for that which is to enjoy light. Verily in Psalm 88 is read with a figure not much unlike: Lord, they shall walk in the light of thy countenance. And again: Thou shalt make known unto me the path of life, the fulfillment of joys is in thy sight, and gladness in thy right hand forevermore.

### Paragraph 1355

But especially the places in heaven, and in that divine Palace, are for Kings. Kings are governors and captains of the people, as they are called Kings and Princes, governors, Magistrates, rulers, both of the political and ecclesiastical government, Doctors, masters, teachers, Artificers, and Parents. For their duty is virtuously to govern their subjects, scholars, or children, to keep them under awe or discipline, to chastise, and direct them to the duties of life and all godliness. This if they do, they shall have a worthy place prepared in Heaven. For Daniel also says in the twelfth chapter: “But the teachers shall shine as the brightness of the firmament; and they that bring many to righteousness, as the stars everlastingly.” O therefore, O happy are you if you bring many to execute the office of righteousness. But woe to you Princes, and teachers, and masters, and Parents, if herein you be negligent. There is prepared for you in hell a place most horrible and miserable, as also Ezekiel has testified.

But if Kings have their place, and the same right honorable, in Heaven, wherefore do the Anabaptists teach, nay why do they lie, that a Christian cannot execute the office of a Magistrate? For here are Kings mentioned to be in heaven, not only as men, but as they were kings, that is, as they were good kings, and executed their office duly, and not forsaking their place, have lived a private life. For it follows, they shall bring their glory and honour unto it. And what is that glory, and what is the honour? it follows again: and they shall bring the glory and honour of nations into it: that is to say, they shall bring into heaven with them, the very nations, their people and subjects, whom they have helped in true godliness and salvation, in teaching, correcting, defending, alluring or drawing, &c. And these be their glory and honour, for Paul in the second chapter of the first to the Corinthians says, “for we are your glory, as you shall be ours also in the day of our Lord Jesus.” And again in the first to the Thessalonians, the same Apostle says: “for what is our hope, joy, or crown of rejoicing? Are not you, in the sight of our Lord Jesus Christ, at his coming? For you are our glory and joy.”

Aquinas rightly says: Paul speaks after the manner of conquerors, who bring their spoils into cities. Therefore he feigns that Princes, preachers, and parents bring with them into heaven such as they have won, which to them shall be an honor and glory. These things always let us think upon, and do our duty enjoined us of God, which we perceive in the everlasting country to have so great a reward. For it shall be the greatest glory that may be, to stand with so many won in the presence of the eternal God, Lamb, and all saints. Contrariwise the greatest shame to stand with so great a multitude of men lost, and that lost through our fault and negligence. Read what things are written in the first chapter of the book of Wisdom, &c.

### Paragraph 1356

In the tenth place follows the guarding of the heavenly gates.[note 95.6] Certainly, in great cities, great and diligent watching, guarding, and attention are taken to ensure the gates are shut and opened at the proper time. But in heaven, such carefulness is not needed. The reason is: the gates are not usually shut during the day, but at night. But in the eternal realm, there is no night, therefore the gates are never shut. There is undoubtedly no night, but continual day. There is no treachery, no ambushes or plots, no perils or dangers: all things, in general, are safe, peaceful, quiet, secure, and sure. The same things are also read in Isaiah, but in a different sense. Arethas says: here is a double understanding, for either he means that there will be peace and security, so great that it will not need to protect the city by shutting its gates. Or else that there also the godly gates of the apostolic doctrine are open to all, for their learning, and so on. Certainly, they will need no teachers nor guides, who see all mysteries clearly now and are brought into heaven itself.

### Paragraph 1357

And especially cleanliness in cities is highly commended, if nothing appears that offends the sight, hearing, and smelling—nothing loathsome to look upon and to be abhorred. And in private houses, the chief praise is if all things shine, and everything stands in order, and does not lie scattered and stink.

### Paragraph 1358

Now therefore, in the eleventh place, he shows that there will be nothing in heaven that is offensive, that is to say, which will not be pleasing and delightful, complete and absolute. This same passage must also be understood in reference to individuals. For it follows: “except those who are written in the Lamb’s book of life.” We understand therefore, that no whoremongers, idolaters, liars, deceivers, or anything unclean, and not purged by the blood of the Son of God through faith, will enter the kingdom of heaven. The Apostle affirms this in 1 Corinthians 5 and 6, and in Ephesians 5. David also asks: “Lord, who shall dwell in your tabernacle, or who shall rest in your holy hill?” And he answers immediately: “He who walks without spot, and does righteousness,” and what follows in Psalm 15. Finally, here will be fulfilled those things written in Deuteronomy 23, concerning those who are prohibited from entering the church. Therefore, this passage has a secret doctrine and a private admonishment, instructing us that if we desire to be heirs of the everlasting kingdom, we should all, while we live on Earth, apply ourselves to righteousness and innocence. For it will follow in Revelation 22: “Outside are dogs and sorcerers and whoremongers,” and so forth. May the Lord bring us by the way of righteousness to everlasting life.

## Sermon 96: ¶ He continues yet in describing the blessed seats.

### Paragraph 1359

.

### Paragraph 1360

And he showed me a pure river of water of life, clear as crystal, proceeding out of the seat of God and of the Lamb. In the midst of the street of it, and on either side of the river, was wood of life, bearing twelve kinds of fruit, and yielding fruit every month. The leaves of the wood served to heal the people. And there shall be no more curse, but the seat of God and the Lamb shall be in it, and his servants shall serve him. They shall see his face, and his name shall be on their foreheads. And there shall be no night there; they need no lamp, nor light of the sun, for the Lord God gives them light, and they shall reign forevermore.

### Paragraph 1361

In the twelfth place is described John’s vision of the pleasantness, abundance, and richness of food in the City of God. Rivers make cities pleasant and delightful. Without springs, fountains, and wholesome waters, cities decay and are scarcely worthy of the name. But if they lack sustenance, they are utterly lost. Therefore, our heavenly City excels in all these things and is most noble. It not only provides sustenance but gives it with great pleasure and delightful finesse. For the trees in this City not only bear fruit but also give an unspeakable and inestimable joy. The river moreover runs through the midst of the streets; on the banks of either side are trees most beautiful to behold, bearing the fruits of life. As I have often intimated in this description, I repeat now that these things are not to be understood literally, as the Millenarians take them. For the Lord speaks to us, even “lisps,” so that we, with our limited understanding, might grasp these things. If anyone should desire earthly things, I believe he could desire nothing greater than what is described here. We should therefore consider that if the Lord could give these earthly things, and would, why could he not give greater things to the souls of the godly and bodies glorified? Yes, the Lord desires that, being withdrawn from the contemplation of earthly things, we should look altogether for celestial and divine things, worthy of blessed souls and glorified bodies. Which, indeed, how great they are and what they shall be, no human tongue can express, however eloquent. For the Lord has prepared greater things for his servants than we can comprehend. Therefore, he brings forth here ample material so that, in a certain manner, we might conceive of heavenly things much more excellent than they are. Therefore, the sense and meaning of all those things spoken here of the river of life and the wood of life, by a most excellent amplification, is none other than that the blessed in the heavenly country shall be quickened by God and preserved in that happy life with unending delight. And there is no doubt that John has borrowed these things, as he does all the rest, since he is the interpreter of the Prophets, from the Scriptures. Therefore, he alludes to Paradise, whose description is set forth in the second chapter of Genesis, which agrees very well with this description of Heaven. For there also springs a river in Paradise, which immediately is divided into four heads and waters the garden of pleasure most pleasantly. In that same Paradise is the wood, that is, the tree of life, bringing forth lively fruit to the eaters, as explained by Augustine in book 13 of *The City of God*, chapter 20. But for the sin of our first parent, we were cast out of that Paradise; and Christ came to bring us again into Paradise, that is to say, into high felicity. Now, therefore, that true paradise, prepared for us by Christ, is shown in heaven and is here described. Into this Paradise entered the Lord after death, and brought with him into the same also the faithful thief, to whom he said, “Verily I say unto thee, this day shalt thou be with me in paradise.” Therefore, we ought not here to imagine the gardens of Hesperides in Earth, or in the air above the globe of the Moon, and reason of Paradise terrestrial. Our Paradise is celestial, which is prepared for us in heaven, as Paul has said in Philippians 3. Paradise is called a garden of pleasure, as at present it is called a golden City or of precious stones, verily by a trope on either side. Hereunto appertains also a passage from Zechariah in chapter 14. There is also another passage of Scripture in chapter 47 of Ezekiel, which John has translated or written out into this place in a manner word for word: “By the river,” he says, “on either side of it shall grow up all manner of trees that bear fruit, whose leaves shall not fall, nor the fruits fail, but every month shall they bring forth new fruit. For their waters flow out of the sanctuary, and their fruits shall be meat, and their leaves medicinable.” And Ezekiel, under a figure, sees that same blessed life and happy seats which John at present sees, by the showing of the Angel. And either of them both sees the happy seats after the same sort, and under the like figure. For there is one only blessedness, common to all the faithful of the whole world. The Patriarchs, Prophets, Apostles, and Martyrs achieve all one felicity. They see the River on either side, and the same flowing out of the Sanctuary, or seat of God. They see on either side the river, trees planted that bring forth the fruits of life. They bring forth fruit every month, fresh and new; and the leaves of either do heal. I suppose the old poets borrowed from the Scriptures such things as they wrote in verses concerning Ambrosia and Nectar, the meat and drink of the Gods. That short verse of Martial is known.

### Paragraph 1362

And Grammarians derive those terms of immortality. But our Saint John here reasoning more elegantly and better concerning these matters, says how the angel showed him a river, which he commends to him by such properties as water is wont to be commended: for its purity, brightness, and clearness. He adds a parable, which sheds light on what he has said, and says, clear as a crystal. After this, he adds that this river is the river of the water of life, that is, lively water, which preserves those who drink it in life. Finally, he shows also the origin or springhead of this river, deriving it from the seat of God, of which seat or throne I have spoken in the fourth and fifth chapters of this book. And by all these things is signified nothing else but that life proceeds from God alone, which he gives to those who serve him in that blessed country, pure, clear, bright, most tested and most perfect, and altogether divine. Concerning the lively springs and fountains of waters, we have touched somewhat in the end of the seventh chapter of this book. Note again that God and the Lamb are so joined together that no man unless he be mad will deny the Son to be of the same substance with the Father.

### Paragraph 1363

Now follows the description of this divine city. The sustenance in the eternal land is the tree of life. It is the Hebrew phrase, “the wood of life,” for the tree of life, or lively sustenance. For there is added bearing fruit. Whether you understand that John saw only one tree, as also there was only one tree in paradise: or more, as in Ezekiel, so that by the general word we may understand the particular kinds of trees, it shall be all one. The location of the tree he shows diligently, to be set in the midst of the street of the city, and on either side the river (whereby it is doubtless gathered that there were many trees) to wit, on the banks of the river, that they may draw lively nourishment out of the river, which flows from the throne. And hereby I suppose is signified that the heavenly food is common and free for all, not locked up, or kept for a few. It is found in the midst of the street of the city: thus the sustenance stands open, and is not hidden. And it draws a lively force out of the river, which springs out of the seat. For that heavenly life is from God, and flows to all his elect. Moreover, it is also declared most diligently, what manner of fruit this shall be: the tree of life says he, does bear fruit twelve times in the year, so that every month it bears fruit fresh and new. The first fruits are delightful: and those who commonly abhor old fruit, had rather have new. Therefore in that blessed land there shall be nothing tedious, unpleasant, loathsome, or in any wise to be rejected, but all things shall be most pleasant, most delicate, fresh, and delectable.

### Paragraph 1364

Now he does not overlook the leaves, and, as Ezekiel did, demonstrates their use. He says they serve as a medicine to heal the people. This is not because there will be diseases or sores in that blessed land, but so that the blessed may have continual and perpetual health. He calls these people Gentiles, not that the Gentiles are presently unclean, but because they once were, but now, being purged by Christ, live whole and sound forevermore.

### Paragraph 1365

And by these allegories he has hitherto depicted in parts those blessed dwellings [note 96.8], prepared for the faithful in that everlasting realm, under the image of a most noble city. Which after he has shown us, he seems as it were to have opened heaven itself, and set forth the eternal felicity to be seen in a manner with mortal eyes, and even to have pointed with the finger, to no other end than that we should be strong and steadfast in the faith of our Lord Jesus Christ. And should never think that one who has ever seen those blessed dwellings, to which we are called by the denying of all pleasures? What if you should despise the present pleasures and should obtain none in time to come? This thought is wicked. Faith teaches you otherwise. But what say you more? Do you desire to know and see such things as God has shown you? You have seen enough and abundantly at this present. The Lord has shown you abundantly enough of celestial life and pleasure at this present. Endeavor now only, that the devil, the world, and Antichrist trodden under, you may aspire and be lifted up into those heavenly dwellings. Moreover, beware that you be not more curious than is fitting or required, and that you seek not to know more, and more exact things of the heavenly tower and perpetual joys, than the Lord himself, who alone knows these things, has revealed to you. Let this evident demonstration of eternal life suffice us. I believe never anyone has disputed better or more rightly, more elegantly and more evidently of the blessed life than here Saint John has done. Let us therefore repose ourselves in God, let us believe his words, let his revelation suffice us, and let us desire to be joined with him in this heavenly court, in all felicity and eternal life most perfect.

### Paragraph 1366

And now John, recalling the main points of this matter, and concluding this passage concerning eternal life, finishes this everlasting felicity in seven parts. We will touch on these briefly, for many believe we have already spoken of them sufficiently and abundantly. First, there will be no curse, no condemnation, no malediction, neither war, nor famine, nor disease, nor anything as is recited by Moses among the curses in Deuteronomy 27:28. It is not that all who are subject to these things are accursed. Job and other holy men were afflicted with sickness; but commonly, the unbelievers and wicked are plagued with them. Not that they should be disciplined and profit in godliness, but that they should first be afflicted here, and by certain degrees pass on to greater torments. What then?

### Paragraph 1367

The second part adds: but the throne of God and of the Lamb will be in that city. That is, the kingdom of God will be there, and God will reign, and all blessing, no malediction, will be among the elect. Therefore, whatever joyful things the Prophets, Christ, and the Apostles have spoken of the kingdom of God, those things will be in heaven, and the blessed will experience them. And again, the Father and the Son are joined together inseparably in the unity of essence, which nevertheless, in the distinction of persons, are exceedingly well not divided, but discerned. These mysteries of the blessed Trinity are known to the faithful.

### Paragraph 1368

Here follows the third point. Some may wonder what the blessed will do in the world everlasting. Therefore John and his associates say that they shall serve God, whom I say is the Lamb, and they shall worship him in honoring, praising, and magnifying him forever. Therefore they shall wholly devote themselves to godly worship. This thing will indeed be to him great pleasure, as Augustine also shows in another place.

### Paragraph 1369

Fourthly, they shall see the face of God.[note 96.13] Augustine discusses much of what is said about God to Paulina, and warns the faithful not to imagine carnal things here. Moses in Exodus 33 and John the Apostle in John 14 have considered it the greatest blessing to see God as he is, and as is commonly said, face to face. And there is undoubtedly in this sight and enjoyment, great felicity, everlasting and most complete joy; however, in this present world, as the Lord said to Moses, it happens to no one. The holy fathers have indeed seen God, but by a form, and so far as he has deigned to reveal and show himself to be seen. As Tertullian shows in the book against Praxeas, but to see the full glory of God with inestimable joy will be granted to us first when, being delivered from misery and purged from corruption, we shall also be clarified in body; then, as John also said in 1 John 3, we shall see him as he is. Job, the most righteous man, speaking of this vision of God,[note 96.14] said: when they shall have encircled or clothed (that is, the Father, Son, and Holy Spirit) this (namely, my body), I shall behold God outside of my flesh; whom I shall see for myself, and my eyes shall look upon, and no other. This is my only desire. Paul, the Doctor of the Gentiles, also spoke of this saying, and said: Now we see in a mirror,[note 96.15] even in a dark speaking; but then shall we see face to face, &c. And Augustine has also discussed this vision in his book *The City of God*, toward the end.

### Paragraph 1370

Finally, they shall have the name of God in their foreheads, either because they will be the children of God, as we have heard in the Epistle to the Philadelphians, in chapter 3 of this book. Indeed, in the heavenly realm it will be plainly known to all who are the children of God. In this world, they are commonly regarded as the children of the devil, yet in truth are the children of God. But this will clearly appear in another world, to the great glory of the elect. And indeed, the brightness of God will shine from the foreheads, or countenances, of the elect, as in times past the brightness of the Lord shone from the face of Moses and Christ. Or because all saints will know one another, since the virtue of God rests in their countenances—a sense I perceive has pleased Primasius. Or because they will be priests before the Lord forever, as the prophets have taught of the chosen. In ancient times, the high priest bore the very name of God on his forehead on a plate of gold, bound to his head with a lace. Undoubtedly, in the heavenly realm, the glory of the children of God will be wonderfully great, especially among those who have confessed the name of Christ on earth; for these, the heavenly Father will glorify.

### Paragraph 1371

In the sixteenth member, what has been said once or twice before is repeated: that the elect in heaven are illuminated with divine glory, of which enough has been spoken previously.

### Paragraph 1372

In the last and seventh part, encompassing as it were all things of life and felicity, and uttering with one word: they shall reign,[note 96.18] he says, forevermore. The Lord Jesus grant to us his faithful, that such things as we have now heard plentifully from his mouth, we may shortly experience in our souls and bodies, and may cry with joy, to God the Father most merciful, and to Jesus Christ the Redeemer most mighty and benign, and to the Holy Spirit the most sweet Comforter be praise and glory forevermore. Amen.

## Sermon 97: ¶ The conclusion of this work, wherein is established the autoritie of the same, and the some collected briefly.

### Paragraph 1373

.

### Paragraph 1374

And he said to me, “These sayings are faithful and true.” And the Lord God of the holy prophets sent his Angel to show to his servants the things which must be shortly fulfilled. Behold, I am coming shortly. Happy is he who keeps the prophecy of this book. I am John who saw these things, and I heard them. And when I had heard and seen, I fell down to worship before the feet of the Angel who showed me these things. And he said to me, “Do not do this.” For I am your fellow servant, and of your brethren the prophets, and of those who keep the sayings of this book. Worship God.

### Paragraph 1375

The sixth and last part of this work contains the conclusion,[note 97.1] which affirms the things we have heard to be divine, certain, and undoubted: for he collects the chiefest things, and moves all men to faith, and study of godliness, that in steadfast hope we should look for the judge of all to come shortly, and to judge the living and the dead. And in goodly order this last book of the Canonical scripture finishes the godly narration and doctrine, with the judgment and end of all things.[note 97.2] For the holy Scripture begins at the first origin of all things, and continues a narration until the end of all things, containing in itself the universality of things, and all such things as are requisite to be known about matters needful and profitable. And all those things our good Lord has given us to be known in the holy scripture, that is to say, in the Canonical books. For they are false harlots who say that all things which appertain to the true and full godliness and salvation of the faithful are not set forth in holy writings, and therefore have need of traditions. They indeed have need of those traditions, which will utter their crafty wares: we need none, who esteem all their wares not worth a farthing to be bought of any man. For Isaiah has sufficiently dissuaded us from their deceivable and crafty jugglings, in chapter 55. And this conclusion contains about 16 articles, which we shall discuss in order.

### Paragraph 1376

Immediately after the beginning is set a grave asseveration, that the things which he has said or written hitherto are true, sure, certain, and undoubted, he has in a manner the same sentence also in the nineteenth chapter of this book. And he calls faithful sayings, which are stable, ratified, and undoubted. And the sentence is referred in the things which he has spoken of the blessed life to the world to come, lest we should be left in any doubtfulness. Again it is referred to the whole narration of this book. And this sentence seems to be a clause of assertion, confirming the certainty of the matter proposed, as are those also in the prophets: “For the Lord has spoken,” and again, “Thus says the Lord of Hosts,” and that same most used in the gospel, “Verily, verily, I say unto you.” And in the apostolic Epistles, “God is my witness that I lie not.” And the goodness of God does succor our infirmity, whereby many times we, doubting of the verity of God’s words, do waver, and confirms our hope with these as it were anchors. Wherefore these must be diligently beaten in and urged in the ecclesiastical doctrine. Arethas expounding this place: as the wonted manner of this holy Evangelist is always, so is it here also. For like as in his gospel, in token of loyalty, he says, “And we know that his testimony is true,” so in this place also, setting to his seal, he says, “These sayings are faithful and true.” Hitherto he. Therefore it shall be an unworthy thing to doubt, be it ever so little, of the things that are written in this book, and in other books of the canonical Scripture.

### Paragraph 1377

Secondly, he repeats who is the author of this work, and that all these things are revealed to him. And truly, there is no other author but the Lord God himself, and that the God of the holy Prophets. This has great power, for he shows not only that he is one and the same God of both Testaments, who by his Spirit has inspired the prophets and Apostles; but also bids us secretly to esteem the truth and certainty of this book of the prophetic matters. For if he could in old time tell his people beforehand of things to come, and utter all things by the prophets, what marvel is it, if he now also performs the same by Saint John? And if all those things came to pass which the prophets did prophesy to come, neither did any word, no nor one jot fall to the ground which was not fulfilled, then there is no man also that will doubt of the truth of this book, if at least he consider that the same God which in times past was with the prophets, is now also with blessed John. The prophets said how the land of Canaan should be delivered into the possession of the children of Israel: it was delivered. The selfsame prophesied that the people of Israel should for their sins be cast out again of the same land into Babylon: they were cast out. After they prophesied again that they should be delivered, should repair the City, to which Christ would come, which should redeem mankind, and call into fellowship of life and bless all nations: they were delivered, they repaired their City: Christ came, and redeemed mankind, and the gospel was preached throughout the whole world. What thing then remains, but that the church should be troubled, Antichrist should come, and reign, and that the true Christians and he should wage battle together, and the Judge come at the last unto judgment, and reward every one according to his doings? And this place proves the divinity of Christ infallibly. For what can be spoken more plainly than was said? The Lord God of holy Prophets sent forth his Angel. So in the first chapter is said: The revelation of Jesus Christ, which God gave him. And a little after he says: I Jesus sent my Angel, that he might testify unto you, &c. Herein therefore is shown the unity of the divine substance, and the destruction of persons.

### Paragraph 1378

And the manner of the revelation is shown, or repeated and collected rather: he sent his angel. Christ therefore, by his angel, shows all things to John. For no one has seen God at any time; neither shall the Lord come down again from heaven before the judgment. Wherefore this whole vision was exhibited and declared by the angel, which was the messenger of Christ the Lord. Wherefore all things are properly referred to Christ, who sent the angel. But to whom did he show or reveal these things? To his servant. For the scorners of God laugh at these things, and take them for fables. But God loves his worshippers, and warns them of all things in due season.

### Paragraph 1379

Now he gathers together some of the things he has treated hitherto. These are chiefly contained in two points. For he shows what must be done, briefly. For this book contains the destinies of the church from the time of the Apostles to the end of the world. Therefore, he did not prophesy far off, but of things that began in the very time of Saint John. And if these things must be done, who shall resist? Not that I will establish the necessity of the Stoics, but that I acknowledge the mighty working of God, according to his providence and righteousness. After this, he adds another point: “Behold, I come quickly.” For this book comprehends many things which concern the judgment itself, and the last judgment, to which I will come so swiftly and unexpectedly that the wicked and ungodly shall expect nothing less. For the Lord says in the Gospel: “It shall be as in the days of Noah and Lot.” And in the hour that you think not, the Son of Man will come. Moreover, as brightness comes forth from the East and shines in the West, so shall be the coming of the Son of Man. And therefore, the Lord says now also at this present: “Behold, I come quickly.” For suddenly, while he seems in the meantime to do another thing, at unawares he brings in the Lord speaking, and that a wonderful matter, as this particle indicates. For Saint Paul has also written, while they shall say peace and security, sudden destruction shall come upon them.

### Paragraph 1380

4. But what benefit should God’s servants expect from this book? He encompasses much in a few words, and says: “Happy is he who keeps the words of the prophecy of this book.” Felicity and blessedness are the fruits to be gained from this book. In this present world, being united with Christ, we shall walk in the way of righteousness, and avoid the schemes of Antichrist; and shall not feel the torments that arise in the conscience, from the corruption of depraved religion. And when we depart from here, we shall go directly to those blessed places. This is the ultimate blessedness and felicity. And let us note that it is not enough either to have seen, or heard, or read this book; it must needs be kept. For we must beware that it does not enter one ear and exit the other, so that we forget the things that are told us, but rather that we shape our entire life according to the doctrine of this book. And he attributed to it the title of prophecy. All Scripture is called prophecy, meaning divine. But considering how this book, for the most part, reveals things to come for the church, it is rightly called a prophecy.

### Paragraph 1381

5. He repeats again and emphasizes, [?beats upon], both his name and also that he is a witness who saw and heard, who may surely be credited. And thus he will gain authority for this book. For it must needs be held in great estimation that which was conceived and written by the Apostle and Evangelist John. Many consider it a fault that John so diligently expresses his name. But it is marvelous that they do not observe the same elsewhere, and in the writings of others not without praise. Does not the same John repeat and impress the name of a disciple in the story of the Gospel? Who should reprehend this? I see not therefore what he has offended herein: but rather he foresaw in the spirit that many would speak against this book, not without great cause, and with much fruit. And also of extreme necessity he importuned his name. And the Apostle Paul also to the Galatians: “Behold, I Paul say unto you,” says he, “if you be circumcised, Christ shall profit you nothing.” The same also, to move affection, impresses his name to Philemon. Arethas therefore very aptly expounds this place. And this, says he, is a certain propriety of speech in this Apostolic soul. For even as he did in the Gospel also, where he says: “And he that saw hath borne witness, and his testimony is true,” the same doth he in this place also, testifying that he was both a hearer and beholder of these things which are prophesied. For hereby he wins credit to the things which had been seen. Thus much he. Others have thought that not without cause, John has in this book repeated his name oftener than in his story, for that men will more hardly believe a prophecy speaking of things yet to come, than a story which tells of matters past.

### Paragraph 1382

6. In the sixth place, he adds what happened to him again with the angel, revealing to him these great mysteries. A similar story would apply to all the world, as we have in the nineteenth chapter, where we also explained it; and anyone who wishes may see. And yet, the commentators ask how it happens that John again does what he did before and was prohibited by the angel? Thomas Aquinas suggests that John, being beside himself because of the splendor of the visions, did this as one who is astonished. The gloss says that before, the angel forbade him from worshipping with *latria*, but here he prohibits him from worshipping with *dulia*. But to me it appears—preferring always the better judgment of others—that in John is shown to all the godly how great is the frailty of humankind to fall, unless it is restrained and drawn back by the mighty hand of God. The angel had shown John expressly before that he should not do what he now did, and now repeats it again. For having, as it were, forgotten those things because of the splendor of the angel, he would surely have given him some worship. For so we permit to ourselves more than is proper, especially toward noble persons, whom for their excellent gifts of God we deem worthy, whom we may also without offense to God even worship. That opinion deceives most people in our time, those who against the comeliness of sincere religion worship and honor saints. But the angel of the Lord here neither invents nor brings forth any new doctrine, but that old in form, as they term it, so that we should understand that the will of God is always one and perpetual, which will does not want the most excellent creatures to be worshipped, but one God alone to be honored. He repeats therefore the same reasons, which he also objected before. Therefore be they always of force, with all, and at all times. John meanwhile seems to want to commend to us the splendor of this vision or revelation; and that the angel constantly admonished him of his duty, and us all by him, that the thing which is proper to God, we should not transfer to any creatures, and it deserves exceeding great praise here that John here conceals nothing, but by express words commits to writing his fall and rebuking of the angel most evidently. For by his fall he would admonish that the godly should not fall in like cases, but give all glory to God. Here seems also to be observed a marvelous affection in the manner of speaking. For the angel cries out to John, being ready to fall down now, yea prostrate already, and now about to worship: See thou do it not, that thou verily intend to do. Here is expressed the carefulness of mind and haste with which he goes about to prevent the enterprise of John. And thus diligent are the Holy Spirits in heaven in letting all things, that by any means do turn us from God, to the worshipping of creatures, much less would they themselves be worshipped, or have the things attributed to them, which the Papists at this day attribute by force of arms. The Lord of clemency and mercy convert them to a right mind, that they may attribute all glory to God.

## Sermon 98: ¶ S. John is commanded not to seal this book, but to publish it, having respect to no man.

### Paragraph 1383

.

### Paragraph 1384

And he said to me: Seal not the sayings of the prophecy of this book. For the time is at hand. He who does evil, let him do evil still; and he who is defiled, let him be defiled still; and he who is righteous, let him be more righteous; and he who is holy, let him be more holy.

### Paragraph 1385

7 The seventh point in this conclusion forbids John from sealing the book that was written. But the angel says, “Seal it not.” Certainly, letters and books are often sealed either for the sake of credibility and confirmation, or so that they are not openly read by all, but only by those to whom they are assigned. An angel says to Daniel in chapter 12, “You, Daniel, seal the words and the book until the last time.” He is commanded to shut the book, that is to say, to make an end, and to expect no further revelation. Finally, he is commanded to shut it for the ungodly, to whom this book will assuredly seem dark and closed. For it follows: “Many shall err, and knowledge shall be manifold.” For those not ruled by the certain and sure word of God have nothing certainly tried and known, but wander through manifold, sundry, and uncertain opinions, judgments, and traditions of men. For Daniel says that knowledge shall be variable: that is to say, there shall be innumerable opinions and sects concerning the religion and service of God; whereas there is but one only true opinion, doctrine, faith, or religion—the same that Daniel set forth in his book, which book also he sealed, that is to say, confirmed it as if with godly seals, as authentic and authorized, and worthy to be credited. However, at this present, John is not commanded in the same sense and meaning not to seal his book, which we know to be altogether authentic. But what the angel means by such a thing is to conceal or cover not, and hide not this book: so that it may be written, that it might be a public doctrine throughout the world, whereby all men might be instructed in the things revealed from heaven, that they be not, through the schemes and tyranny of Antichrist, withdrawn from the kingdom of Christ to the kingdom of Antichrist. For God would that all these things should be most common and manifestly known to all men. This sense has Arethas also explained, saying: “Seal them not,” that is, keep them not sealed to yourself, but publish them to all. The reason is added: for the time is at hand, wherein verily these things which I have said shall come to pass. Wherefore the faithful had need of warning, confirming, and comfort. Considering therefore that this book is set forth that it might admonish, strengthen, and comfort the faithful, the same ought not to be shut but wide open. For this is the good will of God, that this his word should be preached in his church to the profit of all faithful. Let them look, therefore, what they do, who would have this book not only shut up, but clean taken away; neither think it can be understood as obscure and full of dark speakings. But to God be praise and thanks giving, who has vouchsafed to provide for us faithfully and in time by this most profitable and most necessary book.

### Paragraph 1386

The eighth point of this conclusion seems to address a certain objection. For someone might say: You want this book to be open and accessible to all people of all ranks, genders, and ages. But there will be some who will utterly despise it. Therefore, preaching it will be in vain, and we will urge these writings in vain, especially those who will mock it and interpret it as they please. But he seems to anticipate this, saying: Surely, there will be countless unrighteous people who will proceed unchecked in their wickedness and will increasingly exceed and surpass themselves. Yet, there will also be righteous people who, persevering in all righteousness, will grow in holy virtues and will excel in this as well. Therefore, do not hesitate to proclaim to all what I have commanded you in this book, without being concerned about the outcome. Leave that to me; fulfill the office of preaching. I will ensure that your faithful preaching is not in vain. And leave them alone, if you see that some will be entirely depraved and perish in their depravity, saying they despise all your faithful labor. For you have done your duty and are blameless. They perish through their own fault. Therefore, I do not want you or anyone else to be overly concerned when you see many rejecting the purity of God’s word and preferring to wallow in filth. And we read elsewhere,[?], that the Gospel is preached to many for their condemnation, and the fragrance of the gospel is sweet to some for salvation, but to others an intolerable sting leading to destruction. A similar passage is found in Ezekiel 2, where we read that the Lord said to the prophet: “Son of man, I am sending you to the people of Israel, a rebellious people who have rebelled against me, they and their fathers, until this day. They are a people of stubborn heart and a perverse mind. I am sending you to them, and you shall tell them: Thus says the Lord God, whether they hear or not, for it is a rebellious house, so that they may know that there has been a prophet among them. And you, son of man, do not fear them or be afraid of their words, for they are contentious and sharp like thorns, and you live among scorpions. But because of this, you must not be afraid of them; you must speak my words to them, whether they hear or not.”

### Paragraph 1387

However, we must take heed here that we do not understand that God commands the ungodly to proceed to be more ungodly, where the angel says, “Let him who is unjust be unjust still,” and so forth. For it seems to be a saying in a manner like this, as it is in the gospel: “Whatever you are doing, do it more speedily.” For he commands him to do that which he knew he would do. In the same way here also, that he knew the wicked would do, he says they shall do; neither does he will that their doings should trouble John and the faithful preacher, saying that there shall also be many good people who shall also apply themselves to righteousness. We are wont also to say with a much like phrase: if it will not be otherwise, we must be content. Not that we bid him who perishes to perish; but that so we reproach his madness to him, and signify that he perishes through his own fault, willingly and knowingly. Arethas says: “It is no exhortation,” but rather a rebuking of every one to whom he applies himself. And Thomas of Aquin says: “The sense is, let him who hurts, hurt still.” That is, he will hurt by doing other evils, so that the angel be understood to have said these things in prophesying, not in wishing, and so forth. And so the meaning would be: the wicked confirming the prophecy shall continue to be wicked, the godly again shall grow in the holy study of righteousness. This sense truly seems most plain of all. Neither do they differ much from these that are read in the twelfth chapter of Daniel by these words: “Go, Daniel,” says the angel, “and search not curiously over the instant of the last time; for the sayings are closed and sealed until the last time. Many shall be purified and made white and cast anew. But the wicked shall do wickedly, and all ungodly shall not understand, but the learned shall teach.” From these things, nothing at all swerves from the words of the Apostle, speaking and prophesying of the later times: “All who will live godly in Christ Jesus shall suffer persecution for righteousness. Nevertheless, evil men and deceivers grow worse and worse, while they both lead others into error and err themselves.” Therefore, saying that the later age of this world shall be such, let us, who are called to this function, proceed constantly to advance, set forth, and beat in, the very word of God and revelation of Jesus Christ unto all men, regarding nothing what the world and worldly men speak against it.

### Paragraph 1388

And he skillfully contrasts two kinds of men against two others, the unrighteous against the righteous, and the defiled against the holy. “Let him who does evil,” he says, “continue to do evil; let him who is unrighteous continue to be unrighteous; let him who harms by persecuting the godly continue to harm, and indeed, to intensify it.” Against this, he sets: “Let him who is righteous be even more righteous, let him advance further, and grow increasingly in all godliness, and go beyond himself in righteousness, both in faith and works.” For by the righteousness of faith we are justified; by the righteousness of works, we are declared to be righteous. And those who are righteous not only harm no one but also benefit and do good to all. Conversely, the unrighteous, who lack true faith, lack light, and therefore walk in darkness, doing the works of darkness, persecuting both the righteous, and righteousness, and troubling all. And that there should be such men in the later days, the Lord also prophesied in the twenty-fourth chapter of the Gospel according to Matthew.

### Paragraph 1389

Another kind of men are unclean, polluted, filthy, and vile, and so forth. He says, “Let the filthy be silent in their filth.” And the interpreters of the Greek language advise that [illegible] refers to that filth which we gather with the ends of our nails. He signified unclean persons in body and soul: idolaters, fornicators, gluttons, and the like. Against them he has placed the holy, pure, and clean—that is to say, purified by faith and diligently applying themselves to sanctification. Therefore, just as the filthy wallow themselves more and more in the mire, and array and defile themselves vilely, so the godly apply themselves more and more daily to cleanliness and holiness of life. The Lord Jesus justifies and sanctifies us forevermore.

## Sermon 99: ¶ He gathers such things as he hath taught of the last Judgment, and of the rewards of the godly, and of the torments of the wicked.

### Paragraph 1390

.

### Paragraph 1391

And behold, I am coming quickly, and my recompense is with me, to give to every man according to his works. I am Alpha and Omega, the beginning and the end, the first and the last. Blessed are those who keep his commandments, that their power may be in the tree of life, and that they may enter through the gates into the city. For outside will be dogs and sorcerers, and adulterers, and murderers, and idolaters, and whoever loves and practices falsehood [? leasings].

### Paragraph 1392

The ninth point of this conclusion concerns the Lord’s coming for judgment, and the rewards prepared for the righteous, and the punishments appointed for the unrepentant and wicked. He gathers here that he treated this matter more thoroughly in chapters 19 and 20, and in other passages of this book. And this point, above all others, he emphasizes and urges most earnestly, for it is of great importance if we both understand it correctly and consider it often in our minds. For we will sin less freely, but will watch more diligently.

### Paragraph 1393

And in this conclusion of St. John, the speakers are often changed. For now John speaks himself, and immediately he introduces the Lord speaking. Indeed, at this point he makes the Lord Christ himself speak, saying, “Behold, I come quickly.” For the word pronounced from Christ’s mouth has more authority and carries more weight with all than what the Apostle speaks. And in saying that he will come shortly, he intends to stir up all people to watch, repent, and pray. For in the Gospel he said, “Watch, for you know neither the day nor the hour.” Your Lord will come at an hour when you think least. He fears the slothful and unclean people who comfort themselves, believing that the Lord will not come at all, and if he does come, that it will be long first, and perhaps never. Against them, pleading, he says how he will come quickly. Against the same, Malachi reasoned in chapters 3 and 4. And St. Peter reasoned in chapters 2 and 3. Moreover, in affirming that he will come shortly, he comforts the godly who are tempted and tossed diversely in this world. For the godly sometimes cry out that the Lord delays his coming for a long time, that he is too benign to his enemies. Wherefore he says that he will come soon enough, that is to say, in due time, that he may both deliver his servants and destroy and root out his enemies and opponents.

### Paragraph 1394

For it follows, what manner of person he is, how, and to what end he will come: he will come gloriously with great majesty and power to deliver and save the faithful, and condemn the ungodly, for he says, “My reward is with me.” These words seem to be taken from the fortieth chapter of Isaiah. And they signify that God is abundantly furnished with all the instruments with which a rewarder and revenger must be furnished. Therefore he says, “The reward which I shall give to every one, after his doings, I have presently with me, and that ready and plentiful.” For our king and judge lacks not power and treasure; as often the kings of this world either cannot pay their soldiers’ wages as they ought, or do not have it ready, and delay the payment a long time. But this our Captain says, “My reward is with me.” And immediately explaining himself, he says that he will reward every one according as his doing shall be. For so the Apostle also in 2 Corinthians 5 says, how we must all appear before the judgment seat of Christ, that every one may receive such things as are done by the body, according as he has done, whether it be good or evil. For in the sixteenth chapter of the gospel of Saint Matthew, the Lord said likewise that there would come a time when the Son of Man should come in the glory of his Father, with his angels, and then he shall render to every one after his doings. The same is taught by the Apostle in Romans 2.

### Paragraph 1395

And to ensure that no one should doubt that our judge can accomplish in deed what he said he would do, namely to render to every man according to his works, he adds and says, “I am Alpha and Omega; the beginning and the end,” and so forth. By these words, he signifies that he is truly God, eternal, and almighty. The sentence is taken from the 43rd and 45th chapters of Isaiah and has been explained previously. These things teach us that Jesus Christ is truly God, and therefore the rewarder of all, most bountiful and most righteous.

### Paragraph 1396

Consequently again, expressly, more plainly, and by a division, John declares what and to whom the Lord will give. First, indeed, he treats of the reward prepared for the good, then of the punishment appointed for the evil by the just judgment of God. And the reward is paid, or given rather, as Paul says, to those who keep his commandments, namely Christ’s. For not those who read or hear the commandments of God, or boast and preach them, are blessed, but those who keep and perform them in deed. For so our Lord and Savior Christ taught us in the Gospel according to Matthew, chapter 7, and Luke, chapter 11. And his commandments are those that are expounded in the ten precepts, or in the Gospel restrained to the love of God and our neighbor, or which are named by John the Apostle faith and love. It behooves us therefore to be religious, in case we look to receive a reward from God. And what is the reward that is given by the judge to the godly worshippers of God? That is taken in three ways. First, they are called happy and blessed. Secondly, they shall have power over the tree of life, that is to say, the fruits of the tree of life shall be in their power: that is to wit, they shall live an eternal life with Christ, as before is declared. For he alludes to the former things. Lastly, they shall enter in, he says, by the gates into the city (to wit, before also described) into the everlasting country.

### Paragraph 1397

After this, he addresses the punishments appointed for the wicked, and in one word encompasses them all, when he says, “without.” For by this single word, he excludes the wicked from the heavenly realm and confines them to hell, where torments are unspeakable, endless, and innumerable. And Saint John here follows the Lord in the Gospel, saying: “I say to you, that many shall come from the East and from the West, and shall rest with Abraham, Isaac, and Jacob in the kingdom of heaven; but the children of the kingdom shall be cast out into the outer darkness, where there will be weeping and gnashing of teeth.” Similarly, in the parable of the ten virgins, the gate is said to be shut, and the foolish virgins shut out from celestial joys. He also commands the unprofitable servant to be cast out into the outer darkness. Likewise, in Luke 13, the Lord says how the unbelievers shall be expelled.

### Paragraph 1398

And who, I pray, shall be cast out in that final judgment? Dogs, and those listed in the register of the condemned. The term “dogs” is not always used in Scripture to denote the wicked, but generally it does. Abner, the prince of King Saul’s wars, says to Ishbosheth, “Am I the head of a dog, to defend the house of Saul against Judah?” This signifies that he had incurred the displeasure of the tribe of Judah, because he had retained ten tribes still in their duty and under the dominion of the house of King Saul. Elsewhere, as in Matthew 15, Gentiles, or pagans, or those estranged from the people of God, seem to be called dogs. Some today call the Turks “Turkish dogs,” meaning Turkish unbelievers. Furthermore, the prophet Isaiah calls the false prophets dogs: shameless, ravenous, insatiable, unable to bark and defend the Lord’s flock, or else unwilling and sleepy. With the same meaning, the Apostle says to the Philippians: “Beware of dogs, beware of evil workers, etc.” Moreover, in Scripture, angry, fierce, cruel men, those who despise godly things, who bark at the truth, slanderers, persecutors, and blasphemers are called dogs. For in Psalm 22, David, a figure of Christ the Lord, cries: “Dogs have surrounded me; the counsel of the wicked has encompassed me.” Whom he now calls dogs, he immediately names wicked. And when Shimei cursed David, Abishai the son of Zeruiah says: “Why does this dog, who is to die, curse my lord the king?” And the Lord in the Gospel forbids giving what is holy to dogs, or pearls to swine. Finally, these filthy men, the unclean, are called dogs; they are without repentance, wallowing themselves in the filth of sin and wickedness.

### Paragraph 1399

For St. Peter calls such people dogs returning to their vomit. And the Lord forbids that anyone bring the price of a prostitute or a dog into the Temple. Even then, the Jewish priests rejected the price of blood offered by Judas. Therefore, under the name of dogs, we understand pagans or unbelievers, false prophets or deceivers, cruel men, blasphemers, persecutors of the truth, cursed speakers, contemners of the truth, unclean and filthy, and so on.

### Paragraph 1400

And as for the matters that have been explained before, namely in chapter 9 and at the beginning and end of chapter 21. He adds a lie here, he who loves and promotes. For many do not openly create them, but they love, favor, and advance them. Many both love and create them. They love chiefly a falsehood, which supports false doctrine and delights therein. But concerning this, Primasius, Bishop of Utica, speaks purposefully: he says that to all these things, one must give not diligence of explanation, but carefulness of avoiding the evils. The Lord Jesus save us from all evil. Amen.

## Sermon 100: ¶ Christ is showed again to be Author of this book, how great he is here. Here is also declared the desire of the church, wiſſhyng for the coming of Christ, and the liberal promise of the Lord.

### Paragraph 1401

.

### Paragraph 1402

I Jesus sent my angel to testify to you these things in the congregations. I am the root and generation of David, and the bright morning star. And the Spirit and the bride said, “Come.” And let him who hears say also, “Come.” And let him who is thirsty come. And let whoever wills take of the water of life, freely.

### Paragraph 1403

The tenth point of this conclusion again demonstrates that the author of this work is Jesus Christ, as Saint John states, to ensure that what is spoken carries greater authority and is more readily accepted by the audience. Therefore, nothing should be attributed to Saint John except for the writing of the work, that is to say, that he first saw these things and committed them to writing. The manner of the revelation is also repeated. Christ himself did not come down to earth, but sent forth his angel, who, from Christ and in Christ’s name, opened and showed these things to Saint John. The purpose of the angel’s sending or revelation is also specified: that he should testify these things in congregations and to all people throughout the world, until the end of the world. From these few words, we learn that credit must be given to this book, as it was proposed by the very Son of God through his angel and apostle, and indeed proposed to all within the church. Again, this demonstrates that Jesus Christ is very God, the Lord of angels, as Saint Paul affirms in the first chapter of Hebrews. [note 100.2] This has also been spoken of before. These clear testimonies of Scripture ought to move the faithful more than the foolishness of Servetus the Spaniard and the Arian and Jewish practices of Servetans. Let us observe moreover that Christ sent his angel, not to judge or to teach, but to testify. Testimonies lawfully taken or committed to writing and sealed are not lawful to contradict; for they are entirely considered authentic. But this entire book was written by Saint John and is a witness or the testimony of God’s angel. Therefore, it is unlawful to doubt anything within it, and we ought to hold the same opinion of all other books of the Old and New Testaments. For the prophets and apostles are called the witnesses of God, and the gospel and doctrine of the prophets and apostles is the witness or testimony. He is mad who thinks that the canonical Scripture is of itself authentic unless it is first made authentic by the approval of the church and councils. Moreover, we understand that the doctrine of this entire book does not belong only to the seven churches of Asia, but to all who are scattered throughout the world; and therefore it pertains especially and singularly to us, who live today in Zurich or Switzerland, England, France, or Germany. Arethas, Bishop of Caesarea, says that he should testify, that is to say, that he should protest not privately or obscurely, but in the hearing of all churches dispersed throughout the world, so that no one pretending willful ignorance remains uncorrected.

### Paragraph 1404

And immediately the Lord himself also shows and declares, who, and how great he is, and what we faithful have laid up in store in him. And he uses again parables and allusions for the more clarity. First, he calls himself the root and offspring of David, that is to say, a true and natural man. For we have heard before that he is very and natural God. And he cuts off from all heretics denying and impugning the true flesh of Christ, all witnesses. Strongly proving that he, after the flesh, is of our own nature. Whereof he is also called in Scripture the fruit of the womb of David, and he that is risen from his loins. Moreover, it is said to the Davidic virgin and mother of God: “You shall conceive in your womb and bring forth a son.” Therefore, he calls himself both the root and offspring of David. And the phrase of speech is to be marked. For the like is read in Ezekiel 16. “Your root and your offspring is of the land of Canaan”: that is to say, your birth is of the Canaanites, or your origin is of a polluted people. Yet, it seems here nevertheless also another certain thing is signified. For the root bears a tree and nourishes it, or quickens it. The root is not borne or nourished by the tree. And Christ the Lord is the foundation and preservation of the house of David and the church of the faithful. That David is preserved, that the offspring of David is not rooted out, which often times has deserved to be, it is done in respect or merit of Christ the Lord. Christ has saved them, the same saves also, so many as are saved, as he is the head of all the promises made unto David, adding virtue and even perfection, as in whom is perfect salvation and all fullness, as the clear testimonies of the Prophet Isaiah bear witness in chapters 7 and 37, and elsewhere, also in Hosea 3 and Ezekiel 34 and 37. And not an unlike place is in 1 Kings, chapter 15. John also in 1 [?John 1:1] of this book names Christ the root of David, and so forth.

### Paragraph 1405

Again the Lord calls himself a Star, and that not obscure, but shining and bright, even the morning Star. When he called himself a Star, he had respect to the most ancient oracle of Balaam, that most wise prophet in the East. He prophesied that a star should arise out of Israel, that is to say a celestial star, and even the very Son of God should be born of a woman. And that the same star did arise, the magicians testify in the second chapter of Saint Matthew. And it is called bright, because Christ is the light, illuminating all men that come into the world. Of this matter the same Saint John has treated much in the first, eighth, and ninth chapters of his Evangelical story. The same our Lord is also the morning star, so called by Saint Peter, 2 Peter 1, and by this our Saint John in the second chapter of the Apocalypse. For just as Lucifer rising, draws the day star after him, so Christ shining in the hearts of the faithful, does lighten them more and more in this present world also, and in the life to come does clothes them whole with the celestial light. Thomas of Aquin expounding this place says, “The morning star is to wit the messenger of the day, that is the everlasting felicity through his resurrection.” And these things we have heard hitherto from the mouth of Christ, concerning Christ, who and how great he is, and what treasures we have laid up in store in him. He is very God and man, was incarnate for us, that he might be our root, virtue, life, light, and salvation. Therefore have we reposed in him all fullness of salvation. And so we see again, that this book is written with the apostolic spirit, which spirit verily so often as occasion serves, reasons excellently of Christ, and preaches his salvation, and commends the faith in him, unto all the faithful. The same spirit therefore has inspired both book of the Gospel and the Apocalypse of Saint John, and caused them to be written by the same author.

### Paragraph 1406

11. In the eleventh place, the church is introduced, desiring the coming of Christ in judgment. For our Lord Jesus Christ is so good, so benign, and wholesome, that this entire book promises he will come and deliver the church of saints afflicted in this world. Now is recited the desire of this same church, wishing and calling the Lord, saying, “Come.” For soon we shall hear the Lord promising and saying, “I will come quickly.” And the church again responds, “Amen. Even so, come, Lord Jesus.” To show that the spirit within our body earnestly cries to the Lord for our deliverance and glorification, the Apostle mentions this much in Romans 8. By “the spirit,” however, we may understand every spiritual person also. And therefore Arethas, who calls them “spirit,” says that they are those accounted worthy of the spiritual marriage: and the bride, the church itself. Thus he says. We have spoken of the bride many times in this work, so we need not be tedious in repeating the same. Nevertheless, with a wonderful desire, all the godly long for the Lord to come in judgment: that day is terrible and abhorred for the wicked, but most joyful and wished for by the godly. For the godly perceive that they will once be delivered from all evils and abundantly rewarded with all good things, that the glory and truth of God will be advanced and established, that all ungodliness will be abolished, and the wicked will be tormented by the just judgment of God. Whereupon Peter in Acts 3 calls this day the restoration and fulfillment of all things that God has at any time spoken by the mouth of his prophets. In that same day, therefore, all the promises of God, even the greatest matters, will be fulfilled completely. Therefore, the Lord says in the gospel, “Lift up your heads, for your redemption draws near.” Those who mourn and are desperate cast down their heads; but the Lord bids us lift up our heads, to be cheerful and of good hope. For we shall certainly be delivered and glorified, who have been in the world a laughingstock and had in derision by all men. Therefore, the passages must be interpreted figuratively, which portray the exceedingly great lamentation and howling that will be in that day. For the wicked, in anguish, pain, and utter desperation, will cry out and tear themselves. The godly will rejoice in him whom they see coming, showing the wounds by which they are redeemed. Just as the desire of the saints was greatest when the first coming of our savior approached, as is evident in Simeon alone (Luke 2), so at the second coming of Christ in judgment, all the saints with unceasing voices will cry and continually do cry, “Come, Lord Jesus, come and deliver us, come and maintain your glory and church, almost brought to naught; come, our redeemer and savior, so wished and looked for, dispatch us from evils, grant us the promised good things, and so forth.”

### Paragraph 1407

Therefore, the following things [note 100.7] may be referred either to the church or to John, that either the church or John should say: “Let him who hears say, ‘Come.’” Arethas explains this passage briefly and well: by these words, he suggests those who are not yet admitted to the flock, yet are ready to hear godly matters and diligently seek to know the Lord. So much he says. And without doubt, the desire of the godly is so great that they long for all creation to pray the Lord to come in judgment, as we often see in the Psalms, where the godly exhort the sun and moon and all creation to praise and speak well of the Lord.

### Paragraph 1408

12. The twelfth place of the conclusion contains a most large promise and comfort of Christ. For he promises again plainly. As though he should say: I know what things the faithful will suffer under Antichrist, what also and how great craft the same will practice. All things he will sell for money, Heaven and Earth, and those things also which are not in his power. And he will deceive many, and will spoil many. And all the godly he will vex and oppress with grievous persecution. Therefore if I tarry long, and come not speedily, in as much as the desires of the saints covet the same, you that love and believe in me, flee Antichrist, give not yourselves to be spoiled by him. Look for me, have recourse unto me. He that is thirsty, that is, he that desires an heavenly gift, or he that is in anguish or tormented with cares and sundry evils, let him come to me; to me I say, let him come. I shall fill him with good things, deliver from evil, and will comfort him, and strengthen him with my spirit, in all manner dangers, that he may patiently bear and overcome all evils. And he seems to have borrowed these wholesome and most full of consolation words from the doctrine of Isaiah, which is in the fifty-fifth chapter, and in the seventh chapter of John. Hereof are spoken certain things about the beginning of the twenty-first chapter. Where we read the Lord to have said: And to him that is thirsty will I give of the well of the water of life freely.

### Paragraph 1409

[note 100.9]But where he says, “and he who wills,” he does not mean, as many mistakenly believe, that it rests with our will to be saved. For we know that the Apostle has said, “It is not by will, nor by effort, but by God’s mercy.” The Lord saves us of his own good will, yet he does not save the unwilling, but the willing. But he gives us the ability to will, according to the Apostle’s saying, “It is God who works in us both to will and to accomplish.” Primasius says: by no good gifts does one go before receiving the water of life freely. For what do you have, says the Apostle, that you did not receive? Therefore, we have freely received from God the will to come also; to whom we gave nothing first, that we should be, much less that we should be made righteous from sinners. Thus he says. Nevertheless, it might seem to be a manner of speaking common among the Germans: which is, “I make it free for all to come,” I clearly exclude no one, I bid all come; so, “he who wills,” that is to say, “come all, and receive water,” and so forth. To the Lord be glory.

## Sermon 101: ¶ Punishment is decreed to the corrupters of this book. The lord says, that he will certainly come to Judgment. The church wiſſheth for his coming.

### Paragraph 1410

.

### Paragraph 1411

I testify to every person who hears the words of the prophecy of this book: if anyone adds to these things, God will add to him the plagues that are written in this book. And if anyone diminishes from the words of the book of this prophecy, God will take away his part from the book of life, and from the holy city, and from the things which are written in this book. He who testifies these things says, “Surely, it is.” I am coming quickly. Amen. Even so, come, Lord Jesus. The grace of our Lord Jesus Christ be with you all. Amen.

### Paragraph 1412

13. In the concluding section of this decree, a penalty is established for those who despise this book, but especially for counterfeiters or forgers, who (as Dr. Bibliander has said aptly and piously) dare attempt to corrupt or falsify this godly instrument and holy charter of Christ’s empire and bishopric, by adding anything or taking away, or altering the true meaning and sense thereof. This passage is taken outside of the common practice of men. For princes are accustomed at the end of their writings to secure them against detractors by means of threats and warnings. Antichrist, the mimic of our Lord Christ, at the end of his bull adds: “If anyone rashly presumes to go against this our commandment, or presumptuously to infringe the same, let him know that he shall incur the indignation of Almighty God and the blessed Apostles Peter and Paul, and our high displeasure.” And likewise in safeguarding treasures and public things, where danger is feared, they affix writings and seals with wax. For this reason, indeed, knowing that there would be some who would seek to oppress and abolish this book, the Lord sends it well armed to all future generations. We read in ancient authors that certain heretics in the early church greatly asserted themselves in corrupting the scriptures; and that some of them rejected entire books of the holy scripture. And Tertullian attributes the same to Marcion, who also corrupted holy books. Nevertheless, through God’s goodness, it came to pass that we have received the holy books whole and uncorrupted. This is plainly shown by St. Jerome in his commentaries on Isaiah, book 3. And Erasmus of Rotterdam in the Apology of the New Testament, and also in his Apology against James Latomus, and so forth.

### Paragraph 1413

However, the Lord at this present time does no new thing, while he commands that nothing should be added or taken away. For once or twice he commanded by Moses: “You shall add nothing to my word, nor take anything from it.” And Solomon in the thirtieth chapter of the Proverbs commands the same. But many marvel and find fault that he has threatened so many plagues to the corrupters. Why then do they not blame and reprove Paul as well, who in one word has encompassed as many plagues and displeasures as John has recited here, where he said to the Galatians: “Although I, or an angel from heaven, should preach to you a gospel other than the one we have preached to you, let him be an outcast or accursed.” And he repeats these same words. Wherefore, if they grant that Paul has offended so little in this regard, that he deserves praise also, let them cease blaming those things which are here most purposefully placed by our Lord Jesus Christ himself through John in their place and time due. Anathema (a word Paul used) is he who is cut off from the fellowship of all good men, devoted to extreme punishment, and even subject to all the evils of both this present life and the life to come. This may be gathered from Deuteronomy and other holy books.

### Paragraph 1414

And he adds and takes away, not one who uses other diverse and plain words in explaining a sentence of the Apocalypse; but he who puts in anything contrary to the true sense, or varies from those things which are here expressed by the Lord: or he who takes away anything, obscures or perverts what the Lord himself has expressly signified. Thomas of Aquin says he adds, which is to put in a lie; he diminishes, which takes away anything from that which is written therein, or also denies gain, and says the same. Thus much he says. Therefore, this addition and subtraction consists not in words only, but rather in sense. For neither are the prophets, in explaining the law at large, thought to have added anything to God’s word; neither are the Apostles, preaching the liberty of the gospel, said to have taken anything away from the law.

### Paragraph 1415

And to testify—that is, to affirm something definite by a declaration or to urge something earnestly, as if to bind the listener—so that he should certainly know that these things spoken have authority over him, and that God will punish unless he obeys.

### Paragraph 1416

Concerning the pains or plagues which he threatens at this time, this is spoken of in chapters 15, 16, 17, and 18, and so forth. Likewise, it is declared beforehand that what might be spoken here concerning the book of life and the holy city. Moreover, he comprehends here also all good things in like manner, which are promised in this book to the godly and obedient servants of God, of all which things the contemner, falsifier, and corrupter of this book shall be deprived. With how great evils and dangers than do they entangle themselves, who would have this book utterly suppressed, and prevented from being expounded openly, and coming into the hands of all men? Again, it is most certain that they shall obtain of God all manner of blessing, so many as have a good opinion, and think devoutly of this book, and will set forth and commend unto all men the things that are written in the same to the glory of God, and salvation of the faithful.

### Paragraph 1417

14. In the fourteenth place, the authority of this book is sealed, and as it were signed with a subscription. For it follows, he says that bears witness of these things, or he that testifies these things. In general, all the commentators suppose these to be Christ’s words, as though he himself, for a confirmation, had added to it and said: “I Jesus have revealed all these things, and especially those concerning the threatenings against the corrupter, as a true witness, and the same to be undoubted.” For Aquinas says: “Here Christ is brought in, affirming the aforementioned threatening, and approving all things that are written in this book.” But I, for my part, reserving the judgments of others safe, suppose this to be the subscription of Saint John, the writer of this book. For notaries, secretaries, or chancellors of princes are wont at the end of the kings or emperors’ letters or writings to subscribe their name. And indeed Saint John in the History of the Gospel has observed the same manner. For in chapter 19, he says: “And he that saw bore witness, and his testimony is true.” And at the end of the History he subscribes and signs underneath with these words: “This is that disciple which bears witness of these things, and wrote them, and we know that his testimony is true.” For all the church knew and confessed this. In like manner, he seems at this present to have subscribed these things also in his own name and to have said, “He that testifies these things, says.”

### Paragraph 1418

15. Next, he brings in again the Lord Jesus himself speaking and promising that he will certainly come for judgment, truly to redeem and glorify the righteous, and to punish the wicked. Therefore, with a strong affirmation, he says, “Even so, surely, and without doubt I am coming.” Although I seem to delay, and to some I do not appear to be coming at all. Nevertheless, I will certainly come in due time, as has been said and declared before. And this same thing is repeated, in similar manner with the same words, more often as a matter most worthy to be noted and known.

### Paragraph 1419

He immediately adds the faith and wish, and great desire of Saint John, and of the faithful church, or of any godly person submitting himself to the promise, and saying, “Amen, even so.” That is to say, we acknowledge it to be most certain and undoubted that you promise that you will come. Therefore, we look for you, the Judge of the living and the dead; yea, and we pray with our innermost being, “Come, Lord Jesus.” For elsewhere also we pray daily, “Thy kingdom come.” And all the godly, with sighs unspeakable, wish for the coming of the Judge, for glory, whereof the Apostle treats in the eighth chapter to the Romans, and we have touched the same matter before.

### Paragraph 1420

Finally, in his conclusion, he wishes, in the apostolic manner, the grace of our Lord Jesus Christ to all who hear and read this book. Paul writes in his second letter to the Thessalonians, chapter three: “Thus I write in every letter: May the grace of our Lord Jesus Christ be with you all.” Amen. Thus he aligns himself with the apostolic spirit, which is present everywhere. Grace encompasses the entire matter of redemption and the gifts of Christ. Therefore, he wishes to us all the blessing that we have in Christ Jesus our Lord. Of this, the vessel of election has abundantly been endowed. Paul, in chapter one to the Ephesians, prays that the Lord Jesus, who has revealed these holy mysteries to us, may write them in our minds, and deliver us from Antichrist and from all evils; and keep us in the true faith and in his grace. To him be honor and glory, praise and thanks giving, together with the Father and the Holy Spirit, forever: Amen. Come, Lord Jesus, our redeemer and only savior, and glorify those who eagerly await your coming, that we may glorify you forever. Amen.
